\documentclass[
  justified,
  openright,
]{tufte-book}

\usepackage{placeins}
\usepackage{xstring}

\newcommand{\smartnoteformat}[1]{%
  % Check if it starts with "Accessibility", "Accessible", or the bold versions
  \IfBeginWith{#1}{Accessibility}{\color{blue}}{%
  \IfBeginWith{#1}{Accessible}{\color{blue}}{%
  \IfBeginWith{#1}{\textbf{Accessibility}}{\color{blue}}{%
  \IfBeginWith{#1}{\textbf{Accessible}}{\color{blue}}{\color{black}}}}}%
  #1%
}

% --------------------------------------------------
% Engine sanity (XeLaTeX or LuaLaTeX required)
% --------------------------------------------------
\usepackage{iftex}
\ifPDFTeX
  \PackageError{template}{This template requires XeLaTeX or LuaLaTeX}{}
\fi

% --------------------------------------------------
% Core math
% --------------------------------------------------
\usepackage{amsmath,amssymb,amsfonts}

% --------------------------------------------------
% Tufte: subsection is forbidden --- remap safely
% --------------------------------------------------
\makeatletter
\let\subsection\paragraph
\let\paragraph\paragraph
\makeatother

% --------------------------------------------------
% Fix Pandoc tabularray counter=none bug
% --------------------------------------------------
\makeatletter
\ExplSyntaxOn
\cs_if_exist:NF \c@none { \newcounter{none} }
\ExplSyntaxOff
\makeatother

\usepackage{fontspec}

\IfFontExistsTF{TeX Gyre Pagella}
  {\setmainfont{TeX Gyre Pagella}}
  {\setmainfont{Latin Modern Roman}}

\IfFontExistsTF{TeX Gyre Heros}
  {\setsansfont[Scale=0.92]{TeX Gyre Heros}}
  {\setsansfont{Latin Modern Sans}}

\IfFontExistsTF{TeX Gyre Cursor}
  {\setmonofont[Scale=0.82]{TeX Gyre Cursor}}
  {\setmonofont{Latin Modern Mono}}
% --------------------------------------------------
% Graphics and tables
% --------------------------------------------------
\usepackage{graphicx}
\usepackage{booktabs}
\usepackage{longtable}
\usepackage{array}
\usepackage{tabularray}
\UseTblrLibrary{booktabs}

\setkeys{Gin}{width=\maxwidth,height=\maxheight,keepaspectratio}

% --------------------------------------------------
% Colors
% --------------------------------------------------
\usepackage[table,xcdraw,svgnames]{xcolor}

\definecolor{DarkGray}{RGB}{55,60,70}
\definecolor{CalloutBlue}{HTML}{DCEEFF}
\definecolor{CalloutBorder}{HTML}{5B8BD9}
\definecolor{BrandBlue}{RGB}{18,65,124}
\definecolor{BrandAccent}{RGB}{0,153,204}
\definecolor{LightGray}{RGB}{245,246,248}
\definecolor{MidGray}{RGB}{175,180,190}
\definecolor{WarningAmber}{RGB}{230,100,0}
\definecolor{WarningBg}{RGB}{255,248,235}
\definecolor{CautionYellow}{RGB}{180,140,0}
\definecolor{CautionBg}{RGB}{255,252,230}
\definecolor{TipGreen}{RGB}{36,110,44}
\definecolor{TipBg}{RGB}{237,248,239}
\definecolor{NoteBg}{RGB}{235,243,255}
\definecolor{NoteBlue}{RGB}{18,90,160}
\definecolor{DangerRed}{RGB}{185,28,28}
\definecolor{DangerBg}{RGB}{255,237,237}
\definecolor{TableHeader}{RGB}{18,65,124}
\definecolor{TableRowAlt}{RGB}{235,241,252}
\definecolor{MarginRule}{RGB}{210,215,225}
\definecolor{SpineAccent}{RGB}{0,153,204}
% --------------------------------------------------
% --------------------------------------------------
% Syntax-highlighting colours and token commands
%
% Nine semantic commands cover all token categories.
% Identity tokens (NormalTok, BuiltInTok, etc.) were
% stripped from the document body; no macro needed.
% --------------------------------------------------
\definecolor{CodeBg}{RGB}{250,250,248}        % warm near-white background
\definecolor{CodeRule}{RGB}{18,65,124}         % brand blue left-rule accent
\definecolor{CodeComment}{RGB}{140,148,163}    % muted blue-gray
\definecolor{CodeString}{RGB}{0,120,100}       % deep teal
\definecolor{CodeNumber}{RGB}{110,60,160}      % soft violet
\definecolor{CodeKeyword}{RGB}{18,65,124}      % brand blue
\definecolor{CodeType}{RGB}{55,90,140}         % mid blue-gray
\definecolor{CodeAccent}{RGB}{0,140,185}       % brand accent (imports, preprocessor)
\definecolor{CodeError}{RGB}{180,28,28}        % Tufte red

% Semantic token commands (only these appear in the document body)
\newcommand{\KeywordTok}[1]{\textcolor{CodeKeyword}{\textbf{#1}}}  % keywords, control flow
\newcommand{\StringTok}[1]{\textcolor{CodeString}{#1}}              % strings, chars, verbatim
\newcommand{\CommentTok}[1]{\textcolor{CodeComment}{\textit{#1}}}  % all comment variants
\newcommand{\CodeNumTok}[1]{\textcolor{CodeNumber}{#1}}             % integers, floats, constants
\newcommand{\CodeFuncTok}[1]{\textcolor{CodeKeyword}{#1}}           % function/method names
\newcommand{\CodeTypeTok}[1]{\textcolor{CodeType}{#1}}              % types, attributes
\newcommand{\CodeAccentTok}[1]{\textcolor{CodeAccent}{#1}}          % imports, preprocessor
\newcommand{\OperatorTok}[1]{\textcolor{DarkGray}{#1}}              % operators, punctuation
\newcommand{\ErrorTok}[1]{\textcolor{CodeError}{\underline{#1}}}    % errors
\newcommand{\WarningTok}[1]{\textcolor{WarningAmber}{\textbf{#1}}}  % warnings

% --------------------------------------------------
% Code-block environments  (Tufte-native style)
%
% Visual design:
%   • 3 pt brand-blue left rule (Tufte accent, not a box)
%   • Warm near-white background (CodeBg)
%   • Modest left/right/top/bottom padding via mdframed
%   • Monospace font at \footnotesize (readable in Tufte text width)
%   • No heavy borders; top/bottom spacing handled by mdframed skip
% --------------------------------------------------
\usepackage{fancyvrb}
\usepackage[framemethod=default]{mdframed}

\mdfdefinestyle{TufteCode}{%
  backgroundcolor=CodeBg,
  linecolor=CodeRule,
  linewidth=2.5pt,
  leftline=true,
  rightline=false,
  topline=false,
  bottomline=false,
  innerleftmargin=10pt,
  innerrightmargin=8pt,
  innertopmargin=6pt,
  innerbottommargin=6pt,
  skipabove=\medskipamount,
  skipbelow=\medskipamount,
  splittopskip=\topskip,
  splitbottomskip=0pt,
  nobreak=false,
}

\newenvironment{Shaded}{%
  \begin{mdframed}[style=TufteCode]%
}{%
  \end{mdframed}%
}

% Highlighting: verbatim env inside Shaded; footnotesize suits Tufte column
\DefineVerbatimEnvironment{Highlighting}{Verbatim}{%
  commandchars=\\\{\},%
  fontsize=\footnotesize,%
  baselinestretch=1.15,%
  xleftmargin=0pt%
}

% --------------------------------------------------
% Tufte: subsection is forbidden --- remap safely
% --------------------------------------------------
\makeatletter
\let\subsection\paragraph
\let\paragraph\paragraph
\makeatother

% --------------------------------------------------
% Utilities
% --------------------------------------------------
\usepackage{etoolbox}
\usepackage{xparse}
\usepackage{ifthen}
\usepackage{textcase}
\usepackage{titlesec}
% --------------------------------------------------
% Fix Pandoc tabularray counter=none bug
% --------------------------------------------------
\makeatletter
\ExplSyntaxOn
\cs_if_exist:NF \c@none { \newcounter{none} }
\ExplSyntaxOff
\makeatother
% --------------------------------------------------
% Paragraph & section spacing
% --------------------------------------------------
\setlength{\parindent}{1em}
\setlength{\parskip}{0pt}

\titlespacing*{\section}{0pt}{2.0ex plus .2ex}{1.0ex}
\titlespacing*{\subsection}{0pt}{1.5ex plus .2ex}{0.8ex}

% --------------------------------------------------
% Margin notes (tufte)
% --------------------------------------------------
\setlength{\marginparwidth}{2.2in}

\newcommand{\mn}[1]{%
  \marginnote{\raggedright\footnotesize\color{DarkGray}#1}
}

\NewDocumentCommand{\margincallout}{m}{%
  \marginnote{%
    \setlength{\fboxsep}{6pt}%
    \colorbox{CalloutBlue}{%
      \begin{minipage}{\marginparwidth}
      \footnotesize #1
      \end{minipage}
    }
  }
}

\newcommand{\amn}[1]{\mn{\color{blue}#1}}

\usepackage{chngcntr}
%\counterwithin{sidenote}{chapter}

% --------------------------------------------------
% Soul / tabularray / marking fixes
% --------------------------------------------------
\usepackage{soul}
\AfterEndEnvironment{longtblr}{\soulsuspend}

\makeatletter
\providecommand{\soulsuspend}{}
\providecommand{\soulsuspend*}{}
\providecommand{\soulsuspend@}{}
\makeatother

\DeclareRobustCommand{\checkbox}{\(\square)}

\soulregister\checkbox{0}
\soulregister\textbf{1}
\soulregister\emph{1}
\soulregister\href{2}
\soulregister\label{0}
\soulregister\ref{0}
\soulregister\pageref{0}
\soulregister\footnote{1}
\soulregister\cite{0}

% --------------------------------------------------
% Headers / marks (no case mangling)
% --------------------------------------------------
\makeatletter
\renewcommand{\chaptermark}[1]{\markboth{#1}{}}
\renewcommand{\sectionmark}[1]{\markright{#1}}

\pretocmd\markboth{\let\allcapsspacing\relax\let\smallcapsspacing\relax}{}{}
\pretocmd\markright{\let\allcapsspacing\relax\let\smallcapsspacing\relax}{}{}
\makeatother

% --------------------------------------------------
% Blank page helper
% --------------------------------------------------
\newcommand{\blankpage}{%
  \clearpage
  \thispagestyle{empty}
  \null\vfill
  \begin{center}
    \color{MidGray}\sffamily\small\itshape
    (This page intentionally left blank.)
  \end{center}
  \vfill
  \clearpage
}

% --------------------------------------------------
% Fancyhdr
% --------------------------------------------------
\usepackage{fancyhdr}
\pagestyle{fancy}
\fancyhf{}
\fancyhead[LE]{\small\scshape\leftmark}
\fancyhead[RO]{\small\scshape\rightmark}
\fancyfoot[C]{\thepage}
\AtBeginDocument{\pagestyle{fancy}}

% --------------------------------------------------
% Hyperlinks
% --------------------------------------------------
\usepackage{hyperref}
\hypersetup{
  colorlinks=true,
  linkcolor=blue,
  urlcolor=blue,
  citecolor=blue,
  hypertexnames=false
}

% --------------------------------------------------
% Pandoc helpers
% --------------------------------------------------
\providecommand{\tightlist}{%
  \setlength{\itemsep}{0pt}\setlength{\parskip}{0pt}}

% Show only parts and chapters in TOC; number only parts
\setcounter{tocdepth}{1}
\setcounter{secnumdepth}{-1}

% --------------------------------------------------
% Metadata
% --------------------------------------------------



% --------------------------------------------------
% ============================================================
%  DOCUMENT
% ============================================================
\begin{document}
% -- Front matter -------------------------------------------
\frontmatter
\pagestyle{fancy}

% 1) Blank half title page
\thispagestyle{empty}
\null
\clearpage

% 2) Frontispiece (placeholder)
\thispagestyle{empty}
\begin{center}
  \vspace*{0.35\textheight}
  {\sffamily\large Frontispiece}
\end{center}
\clearpage

% 3) Full title page
\begin{titlepage}
    \vspace*{1.6in}
    {\sffamily\fontsize{9}{11}\selectfont\MakeUppercase{Training Manual}}
    \vspace{0.3in}
    {\color{BrandAccent}\rule{4.5in}{2pt}}
    \vspace{0.25in}
    % ADD YOUR TITLE TEXT HERE:
    {\sffamily\bfseries\fontsize{32}{38}\selectfont Coding With \\ScreenReaders}
    \vspace{0.2in}
    {\color{BrandAccent}\rule{4.5in}{0.6pt}}
    \vspace{0.5in}
    \vfill
    {\sffamily\fontsize{7.5}{9}\selectfont\color{MidGray}%
      CONFIDENTIAL --- For Internal Use Only%
    }
  \clearpage
  \nopagecolor
\end{titlepage}

% 4) Copyright page (placeholder)
\thispagestyle{empty}
\vspace*{0.25\textheight}
{\small
\noindent Copyright \textcopyright\ \the\year\ Michael Ryan Hunsaker, M.Ed., Ph.D. Coding With ScreenReaders is licensed under a Creative Commons Attribution-NonCommerical-NoDerivatives 4.0 International: https://creativecommons.org/licenses/by-nc-nd/4.0/  All rights reserved.

\noindent NO AI TRAINING: Without in any way limiting the author's rights under copyright, any use of this publication to “train” generative artificial intelligence (AI) technologies to generate text is expressly prohibited.
}

\clearpage

% 5) Contents
{
\hypersetup{linkcolor=BrandBlue}
\tableofcontents
\clearpage
}

% 6) Dedication (placeholder)
\chapter*{Dedication}
\thispagestyle{empty}
\begin{center}
  \vspace*{0.3\textheight}
  {\itshape To the countless individuals who have been told they lack the capability to code, yet remain determined to learn and create through accessible coding practices, this book is dedicated.}\\[1em]
\end{center}
\clearpage

% 7) Epigraph (placeholder)
\chapter*{Epigraph}
\thispagestyle{empty}
\begin{flushright}
  \vspace*{0.25\textheight}
  {\itshape The power of the Web is in its universality. Access by everyone regardless of disability is an essential aspect}\\[1em]
  --- Tim Berners-Lee, Inventor of the World Wide Web
\end{flushright}
\clearpage

% 8) Introduction (placeholder)
\newpage
\chapter*{Introduction}
\addcontentsline{toc}{chapter}{Introduction}
\section{Curriculum Package --- Python, Java, and
JavaScript}\label{curriculum-package-python-java-and-javascript}

This package contains three complete introductory programming curricula,
each adapted for visually impaired students who use screen readers.

\begin{center}\rule{0.5\linewidth}{0.5pt}\end{center}

\section{Package Contents}\label{package-contents}

\begin{verbatim}
curriculum/
|--- README.md                          (this file)
|--- python/
|   |--- unit00-setup.md
|   |--- unit01-foundations.md
|   |--- unit02-numeric-computations.md
|   |--- unit03-control-flow.md
|   |--- unit04-06-data-functions.md
|   |--- unit07-09-oop-files.md
|   +--- unit10-11-projects-capstone.md
|--- java/
|   |--- README.md
|   |--- unit00-01-setup-foundations.md
|   |--- unit02-04-io-control-arrays.md
|   |--- unit05-08-strings-methods-oop-files.md
|   +--- unit09-11-collections-projects.md
+--- javascript/
    |--- README.md
    |--- unit00-01-setup-foundations.md
    |--- unit02-04-io-control-data.md
    |--- unit05-08-strings-functions-oop-files.md
    +--- unit09-11-modules-collections-projects.md
\end{verbatim}

\begin{center}\rule{0.5\linewidth}{0.5pt}\end{center}

\section{Curriculum Comparison}\label{curriculum-comparison}

{\def\LTcaptype{none} % do not increment counter
\begin{longtable}[]{@{}
  >{\raggedright\arraybackslash}p{(\linewidth - 6\tabcolsep) * \real{0.1604}}
  >{\raggedright\arraybackslash}p{(\linewidth - 6\tabcolsep) * \real{0.3019}}
  >{\raggedright\arraybackslash}p{(\linewidth - 6\tabcolsep) * \real{0.2830}}
  >{\raggedright\arraybackslash}p{(\linewidth - 6\tabcolsep) * \real{0.2547}}@{}}
\toprule\noalign{}
\begin{minipage}[b]{\linewidth}\raggedright
Feature
\end{minipage} & \begin{minipage}[b]{\linewidth}\raggedright
Python
\end{minipage} & \begin{minipage}[b]{\linewidth}\raggedright
Java
\end{minipage} & \begin{minipage}[b]{\linewidth}\raggedright
JavaScript
\end{minipage} \\
\midrule\noalign{}
\endhead
\bottomrule\noalign{}
\endlastfoot
\textbf{Difficulty} & Beginner & Intermediate & Beginner-Intermediate \\
\textbf{Typing} & Dynamic & Static & Dynamic \\
\textbf{Syntax} & Minimal & Verbose & Moderate \\
\textbf{Best For} & First language, data science, AI & Enterprise,
Android, CS theory & Web development, full-stack \\
\textbf{Terminal Tool} & \texttt{python3\ file.py} & \texttt{javac} +
\texttt{java} & \texttt{node\ file.js} \\
\textbf{Testing} & \texttt{unittest} & JUnit 5 & Jest \\
\textbf{JSON} & \texttt{json} module & Gson library & Built-in
\texttt{JSON} object \\
\textbf{Input} & \texttt{input()} & \texttt{Scanner} & \texttt{readline}
module \\
\end{longtable}
}

\begin{center}\rule{0.5\linewidth}{0.5pt}\end{center}

\section{Accessibility Notes (All
Languages)}\label{accessibility-notes-all-languages}

All three curricula share these accessibility principles:

\begin{enumerate}
\def\labelenumi{\arabic{enumi}.}
\tightlist
\item
  \textbf{Screen Reader Optimized Mode in VS Code:} Press
  \texttt{Ctrl\ +\ Shift\ +\ P} -\textgreater{} ``Toggle Screen Reader
  Optimized''.
\item
  \textbf{Indentation Announcement:} Configure your screen reader to
  announce indentation changes --- critical for Python's block
  structure; helpful for Java and JavaScript too.
\item
  \textbf{Error Navigation:} Use \texttt{F8} / \texttt{Shift+F8} to jump
  between errors without scanning visually.
\item
  \textbf{All Output is Text-Based:} No graphical output required for
  any core exercises.
\item
  \textbf{Terminal First:} All programs run from the integrated
  terminal, which is accessible with screen readers.
\end{enumerate}

\begin{center}\rule{0.5\linewidth}{0.5pt}\end{center}

\section{Recommended Learning Paths}\label{recommended-learning-paths}

\subsection{Complete Beginner}\label{complete-beginner}

Start with \textbf{Python} -\textgreater{} then \textbf{JavaScript}
(many concepts transfer directly).

\subsection{Web Developer Track}\label{web-developer-track}

Start with \textbf{JavaScript} -\textgreater{} deepen with Node.js and
frameworks.

\subsection{Computer Science Theory
Track}\label{computer-science-theory-track}

Start with \textbf{Python} -\textgreater{} move to \textbf{Java} for
stronger OOP and typing foundations.

\subsection{All-Rounder}\label{all-rounder}

\textbf{Python} (learn concepts) -\textgreater{} \textbf{Java} (learn
strict typing and OOP) -\textgreater{} \textbf{JavaScript} (apply to the
web).

\begin{center}\rule{0.5\linewidth}{0.5pt}\end{center}

\section{Credits}\label{credits}

Curriculum designed for accessibility-first computer science
education.\\
Adapted from:

\begin{itemize}
\tightlist
\item
  Matthes, E. (2023). \emph{Python Crash Course} (3rd ed.). No Starch
  Press.
\item
  Zelle, J.M. (2017). \emph{Python Programming: An Introduction to
  Computer Science} (3rd ed.). Franklin, Beedle \& Associates.
\item
  Eck, D.J. (2022). \emph{Introduction to Programming Using Java} (9th
  ed.). {[}Free at math.hws.edu/javanotes{]}
\item
  MDN Web Docs. {[}developer.mozilla.org{]}
\end{itemize}

Original Python curriculum by Michael Ryan Hunsaker, M.Ed., Ph.D.~---
Davis School District.

\clearpage

% -- Main content -------------------------------------------
\mainmatter
\pagenumbering{arabic} % Switch to numbers (1, 2, 3...) for the main content
\raggedbottom
\part{Python Curriculum}

\section{Teacher's Guide: Introduction to Computer Science with
Python}\label{teachers-guide-introduction-to-computer-science-with-python}

\subsection{For Students with Visual Impairments --- Complete Instructor
Edition}\label{for-students-with-visual-impairments-complete-instructor-edition}

\begin{center}\rule{0.5\linewidth}{0.5pt}\end{center}

\subsection{Table of Contents}\label{table-of-contents}

\begin{enumerate}
\def\labelenumi{\arabic{enumi}.}
\tightlist
\item
  \hyperref[philosophy-and-foundational-principles]{Philosophy and
  Foundational Principles}
\item
  \hyperref[understanding-your-students]{Understanding Your Students}
\item
  \hyperref[classroom-setup-and-technology-configuration]{Classroom
  Setup and Technology Configuration}
\item
  \hyperref[universal-spoken-code-conventions]{Universal Spoken Code
  Conventions}
\item
  \hyperref[unit-0-setup-and-environment]{Unit 0 --- Setup and
  Environment}
\item
  \hyperref[unit-1-foundations-and-getting-started]{Unit 1 ---
  Foundations and Getting Started}
\item
  \hyperref[unit-2-numeric-computations-and-io]{Unit 2 --- Numeric
  Computations and I/O}
\item
  \hyperref[unit-3-control-flow-decisions-and-loops]{Unit 3 --- Control
  Flow, Decisions, and Loops}
\item
  \hyperref[unit-4-data-structures-lists-and-dictionaries]{Unit 4 ---
  Data Structures: Lists and Dictionaries}
\item
  \hyperref[unit-5-user-input-and-while-loops]{Unit 5 --- User Input and
  While Loops}
\item
  \hyperref[unit-6-functions]{Unit 6 --- Functions}
\item
  \hyperref[unit-7-classes-and-object-oriented-programming]{Unit 7 ---
  Classes and Object-Oriented Programming}
\item
  \hyperref[unit-8-files-and-exceptions]{Unit 8 --- Files and
  Exceptions}
\item
  \hyperref[unit-9-testing-and-data-analysis]{Unit 9 --- Testing and
  Data Analysis}
\item
  \hyperref[unit-10-project-based-learning]{Unit 10 --- Project-Based
  Learning}
\item
  \hyperref[unit-11-capstone-project]{Unit 11 --- Capstone Project}
\item
  \hyperref[assessment-strategies]{Assessment Strategies}
\item
  \hyperref[supporting-diverse-needs-within-the-visually-impaired-population]{Supporting
  Diverse Needs Within the Visually Impaired Population}
\item
  \hyperref[teaching-ai-tools-responsibly]{Teaching AI Tools
  Responsibly}
\item
  \hyperref[alternative-and-supplementary-resources]{Alternative and
  Supplementary Resources}
\item
  \hyperref[quick-reference-sheets]{Quick Reference Sheets}
\end{enumerate}

\begin{center}\rule{0.5\linewidth}{0.5pt}\end{center}

\subsection{1. Philosophy and Foundational
Principles}\label{philosophy-and-foundational-principles}

\subsection{Why Python for Blind
Students?}\label{why-python-for-blind-students}

Python is widely recognized as one of the most accessible first
programming languages for visually impaired learners for several
concrete reasons:

\begin{itemize}
\tightlist
\item
  \textbf{Text-based I/O is the default.} Unlike languages that push
  beginners toward GUI frameworks or graphical plotting, Python's most
  natural outputs are plain text to the terminal --- fully accessible to
  any screen reader.
\item
  \textbf{Indentation enforces structure.} While indentation introduces
  a learning challenge (addressed in detail below), it also means that
  code read aloud by a screen reader carries structural meaning in the
  form of announced indent levels --- something curly-brace languages do
  not provide.
\marginnote{\textbf{Indentation announcement} must be configured in the screen reader before the first Python lesson---not as an afterthought. See Section~3 (VS Code Configuration) for step-by-step setup for NVDA, JAWS, VoiceOver, and Orca.}
\item
  \textbf{Error messages are descriptive.} Python's tracebacks include
  the file name, line number, error type, and a message. A screen reader
  user can navigate directly to the problem without visual scanning.
\item
  \textbf{The standard library is enormous.} Python ships with modules
  for file I/O, JSON, math, random numbers, CSV parsing, unit testing,
  and more --- reducing the complexity of dependency management in a
  learning context.
\item
  \textbf{The REPL is an accessible sandbox.} Python's interactive
  interpreter provides immediate spoken feedback for any expression,
  enabling rapid experimentation without the overhead of creating and
  running files.
\sidenote{The REPL is the single most accessible Python environment for early learners: no file navigation, no run command, immediate audio response. Start every new concept here before moving to files.}
\end{itemize}

\subsection{Core Pedagogical
Commitments}\label{core-pedagogical-commitments}

\textbf{Commit to ears-first instruction.} Every demonstration, every
example, every debugging session should be conducted as if the teacher
is also using a screen reader. Ask yourself before every lesson:  ``right
now?''
\marginnote{This self-check is the single most useful habit an instructor can build. Apply it to every demo, every error message, every output format. If the answer is unclear, redesign the demonstration.}

\textbf{Commit to explicit navigation.} Sighted students can glance at
code and immediately locate structures. Blind students must navigate
deliberately using keyboard shortcuts and screen reader commands. Teach
navigation skills as primary content, not auxiliary support.

\textbf{Commit to verbal precision.} Vague instructions like ``look at
the next line'' or ``you can see here that\ldots{}'' create invisible
barriers. Replace all visual language with navigation language: ``Move
down two lines,'' ``Jump to line 14 with Ctrl+G,'' ``Navigate to the end
of the file with Ctrl+End.''
\sidenote{Audit your own speech during the first three lessons. Most sighted instructors use spatial language constantly without noticing: ``you can see,'' ``look here,'' ``over on the right.'' Recording a lesson and reviewing it is the fastest way to identify and correct these habits.}

\textbf{Commit to error-positive culture.} For a sighted student, an
error is a red underline they can ignore or address at leisure. For a
screen reader user, an unaddressed error is noise that clutters the
output of every subsequent navigation. Treat errors as immediate
priorities and build the habit of reading and resolving them before
moving on.

\textbf{Commit to redundancy.} Provide concepts in multiple formats ---
spoken explanation, written pseudocode, physical analogy, and code.
Blind students cannot fall back on visual re-reading; redundant
representation builds retention.

\begin{center}\rule{0.5\linewidth}{0.5pt}\end{center}

\subsection{2. Understanding Your
Students}\label{understanding-your-students}

\subsection{The Heterogeneity of Visual
Impairment}\label{the-heterogeneity-of-visual-impairment}

``Visually impaired'' encompasses a wide spectrum. Your students may
include:

\begin{fullwidth}{\def\LTcaptype{none} % do not increment counter
\begin{longtable}[]{@{}
  >{\raggedright\arraybackslash}p{(\linewidth - 2\tabcolsep) * \real{0.2533}}
  >{\raggedright\arraybackslash}p{(\linewidth - 2\tabcolsep) * \real{0.7467}}@{}}
\toprule\noalign{}
\begin{minipage}[b]{\linewidth}\raggedright
Profile
\end{minipage} & \begin{minipage}[b]{\linewidth}\raggedright
Implications for Teaching
\end{minipage} \\
\midrule\noalign{}
\endhead
\bottomrule\noalign{}
\endlastfoot
\textbf{Congenitally blind} (blind from birth) & May have no visual
spatial metaphors. Analogies must be tactile, auditory, or conceptual.
May be faster with screen readers from years of experience. \\
\textbf{Late-onset blindness} (lost vision after sighted years) & May
retain strong visual mental models. Can use these as bridges but should
not be required to. May still be learning screen reader skills. \\
\textbf{Low vision} (functional but impaired vision) & May use a
combination of magnification and screen reading. May be able to see some
screen content but cannot rely on fine detail. Do not assume screen-only
or vision-only. \\
\textbf{Color blindness} & Primarily relevant for syntax highlighting.
Ensure themes are not relied upon for meaning. \\
\textbf{Progressive conditions} & Vision may change during the course.
Build habits that remain useful at any level of vision. \\
\end{longtable}
}\end{fullwidth}

Never assume a student's profile without asking. Some students are
private about their diagnoses; frame questions around what tools they
use and what works best, not what their condition is.

\subsection{Screen Reader Literacy
Varies}\label{screen-reader-literacy-varies}

Some students arrive with years of fluent screen reader experience.
Others are relatively new to their screen reader and may still be
developing efficiency. Assess screen reader fluency separately from
programming aptitude --- a student who struggles to navigate a file is
not struggling with Python; they need more time with their tool.

Suggest that students spend time with their screen reader in
non-programming contexts (browsing documentation, reading formatted
text) before or alongside the first programming units.

\subsection{Braille Display Users}\label{braille-display-users}

Some students use a refreshable braille display in addition to or
instead of audio output. Key considerations:

\begin{itemize}
\tightlist
\item
  Braille displays have limited width (typically 40 or 80 cells). Long
  lines of code may need to be panned across. Encourage keeping line
  length under 79 characters (PEP 8 standard) for this reason.
\sidenote{PEP~8'slength under 79 characters (PEP 8 standard) for this reason}
\amn{\textbf{Accessibility Tip:} Set \texttt{editor.rulers: [79]} in VS Code settings to display a vertical guide at column 79. With Screen Reader Optimized Mode on, VS Code will announce when the cursor crosses this boundary, giving auditory feedback about line length without visual scanning.}
\sidenote{PEP~8's 79-character line limit was originally a terminal width convention, but it maps almost perfectly to an 80-cell braille display. Frame this limit to students as a practical accessibility constraint, not an arbitrary style rule.}
\item
  Code rendered in braille uses Grade 1 (uncontracted) braille because
  contractions would misrepresent code characters. Confirm this setting
  is active.
\marginnote{Confirm Grade~1 (uncontracted) braille is active \emph{before} the first coding session. A student unknowingly reading contracted braille will misread operators, brackets, and keywords.}
\item
  Navigation with a braille display typically involves routing cursor
  buttons --- physical buttons above each braille cell that move the
  cursor to that position. Students using these may navigate differently
  than audio-only users.
\item
  Braille display users often prefer lower verbosity audio output since
  they read braille silently. Do not require audio announcements as the
  only feedback mechanism.
\end{itemize}

\subsection{Students Who Are Also
Deaf-Blind}\label{students-who-are-also-deaf-blind}

Students who are deaf-blind rely exclusively on braille for all
information. Audio feedback from the screen reader is not available.
Considerations:

\begin{itemize}
\tightlist
\item
  All instructional communication must be via tactile means, text-based
  communication, or tactile sign language interpreters.
\item
  Braille output of code is the primary interface. Pay special attention
  to line length and symbol representation.
\item
  Terminal output must be navigable via braille display --- avoid
  programs that produce output faster than a braille display can be
  read.
\item
   this guide.
\marginnote{The HKNC (Helen Keller National Center) and Perkins School for the Blind both publish guidance on programming education for deafblind learners. \texttt{hknc.org} and \texttt{perkins.org}.}
\end{itemize}

\begin{center}\rule{0.5\linewidth}{0.5pt}\end{center}

\subsection{3. Classroom Setup and Technology
Configuration}\label{classroom-setup-and-technology-configuration}

\subsection{Physical Environment}\label{physical-environment}

\begin{itemize}
\tightlist
\item
  Ensure the student's workstation does not have excessive background
  noise that interferes with screen reader audio. Headphones are often
  preferable in a shared classroom.
\item
  Position the workstation so the teacher can sit beside (not in front
  of) the student, allowing them to listen to the screen reader output
  together.
\item
  If a classroom assistant or orientation and mobility specialist is
  present, coordinate roles in advance. The programming teacher leads on
  programming content; the VI specialist leads on accessibility tool
  configuration.
\end{itemize}

\subsection{VS Code Configuration
Checklist}\label{vs-code-configuration-checklist}

Work through this checklist with the student on day one:

\textbf{Screen Reader}

\begin{itemize}
\tightlist
\item
  \begin{itemize}
  \tightlist
  \item[$\square$]
    Screen Reader Optimized Mode ON (Ctrl+Shift+P -\textgreater{} Toggle
    Screen Reader Optimized)
\marginnote{Without Screen Reader Optimized Mode, VS Code suppresses many accessibility announcements to reduce ``noise'' for sighted users. This mode must be ON for the editor to be functional with a screen reader.}
  \end{itemize}
\item
  \begin{itemize}
  \tightlist
  \item[$\square$]
    Accessibility help panel reviewed (Shift+Alt+F1)
  \end{itemize}
\end{itemize}

\textbf{Editor Settings (Ctrl+,)}

\begin{itemize}
\tightlist
\item
  \begin{itemize}
  \tightlist
  \item[$\square$]
    editor.lineNumbers = ``on''
  \end{itemize}
\item
  \begin{itemize}
  \tightlist
  \item[$\square$]
    editor.renderWhitespace = ``all'' or ``boundary''
  \end{itemize}
\item
  \begin{itemize}
  \tightlist
  \item[$\square$]
    editor.wordWrap = ``on'' (for long lines)
  \end{itemize}
\item
  \begin{itemize}
  \tightlist
  \item[$\square$]
    workbench.reduceMotion = ``on'' (reduces distracting animations)
  \end{itemize}
\item
  \begin{itemize}
  \tightlist
  \item[$\square$]
    editor.tabSize = 4
  \end{itemize}
\item
  \begin{itemize}
  \tightlist
  \item[$\square$]
    editor.insertSpaces = true (spaces not tabs for Python)
\sidenote{Mixing spaces and tabs is the most common source of \texttt{IndentationError} for new Python students. Setting \texttt{insertSpaces~=~true} at the editor level prevents the problem entirely and should be verified on day one.}
  \end{itemize}
\end{itemize}

\textbf{Screen Reader Indentation (in screen reader settings, not VS
Code)}

\begin{itemize}
\tightlist
\item
  \begin{itemize}
  \tightlist
  \item[$\square$]
    NVDA: Document Formatting -\textgreater{} Indentation
    -\textgreater{} ``Say all'' or ``Beeps''
  \end{itemize}
\item
  \begin{itemize}
  \tightlist
  \item[$\square$]
    JAWS: Configure speech for VS Code environment
  \end{itemize}
\item
  \begin{itemize}
  \tightlist
  \item[$\square$]
    VoiceOver: Verbosity -\textgreater{} Text -\textgreater{} announce
    indentation
  \end{itemize}
\item
  \begin{itemize}
  \tightlist
  \item[$\square$]
    Narrator: No built-in indent announcement; use Beep plugin or
    extension
  \end{itemize}
\end{itemize}

\textbf{Extensions Installed}

\begin{itemize}
\tightlist
\item
  \begin{itemize}
  \tightlist
  \item[$\square$]
    Python (Microsoft)
  \end{itemize}
\item
  \begin{itemize}
  \tightlist
  \item[$\square$]
    Pylance
  \end{itemize}
\item
  \begin{itemize}
  \tightlist
  \item[$\square$]
    Python Indent (for consistent auto-indentation)
  \end{itemize}
\item
  \begin{itemize}
  \tightlist
  \item[$\square$]
    GitLens or Git Graph (for version control visibility)
  \end{itemize}
\item
  \begin{itemize}
  \tightlist
  \item[$\square$]
    Error Lens (inline error annotations --- screen-reader-readable)
  \end{itemize}
\end{itemize}

\textbf{Keyboard Shortcuts Practiced}

\begin{itemize}
\tightlist
\item
  \begin{itemize}
  \tightlist
  \item[$\square$]
    Ctrl+G --- Go to line
  \end{itemize}
\item
  \begin{itemize}
  \tightlist
  \item[$\square$]
    F8 --- Next error
  \end{itemize}
\item
  \begin{itemize}
  \tightlist
  \item[$\square$]
    Shift+F8 --- Previous error
  \end{itemize}
\item
  \begin{itemize}
  \tightlist
  \item[$\square$]
    Ctrl+Shift+M --- Problems panel
\amn{\textbf{Accessibility Tip:} The Problems panel (\texttt{Ctrl+Shift+M}) shows every error and warning as a list you can navigate with arrow keys. Use \texttt{F8} to jump to the next error location in the editor. This is the most reliable way to find errors without visually scanning the code.}
    Ctrl+Shift+M --- Problems panel
  \end{itemize}
\item
  \begin{itemize}
  \tightlist
  \item[$\square$]
    Shift+Alt+F2 --- Accessible View
  \end{itemize}
\item
  \begin{itemize}
  \tightlist
  \item[$\square$]
    Ctrl+Shift+P --- Command Palette
  \end{itemize}
\item
  \begin{itemize}
  \tightlist
  \item[$\square$]
    Ctrl+` --- Toggle terminal
  \end{itemize}
\item
  \begin{itemize}
  \tightlist
  \item[$\square$]
    Ctrl+Shift+E --- Explorer
  \end{itemize}
\item
  \begin{itemize}
  \tightlist
  \item[$\square$]
    Ctrl+Shift+O --- Go to symbol (function/class navigation)
  \end{itemize}
\item
  \begin{itemize}
  \tightlist
  \item[$\square$]
    Alt+F12 --- Peek Definition
  \end{itemize}
\item
  \begin{itemize}
  \tightlist
  \item[$\square$]
    F12 --- Go to Definition
  \end{itemize}
\item
  \begin{itemize}
  \tightlist
  \item[$\square$]
    Shift+F12 --- Find all References
  \end{itemize}
\end{itemize}

\subsection{Terminal Configuration}\label{terminal-configuration}

\begin{itemize}
\tightlist
\item
  \textbf{Windows Terminal} (preferred over legacy cmd.exe on Windows)
  has better accessibility with NVDA and JAWS.
\item
  \textbf{NVDA + Windows Terminal:} Enable ``Use UI Automation'' in NVDA
  settings for better terminal reading.
\item
  \textbf{VoiceOver + macOS Terminal:} Works well by default. Use
  \texttt{VO+Shift+Down} to interact with terminal content.
\item
  \textbf{Orca + Linux terminal:} Orca reads terminal output well. Use
  \texttt{Orca+T} (flat review) to re-read recent output.
\item
  Configure the terminal font to a monospace font at a size comfortable
  for any residual vision.
\end{itemize}

\subsection{Python Virtual Environment Setup for
Accessibility}\label{python-virtual-environment-setup-for-accessibility}

\begin{Shaded}
\begin{Highlighting}[]
\CommentTok{\# Create environment with a consistent name the student will always type}
python3 \CodeTypeTok{{-}m} venv .venv
 
\CommentTok{\# Activation {-}{-}{-} save this as a text file the student can open and copy from}
\CommentTok{\# Windows: .venv\textbackslash{}Scripts\textbackslash{}activate}
\CommentTok{\# macOS/Linux: source .venv/bin/activate}
 
\CommentTok{\# Verify with a spoken command}
python3 \CodeTypeTok{{-}{-}version}
pip list
\end{Highlighting}
\end{Shaded}

Consider creating a startup script (\texttt{start.bat} on Windows,
\texttt{start.sh} on macOS/Linux) that activates the environment and
opens VS Code.
\marginnote{Store the activation command in a plain \texttt{.txt} file the student can open and copy from. Never require them to recall multi-step activation paths from memory during the first weeks of the course.} Students can run one command to get their environment
ready, reducing setup cognitive load.

\begin{center}\rule{0.5\linewidth}{0.5pt}\end{center}

\subsection{4. Universal Spoken Code
Conventions}\label{universal-spoken-code-conventions}

Establish this vocabulary on the first day and reinforce it
consistently. When a student is reading code aloud (for debugging, peer
review, or just narrating their work), they should use these terms
without hesitation.

\subsection{Symbol Pronunciation
Table}\label{symbol-pronunciation-table}

\begin{fullwidth}{\def\LTcaptype{none} % do not increment counter
\begin{longtable}[]{@{}
  >{\raggedright\arraybackslash}p{(\linewidth - 4\tabcolsep) * \real{0.2500}}
  >{\raggedright\arraybackslash}p{(\linewidth - 4\tabcolsep) * \real{0.3696}}
  >{\raggedright\arraybackslash}p{(\linewidth - 4\tabcolsep) * \real{0.3804}}@{}}
\toprule\noalign{}
\begin{minipage}[b]{\linewidth}\raggedright
Written Symbol
\end{minipage} & \begin{minipage}[b]{\linewidth}\raggedright
Spoken Form
\end{minipage} & \begin{minipage}[b]{\linewidth}\raggedright
Notes
\end{minipage} \\
\midrule\noalign{}
\endhead
\bottomrule\noalign{}
\endlastfoot
\texttt{=} & ``gets'' & Assignment only \\
\texttt{==} & ``equals equals'' or ``is equal to'' & Comparison \\
\texttt{!=} & ``not equal to'' or ``bang equals'' & \\
\texttt{\textgreater{}=} & ``greater than or equal to'' & \\
\texttt{\textless{}=} & ``less than or equal to'' & \\
\texttt{\textgreater{}} & ``greater than'' & \\
\texttt{\textless{}} & ``less than'' & \\
\texttt{+=} & ``plus equals'' or ``add and assign'' & \\
\texttt{-=} & ``minus equals'' & \\
\texttt{*=} & ``times equals'' & \\
\texttt{**} & ``to the power of'' or ``double star'' & \\
\texttt{//} & ``floor divide'' or ``integer divide'' & \\
\texttt{\%} & ``modulo'' or ``remainder'' & \\
\texttt{()} & ``open paren \ldots{} close paren'' & Always say both \\
\texttt{{[}{]}} & ``open bracket \ldots{} close bracket'' & \\
\texttt{\{\}} & ``open brace \ldots{} close brace'' & \\
\texttt{:} & ``colon'' & End of if/for/def/class line \\
\texttt{"text"} & ``quote text quote'' & \\
\texttt{\textquotesingle{}text\textquotesingle{}} & ``single-quote text
single-quote'' & \\
\texttt{f"..."} & ``f-string'' & Mention the f prefix \\
\texttt{\{var\}} inside f-string & ``open brace var close brace'' & \\
\texttt{\#} & ``hash'' & Followed by comment text \\
\texttt{\textbackslash{}n} & ``backslash n'' or ``newline'' & \\
\texttt{\textbackslash{}t} & ``backslash t'' or ``tab'' & \\
\texttt{\_} & ``underscore'' & In names: ``my underscore name'' \\
\texttt{\_\_init\_\_} & ``dunder init'' & Double-underscore methods \\
\texttt{\_\_str\_\_} & ``dunder str'' & \\
\texttt{.} & ``dot'' & In \texttt{math.sqrt}, say ``math dot sqrt'' \\
\texttt{-\textgreater{}} & ``arrow'' & Type hints \\
\texttt{*args} & ``star args'' & \\
\texttt{**kwargs} & ``double-star kwargs'' & \\
\texttt{None} & ``None'' (capital N) & Distinct from the word
``none'' \\
\texttt{True} / \texttt{False} & ``True'' / ``False'' (capital) & \\
\texttt{import} & ``import'' & \\
\texttt{from\ x\ import\ y} & ``from x import y'' & Read as English
sentence \\
\end{longtable}
}\end{fullwidth}

\marginnote{The \emph{spoken} distinction between \texttt{=} (``gets'') and \texttt{==} (``equals equals'') is one of the most important conventions to establish early. Many bugs in beginner Python trace directly to mixing these up. Consistent spoken vocabulary prevents this class of error.}

\subsection{Reading a Line of Code
Aloud}\label{reading-a-line-of-code-aloud}\amn{\textbf{Accessibility Tip:} Use f-strings (\texttt{f"Value: \{x\}"}) rather than string concatenation (\texttt{"Value: " + str(x)}) for output. F-strings are read more naturally by screen readers and make the intent of each print statement immediately clear when navigating code.}
\subsection{Reading a Line of Code
Aloud}\label{reading-a-line-of-code-aloud}

Teach students to read full lines this way, left to right, then explain:

\begin{Shaded}
\begin{Highlighting}[]
total\_price \OperatorTok{=} unit\_price \OperatorTok{*} quantity \OperatorTok{+} tax
\end{Highlighting}
\end{Shaded}

\emph{``total underscore price gets unit underscore price times quantity
plus tax''}

Then explain: \emph{``This line calculates the total price by
multiplying the unit price by the quantity, then adding tax, and stores
the result in total\_price.''}


\marginnote{The two-pass protocol mirrors how experienced blind developers actually work: one pass for syntax verification, a second for semantic reasoning. Building this habit early prevents the tendency to skip structural details.}

\begin{center}\rule{0.5\linewidth}{0.5pt}\end{center}

\subsection{5. Unit 0 --- Setup and
Environment}\label{unit-0-setup-and-environment}

\subsection{Lesson 0.1: Installing Python and VS
Code}\label{lesson-0.1-installing-python-and-vs-code}

\paragraph{Before This Lesson}\label{before-this-lesson}

Survey the student's current technology setup:

\begin{itemize}
\tightlist
\item
  What screen reader do they use, and what version?
\item
  Do they use a braille display?
\item
  What operating system?
\item
  Have they used a terminal before?
\item
  Do they have administrator access to their computer?
\end{itemize}

Confirm administrator access before scheduling this lesson. Installation
without admin rights will fail.
\amn{\textbf{Accessibility Tip:} Always verify Python is on PATH immediately after installation by opening a fresh terminal and running \texttt{python3 -{}-version}. A PATH failure produces a ``command not found'' error that is easy to miss if the student navigates away from the terminal before reading the output.}

\paragraph{Instructional Cues}\label{instructional-cues}

Walk through every installer step verbally as the student performs it.
Do not say ``click the next button'' --- say ``Tab until you hear `Next
button,' then press Enter.'' Treat this as teaching screen reader
navigation, not just software installation.

before anything else:
\marginnote{\textbf{PATH failures} are the most common first-day blocker. Symptoms: \texttt{python3} gives ``command not found'' or opens the Microsoft Store on Windows. Fix: rerun installer and explicitly check ``Add Python to PATH,'' or manually add the Python directory to the system PATH via Settings.} \texttt{python3\ -\/-version}. If this fails, stop
and address the PATH issue before proceeding. A student who continues
without Python on PATH will encounter mysterious failures throughout the
course.

\paragraph{Non-Obvious Accessibility
Strategies}\label{non-obvious-accessibility-strategies}

\begin{itemize}
\tightlist
\item
  \textbf{Windows Defender SmartScreen:} On Windows, running a
  downloaded \texttt{.exe} from Python.org may trigger a SmartScreen
  dialog. This dialog can be navigated with Tab and Enter but is not
  always announced immediately by NVDA. Tell the student what to expect
  before they run the installer.
\item
  \textbf{macOS Gatekeeper:} Downloaded packages may require right-click
  -\textgreater{} Open to bypass security. On a screen reader,
  right-click is VO+Shift+M (VoiceOver) or App Key (NVDA with keyboard
  access).
\item
  \textbf{Linux package managers as an alternative:} For Linux students,
  \texttt{sudo\ apt\ install\ python3\ python3-pip\ python3-venv} is
  often more reliable than the Python.org download and is fully
  scriptable.
\item
  \textbf{Portable Python for school computers:} Some schools restrict
  software installation.  reasonably accessible.
\sidenote{GitHub Codespaces provides a full VS Code environment in the browser with no local installation. It inherits all of VS Code's screen reader support and is often more reliable in locked-down school IT environments than local installs.}
\end{itemize}

\paragraph{Alternative Installation
Methods}\label{alternative-installation-methods}

\begin{fullwidth}{\def\LTcaptype{none} % do not increment counter
\begin{longtable}[]{@{}
  >{\raggedright\arraybackslash}p{(\linewidth - 4\tabcolsep) * \real{0.2115}}
  >{\raggedright\arraybackslash}p{(\linewidth - 4\tabcolsep) * \real{0.4327}}
  >{\raggedright\arraybackslash}p{(\linewidth - 4\tabcolsep) * \real{0.3558}}@{}}
\toprule\noalign{}
\begin{minipage}[b]{\linewidth}\raggedright
Method
\end{minipage} & \begin{minipage}[b]{\linewidth}\raggedright
Accessibility
\end{minipage} & \begin{minipage}[b]{\linewidth}\raggedright
Best For
\end{minipage} \\
\midrule\noalign{}
\endhead
\bottomrule\noalign{}
\endlastfoot
Python.org installer & High --- standard dialog flow & Most
installations \\
Chocolatey (Windows) & High --- command-line, screen-reader friendly &
Windows power users \\
Homebrew (macOS) & High --- command-line & macOS users comfortable with
terminal \\
apt/dnf/pacman (Linux) & High --- command-line & Linux users \\
Replit.com & Medium --- web-based, some accessibility gaps & No local
install possible \\
Thonny IDE & Medium --- simpler UI, bundled Python & Beginners who
struggle with VS Code \\
\end{longtable}
}\end{fullwidth}

\textbf{On Thonny:} Thonny is a Python IDE specifically designed for
beginners. It has a simpler interface than VS Code and is sometimes more
accessible for students who are new to both screen readers and
programming. The trade-off is fewer advanced accessibility features
long-term. 
\marginnote{Students who start on Thonny sometimes resist the VS Code migration. Frame it as a deliberate progression: Thonny is a learning environment, VS Code is a professional one. The transition is part of the curriculum.}

\paragraph{Thought Experiments}\label{thought-experiments}

\begin{enumerate}
\def\labelenumi{\arabic{enumi}.}
\item
  \emph{``Why do you think the Python installer has a checkbox called
  `Add Python to PATH'? What is a PATH, and what would happen if Python
  weren't on it? Think about it like a phonebook --- your operating
  system looks up programs by name. If Python isn't in the phonebook,
  what happens when you type \texttt{python3}?''}
\item
  \emph{``We're installing two things: Python and VS Code. What role
  does each one play? If you had only Python and no VS Code, could you
  still write programs? How? What would you lose?''}
\item
  \emph{``Virtual environments keep each project's dependencies
  separate. Imagine two school projects --- one requiring a specific
  edition of a textbook that contradicts the edition another class uses.
  What would happen if you tried to use both editions at the same time?
  How does a virtual environment solve the programming equivalent of
  this problem?''}
\item
  \emph{``Why do we verify every installation immediately after
  completing it? What could go wrong between downloading something and
  confirming it works? Name three things that could silently fail.''}
\end{enumerate}

\begin{center}\rule{0.5\linewidth}{0.5pt}\end{center}

\subsection{Lesson 0.2: Configuring Your Screen Reader for VS
Code}\label{lesson-0.2-configuring-your-screen-reader-for-vs-code}

\paragraph{Instructional Cues}\label{instructional-cues-1}

Before showing any configuration screens, do an unguided exploration:
``Open VS Code, navigate around with only your keyboard and screen
reader, and come back in ten minutes. Tell me what you found, what
worked, what was confusing.'' This builds genuine familiarity and
identifies that student's specific challenges.

Map each configuration step to a problem the student already identified.
This makes the configuration feel purposeful rather than arbitrary.

\paragraph{Screen Reader Specific
Guidance}\label{screen-reader-specific-guidance}

\textbf{NVDA users:}

\begin{itemize}
\tightlist
\item
  Enable ``Use UI Automation to access VS Code'' in NVDA preferences
  (Tools -\textgreater{} Manage Plugins, or via NVDA+N -\textgreater{}
  Preferences -\textgreater{} Settings -\textgreater{} Advanced
  -\textgreater{} Use UI Automation where available). This dramatically
  improves reading accuracy in VS Code.
\item
  Set ``Report dynamic content changes'' to ``All'' to hear error
  highlights and IntelliSense popups.
\item
  Configure Browse Mode to be OFF by default in VS Code --- programming
  requires forms mode (interactive mode) at all times.
\item
  Speech rate: many experienced NVDA users run at 70--90\% of maximum
  speed. For new students, 40--60\% allows better comprehension of code
  read aloud.
\end{itemize}

\textbf{JAWS users:}

\begin{itemize}
\tightlist
\item
  JAWS has VS Code-specific scripts maintained by the JAWS community.
  Check \texttt{FreedomScientific.com} and GitHub for the latest VS Code
  JAWS scripts.
\item
  PlaceMarkers in JAWS can mark frequently visited locations in a file
  --- teach \texttt{Ctrl+Win+K} to set and \texttt{Ctrl+Win+J} to jump.
\item
  JAWS Verbosity settings: set ``Announce line numbers'' to ON for code
  files.
\end{itemize}

\textbf{VoiceOver (macOS) users:}

\begin{itemize}
\tightlist
\item
  VoiceOver interacts with VS Code via the macOS Accessibility API.
  Ensure VS Code is not run with elevated permissions (which can break
  accessibility API access).
\item
  Use \texttt{VO+Shift+Down} to enter interactive elements (the editor,
  terminal), then navigate normally.
\item
  Enable ``Speak text under mouse'' in VoiceOver Utility -\textgreater{}
  Verbosity if the student uses residual vision with a mouse.
\item
  VoiceOver's ``item chooser'' (VO+I) can provide a quick list of
  interactive elements in the current window.
\end{itemize}

\textbf{Windows Narrator users:}

\begin{itemize}
\tightlist
\item
  Narrator is less feature-rich than NVDA or JAWS for programming. It
  works but may require more keystrokes.
\item
  Narrator scan mode must be OFF in VS Code for normal keyboard input to
  work. Toggle with Caps Lock+Space.
\item
   students using Narrator
Consider recommending the free NVDA screen reader as an upgrade for
  students using Narrator, as NVDA's VS Code support is significantly
  better.
\marginnote{NVDA is free, open-source, and maintained by NV Access. Download at \texttt{nvaccess.org}. Its VS Code support is substantially better than Narrator's for programming tasks.}
\end{itemize}

\textbf{Orca (Linux) users:}

\begin{itemize}
\tightlist
\item
  Orca works well with VS Code on Linux via the AT-SPI2 accessibility
  interface.
\item
  Use \texttt{Orca+H} for ``say all'' from cursor position.
\item
  Ensure the \texttt{at-spi2-core} package is installed:
  \texttt{sudo\ apt\ install\ at-spi2-core}.
\item
  Orca's ``flat review'' mode (\texttt{Orca+KP\_Enter} or
  \texttt{Orca+Insert+KP\_Enter}) is useful for reading lines that are
  not focus-accessible.
\end{itemize}

\paragraph{Indentation Configuration --- Step by
Step}\label{indentation-configuration-step-by-step}


\marginnote{If a student's screen reader is not announcing indentation, do not proceed. Every Python control structure depends on indentation. A student writing Python without auditory indentation feedback is working blind in both senses.}

\textbf{Test procedure:}

\begin{enumerate}
\def\labelenumi{\arabic{enumi}.}
\item
  Create a temporary file \texttt{test\_indent.py}.
\item
  Write:

\begin{Shaded}
\begin{Highlighting}[]
\KeywordTok{if} True:
    print(\StringTok{"indented"})
print(\StringTok{"not indented"})
\end{Highlighting}
\end{Shaded}
\item
  Navigate line by line with Down arrow.
\item
  The screen reader should announce the indent level on line 2 before
  reading its content.
\item
  On line 3, the indent should return to zero.
\end{enumerate}

If announcement is not happening, troubleshoot screen reader settings
before proceeding.\amn{\textbf{Accessibility Tip:} Test indentation announcement with a two-line file before writing any real Python. If the screen reader does not announce indent level changes, every control structure will be structurally silent. This is a hard blocker---do not proceed until it is resolved.}
If announcement is not happening, troubleshoot screen reader settings
before proceeding.

\paragraph{Thought Experiments}\label{thought-experiments-1}

\begin{enumerate}
\def\labelenumi{\arabic{enumi}.}
\item
  \emph{``A sighted programmer can look at a Python file and immediately
  understand its structure from indentation --- like an outline. How
  does your screen reader translate that same 2D visual structure into a
  1D audio stream? What information might be lost in that
  translation?''}
\item
  \emph{``VS Code has hundreds of keyboard shortcuts. No sighted
  programmer uses all of them --- they use the 15 or 20 that matter for
  their workflow. If you could only keep 10 shortcuts, which 10 would
  make the biggest difference for writing Python? Why?''}
\item
  \emph{``The Accessible View (Shift+Alt+F2) converts complex visual
  elements into flat text. What kind of information in VS Code might
  require this conversion? Can you think of a scenario where the
  `visual' presentation of something would be unusable without
  translation?''}
\end{enumerate}

\begin{center}\rule{0.5\linewidth}{0.5pt}\end{center}

\subsection{Lesson 0.3: Setting Up Your VS Code
Workspace}\label{lesson-0.3-setting-up-your-vs-code-workspace}

\paragraph{Alternative Workspace
Tools}\label{alternative-workspace-tools}

If VS Code proves too complex initially, consider these alternatives as
temporary or permanent solutions:

\begin{fullwidth}{\def\LTcaptype{none} % do not increment counter
\begin{longtable}[]{@{}
  >{\raggedright\arraybackslash}p{(\linewidth - 4\tabcolsep) * \real{0.2263}}
  >{\raggedright\arraybackslash}p{(\linewidth - 4\tabcolsep) * \real{0.0584}}
  >{\raggedright\arraybackslash}p{(\linewidth - 4\tabcolsep) * \real{0.7153}}@{}}
\toprule\noalign{}
\begin{minipage}[b]{\linewidth}\raggedright
Tool
\end{minipage} & \begin{minipage}[b]{\linewidth}\raggedright
Platform
\end{minipage} & \begin{minipage}[b]{\linewidth}\raggedright
Accessibility Notes
\end{minipage} \\
\midrule\noalign{}
\endhead
\bottomrule\noalign{}
\endlastfoot
\textbf{IDLE} (bundled with Python) & All & Very simple UI. Limited
features but highly keyboard accessible. Good for first programs. \\
\textbf{Thonny} & All & Beginner-focused. Has a ``simple mode.'' Some
screen reader users find it more approachable. \\
\textbf{Mu Editor} & All & Designed for education. Simpler than VS Code.
Python 3 only. \\
\textbf{Jupyter Notebook (terminal)} & All & Text-based notebook.
Accessible for structured exploration. \\
\textbf{Braille Note Apex/Touch} & Device & For students using these
notetakers, Python can be run via their built-in terminals. \\
\textbf{Replit} & Web & No install required. Accessibility varies; test
before recommending. \\
\textbf{NVDA Remote + VS Code} & Windows & Student uses their own screen
reader on their own machine while teacher controls a shared session. \\
\end{longtable}
}
\end{fullwidth}
\paragraph{Thought Experiments}\label{thought-experiments-2}

\begin{enumerate}
\def\labelenumi{\arabic{enumi}.}
\item
  \emph{``Programmers organize their code into folders by project. What
  happens if every programmer at a company kept all their files in one
  shared folder with no organization? How does this analogy apply to
  organizing your own Python projects?''}
\item
  \emph{``VS Code Extensions add features to the editor the same way
  apps add features to a phone. If you could design one VS Code
  extension specifically for blind Python programmers that doesn't exist
  yet, what would it do?''}
\end{enumerate}

\begin{center}\rule{0.5\linewidth}{0.5pt}\end{center}

\subsection{Lesson 0.4: Python Environments and Terminal
Usage}\label{lesson-0.4-python-environments-and-terminal-usage}

\paragraph{Terminal Proficiency as a
Prerequisite}\label{terminal-proficiency-as-a-prerequisite}

Before writing any Python, ensure the student can fluently:

\begin{itemize}
\tightlist
\item
  Navigate to a directory with \texttt{cd}
\item
  List files with \texttt{ls} (macOS/Linux) or \texttt{dir} (Windows)
\item
  Create a directory with \texttt{mkdir}
\item
  Print the current directory with \texttt{pwd} (macOS/Linux) or
  \texttt{cd} with no arguments (Windows)
\item
  Use Tab completion for filenames
\item
  Use Up/Down arrow for command history
\amn{\textbf{Accessibility Tip:} Tab completion is a screen reader user's best friend in the terminal. It reduces typing, eliminates spelling errors in paths, and produces a spoken confirmation of the completed name. Teach it as the default way to enter any filename or directory name.}
\item
  Use Up/Down arrow for command history
\item
  Use \texttt{Ctrl+C} to stop a running program
\end{itemize}

If these are not automatic, spend additional time on terminal basics.
Terminal fluency is foundational to everything else.

\paragraph{Alternative Terminals for
Accessibility}\label{alternative-terminals-for-accessibility}

\begin{fullwidth}{\def\LTcaptype{none} % do not increment counter
\begin{longtable}[]{@{}
  >{\raggedright\arraybackslash}p{(\linewidth - 4\tabcolsep) * \real{0.2597}}
  >{\raggedright\arraybackslash}p{(\linewidth - 4\tabcolsep) * \real{0.1039}}
  >{\raggedright\arraybackslash}p{(\linewidth - 4\tabcolsep) * \real{0.6364}}@{}}
\toprule\noalign{}
\begin{minipage}[b]{\linewidth}\raggedright
Terminal
\end{minipage} & \begin{minipage}[b]{\linewidth}\raggedright
Platform
\end{minipage} & \begin{minipage}[b]{\linewidth}\raggedright
Screen Reader Compatibility
\end{minipage} \\
\midrule\noalign{}
\endhead
\bottomrule\noalign{}
\endlastfoot
\textbf{Windows Terminal} & Windows & Excellent with NVDA/JAWS (use UI
Automation mode) \\
\textbf{cmd.exe} & Windows & Good; legacy but very reliable \\
\textbf{PowerShell} & Windows & Good with NVDA; some JAWS quirks \\
\textbf{macOS Terminal} & macOS & Excellent with VoiceOver \\
\textbf{iTerm2} & macOS & Good; more features, slightly more complex \\
\textbf{GNOME Terminal} & Linux & Good with Orca \\
\textbf{Konsole} & Linux & Good with Orca \\
\textbf{Terminator} & Linux & Good; split-pane support \\
\end{longtable}
}\end{fullwidth}

\paragraph{Thought Experiments}\label{thought-experiments-3}

\begin{enumerate}
\def\labelenumi{\arabic{enumi}.}
\item
  \emph{``Virtual environments isolate each project's dependencies. Why
  would this matter for a professional programmer with 20 active
  projects? What could go wrong without them?''}
\item
  \emph{``When you type \texttt{python3\ hello.py} and press Enter, what
  is the operating system doing? What is Python doing? Walk through it
  step by step as if explaining to someone who has never used a
  computer.''}
\item
  \emph{``\texttt{pip\ install\ requests} downloads and installs code
  written by other people. What are the risks of running code written by
  strangers? What does PyPI (the Python Package Index) do to reduce
  these risks?''}
\end{enumerate}

\begin{center}\rule{0.5\linewidth}{0.5pt}\end{center}

\subsection{6. Unit 1 --- Foundations and Getting
Started}\label{unit-1-foundations-and-getting-started}

\subsection{Lesson 1.1: Welcome to Computer
Science}\label{lesson-1.1-welcome-to-computer-science}

\paragraph{Instructional Cues}\label{instructional-cues-2}

Open with \textbf{the ``Human Computer'' activity} ---  for blind students.
The student
\sidenote{The Human Computer activity works especially well for blind students because they already have deep experience giving precise verbal instructions to navigate the physical world. Orient-and-mobility training is, in a real sense, an exercise in algorithmic thinking.}
gives verbal, step-by-step instructions for a simple physical task. The
teacher follows them literally, making no inferences. This demonstrates
the core challenge of programming: computers are perfectly literal and
do not fill in gaps.

Suggest tasks that connect to the student's daily experience:

\begin{itemize}
\tightlist
\item
  Instructions for making tea or coffee
\item
  Instructions for navigating from the classroom door to their desk
\item
  Instructions for finding a contact in their phone using a screen
  reader
\end{itemize}

After the activity, debrief: \emph{``Where did your instructions assume
I could see? Where did they assume I already knew something? A computer
makes neither assumption.''}

\paragraph{Connecting Computing to Accessibility
Technology}\label{connecting-computing-to-accessibility-technology}

Begin with what the student already knows:

\begin{itemize}
\tightlist
\item
  Their screen reader is software. It runs on hardware. It was written
  by programmers. It converts visual information (text rendering) into
  audio output.
\item
  Ask: ``What are the inputs to your screen reader? What are the
  outputs? What would happen if a website used images of text instead of
  real text --- how would that change what your screen reader can do?''
\item
  
\marginnote{\textbf{ARIA} (Accessible Rich Internet Applications) is a W3C specification that adds semantic meaning to web elements. Poor ARIA implementation is a leading cause of inaccessible websites. A blind student who learns to code is uniquely positioned to understand and advocate for correct ARIA usage.}
\end{itemize}

\paragraph{Thought Experiments}\label{thought-experiments-4}

\begin{enumerate}
\def\labelenumi{\arabic{enumi}.}
\item
  \emph{``Computer science is described as the study of problem-solving,
  not just programming. What non-programming problems have you solved
  using computational thinking --- breaking a problem down into steps,
  looking for patterns, and abstracting the details? Describe one
  example from daily life.''}
\item
  \emph{``Your screen reader converts text to audio in real time. What
  `algorithm' do you think it uses? What decisions does it make --- when
  to pause, how to handle punctuation, what to do with symbols? What
  would a flowchart of those decisions look like?''}
\item
  \emph{``Early computers could only run one program at a time. Modern
  computers run hundreds simultaneously. Without changing the physical
  hardware, what concept makes this possible? (Hint: the answer is
  software --- specifically, an operating system.)''}
\item
  \emph{``Software has often been described as `making the hardware do
  something useful.' But hardware without software is just metal and
  silicon. In what sense is software more important than hardware? In
  what sense is hardware more important?''}
\item
  \emph{``Accessibility technology like screen readers, braille
  displays, and audio description didn't exist when computers were first
  invented. Who decided these should be built? What does this history
  tell you about who gets to shape what technology does?''}
\end{enumerate}

\paragraph{Additional Resources for This
Lesson}\label{additional-resources-for-this-lesson}

\begin{itemize}
\tightlist
\item
  \textbf{``The Computer Boys Take Over'' by Nathan Ensmenger} ---
  accessible history of how programming became a profession.
\item
  \textbf{``Technically Wrong'' by Sara Wachter-Boettcher} --- explores
  how biases in tech affect users, including disabled users. Good for
  older students.
\item
  \textbf{American Foundation for the Blind (AFB) CareerConnect} ---
  profiles of blind professionals in STEM fields, including software
  engineers. Good for motivation and role models.
\item
  \textbf{``Haben: The Deafblind Woman Who Conquered Harvard Law'' by
  Haben Girma} --- autobiography of a deafblind lawyer and disability
  rights advocate. Haben's work directly involves technology access.
\item
  \textbf{Code.org's ``How Computers Work'' video series} ---
  audio-described versions available. Covers hardware and software
  basics accessibly.
\end{itemize}

\begin{center}\rule{0.5\linewidth}{0.5pt}\end{center}

\subsection{Lesson 1.2: Your First Python
Program}\label{lesson-1.2-your-first-python-program}

\paragraph{Instructional Cues}\label{instructional-cues-3}

\textbf{Before any code:} Ask the student to predict. ``If I write
\texttt{print("Hello,\ World!")} and run it, what do you think will
appear in the terminal? Where exactly will it appear?'' Let them answer,
then verify together.

Navigate to the error message together. Establish
the error-reading protocol:
\marginnote{Reading errors from the \emph{bottom up} is counter-intuitive but essential. Python's traceback lists the call chain from outermost to innermost, so the most specific and actionable message is always at the bottom. Establish this habit in the first session.}

\begin{enumerate}
\def\labelenumi{\arabic{enumi}.}
\tightlist
\item
  Terminal output appears after the code.
\item
  Read from the bottom of the error upward --- the bottom is the most
  specific message.
\item
  Line number tells you where to go. Use \texttt{Ctrl+G}.
\item
  Error type (\texttt{SyntaxError}, \texttt{NameError}) narrows the
  category of problem.
\amn{\textbf{Accessibility Tip:} Learn the six most common Python error types and what each means: \texttt{SyntaxError} (malformed code), \texttt{NameError} (undefined variable), \texttt{TypeError} (wrong type), \texttt{ValueError} (right type, bad value), \texttt{IndexError} (list index out of range), \texttt{KeyError} (missing dictionary key). Recognising the type before reading the message cuts debugging time in half.}
\end{enumerate}

Before creating any files, start with \texttt{python3} in the terminal
and type \texttt{print("Hello,\ World!")} directly.
\amn{\textbf{Accessibility Tip:} Open the Python REPL by typing \texttt{python3} and pressing Enter. Every expression you type is immediately evaluated and spoken. Type \texttt{exit()} or press \texttt{Ctrl+D} to quit. The REPL is the fastest way to experiment with Python syntax without creating and running a file.}
\sidenote{\texttt{python3} launches the REPL. On Windows, the command may be \texttt{python} (without the 3) depending on installation. Verify which command works on the student's system and stick to it consistently.}
The immediate audio feedback of the REPL establishes cause and effect
more directly than the file -\textgreater{} run -\textgreater{} find output workflow.

\paragraph{Common First Errors and How to Explain
Them}\label{common-first-errors-and-how-to-explain-them}

\begin{fullwidth}{\def\LTcaptype{none} % do not increment counter
\begin{longtable}[]{@{}
  >{\raggedright\arraybackslash}p{(\linewidth - 4\tabcolsep) * \real{0.2038}}
  >{\raggedright\arraybackslash}p{(\linewidth - 4\tabcolsep) * \real{0.1754}}
  >{\raggedright\arraybackslash}p{(\linewidth - 4\tabcolsep) * \real{0.6209}}@{}}
\toprule\noalign{}
\begin{minipage}[b]{\linewidth}\raggedright
Error
\end{minipage} & \begin{minipage}[b]{\linewidth}\raggedright
Likely Cause
\end{minipage} & \begin{minipage}[b]{\linewidth}\raggedright
Explanation for Student
\end{minipage} \\
\midrule\noalign{}
\endhead
\bottomrule\noalign{}
\endlastfoot
\texttt{SyntaxError:\ unexpected\ EOF\ while\ parsing} & Missing closing
\texttt{"}, \texttt{)}, or \texttt{{]}} & ``Python reached the end of
the file before it expected to. It was waiting for something ---
probably a closing character.'' \\
\texttt{NameError:\ name\ \textquotesingle{}print\textquotesingle{}\ is\ not\ defined}
& Capital P in \texttt{Print} & ``Python is case-sensitive.
\texttt{Print} and \texttt{print} are completely different names. Only
\texttt{print} (lowercase) is the built-in function.'' \\
\texttt{IndentationError} & Extra space before first line & ``Python's
first rule: top-level code cannot be indented. Every line that is not
inside a function or class must start at column 0.'' \\
\texttt{SyntaxError:\ invalid\ syntax} at \texttt{:} & Using = instead
of == in a comparison & ``The assignment operator = doesn't belong
inside an if condition. For comparison, you need the double equals
==.'' \\
\end{longtable}
}
\end{fullwidth}
\paragraph{Thought Experiments}\label{thought-experiments-5}

\begin{enumerate}
\def\labelenumi{\arabic{enumi}.}
\item
  \emph{``\texttt{print()} produces output in the terminal. What would a
  program be like with NO output at all? Could such a program still do
  something useful? (Hint: think about programs that read files, send
  network requests, or modify a database.)''}
\item
  \emph{``Python is an interpreted language --- each line is read and
  executed one at a time. Other languages are compiled --- the entire
  program is translated before any of it runs. How does this affect the
  kind of errors you encounter, and when you encounter them?''}
\item
  \emph{``The \texttt{\#} symbol starts a comment --- text that Python
  ignores entirely. Why do comments exist? If the computer ignores them,
  who are they for? Could you write good programs with no comments?
  Should you?''}
\item
  \emph{``\,`Hello, World!' is the traditional first program in almost
  every programming language. Why do you think this tradition exists?
  What does this simple program prove about a working environment?''}
\end{enumerate}

\paragraph{Alternative Resources for This
Lesson}\label{alternative-resources-for-this-lesson}

\begin{itemize}
\tightlist
\item
  \textbf{``Invent Your Own Computer Games with Python'' by Al Sweigart
  (free online):} Text-based games that produce only terminal output ---
  perfectly accessible. Chapter 1 is a good companion to this lesson.
\item
  \textbf{CS Unplugged --- ``Parity Magic'' activity:} Teaches binary
  representation and algorithms without a computer. Can precede or
  follow this lesson.
\item
  \textbf{Python REPL via iOS/Android with screen reader:} Several apps
  provide a Python REPL on mobile. For students whose phone screen
  reader is more familiar than their computer setup, this may be a
  comfortable way to experiment.
\end{itemize}

\begin{center}\rule{0.5\linewidth}{0.5pt}\end{center}

\subsection{Lesson 1.3: Variables and Simple Data
Types}\label{lesson-1.3-variables-and-simple-data-types}

\paragraph{Instructional Cues}\label{instructional-cues-4}

introduce it verbally and reinforce with tactile analogy if helpful.
\marginnote{For congenitally blind students without strong spatial metaphors, consider an auditory alternative: a variable is like a voice message saved under a contact name. The name stays the same; the message (value) can be re-recorded (reassigned).}
Provide a physical box (a small container) with a label stuck to it and
a slip of paper inside representing the value. The box remains; the
paper changes. The label stays the same; the contents change with
reassignment.

Begin every data type
lesson with \texttt{type()}. \texttt{print(type(42))} -\textgreater{} \texttt{int}.
\marginnote{\texttt{type()} is more important for blind students than for sighted ones. A sighted programmer can infer type from context and syntax highlighting. A screen reader user hears ``42'' whether the value is an integer or a string. \texttt{type()} is the equalizer.} \texttt{print(type("hello"))} -\textgreater{}
\texttt{str}. Make this a reflex the student uses whenever they're
unsure what kind of data they have. For a student who cannot visually
scan code for context clues, \texttt{type()} is often the only reliable
way to confirm what a variable contains.

\textbf{The string/number audio confusion:} When a screen reader reads
\texttt{"42"} and \texttt{42}, they sound identical.
\amn{\textbf{Accessibility Tip:} After any \texttt{input()} call, immediately call \texttt{type()} on the result before doing anything with it. \texttt{"42"} and \texttt{42} are read identically by a screen reader but behave completely differently: \texttt{"42" + 1} raises \texttt{TypeError}, \texttt{42 + 1} gives \texttt{43}.} This is a critical
accessibility issue. Address it directly:
\sidenote{This is the accessibility gotcha that surprises even experienced VI educators. Build in a \texttt{type()} check as a classroom ritual: after every \texttt{input()} call, immediately print \texttt{type()} of the result before using it.}

\begin{itemize}
\tightlist
\item
  \emph{``Your screen reader will read these the same way. The only
  reliable way to tell them apart is \texttt{type()}. Use it.
  Constantly.''}
\item
  Demonstrate the consequence: \texttt{"42"\ +\ 1} raises
  \texttt{TypeError}. \texttt{42\ +\ 1} gives \texttt{43}.
\end{itemize}

\paragraph{Multiple Explanations for Variable Naming
Rules}\label{multiple-explanations-for-variable-naming-rules}

Provide all of these; let the student choose what sticks:

\begin{itemize}
\tightlist
\item
  \emph{Rule-based:} Must start with letter or underscore. May contain
  letters, digits, underscores. May not be a Python keyword.
\item
  \emph{Analogy:} ``Variable names are like legal names. They can't
  start with a number (no one's legal name is `2Fast'), can't contain
  hyphens (the name Mc-Donald would confuse paperwork), and can't be
  reserved words like `True' or `if' that already mean something
  official.''
\item
  \emph{Mnemonic for invalid starts:} ``Numbers first? Never. Hyphens?
  Horrible. Spaces? Surely not.''
\end{itemize}

\paragraph{Thought Experiments}\label{thought-experiments-6}

\begin{enumerate}
\def\labelenumi{\arabic{enumi}.}
\item
  \emph{``If variables are labeled boxes, what happens when you do
  \texttt{x\ =\ 5} and then \texttt{y\ =\ x} and then \texttt{x\ =\ 10}?
  What is \texttt{y} now? Does reassigning \texttt{x} affect \texttt{y}?
  Why or why not?''}
\item
  \emph{``Python is `dynamically typed,' meaning a variable can change
  its type at any time: \texttt{x\ =\ 5} then
  \texttt{x\ =\ \textquotesingle{}hello\textquotesingle{}}. Many
  languages forbid this. What bugs could dynamic typing accidentally
  create? What flexibility does it offer?''}
\item
  \emph{``The integer \texttt{42} and the string
  \texttt{\textquotesingle{}42\textquotesingle{}} look identical when
  spoken aloud. How many other pairs of values can you think of that
  might be ambiguous when a screen reader reads them? What coding habits
  help prevent confusion?''}
\item
  \emph{``Constants in programming represent values that never change
  --- like \texttt{SPEED\_OF\_LIGHT} or \texttt{MAX\_USERS}. Python has
  no true constants (you can always reassign them). Why do programmers
  use \texttt{CAPS\_SNAKE\_CASE} as a convention for constants even
  though Python can't enforce them?''}
\end{enumerate}

\paragraph{Alternative Resources}\label{alternative-resources}

\begin{itemize}
\tightlist
\item
  \textbf{CS50 Python (Harvard, edX):} David Malan's accessible video
  lectures. Closed-captions and transcripts available for all lessons.
  Covers variables and types in Week 0.
\item
  \textbf{Python.org official tutorial:} Plain HTML, navigable by screen
  reader. Section 3 covers numbers, strings, and basic types.
\item
  \textbf{``Python for Everybody'' by Dr.~Chuck Severance (free
  online):} Audio-described video lectures available. Chapter 2 covers
  variables and expressions.
\item
  \textbf{Codecademy ``Learn Python 3'':} Web-based coding environment.
  Accessibility is partial --- code editor is accessible but some UI
  elements are not. Worth testing with the student's specific setup.
\end{itemize}

\begin{center}\rule{0.5\linewidth}{0.5pt}\end{center}

\subsection{7. Unit 2 --- Numeric Computations and
I/O}\label{unit-2-numeric-computations-and-io}

\subsection{Lesson 2.1: The Software Development
Process}\label{lesson-2.1-the-software-development-process}

\paragraph{Instructional Cues}\label{instructional-cues-5}

This is a conceptual lesson. Resist the urge to rush to code. Spend the
full session on design and pseudocode.

Before writing pseudocode
ask the
student to state the problem in one sentence. Then ask ``why?'' five
times to reach the root of what the program really needs to do.
\marginnote{The Five Whys method originated in Toyota's manufacturing quality system (kaizen). Its application to software design was popularized by the lean startup movement. The technique is equally effective for VI students because it produces verbal, non-visual problem specifications.} Then ask ``why?'' five
times to reach the root of what the program really needs to do. This
technique from Lean manufacturing translates well to programming.

Example: ``My program calculates a restaurant tip.''

\begin{itemize}
\tightlist
\item
  Why? ``So the user knows how much to pay.''
\item
  Why do they need to calculate it? ``The tip amount isn't printed on
  the bill.''
\item
  Why? ``The restaurant leaves the amount to the customer's
  discretion.''
\item
  So the real problem is: ``Given a bill amount, calculate 15\%, 18\%,
  and 20\% tip options, plus the total.''
\end{itemize}

This more complete problem statement produces a better program design.

\textbf{Pseudocode format for screen reader users:} Plain \texttt{.txt}
files work best for pseudocode. Avoid word processors or special
formatting tools that add accessibility overhead. The student should be
able to read their pseudocode fluently while writing their Python.

\paragraph{Thought Experiments}\label{thought-experiments-7}

\begin{enumerate}
\def\labelenumi{\arabic{enumi}.}
\item
  \emph{``In engineering, a bridge is designed many times before a
  single beam is placed. The cost of redesigning a blueprint is near
  zero; the cost of redesigning a half-built bridge is enormous. How
  does this ratio (cheap to redesign early, expensive to redesign later)
  apply to software?''}
\item
  \emph{``The design phase involves thinking about edge cases ---
  unusual inputs. For a program that calculates how many buses are
  needed to transport N students, with each bus holding 40 students,
  what edge cases exist? (Consider: What if N=0? What if N=40? What if
  N=41? What if N is negative?)''}
\item
  \emph{``Some agile software development methodologies say `working
  software over comprehensive documentation.' Does this conflict with
  the value of design? How do you reconcile the two?''}
\item
  \emph{``Testing is described as part of the development process. But
  some teams test after development is complete. What is lost by testing
  at the end rather than continuously throughout? What is gained?''}
\end{enumerate}

\begin{center}\rule{0.5\linewidth}{0.5pt}\end{center}

\subsection{Lessons 2.2--2.5: Expressions, Input, Types,
Math}\label{lessons-2.22.5-expressions-input-types-math}

\paragraph{Instructional Cues}\label{instructional-cues-6}


\sidenote{This is one of two facts about \texttt{input()} that every beginner gets wrong at least once. The other is that it pauses program execution until the user presses Enter. Both are invisible to a reader scanning the code without running it.}

\begin{Shaded}
\begin{Highlighting}[]
user\_input \OperatorTok{=} input(\StringTok{"Type the number 5: "})
print(type(user\_input))   \CommentTok{\# \textless{}class \textquotesingle{}str\textquotesingle{}\textgreater{}}
print(user\_input \OperatorTok{+} \CodeNumTok{1})      \CommentTok{\# TypeError}
print(int(user\_input) \OperatorTok{+} \CodeNumTok{1}) \CommentTok{\# 6}
\end{Highlighting}
\end{Shaded}

Run this sequence and let the student hear both the error and the fix.
The \texttt{TypeError} message should be read aloud and analyzed before
the fix is shown.

\textbf{Order of operations --- the hearing test:} Read
\texttt{10\ +\ 5\ *\ 2} aloud and ask for the answer before Python is
consulted. Most students say 30; Python says 20. This is a memorable
demonstration of why precise symbol reading matters. ``Your ear parsed
it as `add 10 and 5, then multiply by 2.' Python parses it as `10 plus
the result of 5 times 2.'\,''

\textbf{\texttt{math} module exploration via REPL:}

\begin{Shaded}
\begin{Highlighting}[]
\CodeAccentTok{import} math
dir(math)  \CommentTok{\# Lists all functions in the module}
help(math.sqrt)  \CommentTok{\# Reads the docstring aloud}
\end{Highlighting}
\end{Shaded}

\texttt{dir()} and \texttt{help()} are screen-reader-friendly tools for
exploring any module. Teach them as first-resort documentation tools.

\paragraph{Thought Experiments}\label{thought-experiments-8}

\begin{enumerate}
\def\labelenumi{\arabic{enumi}.}
\item
  \emph{``Type conversion functions like \texttt{int()} and
  \texttt{float()} can fail:
  \texttt{int(\textquotesingle{}hello\textquotesingle{})} raises a
  \texttt{ValueError}. Who is responsible --- the programmer who called
  \texttt{int()} on untrusted input, or the user who typed `hello'? How
  should a well-designed program handle this?''}
\item
  \emph{``\texttt{math.sqrt(25)} returns \texttt{5.0}, not \texttt{5}.
  Why does a function that logically produces a whole number return a
  float? What would break if it returned an integer instead?''}
\item
  \emph{``Floating-point numbers in computers are approximations ---
  \texttt{0.1\ +\ 0.2} gives \texttt{0.30000000000000004}. Describe a
  real-world application where this imprecision could cause serious
  problems. What strategies exist to avoid it?''}
\item
  \emph{``The \texttt{//} operator performs floor division:
  \texttt{7\ //\ 2} = \texttt{3}. The \texttt{\%} operator gives the
  remainder: \texttt{7\ \%\ 2} = \texttt{1}. Together they completely
  describe integer division. What real-world problem does each solve?
  Describe a program that would need both.''}
\end{enumerate}

\paragraph{Alternative Resources}\label{alternative-resources-1}

\begin{itemize}
\tightlist
\item
  \textbf{``Automate the Boring Stuff with Python'' Chapter 1--3 (free
  online):} Covers expressions, data types, and input with practical
  exercises. Plain HTML, accessible.
\item
  \textbf{Khan Academy ``Introduction to Programming'':} Video series
  with transcripts. Good for visual-thinker students who benefit from
  hearing concepts explained differently.
\item
  \textbf{SoloLearn Python course:} Mobile app with accessible
  exercises. Works well with iOS VoiceOver and Android TalkBack for
  students who prefer mobile.
\item
  \textbf{``Python Crash Course'' by Eric Matthes:} The primary textbook
  for this curriculum. Chapter 2 aligns directly with this unit.
\end{itemize}

\begin{center}\rule{0.5\linewidth}{0.5pt}\end{center}

\subsection{8. Unit 3 --- Control Flow, Decisions, and
Loops}\label{unit-3-control-flow-decisions-and-loops}

\subsection{Lessons 3.1--3.2:
Conditionals}\label{lessons-3.13.2-conditionals}

\paragraph{Instructional Cues}\label{instructional-cues-7}

\textbf{Verbal simulation protocol:} Before running any conditional
code, have the student perform the simulation verbally:

\begin{enumerate}
\def\labelenumi{\arabic{enumi}.}
\tightlist
\item
  ``What is the value of \texttt{score}?''
\item
  ``Evaluate the condition: is \texttt{85\ \textgreater{}=\ 90}? That's
  False.''
\item
  ``Because the condition is False, the if block is skipped.''
\item
  ``Evaluate the next condition: is \texttt{85\ \textgreater{}=\ 80}?
  That's True.''
\item
  ``The elif block runs. We print `Grade: B'.''
\end{enumerate}

This 5-step protocol should become automatic before the student runs any
code. It builds the mental simulation ability that is the core of
debugging.
\marginnote{Mental simulation---tracing code execution step by step without running it---is more important for blind students than for sighted ones. A sighted programmer can run code and scan output visually in seconds. A blind programmer benefits from knowing what \emph{should} happen before navigating the terminal output to verify it.}

\textbf{The assignment-vs-comparison trap:}

\begin{Shaded}
\begin{Highlighting}[]
\CommentTok{\# This is a bug that is SILENT in some languages}
\KeywordTok{if} score \OperatorTok{=} \CodeNumTok{90}:   \CommentTok{\# SyntaxError in Python (good!)}
    print(\StringTok{"Perfect!"})
\end{Highlighting}
\end{Shaded}

Python raises a \texttt{SyntaxError} here, which is actually a safety
feature. Other languages (C, JavaScript) allow this and it becomes a
silent logic bug. Appreciate Python's strictness here.
\sidenote{JavaScript's version of this bug is \texttt{if (score = 90)} which silently assigns 90 to \texttt{score} and evaluates as truthy---no error, wrong behavior. Python's refusal to allow assignment inside a condition is a genuine accessibility feature: it converts a silent semantic error into a loud syntactic one.}

\paragraph{Thought Experiments}\label{thought-experiments-9}

\begin{enumerate}
\def\labelenumi{\arabic{enumi}.}
\item
  \emph{``If Python didn't have \texttt{elif}, only \texttt{if} and
  \texttt{else}, could you still write a grading program? Show how. What
  is the cost of doing it without \texttt{elif}?''}
\item
  \emph{``The \texttt{else} clause at the end of an if-elif chain is a
  catch-all --- it runs if nothing else matched. What is the programmer
  asserting when they write
  \texttt{else:\ print(\textquotesingle{}Invalid\ input\textquotesingle{})}?
  What assumptions about the universe of possible inputs does this make?
  When might those assumptions be wrong?''}
\item
  \emph{``Boolean logic (\texttt{and}, \texttt{or}, \texttt{not}) comes
  from mathematics (George Boole, 1847). Boolean logic also describes
  how circuits in a computer work. In what sense is every program you
  write built on top of millions of tiny \texttt{and}, \texttt{or}, and
  \texttt{not} operations happening in hardware?''}
\item
  \emph{``Short-circuit evaluation: in \texttt{if\ A\ and\ B}, Python
  doesn't evaluate \texttt{B} if \texttt{A} is False (because the result
  is already False). In \texttt{if\ A\ or\ B}, Python doesn't evaluate
  \texttt{B} if \texttt{A} is True. When is this behavior a useful
  feature rather than just an optimization?''}
\item
  \emph{``Consider a program that decides whether to grant access based
  on two conditions: correct password AND recognized device. Map this to
  Python \texttt{and} logic. Now consider: unlock if either condition
  passes (maybe for emergency access). Map this to \texttt{or}. What are
  the security implications of each design?''}
\end{enumerate}

\paragraph{Alternative Resources}\label{alternative-resources-2}

\begin{itemize}
\tightlist
\item
  \textbf{CS Unplugged ``Binary Numbers'' activity:} Physical,
  non-computer activity that illustrates Boolean thinking. Can be
  adapted for tactile use.
\item
  \textbf{``Think Python'' by Allen B. Downey (free online):} Chapter 5
  covers conditionals with clear explanation and exercises. HTML version
  is screen-reader-friendly.
\item
  \textbf{Parsons Problems (online):} Code ordering exercises --- the
  student arranges given lines into correct order. Does not require
  typing code from scratch. Available at parsons.problemsolving.org
  (accessibility varies; test first).
\end{itemize}

\begin{center}\rule{0.5\linewidth}{0.5pt}\end{center}

\subsection{Lessons 3.3--3.4: Lists and For
Loops}\label{lessons-3.33.4-lists-and-for-loops}

\paragraph{Instructional Cues}\label{instructional-cues-8}

\textbf{Mental models for lists --- choose one per student:}

\begin{itemize}
\tightlist
\item
  \emph{The train:} Each car has a number starting at 0. Car 0 is at the
  front. You can add cars to the end, remove from anywhere, and swap
  contents.
\item
  \emph{The playlist:} Songs in order, starting at position 0. You can
  shuffle, sort, add, or remove tracks.
\item
  \emph{The physical stack of cards:} A sequence of items that can be
  addressed by position.
\end{itemize}

\textbf{The for loop as a bookmark reader:} \emph{``Imagine Python has a
finger that moves along a list. It starts at position 0. It reads what's
there, runs the loop body, moves the finger to position 1, reads, runs
the body, and continues until the finger falls off the end.''}

\textbf{Tracing loops audibly --- required exercise:} For the first five
loops students write, require them to verbally trace every iteration
before running the code:
\marginnote{Verbal tracing is not just a learning exercise---it is a professional debugging technique. Experienced blind developers often trace loops aloud or on paper before running code because navigating loop output line-by-line with a screen reader is time-consuming. Building this habit early pays dividends.}

\begin{itemize}
\tightlist
\item
  ``Iteration 1: \texttt{fruit} gets `apple'. We print `I like
  apple.'\,''
\item
  ``Iteration 2: \texttt{fruit} gets `banana'. We print `I like
  banana.'\,''
\item
  etc.
\end{itemize}

\paragraph{Non-Obvious Accessibility
Strategies}\label{non-obvious-accessibility-strategies-1}

\begin{itemize}
\item
  \textbf{Printing list contents accessibly:} A bare
  \texttt{print(my\_list)} produces
  \texttt{{[}\textquotesingle{}apple\textquotesingle{},\ \textquotesingle{}banana\textquotesingle{},\ \textquotesingle{}cherry\textquotesingle{}{]}}
  with brackets and quotes. This is cluttered when read aloud. Teach the
  clean print pattern early:

\begin{Shaded}
\begin{Highlighting}[]
\KeywordTok{for} i, item \KeywordTok{in} enumerate(my\_list):
    print(\StringTok{f"  }\StringTok{\{}i\StringTok{\}}\StringTok{: }\StringTok{\{}item\StringTok{\}}\StringTok{"})
\end{Highlighting}
\end{Shaded}

  This gives numbered output, one item per line, that is easy to
  navigate and understand when heard.
\sidenote{Number-then-item output (\texttt{1: apple}, \texttt{2: banana}) is the terminal equivalent of a screen reader's list navigation mode. The number gives position context that raw list notation \texttt{['apple', 'banana']} does not provide when spoken sequentially.}
\item
  \textbf{\texttt{len()} as a safety check before index access:} Teach
  students to always know the length before indexing into a list. This
  prevents \texttt{IndexError}. Make \texttt{len(my\_list)} the first
  thing printed when debugging a list problem.
\item
  \textbf{Index arithmetic verbalized:} When accessing
  \texttt{my\_list{[}-1{]}}, say ``the last element, which is at index
  \texttt{len(my\_list)\ minus\ 1}, which Python shortens to
  \texttt{-1}.'' Connecting the shorthand to the underlying meaning
  prevents confusion.
\end{itemize}

\paragraph{Thought Experiments}\label{thought-experiments-10}

\begin{enumerate}
\def\labelenumi{\arabic{enumi}.}
\item
  \emph{``Zero-based indexing means the first element is at position 0.
  This feels unnatural but is standard across almost all programming
  languages. Why? (Hint: think about memory addresses --- the address of
  the first element, and the address of any element as an offset from
  there.)''}
\item
  \emph{``A list allows duplicate values. A set (which you'll learn
  later) does not. When would duplicates in a list be meaningful data?
  When would they indicate an error in your program?''}
\item
  \emph{``A \texttt{for} loop iterates over every element in a list. Is
  there any computation a \texttt{for} loop can do that a \texttt{while}
  loop cannot? Is there any computation a \texttt{while} loop can do
  that a \texttt{for} loop cannot?''}
\item
  \emph{``What happens when you iterate over an empty list:
  \texttt{for\ item\ in\ {[}{]}:}? Does Python crash, print an error, or
  do nothing? Why is `do nothing' the correct behavior rather than an
  error?''}
\end{enumerate}

\begin{center}\rule{0.5\linewidth}{0.5pt}\end{center}

\subsection{Lesson 3.5: String
Operations}\label{lesson-3.5-string-operations}

\paragraph{Instructional Cues}\label{instructional-cues-9}

\textbf{Immutability demonstration:}

\begin{Shaded}
\begin{Highlighting}[]
text \OperatorTok{=} \StringTok{"hello"}
text.upper()        \CommentTok{\# Creates a new string {-}{-}{-} doesn\textquotesingle{}t modify text}
print(text)         \CommentTok{\# Still "hello"}
text \OperatorTok{=} text.upper() \CommentTok{\# Now reassigned}
print(text)         \CommentTok{\# "HELLO"}
\end{Highlighting}
\end{Shaded}

This demonstration catches a persistent beginner misunderstanding. Ask
the student to predict the output of \texttt{print(text)} after
\texttt{text.upper()} before running it.

\textbf{Method chaining verbalized:}
\texttt{"\ \ hello\ world\ \ ".strip().title()} --- read this as:
``Start with the string. Strip whitespace from both ends. Then convert
the result to title case.''

\paragraph{Thought Experiments}\label{thought-experiments-11}

\begin{enumerate}
\def\labelenumi{\arabic{enumi}.}
\item
  \emph{``Strings are immutable --- they cannot be changed, only
  replaced. Lists are mutable --- they can be changed in-place. What
  architectural difference between strings and lists makes one mutable
  and the other not? (Think about how they're stored in memory.)''}
\item
  \emph{``\texttt{.split()} splits a string into a list of words.
  \texttt{.join()} combines a list of strings into one string. They are
  inverses. What is the inverse of \texttt{.replace()}? Does it have a
  direct inverse? If not, how would you reverse a replacement?''}
\item
  \emph{``String formatting with f-strings:
  \texttt{f\textquotesingle{}Value\ is\ \{x:.2f\}\textquotesingle{}}
  formats \texttt{x} as a float with 2 decimal places. Why does a
  programming language need a mini-language for formatting numbers? What
  would programs look like if we had to format every number manually
  with division and modulo?''}
\end{enumerate}

\paragraph{Alternative Resources}\label{alternative-resources-3}

\begin{itemize}
\tightlist
\item
  \textbf{Python String Methods documentation}
  (docs.python.org/3/library/stdtypes.html\#string-methods): Plain HTML
  table of all string methods. Screen-reader navigable. Essential
  reference.
\item
  \textbf{``Python for Data Analysis'' by Wes McKinney:} Chapter on
  string manipulation provides real-world examples beyond what this
  curriculum covers. For motivated students.
\end{itemize}

\begin{center}\rule{0.5\linewidth}{0.5pt}\end{center}

\subsection{9. Unit 4 --- Data Structures: Lists and
Dictionaries}\label{unit-4-data-structures-lists-and-dictionaries}

\emph{(See all lessons above for list material; this section adds
dictionary-specific guidance.)}

\subsection{Dictionary Lessons}\label{dictionary-lessons}

\paragraph{Instructional Cues}\label{instructional-cues-10}

\textbf{Multiple mental models for dictionaries:}

\begin{itemize}
\tightlist
\item
  \emph{Real dictionary:} You look up a word (key) to find its
  definition (value). Words are unique; definitions can be anything.
\item
  \emph{Contact list:} Each name (key) maps to a phone number (value).
  No two contacts have the same name in a well-maintained list.
\item
  \emph{JSON configuration file:} Settings names map to their values.
  This foreshadows Unit 8.
\end{itemize}

\textbf{The \texttt{.get()} method vs.~direct access:}

\begin{Shaded}
\begin{Highlighting}[]
\CommentTok{\# This crashes on missing key}
print(my\_dict[\StringTok{"nonexistent\_key"}])   \CommentTok{\# KeyError}
 
\CommentTok{\# This returns a default {-}{-}{-} never crashes}
print(my\_dict.get(\StringTok{"nonexistent\_key"}, \StringTok{"Not found"}))  \CommentTok{\# "Not found"}
\end{Highlighting}
\end{Shaded}

For a screen reader user who may not immediately catch a crash in
terminal output, teaching \texttt{.get()} as the default access pattern
reduces debugging overhead.
\marginnote{A \texttt{KeyError} crash produces output that looks similar to normal output} \amn{\textbf{Accessibility Tip:} Use \texttt{dict.get(key, default)} instead of \texttt{dict[key]} whenever a key might not exist. A \texttt{KeyError} crash is easy to miss when navigating terminal output with a screen reader, since the error appears in the same buffer as normal output. \texttt{.get()} prevents the crash entirely.}
\marginnote{A \texttt{KeyError} crash produces output that looks similar to normal output when navigating with a screen reader---it appears in the same terminal buffer. Using \texttt{.get()} with a default eliminates an entire category of crash that might otherwise go unnoticed mid-session.}

\paragraph{Thought Experiments}\label{thought-experiments-12}

\begin{enumerate}
\def\labelenumi{\arabic{enumi}.}
\item
  \emph{``Dictionaries use keys to retrieve values in O(1) time ---
  essentially instant regardless of dictionary size. Lists use indices,
  also O(1). But searching a list for a specific value (not by index)
  takes O(n) time --- proportional to list length. When does this
  difference matter in practice?''}
\item
  \emph{``Your dictionary cannot have duplicate keys --- adding a second
  value for an existing key replaces the first. What real-world data
  problems does this property solve? What problems does it create?''}
\item
  \emph{``JSON (JavaScript Object Notation) is essentially Python
  dictionaries written in a slightly different format. What does it mean
  that a data format invented for JavaScript has become the universal
  language of data exchange? What does this tell you about the value of
  common standards?''}
\item
  \emph{``A `nested dictionary' contains other dictionaries as values.
  Model a university course catalog as a nested dictionary. What are the
  outer keys? What are the inner keys? What would
  \texttt{catalog{[}\textquotesingle{}CS\textquotesingle{}{]}{[}\textquotesingle{}101\textquotesingle{}{]}{[}\textquotesingle{}instructor\textquotesingle{}{]}}
  represent?''}
\end{enumerate}

\begin{center}\rule{0.5\linewidth}{0.5pt}\end{center}

\subsection{10. Unit 5 --- User Input and While
Loops}\label{unit-5-user-input-and-while-loops}

\subsection{Thought Experiments}\label{thought-experiments-13}

\begin{enumerate}
\def\labelenumi{\arabic{enumi}.}
\item
  \emph{``A while loop with a flag variable and a while loop with
  \texttt{break} can often accomplish the same thing. Compare these two
  equivalent programs and explain when you would prefer each:}

\begin{Shaded}
\begin{Highlighting}[]
\CommentTok{\# Version A: flag}
active \OperatorTok{=} True
\KeywordTok{while} active:
    cmd \OperatorTok{=} input(\StringTok{"\textgreater{} "})
    \KeywordTok{if} cmd \OperatorTok{==} \StringTok{"quit"}: active \OperatorTok{=} False
    \KeywordTok{else}: print(\StringTok{"You said:"}, cmd)
 
\CommentTok{\# Version B: break}
\KeywordTok{while} True:
    cmd \OperatorTok{=} input(\StringTok{"\textgreater{} "})
    \KeywordTok{if} cmd \OperatorTok{==} \StringTok{"quit"}: \KeywordTok{break}
    print(\StringTok{"You said:"}, cmd)
\end{Highlighting}
\end{Shaded}

  \emph{Which is clearer? Which is easier to extend if you add more exit
  conditions?{\kern0pt}``}
\item
  \emph{``An infinite loop that waits for user input is the core pattern
  of every interactive program --- email clients, video games, operating
  systems. Describe the `event loop' of your screen reader in terms of
  while loops and if statements.while loops and if statements.
What is it waiting for? What does it do}
\marginnote{This thought experiment has real depth for VI students: they \emph{use} an event loop every day. A screen reader's main loop listens for keyboard events, focus changes, text insertions, and spoken output requests. Understanding that their screen reader is ``just'' a while loop with a lot of if statements demystifies the tool and connects programming concepts to lived experience.}
  \emph{when it gets an event?''}
\item
  \emph{``The \texttt{continue} statement skips the rest of the current
  iteration. An alternative is to wrap the loop body in an \texttt{if}
  statement. When is \texttt{continue} cleaner? When does the
  \texttt{if} approach produce clearer code?''}
\end{enumerate}

\begin{center}\rule{0.5\linewidth}{0.5pt}\end{center}

\subsection{11. Unit 6 --- Functions}\label{unit-6-functions}

\subsection{Thought Experiments (in addition to existing
guide)}\label{thought-experiments-in-addition-to-existing-guide}

\begin{enumerate}
\def\labelenumi{\arabic{enumi}.}
\item
  \emph{``A function is described as a `black box' --- you give it
  inputs and get outputs without knowing the internal details.the internal details.
What are
\marginnote{The black-box metaphor is particularly resonant for blind students: they already interact with most technology as a black box, trusting that inputs produce the correct outputs without being able to inspect the internals. Programming inverts this: the student is now the author of the black box, responsible for what comes out.}
  the benefits of this abstraction for the programmer who writes the
  function? For the programmer who uses it? What are the limits of this
  metaphor?''}
\item
  \emph{``Default parameter values:
  \texttt{def\ greet(name=\textquotesingle{}World\textquotesingle{})}.
  This allows the function to be called with or without an argument. How
  does this enable backward compatibility --- adding a new feature
  without breaking existing code that calls the function?''}
\item
  \emph{``Functions should `do one thing.' What does `one thing' mean?
  Is \texttt{calculate\_and\_print\_tax()} one thing or two? If you
  separate it into \texttt{calculate\_tax()} and \texttt{print\_tax()},
  what becomes possible that wasn't before?''}
\item
  \_``Recursive functions call themselves.
  \texttt{factorial(n)\ =\ n\ \_\ factorial(n-1)}. What is the base case
  and why is it required? What goes wrong if a recursive function has no
  base case?''*
\item
  \_``A function is a way to name a piece of code. Why does naming
  matter in programming? What is the cognitive difference between
  \texttt{x\ =\ a\ \_\ b\ +\ c\ -\ d}and\texttt{x\ =\ calculate\_net\_profit(revenue,\ cost,\ overhead,\ taxes)}?''*
\end{enumerate}

\paragraph{Alternative Resources}\label{alternative-resources-4}

\begin{itemize}
\tightlist
\item
  \textbf{``Structure and Interpretation of Computer Programs'' (SICP)
  --- MIT OpenCourseWare:} Conceptually deep treatment of functions as
  first-class objects. Lectures have transcripts. For advanced students.
\item
  \textbf{Exercism.io Python Track:} Mentor-guided exercises with
  community feedback. Web platform with reasonable accessibility.
  Students can complete exercises at their own pace.
\item
  \textbf{Python Tutor (pythontutor.com):} Step-by-step execution with
  variable state display. The text-based ``state panel'' below the code
  is screen-reader accessible and shows variable values at each step.
\end{itemize}

\begin{center}\rule{0.5\linewidth}{0.5pt}\end{center}

\subsection{12. Unit 7 --- Classes and Object-Oriented
Programming}\label{unit-7-classes-and-object-oriented-programming}

\subsection{Extended Instructional
Cues}\label{extended-instructional-cues}

\textbf{Staged introduction to OOP --- do not rush to syntax:}

Stage 1 (no code): Discuss the concept. What is an object? What does
``bundling data and behavior'' mean? Use tangible examples the student
knows.

Stage 2 (pseudocode): Write a class in plain English before Python
syntax:

\begin{verbatim}
Class: BankAccount
  Has: owner_name, balance, account_number
  Can do: deposit(amount), withdraw(amount), get_balance()
\end{verbatim}

Stage 3 (Python skeleton): Write the class structure with empty methods
and docstrings.

Stage 4 (full implementation): Fill in method bodies.

This staging prevents syntax confusion from obscuring the conceptual
learning.
\marginnote{The four-stage OOP introduction (verbal, pseudocode, skeleton, implementation) maps directly to the Concrete-Representational-Abstract (CRA) instructional framework used in special education. VI students benefit from spending longer at the verbal and pseudocode stages than sighted students, who can re-read syntax examples visually.}

\textbf{\texttt{\_\_str\_\_} is mandatory for accessibility:} Every
class a student creates must define \texttt{\_\_str\_\_}. define \texttt{\_\_str\_\_}.
Without it,
\marginnote{Make \texttt{\_\_str\_\_} a graded requirement on every assignment involving classes. Frame it as an accessibility contract: the programmer is responsible for ensuring their objects produce meaningful audio output. This connects programming practice to the student's lived experience of inaccessible software.}
\texttt{print(my\_object)} produces
\texttt{\textless{}\_\_main\_\_.BankAccount\ object\ at\ 0x7f...\textgreater{}}
--- meaningless when heard. The requirement to define
\texttt{\_\_str\_\_} also reinforces thinking about what the object's
most important information is.

\begin{Shaded}
\begin{Highlighting}[]
\KeywordTok{def} \CodeFuncTok{\_\_str\_\_}(self):
    \KeywordTok{return} \StringTok{f"BankAccount(owner=}\StringTok{\{}self\StringTok{.}owner\_name\StringTok{\}}\StringTok{, balance=$}\StringTok{\{}self\StringTok{.}balance\StringTok{:.2f\}}\StringTok{)"}
\end{Highlighting}
\end{Shaded}

\paragraph{Thought Experiments}\label{thought-experiments-14}

\begin{enumerate}
\def\labelenumi{\arabic{enumi}.}
\item
  \emph{``Encapsulation means hiding internal details. A car's engine is
  encapsulated under the hood --- you don't need to know how fuel
  injection works to drive. What are the analogous `hidden internals' in
  a BankAccount class? What should be private? What should be public?
  Why?''}
\item
  \emph{``Inheritance creates `is-a' relationships. Composition creates
  `has-a' relationships. A \texttt{SavingsAccount} is-a
  \texttt{BankAccount}. A \texttt{Bank} has-a list of
  \texttt{BankAccount} objects. Design a simple banking system using
  both relationships. What would go wrong if you tried to use only
  inheritance?''}
\item
  \emph{``Method overriding allows a subclass to replace a parent's
  method. When is this the right design? When should you call
  \texttt{super().method()} inside the override vs.~replace the method
  entirely?''}
\item
  \emph{``Python allows multiple inheritance (inheriting from two parent
  classes). Most languages do not permit this. What problems does
  multiple inheritance create? What problems does it solve? Would you
  use it?''}
\item
  \emph{``Object-oriented programming was supposed to make code more
  reusable. Has it delivered on this promise? What kinds of code reuse
  are well-served by OOP? What kinds aren't?''}
\end{enumerate}

\paragraph{Alternative Resources}\label{alternative-resources-5}

\begin{itemize}
\tightlist
\item
  \textbf{``Object-Oriented Programming in Python'' by Michael
  Goldwasser and David Letscher:} Textbook specifically pairing Python
  and OOP concepts. More formal treatment for students who want depth.
\item
  \textbf{``Fluent Python'' by Luciano Ramalho:} Advanced but deeply
  accessible to screen readers. Covers Python's data model including
  \texttt{\_\_str\_\_}, \texttt{\_\_repr\_\_}, and special methods. For
  motivated students.
\item
  \textbf{Real Python ``Object-Oriented Programming in Python'' tutorial
  (realpython.com/python3-object-oriented-programming/):} Long but
  thorough. Plain HTML.
\end{itemize}

\begin{center}\rule{0.5\linewidth}{0.5pt}\end{center}

\subsection{13. Unit 8 --- Files and
Exceptions}\label{unit-8-files-and-exceptions}

\subsection{Extended Instructional
Cues}\label{extended-instructional-cues-1}

\textbf{Files as persistence --- make it tangible:} Run a program that
stores a number in a variable, quits, and starts again --- the variable
is gone. Then run the same program storing the number in a file --- it
persists. Ask: ``What changed? What is different about a file vs.~a
variable?'' Let the student reason through it before explaining.

\textbf{The \texttt{with} statement --- demonstrate the alternative:}
\marginnote{The \texttt{with} statement is a context manager.} \amn{\textbf{Accessibility Tip:} Always use \texttt{with open(filename) as f:} rather than \texttt{f = open(...)} / \texttt{f.close()}. The \texttt{with} block guarantees the file is closed even if the program crashes mid-read. Unclosed files can cause silent data loss that produces no error message.}
\marginnote{The \texttt{with} statement is a context manager. Python's context manager protocol (\texttt{\_\_enter\_\_} and \texttt{\_\_exit\_\_}) guarantees cleanup even when exceptions are raised. For files, databases, and network connections, unclosed resources cause silent data loss or corruption---particularly hard to debug with a screen reader.}

\begin{Shaded}
\begin{Highlighting}[]
\CommentTok{\# Without \textquotesingle{}with\textquotesingle{} {-}{-}{-} dangerous}
f \OperatorTok{=} open(\StringTok{"data.txt"}, \StringTok{"r"})
contents \OperatorTok{=} f.read()
\CommentTok{\# If an error occurs here, f.close() is never called}
f.close()
 
\CommentTok{\# With \textquotesingle{}with\textquotesingle{} {-}{-}{-} safe}
\KeywordTok{with} open(\StringTok{"data.txt"}, \StringTok{"r"}) \CodeAccentTok{as} f:
    contents \OperatorTok{=} f.read()
\CommentTok{\# f is closed automatically even if an error occurs}
\end{Highlighting}
\end{Shaded}

\textbf{Exception types --- teach by demonstration, not memorization:}
Create a ``crash collection'' --- a series of small programs that each
produce a specific exception. Let the student run each one, read the
error, and discuss why it occurred before seeing the solution. before seeing the solution.
\sidenote{A crash collection serves double duty: it teaches exception recognition \emph{and} builds screen reader navigation fluency. Each crash is an opportunity to practice the bottom-up error reading protocol: terminal, bottom line first, error type, line number, Ctrl+G to navigate.} The goal
is pattern recognition, not memorization.

\paragraph{Thought Experiments}\label{thought-experiments-15}

\begin{enumerate}
\def\labelenumi{\arabic{enumi}.}
\item
  \emph{``Files are persistent storage --- they survive when your
  program exits. Variables are temporary --- they're gone when the
  program ends. What determines which data should be stored in files
  vs.~kept only in memory? What are the tradeoffs?''}
\item
  \emph{``The \texttt{with} statement uses Python's `context manager'
  protocol. This protocol says: `when this block starts, do X; when it
  ends (even if it ends with an error), do Y.' What other situations in
  programming require guaranteed cleanup? What happens in real systems
  when cleanup is skipped?''}
\item
  \emph{``Exceptions propagate up the call stack --- if function A calls
  B which calls C, and C raises an exception, it moves up to B, then A,
  until something catches it. If nothing catches it, the program
  crashes. Design the exception handling strategy for a program that
  reads a config file, parses it, and then connects to a database. Where
  should exceptions be caught?''}
\item
  \emph{``JSON stores data as text. Binary formats (like \texttt{pickle}
  in Python) store data in a compact binary form. JSON is
  human-readable; binary is not. JSON is language-independent; pickle is
  Python-only. For each of these use cases, which format would you
  choose: (a) a config file the user might edit by hand, (b) a large
  dataset processed only by your own Python code, (c) data shared
  between a Python backend and a JavaScript frontend?''}
\end{enumerate}

\begin{center}\rule{0.5\linewidth}{0.5pt}\end{center}

\subsection{14. Unit 9 --- Testing and Data
Analysis}\label{unit-9-testing-and-data-analysis}

\subsection{Extended Instructional
Cues}\label{extended-instructional-cues-2}

\textbf{Tests as documentation --- make this concrete:} Show two
identical functions --- one with tests that describe the expected
behavior, one without. one without.
Ask: ``Which function is easier to understand
\marginnote{For VI students, tests-as-documentation has a concrete accessibility payoff: a well-written test suite is a specification written in plain text that can be read aloud. When a student is handed unfamiliar code, running the tests and reading their names (\texttt{test\_add\_negative\_numbers}, \texttt{test\_empty\_list}) is a faster way to understand the code's intent than reading the implementation.}
without reading the implementation?'' The tests are a specification in
executable form.

introduce TDD: write the test before the implementation.
\sidenote{TDD has an unexpected accessibility benefit: it forces the programmer to specify behavior in text before implementing it, creating an audio-friendly specification document. The test file itself serves as a checklist of requirements.}
This inverts the usual workflow:

\begin{enumerate}
\def\labelenumi{\arabic{enumi}.}
\tightlist
\item
  Write a test that fails (because the function doesn't exist yet).
\item
  Write the minimum implementation to make the test pass.
\item
  Refactor.
\end{enumerate}

This discipline is genuinely useful for beginners because it forces them
to think about what a function should do before worrying about how.

\paragraph{Thought Experiments}\label{thought-experiments-16}

\begin{enumerate}
\def\labelenumi{\arabic{enumi}.}
\item
  \emph{``Unit testing tests one function in isolation. Integration
  testing tests multiple components working together. System testing
  tests the entire program. Why have three levels rather than just one?
  What does each level catch that the others might miss?''}
\item
  \emph{``A test can only prove that a program behaves correctly for the
  inputs tested. It cannot prove correctness for all possible inputs.
  What does this mean for the value of testing? How do you decide which
  inputs to test?''}
\item
  \emph{``Random numbers make testing difficult --- if the output
  changes every run, how do you write an assertion? Describe a strategy
  for testing a function that uses \texttt{random.randint()}. (Hint:
  think about seeding.)''}
\item
  \emph{``CSV data from the real world is often `dirty' --- missing
  values, extra commas, inconsistent formatting, encoding issues. What
  does a robust CSV parser need to handle? What assumptions would a
  naive parser make that would cause it to fail on real-world data?''}
\end{enumerate}

\begin{center}\rule{0.5\linewidth}{0.5pt}\end{center}

\subsection{15. Unit 10 --- Project-Based
Learning}\label{unit-10-project-based-learning}

\subsection{Extended Instructional
Cues}\label{extended-instructional-cues-3}

\textbf{Design document requirements:} Require a written design document
before any code is written. For a screen reader user, a plain text file
works best.
\marginnote{Plain \texttt{.txt} design documents are more accessible than Word, Google Docs, or Markdown-rendered HTML. The student can keep the design doc open in a second terminal pane and navigate between it and their code file with \texttt{Ctrl+`} and \texttt{Ctrl+Shift+`}.} The document should include:

\begin{verbatim}
Project Name: [title]
One-sentence purpose: [what problem does this solve?]
Target user: [who will use this?]

Features (ranked by priority):
1. [most important feature]
2. ...

Data I will store: [what information needs to persist?]
Data structures: [which Python structures will hold this?]

Functions I will write:
- [function name](inputs) -> output: [brief description]
- ...

Main loop pseudocode:
[plain English description of program flow]

Edge cases I have considered:
- [what happens if the user enters nothing?]
- [what happens if the file doesn't exist?]
- ...
\end{verbatim}

\textbf{Iterative development enforcement:} Require demonstrating a
working version at the end of every session, no matter how minimal.
``Working'' means runnable and producing correct output for at least one
input. A student who cannot demonstrate working code at session end
needs to reduce scope, not work harder overnight.
\sidenote{The ``working at every session'' requirement has an accessibility motivation: a non-working program is invisible to a screen reader. A student who has written fifty lines but none of them run has nothing to navigate, nothing to test, and no audio feedback. Iterative development---even one working function per session---provides continuous audio checkpoints.}

\paragraph{Thought Experiments}\label{thought-experiments-17}

\begin{enumerate}
\def\labelenumi{\arabic{enumi}.}
\item
  \emph{``A program designed for sighted users might use color to
  indicate status (red for error, green for success), icons to represent
  actions, and layout to group related information. How would you
  replicate each of these design patterns in a screen-reader-accessible
  terminal program?''}
\item
  \emph{``Your contact manager has a search function. How many different
  search behaviors could be useful? Search by exact name match, partial
  match, phone prefix, email domain, notes keywords? What should the
  program do when there are 0 results? 1 result? Multiple results?''}
\item
  \emph{``Error messages should be actionable --- they should tell the
  user not just what went wrong, but what to do next. Rewrite these
  unhelpful messages to be actionable: `Error: invalid input', `File not
  found', `Operation failed.'\,''}
\end{enumerate}

\begin{center}\rule{0.5\linewidth}{0.5pt}\end{center}

\subsection{16. Unit 11 --- Capstone
Project}\label{unit-11-capstone-project}

\subsection{Assessment Guidance}\label{assessment-guidance}

\paragraph{Process Checkpoints}\label{process-checkpoints}

\begin{fullwidth}{\def\LTcaptype{none} % do not increment counter
\begin{longtable}[]{@{}
  >{\raggedright\arraybackslash}p{(\linewidth - 4\tabcolsep) * \real{0.1917}}
  >{\raggedright\arraybackslash}p{(\linewidth - 4\tabcolsep) * \real{0.1250}}
  >{\raggedright\arraybackslash}p{(\linewidth - 4\tabcolsep) * \real{0.6833}}@{}}
\toprule\noalign{}
\begin{minipage}[b]{\linewidth}\raggedright
Checkpoint
\end{minipage} & \begin{minipage}[b]{\linewidth}\raggedright
Timing
\end{minipage} & \begin{minipage}[b]{\linewidth}\raggedright
What to Evaluate
\end{minipage} \\
\midrule\noalign{}
\endhead
\bottomrule\noalign{}
\endlastfoot
Written design document & Before any code & Is the scope realistic? Are
features prioritized? Are data structures appropriate? \\
Data model working & End of week 1 & Can items be created, stored in
memory, and printed? \\
Core CRUD functional & End of week 2 & Add, view, edit, delete all
work? \\
Persistence working & End of week 3 & Data survives program restart? \\
Error handling complete & End of week 4 & Program handles all invalid
inputs without crashing? \\
Is all output meaningful when
read aloud? \\
Final demo & End & Full demonstration with screen reader active \\
\end{longtable}
}
\end{fullwidth}
\marginnote{The accessibility review checkpoint is not optional or cosmetic. A program whose output is inaccessible when read aloud has failed its primary user requirement, regardless of whether it compiles and runs correctly.}

\paragraph{The Screen Reader Demo
Rubric}\label{the-screen-reader-demo-rubric}


\marginnote{Requiring the screen reader to be \emph{active} during the final demo---not just available---ensures the student has tested their own program under real conditions. Many accessibility issues only surface when the screen reader is actually speaking the output.} Evaluate each:

\begin{fullwidth}{\def\LTcaptype{none} % do not increment counter
\begin{longtable}[]{@{}
  >{\raggedright\arraybackslash}p{(\linewidth - 6\tabcolsep) * \real{0.1316}}
  >{\raggedright\arraybackslash}p{(\linewidth - 6\tabcolsep) * \real{0.3474}}
  >{\raggedright\arraybackslash}p{(\linewidth - 6\tabcolsep) * \real{0.2737}}
  >{\raggedright\arraybackslash}p{(\linewidth - 6\tabcolsep) * \real{0.2474}}@{}}
\toprule\noalign{}
\begin{minipage}[b]{\linewidth}\raggedright
Criterion
\end{minipage} & \begin{minipage}[b]{\linewidth}\raggedright
Excellent
\end{minipage} & \begin{minipage}[b]{\linewidth}\raggedright
Satisfactory
\end{minipage} & \begin{minipage}[b]{\linewidth}\raggedright
Needs Work
\end{minipage} \\
\midrule\noalign{}
\endhead
\bottomrule\noalign{}
\endlastfoot
\textbf{Menu navigation} & Student navigates menu by listening; confirms
choice before acting & Student navigates but occasionally needs visual
cues & Student needs visual confirmation of menu state \\
\textbf{Input prompts} & All prompts clearly describe expected format
and type & Most prompts are descriptive & Some prompts are vague or
missing \\
\textbf{Confirmation messages} & Every action produces a clear spoken
confirmation & Most actions confirmed & Some actions succeed silently \\
\textbf{Error messages} & All error messages are specific and actionable
& Most errors give useful messages & Some errors produce generic or
missing messages \\
\textbf{Data persistence} & Student demonstrates save and load by
quitting and restarting & Save/load works but student doesn't
demonstrate live & Save or load not fully functional \\
\textbf{Code explanation} & Student can explain any line of their code
when asked & Student can explain most of their code & Student cannot
explain sections of their code \\
\end{longtable}
}
\end{fullwidth}

\begin{center}\rule{0.5\linewidth}{0.5pt}\end{center}

\subsection{17. Assessment Strategies}\label{assessment-strategies}
\marginnote{Assessment for VI students requires rethinking the standard model. Timed written tests disadvantage students who navigate code via keyboard. Oral defenses, portfolio reviews, and screen-reader-active demonstrations are more equitable and reveal deeper understanding.}

\subsection{Accessible Assessment
Design}\label{accessible-assessment-design}

Traditional programming assessments (written tests, timed coding) create
unnecessary barriers for blind students. Use alternatives that assess
the same competencies:

\textbf{Think-aloud problems:} Give the student a short piece of code
and ask them to read it aloud and explain what it does, step by step.
This assesses reading comprehension and debugging ability without
requiring written answers.

\textbf{Debugging exercises:} Provide a working program and a broken
version. Ask the student to identify the error using only their screen
reader and debugging tools. This directly tests real-world skill.

\textbf{Feature extension:} Give a working program and ask the student
to add a specific feature. This is a more realistic assessment than
writing from scratch and allows scaffolding.

\textbf{Verbal design review:} Present a programming problem and ask the
student to describe their approach --- data structures, functions,
algorithm --- without writing any code. Assess reasoning quality.

\textbf{Portfolio review:} Maintain a collection of all programs written
during the course. Assess growth, consistency, commenting quality, and
the student's ability to explain their own work.

\subsection{Grading Considerations}\label{grading-considerations}

\begin{itemize}
\tightlist
\item
  Do not penalize for missing visual formatting (e.g., a missing newline
  that only matters visually).
\item
  Do reward for accessibility features the student adds above
  requirements: clear prompts, confirmation messages, help text.
\item
  Recognize that a blind student may spend significantly more time
  debugging than a sighted peer on identical problems. Factor this into
  time expectations.
\item
  Consider allowing oral explanations as alternatives to written code
  comments for students who find comment-writing slow.
\end{itemize}

\begin{center}\rule{0.5\linewidth}{0.5pt}\end{center}

\subsection{18. Supporting Diverse Needs Within the Visually Impaired
Population}\label{supporting-diverse-needs-within-the-visually-impaired-population}

\subsection{For Students with Low
Vision}\label{for-students-with-low-vision}

\begin{itemize}
\tightlist
\item
  Use high-contrast VS Code themes: Dark High Contrast, Light High
  Contrast (available under Color Themes in settings).
\item
  Increase editor font size: \texttt{Ctrl+Shift+P} -\textgreater{}
  ``Preferences: Open Settings'' -\textgreater{}
  \texttt{editor.fontSize}.
\item
  Enable zoom magnification: \texttt{Ctrl++} to zoom the VS Code window.
\item
  Test whether syntax highlighting is visible with their specific visual
  condition. Some students benefit from; others are confused by color
  information they cannot fully perceive.
\end{itemize}

\subsection{For Students Who Are
Deaf-Blind}\label{for-students-who-are-deaf-blind}

\begin{itemize}
\tightlist
\item
  All communication with the teacher must be through the braille display
  or tactile signing --- no spoken instruction.
\item
  Screen reader audio verbosity settings must be managed by the student
  (since they cannot hear audio feedback during class).
\item
  Consider pairing with a trained interpreter who knows both tactile
  communication and computing concepts.
\item
  Review the Perkins School for the Blind's resources on deaf-blind
  education technology.
\end{itemize}

\subsection{For Students with Physical Disabilities in Addition to
Visual
Impairment}\label{for-students-with-physical-disabilities-in-addition-to-visual-impairment}

\begin{itemize}
\tightlist
\item
  Switch access: VS Code is fully keyboard-navigable and can be operated
  with adapted keyboards, switch interfaces, or eye-gaze systems paired
  with virtual keyboards.
\item
  Voice control: Dragon NaturallySpeaking and macOS Voice Control can
  both operate VS Code. Voice coding requires additional voice commands
  for symbols.
\item
  Sip-and-puff devices: Can control keyboard input. Students using these
  will have different typing rates --- extend all time-based tasks.
\end{itemize}

\subsection{For Students with Cognitive Differences in Addition to
Visual
Impairment}\label{for-students-with-cognitive-differences-in-addition-to-visual-impairment}

\begin{itemize}
\tightlist
\item
  Reduce the amount of new vocabulary per session.
\item
  Use consistent metaphors across the entire course --- do not introduce
  a new metaphor for the same concept.
\item
  Provide written summaries of every lesson that the student can review
  with their screen reader later.
\item
  Break complex concepts into smaller steps; complete each step before
  introducing the next.
\item
  Build in regular review sessions rather than only forward-moving
  instruction.
\end{itemize}

\begin{center}\rule{0.5\linewidth}{0.5pt}\end{center}


\sidenote{AI code assistants (GitHub Copilot, Claude, ChatGPT) can generate code that \emph{appears} correct but contains subtle logic errors or accessibility oversights. For VI students learning to code, the risk of uncritically accepting AI output is high. Teach a ``trust but verify'' workflow: read every line the AI generates before running it.}

\subsection{Framing AI for the
Classroom}\label{framing-ai-for-the-classroom}

Establish a clear classroom policy before students begin using AI tools:

\textbf{Permitted:} Using AI to explain error messages, generate a first
draft for the student to critique and improve, explain a concept a
different way, suggest test cases to consider.

\textbf{Not permitted:} Submitting AI-generated code as completed work
without reading and understanding every line, using AI to complete
assessments that are meant to evaluate individual skill.

\textbf{Required when AI is used:} The student must be able to explain
any AI-generated code they submit as if they wrote it themselves.
Inability to explain code is evidence it was not understood.

\subsection{AI Tools and Blind Students --- Specific
Considerations}\label{ai-tools-and-blind-students-specific-considerations}

\textbf{Potential accessibility benefit:} AI tools like GitHub Copilot
are text-based --- suggestions appear as text that a screen reader can
read. Copilot's Accessible View integration in VS Code is improving.

\textbf{Potential accessibility risk:} AI may generate code that is
visually clear but auditorily confusing --- complex nested structures,
long lines, or dense comprehensions that are harder to navigate with a
screen reader than to read with eyes.

\textbf{Teaching point:} Ask students to evaluate AI-generated code for
screen-reader accessibility:

\begin{itemize}
\tightlist
\item
  Are variable names descriptive?
\item
  Are long lines split into readable pieces?
\item
  Are print statements and prompts clear and complete?
\item
  Does the code have docstrings?
\end{itemize}

If not, the AI-generated code needs to be modified before use ---
regardless of whether it runs correctly.

\subsection{Thought Experiments on
AI}\label{thought-experiments-on-ai}

\begin{enumerate}
\def\labelenumi{\arabic{enumi}.}
\item
  \emph{``An AI writes code that solves your problem but you don't
  understand it. You submit it. It passes tests. Have you learned
  anything? What is the difference between `having a solution' and
  `knowing how to solve problems'?''}
\item
  \emph{``AI code generators are trained on billions of lines of code
  written mostly by sighted programmers. What biases might this
  introduce in the code they generate? Think specifically about:
  variable names, output design, prompts, error messages,
  documentation.''}
\item
  \emph{``If AI can write most code, which skills become more valuable
  for future programmers? Consider: understanding requirements,
  debugging, testing, communicating about code, evaluating solutions,
  and accessibility design.''}
\end{enumerate}

\begin{center}\rule{0.5\linewidth}{0.5pt}\end{center}

\subsection{20. Alternative and Supplementary
Resources}\label{alternative-and-supplementary-resources}

\subsection{Textbooks and Reading}\label{textbooks-and-reading}

\begin{fullwidth}{\def\LTcaptype{none} % do not increment counter
\begin{longtable}[]{@{}
  >{\raggedright\arraybackslash}p{(\linewidth - 6\tabcolsep) * \real{0.3706}}
  >{\raggedright\arraybackslash}p{(\linewidth - 6\tabcolsep) * \real{0.1538}}
  >{\raggedright\arraybackslash}p{(\linewidth - 6\tabcolsep) * \real{0.1189}}
  >{\raggedright\arraybackslash}p{(\linewidth - 6\tabcolsep) * \real{0.3566}}@{}}
\toprule\noalign{}
\begin{minipage}[b]{\linewidth}\raggedright
Resource
\end{minipage} & \begin{minipage}[b]{\linewidth}\raggedright
Level
\end{minipage} & \begin{minipage}[b]{\linewidth}\raggedright
Format
\end{minipage} & \begin{minipage}[b]{\linewidth}\raggedright
Accessibility
\end{minipage} \\
\midrule\noalign{}
\endhead
\bottomrule\noalign{}
\endlastfoot
\textbf{Python Crash Course, 3rd Ed.} (Matthes) & Beginner & Print/ebook
& Good; primary textbook for this curriculum \\
\textbf{Python Programming: An Introduction to CS} (Zelle) & Beginner &
Print/ebook & Good; companion textbook \\
\textbf{Automate the Boring Stuff with Python} (Sweigart) &
Beginner--Intermediate & Free HTML / print & Excellent; HTML version
fully accessible \\
\textbf{Think Python} (Downey) & Beginner & Free HTML / print &
Excellent HTML version \\
\textbf{Learn Python 3 the Hard Way} (Shaw) & Beginner & Print/ebook &
Plain text exercises, highly accessible \\
\textbf{Python for Everybody} (Severance) & Beginner & Free HTML + video
& Video lectures have transcripts; good accessibility \\
\textbf{Fluent Python} (Ramalho) & Advanced & Print/ebook & Dense but
thorough; for motivated advanced students \\
\textbf{Python Docs} (python.org/3) & Reference & Free HTML & Excellent
screen-reader accessibility; definitive \\
\end{longtable}
}
\end{fullwidth}

\subsection{Online Courses and
Platforms}\label{online-courses-and-platforms}

\begin{fullwidth}{\def\LTcaptype{none} % do not increment counter
\begin{longtable}[]{@{}
  >{\raggedright\arraybackslash}p{(\linewidth - 6\tabcolsep) * \real{0.1524}}
  >{\raggedright\arraybackslash}p{(\linewidth - 6\tabcolsep) * \real{0.2378}}
  >{\raggedright\arraybackslash}p{(\linewidth - 6\tabcolsep) * \real{0.3293}}
  >{\raggedright\arraybackslash}p{(\linewidth - 6\tabcolsep) * \real{0.2805}}@{}}
\toprule\noalign{}
\begin{minipage}[b]{\linewidth}\raggedright
Platform
\end{minipage} & \begin{minipage}[b]{\linewidth}\raggedright
Course
\end{minipage} & \begin{minipage}[b]{\linewidth}\raggedright
Screen Reader Compatibility
\end{minipage} & \begin{minipage}[b]{\linewidth}\raggedright
Notes
\end{minipage} \\
\midrule\noalign{}
\endhead
\bottomrule\noalign{}
\endlastfoot
\textbf{edX / CS50P} (Harvard) & Introduction to Programming with Python
& High --- video transcripts, accessible exercises & Free audit. Very
high quality. \\
\textbf{Coursera} & Python for Everybody (Michigan) & Medium --- some
interactive elements inaccessible & Free audit available \\
\textbf{Codecademy} & Learn Python 3 & Partial --- editor accessible,
some UI not & Requires testing with specific screen reader \\
\textbf{SoloLearn} & Python Core & High --- mobile app works well with
VoiceOver/TalkBack & Good for mobile-first learners \\
\textbf{freeCodeCamp} & Python curriculum & Medium --- varies by module
& Free; large community \\
Excellent for structured practice after basics
\marginnote{Exercism is particularly well-suited for blind students: all exercises are solved locally in VS Code with a CLI tool, test-driven, and submitted via terminal. No web-based code editor required.} \\
\textbf{CodingBat} & Python problems & High --- plain HTML,
browser-based & Good for logic practice; no complex UI \\
\textbf{replit.com} & Any Python project & Medium --- test with
student's setup & No-install option; accessibility variable \\
\end{longtable}
}
\end{fullwidth}

\subsection{Video Resources}\label{video-resources}

\begin{fullwidth}{\def\LTcaptype{none} % do not increment counter
\begin{longtable}[]{@{}
  >{\raggedright\arraybackslash}p{(\linewidth - 4\tabcolsep) * \real{0.2993}}
  >{\raggedright\arraybackslash}p{(\linewidth - 4\tabcolsep) * \real{0.3061}}
  >{\raggedright\arraybackslash}p{(\linewidth - 4\tabcolsep) * \real{0.3946}}@{}}
\toprule\noalign{}
\begin{minipage}[b]{\linewidth}\raggedright
Resource
\end{minipage} & \begin{minipage}[b]{\linewidth}\raggedright
Description
\end{minipage} & \begin{minipage}[b]{\linewidth}\raggedright
Accessibility
\end{minipage} \\
\midrule\noalign{}
\endhead
\bottomrule\noalign{}
\endlastfoot
\textbf{CS50P (Harvard)} & Full course lectures, free on YouTube and edX
& Closed captions, transcripts; very high production quality \\
\textbf{Corey Schafer Python tutorials} (YouTube) & In-depth concept
tutorials & Closed captions available \\
\textbf{Tech With Tim Python} (YouTube) & Project-based tutorials &
Closed captions available \\
\textbf{Sentdex Python tutorials} (YouTube) & Large series from beginner
to advanced & Captions available \\
\textbf{Python.org Getting Started} & Short official tutorials & Plain
text, no video \\
\end{longtable}
}
\end{fullwidth}

\subsection{Tools Specifically for Accessibility in Programming
Education}\label{tools-specifically-for-accessibility-in-programming-education}

\begin{fullwidth}{\def\LTcaptype{none} % do not increment counter
\begin{longtable}[]{@{}
  >{\raggedright\arraybackslash}p{(\linewidth - 4\tabcolsep) * \real{0.2329}}
  >{\raggedright\arraybackslash}p{(\linewidth - 4\tabcolsep) * \real{0.3151}}
  >{\raggedright\arraybackslash}p{(\linewidth - 4\tabcolsep) * \real{0.4521}}@{}}
\toprule\noalign{}
\begin{minipage}[b]{\linewidth}\raggedright
Tool
\end{minipage} & \begin{minipage}[b]{\linewidth}\raggedright
Purpose
\end{minipage} & \begin{minipage}[b]{\linewidth}\raggedright
Notes
\end{minipage} \\
\midrule\noalign{}
\endhead
\bottomrule\noalign{}
\endlastfoot
\textbf{Python Tutor} (pythontutor.com) & Step-by-step code execution
with state display & Text-based state panel is screen-reader
accessible \\
\textbf{Mu Editor} & Beginner-focused Python IDE & Simpler than VS Code;
reasonable screen reader support \\
\textbf{Thonny} & Beginner IDE with step debugger & Variable viewer is
navigable via keyboard \\
\textbf{NVDA + Remote} & Remote screen sharing for teacher support &
Student's screen reader on their machine; teacher can see output \\
\textbf{GitHub Codespaces} & Cloud VS Code environment & Same
accessibility as local VS Code; useful when local setup fails \\
\textbf{Braille Institute Resources} & General adaptive technology
guidance & brailleinstitute.org \\
\textbf{APH (American Printing House)} & Educational materials for blind
students & aph.org \\
\end{longtable}
}
\end{fullwidth}

\subsection{Communities and Professional
Organizations}\label{communities-and-professional-organizations}

\begin{fullwidth}{\def\LTcaptype{none} % do not increment counter
\begin{longtable}[]{@{}
  >{\raggedright\arraybackslash}p{(\linewidth - 2\tabcolsep) * \real{0.5168}}
  >{\raggedright\arraybackslash}p{(\linewidth - 2\tabcolsep) * \real{0.4832}}@{}}
\toprule\noalign{}
\begin{minipage}[b]{\linewidth}\raggedright
Organization
\end{minipage} & \begin{minipage}[b]{\linewidth}\raggedright
Relevance
\end{minipage} \\
\midrule\noalign{}
\endhead
\bottomrule\noalign{}
\endlastfoot
\textbf{National Federation of the Blind (NFB) --- Science \&
Engineering Division} & Blind scientists and engineers; mentorship
opportunities \\
\textbf{American Council of the Blind (ACB) --- Technology Interest
Group} & Accessible technology advocacy and resources \\
\textbf{CSUN Assistive Technology Conference} & Annual conference; many
papers on computing education for blind students \\
\textbf{SIGACCESS} (ACM SIG on Computers and Accessibility) & Academic
research on accessible computing \\
\textbf{Ability Central} & Technology resources for people with
disabilities \\
\textbf{Hadley Institute for the Blind} & Professional development for
educators of blind students \\
\textbf{Perkins School for the Blind} & Resources including STEM
education for blind students \\
\textbf{Carroll Center for the Blind} & Technology training for blind
adults; may have curriculum resources \\
\textbf{Stack Overflow Blind/VI Community} & Informal community; blind
developers helping each other \\
\end{longtable}
}
\end{fullwidth}

\begin{center}\rule{0.5\linewidth}{0.5pt}\end{center}

\subsection{21. Quick Reference Sheets}\label{quick-reference-sheets}

\subsection{Error Message Reading
Protocol}\label{error-message-reading-protocol}

When an error occurs, guide the student through this sequence:

\marginnote{The bottom-up reading rule is the single most important}\amn{\textbf{Accessibility Tip:} When an error occurs, navigate to the \emph{bottom} of the terminal output first---the last line of a Python traceback contains the most specific error message. Use \texttt{Ctrl+End} to jump to the end of the terminal buffer, then read upward to the line number.}
\marginnote{The bottom-up reading rule is the single most important error-debugging habit for screen reader users. Python tracebacks can be many lines long; the first line heard is always the least specific. Jump to the bottom first.}

\begin{verbatim}
1. Navigate to terminal output (Ctrl+` or click panel)
2. Read from the BOTTOM of the error output upward
3. Note the ERROR TYPE (last line): SyntaxError, NameError, TypeError, etc.
4. Note the ERROR MESSAGE: specific description of what went wrong
5. Note the LINE NUMBER in the traceback
6. Jump to that line with Ctrl+G -> [line number] -> Enter
7. Read the line and the lines immediately before and after it
8. Categorize the error (see table below)
9. Fix and re-run
\end{verbatim}

\subsection{Error Type Quick
Reference}\label{error-type-quick-reference}

\begin{fullwidth}{\def\LTcaptype{none} % do not increment counter
\begin{longtable}[]{@{}
  >{\raggedright\arraybackslash}p{(\linewidth - 4\tabcolsep) * \real{0.1387}}
  >{\raggedright\arraybackslash}p{(\linewidth - 4\tabcolsep) * \real{0.3650}}
  >{\raggedright\arraybackslash}p{(\linewidth - 4\tabcolsep) * \real{0.4964}}@{}}
\toprule\noalign{}
\begin{minipage}[b]{\linewidth}\raggedright
Error Type
\end{minipage} & \begin{minipage}[b]{\linewidth}\raggedright
Common Cause
\end{minipage} & \begin{minipage}[b]{\linewidth}\raggedright
First Thing to Check
\end{minipage} \\
\midrule\noalign{}
\endhead
\bottomrule\noalign{}
\endlastfoot
\texttt{SyntaxError} & Missing \texttt{:}, \texttt{)}, \texttt{"},
\texttt{{]}}, or \texttt{=} instead of \texttt{==} & The exact column
number in the error message \\
\texttt{IndentationError} & Inconsistent spaces/tabs & Line indentation;
use \texttt{cat\ -A} or \texttt{show\ whitespace} in VS Code \\
\texttt{NameError} & Variable not defined, or misspelled & Check
spelling; check if variable was defined before use \\
\texttt{TypeError} & Wrong data type & Use \texttt{type()} on each
variable involved in the operation \\
\texttt{ValueError} & Right type, wrong value & Check what you're
converting: \texttt{int(\textquotesingle{}hello\textquotesingle{})}
fails; \texttt{int(\textquotesingle{}5\textquotesingle{})} works \\
\texttt{IndexError} & List index out of range & Print \texttt{len(list)}
and check what index you're trying to use \\
\texttt{KeyError} & Dictionary key doesn't exist & Use \texttt{.get()}
or check with \texttt{if\ key\ in\ dict} first \\
\texttt{FileNotFoundError} & File doesn't exist at that path & Print
\texttt{os.getcwd()} and verify the file path \\
\texttt{AttributeError} & Called a method that doesn't exist & Check
spelling and whether the object is the type you think it is \\
\texttt{ZeroDivisionError} & Dividing by zero & Add a check before the
division: \texttt{if\ denominator\ !=\ 0} \\
\texttt{RecursionError} & Infinite recursion & Verify base case exists
and is reachable \\
\end{longtable}
}
\end{fullwidth}

\subsection{Accessible Output
Patterns}\label{accessible-output-patterns}

\amn{\textbf{Accessible output rule:} Never print a raw list}\amn{\textbf{Accessibility Tip:} When printing any collection, use \texttt{enumerate()} to number each item and print one per line. Numbered, labelled output is far easier to navigate with a screen reader than bracket-notation list output. Make this a habit from the first list lesson.}
\amn{\textbf{Accessible output rule:} Never print a raw list or dict. Always iterate with \texttt{enumerate} and label each item. Raw collection output (\texttt{[item1, item2, ...]}) is structurally noisy when read aloud by a screen reader.}

\begin{Shaded}
\begin{Highlighting}[]
print(\StringTok{"Items in collection:"})
\KeywordTok{for} i, item \KeywordTok{in} enumerate(my\_list, \CodeNumTok{1}):
    print(\StringTok{f"  }\StringTok{\{}i\StringTok{\}}\StringTok{. }\StringTok{\{}item\StringTok{\}}\StringTok{"})
print(\StringTok{f"Total: }\StringTok{\{}len(my\_list)\StringTok{\}}\StringTok{ items"})
 
\CommentTok{\# Instead of: print(my\_dict)}
\CommentTok{\# Use:}
print(\StringTok{"Current settings:"})
\KeywordTok{for} key, value \KeywordTok{in} my\_dict.items():
    print(\StringTok{f"  }\StringTok{\{}key\StringTok{\}}\StringTok{: }\StringTok{\{}value\StringTok{\}}\StringTok{"})
 
\CommentTok{\# Menu display:}
print(\StringTok{"}\StringTok{\textbackslash{}n}\StringTok{{-}{-}{-} Main Menu {-}{-}{-}"})
print(\StringTok{"1. Add item"})
print(\StringTok{"2. View all items"})
print(\StringTok{"3. Search"})
print(\StringTok{"4. Delete item"})
print(\StringTok{"5. Save and quit"})
print(\StringTok{"{-}{-}{-}{-}{-}{-}{-}{-}{-}{-}{-}{-}{-}{-}{-}{-}{-}"})
choice \OperatorTok{=} input(\StringTok{"Enter option number (1{-}5): "})
 
\CommentTok{\# Confirmation messages:}
print(\StringTok{f"Contact \textquotesingle{}}\StringTok{\{}name\StringTok{\}}\StringTok{\textquotesingle{} added successfully."})
print(\StringTok{f"Found }\StringTok{\{}len(results)\StringTok{\}}\StringTok{ matching records."})
print(\StringTok{"Data saved to contacts.json."})
print(\StringTok{"No matching records found."})
print(\StringTok{"Error: Name field cannot be empty. Please try again."})
\end{Highlighting}
\end{Shaded}

\subsection{VS Code Keyboard Shortcut Reference
Card}\label{vs-code-keyboard-shortcut-reference-card}

\begin{verbatim}
NAVIGATION
  Ctrl+G          Go to specific line number
  Ctrl+Home       Jump to start of file
  Ctrl+End        Jump to end of file
  Ctrl+Shift+O    Go to function/class (symbol picker)
  F12             Go to Definition
  Alt+F12         Peek Definition
  Shift+F12       Find all References
  Ctrl+P          Quick open file by name

ERRORS AND DIAGNOSTICS
  F8              Next error/warning
  Shift+F8        Previous error/warning
  Ctrl+Shift+M    Open Problems panel
  Shift+Alt+F2    Accessible View (for hover/tooltip content)

EDITING
  Shift+Alt+F     Format document (auto-indent)
  Ctrl+/          Toggle comment on selected lines
  Tab             Indent selected lines
  Shift+Tab       Outdent selected lines
  Ctrl+D          Select next occurrence of word
  Ctrl+H          Find and replace

TERMINAL
  Ctrl+`          Toggle integrated terminal
  Ctrl+Shift+`    New terminal
  Up/Down arrows  Navigate command history
  Tab             Auto-complete filenames
  Ctrl+C          Stop running program

SCREEN READER
  Ctrl+Shift+P    Command Palette
  Ctrl+Shift+P -> "Toggle Screen Reader Optimized"
  Shift+Alt+F1    Accessibility help
\end{verbatim}

\newpage

\section{Unit 0: Setting Up Your Accessible Python Development
Environment}\label{unit-0-setting-up-your-accessible-python-development-environment}

\subsection{Overview}\label{overview}

A detailed environment setup is not merely a technical prerequisite; it
empowers students and reduces dependency on external assistance. For
visually impaired students, relying on external help for every setup or
environment issue can be disempowering and slow down the learning
process. This unit provides extremely detailed, screen-reader-friendly
instructions for environment setup.

\textbf{Target Audience:} Students with no prior programming experience,
relying on screen readers for all computer interactions.

\begin{center}\rule{0.5\linewidth}{0.5pt}\end{center}

\subsection{Lesson 0.1: Installing Python and Visual Studio
Code}\label{lesson-0.1-installing-python-and-visual-studio-code}

\subsection{Reading References}\label{reading-references}

\begin{itemize}
\tightlist
\item
  Matthes: Chapter 1, pages 3--14
\item
  Matthes: Appendix A, pages 463--468
\end{itemize}

\subsection{Installing Python 3}\label{installing-python-3}

Python is not always installed on a computer by default. Download the
latest version from \url{https://www.python.org/downloads/}.

\paragraph{Windows}\label{windows}

\begin{enumerate}
\def\labelenumi{\arabic{enumi}.}
\tightlist
\item
  Locate the downloaded \texttt{.exe} file (usually in your Downloads
  folder). Activate it.
\item
  \textbf{CRUCIAL STEP:} When the installer window appears, check the
  \textbf{``Add Python to PATH''} checkbox. Use Tab to navigate to this
  checkbox and press Spacebar to check it.
\item
  Navigate to ``Install Now'' and activate it.
\item
  Follow any further prompts. Activate ``Close'' when finished.
\end{enumerate}

\paragraph{macOS}\label{macos}

\begin{enumerate}
\def\labelenumi{\arabic{enumi}.}
\tightlist
\item
  Locate the downloaded \texttt{.pkg} file. Activate it.
\item
  Activate ``Continue'' through the steps.
\item
  Agree to the license terms. Choose the installation location (default
  is fine).
\item
  Enter your administrator password when prompted.
\item
  Activate ``Install'' and then ``Close'' when finished.
\end{enumerate}

\paragraph{Linux}\label{linux}

Python 3 is often pre-installed. Verify in a terminal:

\begin{Shaded}
\begin{Highlighting}[]
python3 \CodeTypeTok{{-}{-}version}
\end{Highlighting}
\end{Shaded}

If not installed (Debian/Ubuntu):

\begin{Shaded}
\begin{Highlighting}[]
\CodeFuncTok{sudo} apt update
\CodeFuncTok{sudo} apt install python3
\CodeFuncTok{sudo} apt install python3{-}pip
\end{Highlighting}
\end{Shaded}

\subsection{Verifying Python
Installation}\label{verifying-python-installation}

Open a terminal and type:

\begin{Shaded}
\begin{Highlighting}[]
python3 \CodeTypeTok{{-}{-}version}
\end{Highlighting}
\end{Shaded}

You should hear the version number (e.g., ``Python 3.10.12''). Also
check pip:

\begin{Shaded}
\begin{Highlighting}[]
pip3 \CodeTypeTok{{-}{-}version}
\end{Highlighting}
\end{Shaded}

\subsection{Installing Visual Studio
Code}\label{installing-visual-studio-code}

\begin{enumerate}
\def\labelenumi{\arabic{enumi}.}
\tightlist
\item
  Navigate to \url{https://code.visualstudio.com/download}.
\item
  Download the installer for your OS.
\end{enumerate}

\paragraph{Windows}\label{windows-1}

\begin{itemize}
\tightlist
\item
  \textbf{CRUCIAL STEP:} On the ``Select Additional Tasks'' screen,
  check \textbf{``Add to PATH''}.
\item
  Activate ``Next'' and then ``Install''.
\end{itemize}

\paragraph{macOS}\label{macos-1}

\begin{itemize}
\tightlist
\item
  Drag the ``Visual Studio Code'' application to your ``Applications''
  folder.
\end{itemize}

\paragraph{Linux (Debian/Ubuntu)}\label{linux-debianubuntu}

\begin{Shaded}
\begin{Highlighting}[]
cd Downloads
\CodeFuncTok{sudo} dpkg \CodeTypeTok{{-}i} code\_\CodeAccentTok{*}.deb
\end{Highlighting}
\end{Shaded}

\subsection{Verify VS Code}\label{verify-vs-code}

\begin{Shaded}
\begin{Highlighting}[]
code \CodeTypeTok{{-}{-}version}
\end{Highlighting}
\end{Shaded}

\subsection{Skills Checklist}\label{skills-checklist}

\begin{itemize}
\tightlist
\item[$\square$]
  I can successfully download the Python 3 installer for my operating
  system.
\item[$\square$]
  I can run the Python installer and correctly select the ``Add Python
  to PATH'' option.
\item[$\square$]
  I can successfully download the VS Code installer for my operating
  system.
\item[$\square$]
  I can run the VS Code installer and correctly select the ``Add to
  PATH'' option.
\item[$\square$]
  I can open my computer's system terminal.
\item[$\square$]
  I can verify Python's installation by running
  \texttt{python3\ -\/-version} in the terminal.
\item[$\square$]
  I can verify VS Code's installation by running
  \texttt{code\ -\/-version} in the terminal.
\end{itemize}

\begin{center}\rule{0.5\linewidth}{0.5pt}\end{center}

\subsection{Lesson 0.2: Configuring Your Screen Reader for VS
Code}\label{lesson-0.2-configuring-your-screen-reader-for-vs-code-1}

\subsection{Reading References}\label{reading-references-1}

\begin{itemize}
\tightlist
\item
  Matthes: Appendix B, pages 469--476
\item
  Matthes: Appendix C, pages 477--482
\end{itemize}

\subsection{General Screen Reader Tips for VS
Code}\label{general-screen-reader-tips-for-vs-code}

\textbf{1. Use Screen Reader Mode}

\begin{itemize}
\tightlist
\item
  Press \texttt{Ctrl\ +\ Shift\ +\ P} to open the Command Palette.
\item
  Type ``Toggle Screen Reader Optimized'' and press Enter.
\end{itemize}

\textbf{2. Accessible View}

\begin{itemize}
\tightlist
\item
  Press \texttt{Shift\ +\ Alt\ +\ F2} (Windows/Linux) or
  \texttt{Shift\ +\ Option\ +\ F2} (macOS).
\item
  Opens a flat, text-based representation of complex UI elements.
\item
  Press \texttt{Esc} to close.
\end{itemize}

\textbf{3. Keyboard-First Navigation}

\begin{itemize}
\tightlist
\item
  \texttt{Tab} and \texttt{Shift+Tab}: Move focus between UI elements.
\item
  Arrow Keys: Navigate within active panes.
\end{itemize}

\textbf{4. Status Bar} The status bar at the bottom provides current
line/column number, file type, active Python environment, and
error/warning counts.

\subsection{Screen Reader Indentation
Configuration}\label{screen-reader-indentation-configuration}

Python relies heavily on indentation. Configure your screen reader to
announce indentation changes.

\paragraph{NVDA (Windows)}\label{nvda-windows}

\begin{enumerate}
\def\labelenumi{\arabic{enumi}.}
\tightlist
\item
  Press \texttt{NVDA\ Key\ +\ N} -\textgreater{} Preferences
  -\textgreater{} Settings.
\item
  Find ``Document Formatting''.
\item
  Locate ``Indentation'' and set it to \textbf{``Say all''} or
  \textbf{``Beeps''}.
\item
  Activate ``OK''.
\end{enumerate}

\paragraph{JAWS (Windows)}\label{jaws-windows}

\begin{enumerate}
\def\labelenumi{\arabic{enumi}.}
\tightlist
\item
  Press \texttt{JAWS\ Key\ +\ J} -\textgreater{} Alt+O for Options
  -\textgreater{} B for Basics.
\item
  Find ``Text Processing'' settings.
\item
  Configure indentation announcement.
\end{enumerate}

\paragraph{VoiceOver (macOS)}\label{voiceover-macos}

VoiceOver generally handles indentation well by default, announcing ``4
spaces'' or ``tab'' as you navigate.

\subsection{Skills Checklist}\label{skills-checklist-1}

\begin{itemize}
\tightlist
\item[$\square$]
  I can identify and toggle VS Code's ``Screen Reader Optimized Mode''.
\item[$\square$]
  I can open and close the ``Accessible View''.
\item[$\square$]
  I can navigate between VS Code panes using Tab and Shift+Tab.
\item[$\square$]
  I can navigate line by line and character by character within the
  editor.
\item[$\square$]
  I can locate and listen to information in the VS Code Status Bar.
\item[$\square$]
  I can configure my screen reader to announce indentation changes.
\end{itemize}

\begin{center}\rule{0.5\linewidth}{0.5pt}\end{center}

\subsection{Lesson 0.3: Setting Up Your VS Code
Workspace}\label{lesson-0.3-setting-up-your-vs-code-workspace-1}

\subsection{Reading References}\label{reading-references-2}

\begin{itemize}
\tightlist
\item
  Matthes: Appendix B, pages 469--476
\item
  Matthes: Chapter 8, pages 149--152
\end{itemize}

\subsection{Installing VS Code
Extensions}\label{installing-vs-code-extensions}

\begin{enumerate}
\def\labelenumi{\arabic{enumi}.}
\tightlist
\item
  Open the Extensions View: \texttt{Ctrl\ +\ Shift\ +\ X}
  (Windows/Linux) or \texttt{Command\ +\ Shift\ +\ X} (macOS).
\item
  Search and install \textbf{``Python''} by Microsoft --- the official
  and most important Python extension.
\item
  Search and install \textbf{``Pylance''} --- provides excellent
  IntelliSense and type checking.
\item
  Optional: \textbf{``Code Runner''} --- allows running code files with
  \texttt{Ctrl+Alt+N}.
\end{enumerate}

\subsection{Creating Your First Workspace
Folder}\label{creating-your-first-workspace-folder}

\textbf{Windows:} Open File Explorer (\texttt{Windows\ Key\ +\ E}).
Navigate to Documents. Press \texttt{Ctrl\ +\ Shift\ +\ N} to create a
new folder named \texttt{PythonProjects}.

\textbf{macOS:} Open Finder. Navigate to Documents. Press
\texttt{Command\ +\ Shift\ +\ N} to create \texttt{PythonProjects}.

\subsection{Opening the Folder in VS
Code}\label{opening-the-folder-in-vs-code}

\textbf{Method 1 (Recommended):} Open terminal, navigate to the folder,
then:

\begin{Shaded}
\begin{Highlighting}[]
code .
\end{Highlighting}
\end{Shaded}

\textbf{Method 2:} In VS Code, press \texttt{Ctrl\ +\ K,\ Ctrl\ +\ O} to
``Open Folder''.

\subsection{VS Code Navigation
Shortcuts}\label{vs-code-navigation-shortcuts}

\begin{fullwidth}{\def\LTcaptype{none} % do not increment counter
\begin{longtable}[]{@{}ll@{}}
\toprule\noalign{}
Area & Shortcut \\
\midrule\noalign{}
\endhead
\bottomrule\noalign{}
\endlastfoot
Explorer & \texttt{Ctrl\ +\ Shift\ +\ E} \\
Search & \texttt{Ctrl\ +\ Shift\ +\ F} \\
Extensions & \texttt{Ctrl\ +\ Shift\ +\ X} \\
Toggle Sidebar & \texttt{Ctrl\ +\ B} \\
Command Palette & \texttt{Ctrl\ +\ Shift\ +\ P} \\
Problems View & \texttt{Ctrl\ +\ Shift\ +\ M} \\
Next Error & \texttt{F8} \\
Previous Error & \texttt{Shift\ +\ F8} \\
Integrated Terminal & `Ctrl + `` \\
\end{longtable}
}
\end{fullwidth}

\subsection{Skills Checklist}\label{skills-checklist-2}

\begin{itemize}
\tightlist
\item[$\square$]
  I can open the Extensions View in VS Code.
\item[$\square$]
  I can search for and install the ``Python'' extension.
\item[$\square$]
  I can search for and install the ``Pylance'' extension.
\item[$\square$]
  I can create a new folder for Python projects.
\item[$\square$]
  I can open a specific folder as a workspace in VS Code.
\item[$\square$]
  I can navigate the Activity Bar, Explorer, Editor, and Panel using
  keyboard shortcuts.
\item[$\square$]
  I can effectively use the Command Palette
  (\texttt{Ctrl\ +\ Shift\ +\ P}).
\end{itemize}

\begin{center}\rule{0.5\linewidth}{0.5pt}\end{center}

\subsection{Lesson 0.4: Python Environments and Terminal
Usage}\label{lesson-0.4-python-environments-and-terminal-usage-1}

\subsection{Reading References}\label{reading-references-3}

\begin{itemize}
\tightlist
\item
  Matthes: Chapter 1, pages 3--14
\item
  Matthes: Chapter 8, pages 149--152
\item
  Matthes: Chapter 11, pages 210--211
\end{itemize}

\subsection{Creating Your First Python
File}\label{creating-your-first-python-file}

\begin{enumerate}
\def\labelenumi{\arabic{enumi}.}
\tightlist
\item
  Ensure Explorer is focused: \texttt{Ctrl\ +\ Shift\ +\ E}.
\item
  Create new file: \texttt{Ctrl\ +\ N}.
\item
  Save: \texttt{Ctrl\ +\ S}. Name it \texttt{hello.py}.
\item
  Type:
\end{enumerate}

\begin{Shaded}
\begin{Highlighting}[]
print(\StringTok{"Hello from VS Code!"})
\end{Highlighting}
\end{Shaded}

\subsection{Running Python Scripts}\label{running-python-scripts}

\textbf{From System Terminal:}

\begin{Shaded}
\begin{Highlighting}[]
cd \textasciitilde{}/Documents/PythonProjects
python3 hello.py
\end{Highlighting}
\end{Shaded}

\textbf{From VS Code Integrated Terminal:}

\begin{enumerate}
\def\labelenumi{\arabic{enumi}.}
\tightlist
\item
  Open terminal: `Ctrl + ``
\item
  The terminal opens in your project folder automatically.
\item
  Run: \texttt{python3\ hello.py}
\end{enumerate}

\subsection{Python Virtual
Environments}\label{python-virtual-environments}

Virtual environments keep project dependencies separate, preventing
conflicts.

\textbf{Create a virtual environment:}

\begin{Shaded}
\begin{Highlighting}[]
python3 \CodeTypeTok{{-}m} venv .venv
\end{Highlighting}
\end{Shaded}

\textbf{Activate:}

\begin{itemize}
\tightlist
\item
  Windows (Command Prompt):
  \texttt{.venv\textbackslash{}Scripts\textbackslash{}activate.bat}
\item
  Windows (PowerShell):
  \texttt{.venv\textbackslash{}Scripts\textbackslash{}Activate.ps1}
\item
  macOS/Linux: \texttt{source\ ./.venv/bin/activate}
\end{itemize}

You will hear \texttt{(.venv)} at the beginning of your prompt when
active.

\textbf{Deactivate:}

\begin{Shaded}
\begin{Highlighting}[]
deactivate
\end{Highlighting}
\end{Shaded}

\textbf{Select Interpreter in VS Code:}

\begin{enumerate}
\def\labelenumi{\arabic{enumi}.}
\tightlist
\item
  Press \texttt{Ctrl\ +\ Shift\ +\ P}.
\item
  Type ``Python: Select Interpreter'' and press Enter.
\item
  Select the one that says ``Python 3.x.x (.venv)''.
\end{enumerate}

\textbf{Install packages:}

\begin{Shaded}
\begin{Highlighting}[]
pip install requests
pip list
\end{Highlighting}
\end{Shaded}

\subsection{Skills Checklist}\label{skills-checklist-3}

\begin{itemize}
\tightlist
\item[$\square$]
  I can create a new Python file (.py) in VS Code.
\item[$\square$]
  I can write a simple \texttt{print()} statement.
\item[$\square$]
  I can run a Python script from the system terminal.
\item[$\square$]
  I can run a Python script from the integrated terminal in VS Code.
\item[$\square$]
  I understand why virtual environments are useful.
\item[$\square$]
  I can create, activate, and deactivate a virtual environment.
\item[$\square$]
  I can select the correct Python interpreter in VS Code.
\item[$\square$]
  I can install Python packages using pip.
\end{itemize}

\begin{center}\rule{0.5\linewidth}{0.5pt}\end{center}

\subsection{Lesson 0.5: General Accessibility Tools and
Strategies}\label{lesson-0.5-general-accessibility-tools-and-strategies}

\subsection{VS Code Accessibility
Features}\label{vs-code-accessibility-features}

\begin{itemize}
\tightlist
\item
  \textbf{Accessible View} (\texttt{Shift+Alt+F2}): Reveals hover
  information and error details.
\item
  \textbf{Focus on Errors} (\texttt{F8} / \texttt{Shift+F8}): Jump
  between errors without visual scanning.
\item
  \textbf{Go to Definition} (\texttt{F12}): Navigate to where a function
  or variable is defined.
\item
  \textbf{Go to Reference} (\texttt{Shift+F12}): Find all usages of a
  symbol.
\item
  \textbf{Go to Line/Column} (\texttt{Ctrl+G}): Jump to a specific line
  number.
\end{itemize}

\subsection{Python Documentation}\label{python-documentation}

The official Python documentation is well-structured and accessible:
\url{https://docs.python.org/3/}

\subsection{Unit 0 Completion
Checklist}\label{unit-0-completion-checklist}

\begin{itemize}
\tightlist
\item[$\square$]
  Python 3 is installed and \texttt{python3\ -\/-version} works in your
  system terminal.
\item[$\square$]
  VS Code is installed and \texttt{code\ -\/-version} works in your
  system terminal.
\item[$\square$]
  Your screen reader is configured to announce indentation changes in VS
  Code.
\item[$\square$]
  You can open your PythonProjects folder in VS Code.
\item[$\square$]
  You have installed the official ``Python'' VS Code extension.
\item[$\square$]
  You can create a new .py file, write \texttt{print("Hello!")}, and
  save it.
\item[$\square$]
  You can run your hello.py file from both terminals.
\item[$\square$]
  You can create and activate a virtual environment (.venv).
\item[$\square$]
  You can select the virtual environment interpreter in VS Code.
\item[$\square$]
  You can install a package using pip.
\item[$\square$]
  You know how to use \texttt{Ctrl\ +\ Shift\ +\ P} (Command Palette)
  and \texttt{F8} (Go to next error).
\end{itemize}

\textbf{Congratulations! You now have a powerful and accessible
development environment ready for your Python journey.}

\newpage

\section{Unit 1: Foundations of Computing and Getting Started with
Python}\label{unit-1-foundations-of-computing-and-getting-started-with-python}

\subsection{Overview}\label{overview-1}

This unit covers basic computer science concepts, setting up the Python
environment, and writing and executing simple Python programs.

\textbf{Objectives:}

\begin{itemize}
\tightlist
\item
  Understand what computer science is and how hardware relates to
  software.
\item
  Write and run your first Python program.
\item
  Use variables to store and manipulate simple data types.
\end{itemize}

\begin{center}\rule{0.5\linewidth}{0.5pt}\end{center}

\subsection{Lesson 1.1: Welcome to Computer
Science}\label{lesson-1.1-welcome-to-computer-science-1}

\subsection{Reading References}\label{reading-references-4}

\begin{itemize}
\tightlist
\item
  Zelle: Chapter 1, pages 1--7
\item
  Matthes: Chapter 1, pages 3--14
\end{itemize}

\subsection{What Is Computer
Science?}\label{what-is-computer-science}

A computer is a \textbf{universal machine} --- it can be programmed to
solve any problem that can be broken down into a series of precise,
logical steps. The true power of a computer doesn't come from its
hardware but from the software we create.

\textbf{Computer science} is the study of problem-solving, not just
programming. It's an interdisciplinary field involving breaking down
complex problems, designing algorithms, and implementing efficient
solutions.

\subsection{Hardware vs.~Software}\label{hardware-vs.-software}

\begin{fullwidth}{\def\LTcaptype{none} % do not increment counter
\begin{longtable}[]{@{}
  >{\raggedright\arraybackslash}p{(\linewidth - 2\tabcolsep) * \real{0.1860}}
  >{\raggedright\arraybackslash}p{(\linewidth - 2\tabcolsep) * \real{0.8140}}@{}}
\toprule\noalign{}
\begin{minipage}[b]{\linewidth}\raggedright
Component
\end{minipage} & \begin{minipage}[b]{\linewidth}\raggedright
Description
\end{minipage} \\
\midrule\noalign{}
\endhead
\bottomrule\noalign{}
\endlastfoot
\textbf{Hardware} & Physical components you can touch: keyboard, screen,
processor, memory \\
\textbf{Software} & Instructions that direct hardware: programs,
operating systems, apps \\
\textbf{CPU} & The processor that executes instructions \\
\textbf{Memory (RAM)} & Temporary storage for active programs and
data \\
\textbf{Storage} & Permanent storage for files \\
\end{longtable}
}
\end{fullwidth}

\textbf{Analogy:} Think of a kitchen. The stove, refrigerator, and pots
are \textbf{hardware}. A cookie recipe is \textbf{software} --- it
provides precise steps for the ``hardware'' (the cook) to follow.

\subsection{Why Python?}\label{why-python}

Python is an ideal language for beginners because of its:

\begin{itemize}
\tightlist
\item
  \textbf{Readability} --- English-like syntax makes code easy to
  understand.
\item
  \textbf{Simplicity} --- Less ``boilerplate'' code than other
  languages.
\item
  \textbf{Versatility} --- Used in web development, data science, AI,
  and more.
\end{itemize}

\subsection{The Problem-Solving
Approach}\label{the-problem-solving-approach}

\begin{enumerate}
\def\labelenumi{\arabic{enumi}.}
\tightlist
\item
  Understand the problem.
\item
  Break it down into smaller pieces.
\item
  Design a step-by-step solution (an algorithm).
\item
  Implement the solution in Python.
\item
  Test and refine your code.
\end{enumerate}

\subsection{Activities}\label{activities}

\textbf{Activity 1.1.1: ``Human Computer'' Role-Play}

\begin{itemize}
\tightlist
\item
  In pairs, one student is the ``computer'' and the other is the
  ``programmer.''
\item
  The programmer provides simple, precise instructions for a common task
  (e.g., ``make a peanut butter and jelly sandwich'').
\item
  The ``computer'' follows only what is literally said.
\item
  Purpose: Illustrate the need for precise, unambiguous instructions.
\end{itemize}

\textbf{Activity 1.1.2: Identify Hardware and Software}

\begin{itemize}
\tightlist
\item
  List five items you use daily.
\item
  For each, identify whether it is hardware or software and explain your
  choice.
\end{itemize}

\subsection{Skills Checklist}\label{skills-checklist-4}

\begin{itemize}
\tightlist
\item[$\square$]
  I can define what computer science is.
\item[$\square$]
  I can distinguish between computer hardware and software.
\item[$\square$]
  I can explain what a computer program is.
\item[$\square$]
  I can identify Python as a programming language.
\end{itemize}

\begin{center}\rule{0.5\linewidth}{0.5pt}\end{center}

\subsection{Lesson 1.2: Your First Python Program --- ``Hello
World!''}\label{lesson-1.2-your-first-python-program-hello-world}

\subsection{Reading References}\label{reading-references-5}

\begin{itemize}
\tightlist
\item
  Zelle: Chapter 1, pages 7--13
\item
  Matthes: Chapter 1, pages 3--14
\end{itemize}

\subsection{\texorpdfstring{The \texttt{print()}
Function}{The print() Function}}\amn{\textbf{Accessibility Tip:} Every \texttt{print()} call produces one line of terminal output. Navigate output with the arrow keys in your terminal. Use \texttt{Ctrl+End} to jump to the most recent output, then arrow up to read previous lines. Always add a descriptive label: \texttt{print("Name:", name)} rather than just \texttt{print(name)}.}
\subsection{\texorpdfstring{The \texttt{print()}
Function}{The print() Function}}\label{the-print-function}

The \texttt{print()} function is Python's primary method for displaying
output to the user.

\begin{Shaded}
\begin{Highlighting}[]
print(\StringTok{"Hello, World!"})
\end{Highlighting}
\end{Shaded}

\textbf{Expected output:}

\begin{verbatim}
Hello, World!
\end{verbatim}

\subsection{Key Concepts}\label{key-concepts}

\begin{itemize}
\tightlist
\item
  \textbf{Strings} are text data enclosed in quotation marks (single or
  double).
\item
  \textbf{Comments} start with \texttt{\#} and are ignored by Python.
\item
  Python executes code \textbf{line by line}, from top to bottom.
\end{itemize}

\subsection{Examples}\label{examples}

\textbf{Example 1.2.1: Simple Greeting}

\begin{Shaded}
\begin{Highlighting}[]
\CommentTok{\# This program prints a classic greeting.}
print(\StringTok{"Hello, World!"})
\end{Highlighting}
\end{Shaded}

\textbf{Example 1.2.2: Personalized Message}

\begin{Shaded}
\begin{Highlighting}[]
\CommentTok{\# This program prints a message with your name.}
print(\StringTok{"Hello, Alice!"})
print(\StringTok{"Welcome to Python programming!"})
\end{Highlighting}
\end{Shaded}

Expected output:

\begin{verbatim}
Hello, Alice!
Welcome to Python programming!
\end{verbatim}

\subsection{Common Mistakes}\label{common-mistakes}

\begin{Shaded}
\begin{Highlighting}[]
print(\StringTok{"Hello, World"}        \CommentTok{\# Missing closing parenthesis {-}{-}{-} SyntaxError}
print \StringTok{"Hello"}               \CommentTok{\# Python 2 style {-}{-}{-} SyntaxError in Python 3}
Print(\StringTok{"Hello"})              \CommentTok{\# Capital P {-}{-}{-} NameError}
\end{Highlighting}
\end{Shaded}

\subsection{Activities}\label{activities-1}

\textbf{Activity 1.2.1: My First Message}

\begin{enumerate}
\def\labelenumi{\arabic{enumi}.}
\item
  Open VS Code and confirm the PythonProjects folder is open.
\item
  Create a new file as \texttt{my\_message.py}.
\item
  Type:

\begin{Shaded}
\begin{Highlighting}[]
print(\StringTok{"My favorite color is blue."})
print(\StringTok{"My favorite animal is a dog."})
\end{Highlighting}
\end{Shaded}
\item
  Save and run the file.
\end{enumerate}

\textbf{Activity 1.2.2: Exploring Print Statements}

\begin{enumerate}
\def\labelenumi{\arabic{enumi}.}
\tightlist
\item
  Experiment with:

  \begin{itemize}
  \tightlist
  \item
    Printing numbers: \texttt{print(123)}
  \item
    Printing sentences with punctuation.
  \item
    Omitting quotes: \texttt{print(Hello)} --- expect a
    \texttt{NameError}.
  \end{itemize}
\end{enumerate}

\subsection{Skills Checklist}\label{skills-checklist-5}

\begin{itemize}
\tightlist
\item[$\square$]
  I can write a \texttt{print()} statement in Python.
\item[$\square$]
  I can save a Python file with the \texttt{.py} extension.
\item[$\square$]
  I can execute a Python script from the terminal.
\item[$\square$]
  I understand that the output of \texttt{print()} appears in the
  terminal.
\end{itemize}

\begin{center}\rule{0.5\linewidth}{0.5pt}\end{center}

\subsection{Lesson 1.3: Variables and Simple Data
Types}\label{lesson-1.3-variables-and-simple-data-types-1}

\subsection{Reading References}\label{reading-references-6}

\begin{itemize}
\tightlist
\item
  Zelle: Chapter 2, pages 24--26
\item
  Matthes: Chapter 2, pages 15--32
\end{itemize}

\subsection{What Is a Variable?}\label{what-is-a-variable}\amn{\textbf{Accessibility Tip:} Give every variable a descriptive name that says what it contains: \texttt{student\_age} rather than \texttt{x}. When your screen reader announces variable names during debugging, descriptive names make the output immediately meaningful.}
\subsection{What Is a Variable?}\label{what-is-a-variable}\amn{\textbf{Accessibility Tip:} Give every variable a descriptive name that says what it contains: \texttt{student\_age} rather than \texttt{x}. When your screen reader announces variable names during debugging, descriptive names make the output immediately meaningful.}
\subsection{What Is a Variable?}\label{what-is-a-variable}

A \textbf{variable} is a name that refers to a value stored in the
computer's memory. The \textbf{assignment operator} \texttt{=} creates a
connection between a name and a value.

\begin{Shaded}
\begin{Highlighting}[]
age \OperatorTok{=} \CodeNumTok{25}
name \OperatorTok{=} \StringTok{"Alice"}
height \OperatorTok{=} \CodeNumTok{5.6}
\end{Highlighting}
\end{Shaded}

\textbf{Rules for variable names:}

\begin{itemize}
\tightlist
\item
  Must start with a letter or underscore
\item
  Can only contain letters, digits, and underscores
\item
  Are case-sensitive (\texttt{age} is different from \texttt{Age})
\item
  Cannot be a Python keyword (\texttt{if}, \texttt{for}, \texttt{while},
  etc.)
\end{itemize}

\textbf{Best practices:}

\begin{itemize}
\tightlist
\item
  Use descriptive names: \texttt{student\_count}, not \texttt{sc}
\item
  Use lowercase with underscores: \texttt{first\_name}
\end{itemize}

\subsection{Essential Data Types}\label{essential-data-types}

\begin{fullwidth}{\def\LTcaptype{none} % do not increment counter
\begin{longtable}[]{@{}lll@{}}
\toprule\noalign{}
Type & Description & Examples \\
\midrule\noalign{}
\endhead
\bottomrule\noalign{}
\endlastfoot
\texttt{str} (String) & Text data & \texttt{"hello"},
\texttt{"Python"} \\
\texttt{int} (Integer) & Whole numbers & \texttt{10}, \texttt{-5},
\texttt{0} \\
\texttt{float} (Float) & Numbers with decimals & \texttt{3.14},
\texttt{9.99} \\
\texttt{bool} (Boolean) & True or False & \texttt{True},
\texttt{False} \\
\end{longtable}
}
\end{fullwidth}

\subsection{Examples}\label{examples-1}

\textbf{Example 1.3.1: Storing and Printing a String}

\begin{Shaded}
\begin{Highlighting}[]
student\_name \OperatorTok{=} \StringTok{"Charlie"}
print(\StringTok{"Hello, "} \OperatorTok{+} student\_name \OperatorTok{+} \StringTok{"!"})
\end{Highlighting}
\end{Shaded}

Output:

\begin{verbatim}
Hello, Charlie!
\end{verbatim}

\textbf{Example 1.3.2: Storing and Printing Numbers}

\begin{Shaded}
\begin{Highlighting}[]
apple\_count \OperatorTok{=} \CodeNumTok{5}
orange\_count \OperatorTok{=} \CodeNumTok{3}
total\_fruit \OperatorTok{=} apple\_count \OperatorTok{+} orange\_count
print(\StringTok{"You have"}, total\_fruit, \StringTok{"fruits in total."})
 
price\_per\_item \OperatorTok{=} \CodeNumTok{2.50}
quantity \OperatorTok{=} \CodeNumTok{4}
total\_cost \OperatorTok{=} price\_per\_item \OperatorTok{*} quantity
print(\StringTok{"The total cost is:"}, total\_cost)
\end{Highlighting}
\end{Shaded}

Output:

\begin{verbatim}
You have 8 fruits in total.
The total cost is: 10.0
\end{verbatim}

\subsection{Common Errors}\label{common-errors}

\begin{Shaded}
\begin{Highlighting}[]
\CodeNumTok{2}\ErrorTok{nd\_number} \OperatorTok{=} \CodeNumTok{10}         \CommentTok{\# Invalid variable name {-}{-}{-} starts with digit}
print(\StringTok{"Result: "} \OperatorTok{+} \CodeNumTok{42})  \CommentTok{\# TypeError {-}{-}{-} can\textquotesingle{}t concatenate str and int}
\end{Highlighting}
\end{Shaded}

\subsection{Activities}\label{activities-2}

\textbf{Activity 1.3.1: Personal Information Variables}

\begin{Shaded}
\begin{Highlighting}[]
my\_name \OperatorTok{=} \StringTok{"Jamie"}
favorite\_number \OperatorTok{=} \CodeNumTok{7}
height\_in\_meters \OperatorTok{=} \CodeNumTok{1.75}
favorite\_food \OperatorTok{=} \StringTok{"pizza"}
 
print(\StringTok{"My name is:"}, my\_name)
print(\StringTok{"My favorite number is:"}, favorite\_number)
print(\StringTok{"My height is:"}, height\_in\_meters, \StringTok{"meters"})
print(\StringTok{"My favorite food is:"}, favorite\_food)
\end{Highlighting}
\end{Shaded}

\textbf{Activity 1.3.2: Simple Math Challenge}

\begin{Shaded}
\begin{Highlighting}[]
num1 \OperatorTok{=} \CodeNumTok{15}
num2 \OperatorTok{=} \CodeNumTok{7}
sum\_result \OperatorTok{=} num1 \OperatorTok{+} num2
difference \OperatorTok{=} num1 \OperatorTok{{-}} num2
product \OperatorTok{=} num1 \OperatorTok{*} num2
quotient \OperatorTok{=} num1 \OperatorTok{/} num2
 
print(\StringTok{"Sum:"}, sum\_result)
print(\StringTok{"Difference:"}, difference)
print(\StringTok{"Product:"}, product)
print(\StringTok{"Quotient:"}, quotient)
\end{Highlighting}
\end{Shaded}

\subsection{Key Table: Python Data Types and
Operations}\label{key-table-python-data-types-and-operations}

\begin{fullwidth}{\def\LTcaptype{none} % do not increment counter
\begin{longtable}[]{@{}
  >{\raggedright\arraybackslash}p{(\linewidth - 6\tabcolsep) * \real{0.0978}}
  >{\raggedright\arraybackslash}p{(\linewidth - 6\tabcolsep) * \real{0.2391}}
  >{\raggedright\arraybackslash}p{(\linewidth - 6\tabcolsep) * \real{0.2283}}
  >{\raggedright\arraybackslash}p{(\linewidth - 6\tabcolsep) * \real{0.4348}}@{}}
\toprule\noalign{}
\begin{minipage}[b]{\linewidth}\raggedright
Data Type
\end{minipage} & \begin{minipage}[b]{\linewidth}\raggedright
Description
\end{minipage} & \begin{minipage}[b]{\linewidth}\raggedright
Example Value
\end{minipage} & \begin{minipage}[b]{\linewidth}\raggedright
Common Operations
\end{minipage} \\
\midrule\noalign{}
\endhead
\bottomrule\noalign{}
\endlastfoot
\texttt{str} & Text, words, sentences & \texttt{"Hello"},
\texttt{"Python"} & \texttt{+} (concatenation), \texttt{len()},
\texttt{.upper()} \\
\texttt{int} & Whole numbers & \texttt{10}, \texttt{-5}, \texttt{0} &
\texttt{+}, \texttt{-}, \texttt{*}, \texttt{/}, \texttt{\%} \\
\texttt{float} & Numbers with decimals & \texttt{3.14}, \texttt{9.99},
\texttt{0.0} & \texttt{+}, \texttt{-}, \texttt{*}, \texttt{/} \\
\end{longtable}
}
\end{fullwidth}

\subsection{Skills Checklist}\label{skills-checklist-6}

\begin{itemize}
\tightlist
\item[$\square$]
  I can define what a variable is in programming.
\item[$\square$]
  I can create a valid variable name in Python.
\item[$\square$]
  I can assign string, integer, and float values to variables.
\item[$\square$]
  I can use \texttt{print()} to display the value of a variable.
\item[$\square$]
  I can perform basic arithmetic operations with numbers in Python.
\end{itemize}

\newpage

\section{Unit 2: Writing Simple Programs and Numeric
Computations}\label{unit-2-writing-simple-programs-and-numeric-computations}

\subsection{Overview}\label{overview-2}

This unit explores the software development process, more complex
numeric operations, and how to obtain and use input from users.

\textbf{Objectives:}

\begin{itemize}
\tightlist
\item
  Understand the phases of software development.
\item
  Work with different numeric data types (integers, floats).
\item
  Perform type conversions.
\item
  Use mathematical operators and the math library.
\item
  Get user input with \texttt{input()}.
\end{itemize}

\begin{center}\rule{0.5\linewidth}{0.5pt}\end{center}

\subsection{Lesson 2.1: The Software Development
Process}\label{lesson-2.1-the-software-development-process-1}

\subsection{Reading References}\label{reading-references-7}

\begin{itemize}
\tightlist
\item
  Zelle: Chapter 2, pages 21--22 and pages 35--36
\end{itemize}

\subsection{The Three-Phase Development
Process}\label{the-three-phase-development-process}

\textbf{1. Design (Plan first)} Before writing any code, fully
understand the problem:

\begin{itemize}
\tightlist
\item
  Identify \textbf{inputs} (information the program needs) and
  \textbf{outputs} (results it must produce).
\item
  Write an \textbf{algorithm} --- a step-by-step plan in plain English
  (\textbf{pseudocode}).
\item
  Consider \textbf{edge cases} --- unusual inputs that might break your
  program.
\end{itemize}

\textbf{2. Implement (Write the code)}

\begin{itemize}
\tightlist
\item
  Translate your design into Python syntax.
\item
  Code \textbf{incrementally} --- build in small pieces and test as you
  go.
\item
  Follow coding standards and use descriptive variable names.
\end{itemize}

\textbf{3. Test and Debug (Verify it works)}

\begin{itemize}
\tightlist
\item
  Run with a variety of inputs, including normal values and edge cases.
\item
  Read error messages carefully --- they tell you the type of error and
  the line number.
\item
  Use \texttt{print()} statements to check variable values during
  debugging.
\end{itemize}

\subsection{Why This Process Matters}\label{why-this-process-matters}

\begin{itemize}
\tightlist
\item
  \textbf{Reduces frustration} --- Having a plan prevents getting stuck.
\item
  \textbf{Improves code quality} --- A thoughtful design leads to
  cleaner code.
\item
  \textbf{Simplifies debugging} --- Test each component in isolation.
\item
  \textbf{Saves time} --- Good design prevents time-consuming rewrites.
\end{itemize}

\subsection{Activities}\label{activities-3}

\textbf{Activity 2.1.1: Planning a Tip Calculator} Write pseudocode
before coding:

\begin{verbatim}
1. GET bill amount from user
2. CALCULATE tip as 15% of bill amount
3. CALCULATE total amount (bill + tip)
4. PRINT tip amount
5. PRINT total amount
\end{verbatim}

\textbf{Activity 2.1.2: Tracing a Simple Program} Without running the
code, trace the execution and predict output:

\begin{Shaded}
\begin{Highlighting}[]
a \OperatorTok{=} \CodeNumTok{10}
b \OperatorTok{=} \CodeNumTok{5}
c \OperatorTok{=} a \OperatorTok{*} b
a \OperatorTok{=} a \OperatorTok{+} \CodeNumTok{2}
print(c)
print(a)
\end{Highlighting}
\end{Shaded}

\subsection{Skills Checklist}\label{skills-checklist-7}

\begin{itemize}
\tightlist
\item[$\square$]
  I can describe the purpose of the software development process.
\item[$\square$]
  I can list and briefly explain the main phases of software
  development.
\item[$\square$]
  I understand that problem-solving and design happen before writing
  code.
\end{itemize}

\begin{center}\rule{0.5\linewidth}{0.5pt}\end{center}

\subsection{Lesson 2.2: Elements of Programs --- Expressions and
Assignment}\label{lesson-2.2-elements-of-programs-expressions-and-assignment}

\subsection{Reading References}\label{reading-references-8}

\begin{itemize}
\tightlist
\item
  Zelle: Chapter 2, pages 24--26
\end{itemize}

\subsection{Expressions and
Assignment}\label{expressions-and-assignment}

An \textbf{expression} is any combination of values, variables, and
operators that evaluates to a single result. Python follows the standard
mathematical order of operations (PEMDAS).

\begin{Shaded}
\begin{Highlighting}[]
result1 \OperatorTok{=} \CodeNumTok{10} \OperatorTok{/} \CodeNumTok{2} \OperatorTok{+} \CodeNumTok{3}  \CommentTok{\# 10/2 is 5, then 5+3 = 8.0}
message \OperatorTok{=} \StringTok{"Python"} \OperatorTok{+} \StringTok{" "} \OperatorTok{+} \StringTok{"is fun!"}
\end{Highlighting}
\end{Shaded}

\textbf{Assignment:}

\begin{Shaded}
\begin{Highlighting}[]
score \OperatorTok{=} \CodeNumTok{50}
score \OperatorTok{=} score \OperatorTok{+} \CodeNumTok{20}  \CommentTok{\# score is now 70}
\end{Highlighting}
\end{Shaded}

\subsection{Examples}\label{examples-2}

\textbf{Example 2.2.1: Expression Evaluation}

\begin{Shaded}
\begin{Highlighting}[]
result1 \OperatorTok{=} \CodeNumTok{10} \OperatorTok{/} \CodeNumTok{2} \OperatorTok{+} \CodeNumTok{3}
print(\StringTok{"Result 1:"}, result1)   \CommentTok{\# Output: Result 1: 8.0}
 
message \OperatorTok{=} \StringTok{"Python"} \OperatorTok{+} \StringTok{" "} \OperatorTok{+} \StringTok{"is fun!"}
print(\StringTok{"Message:"}, message)    \CommentTok{\# Output: Message: Python is fun!}
\end{Highlighting}
\end{Shaded}

\textbf{Example 2.2.2: Reassigning Variables}

\begin{Shaded}
\begin{Highlighting}[]
score \OperatorTok{=} \CodeNumTok{50}
print(\StringTok{"Initial score:"}, score)
score \OperatorTok{=} score \OperatorTok{+} \CodeNumTok{20}
print(\StringTok{"New score:"}, score)     \CommentTok{\# Output: New score: 70}
\end{Highlighting}
\end{Shaded}

\subsection{Skills Checklist}\label{skills-checklist-8}

\begin{itemize}
\tightlist
\item[$\square$]
  I can identify an expression in Python code.
\item[$\square$]
  I can predict the result of simple Python expressions.
\item[$\square$]
  I can explain how the assignment operator \texttt{=} works.
\item[$\square$]
  I can demonstrate how to update a variable's value.
\end{itemize}

\begin{center}\rule{0.5\linewidth}{0.5pt}\end{center}

\subsection{Lesson 2.3: Input and Output
Revisited}\label{lesson-2.3-input-and-output-revisited}

\subsection{Reading References}\label{reading-references-9}

\begin{itemize}
\tightlist
\item
  Zelle: Chapter 2, pages 27--33
\item
  Matthes: Chapter 7, pages 114--116
\end{itemize}

\subsection{\texorpdfstring{The \texttt{input()}
Function}{The input() Function}}\amn{\textbf{Accessibility Tip:} \texttt{input()} pauses your program and waits for you to type. Include a clear prompt describing exactly what to type: \texttt{input("Enter your age as a whole number: ")}. After pressing Enter, immediately call \texttt{print(type(variable))} to verify the value is the type you expect.}
\subsection{\texorpdfstring{The \texttt{input()}
Function}{The input() Function}}\label{the-input-function}

The \texttt{input()} function pauses program execution and waits for the
user to type something and press Enter. \textbf{It always returns a
string.}

\begin{Shaded}
\begin{Highlighting}[]
user\_name \OperatorTok{=} input(\StringTok{"What is your name? "})
print(\StringTok{f"Hello, }\StringTok{\{}user\_name\StringTok{\}}\StringTok{! Nice to meet you."})
\end{Highlighting}
\end{Shaded}

\subsection{F-Strings (Formatted String
Literals)}\label{f-strings-formatted-string-literals}

F-strings provide a clean way to embed variable values directly into a
string:

\begin{Shaded}
\begin{Highlighting}[]
name \OperatorTok{=} \StringTok{"Alice"}
age \OperatorTok{=} \CodeNumTok{30}
print(\StringTok{f"Hello, }\StringTok{\{}name\StringTok{\}}\StringTok{! You are }\StringTok{\{}age\StringTok{\}}\StringTok{ years old."})
\end{Highlighting}
\end{Shaded}

\subsection{Converting Input for
Calculations}\label{converting-input-for-calculations}

Since \texttt{input()} always returns a string, you must convert it for
math:

\begin{Shaded}
\begin{Highlighting}[]
age\_str \OperatorTok{=} input(\StringTok{"How old are you? "})
age\_int \OperatorTok{=} int(age\_str)
print(\StringTok{f"Next year you will be }\StringTok{\{}age\_int \OperatorTok{+} \CodeNumTok{1}\StringTok{\}}\StringTok{ years old."})
\end{Highlighting}
\end{Shaded}

\subsection{Examples}\label{examples-3}

\textbf{Example 2.3.1: Basic Input}

\begin{Shaded}
\begin{Highlighting}[]
user\_name \OperatorTok{=} input(\StringTok{"What is your name? "})
print(\StringTok{f"Hello, }\StringTok{\{}user\_name\StringTok{\}}\StringTok{! Nice to meet you."})
\end{Highlighting}
\end{Shaded}

\textbf{Example 2.3.2: Numeric Input and Calculation}

\begin{Shaded}
\begin{Highlighting}[]
num1\_str \OperatorTok{=} input(\StringTok{"Enter the first number: "})
num1 \OperatorTok{=} int(num1\_str)
 
num2\_str \OperatorTok{=} input(\StringTok{"Enter the second number: "})
num2 \OperatorTok{=} int(num2\_str)
 
sum\_result \OperatorTok{=} num1 \OperatorTok{+} num2
print(\StringTok{f"The sum of }\StringTok{\{}num1\StringTok{\}}\StringTok{ and }\StringTok{\{}num2\StringTok{\}}\StringTok{ is }\StringTok{\{}sum\_result\StringTok{\}}\StringTok{."})
\end{Highlighting}
\end{Shaded}

\subsection{Activities}\label{activities-4}

\textbf{Activity 2.3.1: Temperature Converter}

\begin{Shaded}
\begin{Highlighting}[]
fahrenheit\_str \OperatorTok{=} input(\StringTok{"Enter temperature in Fahrenheit: "})
fahrenheit \OperatorTok{=} float(fahrenheit\_str)
celsius \OperatorTok{=} (fahrenheit \OperatorTok{{-}} \CodeNumTok{32}) \OperatorTok{*} \CodeNumTok{5} \OperatorTok{/} \CodeNumTok{9}
print(\StringTok{f"}\StringTok{\{}fahrenheit\StringTok{\}}\StringTok{°F is }\StringTok{\{}celsius\StringTok{:.2f\}}\StringTok{°C"})
\end{Highlighting}
\end{Shaded}

\subsection{Skills Checklist}\label{skills-checklist-9}

\begin{itemize}
\tightlist
\item[$\square$]
  I can use the \texttt{input()} function to ask the user for
  information.
\item[$\square$]
  I can provide a clear prompt message when using \texttt{input()}.
\item[$\square$]
  I remember that \texttt{input()} returns a string, even for numbers.
\item[$\square$]
  I can combine text and variable values in a \texttt{print()} statement
  using f-strings.
\end{itemize}

\begin{center}\rule{0.5\linewidth}{0.5pt}\end{center}

\subsection{Lesson 2.4: Computing with Numbers and Type
Conversions}\label{lesson-2.4-computing-with-numbers-and-type-conversions}

\subsection{Reading References}\label{reading-references-10}

\begin{itemize}
\tightlist
\item
  Zelle: Chapter 3, pages 41--48
\item
  Matthes: Chapter 2, pages 26--28
\end{itemize}

\subsection{Integers and Floats}\label{integers-and-floats}

\begin{fullwidth}{\def\LTcaptype{none} % do not increment counter
\begin{longtable}[]{@{}lll@{}}
\toprule\noalign{}
Type & Description & Example \\
\midrule\noalign{}
\endhead
\bottomrule\noalign{}
\endlastfoot
\texttt{int} & Whole numbers (exact) & \texttt{5}, \texttt{-20},
\texttt{1000000} \\
\texttt{float} & Numbers with decimals & \texttt{3.14}, \texttt{-2.5},
\texttt{6.0} \\
\end{longtable}
}
\end{fullwidth}

\textbf{Key behaviors:}

\begin{itemize}
\tightlist
\item
  Division \texttt{/} always returns a float: \texttt{10\ /\ 2}
  -\textgreater{} \texttt{5.0}
\item
  Floor division \texttt{//} returns the integer part:
  \texttt{10\ //\ 3} -\textgreater{} \texttt{3}
\item
  Mixed operations promote to float: \texttt{10\ +\ 3.5} -\textgreater{}
  \texttt{13.5}
\end{itemize}

\subsection{Type Conversion
Functions}\label{type-conversion-functions}

\begin{Shaded}
\begin{Highlighting}[]
int(\StringTok{"123"})      \CommentTok{\# {-}\textgreater{} 123  (truncates: int(3.9) {-}\textgreater{} 3)}
float(\StringTok{"3.14"})   \CommentTok{\# {-}\textgreater{} 3.14}
str(\CodeNumTok{42})         \CommentTok{\# {-}\textgreater{} "42"}
\end{Highlighting}
\end{Shaded}

\subsection{Examples}\label{examples-4}

\textbf{Example 2.4.1: Type Conversion}

\begin{Shaded}
\begin{Highlighting}[]
num\_str \OperatorTok{=} \StringTok{"123"}
num\_int \OperatorTok{=} int(num\_str)
print(\StringTok{f"String \textquotesingle{}}\StringTok{\{}num\_str\StringTok{\}}\StringTok{\textquotesingle{} as integer: }\StringTok{\{}num\_int\StringTok{\}}\StringTok{, type: }\StringTok{\{}type(num\_int)\StringTok{\}}\StringTok{"})
 
num\_float \OperatorTok{=} float(num\_str)
print(\StringTok{f"String \textquotesingle{}}\StringTok{\{}num\_str\StringTok{\}}\StringTok{\textquotesingle{} as float: }\StringTok{\{}num\_float\StringTok{\}}\StringTok{, type: }\StringTok{\{}type(num\_float)\StringTok{\}}\StringTok{"})
 
my\_number \OperatorTok{=} \CodeNumTok{42}
my\_text \OperatorTok{=} str(my\_number)
print(\StringTok{f"Integer }\StringTok{\{}my\_number\StringTok{\}}\StringTok{ as string: \textquotesingle{}}\StringTok{\{}my\_text\StringTok{\}}\StringTok{\textquotesingle{}, type: }\StringTok{\{}type(my\_text)\StringTok{\}}\StringTok{"})
\end{Highlighting}
\end{Shaded}

\textbf{Common Type Error:}

\begin{Shaded}
\begin{Highlighting}[]
\CommentTok{\# This will cause a TypeError:}
\CommentTok{\# print("My age is " + 14)}
\CommentTok{\# Fix: print("My age is " + str(14))}
\CommentTok{\# Or: print(f"My age is \{14\}")}
\end{Highlighting}
\end{Shaded}

\subsection{Activities}\label{activities-5}

\textbf{Activity 2.4.1: Mixed Bag Calculator}

\begin{Shaded}
\begin{Highlighting}[]
age\_str \OperatorTok{=} input(\StringTok{"Enter your age: "})
age \OperatorTok{=} int(age\_str)
 
weight\_str \OperatorTok{=} input(\StringTok{"Enter your weight in kg: "})
weight \OperatorTok{=} float(weight\_str)
 
phrase \OperatorTok{=} input(\StringTok{"Enter a short phrase: "})
 
print(\StringTok{f"Age doubled: }\StringTok{\{}age \OperatorTok{*} \CodeNumTok{2}\StringTok{\}}\StringTok{"})
print(\StringTok{f"Weight + age: }\StringTok{\{}weight \OperatorTok{+} age\StringTok{\}}\StringTok{"})
print(\StringTok{f"Phrase length: }\StringTok{\{}len(phrase)\StringTok{\}}\StringTok{ characters"})
\end{Highlighting}
\end{Shaded}

\subsection{Skills Checklist}\label{skills-checklist-10}

\begin{itemize}
\tightlist
\item[$\square$]
  I can identify when to use an integer versus a float.
\item[$\square$]
  I can convert a string to an integer using \texttt{int()}.
\item[$\square$]
  I can convert a string to a float using \texttt{float()}.
\item[$\square$]
  I can convert a number to a string using \texttt{str()}.
\item[$\square$]
  I can apply type conversions to perform calculations on user input.
\item[$\square$]
  I understand and can apply the order of operations.
\end{itemize}

\begin{center}\rule{0.5\linewidth}{0.5pt}\end{center}

\subsection{Lesson 2.5: Using the Math
Library}\label{lesson-2.5-using-the-math-library}

\subsection{Reading References}\label{reading-references-11}

\begin{itemize}
\tightlist
\item
  Zelle: Chapter 3, pages 49--51
\end{itemize}

\subsection{Importing the Math
Module}\label{importing-the-math-module}

\begin{Shaded}
\begin{Highlighting}[]
\CodeAccentTok{import} math
\end{Highlighting}
\end{Shaded}

\subsection{Key Math Functions and
Constants}\label{key-math-functions-and-constants}

\begin{fullwidth}{\def\LTcaptype{none} % do not increment counter
\begin{longtable}[]{@{}
  >{\raggedright\arraybackslash}p{(\linewidth - 4\tabcolsep) * \real{0.2533}}
  >{\raggedright\arraybackslash}p{(\linewidth - 4\tabcolsep) * \real{0.3467}}
  >{\raggedright\arraybackslash}p{(\linewidth - 4\tabcolsep) * \real{0.4000}}@{}}
\toprule\noalign{}
\begin{minipage}[b]{\linewidth}\raggedright
Function/Constant
\end{minipage} & \begin{minipage}[b]{\linewidth}\raggedright
Description
\end{minipage} & \begin{minipage}[b]{\linewidth}\raggedright
Example
\end{minipage} \\
\midrule\noalign{}
\endhead
\bottomrule\noalign{}
\endlastfoot
\texttt{math.pi} & π ≈ 3.14159 & \texttt{math.pi} \\
\texttt{math.e} & e ≈ 2.71828 & \texttt{math.e} \\
\texttt{math.sqrt(x)} & Square root of x & \texttt{math.sqrt(25)}
-\textgreater{} \texttt{5.0} \\
\texttt{math.pow(x,\ y)} & x raised to power y &
\texttt{math.pow(2,\ 4)} -\textgreater{} \texttt{16.0} \\
\texttt{math.ceil(x)} & Round up & \texttt{math.ceil(4.2)}
-\textgreater{} \texttt{5} \\
\texttt{math.floor(x)} & Round down & \texttt{math.floor(4.8)}
-\textgreater{} \texttt{4} \\
\texttt{math.log(x,\ base)} & Logarithm & \texttt{math.log(100,\ 10)}
-\textgreater{} \texttt{2.0} \\
\texttt{math.sin(x)} & Sine (in radians) & \texttt{math.sin(math.pi/2)}
-\textgreater{} \texttt{1.0} \\
\texttt{math.radians(deg)} & Convert degrees to radians &
\texttt{math.radians(90)} \\
\end{longtable}
}
\end{fullwidth}

\subsection{Examples}\label{examples-5}

\textbf{Example 2.5.1: Circle Calculations}

\begin{Shaded}
\begin{Highlighting}[]
\CodeAccentTok{import} math
 
radius \OperatorTok{=} \CodeNumTok{5}
area \OperatorTok{=} math.pi \OperatorTok{*} radius \OperatorTok{**} \CodeNumTok{2}
circumference \OperatorTok{=} \CodeNumTok{2} \OperatorTok{*} math.pi \OperatorTok{*} radius
 
print(\StringTok{f"Area of circle with radius }\StringTok{\{}radius\StringTok{\}}\StringTok{: }\StringTok{\{}area\StringTok{:.2f\}}\StringTok{"})
print(\StringTok{f"Circumference: }\StringTok{\{}circumference\StringTok{:.2f\}}\StringTok{"})
 
number \OperatorTok{=} \CodeNumTok{81}
square\_root \OperatorTok{=} math.sqrt(number)
print(\StringTok{f"Square root of }\StringTok{\{}number\StringTok{\}}\StringTok{: }\StringTok{\{}square\_root\StringTok{\}}\StringTok{"})
\end{Highlighting}
\end{Shaded}

\subsection{Activities}\label{activities-6}

\textbf{Activity 2.5.1: Circle Calculator}

\begin{Shaded}
\begin{Highlighting}[]
\CodeAccentTok{import} math
 
radius\_str \OperatorTok{=} input(\StringTok{"Enter the radius: "})
radius \OperatorTok{=} float(radius\_str)
 
area \OperatorTok{=} math.pi \OperatorTok{*} radius \OperatorTok{**} \CodeNumTok{2}
circumference \OperatorTok{=} \CodeNumTok{2} \OperatorTok{*} math.pi \OperatorTok{*} radius
 
print(\StringTok{f"Area: }\StringTok{\{}area\StringTok{:.2f\}}\StringTok{"})
print(\StringTok{f"Circumference: }\StringTok{\{}circumference\StringTok{:.2f\}}\StringTok{"})
\end{Highlighting}
\end{Shaded}

\textbf{Activity 2.5.2: Distance Finder}

\begin{Shaded}
\begin{Highlighting}[]
\CodeAccentTok{import} math
 
x \OperatorTok{=} float(input(\StringTok{"Enter x coordinate: "}))
y \OperatorTok{=} float(input(\StringTok{"Enter y coordinate: "}))
distance \OperatorTok{=} math.sqrt(x\OperatorTok{**}\CodeNumTok{2} \OperatorTok{+} y\OperatorTok{**}\CodeNumTok{2})
print(\StringTok{f"Distance from origin: }\StringTok{\{}distance\StringTok{:.2f\}}\StringTok{"})
\end{Highlighting}
\end{Shaded}

\subsection{Skills Checklist}\label{skills-checklist-11}

\begin{itemize}
\tightlist
\item[$\square$]
  I can explain what a Python module is.
\item[$\square$]
  I can import the \texttt{math} module into my program.
\item[$\square$]
  I can use \texttt{math.pi} and \texttt{math.sqrt()}.
\item[$\square$]
  I can use \texttt{math.ceil()} to round up and \texttt{math.floor()}
  to round down.
\end{itemize}

\newpage

\section{Unit 3: Control Flow --- Decisions and
Loops}\label{unit-3-control-flow-decisions-and-loops-1}

\subsection{Overview}\label{overview-3}

This unit introduces how programs make decisions using \texttt{if}
statements and perform repetitive tasks with \texttt{for} loops, along
with fundamental string manipulation techniques.

\textbf{Objectives:}

\begin{itemize}
\tightlist
\item
  Control program flow with \texttt{if}, \texttt{elif}, and
  \texttt{else}.
\item
  Iterate over sequences using \texttt{for} loops.
\item
  Understand and manipulate strings.
\end{itemize}

\begin{center}\rule{0.5\linewidth}{0.5pt}\end{center}

\subsection{\texorpdfstring{Lesson 3.1: Boolean Expressions and
\texttt{if}
Statements}{Lesson 3.1: Boolean Expressions and if Statements}}\label{lesson-3.1-boolean-expressions-and-if-statements}

\subsection{Reading References}\label{reading-references-12}

\begin{itemize}
\tightlist
\item
  Zelle: Chapter 7, pages 165--173
\item
  Matthes: Chapter 5, pages 71--78
\end{itemize}

\subsection{Comparison and Boolean
Operators}\amn{\textbf{Accessibility Tip:} Pronounce comparison operators as words: \texttt{>} is ``greater than,'' \texttt{==} is ``equals equals'' (never just ``equals'' --- that is assignment). When reading a condition aloud: \texttt{if age >= 18 and verified:} becomes ``if age is greater than or equal to 18 and verified.''}
\subsection{Comparison and Boolean
Operators}\label{comparison-and-boolean-operators}

\begin{fullwidth}{\def\LTcaptype{none} % do not increment counter
\begin{longtable}[]{@{}llll@{}}
\toprule\noalign{}
Operator & Meaning & Example & Result \\
\midrule\noalign{}
\endhead
\bottomrule\noalign{}
\endlastfoot
\texttt{==} & Equal to & \texttt{5\ ==\ 5} & \texttt{True} \\
\texttt{!=} & Not equal to & \texttt{5\ !=\ 10} & \texttt{True} \\
\texttt{\textless{}} & Less than & \texttt{5\ \textless{}\ 10} &
\texttt{True} \\
\texttt{\textgreater{}} & Greater than & \texttt{10\ \textgreater{}\ 5}
& \texttt{True} \\
\texttt{\textless{}=} & Less than or equal to &
\texttt{5\ \textless{}=\ 5} & \texttt{True} \\
\texttt{\textgreater{}=} & Greater than or equal &
\texttt{10\ \textgreater{}=\ 5} & \texttt{True} \\
\end{longtable}
}
\end{fullwidth}

\begin{fullwidth}{\def\LTcaptype{none} % do not increment counter
\begin{longtable}[]{@{}
  >{\raggedright\arraybackslash}p{(\linewidth - 6\tabcolsep) * \real{0.2388}}
  >{\raggedright\arraybackslash}p{(\linewidth - 6\tabcolsep) * \real{0.4179}}
  >{\raggedright\arraybackslash}p{(\linewidth - 6\tabcolsep) * \real{0.2388}}
  >{\raggedright\arraybackslash}p{(\linewidth - 6\tabcolsep) * \real{0.1045}}@{}}
\toprule\noalign{}
\begin{minipage}[b]{\linewidth}\raggedright
Boolean Operator
\end{minipage} & \begin{minipage}[b]{\linewidth}\raggedright
Description
\end{minipage} & \begin{minipage}[b]{\linewidth}\raggedright
Example
\end{minipage} & \begin{minipage}[b]{\linewidth}\raggedright
Result
\end{minipage} \\
\midrule\noalign{}
\endhead
\bottomrule\noalign{}
\endlastfoot
\texttt{and} & True if both are true & \texttt{True\ and\ False} &
\texttt{False} \\
\texttt{or} & True if at least one is true & \texttt{True\ or\ False} &
\texttt{True} \\
\texttt{not} & Inverts the truth value & \texttt{not\ True} &
\texttt{False} \\
\end{longtable}
}
\end{fullwidth}

\subsection{\texorpdfstring{The \texttt{if}
Statement}{The if Statement}}\amn{\textbf{Accessibility Tip:} Every \texttt{if}, \texttt{elif}, \texttt{else}, \texttt{for}, \texttt{while}, \texttt{def}, and \texttt{class} line must end with a colon. When your screen reader navigates into an indented block, it announces the indent level. If it does not, your indentation announcement is not configured---stop and fix this before writing any more code.}
\subsection{\texorpdfstring{The \texttt{if}
Statement}{The if Statement}}\label{the-if-statement}

\begin{Shaded}
\begin{Highlighting}[]
age \OperatorTok{=} \CodeNumTok{15}
\KeywordTok{if} age \OperatorTok{\textgreater{}=} \CodeNumTok{18}:
    print(\StringTok{"You are old enough to vote."})
print(\StringTok{"Program finished."})  \CommentTok{\# This runs regardless}
\end{Highlighting}
\end{Shaded}

\textbf{Indentation defines which lines belong to the \texttt{if}
block.}

\subsection{Examples}\label{examples-6}

\textbf{Example 3.1.1: Simple Age Check}

\begin{Shaded}
\begin{Highlighting}[]
age \OperatorTok{=} \CodeNumTok{17}
\KeywordTok{if} age \OperatorTok{\textgreater{}=} \CodeNumTok{18}:
    print(\StringTok{"You are an adult."})
\KeywordTok{else}:
    print(\StringTok{"You are a minor."})
\end{Highlighting}
\end{Shaded}

\textbf{Example 3.1.2: Password Length Check}

\begin{Shaded}
\begin{Highlighting}[]
password \OperatorTok{=} \StringTok{"mysecret"}
\KeywordTok{if} len(password) \OperatorTok{\textless{}} \CodeNumTok{8}:
    print(\StringTok{"Weak password: Must be at least 8 characters."})
\end{Highlighting}
\end{Shaded}

\subsection{Activities}\label{activities-7}

\textbf{Activity 3.1.1: Number Comparison Game}

\begin{Shaded}
\begin{Highlighting}[]
num1 \OperatorTok{=} int(input(\StringTok{"Enter the first number: "}))
num2 \OperatorTok{=} int(input(\StringTok{"Enter the second number: "}))
 
\KeywordTok{if} num1 \OperatorTok{\textgreater{}} num2:
    print(\StringTok{f"}\StringTok{\{}num1\StringTok{\}}\StringTok{ is greater than }\StringTok{\{}num2\StringTok{\}}\StringTok{."})
\KeywordTok{if} num1 \OperatorTok{==} num2:
    print(\StringTok{"The numbers are equal."})
\KeywordTok{if} num1 \OperatorTok{\textless{}} num2:
    print(\StringTok{f"}\StringTok{\{}num1\StringTok{\}}\StringTok{ is less than }\StringTok{\{}num2\StringTok{\}}\StringTok{."})
\end{Highlighting}
\end{Shaded}

\subsection{Skills Checklist}\label{skills-checklist-12}

\begin{itemize}
\tightlist
\item[$\square$]
  I can identify \texttt{True} and \texttt{False} as Boolean values.
\item[$\square$]
  I can use comparison operators to create Boolean expressions.
\item[$\square$]
  I can write a basic \texttt{if} statement.
\item[$\square$]
  I can correctly indent the code block belonging to an \texttt{if}
  statement.
\end{itemize}

\begin{center}\rule{0.5\linewidth}{0.5pt}\end{center}

\subsection{\texorpdfstring{Lesson 3.2: Chained Conditionals
(\texttt{elif},
\texttt{else})}{Lesson 3.2: Chained Conditionals (elif, else)}}\label{lesson-3.2-chained-conditionals-elif-else}

\subsection{Reading References}\label{reading-references-13}

\begin{itemize}
\tightlist
\item
  Zelle: Chapter 7, pages 173--176
\item
  Matthes: Chapter 5, pages 79--84
\end{itemize}

\subsection{\texorpdfstring{\texttt{if-elif-else}
Structure}{if-elif-else Structure}}\label{if-elif-else-structure}

Python checks each condition in \textbf{sequential order}. As soon as a
condition is \texttt{True}, it runs that block and skips the rest.

\begin{Shaded}
\begin{Highlighting}[]
score \OperatorTok{=} \CodeNumTok{85}
 
\KeywordTok{if} score \OperatorTok{\textgreater{}=} \CodeNumTok{90}:
    print(\StringTok{"Grade: A"})
\KeywordTok{elif} score \OperatorTok{\textgreater{}=} \CodeNumTok{80}:
    print(\StringTok{"Grade: B"})
\KeywordTok{elif} score \OperatorTok{\textgreater{}=} \CodeNumTok{70}:
    print(\StringTok{"Grade: C"})
\KeywordTok{else}:
    print(\StringTok{"Grade: F"})
\end{Highlighting}
\end{Shaded}

Output: \texttt{Grade:\ B}

\subsection{Examples}\label{examples-7}

\textbf{Example 3.2.1: Stages of Life}

\begin{Shaded}
\begin{Highlighting}[]
age \OperatorTok{=} \CodeNumTok{25}
\KeywordTok{if} age \OperatorTok{\textless{}} \CodeNumTok{2}:
    print(\StringTok{"You are a baby."})
\KeywordTok{elif} age \OperatorTok{\textless{}} \CodeNumTok{4}:
    print(\StringTok{"You are a toddler."})
\KeywordTok{elif} age \OperatorTok{\textless{}} \CodeNumTok{13}:
    print(\StringTok{"You are a kid."})
\KeywordTok{elif} age \OperatorTok{\textless{}} \CodeNumTok{20}:
    print(\StringTok{"You are a teenager."})
\KeywordTok{else}:
    print(\StringTok{"You are an adult."})
\end{Highlighting}
\end{Shaded}

\subsection{Activities}\label{activities-8}

\textbf{Activity 3.2.1: Movie Ticket Pricing}

\begin{Shaded}
\begin{Highlighting}[]
age \OperatorTok{=} int(input(\StringTok{"Enter your age: "}))
 
\KeywordTok{if} age \OperatorTok{\textless{}=} \CodeNumTok{5}:
    print(\StringTok{"Ticket price: Free"})
\KeywordTok{elif} age \OperatorTok{\textless{}=} \CodeNumTok{12}:
    print(\StringTok{"Ticket price: $7"})
\KeywordTok{elif} age \OperatorTok{\textless{}=} \CodeNumTok{64}:
    print(\StringTok{"Ticket price: $12"})
\KeywordTok{else}:
    print(\StringTok{"Ticket price: $5"})
\end{Highlighting}
\end{Shaded}

\subsection{Skills Checklist}\label{skills-checklist-13}

\begin{itemize}
\tightlist
\item[$\square$]
  I can use an \texttt{else} statement to provide an alternative path.
\item[$\square$]
  I can use \texttt{elif} to check multiple conditions in order.
\item[$\square$]
  I can use \texttt{and}, \texttt{or}, and \texttt{not} operators to
  combine conditions.
\end{itemize}

\begin{center}\rule{0.5\linewidth}{0.5pt}\end{center}

\subsection{Lesson 3.3: Introduction to
Lists}\label{lesson-3.3-introduction-to-lists}

\subsection{Reading References}\label{reading-references-14}

\begin{itemize}
\tightlist
\item
  Zelle: Chapter 5, pages 105--106
\item
  Matthes: Chapter 3, pages 33--47
\end{itemize}

\subsection{Creating and Accessing
Lists}\amn{\textbf{Accessibility Tip:} Never use \texttt{print(my\_list)} to inspect a list. Use \texttt{for i, item in enumerate(my\_list, 1): print(f"\{i\}. \{item\}")} instead. This gives numbered, one-item-per-line output that is easy to navigate with arrow keys in the terminal.}
\subsection{Creating and Accessing
Lists}\label{creating-and-accessing-lists}

\begin{Shaded}
\begin{Highlighting}[]
colors \OperatorTok{=} [\StringTok{"red"}, \StringTok{"blue"}, \StringTok{"green"}, \StringTok{"yellow"}]
 
print(colors[\CodeNumTok{0}])   \CommentTok{\# "red"   (first item)}
print(colors[\OperatorTok{{-}}\CodeNumTok{1}])  \CommentTok{\# "yellow" (last item)}
print(len(colors)) \CommentTok{\# 4}
\end{Highlighting}
\end{Shaded}

\textbf{Zero-based indexing:} The first item is at index \texttt{0}.

\subsection{Modifying Lists}\label{modifying-lists}

\begin{Shaded}
\begin{Highlighting}[]
fruits \OperatorTok{=} [\StringTok{"apple"}, \StringTok{"banana"}, \StringTok{"cherry"}]
fruits[\CodeNumTok{1}] \OperatorTok{=} \StringTok{"grape"}  \CommentTok{\# Change "banana" to "grape"}
print(fruits)         \CommentTok{\# ["apple", "grape", "cherry"]}
\end{Highlighting}
\end{Shaded}

\subsection{Activities}\label{activities-9}

\textbf{Activity 3.3.1: My Top 3 Hobbies}

\begin{Shaded}
\begin{Highlighting}[]
hobbies \OperatorTok{=} [\StringTok{"reading"}, \StringTok{"coding"}, \StringTok{"hiking"}]
print(\StringTok{"My hobbies:"}, hobbies)
print(\StringTok{"My second hobby:"}, hobbies[\CodeNumTok{1}])
hobbies[\CodeNumTok{2}] \OperatorTok{=} \StringTok{"swimming"}
print(\StringTok{"Updated hobbies:"}, hobbies)
\end{Highlighting}
\end{Shaded}

\subsection{Skills Checklist}\label{skills-checklist-14}

\begin{itemize}
\tightlist
\item[$\square$]
  I can create a list in Python.
\item[$\square$]
  I can access individual items in a list using their index.
\item[$\square$]
  I can access list items from the end using negative indices.
\item[$\square$]
  I can change the value of an item at a specific index.
\end{itemize}

\begin{center}\rule{0.5\linewidth}{0.5pt}\end{center}

\subsection{\texorpdfstring{Lesson 3.4: Iterating with \texttt{for}
Loops}{Lesson 3.4: Iterating with for Loops}}\label{lesson-3.4-iterating-with-for-loops}

\subsection{Reading References}\label{reading-references-15}

\begin{itemize}
\tightlist
\item
  Zelle: Chapter 2, pages 34--36
\item
  Matthes: Chapter 4, pages 49--60
\end{itemize}

\amn{\textbf{Accessibility Tip:} Before running a \texttt{for} loop, trace it aloud: state what variable changes, what range or list it covers, and what the body does each iteration. This verbal trace catches errors before you hear many lines of wrong output.}
\subsection{\texorpdfstring{The \texttt{for} Loop
Structure}{The for Loop Structure}}\label{the-for-loop-structure}

\begin{Shaded}
\begin{Highlighting}[]
\KeywordTok{for} item \KeywordTok{in} sequence:
    \CommentTok{\# code that runs for each item}
\end{Highlighting}
\end{Shaded}

\subsection{Examples}\label{examples-8}

\textbf{Example 3.4.1: Looping Through a List}

\begin{Shaded}
\begin{Highlighting}[]
fruits \OperatorTok{=} [\StringTok{"apple"}, \StringTok{"banana"}, \StringTok{"cherry"}]
\KeywordTok{for} fruit \KeywordTok{in} fruits:
    print(\StringTok{f"I like }\StringTok{\{}fruit\StringTok{\}}\StringTok{."})
print(\StringTok{"Finished listing fruits."})
\end{Highlighting}
\end{Shaded}

\textbf{Example 3.4.2: Looping Through a String}

\begin{Shaded}
\begin{Highlighting}[]
word \OperatorTok{=} \StringTok{"Python"}
\KeywordTok{for} char \KeywordTok{in} word:
    print(char)
\end{Highlighting}
\end{Shaded}

\textbf{Example 3.4.3: Using \texttt{range()}}

\begin{Shaded}
\begin{Highlighting}[]
\KeywordTok{for} i \KeywordTok{in} range(\CodeNumTok{5}):      \CommentTok{\# 0, 1, 2, 3, 4}
    print(i)
 
\KeywordTok{for} i \KeywordTok{in} range(\CodeNumTok{2}, \CodeNumTok{8}):   \CommentTok{\# 2, 3, 4, 5, 6, 7}
    print(i)
 
\KeywordTok{for} i \KeywordTok{in} range(\CodeNumTok{0}, \CodeNumTok{10}, \CodeNumTok{2}):  \CommentTok{\# 0, 2, 4, 6, 8}
    print(i)
\end{Highlighting}
\end{Shaded}

\subsection{Activities}\label{activities-10}

\textbf{Activity 3.4.1: My Favorite Things}

\begin{Shaded}
\begin{Highlighting}[]
favorite\_things \OperatorTok{=} [\StringTok{"sunshine"}, \StringTok{"coffee"}, \StringTok{"books"}, \StringTok{"music"}, \StringTok{"mountains"}]
\KeywordTok{for} thing \KeywordTok{in} favorite\_things:
    print(\StringTok{f"I love }\StringTok{\{}thing\StringTok{\}}\StringTok{!"})
\end{Highlighting}
\end{Shaded}

\textbf{Activity 3.4.2: Countdown Timer}

\begin{Shaded}
\begin{Highlighting}[]
start \OperatorTok{=} int(input(\StringTok{"Enter a starting number: "}))
\KeywordTok{for} i \KeywordTok{in} range(start, \CodeNumTok{0}, \OperatorTok{{-}}\CodeNumTok{1}):
    print(i)
print(\StringTok{"Blast off!"})
\end{Highlighting}
\end{Shaded}

\subsection{Skills Checklist}\label{skills-checklist-15}

\begin{itemize}
\tightlist
\item[$\square$]
  I can explain the purpose of a \texttt{for} loop.
\item[$\square$]
  I can write a \texttt{for} loop to iterate over items in a list.
\item[$\square$]
  I can write a \texttt{for} loop to iterate over characters in a
  string.
\item[$\square$]
  I can use the \texttt{range()} function to control how many times a
  loop runs.
\end{itemize}

\begin{center}\rule{0.5\linewidth}{0.5pt}\end{center}

\subsection{Lesson 3.5: Basic String
Operations}\label{lesson-3.5-basic-string-operations}

\subsection{Reading References}\label{reading-references-16}

\begin{itemize}
\tightlist
\item
  Zelle: Chapter 5, pages 99--115
\item
  Matthes: Chapter 2, pages 19--25
\end{itemize}

\subsection{Common String Methods}\label{common-string-methods}

\begin{Shaded}
\begin{Highlighting}[]
message \OperatorTok{=} \StringTok{"  Hello Python!  "}
 
message.strip()          \CommentTok{\# "Hello Python!"   {-}{-}{-} removes whitespace}
message.upper()          \CommentTok{\# "  HELLO PYTHON!  "}
message.lower()          \CommentTok{\# "  hello python!  "}
message.replace(\StringTok{"Hello"}, \StringTok{"Hi"})  \CommentTok{\# "  Hi Python!  "}
len(message)             \CommentTok{\# 17}
\end{Highlighting}
\end{Shaded}

\textbf{Important:} String methods return a \textbf{new string} --- they
do not modify the original.

\subsection{Examples}\label{examples-9}

\textbf{Example 3.5.1: String Manipulation}

\begin{Shaded}
\begin{Highlighting}[]
message \OperatorTok{=} \StringTok{"  Hello Python  "}
print(\StringTok{f"Original: \textquotesingle{}}\StringTok{\{}message\StringTok{\}}\StringTok{\textquotesingle{}"})
print(\StringTok{f"Length: }\StringTok{\{}len(message)\StringTok{\}}\StringTok{"})
print(\StringTok{f"Uppercase: \textquotesingle{}}\StringTok{\{}message\StringTok{.}upper()\StringTok{\}}\StringTok{\textquotesingle{}"})
print(\StringTok{f"Stripped: \textquotesingle{}}\StringTok{\{}message\StringTok{.}strip()\StringTok{\}}\StringTok{\textquotesingle{}"})
\end{Highlighting}
\end{Shaded}

\textbf{Example 3.5.2: Combining Strings}

\begin{Shaded}
\begin{Highlighting}[]
first\_name \OperatorTok{=} \StringTok{"Alice"}
last\_name \OperatorTok{=} \StringTok{"Smith"}
full\_name \OperatorTok{=} first\_name \OperatorTok{+} \StringTok{" "} \OperatorTok{+} last\_name
print(\StringTok{f"Full name: }\StringTok{\{}full\_name\StringTok{\}}\StringTok{"})
\end{Highlighting}
\end{Shaded}

\subsection{Activities}\label{activities-11}

\textbf{Activity 3.5.1: Name Formatter}

\begin{Shaded}
\begin{Highlighting}[]
first\_name \OperatorTok{=} input(\StringTok{"Enter your first name: "})
last\_name \OperatorTok{=} input(\StringTok{"Enter your last name: "})
full\_name \OperatorTok{=} first\_name \OperatorTok{+} \StringTok{" "} \OperatorTok{+} last\_name
print(\StringTok{f"Full name: }\StringTok{\{}full\_name\StringTok{\}}\StringTok{"})
print(\StringTok{f"Uppercase: }\StringTok{\{}full\_name\StringTok{.}upper()\StringTok{\}}\StringTok{"})
print(\StringTok{f"Total characters: }\StringTok{\{}len(full\_name)\StringTok{\}}\StringTok{"})
\end{Highlighting}
\end{Shaded}

\subsection{Skills Checklist}\label{skills-checklist-16}

\begin{itemize}
\tightlist
\item[$\square$]
  I can combine strings using concatenation.
\item[$\square$]
  I can find the length of a string using \texttt{len()}.
\item[$\square$]
  I can convert a string to uppercase/lowercase.
\item[$\square$]
  I can remove extra spaces using \texttt{.strip()}.
\item[$\square$]
  I can replace parts of a string using \texttt{.replace()}.
\end{itemize}

\newpage

\section{Unit 4: Data Structures --- Lists and
Dictionaries}\label{unit-4-data-structures-lists-and-dictionaries-1}

\subsection{Overview}\label{overview-4}

This unit deepens understanding of lists and introduces dictionaries as
powerful mechanisms for organizing data.

\textbf{Objectives:}

\begin{itemize}
\tightlist
\item
  Add, remove, and organize list elements.
\item
  Work with tuples (immutable lists).
\item
  Use dictionaries for key-value pair storage.
\item
  Loop through dictionaries.
\end{itemize}

\begin{center}\rule{0.5\linewidth}{0.5pt}\end{center}

\subsection{Lesson 4.1: More on Lists}\label{lesson-4.1-more-on-lists}

\subsection{Adding and Removing
Elements}\label{adding-and-removing-elements}

\begin{Shaded}
\begin{Highlighting}[]
my\_list \OperatorTok{=} [\CodeNumTok{10}, \CodeNumTok{20}, \CodeNumTok{30}]
 
my\_list.append(\CodeNumTok{40})          \CommentTok{\# Add to end {-}\textgreater{} [10, 20, 30, 40]}
my\_list.insert(\CodeNumTok{1}, \CodeNumTok{15})       \CommentTok{\# Insert at index 1 {-}\textgreater{} [10, 15, 20, 30, 40]}
\KeywordTok{del} my\_list[\CodeNumTok{0}]              \CommentTok{\# Delete by index {-}\textgreater{} [15, 20, 30, 40]}
removed \OperatorTok{=} my\_list.pop()     \CommentTok{\# Remove and return last {-}\textgreater{} removed=40, list=[15, 20, 30]}
my\_list.remove(\CodeNumTok{20})          \CommentTok{\# Remove by value {-}\textgreater{} [15, 30]}
\end{Highlighting}
\end{Shaded}

\subsection{Example}\label{example}

\begin{Shaded}
\begin{Highlighting}[]
todo\_list \OperatorTok{=} [\StringTok{"Buy groceries"}, \StringTok{"Walk the dog"}]
 
todo\_list.append(\StringTok{"Call grandma"})
todo\_list.insert(\CodeNumTok{0}, \StringTok{"Wake up early"})
\KeywordTok{del} todo\_list[\CodeNumTok{1}]
completed\_task \OperatorTok{=} todo\_list.pop()
todo\_list.remove(\StringTok{"Walk the dog"})
print(todo\_list)
\end{Highlighting}
\end{Shaded}

\subsection{Skills Checklist}\label{skills-checklist-17}

\begin{itemize}
\tightlist
\item[$\square$]
  I can add an item to the end of a list using \texttt{.append()}.
\item[$\square$]
  I can insert an item at a specific position using \texttt{.insert()}.
\item[$\square$]
  I can remove an item by index using \texttt{del}.
\item[$\square$]
  I can remove an item and retrieve its value using \texttt{.pop()}.
\item[$\square$]
  I can remove the first occurrence of a value using \texttt{.remove()}.
\end{itemize}

\begin{center}\rule{0.5\linewidth}{0.5pt}\end{center}

\subsection{Lesson 4.2: Organizing
Lists}\label{lesson-4.2-organizing-lists}

\subsection{Sorting and Reversing}\label{sorting-and-reversing}

\begin{Shaded}
\begin{Highlighting}[]
numbers \OperatorTok{=} [\CodeNumTok{5}, \CodeNumTok{2}, \CodeNumTok{9}, \CodeNumTok{1}]
numbers.sort()                          \CommentTok{\# Permanent sort in{-}place}
print(numbers)                          \CommentTok{\# [1, 2, 5, 9]}
 
fruits \OperatorTok{=} [\StringTok{"orange"}, \StringTok{"apple"}, \StringTok{"banana"}]
sorted\_fruits \OperatorTok{=} sorted(fruits)          \CommentTok{\# Temporary sort {-}{-}{-} original unchanged}
print(fruits)                           \CommentTok{\# [\textquotesingle{}orange\textquotesingle{}, \textquotesingle{}apple\textquotesingle{}, \textquotesingle{}banana\textquotesingle{}]}
print(sorted\_fruits)                    \CommentTok{\# [\textquotesingle{}apple\textquotesingle{}, \textquotesingle{}banana\textquotesingle{}, \textquotesingle{}orange\textquotesingle{}]}
 
numbers.reverse()                       \CommentTok{\# Reverse in{-}place}
numbers.sort(reverse\OperatorTok{=}True)              \CommentTok{\# Sort descending}
print(len(numbers))                     \CommentTok{\# Length}
\end{Highlighting}
\end{Shaded}

\subsection{Skills Checklist}\label{skills-checklist-18}

\begin{itemize}
\tightlist
\item[$\square$]
  I can sort a list permanently using \texttt{.sort()}.
\item[$\square$]
  I can get a temporarily sorted version using \texttt{sorted()}.
\item[$\square$]
  I can reverse a list using \texttt{.reverse()}.
\item[$\square$]
  I can find how many items are in a list using \texttt{len()}.
\end{itemize}

\begin{center}\rule{0.5\linewidth}{0.5pt}\end{center}

\subsection{Lesson 4.3: Tuples --- Immutable
Lists}\label{lesson-4.3-tuples-immutable-lists}

\subsection{Tuples}\label{tuples}

Tuples are defined with parentheses and \textbf{cannot be changed} after
creation.

\begin{Shaded}
\begin{Highlighting}[]
coordinates \OperatorTok{=} (\CodeNumTok{10}, \CodeNumTok{20}, \CodeNumTok{30})
print(coordinates[\CodeNumTok{0}])     \CommentTok{\# 10}
print(coordinates[\OperatorTok{{-}}\CodeNumTok{1}])    \CommentTok{\# 30}
 
\CommentTok{\# This would cause a TypeError:}
\CommentTok{\# coordinates[1] = 25}
\end{Highlighting}
\end{Shaded}

\textbf{When to use tuples:} For data that should remain constant
(coordinates, RGB values, fixed settings).

\subsection{Skills Checklist}\label{skills-checklist-19}

\begin{itemize}
\tightlist
\item[$\square$]
  I can create a tuple in Python.
\item[$\square$]
  I can explain the main difference between a list and a tuple
  (mutability).
\item[$\square$]
  I can access items in a tuple by their index.
\end{itemize}

\begin{center}\rule{0.5\linewidth}{0.5pt}\end{center}

\subsection{Lesson 4.4: Introducing
Dictionaries}\label{lesson-4.4-introducing-dictionaries}

\subsection{Reading References}\label{reading-references-17}

\begin{itemize}
\tightlist
\item
  Zelle: Chapter 11, pages 307--309
\item
  Matthes: Chapter 6, pages 91--98
\end{itemize}

\subsection{Dictionary Basics}\label{dictionary-basics}\amn{\textbf{Accessibility Tip:} Use \texttt{my\_dict.get(key, "not found")} instead of \texttt{my\_dict[key]}. If the key does not exist, \texttt{my\_dict[key]} raises \texttt{KeyError} and crashes; \texttt{.get()} returns your default value silently. Crashes mid-session can be hard to detect with a screen reader.}
\subsection{Dictionary Basics}\label{dictionary-basics}

A dictionary stores \textbf{key-value pairs}. Keys must be unique and
immutable (strings, numbers, tuples). Values can be anything.

\begin{Shaded}
\begin{Highlighting}[]
student\_profile \OperatorTok{=} \{\}
    \StringTok{"name"}: \StringTok{"Emily"},
    \StringTok{"age"}: \CodeNumTok{13},
    \StringTok{"grade"}: \StringTok{"7th"},
    \StringTok{"favorite\_subject"}: \StringTok{"Science"}
\
 
print(student\_profile[\StringTok{"name"}])      \CommentTok{\# "Emily"}
student\_profile[\StringTok{"city"}] \OperatorTok{=} \StringTok{"Springfield"}   \CommentTok{\# Add new key}
student\_profile[\StringTok{"age"}] \OperatorTok{=} \CodeNumTok{14}              \CommentTok{\# Modify existing key}
\KeywordTok{del} student\_profile[\StringTok{"grade"}]             \CommentTok{\# Remove key}
\end{Highlighting}
\end{Shaded}

\subsection{Examples}\label{examples-10}

\textbf{Example 4.4.1: Book Information Dictionary}

\begin{Shaded}
\begin{Highlighting}[]
my\_book \OperatorTok{=} \{\}
    \StringTok{"title"}: \StringTok{"The Hitchhiker\textquotesingle{}s Guide to the Galaxy"},
    \StringTok{"author"}: \StringTok{"Douglas Adams"},
    \StringTok{"publication\_year"}: \CodeNumTok{1979},
    \StringTok{"genre"}: \StringTok{"Science Fiction"}
\
 
print(\StringTok{f"Title: }\StringTok{\{}my\_book[\StringTok{\textquotesingle{}title\textquotesingle{}}]\StringTok{\}}\StringTok{"})
print(\StringTok{f"Author: }\StringTok{\{}my\_book[\StringTok{\textquotesingle{}author\textquotesingle{}}]\StringTok{\}}\StringTok{"})
 
my\_book[\StringTok{"number\_of\_pages"}] \OperatorTok{=} \CodeNumTok{193}
my\_book[\StringTok{"genre"}] \OperatorTok{=} \StringTok{"Comedy Science Fiction"}
print(my\_book)
\end{Highlighting}
\end{Shaded}

\subsection{Skills Checklist}\label{skills-checklist-20}

\begin{itemize}
\tightlist
\item[$\square$]
  I can create a dictionary in Python.
\item[$\square$]
  I can explain what a key-value pair is.
\item[$\square$]
  I can access a value in a dictionary using its key.
\item[$\square$]
  I can add, modify, and remove key-value pairs.
\end{itemize}

\begin{center}\rule{0.5\linewidth}{0.5pt}\end{center}

\subsection{Lesson 4.5: Looping Through
Dictionaries}\label{lesson-4.5-looping-through-dictionaries}

\subsection{Dictionary Iteration
Methods}\label{dictionary-iteration-methods}

\begin{Shaded}
\begin{Highlighting}[]
inventory \OperatorTok{=} \{\}\StringTok{"apples"}: \CodeNumTok{50}, \StringTok{"bananas"}: \CodeNumTok{30}, \StringTok{"oranges"}: \CodeNumTok{40}\
 
\CommentTok{\# Loop through keys (default)}
\KeywordTok{for} item\_name \KeywordTok{in} inventory:
    print(item\_name)
 
\CommentTok{\# Loop through values only}
\KeywordTok{for} count \KeywordTok{in} inventory.values():
    print(count)
 
\CommentTok{\# Loop through key{-}value pairs}
\KeywordTok{for} item\_name, count \KeywordTok{in} inventory.items():
    print(\StringTok{f"}\StringTok{\{}item\_name\StringTok{\}}\StringTok{: }\StringTok{\{}count\StringTok{\}}\StringTok{ in stock"})
\end{Highlighting}
\end{Shaded}

\subsection{Skills Checklist}\label{skills-checklist-21}

\begin{itemize}
\tightlist
\item[$\square$]
  I can loop through all the keys in a dictionary.
\item[$\square$]
  I can loop through all the values in a dictionary.
\item[$\square$]
  I can loop through both keys and values at the same time.
\end{itemize}

\begin{center}\rule{0.5\linewidth}{0.5pt}\end{center}

\section{\texorpdfstring{Unit 5: User Input and \texttt{while}
Loops}{Unit 5: User Input and while Loops}}\label{unit-5-user-input-and-while-loops-1}

\subsection{Overview}\label{overview-5}

This unit focuses on obtaining diverse input from users and controlling
program execution through indefinite loops.

\textbf{Objectives:}

\begin{itemize}
\tightlist
\item
  Implement indefinite loops using \texttt{while} statements.
\item
  Control \texttt{while} loop execution with \texttt{break} and
  \texttt{continue}.
\item
  Avoid infinite loops.
\end{itemize}

\begin{center}\rule{0.5\linewidth}{0.5pt}\end{center}

\subsection{\texorpdfstring{Lesson 5.1: How the \texttt{input()}
Function
Works}{Lesson 5.1: How the input() Function Works}}\label{lesson-5.1-how-the-input-function-works}

\subsection{Key Reminder}\label{key-reminder}

\texttt{input()} \textbf{always returns a string}. Convert to numeric
types as needed.

\begin{Shaded}
\begin{Highlighting}[]
age\_input \OperatorTok{=} input(\StringTok{"How old are you? "})
print(\StringTok{f"Type: }\StringTok{\{}type(age\_input)\StringTok{\}}\StringTok{"})    \CommentTok{\# \textless{}class \textquotesingle{}str\textquotesingle{}\textgreater{}}
 
age\_int \OperatorTok{=} int(age\_input)
print(\StringTok{f"In 5 years: }\StringTok{\{}age\_int \OperatorTok{+} \CodeNumTok{5}\StringTok{\}}\StringTok{"})
\end{Highlighting}
\end{Shaded}

\subsection{Skills Checklist}\label{skills-checklist-22}

\begin{itemize}
\tightlist
\item[$\square$]
  I can effectively use \texttt{input()} to get user data.
\item[$\square$]
  I can confirm that \texttt{input()}'s return value is always a string.
\item[$\square$]
  I can convert user input to a numeric type when necessary.
\end{itemize}

\begin{center}\rule{0.5\linewidth}{0.5pt}\end{center}

\subsection{\texorpdfstring{Lesson 5.2: The \texttt{while} Loop in
Action}{Lesson 5.2: The while Loop in Action}}\label{lesson-5.2-the-while-loop-in-action}

\subsection{\texorpdfstring{\texttt{for} vs.~\texttt{while} Loop
Comparison}{for vs.~while Loop Comparison}}\amn{\textbf{Accessibility Tip:} Use \texttt{for} when you know the number of iterations in advance. Use \texttt{while} when you are waiting for a condition to change (user input, file end, sensor value). Every \texttt{while} loop must update the condition variable inside the loop body or it will run forever.}
\subsection{\texorpdfstring{\texttt{for} vs.~\texttt{while} Loop
Comparison}{for vs.~while Loop Comparison}}\label{for-vs.-while-loop-comparison}

\begin{fullwidth}{\def\LTcaptype{none} % do not increment counter
\begin{longtable}[]{@{}
  >{\raggedright\arraybackslash}p{(\linewidth - 4\tabcolsep) * \real{0.2128}}
  >{\raggedright\arraybackslash}p{(\linewidth - 4\tabcolsep) * \real{0.3404}}
  >{\raggedright\arraybackslash}p{(\linewidth - 4\tabcolsep) * \real{0.4468}}@{}}
\toprule\noalign{}
\begin{minipage}[b]{\linewidth}\raggedright
Feature
\end{minipage} & \begin{minipage}[b]{\linewidth}\raggedright
\texttt{for} Loop
\end{minipage} & \begin{minipage}[b]{\linewidth}\raggedright
\texttt{while} Loop
\end{minipage} \\
\midrule\noalign{}
\endhead
\bottomrule\noalign{}
\endlastfoot
\textbf{Purpose} & Definite iteration (known count) & Indefinite
iteration (condition-based) \\
\textbf{When to Use} & Iterating over lists, ranges & Repeating until a
condition is met \\
\textbf{Control Variable} & Automatically updated & Must be manually
updated \\
\textbf{Risk} & Less prone to infinite loops & High risk if condition
never becomes false \\
\end{longtable}
}
\end{fullwidth}

\subsection{Examples}\label{examples-11}

\textbf{Example 5.2.1: Simple Counter}

\begin{Shaded}
\begin{Highlighting}[]
count \OperatorTok{=} \CodeNumTok{1}
\KeywordTok{while} count \OperatorTok{\textless{}=} \CodeNumTok{5}:
    print(count)
    count \OperatorTok{+=} \CodeNumTok{1}
print(\StringTok{"Counting finished!"})
\end{Highlighting}
\end{Shaded}

\textbf{Example 5.2.2: User Input Loop (Sentinel Value)}

\begin{Shaded}
\begin{Highlighting}[]
message \OperatorTok{=} \StringTok{""}
\KeywordTok{while} message \OperatorTok{!=} \StringTok{\textquotesingle{}quit\textquotesingle{}}:
    message \OperatorTok{=} input(\StringTok{"Enter a message, or type \textquotesingle{}quit\textquotesingle{} to exit: "})
    \KeywordTok{if} message \OperatorTok{!=} \StringTok{\textquotesingle{}quit\textquotesingle{}}:
        print(\StringTok{f"You said: }\StringTok{\{}message\StringTok{\}}\StringTok{"})
print(\StringTok{"Goodbye!"})
\end{Highlighting}
\end{Shaded}

\subsection{Skills Checklist}\label{skills-checklist-23}

\begin{itemize}
\tightlist
\item[$\square$]
  I can explain the difference between a \texttt{for} loop and a
  \texttt{while} loop.
\item[$\square$]
  I can write a basic \texttt{while} loop.
\item[$\square$]
  I can make sure my \texttt{while} loop has a way to stop.
\end{itemize}

\begin{center}\rule{0.5\linewidth}{0.5pt}\end{center}

\subsection{\texorpdfstring{Lesson 5.3: Using Flags and \texttt{break}
to Exit
Loops}{Lesson 5.3: Using Flags and break to Exit Loops}}\label{lesson-5.3-using-flags-and-break-to-exit-loops}

\subsection{Flag Variable}\label{flag-variable}

\begin{Shaded}
\begin{Highlighting}[]
active \OperatorTok{=} True
\KeywordTok{while} active:
    message \OperatorTok{=} input(\StringTok{"Say something (or \textquotesingle{}quit\textquotesingle{} to stop): "})
    \KeywordTok{if} message \OperatorTok{==} \StringTok{\textquotesingle{}quit\textquotesingle{}}:
        active \OperatorTok{=} False
    \KeywordTok{else}:
        print(\StringTok{f"You said: }\StringTok{\{}message\StringTok{\}}\StringTok{"})
\end{Highlighting}
\end{Shaded}

\subsection{\texorpdfstring{\texttt{break}
Statement}{break Statement}}\label{break-statement}

\begin{Shaded}
\begin{Highlighting}[]
\KeywordTok{while} True:
    city \OperatorTok{=} input(\StringTok{"Enter a city (or \textquotesingle{}done\textquotesingle{} to finish): "})
    \KeywordTok{if} city \OperatorTok{==} \StringTok{\textquotesingle{}done\textquotesingle{}}:
        \KeywordTok{break}
    print(\StringTok{f"I\textquotesingle{}ve heard }\StringTok{\{}city\StringTok{\}}\StringTok{ is a great place!"})
\end{Highlighting}
\end{Shaded}

\subsection{Skills Checklist}\label{skills-checklist-24}

\begin{itemize}
\tightlist
\item[$\square$]
  I can use a Boolean ``flag'' variable to control a \texttt{while}
  loop.
\item[$\square$]
  I can use the \texttt{break} statement to exit a loop early.
\end{itemize}

\begin{center}\rule{0.5\linewidth}{0.5pt}\end{center}

\subsection{\texorpdfstring{Lesson 5.4: Using \texttt{continue} in
Loops}{Lesson 5.4: Using continue in Loops}}\label{lesson-5.4-using-continue-in-loops}

\subsection{\texorpdfstring{\texttt{continue} --- Skip the Current
Iteration}{continue --- Skip the Current Iteration}}\label{continue-skip-the-current-iteration}

\begin{Shaded}
\begin{Highlighting}[]
number \OperatorTok{=} \CodeNumTok{0}
\KeywordTok{while} number \OperatorTok{\textless{}} \CodeNumTok{10}:
    number \OperatorTok{+=} \CodeNumTok{1}
    \KeywordTok{if} number \OperatorTok{\%} \CodeNumTok{2} \OperatorTok{==} \CodeNumTok{0}:
        \KeywordTok{continue}        \CommentTok{\# Skip even numbers}
    print(number)       \CommentTok{\# Prints 1, 3, 5, 7, 9}
\end{Highlighting}
\end{Shaded}

\subsection{\texorpdfstring{\texttt{continue}
vs.~\texttt{break}}{continue vs.~break}}\label{continue-vs.-break}

\begin{itemize}
\tightlist
\item
  \texttt{break} --- \textbf{exits the loop entirely}
\item
  \texttt{continue} --- \textbf{skips only the current iteration}, loop
  continues
\end{itemize}

\subsection{Skills Checklist}\label{skills-checklist-25}

\begin{itemize}
\tightlist
\item[$\square$]
  I can use \texttt{continue} to skip the rest of the current loop
  iteration.
\item[$\square$]
  I can identify scenarios where \texttt{continue} would be useful.
\end{itemize}

\begin{center}\rule{0.5\linewidth}{0.5pt}\end{center}

\subsection{Lesson 5.5: Avoiding Infinite
Loops}\label{lesson-5.5-avoiding-infinite-loops}

An \textbf{infinite loop} never ends because its condition always
remains \texttt{True}.

\textbf{Common causes:}

\begin{enumerate}
\def\labelenumi{\arabic{enumi}.}
\tightlist
\item
  Forgetting to update the loop variable.
\item
  Updating the variable in the wrong direction.
\item
  Using an incorrect logical condition.
\end{enumerate}

\begin{Shaded}
\begin{Highlighting}[]
\CommentTok{\# INFINITE LOOP {-}{-}{-} Don\textquotesingle{}t run!}
count \OperatorTok{=} \CodeNumTok{1}
\KeywordTok{while} count \OperatorTok{\textless{}=} \CodeNumTok{5}:
    print(count)
    \CommentTok{\# Missing: count += 1}
 
\CommentTok{\# CORRECT}
count \OperatorTok{=} \CodeNumTok{1}
\KeywordTok{while} count \OperatorTok{\textless{}=} \CodeNumTok{5}:
    print(count)
    count \OperatorTok{+=} \CodeNumTok{1}
\end{Highlighting}
\end{Shaded}


\sidenote{\texttt{Ctrl+C} sends a \texttt{SIGINT}}\amn{\textbf{Accessibility Tip:} When a program hangs or enters an infinite loop, press \texttt{Ctrl+C} in the terminal to stop it immediately. If the screen reader does not announce the new prompt after stopping, press Enter once to restore focus. Then use the Up arrow to recall the last command.}
\sidenote{\texttt{Ctrl+C} sends a \texttt{SIGINT} (interrupt signal) to the running process. On Windows terminals with some screen readers, focus may not return to the prompt automatically after \texttt{Ctrl+C}---press Enter once if the cursor seems unresponsive.}

\subsection{Skills Checklist}\label{skills-checklist-26}

\begin{itemize}
\tightlist
\item[$\square$]
  I can define what an infinite loop is.
\item[$\square$]
  I can identify common causes of infinite loops.
\item[$\square$]
  I know how to stop a running program in an infinite loop.
\end{itemize}

\begin{center}\rule{0.5\linewidth}{0.5pt}\end{center}

\section{Unit 6: Functions}\label{unit-6-functions-1}

\subsection{Overview}\label{overview-6}

This unit introduces functions, enabling organization of code into
reusable blocks, passing arguments, and returning values.

\textbf{Objectives:}

\begin{itemize}
\tightlist
\item
  Define and call functions.
\item
  Pass arguments to functions.
\item
  Return values from functions.
\item
  Understand the concept of scope.
\end{itemize}

\begin{center}\rule{0.5\linewidth}{0.5pt}\end{center}

\subsection{Lesson 6.1: Defining
Functions}\label{lesson-6.1-defining-functions}

\subsection{Reading References}\label{reading-references-18}

\begin{itemize}
\tightlist
\item
  Zelle: Chapter 6, pages 137--142
\item
  Matthes: Chapter 8, pages 129--130
\end{itemize}

\subsection{Function Syntax}\label{function-syntax}\amn{\textbf{Accessibility Tip:} Every function you write should have a docstring as its first line: \texttt{"""Describe what this function does."""}. When IntelliSense shows the function, it reads the docstring aloud. This is how you document a function for both yourself and anyone else who uses it.}
\subsection{Function Syntax}\label{function-syntax}

\begin{Shaded}
\begin{Highlighting}[]
\KeywordTok{def} function\_name(parameter1, parameter2):
    \CommentTok{"""Docstring describing what the function does."""}
    \CommentTok{\# Function body}
    print(\StringTok{f"Hello, }\StringTok{\{}parameter1\StringTok{\}}\StringTok{!"})
\end{Highlighting}
\end{Shaded}

\subsection{Examples}\label{examples-12}

\textbf{Example 6.1.1: Simple Greeting Function}

\begin{Shaded}
\begin{Highlighting}[]
\KeywordTok{def} greet\_user():
    \CommentTok{"""Prints a simple greeting message."""}
    print(\StringTok{"Hello, Python programmer!"})
 
greet\_user()   \CommentTok{\# Call the function}
greet\_user()   \CommentTok{\# Call it again}
\end{Highlighting}
\end{Shaded}

\textbf{Example 6.1.2: Function with a Parameter}

\begin{Shaded}
\begin{Highlighting}[]
\KeywordTok{def} greet\_name(name):
    \CommentTok{"""Prints a personalized greeting."""}
    print(\StringTok{f"Hello, }\StringTok{\{}name\StringTok{\}}\StringTok{!"})
 
greet\_name(\StringTok{"Alice"})
greet\_name(\StringTok{"Bob"})
\end{Highlighting}
\end{Shaded}

\subsection{Skills Checklist}\label{skills-checklist-27}

\begin{itemize}
\tightlist
\item[$\square$]
  I can explain why functions are useful.
\item[$\square$]
  I can define a simple function in Python.
\item[$\square$]
  I can use parameters to specify function inputs.
\item[$\square$]
  I can call a function by its name.
\end{itemize}

\begin{center}\rule{0.5\linewidth}{0.5pt}\end{center}

\subsection{Lesson 6.2: Passing
Arguments}\label{lesson-6.2-passing-arguments}

\subsection{Positional vs.~Keyword
Arguments}\label{positional-vs.-keyword-arguments}

\begin{Shaded}
\begin{Highlighting}[]
\KeywordTok{def} describe\_animal(animal\_type, animal\_name):
    print(\StringTok{f"I have a }\StringTok{\{}animal\_type\StringTok{\}}\StringTok{ named }\StringTok{\{}animal\_name\StringTok{\}}\StringTok{."})
 
\CommentTok{\# Positional arguments (order matters)}
describe\_animal(\StringTok{"dog"}, \StringTok{"Buddy"})
 
\CommentTok{\# Keyword arguments (order doesn\textquotesingle{}t matter)}
describe\_animal(animal\_name\OperatorTok{=}\StringTok{"Max"}, animal\_type\OperatorTok{=}\StringTok{"cat"})
\end{Highlighting}
\end{Shaded}

\subsection{Default Values}\label{default-values}

\begin{Shaded}
\begin{Highlighting}[]
\KeywordTok{def} make\_shirt(size, message\OperatorTok{=}\StringTok{"I love Python"}):
    print(\StringTok{f"Making a }\StringTok{\{}size\StringTok{\}}\StringTok{{-}sized shirt: \textquotesingle{}}\StringTok{\{}message\StringTok{\}}\StringTok{\textquotesingle{}"})
 
make\_shirt(\StringTok{"large"})                       \CommentTok{\# Uses default message}
make\_shirt(\StringTok{"medium"}, \StringTok{"Hello World"})       \CommentTok{\# Custom message}
\end{Highlighting}
\end{Shaded}

\subsection{Skills Checklist}\label{skills-checklist-28}

\begin{itemize}
\tightlist
\item[$\square$]
  I can differentiate between a parameter and an argument.
\item[$\square$]
  I can pass arguments using positional order.
\item[$\square$]
  I can pass arguments using keyword arguments.
\end{itemize}

\begin{center}\rule{0.5\linewidth}{0.5pt}\end{center}

\subsection{Lesson 6.3: Return Values}\label{lesson-6.3-return-values}

\subsection{\texorpdfstring{The \texttt{return}
Statement}{The return Statement}}\label{the-return-statement}

\begin{Shaded}
\begin{Highlighting}[]
\KeywordTok{def} add\_numbers(num1, num2):
    \CommentTok{"""Adds two numbers and returns the result."""}
    total \OperatorTok{=} num1 \OperatorTok{+} num2
    \KeywordTok{return} total
 
result \OperatorTok{=} add\_numbers(\CodeNumTok{5}, \CodeNumTok{3})
print(\StringTok{f"The sum is: }\StringTok{\{}result\StringTok{\}}\StringTok{"})   \CommentTok{\# 8}
\end{Highlighting}
\end{Shaded}


\marginnote{This distinction trips up nearly every beginner. A function that \texttt{print()}s its result}\amn{\textbf{Accessibility Tip:} When testing a function, call it and print the result: \texttt{print(my\_function(5))}. If the function uses \texttt{return}, you must capture and print the value yourself. If the function uses \texttt{print()} internally, calling it will display output automatically---but the value cannot be reused.}
\marginnote{This distinction trips up nearly every beginner. A function that \texttt{print()}s its result cannot be composed with other functions---the value disappears into the terminal. A function that \texttt{return}s its result can be stored, passed, or tested.}

\subsection{Skills Checklist}\label{skills-checklist-29}

\begin{itemize}
\tightlist
\item[$\square$]
  I can use \texttt{return} to send a value back from a function.
\item[$\square$]
  I can explain the difference between \texttt{print()} and
  \texttt{return}.
\item[$\square$]
  I can store the returned value of a function in a variable.
\end{itemize}

\begin{center}\rule{0.5\linewidth}{0.5pt}\end{center}

\subsection{Lesson 6.4: Functions with Lists and
Dictionaries}\label{lesson-6.4-functions-with-lists-and-dictionaries}


\sidenote{Mutability is one of Python's most common sources of subtle bugs. If you want a function to work on a \emph{copy} of a list, pass \texttt{my\_list[:]} or \texttt{list(my\_list)} instead of the original. This is especially important in tests where the same list is reused across multiple test cases.}

\begin{Shaded}
\begin{Highlighting}[]
\KeywordTok{def} add\_new\_user(user\_list, new\_user):
    user\_list.append(new\_user)
 
users \OperatorTok{=} [\StringTok{\textquotesingle{}Alice\textquotesingle{}}, \StringTok{\textquotesingle{}Bob\textquotesingle{}}]
add\_new\_user(users, \StringTok{\textquotesingle{}Charlie\textquotesingle{}})
print(users)   \CommentTok{\# [\textquotesingle{}Alice\textquotesingle{}, \textquotesingle{}Bob\textquotesingle{}, \textquotesingle{}Charlie\textquotesingle{}]}
\end{Highlighting}
\end{Shaded}

To prevent modification, pass a copy:

\begin{Shaded}
\begin{Highlighting}[]
add\_new\_user(users[:], \StringTok{\textquotesingle{}Charlie\textquotesingle{}})   \CommentTok{\# Original unchanged}
\end{Highlighting}
\end{Shaded}

\subsection{Skills Checklist}\label{skills-checklist-30}

\begin{itemize}
\tightlist
\item[$\square$]
  I can pass a list or dictionary as an argument to a function.
\item[$\square$]
  I understand that modifying a list inside a function affects the
  original.
\item[$\square$]
  I can create a copy of a list to avoid unintended modifications.
\end{itemize}

\begin{center}\rule{0.5\linewidth}{0.5pt}\end{center}

\subsection{Lesson 6.5: Arbitrary Number of
Arguments}\label{lesson-6.5-arbitrary-number-of-arguments}

\subsection{\texorpdfstring{\texttt{*args} --- Arbitrary Positional
Arguments}{*args --- Arbitrary Positional Arguments}}\label{args-arbitrary-positional-arguments}

\begin{Shaded}
\begin{Highlighting}[]
\KeywordTok{def} make\_pizza(size, \OperatorTok{*}toppings):
    print(\StringTok{f"Making a }\StringTok{\{}size\StringTok{\}}\StringTok{{-}inch pizza with:"})
    \KeywordTok{for} topping \KeywordTok{in} toppings:
        print(\StringTok{f"  {-} }\StringTok{\{}topping\StringTok{\}}\StringTok{"})
 
make\_pizza(\CodeNumTok{12}, \StringTok{\textquotesingle{}pepperoni\textquotesingle{}})
make\_pizza(\CodeNumTok{16}, \StringTok{\textquotesingle{}mushrooms\textquotesingle{}}, \StringTok{\textquotesingle{}green peppers\textquotesingle{}}, \StringTok{\textquotesingle{}extra cheese\textquotesingle{}})
\end{Highlighting}
\end{Shaded}

\subsection{\texorpdfstring{\texttt{**kwargs} --- Arbitrary Keyword
Arguments}{**kwargs --- Arbitrary Keyword Arguments}}\label{kwargs-arbitrary-keyword-arguments}

\begin{Shaded}
\begin{Highlighting}[]
\KeywordTok{def} build\_profile(first, last, \OperatorTok{**}user\_info):
    user\_info[\StringTok{\textquotesingle{}first\_name\textquotesingle{}}] \OperatorTok{=} first
    user\_info[\StringTok{\textquotesingle{}last\_name\textquotesingle{}}] \OperatorTok{=} last
    \KeywordTok{return} user\_info
 
profile \OperatorTok{=} build\_profile(\StringTok{\textquotesingle{}albert\textquotesingle{}}, \StringTok{\textquotesingle{}einstein\textquotesingle{}},
                        location\OperatorTok{=}\StringTok{\textquotesingle{}princeton\textquotesingle{}},
                        field\OperatorTok{=}\StringTok{\textquotesingle{}physics\textquotesingle{}})
print(profile)
\end{Highlighting}
\end{Shaded}

\subsection{Skills Checklist}\label{skills-checklist-31}

\begin{itemize}
\tightlist
\item[$\square$]
  I can define a function using \texttt{*args} for arbitrary positional
  arguments.
\item[$\square$]
  I can define a function using \texttt{**kwargs} for arbitrary keyword
  arguments.
\end{itemize}

\newpage

\section{Unit 7: Classes and Object-Oriented
Programming}\label{unit-7-classes-and-object-oriented-programming-1}

\subsection{Overview}\label{overview-7}

This unit introduces the fundamental concepts of Object-Oriented
Programming (OOP) in Python, focusing on defining and using classes,
attributes, methods, and inheritance.

\textbf{Objectives:}

\begin{itemize}
\tightlist
\item
  Define and instantiate classes.
\item
  Work with attributes and methods.
\item
  Understand inheritance and composition.
\end{itemize}

\begin{center}\rule{0.5\linewidth}{0.5pt}\end{center}

\subsection{Lesson 7.1: Why Learn About
Classes?}\label{lesson-7.1-why-learn-about-classes}

\subsection{Reading References}\label{reading-references-19}

\begin{itemize}
\tightlist
\item
  Zelle: Chapter 4, pages 65--67; Chapter 10, pages 247--248
\item
  Matthes: Chapter 9, page 157
\end{itemize}

\subsection{Classes as Blueprints}\label{classes-as-blueprints}

A \textbf{class} is a blueprint or template for creating objects. An
\textbf{object} is a concrete instance created from that blueprint.

\textbf{Analogy:} A ``Car'' class defines that every car has a
\texttt{color} and a \texttt{speed} (attributes) and can
\texttt{accelerate()} and \texttt{brake()} (methods). Your red sedan and
your friend's blue pickup are two separate \textbf{objects} of the same
``Car'' class.

\subsection{Key OOP Concepts}\label{key-oop-concepts}

\begin{itemize}
\tightlist
\item
  \textbf{Encapsulation} --- Bundling data and methods into a single
  unit.
\item
  \textbf{Abstraction} --- Hiding complexity from the outside.
\item
  \textbf{Inheritance} --- New class reuses functionality of existing
  class.
\item
  \textbf{Polymorphism} --- Same method name behaves differently for
  different classes.
\end{itemize}

\subsection{Skills Checklist}\label{skills-checklist-32}

\begin{itemize}
\tightlist
\item[$\square$]
  I can explain the basic idea of Object-Oriented Programming.
\item[$\square$]
  I can describe a class as a blueprint and an object as an instance.
\end{itemize}

\begin{center}\rule{0.5\linewidth}{0.5pt}\end{center}

\subsection{Lesson 7.2: Creating and Using a
Class}\label{lesson-7.2-creating-and-using-a-class}

\subsection{Reading References}\label{reading-references-20}

\begin{itemize}
\tightlist
\item
  Zelle: Chapter 10, pages 248--254
\item
  Matthes: Chapter 9, pages 158--162
\end{itemize}

\subsection{Class Syntax}\label{class-syntax}

\begin{Shaded}
\begin{Highlighting}[]
\KeywordTok{class} Dog:
    \CommentTok{"""A simple attempt to model a dog."""}
 
    \KeywordTok{def} \CodeFuncTok{\_\_init\_\_}(self, name, age):
        \CommentTok{"""Initialize name and age attributes."""}
        self.name \OperatorTok{=} name
        self.age \OperatorTok{=} age
        print(\StringTok{f"A new dog named }\StringTok{\{}self\StringTok{.}name\StringTok{\}}\StringTok{ has been created!"})
 
\CommentTok{\# Create instances}
my\_dog \OperatorTok{=} Dog(\StringTok{"Buddy"}, \CodeNumTok{3})
your\_dog \OperatorTok{=} Dog(\StringTok{"Lucy"}, \CodeNumTok{1})
 
print(\StringTok{f"My dog\textquotesingle{}s name is }\StringTok{\{}my\_dog\StringTok{.}name\StringTok{\}}\StringTok{."})
print(\StringTok{f"Your dog\textquotesingle{}s name is }\StringTok{\{}your\_dog\StringTok{.}name\StringTok{\}}\StringTok{."})
\end{Highlighting}
\end{Shaded}

The \texttt{\_\_init\_\_()} method runs automatically when a new object
is created. \texttt{self} refers to the specific instance being created.

\subsection{Skills Checklist}\label{skills-checklist-33}

\begin{itemize}
\tightlist
\item[$\square$]
  I can define a basic Python class.
\item[$\square$]
  I can define the \texttt{\_\_init\_\_} method within a class.
\item[$\square$]
  I understand the purpose of the \texttt{self} parameter.
\item[$\square$]
  I can create an instance (object) of a class.
\end{itemize}

\begin{center}\rule{0.5\linewidth}{0.5pt}\end{center}

\subsection{Lesson 7.3: Working with Attributes and
Methods}\label{lesson-7.3-working-with-attributes-and-methods}

\subsection{Reading References}\label{reading-references-21}

\begin{itemize}
\tightlist
\item
  Matthes: Chapter 9, pages 162--166
\end{itemize}

\subsection{Calling Methods}\label{calling-methods}

\begin{Shaded}
\begin{Highlighting}[]
\KeywordTok{class} Dog:
    \KeywordTok{def} \CodeFuncTok{\_\_init\_\_}(self, name, age):
        self.name \OperatorTok{=} name
        self.age \OperatorTok{=} age
 
    \KeywordTok{def} sit(self):
        \CommentTok{"""Simulate a dog sitting."""}
        print(\StringTok{f"}\StringTok{\{}self\StringTok{.}name\StringTok{\}}\StringTok{ is now sitting."})
 
    \KeywordTok{def} roll\_over(self):
        \CommentTok{"""Simulate a dog rolling over."""}
        print(\StringTok{f"}\StringTok{\{}self\StringTok{.}name\StringTok{\}}\StringTok{ rolled over!"})
 
    \KeywordTok{def} get\_older(self, years\_to\_add):
        \CommentTok{"""Add years to the dog\textquotesingle{}s age."""}
        self.age \OperatorTok{+=} years\_to\_add
        print(\StringTok{f"}\StringTok{\{}self\StringTok{.}name\StringTok{\}}\StringTok{ is now }\StringTok{\{}self\StringTok{.}age\StringTok{\}}\StringTok{ years old."})
 
my\_dog \OperatorTok{=} Dog(\StringTok{"Buddy"}, \CodeNumTok{3})
my\_dog.sit()
my\_dog.roll\_over()
my\_dog.age \OperatorTok{=} \CodeNumTok{4}          \CommentTok{\# Direct attribute modification}
my\_dog.get\_older(\CodeNumTok{2})     \CommentTok{\# Method{-}based modification}
\end{Highlighting}
\end{Shaded}

\subsection{Skills Checklist}\label{skills-checklist-34}

\begin{itemize}
\tightlist
\item[$\square$]
  I can access an object's attributes using dot notation.
\item[$\square$]
  I can define and call methods within a class.
\item[$\square$]
  I can use a method to change an object's attribute.
\end{itemize}

\begin{center}\rule{0.5\linewidth}{0.5pt}\end{center}

\subsection{Lesson 7.4: Inheritance}\label{lesson-7.4-inheritance}

\subsection{Reading References}\label{reading-references-22}

\begin{itemize}
\tightlist
\item
  Zelle: Chapter 12, pages 351--352
\item
  Matthes: Chapter 9, pages 167--170
\end{itemize}

\subsection{Creating a Child Class}\label{creating-a-child-class}

\begin{Shaded}
\begin{Highlighting}[]
\KeywordTok{class} Car:
    \KeywordTok{def} \CodeFuncTok{\_\_init\_\_}(self, make, model, year):
        self.make \OperatorTok{=} make
        self.model \OperatorTok{=} model
        self.year \OperatorTok{=} year
        self.odometer\_reading \OperatorTok{=} \CodeNumTok{0}
 
    \KeywordTok{def} get\_descriptive\_name(self):
        \KeywordTok{return} \StringTok{f"}\StringTok{\{}self\StringTok{.}year\StringTok{\}}\StringTok{ }\StringTok{\{}self\StringTok{.}make\StringTok{\}}\StringTok{ }\StringTok{\{}self\StringTok{.}model\StringTok{\}}\StringTok{"}.title()
 
    \KeywordTok{def} read\_odometer(self):
        print(\StringTok{f"This car has }\StringTok{\{}self\StringTok{.}odometer\_reading\StringTok{\}}\StringTok{ miles on it."})
 
 
\KeywordTok{class} ElectricCar(Car):
    \KeywordTok{def} \CodeFuncTok{\_\_init\_\_}(self, make, model, year):
        super().\CodeFuncTok{\_\_init\_\_}(make, model, year)   \CommentTok{\# Call parent\textquotesingle{}s \_\_init\_\_}
        self.battery\_size \OperatorTok{=} \CodeNumTok{75}
 
    \KeywordTok{def} describe\_battery(self):
        print(\StringTok{f"This car has a }\StringTok{\{}self\StringTok{.}battery\_size\StringTok{\}}\StringTok{{-}kWh battery."})
 
 
my\_ev \OperatorTok{=} ElectricCar(\StringTok{\textquotesingle{}ford\textquotesingle{}}, \StringTok{\textquotesingle{}mustang mach e\textquotesingle{}}, \CodeNumTok{2023})
print(my\_ev.get\_descriptive\_name())   \CommentTok{\# Inherited method}
my\_ev.describe\_battery()              \CommentTok{\# New method}
my\_ev.read\_odometer()                 \CommentTok{\# Inherited method}
\end{Highlighting}
\end{Shaded}

\subsection{\texorpdfstring{\texttt{super()}
Function}{super() Function}}\label{super-function}

\texttt{super().\_\_init\_\_()} calls the parent class's
\texttt{\_\_init\_\_()} method, ensuring inherited attributes are
properly initialized.

\subsection{Skills Checklist}\label{skills-checklist-35}

\begin{itemize}
\tightlist
\item[$\square$]
  I can define a child class that inherits from a parent class.
\item[$\square$]
  I can use \texttt{super().\_\_init\_\_()} to properly initialize
  inherited attributes.
\item[$\square$]
  I can add new attributes and methods to a child class.
\item[$\square$]
  I can override a parent method in a child class.
\end{itemize}

\begin{center}\rule{0.5\linewidth}{0.5pt}\end{center}

\subsection{Lesson 7.5: Instances as Attributes
(Composition)}\label{lesson-7.5-instances-as-attributes-composition}

\subsection{Composition --- ``Has-A''
Relationship}\label{composition-has-a-relationship}

Instead of inheriting, a class \textbf{contains an instance of another
class} as an attribute.

\begin{Shaded}
\begin{Highlighting}[]
\KeywordTok{class} Battery:
    \KeywordTok{def} \CodeFuncTok{\_\_init\_\_}(self, battery\_size\OperatorTok{=}\CodeNumTok{75}):
        self.battery\_size \OperatorTok{=} battery\_size
 
    \KeywordTok{def} describe\_battery(self):
        print(\StringTok{f"This car has a }\StringTok{\{}self\StringTok{.}battery\_size\StringTok{\}}\StringTok{{-}kWh battery."})
 
    \KeywordTok{def} get\_range(self):
        \KeywordTok{if} self.battery\_size \OperatorTok{==} \CodeNumTok{75}:
            range\_miles \OperatorTok{=} \CodeNumTok{260}
        \KeywordTok{else}:
            range\_miles \OperatorTok{=} \CodeNumTok{310}
        print(\StringTok{f"About }\StringTok{\{}range\_miles\StringTok{\}}\StringTok{ miles of range."})
 
 
\KeywordTok{class} ElectricCar:
    \KeywordTok{def} \CodeFuncTok{\_\_init\_\_}(self, make, model, year):
        self.make \OperatorTok{=} make
        self.model \OperatorTok{=} model
        self.year \OperatorTok{=} year
        self.battery \OperatorTok{=} Battery()   \CommentTok{\# Composition: ElectricCar HAS A Battery}
 
 
my\_ev \OperatorTok{=} ElectricCar(\StringTok{\textquotesingle{}ford\textquotesingle{}}, \StringTok{\textquotesingle{}mustang mach e\textquotesingle{}}, \CodeNumTok{2023})
my\_ev.battery.describe\_battery()
my\_ev.battery.get\_range()
\end{Highlighting}
\end{Shaded}

\subsection{Skills Checklist}\label{skills-checklist-36}

\begin{itemize}
\tightlist
\item[$\square$]
  I can explain the concept of composition in OOP.
\item[$\square$]
  I can create an instance of one class and assign it as an attribute of
  another.
\item[$\square$]
  I can access attributes and methods of a nested object.
\end{itemize}

\begin{center}\rule{0.5\linewidth}{0.5pt}\end{center}

\section{Unit 8: Files and
Exceptions}\label{unit-8-files-and-exceptions-1}

\subsection{Overview}\label{overview-8}

This unit covers interacting with text files for reading and writing
data, handling common errors gracefully using exceptions, and storing
data using JSON.

\textbf{Objectives:}

\begin{itemize}
\tightlist
\item
  Read from and write to files.
\item
  Handle exceptions with \texttt{try-except} blocks.
\item
  Store and load data using JSON.
\end{itemize}

\begin{center}\rule{0.5\linewidth}{0.5pt}\end{center}

\subsection{Lesson 8.1: Reading from a
File}\label{lesson-8.1-reading-from-a-file}

\subsection{Reading References}\label{reading-references-23}

\begin{itemize}
\tightlist
\item
  Zelle: Chapter 5, pages 125--127
\item
  Matthes: Chapter 10, pages 184--189
\end{itemize}

\subsection{Reading Files}\label{reading-files}\amn{\textbf{Accessibility Tip:} Always use \texttt{with open(filename) as f:} to open files. The \texttt{with} block guarantees the file is closed when you are done, even if an error occurs. Forgetting to close files can cause data loss that produces no error message.}
\subsection{Reading Files}\label{reading-files}

\begin{Shaded}
\begin{Highlighting}[]
\CommentTok{\# Reading entire file}
\KeywordTok{with} open(\StringTok{\textquotesingle{}my\_notes.txt\textquotesingle{}}, \StringTok{\textquotesingle{}r\textquotesingle{}}) \CodeAccentTok{as} file\_object:
    contents \OperatorTok{=} file\_object.read()
print(contents)
 
\CommentTok{\# Reading line by line}
\KeywordTok{with} open(\StringTok{\textquotesingle{}my\_notes.txt\textquotesingle{}}, \StringTok{\textquotesingle{}r\textquotesingle{}}) \CodeAccentTok{as} file\_object:
    \KeywordTok{for} line \KeywordTok{in} file\_object:
        print(line.strip())
 
\CommentTok{\# Reading all lines into a list}
\KeywordTok{with} open(\StringTok{\textquotesingle{}my\_notes.txt\textquotesingle{}}, \StringTok{\textquotesingle{}r\textquotesingle{}}) \CodeAccentTok{as} file\_object:
    lines \OperatorTok{=} file\_object.readlines()
\end{Highlighting}
\end{Shaded}

The \texttt{with} statement \textbf{automatically closes the file} even
if an error occurs.

\subsection{File Modes}\label{file-modes}

\begin{fullwidth}{\def\LTcaptype{none} % do not increment counter
\begin{longtable}[]{@{}ll@{}}
\toprule\noalign{}
Mode & Description \\
\midrule\noalign{}
\endhead
\bottomrule\noalign{}
\endlastfoot
\texttt{\textquotesingle{}r\textquotesingle{}} & Read (default) \\
\texttt{\textquotesingle{}w\textquotesingle{}} & Write (creates new or
overwrites) \\
\texttt{\textquotesingle{}a\textquotesingle{}} & Append (adds to end) \\
\end{longtable}
}
\end{fullwidth}

\subsection{Skills Checklist}\label{skills-checklist-37}

\begin{itemize}
\tightlist
\item[$\square$]
  I can open a file for reading using \texttt{with\ open()}.
\item[$\square$]
  I can read the entire content of a file into a string.
\item[$\square$]
  I can read a file line by line using a \texttt{for} loop.
\item[$\square$]
  I can remove whitespace/newlines using \texttt{.strip()}.
\end{itemize}

\begin{center}\rule{0.5\linewidth}{0.5pt}\end{center}

\subsection{Lesson 8.2: Writing to a
File}\label{lesson-8.2-writing-to-a-file}

\subsection{Reading References}\label{reading-references-24}

\begin{itemize}
\tightlist
\item
  Zelle: Chapter 5, pages 127--128
\item
  Matthes: Chapter 10, pages 190--192
\end{itemize}

\begin{Shaded}
\begin{Highlighting}[]
\CommentTok{\# Write mode (overwrites existing content)}
\KeywordTok{with} open(\StringTok{\textquotesingle{}programming\_message.txt\textquotesingle{}}, \StringTok{\textquotesingle{}w\textquotesingle{}}) \CodeAccentTok{as} f:
    f.write(\StringTok{"I love programming in Python!}\StringTok{\textbackslash{}n}\StringTok{"})
 
\CommentTok{\# Append mode (adds to end)}
\KeywordTok{with} open(\StringTok{\textquotesingle{}programming\_message.txt\textquotesingle{}}, \StringTok{\textquotesingle{}a\textquotesingle{}}) \CodeAccentTok{as} f:
    f.write(\StringTok{"It\textquotesingle{}s a very versatile language.}\StringTok{\textbackslash{}n}\StringTok{"})
    f.write(\StringTok{"I\textquotesingle{}m excited to learn more!}\StringTok{\textbackslash{}n}\StringTok{"})
\end{Highlighting}
\end{Shaded}

\textbf{Important:} The \texttt{.write()} method does not add newlines
automatically. Use \texttt{\textbackslash{}n}.

\subsection{Skills Checklist}\label{skills-checklist-38}

\begin{itemize}
\tightlist
\item[$\square$]
  I can open a file for writing using \texttt{"w"} mode.
\item[$\square$]
  I can open a file for appending using \texttt{"a"} mode.
\item[$\square$]
  I can write a string to a file.
\item[$\square$]
  I can add newlines to a file using \texttt{\textbackslash{}n}.
\end{itemize}

\begin{center}\rule{0.5\linewidth}{0.5pt}\end{center}

\subsection{Lesson 8.3: Exceptions --- Handling
Errors}\label{lesson-8.3-exceptions-handling-errors}

\subsection{Reading References}\label{reading-references-25}

\begin{itemize}
\tightlist
\item
  Zelle: Chapter 7, pages 177--179
\item
  Matthes: Chapter 10, pages 192--199
\end{itemize}

\subsection{\texorpdfstring{\texttt{try-except}
Block}{try-except Block}}\amn{\textbf{Accessibility Tip:} Always catch specific exception types: \texttt{except ValueError:} not bare \texttt{except:}. Print the error in the except block: \texttt{except ValueError as e: print("Error:", e)}. A bare \texttt{except} silences all errors---including ones you did not expect---making bugs invisible.}
\subsection{\texorpdfstring{\texttt{try-except}
Block}{try-except Block}}\label{try-except-block}

\begin{Shaded}
\begin{Highlighting}[]
\KeywordTok{try}:
    result \OperatorTok{=} \CodeNumTok{10} \OperatorTok{/} \CodeNumTok{0}
\KeywordTok{except} \CodeAccentTok{ZeroDivisionError}:
    print(\StringTok{"Error: You cannot divide by zero!"})
print(\StringTok{"Program continues after error handling."})
\end{Highlighting}
\end{Shaded}

\subsection{Common Exception Types}\label{common-exception-types}

\begin{fullwidth}{\def\LTcaptype{none} % do not increment counter
\begin{longtable}[]{@{}ll@{}}
\toprule\noalign{}
Exception & When it occurs \\
\midrule\noalign{}
\endhead
\bottomrule\noalign{}
\endlastfoot
\texttt{ValueError} & Invalid value (e.g., \texttt{int("abc")}) \\
\texttt{ZeroDivisionError} & Division by zero \\
\texttt{FileNotFoundError} & File doesn't exist \\
\texttt{IndexError} & List index out of range \\
\texttt{KeyError} & Dictionary key doesn't exist \\
\texttt{TypeError} & Wrong type (e.g., \texttt{"5"\ +\ 5}) \\
\end{longtable}
}
\end{fullwidth}

\subsection{Examples}\label{examples-13}

\textbf{Example 8.3.1: Handling ZeroDivisionError}

\begin{Shaded}
\begin{Highlighting}[]
\KeywordTok{try}:
    result \OperatorTok{=} \CodeNumTok{10} \OperatorTok{/} \CodeNumTok{0}
\KeywordTok{except} \CodeAccentTok{ZeroDivisionError}:
    print(\StringTok{"Error: Cannot divide by zero!"})
\end{Highlighting}
\end{Shaded}

\textbf{Example 8.3.2: Handling ValueError}

\begin{Shaded}
\begin{Highlighting}[]
\KeywordTok{try}:
    num \OperatorTok{=} int(input(\StringTok{"Enter a number: "}))
    print(\StringTok{f"You entered: }\StringTok{\{}num\StringTok{\}}\StringTok{"})
\KeywordTok{except} \CodeAccentTok{ValueError}:
    print(\StringTok{"Error: That was not a valid number."})
\end{Highlighting}
\end{Shaded}

\subsection{Skills Checklist}\label{skills-checklist-39}

\begin{itemize}
\tightlist
\item[$\square$]
  I can use a \texttt{try-except} block to handle errors.
\item[$\square$]
  I can name common exception types and what causes them.
\item[$\square$]
  I can catch specific exception types.
\end{itemize}

\begin{center}\rule{0.5\linewidth}{0.5pt}\end{center}

\subsection{\texorpdfstring{Lesson 8.4: The \texttt{else} and
\texttt{finally}
Blocks}{Lesson 8.4: The else and finally Blocks}}\label{lesson-8.4-the-else-and-finally-blocks}

\subsection{\texorpdfstring{\texttt{else} Block --- Runs only on
success}{else Block --- Runs only on success}}\label{else-block-runs-only-on-success}

\begin{Shaded}
\begin{Highlighting}[]
\KeywordTok{try}:
    num \OperatorTok{=} int(input(\StringTok{"Enter a number: "}))
\KeywordTok{except} \CodeAccentTok{ValueError}:
    print(\StringTok{"That\textquotesingle{}s not a number."})
\KeywordTok{else}:
    print(\StringTok{f"You entered }\StringTok{\{}num\StringTok{\}}\StringTok{."})   \CommentTok{\# Only runs if no exception}
\end{Highlighting}
\end{Shaded}

\subsection{\texorpdfstring{\texttt{finally} Block --- Always runs
(cleanup)}{finally Block --- Always runs (cleanup)}}\label{finally-block-always-runs-cleanup}

\begin{Shaded}
\begin{Highlighting}[]
file\_handle \OperatorTok{=} None
\KeywordTok{try}:
    file\_handle \OperatorTok{=} open(\StringTok{\textquotesingle{}data.txt\textquotesingle{}}, \StringTok{\textquotesingle{}w\textquotesingle{}})
    file\_handle.write(\StringTok{"Some data."})
\KeywordTok{except} \CodeAccentTok{TypeError}:
    print(\StringTok{"A type error occurred."})
\KeywordTok{finally}:
    \KeywordTok{if} file\_handle:
        file\_handle.close()     \CommentTok{\# Always closes the file}
        print(\StringTok{"File closed."})
\end{Highlighting}
\end{Shaded}

\subsection{Skills Checklist}\label{skills-checklist-40}

\begin{itemize}
\tightlist
\item[$\square$]
  I can use an \texttt{else} block for code that runs only on success.
\item[$\square$]
  I can use a \texttt{finally} block for code that always runs.
\end{itemize}

\begin{center}\rule{0.5\linewidth}{0.5pt}\end{center}

\subsection{Lesson 8.5: Storing Data with
JSON}\label{lesson-8.5-storing-data-with-json}

\subsection{Reading References}\label{reading-references-26}

\begin{itemize}
\tightlist
\item
  Matthes: Chapter 10, pages 200--204
\end{itemize}

\subsection{\texorpdfstring{The \texttt{json}
Module}{The json Module}}\label{the-json-module}

JSON (JavaScript Object Notation) is a lightweight format for storing
and exchanging data. Python's \texttt{json} module converts between
Python objects and JSON.

\begin{Shaded}
\begin{Highlighting}[]
\CodeAccentTok{import} json
 
\CommentTok{\# Save data}
user\_settings \OperatorTok{=} \{\}
    \StringTok{"username"}: \StringTok{"coder\_kid"},
    \StringTok{"theme"}: \StringTok{"dark"},
    \StringTok{"notifications\_on"}: True,
    \StringTok{"favorite\_numbers"}: [\CodeNumTok{7}, \CodeNumTok{13}, \CodeNumTok{42}]
\
 
\KeywordTok{with} open(\StringTok{\textquotesingle{}user\_settings.json\textquotesingle{}}, \StringTok{\textquotesingle{}w\textquotesingle{}}) \CodeAccentTok{as} f:
    json.dump(user\_settings, f, indent\OperatorTok{=}\CodeNumTok{4})
 
\CommentTok{\# Load data}
\KeywordTok{try}:
    \KeywordTok{with} open(\StringTok{\textquotesingle{}user\_settings.json\textquotesingle{}}, \StringTok{\textquotesingle{}r\textquotesingle{}}) \CodeAccentTok{as} f:
        loaded\_settings \OperatorTok{=} json.load(f)
    print(\StringTok{f"Username: }\StringTok{\{}loaded\_settings[\StringTok{\textquotesingle{}username\textquotesingle{}}]\StringTok{\}}\StringTok{"})
    print(\StringTok{f"Theme: }\StringTok{\{}loaded\_settings[\StringTok{\textquotesingle{}theme\textquotesingle{}}]\StringTok{\}}\StringTok{"})
\KeywordTok{except} \CodeAccentTok{FileNotFoundError}:
    print(\StringTok{"Settings file not found."})
\KeywordTok{except} json.JSONDecodeError:
    print(\StringTok{"Error: Could not decode settings file."})
\end{Highlighting}
\end{Shaded}

\subsection{JSON\textless-\textgreater Python
Mapping}\label{json-python-mapping}

\begin{fullwidth}{\def\LTcaptype{none} % do not increment counter
\begin{longtable}[]{@{}ll@{}}
\toprule\noalign{}
JSON & Python \\
\midrule\noalign{}
\endhead
\bottomrule\noalign{}
\endlastfoot
Object \texttt{\{\}} & Dictionary \texttt{\{\}} \\
Array \texttt{{[}{]}} & List \texttt{{[}{]}} \\
String & \texttt{str} \\
Number & \texttt{int} or \texttt{float} \\
Boolean & \texttt{bool} \\
null & \texttt{None} \\
\end{longtable}
}
\end{fullwidth}

\subsection{Skills Checklist}\label{skills-checklist-41}

\begin{itemize}
\tightlist
\item[$\square$]
  I can explain what JSON is and why it's used.
\item[$\square$]
  I can use \texttt{json.dump()} to save Python data to a JSON file.
\item[$\square$]
  I can use \texttt{json.load()} to load JSON data from a file.
\item[$\square$]
  I can save and load lists of dictionaries using JSON.
\end{itemize}

\begin{center}\rule{0.5\linewidth}{0.5pt}\end{center}

\section{Unit 9: Testing and Data
Analysis}\label{unit-9-testing-and-data-analysis-1}

\subsection{Overview}\label{overview-9}

This unit introduces testing code for correctness and explores basic
text-based data analysis.

\begin{center}\rule{0.5\linewidth}{0.5pt}\end{center}

\subsection{Lesson 9.1: Testing a
Function}\label{lesson-9.1-testing-a-function}

\subsection{Reading References}\label{reading-references-27}

\begin{itemize}
\tightlist
\item
  Matthes: Chapter 11, pages 209--217
\end{itemize}

\subsection{\texorpdfstring{Unit Testing with
\texttt{unittest}}{Unit Testing with unittest}}\label{unit-testing-with-unittest}

\begin{Shaded}
\begin{Highlighting}[]
\CommentTok{\# my\_functions.py}
\KeywordTok{def} add\_two\_numbers(a, b):
    \KeywordTok{return} a \OperatorTok{+} b
\end{Highlighting}
\end{Shaded}

\begin{Shaded}
\begin{Highlighting}[]
\CommentTok{\# test\_my\_functions.py}
\CodeAccentTok{import} unittest
\CodeAccentTok{from} my\_functions \CodeAccentTok{import} add\_two\_numbers
 
\KeywordTok{class} TestAddNumbers(unittest.TestCase):
    \KeywordTok{def} test\_positive\_numbers(self):
        result \OperatorTok{=} add\_two\_numbers(\CodeNumTok{5}, \CodeNumTok{3})
        self.assertEqual(result, \CodeNumTok{8})
 
    \KeywordTok{def} test\_negative\_numbers(self):
        result \OperatorTok{=} add\_two\_numbers(\OperatorTok{{-}}\CodeNumTok{5}, \OperatorTok{{-}}\CodeNumTok{3})
        self.assertEqual(result, \OperatorTok{{-}}\CodeNumTok{8})
 
    \KeywordTok{def} test\_mixed\_numbers(self):
        result \OperatorTok{=} add\_two\_numbers(\CodeNumTok{10}, \OperatorTok{{-}}\CodeNumTok{4})
        self.assertEqual(result, \CodeNumTok{6})
 
\KeywordTok{if} \_\_name\_\_ \OperatorTok{==} \StringTok{\textquotesingle{}\_\_main\_\_\textquotesingle{}}:
    unittest.main()
\end{Highlighting}
\end{Shaded}

\subsection{Common Assertion Methods}\label{common-assertion-methods}

\begin{fullwidth}{\def\LTcaptype{none} % do not increment counter
\begin{longtable}[]{@{}ll@{}}
\toprule\noalign{}
Method & Description \\
\midrule\noalign{}
\endhead
\bottomrule\noalign{}
\endlastfoot
\texttt{assertEqual(a,\ b)} & Check a == b \\
\texttt{assertTrue(x)} & Check x is True \\
\texttt{assertFalse(x)} & Check x is False \\
\texttt{assertIn(a,\ b)} & Check a in b \\
\texttt{assertIsNone(x)} & Check x is None \\
\end{longtable}
}
\end{fullwidth}

\subsection{Skills Checklist}\label{skills-checklist-42}

\begin{itemize}
\tightlist
\item[$\square$]
  I can create a basic test class using \texttt{unittest.TestCase}.
\item[$\square$]
  I can define test methods within a test class.
\item[$\square$]
  I can use \texttt{assertEqual()} to check for expected results.
\item[$\square$]
  I can run unit tests from the terminal.
\end{itemize}

\begin{center}\rule{0.5\linewidth}{0.5pt}\end{center}

\subsection{Lesson 9.2: Testing a
Class}\label{lesson-9.2-testing-a-class}

\subsection{\texorpdfstring{The \texttt{setUp()}
Method}{The setUp() Method}}\label{the-setup-method}

\texttt{setUp()} runs \textbf{before every test method}, providing a
fresh instance each time:

\begin{Shaded}
\begin{Highlighting}[]
\KeywordTok{class} TestEmployee(unittest.TestCase):
    \KeywordTok{def} setUp(self):
        \CommentTok{"""Create an employee instance for all test methods."""}
        self.employee \OperatorTok{=} Employee(\StringTok{\textquotesingle{}janis\textquotesingle{}}, \StringTok{\textquotesingle{}joplin\textquotesingle{}}, \CodeNumTok{60000})
 
    \KeywordTok{def} test\_give\_default\_raise(self):
        self.employee.give\_raise()
        self.assertEqual(self.employee.salary, \CodeNumTok{65000})
 
    \KeywordTok{def} test\_give\_custom\_raise(self):
        self.employee.give\_raise(\CodeNumTok{10000})
        self.assertEqual(self.employee.salary, \CodeNumTok{70000})
\end{Highlighting}
\end{Shaded}

\subsection{Skills Checklist}\label{skills-checklist-43}

\begin{itemize}
\tightlist
\item[$\square$]
  I can write tests for a Python class.
\item[$\square$]
  I can use \texttt{setUp()} to prepare my tests.
\end{itemize}

\begin{center}\rule{0.5\linewidth}{0.5pt}\end{center}

\subsection{Lesson 9.3: Introduction to Random
Numbers}\label{lesson-9.3-introduction-to-random-numbers}

\begin{Shaded}
\begin{Highlighting}[]
\CodeAccentTok{import} random
 
\CommentTok{\# Random integer (inclusive)}
die\_roll \OperatorTok{=} random.randint(\CodeNumTok{1}, \CodeNumTok{6})
print(\StringTok{f"You rolled: }\StringTok{\{}die\_roll\StringTok{\}}\StringTok{"})
 
\CommentTok{\# Random float between 0.0 and 1.0}
rand\_float \OperatorTok{=} random.random()
 
\CommentTok{\# Random choice from list}
colors \OperatorTok{=} [\StringTok{"red"}, \StringTok{"blue"}, \StringTok{"green"}]
chosen \OperatorTok{=} random.choice(colors)
print(\StringTok{f"Chosen color: }\StringTok{\{}chosen\StringTok{\}}\StringTok{"})
\end{Highlighting}
\end{Shaded}

\subsection{Skills Checklist}\label{skills-checklist-44}

\begin{itemize}
\tightlist
\item[$\square$]
  I can import the \texttt{random} module.
\item[$\square$]
  I can generate a random integer within a range.
\item[$\square$]
  I can select a random item from a list.
\end{itemize}

\begin{center}\rule{0.5\linewidth}{0.5pt}\end{center}

\subsection{Lesson 9.4: Analyzing Text
Data}\label{lesson-9.4-analyzing-text-data}

\begin{Shaded}
\begin{Highlighting}[]
sentence \OperatorTok{=} \StringTok{"the quick brown fox jumps over the lazy dog the quick brown fox"}
words \OperatorTok{=} sentence.lower().split()
 
word\_frequencies \OperatorTok{=} \{\}
\KeywordTok{for} word \KeywordTok{in} words:
    \KeywordTok{if} word \KeywordTok{in} word\_frequencies:
        word\_frequencies[word] \OperatorTok{+=} \CodeNumTok{1}
    \KeywordTok{else}:
        word\_frequencies[word] \OperatorTok{=} \CodeNumTok{1}
 
print(\StringTok{"Word Frequencies:"})
\KeywordTok{for} word, count \KeywordTok{in} word\_frequencies.items():
    print(\StringTok{f"\textquotesingle{}}\StringTok{\{}word\StringTok{\}}\StringTok{\textquotesingle{}: }\StringTok{\{}count\StringTok{\}}\StringTok{"})
\end{Highlighting}
\end{Shaded}

\subsection{Skills Checklist}\label{skills-checklist-45}

\begin{itemize}
\tightlist
\item[$\square$]
  I can read text content from a file.
\item[$\square$]
  I can split a string of text into individual words.
\item[$\square$]
  I can use a dictionary to count the frequency of words.
\end{itemize}

\begin{center}\rule{0.5\linewidth}{0.5pt}\end{center}

\subsection{Lesson 9.5: Working with CSV
Data}\label{lesson-9.5-working-with-csv-data}

\begin{Shaded}
\begin{Highlighting}[]
\CodeAccentTok{import} csv
 
filename \OperatorTok{=} \StringTok{\textquotesingle{}students.csv\textquotesingle{}}
\KeywordTok{try}:
    \KeywordTok{with} open(filename, \StringTok{\textquotesingle{}r\textquotesingle{}}) \CodeAccentTok{as} file\_object:
        reader \OperatorTok{=} csv.reader(file\_object)
        header\_row \OperatorTok{=} next(reader)
        print(\StringTok{f"Columns: }\StringTok{\{}header\_row\StringTok{\}}\StringTok{"})
        \KeywordTok{for} row \KeywordTok{in} reader:
            print(row)
\KeywordTok{except} \CodeAccentTok{FileNotFoundError}:
    print(\StringTok{f"Error: The file }\StringTok{\{}filename\StringTok{\}}\StringTok{ was not found."})
\end{Highlighting}
\end{Shaded}

\subsection{Manual CSV Parsing}\label{manual-csv-parsing}

\begin{Shaded}
\begin{Highlighting}[]
\KeywordTok{with} open(\StringTok{\textquotesingle{}students.csv\textquotesingle{}}, \StringTok{\textquotesingle{}r\textquotesingle{}}) \CodeAccentTok{as} f:
    header \OperatorTok{=} f.readline().strip()
    \KeywordTok{for} line \KeywordTok{in} f:
        data \OperatorTok{=} line.strip().split(\StringTok{\textquotesingle{},\textquotesingle{}})
        name \OperatorTok{=} data[\CodeNumTok{0}]
        age \OperatorTok{=} data[\CodeNumTok{1}]
        grade \OperatorTok{=} data[\CodeNumTok{2}]
        print(\StringTok{f"Name: }\StringTok{\{}name\StringTok{\}}\StringTok{, Age: }\StringTok{\{}age\StringTok{\}}\StringTok{, Grade: }\StringTok{\{}grade\StringTok{\}}\StringTok{"})
\end{Highlighting}
\end{Shaded}

\subsection{Skills Checklist}\label{skills-checklist-46}

\begin{itemize}
\tightlist
\item[$\square$]
  I can explain what CSV data is.
\item[$\square$]
  I can read a CSV file line by line.
\item[$\square$]
  I can split a line of CSV data into individual values.
\end{itemize}

\newpage

\section{Unit 10: Project-Based Learning --- Building Text-Based
Applications}\label{unit-10-project-based-learning-building-text-based-applications}

\subsection{Overview}\label{overview-10}

This unit guides students through designing and building larger,
interactive text-based applications, integrating all previously learned
concepts.

\textbf{Objectives:}

\begin{itemize}
\tightlist
\item
  Practice modular programming and code organization.
\item
  Develop a text-based interactive program.
\item
  Apply top-down design principles.
\end{itemize}

\begin{center}\rule{0.5\linewidth}{0.5pt}\end{center}

\subsection{Lesson 10.1: Designing Your
Project}\label{lesson-10.1-designing-your-project}

\subsection{Reading References}\label{reading-references-28}

\begin{itemize}
\tightlist
\item
  Zelle: Chapter 9, pages 225--235
\item
  Matthes: Part II: Projects, pages 225--226
\end{itemize}

\subsection{Top-Down Design}\label{top-down-design}

\textbf{Top-down design} (stepwise refinement) breaks a complex problem
into smaller, more manageable subproblems.

\textbf{The Design Process:}

\textbf{1. Analyze the Problem}

\begin{itemize}
\tightlist
\item
  Identify \textbf{inputs} (information the program needs).
\item
  Identify \textbf{outputs} (results it must produce).
\item
  Define the scope of the project.
\end{itemize}

\textbf{2. Determine the Major Components} Write a high-level algorithm
in plain English:

\begin{verbatim}
show_greeting()
get_user_input()
calculate_cost()
print_ticket()
\end{verbatim}

\textbf{3. Refine Each Component} Repeat the design process for each
component, adding detail until simple enough to code.

\subsection{Benefits of Top-Down
Design}\label{benefits-of-top-down-design}

\begin{itemize}
\tightlist
\item
  \textbf{Reduces Complexity} --- Focus on one piece at a time.
\item
  \textbf{Encourages Reusability} --- Functions written for one part may
  be useful elsewhere.
\item
  \textbf{Simplifies Debugging} --- Test each component in isolation.
\item
  \textbf{Improves Collaboration} --- Clear plan for assigning work.
\end{itemize}

\subsection{Accessibility Considerations for I/O
Design}\label{accessibility-considerations-for-io-design}

\begin{itemize}
\tightlist
\item
  Provide \textbf{clear, descriptive prompts} for all input.
\item
  Provide \textbf{explicit, labeled output} --- avoid relying on visual
  formatting.
\item
  Print \textbf{confirmation messages} after every action.
\item
  Use consistent \textbf{output separators} like \texttt{-\/-\/-} to
  group information.
\end{itemize}

\subsection{Skills Checklist}\label{skills-checklist-47}

\begin{itemize}
\tightlist
\item[$\square$]
  I can break down a larger problem into smaller, manageable parts.
\item[$\square$]
  I can identify the key features required for a project.
\item[$\square$]
  I can plan which data structures will be most suitable for my
  project's data.
\item[$\square$]
  I can outline the main steps or flow of my program.
\end{itemize}

\begin{center}\rule{0.5\linewidth}{0.5pt}\end{center}

\subsection{Lesson 10.2: Building a Text-Based Adventure Game (Part
1)}\label{lesson-10.2-building-a-text-based-adventure-game-part-1}

\subsection{Reading References}\label{reading-references-29}

\begin{itemize}
\tightlist
\item
  Matthes: Chapter 6, pages 105--111
\end{itemize}

\subsection{Room Definition and
Movement}\label{room-definition-and-movement}

\begin{Shaded}
\begin{Highlighting}[]
room\_descriptions \OperatorTok{=} \{\}
    \StringTok{"start\_room"}: \{
        \StringTok{"description"}: \StringTok{"You are in a dimly lit cave. There\textquotesingle{}s a narrow passage to the NORTH."},
        \StringTok{"exits"}: \{\}\StringTok{"north"}: \StringTok{"tunnel"}\}
    \,
    \StringTok{"tunnel"}: \{
        \StringTok{"description"}: \StringTok{"The tunnel is damp and cold. You can go SOUTH back to the cave or EAST to a clearing."},
        \StringTok{"exits"}: \{\}\StringTok{"south"}: \StringTok{"start\_room"}, \StringTok{"east"}: \StringTok{"clearing"}\}
    \,
    \StringTok{"clearing"}: \{
        \StringTok{"description"}: \StringTok{"You\textquotesingle{}ve reached a small, sunny clearing. A faint path leads WEST."},
        \StringTok{"exits"}: \{\}\StringTok{"west"}: \StringTok{"tunnel"}\}
    \}
\
 
current\_room \OperatorTok{=} \StringTok{"start\_room"}
 
\KeywordTok{while} True:
    print(\StringTok{"{-}{-}{-}"})
    print(room\_descriptions[current\_room][\StringTok{"description"}])
    command \OperatorTok{=} input(\StringTok{"What do you want to do? "}).lower()
 
    \KeywordTok{if} command \KeywordTok{in} room\_descriptions[current\_room][\StringTok{"exits"}]:
        current\_room \OperatorTok{=} room\_descriptions[current\_room][\StringTok{"exits"}][command]
    \KeywordTok{elif} command \OperatorTok{==} \StringTok{"quit"}:
        print(\StringTok{"Exiting game. Goodbye!"})
        \KeywordTok{break}
    \KeywordTok{else}:
        print(\StringTok{"I don\textquotesingle{}t understand that command or you can\textquotesingle{}t go that way."})
\end{Highlighting}
\end{Shaded}

\subsection{Skills Checklist}\label{skills-checklist-48}

\begin{itemize}
\tightlist
\item[$\square$]
  I can design a way to represent different locations in a text-based
  game.
\item[$\square$]
  I can implement player movement based on user commands.
\item[$\square$]
  I can update and track the player's current location.
\item[$\square$]
  I can parse simple commands from user input.
\end{itemize}

\begin{center}\rule{0.5\linewidth}{0.5pt}\end{center}

\subsection{Lesson 10.3: Building a Text-Based Adventure Game (Part
2)}\label{lesson-10.3-building-a-text-based-adventure-game-part-2}

\subsection{Adding Items, Inventory, and
Puzzles}\label{adding-items-inventory-and-puzzles}

\begin{Shaded}
\begin{Highlighting}[]
rooms \OperatorTok{=} \{
    \StringTok{"start"}: \{
        \StringTok{"description"}: \StringTok{"A dark cave. There\textquotesingle{}s a SHINY KEY on the ground."},
        \StringTok{"exits"}: \{\}\StringTok{"north"}: \StringTok{"locked\_door\_room"}\},
        \StringTok{"items"}: [\StringTok{"shiny key"}]
    \,
    \StringTok{"locked\_door\_room"}: \{
        \StringTok{"description"}: \StringTok{"A small chamber with a heavy locked door to the NORTH."},
        \StringTok{"exits"}: \{\}\StringTok{"south"}: \StringTok{"start"}\},
        \StringTok{"items"}: []
    \}
\
 
current\_room \OperatorTok{=} \StringTok{"start"}
player\_inventory \OperatorTok{=} []
door\_locked \OperatorTok{=} True
 
\KeywordTok{while} True:
    print(\StringTok{"}\StringTok{\textbackslash{}n}\StringTok{"} \OperatorTok{+} rooms[current\_room][\StringTok{"description"}])
    \KeywordTok{if} rooms[current\_room][\StringTok{"items"}]:
        print(\StringTok{f"You see: }\StringTok{\{}\StringTok{\textquotesingle{}, \textquotesingle{}}\StringTok{.}join(rooms[current\_room][\StringTok{\textquotesingle{}items\textquotesingle{}}])\StringTok{\}}\StringTok{"})
    print(\StringTok{f"Your inventory: }\StringTok{\{}\StringTok{\textquotesingle{}, \textquotesingle{}}\StringTok{.}join(player\_inventory) \KeywordTok{if} player\_inventory \KeywordTok{else} \StringTok{\textquotesingle{}empty\textquotesingle{}}\StringTok{\}}\StringTok{"})
 
    command \OperatorTok{=} input(\StringTok{"What do you want to do? "}).lower()
 
    \KeywordTok{if} command \OperatorTok{==} \StringTok{"take shiny key"} \KeywordTok{and} \StringTok{"shiny key"} \KeywordTok{in} rooms[current\_room][\StringTok{"items"}]:
        player\_inventory.append(\StringTok{"shiny key"})
        rooms[current\_room][\StringTok{"items"}].remove(\StringTok{"shiny key"})
        print(\StringTok{"You picked up the shiny key."})
    \KeywordTok{elif} command \OperatorTok{==} \StringTok{"use shiny key"} \KeywordTok{and} current\_room \OperatorTok{==} \StringTok{"locked\_door\_room"} \KeywordTok{and} \StringTok{"shiny key"} \KeywordTok{in} player\_inventory \KeywordTok{and} door\_locked:
        door\_locked \OperatorTok{=} False
        print(\StringTok{"You unlocked the door!"})
        rooms[\StringTok{"locked\_door\_room"}][\StringTok{"exits"}][\StringTok{"north"}] \OperatorTok{=} \StringTok{"treasure\_room"}
    \KeywordTok{elif} command \KeywordTok{in} rooms[current\_room][\StringTok{"exits"}]:
        \KeywordTok{if} command \OperatorTok{==} \StringTok{"north"} \KeywordTok{and} current\_room \OperatorTok{==} \StringTok{"locked\_door\_room"} \KeywordTok{and} door\_locked:
            print(\StringTok{"The door is locked! You need a key."})
        \KeywordTok{else}:
            current\_room \OperatorTok{=} rooms[current\_room][\StringTok{"exits"}][command]
    \KeywordTok{elif} command \OperatorTok{==} \StringTok{"quit"}:
        \KeywordTok{break}
    \KeywordTok{else}:
        print(\StringTok{"I don\textquotesingle{}t understand that command."})
\end{Highlighting}
\end{Shaded}

\subsection{Skills Checklist}\label{skills-checklist-49}

\begin{itemize}
\tightlist
\item[$\square$]
  I can add items to locations and manage a player's inventory.
\item[$\square$]
  I can implement commands for picking up and dropping items.
\item[$\square$]
  I can design and implement a simple puzzle.
\item[$\square$]
  I can define win/lose conditions for the game.
\end{itemize}

\begin{center}\rule{0.5\linewidth}{0.5pt}\end{center}

\subsection{Lesson 10.4: Developing a Command-Line
Tool}\label{lesson-10.4-developing-a-command-line-tool}

\subsection{Simple To-Do List Tool}\label{simple-to-do-list-tool}

\begin{Shaded}
\begin{Highlighting}[]
\CodeAccentTok{import} json
 
TODO\_FILE \OperatorTok{=} \StringTok{\textquotesingle{}todos.json\textquotesingle{}}
todos \OperatorTok{=} []
 
 
\KeywordTok{def} load\_todos():
    \KeywordTok{global} todos
    \KeywordTok{try}:
        \KeywordTok{with} open(TODO\_FILE, \StringTok{\textquotesingle{}r\textquotesingle{}}) \CodeAccentTok{as} f:
            todos \OperatorTok{=} json.load(f)
        print(\StringTok{"Tasks loaded."})
    \KeywordTok{except} \CodeAccentTok{FileNotFoundError}:
        print(\StringTok{"No existing tasks found."})
        todos \OperatorTok{=} []
 
 
\KeywordTok{def} save\_todos():
    \KeywordTok{with} open(TODO\_FILE, \StringTok{\textquotesingle{}w\textquotesingle{}}) \CodeAccentTok{as} f:
        json.dump(todos, f, indent\OperatorTok{=}\CodeNumTok{4})
    print(\StringTok{"Tasks saved."})
 
 
\KeywordTok{def} add\_task(task\_description):
    todos.append(\{\StringTok{"description"}: task\_description, \StringTok{"completed"}: False\}
    print(\StringTok{f"Task \textquotesingle{}}\StringTok{\{}task\_description\StringTok{\}}\StringTok{\textquotesingle{} added."})
 
 
\KeywordTok{def} view\_tasks():
    \KeywordTok{if} \KeywordTok{not} todos:
        print(\StringTok{"No tasks in your list."})
        \KeywordTok{return}
    print(\StringTok{"{-}{-}{-} Your To{-}Do List {-}{-}{-}"})
    \KeywordTok{for} i, task \KeywordTok{in} enumerate(todos):
        status \OperatorTok{=} \StringTok{"[X]"} \KeywordTok{if} task[\StringTok{"completed"}] \KeywordTok{else} \StringTok{"[ ]"}
        print(\StringTok{f"}\StringTok{\{}i \OperatorTok{+} \CodeNumTok{1}\StringTok{\}}\StringTok{. }\StringTok{\{}status\StringTok{\}}\StringTok{ }\StringTok{\{}task[\StringTok{\textquotesingle{}description\textquotesingle{}}]\StringTok{\}}\StringTok{"})
    print(\StringTok{"{-}{-}{-}{-}{-}{-}{-}{-}{-}{-}{-}{-}{-}{-}{-}{-}{-}{-}{-}{-}{-}{-}{-}"})
 
 
\KeywordTok{def} mark\_completed(task\_number):
    \KeywordTok{try}:
        index \OperatorTok{=} int(task\_number) \OperatorTok{{-}} \CodeNumTok{1}
        \KeywordTok{if} \CodeNumTok{0} \OperatorTok{\textless{}=} index \OperatorTok{\textless{}} len(todos):
            todos[index][\StringTok{"completed"}] \OperatorTok{=} True
            print(\StringTok{f"Task \textquotesingle{}}\StringTok{\{}todos[index][\StringTok{\textquotesingle{}description\textquotesingle{}}]\StringTok{\}}\StringTok{\textquotesingle{} marked as complete."})
        \KeywordTok{else}:
            print(\StringTok{"Invalid task number."})
    \KeywordTok{except} \CodeAccentTok{ValueError}:
        print(\StringTok{"Please enter a valid number."})
 
 
\KeywordTok{def} main\_menu():
    print(\StringTok{"}\StringTok{\textbackslash{}n}\StringTok{{-}{-}{-} To{-}Do List Menu {-}{-}{-}"})
    print(\StringTok{"1. Add Task"})
    print(\StringTok{"2. View Tasks"})
    print(\StringTok{"3. Mark Task Complete"})
    print(\StringTok{"4. Save Tasks"})
    print(\StringTok{"5. Load Tasks"})
    print(\StringTok{"6. Quit"})
    print(\StringTok{"{-}{-}{-}{-}{-}{-}{-}{-}{-}{-}{-}{-}{-}{-}{-}{-}{-}{-}{-}{-}{-}{-}{-}"})
 
 
load\_todos()
 
\KeywordTok{while} True:
    main\_menu()
    choice \OperatorTok{=} input(\StringTok{"Enter your choice (1{-}6): "})
 
    \KeywordTok{if} choice \OperatorTok{==} \StringTok{\textquotesingle{}1\textquotesingle{}}:
        desc \OperatorTok{=} input(\StringTok{"Enter task description: "})
        add\_task(desc)
    \KeywordTok{elif} choice \OperatorTok{==} \StringTok{\textquotesingle{}2\textquotesingle{}}:
        view\_tasks()
    \KeywordTok{elif} choice \OperatorTok{==} \StringTok{\textquotesingle{}3\textquotesingle{}}:
        view\_tasks()
        task\_num \OperatorTok{=} input(\StringTok{"Enter task number to mark complete: "})
        mark\_completed(task\_num)
    \KeywordTok{elif} choice \OperatorTok{==} \StringTok{\textquotesingle{}4\textquotesingle{}}:
        save\_todos()
    \KeywordTok{elif} choice \OperatorTok{==} \StringTok{\textquotesingle{}5\textquotesingle{}}:
        load\_todos()
    \KeywordTok{elif} choice \OperatorTok{==} \StringTok{\textquotesingle{}6\textquotesingle{}}:
        save\_todos()
        print(\StringTok{"Goodbye!"})
        \KeywordTok{break}
    \KeywordTok{else}:
        print(\StringTok{"Invalid choice. Please try again."})
\end{Highlighting}
\end{Shaded}

\subsection{Skills Checklist}\label{skills-checklist-50}

\begin{itemize}
\tightlist
\item[$\square$]
  I can design an interactive menu system for a tool.
\item[$\square$]
  I can use a \texttt{while} loop to maintain continuous user
  interaction.
\item[$\square$]
  I can break down the tool's functionality into separate functions.
\item[$\square$]
  I can implement file I/O to save and load data.
\end{itemize}

\begin{center}\rule{0.5\linewidth}{0.5pt}\end{center}

\subsection{Lesson 10.5: Refinement and User
Experience}\label{lesson-10.5-refinement-and-user-experience}

\subsection{Principles for Accessible Text
Programs}\label{principles-for-accessible-text-programs}

\begin{enumerate}
\def\labelenumi{\arabic{enumi}.}
\tightlist
\item
  \textbf{Clear prompts} --- Tell the user exactly what to enter and in
  what format.
\item
  \textbf{Consistent output} --- Use the same format for similar
  information.
\item
  \textbf{Graceful error handling} --- Never let the program crash on
  bad input.
\item
  \textbf{Helpful error messages} --- Tell the user what went wrong and
  how to fix it.
\item
  \textbf{``Help'' commands} --- Provide a way to list available
  commands.
\item
  \textbf{Confirmation messages} --- Confirm every action taken.
\end{enumerate}

\subsection{Robust Input Validation}\label{robust-input-validation}

\begin{Shaded}
\begin{Highlighting}[]
\KeywordTok{def} get\_positive\_number(prompt):
    \CommentTok{"""Prompts user until a positive integer is entered."""}
    \KeywordTok{while} True:
        \KeywordTok{try}:
            num \OperatorTok{=} int(input(prompt))
            \KeywordTok{if} num \OperatorTok{\textgreater{}} \CodeNumTok{0}:
                \KeywordTok{return} num
            \KeywordTok{else}:
                print(\StringTok{"Please enter a positive number."})
        \KeywordTok{except} \CodeAccentTok{ValueError}:
            print(\StringTok{"Invalid input. Please enter a whole number."})
 
 
age \OperatorTok{=} get\_positive\_number(\StringTok{"Enter your age: "})
print(\StringTok{f"You entered age: }\StringTok{\{}age\StringTok{\}}\StringTok{"})
\end{Highlighting}
\end{Shaded}

\subsection{Skills Checklist}\label{skills-checklist-51}

\begin{itemize}
\tightlist
\item[$\square$]
  I can write clear, descriptive prompts for user input.
\item[$\square$]
  I can handle invalid input gracefully using \texttt{try-except}.
\item[$\square$]
  I can provide helpful error messages.
\end{itemize}

\begin{center}\rule{0.5\linewidth}{0.5pt}\end{center}

\section{Unit 11: Capstone Project --- Accessible Personal Data
Manager}\label{unit-11-capstone-project-accessible-personal-data-manager}

\subsection{Overview}\label{overview-11}

Students will develop a fully functional, screen-reader-accessible
graphical application using NiceGUI that allows users to manage a
collection of structured personal data.

\textbf{Project Goal:} Build a simple application to manage a list of
structured items.

\textbf{Project Options:}

\begin{itemize}
\tightlist
\item
  Contact Manager: name, phone, email, notes
\item
  Book/Movie Tracker: title, author/director, genre, status
\item
  Simple To-Do List: task description, due date, status
\item
  Simple Inventory: item name, quantity, location, notes
\end{itemize}

\begin{center}\rule{0.5\linewidth}{0.5pt}\end{center}

\subsection{Core Functionality
Requirements}\label{core-functionality-requirements}

\begin{enumerate}
\def\labelenumi{\arabic{enumi}.}
\tightlist
\item
  \textbf{Add New Item} --- A form to input details for a new item.
\item
  \textbf{View All Items} --- Display items in a clear, accessible list
  format.
\item
  \textbf{Edit Item} --- Select an existing item and modify its details.
\item
  \textbf{Delete Item} --- Remove an item from the collection.
\item
  \textbf{Save/Load Data} --- Persist the data to a JSON file.
\end{enumerate}

\begin{center}\rule{0.5\linewidth}{0.5pt}\end{center}

\subsection{Setup}\label{setup}

\begin{Shaded}
\begin{Highlighting}[]
pip install nicegui
\end{Highlighting}
\end{Shaded}

\begin{center}\rule{0.5\linewidth}{0.5pt}\end{center}

\subsection{Guided Implementation}\label{guided-implementation}

\subsection{Lesson 11.1: Introduction to
NiceGUI}\label{lesson-11.1-introduction-to-nicegui}

\textbf{Task 1:} Run the \texttt{app\_skeleton.py} file and navigate
using the keyboard and screen reader to understand the existing
structure.

\textbf{Task 2:} Define your data model class:

\begin{Shaded}
\begin{Highlighting}[]
\KeywordTok{class} Contact:
    \KeywordTok{def} \CodeFuncTok{\_\_init\_\_}(self, name, phone, email, notes):
        self.name \OperatorTok{=} name
        self.phone \OperatorTok{=} phone
        self.email \OperatorTok{=} email
        self.notes \OperatorTok{=} notes
\end{Highlighting}
\end{Shaded}

\textbf{Task 3:} Implement data persistence using the \texttt{json}
module within the \texttt{load\_data()} and \texttt{save\_data()}
functions.

\subsection{Lesson 11.2: Data Model and
Storage}\label{lesson-11.2-data-model-and-storage}

\textbf{Task 1:} Complete the \texttt{add\_item\_logic()} function:

\begin{enumerate}
\def\labelenumi{\arabic{enumi}.}
\tightlist
\item
  Get values from input fields using \texttt{.value} attribute.
\item
  Validate (e.g., name cannot be empty).
\item
  Provide accessible feedback:
  \texttt{ui.notify("Name\ is\ required!",\ type=\textquotesingle{}negative\textquotesingle{})}.
\end{enumerate}

\textbf{Task 2:} Create a new item object and add it to
\texttt{all\_items}.

\textbf{Task 3:} Clear input fields and provide success notification.

\subsection{Lesson 11.3--11.5: Displaying and Managing
Items}\label{lesson-11.311.5-displaying-and-managing-items}

\textbf{Implement \texttt{update\_item\_display()}:}

\begin{enumerate}
\def\labelenumi{\arabic{enumi}.}
\tightlist
\item
  Clear the display: \texttt{items\_display\_column.clear()}.
\item
  Loop through \texttt{all\_items} and create \texttt{ui.card} for each.
\item
  Update \texttt{item\_selection\_dropdown.options}.
\end{enumerate}

\textbf{Implement \texttt{update\_item\_logic()}:}

\begin{enumerate}
\def\labelenumi{\arabic{enumi}.}
\tightlist
\item
  Find the item in \texttt{all\_items}.
\item
  Update attributes from edit form fields.
\item
  Call \texttt{update\_item\_display()} and \texttt{save\_data()}.
\end{enumerate}

\textbf{Implement \texttt{delete\_selected\_item\_logic()}:}

\begin{enumerate}
\def\labelenumi{\arabic{enumi}.}
\tightlist
\item
  Find and remove the item from \texttt{all\_items}.
\item
  Call \texttt{update\_item\_display()}.
\item
  Provide accessible feedback.
\end{enumerate}

\begin{center}\rule{0.5\linewidth}{0.5pt}\end{center}

\subsection{Assessment Criteria}\label{assessment-criteria}

\subsection{Functionality (60\%)}\label{functionality-60}

\begin{itemize}
\tightlist
\item[$\square$]
  Ability to add new items with all specified attributes.
\item[$\square$]
  Correct display of all current items.
\item[$\square$]
  Correct editing of existing items.
\item[$\square$]
  Correct deletion of items.
\item[$\square$]
  Accurate saving and loading to/from \texttt{data.json}.
\item[$\square$]
  No major bugs or crashes during typical usage.
\end{itemize}

\subsection{Code Quality \& Structure
(20\%)}\label{code-quality-structure-20}

\begin{itemize}
\tightlist
\item[$\square$]
  Clear and concise Python code.
\item[$\square$]
  Appropriate use of functions and/or classes.
\item[$\square$]
  Meaningful variable and function names.
\item[$\square$]
  Adequate comments where complex logic exists.
\item[$\square$]
  Proper indentation.
\end{itemize}

\subsection{Accessibility \& User Experience
(15\%)}\label{accessibility-user-experience-15}

\begin{itemize}
\tightlist
\item[$\square$]
  All interactive elements are correctly labeled for screen readers.
\item[$\square$]
  Clear and consistent user feedback via \texttt{ui.notify()}.
\item[$\square$]
  Logical flow and presentation of information.
\item[$\square$]
  No visual-only cues for critical information.
\end{itemize}

\subsection{README and Reflection
(5\%)}\label{readme-and-reflection-5}

\begin{itemize}
\tightlist
\item[$\square$]
  Complete and clear README.
\item[$\square$]
  Thoughtful reflection on development and accessibility challenges.
\end{itemize}

\begin{center}\rule{0.5\linewidth}{0.5pt}\end{center}

\subsection{Deliverables}\label{deliverables}

\begin{enumerate}
\def\labelenumi{\arabic{enumi}.}
\tightlist
\item
  \texttt{app.py} --- Complete NiceGUI application code.
\item
  \texttt{data.json} --- Example JSON file demonstrating saved data.
\item
  \texttt{README.md} --- Including:

  \begin{itemize}
  \tightlist
  \item
    Project title and student name.
  \item
    Description of the application.
  \item
    Instructions on how to run: \texttt{python\ app.py}.
  \item
    Instructions for interacting with a screen reader.
  \item
    Reflection on challenges faced.
  \end{itemize}
\end{enumerate}

\begin{center}\rule{0.5\linewidth}{0.5pt}\end{center}

\subsection{Supplemental: AI Tools and Version
Control}\label{supplemental-ai-tools-and-version-control}

\subsection{AI-Powered Coding Tools}\label{ai-powered-coding-tools}

\textbf{Benefits:}

\begin{itemize}
\tightlist
\item
  Increased productivity and speed.
\item
  Syntax and idiom discovery.
\item
  Rapid prototyping.
\end{itemize}

\textbf{Critical Warning:} AI tools can produce code that looks correct
but contains subtle bugs. \textbf{Always:}

\begin{enumerate}
\def\labelenumi{\arabic{enumi}.}
\tightlist
\item
  Read every line of AI-generated code.
\item
  Verify the logic and intent.
\item
  Look for familiar error types (indentation, name errors, type errors).
\item
  Run and test incrementally.
\item
  Add your own comments and docstrings.
\end{enumerate}

\textbf{The Golden Rule:} You are responsible for your code. Treat AI
suggestions as a first draft that needs meticulous editing.

\subsection{Git and Version Control}\label{git-and-version-control}

\textbf{Git} is the most widely used distributed version control system,
invented in 2005 by Linus Torvalds for the Linux kernel project.

\textbf{Core workflow:}

\begin{enumerate}
\def\labelenumi{\arabic{enumi}.}
\tightlist
\item
  \textbf{Working directory} --- Make changes to files.
\item
  \textbf{Staging area} --- \texttt{git\ add} --- Stage changes.
\item
  \textbf{Repository} --- \texttt{git\ commit} --- Save a permanent
  snapshot.
\item
  \textbf{Remote} --- \texttt{git\ push} / \texttt{git\ pull} --- Sync
  with GitHub.
\end{enumerate}

\textbf{Why Git is especially valuable with AI code:}
\marginnote{Version control is critical for blind coders using AI, because AI frequently introduces bugs that are difficult to detect aurally.}

\begin{itemize}
\tightlist
\item
  \textbf{Rollback} --- If AI-generated code breaks your app, revert
  instantly.
\item
  \textbf{\texttt{git\ blame}} --- See exactly which commit introduced a
  bug.
\item
  \textbf{Branches} --- Test AI experiments in isolation without
  affecting main code.
\item
  \textbf{Pull requests} --- Review AI-generated code before merging.
\end{itemize}
\amn{\textbf{Accessibility Tip:} Run \texttt{git diff} before every commit to hear exactly what changed: lines prefixed with \texttt{+} were added, lines prefixed with \texttt{-} were removed. This is the most reliable way to review changes.}

\textbf{Basic commands:}
\begin{itemize}
\tightlist
\item
  \texttt{git init} --- Initialize a repository.
\item
  \texttt{git add .} --- Stage all changes.
\item
  \texttt{git commit -m "Add feature"} --- Save a snapshot.
\item
  \texttt{git log} --- View commit history.
\item
  \texttt{git diff} --- See what changed.
\item
  \texttt{git checkout -- file.py} --- Discard changes to a file.
\end{itemize}

\FloatBarrier
\clearpage
\part{JavaScript Curriculum}
\chapter*{Introduction to Computer Science with JavaScript}

\subsection{For Visually Impaired
Students}\label{for-visually-impaired-students}

\subsection{Course Overview}\label{course-overview}

This curriculum introduces computer science concepts using JavaScript,
focusing on Node.js for server-side/terminal execution. All exercises
are text-based and terminal-focused --- no browser required for core
concepts.

\textbf{Target Audience:} Students with no prior programming experience
using screen readers (NVDA, JAWS, VoiceOver, or Windows Narrator).

\begin{center}\rule{0.5\linewidth}{0.5pt}\end{center}

\subsection{Why JavaScript?}\label{why-javascript}

JavaScript is one of the most widely used programming languages in the
world. It powers:

\begin{itemize}
\tightlist
\item
  Web development (front-end and back-end)
\item
  Server applications (Node.js)
\item
  Mobile apps (React Native)
\item
  Desktop apps (Electron)
\end{itemize}

Learning JavaScript opens doors to a vast ecosystem of tools and
opportunities.

\begin{center}\rule{0.5\linewidth}{0.5pt}\end{center}

\subsection{Tools Required}\label{tools-required}

\begin{itemize}
\tightlist
\item
  \textbf{Node.js 18 LTS+} --- \url{https://nodejs.org/}
\item
  \textbf{Visual Studio Code} with the following extensions:

  \begin{itemize}
  \tightlist
  \item
    JavaScript (ES6) code snippets
  \item
    ESLint
  \item
    Prettier
  \end{itemize}
\item
  A screen reader (NVDA, JAWS, or VoiceOver)
\end{itemize}

\subsection{Accessibility Notes}\label{accessibility-notes}

\begin{itemize}
\tightlist
\item
  All program output is text-based (Node.js terminal/REPL).
\item
  JavaScript uses curly braces \texttt{\{\}} for code blocks, like Java.
\item
  VS Code's Screen Reader Optimized Mode works well with JavaScript.
\item
  Use \texttt{Ctrl\ +\ Shift\ +\ P} -\textgreater{} ``Toggle Screen
  Reader Optimized'' in VS Code.
\item
  Use \texttt{F8} / \texttt{Shift+F8} to navigate between errors.
\end{itemize}

\begin{center}\rule{0.5\linewidth}{0.5pt}\end{center}

\subsection{Units}\label{units}

\begin{fullwidth}{\def\LTcaptype{none} % do not increment counter
\begin{longtable}[]{@{}ll@{}}
\toprule\noalign{}
Unit & Title \\
\midrule\noalign{}
\endhead
\bottomrule\noalign{}
\endlastfoot
Unit 0 & Setting Up Your JavaScript Development Environment \\
Unit 1 & Foundations --- Getting Started with JavaScript \\
Unit 2 & Numbers, Input, and the Math Object \\
Unit 3 & Control Flow --- Decisions and Loops \\
Unit 4 & Data Structures --- Arrays and Objects \\
Unit 5 & User Input in Node.js \\
Unit 6 & Functions \\
Unit 7 & Classes and Object-Oriented Programming \\
Unit 8 & Files and Error Handling \\
Unit 9 & Modules, Collections, and Testing \\
Unit 10 & Project-Based Learning \\
Unit 11 & Capstone Project \\
\end{longtable}
}
\end{fullwidth}

\begin{center}\rule{0.5\linewidth}{0.5pt}\end{center}

\subsection{References}\label{references}

\begin{itemize}
\tightlist
\item
  MDN Web Docs:
  \url{https://developer.mozilla.org/en-US/docs/Web/JavaScript}
\item
  Eloquent JavaScript (free online):
  \url{https://eloquentjavascript.net/}
\item
  Node.js Documentation: \url{https://nodejs.org/docs/}
\end{itemize}

\newpage

\section{Teacher's Guide: Introduction to Computer Science with
JavaScript}\label{teachers-guide-introduction-to-computer-science-with-javascript}

\subsection{For Students with Visual Impairments --- Complete Instructor
Edition}\label{for-students-with-visual-impairments-complete-instructor-edition-1}

\begin{center}\rule{0.5\linewidth}{0.5pt}\end{center}

\subsection{Table of Contents}\label{table-of-contents-1}

\begin{enumerate}
\def\labelenumi{\arabic{enumi}.}
\tightlist
\item
  \hyperref[philosophy-and-why-javascript]{Philosophy and Why
  JavaScript}
\item
  \hyperref[javascript-specific-accessibility-challenges]{JavaScript-Specific
  Accessibility Challenges}
\item
  \hyperref[classroom-setup-and-configuration]{Classroom Setup and
  Configuration}
\item
  \hyperref[spoken-javascript-conventions]{Spoken JavaScript
  Conventions}
\item
  \hyperref[unit-01-setup-and-foundations]{Unit 0--1: Setup and
  Foundations}
\item
  \hyperref[unit-2-numbers-input-and-math]{Unit 2: Numbers, Input, and
  Math}
\item
  \hyperref[unit-3-control-flow]{Unit 3: Control Flow}
\item
  \hyperref[unit-4-arrays-and-objects]{Unit 4: Arrays and Objects}
\item
  \hyperref[unit-5-strings]{Unit 5: Strings}
\item
  \hyperref[unit-6-functions]{Unit 6: Functions}
\item
  \hyperref[unit-7-classes-and-oop]{Unit 7: Classes and OOP}
\item
  \hyperref[unit-8-files-modules-and-error-handling]{Unit 8: Files,
  Modules, and Error Handling}
\item
  \hyperref[unit-9-collections-analysis-and-testing]{Unit 9:
  Collections, Analysis, and Testing}
\item
  \hyperref[units-1011-projects-and-capstone]{Units 10--11: Projects and
  Capstone}
\item
  \hyperref[bridging-from-pythonjava-to-javascript]{Bridging from
  Python/Java to JavaScript}
\item
  \hyperref[assessment-strategies]{Assessment Strategies}
\item
  \hyperref[supporting-diverse-needs]{Supporting Diverse Needs}
\item
  \hyperref[alternative-and-supplementary-resources]{Alternative and
  Supplementary Resources}
\item
  \hyperref[quick-reference-sheets]{Quick Reference Sheets}
\end{enumerate}

\begin{center}\rule{0.5\linewidth}{0.5pt}\end{center}

\subsection{1. Philosophy and Why
JavaScript}\label{philosophy-and-why-javascript}

\subsection{The Case for Teaching JavaScript to Blind
Students}\label{the-case-for-teaching-javascript-to-blind-students}

JavaScript has a unique profile of benefits and challenges for visually
impaired learners:

\textbf{Universal relevance.} JavaScript is the only language that runs
natively in every web browser. A student who learns JavaScript can
eventually build tools that improve their own accessibility experience
--- custom browser extensions, accessible widgets, automation scripts.
This self-advocacy potential is meaningful.
\sidenote{Several prominent screen reader users and accessibility advocates are also JavaScript developers who built or extended tools for their own use: custom NVDA add-ons, browser scripts that fix inaccessible pages, and automated testing tools. JavaScript is uniquely positioned as the language of accessibility tooling.}

\textbf{Node.js enables accessible terminal programs.} While JavaScript
began in browsers\amn{\textbf{Accessibility Tip:} Run all JavaScript with \texttt{node filename.js} from the terminal rather than in a browser console. The terminal is fully accessible to any screen reader; browser developer consoles have variable and often poor screen reader support.}
\textbf{Node.js enables accessible terminal programs.} While JavaScript
began in browsers, Node.js allows running JavaScript programs in the
terminal --- the same accessible text-based environment used for Python.
\marginnote{The browser is not an accessible first environment for learning JavaScript. Browser developer tools, the console, and the DOM all require navigating complex, partially-accessible UIs. Node.js keeps everything in the terminal where screen readers excel. This curriculum deliberately avoids browser-based JavaScript until students are fluent with the language itself.}
This curriculum focuses on Node.js to keep all I/O in the terminal where
screen readers excel.

\textbf{Asynchronous programming is a real-world skill.} JavaScript's
\texttt{async}/\texttt{await} model reflects how most modern software
actually works --- programs that wait for I/O, network requests, and
user events. Learning this early prepares students for professional
development.

\textbf{The npm ecosystem.} Over 2 million packages on npm provide
pre-built solutions for almost any problem.
\sidenote{npm package quality and accessibility vary enormously. Before recommending any package, verify it works in a Node.js terminal context (no browser DOM required) and produces text output that is screen-reader navigable. Packages that produce progress bars, spinners, or colour-coded terminal output can be disruptive for screen reader users.} This is a double-edged
concern for beginners (dependency complexity), but the ability to say
\texttt{npm\ install} and gain powerful functionality is professionally
relevant.

\textbf{JSON is native.} \texttt{JSON.parse()} and
\texttt{JSON.stringify()} are built into JavaScript with no import
required. For data persistence projects, this reduces friction.

\subsection{An Honest Account of
Challenges}\label{an-honest-account-of-challenges}

\textbf{JavaScript's permissiveness is a silent-failure risk.} Unlike
Python (which raises exceptions for type mismatches) and Java (which
rejects them at compile time), JavaScript silently converts types.
\texttt{"5"\ +\ 3} = \texttt{"53"}. \texttt{"5"\ -\ 3} = \texttt{2}.
\texttt{0\ ==\ false} = \texttt{true}. Without visual scanning of
intermediate results, these silent failures can propagate through a
program undetected. This must be addressed through defensive programming
habits taught from day one.

\textbf{Asynchronous programming is conceptually hard.} The
\texttt{async}/\texttt{await} pattern required for user input is a
genuine hurdle for beginners. There is no way to eliminate it from a
Node.js input-based curriculum. The approach in this guide is to
scaffold it carefully and defer full explanation of the event loop to
later units.

\textbf{Multiple valid syntaxes.} Function declarations, function
expressions, arrow functions, methods --- all do essentially the same
thing. \texttt{var}, \texttt{let}, \texttt{const} --- three variable
declarations. \texttt{==} and \texttt{===} --- two equality operators.
This variety, which professionals call ``JavaScript fatigue,'' creates
unnecessary cognitive overhead for beginners. The approach in this
curriculum is to teach one style consistently
(\texttt{const}/\texttt{let}, arrow functions, \texttt{===}) and
acknowledge alternatives exist without teaching them in parallel.
\marginnote{The one-style-consistently approach is even more important for VI students than for sighted ones. When a student encounters unfamiliar syntax in tutorial code or Stack Overflow answers, they must navigate it without being able to scan the structure visually. Teach recognition of variants as a separate explicit lesson, not by exposure.}

\begin{center}\rule{0.5\linewidth}{0.5pt}\end{center}

\subsection{2. JavaScript-Specific Accessibility
Challenges}\label{javascript-specific-accessibility-challenges}

\subsection{Silent Type Coercion}\label{silent-type-coercion}

This is the greatest accessibility risk in JavaScript for blind
students. Unlike Python's \texttt{TypeError} and Java's compile-time
type check, JavaScript silently converts types in ways that can produce
wrong answers with no error:

\begin{fullwidth}{\def\LTcaptype{none} % do not increment counter
\begin{longtable}[]{@{}
  >{\raggedright\arraybackslash}p{(\linewidth - 6\tabcolsep) * \real{0.2043}}
  >{\raggedright\arraybackslash}p{(\linewidth - 6\tabcolsep) * \real{0.0753}}
  >{\raggedright\arraybackslash}p{(\linewidth - 6\tabcolsep) * \real{0.2366}}
  >{\raggedright\arraybackslash}p{(\linewidth - 6\tabcolsep) * \real{0.4839}}@{}}
\toprule\noalign{}
\begin{minipage}[b]{\linewidth}\raggedright
Expression
\end{minipage} & \begin{minipage}[b]{\linewidth}\raggedright
Result
\end{minipage} & \begin{minipage}[b]{\linewidth}\raggedright
Expected by Programmer
\end{minipage} & \begin{minipage}[b]{\linewidth}\raggedright
Problem
\end{minipage} \\
\midrule\noalign{}
\endhead
\bottomrule\noalign{}
\endlastfoot
\texttt{"5"\ +\ 3} & \texttt{"53"} & \texttt{8} & \texttt{+}
concatenates when one operand is a string \\
\texttt{"5"\ -\ 3} & \texttt{2} & \texttt{2} & \texttt{-} converts
string to number \\
\texttt{"5"\ *\ "3"} & \texttt{15} & \texttt{15} & Multiplication
converts both \\
\texttt{""\ ==\ false} & \texttt{true} & Maybe? & Loose equality
coercion \\
\texttt{0\ ==\ false} & \texttt{true} & Maybe? & Same \\
\texttt{null\ ==\ undefined} & \texttt{true} & Probably not & \\
\texttt{NaN\ ===\ NaN} & \texttt{false} & \texttt{true} & NaN is not
equal to itself \\
\end{longtable}
}
\end{fullwidth}

\textbf{Mitigation:} Require \texttt{Number()} before any arithmetic on
user input.\amn{\textbf{Accessibility Tip:} Wrap every value from user input in \texttt{Number()} before performing arithmetic: \texttt{const n = Number(await askQuestion("Enter number: "))}. Without this, \texttt{"5" + 3} silently produces \texttt{"53"} instead of \texttt{8}---a wrong answer with no error.}
\textbf{Mitigation:} Require \texttt{Number()} before any arithmetic on
user input. Require \texttt{===} always. Teach \texttt{typeof} and
\texttt{isNaN()} as defensive checks. Make logging intermediate values
an explicit exercise before running any calculation.

\subsection{\texorpdfstring{\texttt{undefined} vs.~\texttt{null}
vs.~\texttt{NaN} vs.~\texttt{false} vs.~\texttt{0}
vs.~\texttt{""}}{undefined vs.~null vs.~NaN vs.~false vs.~0 vs.~""}}\label{undefined-vs.-null-vs.-nan-vs.-false-vs.-0-vs.}

All six of these are ``falsy'' in JavaScript. Aurally, they are
indistinguishable unless you call \texttt{typeof} or check explicitly. A
student who hears \texttt{console.log(result)} and the screen reader
announces ``false'' cannot know whether \texttt{result} is
\texttt{false}, \texttt{0}, \texttt{""}, \texttt{null},
\texttt{undefined}, or \texttt{NaN} without further inspection.

\textbf{Teaching approach:} Establish a diagnostic template:

\begin{Shaded}
\begin{Highlighting}[]
console\OperatorTok{.}\CodeFuncTok{log}(\StringTok{"Value:"}\OperatorTok{,} result\OperatorTok{,} \StringTok{"Type:"}\OperatorTok{,} \KeywordTok{typeof} result)\OperatorTok{;}
\end{Highlighting}
\end{Shaded}

This outputs both the value and its type, providing audio context that
\texttt{console.log(result)} alone does not.
\sidenote{The \texttt{"Value:", result, "Type:", typeof result} diagnostic pattern} \amn{\textbf{Accessibility Tip:} When debugging a JavaScript value, always log both value and type: \texttt{console.log("Value:", x, "Type:", typeof x)}. A screen reader announces both, giving you type information that is invisible when only the value is logged.}
\sidenote{The \texttt{"Value:", result, "Type:", typeof result} diagnostic pattern should become as automatic for JavaScript students as \texttt{type()} is in Python. Both serve the same function: converting invisible type information into audible output.}

\subsection{Callback Hell and Asynchronous
Confusion}\label{callback-hell-and-asynchronous-confusion}

Before \texttt{async}/\texttt{await}, Node.js input used nested
callbacks:

\begin{Shaded}
\begin{Highlighting}[]
\CommentTok{// Confusing nesting {-}{-}{-} avoid in teaching}
rl\OperatorTok{.}\CodeFuncTok{question}(\StringTok{"Name? "}\OperatorTok{,} (name) \KeywordTok{=\textgreater{}} \{\}
  rl\OperatorTok{.}\CodeFuncTok{question}(\StringTok{"Age? "}\OperatorTok{,} (age) \KeywordTok{=\textgreater{}} \{
    rl\OperatorTok{.}\CodeFuncTok{question}(\StringTok{"City? "}\OperatorTok{,} (city) \KeywordTok{=\textgreater{}} \{
      console\OperatorTok{.}\CodeFuncTok{log}(\texttt{`${name}, ${age}, ${city}`})\OperatorTok{;}
    )\}\OperatorTok{;}
  )\}\OperatorTok{;}
)\OperatorTok{;}
\end{Highlighting}
\end{Shaded}

This nesting produces code that is confusing to read AND to hear --- the
indentation deepens with each question. The
\texttt{async}/\texttt{await} pattern linearizes this:
\marginnote{The linearisation that \texttt{async}/\texttt{await} provides is not just aesthetically cleaner---it is a direct accessibility improvement. Deeply nested callbacks produce code with increasing indentation levels that a screen reader announces at every line. \texttt{async}/\texttt{await} keeps indentation flat.}

\begin{Shaded}
\begin{Highlighting}[]
\CommentTok{// Clear and linear {-}{-}{-} use this exclusively}
\KeywordTok{const} name \OperatorTok{=} \KeywordTok{await} \CodeFuncTok{askQuestion}(\StringTok{"Name? "})\OperatorTok{;}
\KeywordTok{const} age \OperatorTok{=} \KeywordTok{await} \CodeFuncTok{askQuestion}(\StringTok{"Age? "})\OperatorTok{;}
\KeywordTok{const} city \OperatorTok{=} \KeywordTok{await} \CodeFuncTok{askQuestion}(\StringTok{"City? "})\OperatorTok{;}
console\OperatorTok{.}\CodeFuncTok{log}(\StringTok{\textasciigrave{}}\StringTok{\$\{}name\StringTok{\}}\StringTok{, }\StringTok{\$\{}age\StringTok{\}}\StringTok{, }\StringTok{\$\{}city\StringTok{\}}\StringTok{\textasciigrave{}})\OperatorTok{;}
\end{Highlighting}
\end{Shaded}

Teach \texttt{async}/\texttt{await} from the start and avoid callback
patterns entirely in this course.\amn{\textbf{Accessibility Tip:} Always use \texttt{async}/\texttt{await} for any operation that involves waiting (user input, file reads, network requests). The resulting linear code reads top-to-bottom the same way it executes, making it straightforward to trace with a screen reader without tracking nested callback indentation.}
Teach \texttt{async}/\texttt{await} from the start and avoid callback
patterns entirely in this course.

\begin{center}\rule{0.5\linewidth}{0.5pt}\end{center}

\subsection{3. Classroom Setup and
Configuration}\label{classroom-setup-and-configuration}

\subsection{VS Code for JavaScript}\label{vs-code-for-javascript}

\textbf{Required extensions:}

\begin{itemize}
\tightlist
\item
  \textbf{ESLint} --- catches type issues, undefined variables, and
  style violations before running code before running code
\marginnote{ESLint is a first-class accessibility tool for blind JavaScript developers.}\amn{\textbf{Accessibility Tip:} Install ESLint and configure it to run automatically. In VS Code with Screen Reader Optimized Mode, ESLint warnings are announced as you type. This gives you real-time audio feedback about type errors, undefined variables, and style violations---equivalent to visual syntax highlighting for sighted developers.}
\marginnote{ESLint is a first-class accessibility tool for blind JavaScript developers. It converts silent runtime failures (wrong type, undefined variable, loose equality) into audible editor warnings. Configure it before writing a single line of JavaScript---not as an optional style tool but as a safety layer.}
\item
  \textbf{Prettier} --- auto-formats code consistently (run on save with
  \texttt{editor.formatOnSave:\ true})
\item
  \textbf{JavaScript (ES6) code snippets} --- tab-complete common
  patterns
\end{itemize}

\textbf{Recommended ESLint configuration} (\texttt{.eslintrc.json} in
project root):

\begin{Shaded}
\begin{Highlighting}[]
\CodeFuncTok{\{}
  \CodeTypeTok{"env"}\CodeFuncTok{:} \CodeFuncTok{\{} \CodeTypeTok{"node"}\CodeFuncTok{:} \KeywordTok{true}\CodeFuncTok{,} \CodeTypeTok{"es2021"}\CodeFuncTok{:} \KeywordTok{true} \CodeFuncTok{\},}
  \CodeTypeTok{"extends"}\CodeFuncTok{:} \StringTok{"eslint:recommended"}\CodeFuncTok{,}
  \CodeTypeTok{"rules"}\CodeFuncTok{:} \CodeFuncTok{\{}
    \CodeTypeTok{"eqeqeq"}\CodeFuncTok{:} \StringTok{"error"}\CodeFuncTok{,}
    \CodeTypeTok{"no{-}unused{-}vars"}\CodeFuncTok{:} \StringTok{"warn"}\CodeFuncTok{,}
    \CodeTypeTok{"no{-}undef"}\CodeFuncTok{:} \StringTok{"error"}\CodeFuncTok{,}
    \CodeTypeTok{"prefer{-}const"}\CodeFuncTok{:} \StringTok{"warn"}\CodeFuncTok{,}
    \CodeTypeTok{"no{-}console"}\CodeFuncTok{:} \StringTok{"off"}
  \CodeFuncTok{\}}
\CodeFuncTok{\}}
\end{Highlighting}
\end{Shaded}

The \texttt{eqeqeq:\ error} rule makes \texttt{==} a linting error ---
it must be \texttt{===}. This is an automated enforcement of the rule
that prevents the most common silent bugs.

\textbf{Node.js REPL accessibility:}

\begin{Shaded}
\begin{Highlighting}[]
node                      \CommentTok{\# Start REPL}
.help                     \CommentTok{\# List REPL commands}
.exit                     \CommentTok{\# Quit (or Ctrl+D)}
\end{Highlighting}
\end{Shaded}

The Node.js REPL is excellent for early exploration. Results are printed
immediately after each expression, providing audio feedback without the
compile-run-read cycle.

\subsection{Alternative JavaScript
Environments}\label{alternative-javascript-environments}

\begin{fullwidth}{\def\LTcaptype{none} % do not increment counter
\begin{longtable}[]{@{}
  >{\raggedright\arraybackslash}p{(\linewidth - 6\tabcolsep) * \real{0.2903}}
  >{\raggedright\arraybackslash}p{(\linewidth - 6\tabcolsep) * \real{0.0860}}
  >{\raggedright\arraybackslash}p{(\linewidth - 6\tabcolsep) * \real{0.1398}}
  >{\raggedright\arraybackslash}p{(\linewidth - 6\tabcolsep) * \real{0.4839}}@{}}
\toprule\noalign{}
\begin{minipage}[b]{\linewidth}\raggedright
Environment
\end{minipage} & \begin{minipage}[b]{\linewidth}\raggedright
Platform
\end{minipage} & \begin{minipage}[b]{\linewidth}\raggedright
Accessibility
\end{minipage} & \begin{minipage}[b]{\linewidth}\raggedright
Notes
\end{minipage} \\
\midrule\noalign{}
\endhead
\bottomrule\noalign{}
\endlastfoot
\textbf{Node.js + VS Code} & All & High & Recommended setup \\
\textbf{Node.js + terminal only} & All & High & For students who prefer
minimal UI \\
\textbf{Replit.com} & Web & Medium & No-install; test with student's
setup \\
\textbf{CodeSandbox} & Web & Medium & Cloud editor; variable
accessibility \\
\textbf{GitHub Codespaces} & Web & High & Cloud VS Code; same
accessibility as local \\
\textbf{Glitch.com} & Web & Medium & Project-based; variable
accessibility \\
\textbf{Observable} & Web & Low & Notebook-style; primarily visual \\
\textbf{Deno} & Runtime & High & Alternative to Node.js; similar
accessibility \\
\end{longtable}
}
\end{fullwidth}

\begin{center}\rule{0.5\linewidth}{0.5pt}\end{center}

\subsection{4. Spoken JavaScript
Conventions}\label{spoken-javascript-conventions}

\begin{fullwidth}{\def\LTcaptype{none} % do not increment counter
\begin{longtable}[]{@{}
  >{\raggedright\arraybackslash}p{(\linewidth - 2\tabcolsep) * \real{0.3333}}
  >{\raggedright\arraybackslash}p{(\linewidth - 2\tabcolsep) * \real{0.6667}}@{}}
\toprule\noalign{}
\begin{minipage}[b]{\linewidth}\raggedright
Symbol/Pattern
\end{minipage} & \begin{minipage}[b]{\linewidth}\raggedright
Spoken Form
\end{minipage} \\
\midrule\noalign{}
\endhead
\bottomrule\noalign{}
\endlastfoot
\texttt{const} & ``const'' \\
\texttt{let} & ``let'' \\
\texttt{var} & ``var'' (avoid; legacy) \\
\texttt{===} & ``triple equals'' or ``strict equal'' \\
\texttt{!==} & ``strict not equal'' \\
\texttt{==} & ``loose equal'' (should hear this rarely) \\
\texttt{=\textgreater{}} & ``arrow'' or ``fat arrow'' \\
\texttt{...arr} & ``spread arr'' \\
\texttt{...args} & ``rest args'' \\
\texttt{async} & ``async'' \\
\texttt{await} & ``await'' \\
\texttt{undefined} & ``undefined'' \\
\texttt{null} & ``null'' \\
\texttt{NaN} & ``NaN'' (say each letter N-A-N) \\
\texttt{typeof} & ``typeof'' \\
\texttt{instanceof} & ``instanceof'' \\
\texttt{new} & ``new'' \\
\texttt{this} & ``this'' \\
\texttt{class} & ``class'' \\
\texttt{extends} & ``extends'' \\
\texttt{super} & ``super'' \\
\texttt{import} & ``import'' \\
\texttt{export} & ``export'' \\
\texttt{require()} & ``require'' \\
\texttt{module.exports} & ``module dot exports'' \\
\texttt{Promise} & ``Promise'' \\
\texttt{async\ function} & ``async function'' \\
\texttt{()\ =\textgreater{}\ \{\}} & ``arrow function'' \\
\texttt{(x)\ =\textgreater{}\ x\ *\ 2} & ``arrow function taking x,
returning x times 2'' \\
\texttt{\#field} & ``hash field'' or ``private field'' \\
\texttt{\textasciigrave{}template\ \$\{var\}\textasciigrave{}} &
``backtick template literal with variable var'' \\
\end{longtable}
}
\end{fullwidth}

\textbf{Reading an arrow function:}

\begin{Shaded}
\begin{Highlighting}[]
\KeywordTok{const} doubleIfPositive \OperatorTok{=} (n) \KeywordTok{=\textgreater{}} (n \OperatorTok{\textgreater{}} \CodeNumTok{0} \OperatorTok{?} n \OperatorTok{*} \CodeNumTok{2} \OperatorTok{:} n)\OperatorTok{;}
\end{Highlighting}
\end{Shaded}

\emph{``const doubleIfPositive gets arrow function taking n, returning:
if n greater than 0, n times 2, else n''}

\begin{center}\rule{0.5\linewidth}{0.5pt}\end{center}

\subsection{5. Unit 0--1: Setup and
Foundations}\label{unit-01-setup-and-foundations}

\subsection{Lesson 0.1: Installing
Node.js}\label{lesson-0.1-installing-node.js}

\paragraph{Instructional Cues}\label{instructional-cues-11}

\textbf{Verify installation at every step:}

\begin{Shaded}
\begin{Highlighting}[]
node \CodeTypeTok{{-}{-}version}    \CommentTok{\# Should print v18.x.x or later}
npm \CodeTypeTok{{-}{-}version}     \CommentTok{\# Should print a version number}
npx \CodeTypeTok{{-}{-}version}     \CommentTok{\# Should print a version number (npx is bundled with npm)}
\end{Highlighting}
\end{Shaded}

All three must work. If \texttt{npm} is missing, the student may have
installed only the Node.js runtime without npm. Reinstall from
nodejs.org (npm is included in the standard installer).

\textbf{REPL first --- files second:} Begin with \texttt{node} in the
terminal. Let the student type expressions and hear results immediately.
This builds confidence before introducing the file -\textgreater{} run
-\textgreater{} read cycle:

\begin{verbatim}
$ node
> 5 + 3
8
> "Hello " + "World"
'Hello World'
> typeof 42
'number'
> typeof "hello"
'string'
> .exit
\end{verbatim}

\textbf{The npm init conversation:} When initializing a project,
\texttt{npm\ init\ -y} generates \texttt{package.json} automatically.
Walk the student through reading the generated file:

\begin{Shaded}
\begin{Highlighting}[]
npm init \CodeTypeTok{{-}y}
\CodeFuncTok{cat} package.json   \CommentTok{\# or type package.json on Windows}
\end{Highlighting}
\end{Shaded}

\texttt{package.json} is plain JSON text --- fully readable with a
screen reader. Understanding its structure is important for later
lessons about dependencies and scripts.

\paragraph{Non-Obvious Accessibility
Strategies}\label{non-obvious-accessibility-strategies-2}

\begin{itemize}
\tightlist
\item
  \textbf{Node.js version managers:} \texttt{nvm} (Node Version Manager)
  allows multiple Node.js versions. For a student who may work on
  multiple machines or projects, \texttt{nvm\ use\ 18} is a consistent
  way to set the Node.js version. \texttt{nvm} is command-line-only and
  fully accessible.
\item
  \textbf{Global vs.~local packages:} \texttt{npm\ install\ -g} installs
  globally; \texttt{npm\ install} installs locally (in
  \texttt{node\_modules}). Teach local installation as the default.
  Global packages can interfere across projects.
\end{itemize}

\paragraph{Thought Experiments}\label{thought-experiments-18}

\begin{enumerate}
\def\labelenumi{\arabic{enumi}.}
\item
  \emph{``JavaScript was created in 10 days in 1995. What marks of that
  hasty design are still visible in modern JavaScript? (Think about
  \texttt{typeof\ null\ ===\ \textquotesingle{}object\textquotesingle{}},
  \texttt{NaN\ !==\ NaN}, \texttt{0\ ==\ false}.) Why haven't these been
  fixed?''}
\item
  \emph{``Node.js uses a single-threaded event loop to handle many
  concurrent operations. A traditional multi-threaded approach uses
  multiple threads running in parallel. For a web server handling
  thousands of requests, what are the tradeoffs of each approach? Which
  tends to be more memory efficient?''}
\item
  \emph{``npm has millions of packages. The average Node.js project
  depends on hundreds of packages, many maintained by single developers
  with no organizational backing. The `left-pad incident' in 2016 broke
  thousands of projects when one 11-line package was unpublished. What
  does this incident reveal about the risks and benefits of the npm
  ecosystem?''}
\item
  \emph{``JavaScript runs in browsers, in Node.js servers, in mobile
  apps via React Native, in desktop apps via Electron, in IoT devices,
  and even in serverless cloud functions. What does this ubiquity tell
  you about the relationship between a language's design and its
  adoption?''}
\end{enumerate}

\begin{center}\rule{0.5\linewidth}{0.5pt}\end{center}

\subsection{Lesson 1.3: Variables and Data
Types}\label{lesson-1.3-variables-and-data-types}

\paragraph{Instructional Cues}\label{instructional-cues-12}

\textbf{\texttt{const} vs.~\texttt{let} --- establish the rule on day
one:}

\emph{``Use \texttt{const} for everything. If you get an error saying
you're trying to reassign a \texttt{const}, THEN change it to
\texttt{let}. This approach produces cleaner code and prevents
accidental reassignment bugs.''}

Demonstrate the error:

\begin{Shaded}
\begin{Highlighting}[]
\KeywordTok{const} x \OperatorTok{=} \CodeNumTok{5}\OperatorTok{;}
x \OperatorTok{=} \CodeNumTok{10}\OperatorTok{;} \CommentTok{// TypeError: Assignment to constant variable.}
\end{Highlighting}
\end{Shaded}

\emph{``This error is helpful. \texttt{const} is protecting you from
changing something you said was fixed.''}

\textbf{Dynamic typing with auditory ambiguity:}

Run this sequence in the REPL and ask the student to predict each output
before pressing Enter:

\begin{Shaded}
\begin{Highlighting}[]
\KeywordTok{typeof} \CodeNumTok{42}\OperatorTok{;} \CommentTok{// "number"}
\KeywordTok{typeof} \StringTok{"42"}\OperatorTok{;} \CommentTok{// "string"}
\KeywordTok{typeof} \KeywordTok{true}\OperatorTok{;} \CommentTok{// "boolean"}
\KeywordTok{typeof} \KeywordTok{null}\OperatorTok{;} \CommentTok{// "object"  \textless{}{-} the famous bug}
\KeywordTok{typeof} \KeywordTok{undefined}\OperatorTok{;} \CommentTok{// "undefined"}
\KeywordTok{typeof} \{\}\OperatorTok{;} \CommentTok{// "object"}
\KeywordTok{typeof} []\OperatorTok{;} \CommentTok{// "object"  \textless{}{-} arrays are objects!}
\KeywordTok{typeof} \KeywordTok{function} () \{\}\OperatorTok{;} \CommentTok{// "function"}
\end{Highlighting}
\end{Shaded}

\emph{``Notice that both \texttt{\{\}} (object) and \texttt{{[}{]}}
(array) produce `object'. How do you tell them apart?
\texttt{Array.isArray({[}{]})} returns \texttt{true}.
\texttt{Array.isArray(\{\})} returns \texttt{false}.''}

\textbf{Template literals vs.~concatenation:}

\begin{Shaded}
\begin{Highlighting}[]
\KeywordTok{const} name \OperatorTok{=} \StringTok{"Alice"}\OperatorTok{;}
\KeywordTok{const} age \OperatorTok{=} \CodeNumTok{15}\OperatorTok{;}
 
\CommentTok{// Concatenation {-}{-}{-} confusing order of operations when read aloud}
console\OperatorTok{.}\CodeFuncTok{log}(\StringTok{"Hello, "} \OperatorTok{+} name \OperatorTok{+} \StringTok{"! You are "} \OperatorTok{+} age \OperatorTok{+} \StringTok{" years old."})\OperatorTok{;}
 
\CommentTok{// Template literal {-}{-}{-} linear, clean, easy to hear}
console\OperatorTok{.}\CodeFuncTok{log}(\StringTok{\textasciigrave{}Hello, }\StringTok{\$\{}name\StringTok{\}}\StringTok{! You are }\StringTok{\$\{}age\StringTok{\}}\StringTok{ years old.\textasciigrave{}})\OperatorTok{;}
\end{Highlighting}
\end{Shaded}

Read both versions aloud with the student. Ask which is easier to follow
aurally. Template literals are almost universally clearer when heard,
because the variable values appear where they belong in the sentence.

\paragraph{Thought Experiments}\label{thought-experiments-19}

\begin{enumerate}
\def\labelenumi{\arabic{enumi}.}
\item
  \emph{``JavaScript has only one number type for integers and decimals.
  Python has separate \texttt{int} and \texttt{float}. Java has
  \texttt{int}, \texttt{long}, \texttt{double}, \texttt{float}. What
  precision problems arise from having a single number type? When does
  \texttt{0.1\ +\ 0.2\ =\ 0.30000000000000004} matter, and when doesn't
  it?''}
\item
  \emph{``JavaScript has six falsy values: \texttt{false}, \texttt{0},
  \texttt{""}, \texttt{null}, \texttt{undefined}, \texttt{NaN}.
  Everything else is truthy --- including \texttt{"false"} (a non-empty
  string), \texttt{{[}{]}} (an empty array), and \texttt{\{\}} (an empty
  object). Design a quiz: ask the student whether ten different
  expressions are truthy or falsy. Which answers are surprising?''}
\item
  \emph{``\texttt{var} has function scope; \texttt{let} and
  \texttt{const} have block scope. What does `scope' mean? How does a
  variable declared inside an \texttt{if} block with \texttt{var}
  vs.~\texttt{let} behave differently? Which behavior is more
  predictable and why?''}
\end{enumerate}

\begin{center}\rule{0.5\linewidth}{0.5pt}\end{center}

\subsection{6. Unit 2: Numbers, Input, and
Math}\label{unit-2-numbers-input-and-math}

\subsection{Lesson 2.2: User Input in
Node.js}\label{lesson-2.2-user-input-in-node.js}

\paragraph{Instructional Cues}\label{instructional-cues-13}

\textbf{The input helper module pattern:}

Provide the \texttt{input.js} module in every project skeleton. Students
should not need to re-create this boilerplate:

\begin{Shaded}
\begin{Highlighting}[]
\CommentTok{// input.js}
\KeywordTok{const} readline \OperatorTok{=} \CodeAccentTok{require}(\StringTok{"readline"})\OperatorTok{;}
\KeywordTok{const} rl \OperatorTok{=} readline\OperatorTok{.}\CodeFuncTok{createInterface}(\{\}
  \CodeTypeTok{input}\OperatorTok{:} process\OperatorTok{.}\CodeTypeTok{stdin}\OperatorTok{,}
  \CodeTypeTok{output}\OperatorTok{:} process\OperatorTok{.}\CodeTypeTok{stdout}\OperatorTok{,}
)\OperatorTok{;}
 
\KeywordTok{function} \CodeFuncTok{askQuestion}(prompt) \{\}
  \KeywordTok{return} \KeywordTok{new} Promise((resolve) \KeywordTok{=\textgreater{}} rl\OperatorTok{.}\CodeFuncTok{question}(prompt\OperatorTok{,} resolve))\OperatorTok{;}
\
 
\KeywordTok{function} \CodeFuncTok{close}() \{\}
  rl\OperatorTok{.}\CodeFuncTok{close}()\OperatorTok{;}
\
 
module\OperatorTok{.}\CodeTypeTok{exports} \OperatorTok{=} \{ askQuestion\}\OperatorTok{,} close \\OperatorTok{;}
\end{Highlighting}
\end{Shaded}

Walk through what each line does at least once so students understand
the structure, but do not require students to write it from scratch
until they're comfortable with async programming.

\textbf{The \texttt{async}/\texttt{await} mental model --- teach this
story:}

\emph{``When your code hits \texttt{await}, imagine it puts up a `please
wait' sign and steps aside. Other things can happen while you wait. When
the answer arrives, your code picks up exactly where it paused. The
\texttt{async} keyword on the function declaration is your promise that
`this function may pause at some point.'\,''}

\begin{Shaded}
\begin{Highlighting}[]
\KeywordTok{async} \KeywordTok{function} \CodeFuncTok{main}() \{\}
  console\OperatorTok{.}\CodeFuncTok{log}(\StringTok{"Starting..."})\OperatorTok{;} \CommentTok{// 1: runs immediately}
  \KeywordTok{const} name \OperatorTok{=} \KeywordTok{await} \CodeFuncTok{askQuestion}(\StringTok{"Name? "})\OperatorTok{;} \CommentTok{// 2: pauses here until user types}
  console\OperatorTok{.}\CodeFuncTok{log}(\StringTok{"Got name:"}\OperatorTok{,} name)\OperatorTok{;} \CommentTok{// 3: runs after user answers}
  \KeywordTok{const} age \OperatorTok{=} \KeywordTok{await} \CodeFuncTok{askQuestion}(\StringTok{"Age? "})\OperatorTok{;} \CommentTok{// 4: pauses again}
  console\OperatorTok{.}\CodeFuncTok{log}(\StringTok{\textasciigrave{}}\StringTok{\$\{}name\StringTok{\}}\StringTok{ is }\StringTok{\$\{}age\StringTok{\}}\StringTok{\textasciigrave{}})\OperatorTok{;} \CommentTok{// 5: runs after user answers again}
  \CodeFuncTok{close}()\OperatorTok{;}
\
 
\CodeFuncTok{main}()\OperatorTok{;} \CommentTok{// Kick off the async function}
\end{Highlighting}
\end{Shaded}

Ask: ``What line is running while the program waits for user input?''
(None --- the function is paused. The event loop can handle other
things.)

\textbf{Why \texttt{main()} must be called at the bottom:}
\texttt{async\ function\ main()} defines the function but doesn't run
it. The \texttt{main()} call at the end actually starts execution.
Without the call, nothing happens. This trips up beginners ---
demonstrate it.

\paragraph{Non-Obvious Accessibility
Strategies}\label{non-obvious-accessibility-strategies-3}

\begin{itemize}
\item
  \textbf{\texttt{rl.close()} is mandatory:} Without
  \texttt{rl.close()}, the Node.js process continues running after
  \texttt{main()} finishes. The terminal appears to hang --- cursor
  blinks but no prompt. For a screen reader user, this is confusing.
  Teach \texttt{close()} as the final step of every program that uses
  input.
\item
  \textbf{Prompt trailing space:} Always include a trailing space in
  prompt strings: \texttt{"Enter\ your\ name:\ "}. Without the space,
  the cursor appears immediately after the colon with no separation,
  which some screen readers render as ``Enter your name:cursor'' --- an
  awkward audio experience.
\item
  \textbf{Input validation loop:} Without validation, a blank Enter or a
  non-numeric string for a number causes silent failure. Provide this
  reusable pattern:
\end{itemize}

\begin{Shaded}
\begin{Highlighting}[]
\KeywordTok{async} \KeywordTok{function} \CodeFuncTok{getNumber}(prompt) \{\}
  \KeywordTok{while} (\KeywordTok{true}) \{\}
    \KeywordTok{const} input \OperatorTok{=} \KeywordTok{await} \CodeFuncTok{askQuestion}(prompt)\OperatorTok{;}
    \KeywordTok{const} num \OperatorTok{=} Number(input\OperatorTok{.}\CodeFuncTok{trim}())\OperatorTok{;}
    \KeywordTok{if} (\OperatorTok{!}\CodeAccentTok{isNaN}(num)) \KeywordTok{return} num\OperatorTok{;}
    console\OperatorTok{.}\CodeFuncTok{log}(\StringTok{"Invalid input: please enter a number."})\OperatorTok{;}
  \}
\
\end{Highlighting}
\end{Shaded}

Teach this pattern and require students to use it whenever they need
numeric input.

\paragraph{Thought Experiments}\label{thought-experiments-20}

\begin{enumerate}
\def\labelenumi{\arabic{enumi}.}
\item
  \emph{``The \texttt{readline} module uses an `event listener' pattern
  --- \texttt{rl.question(prompt,\ callback)} means `when the user
  answers, call this function.' The \texttt{async}/\texttt{await}
  pattern wraps this in a Promise. What problem does the Promise solve
  that the raw callback doesn't? (Think about reading the code
  linearly.)''}
\item
  \emph{``Node.js uses a single-threaded event loop. While your program
  is waiting for user input (\texttt{await\ askQuestion}), the thread is
  free. What else could the thread be doing? In a real application, this
  might mean handling a different request or updating state. What does
  this tell you about why \texttt{async}/\texttt{await} is important for
  scalable programs?''}
\item
  \emph{``Synchronous input (blocking until the user responds) and
  asynchronous input (doing other things while waiting) represent
  fundamentally different programming models. For a single-user
  command-line program, which matters? For a server handling 10,000
  simultaneous users, which matters more? Why?''}
\end{enumerate}

\begin{center}\rule{0.5\linewidth}{0.5pt}\end{center}

\subsection{7. Unit 3: Control Flow}\label{unit-3-control-flow}

\subsection{Lessons 3.1--3.3}\label{lessons-3.13.3}

\paragraph{Instructional Cues}\label{instructional-cues-14}

\textbf{\texttt{===} enforcement --- make it automatic:}

With ESLint configured to flag \texttt{==}, the student will hear a
warning whenever they use loose equality. Walk through the first warning
together:

\begin{Shaded}
\begin{Highlighting}[]
\KeywordTok{if} (userInput \OperatorTok{==} \CodeNumTok{5}) \{   \CommentTok{// ESLint: Expected \textquotesingle{}===\textquotesingle{} and instead saw \textquotesingle{}==\textquotesingle{}}
\end{Highlighting}
\end{Shaded}

\emph{``ESLint is protecting you from a common silent bug. Change this
to \texttt{===}. The only exception is checking for null-or-undefined:
\texttt{x\ ==\ null} is sometimes used as a shorthand for
\texttt{x\ ===\ null\ \textbar{}\textbar{}\ x\ ===\ undefined} --- but
avoid this until you understand it.''}
\sidenote{\texttt{===} vs \texttt{==} is the JavaScript equivalent of Python's type confusion between \texttt{"42"} and \texttt{42}. Just as blind students are trained to use \texttt{type()} in Python, they should use \texttt{===} as the default in JavaScript to avoid coercion surprises that are invisible when reading code without running it.}

\textbf{Truthy/falsy in input validation:}

\begin{Shaded}
\begin{Highlighting}[]
\KeywordTok{const} name \OperatorTok{=} \KeywordTok{await} \CodeFuncTok{askQuestion}(\StringTok{"Enter your name: "})\OperatorTok{;}
\KeywordTok{if} (\OperatorTok{!}name\OperatorTok{.}\CodeFuncTok{trim}()) \{
  console\OperatorTok{.}\CodeFuncTok{log}(\StringTok{"Name cannot be empty."})\OperatorTok{;}
  \CommentTok{// re{-}prompt...}
\}
\end{Highlighting}
\end{Shaded}

\texttt{!name.trim()} is true when the trimmed string is empty
(\texttt{""} is falsy). This is idiomatic JavaScript. Read it aloud:
``if not name dot trim'' --- ask: ``What does `not an empty string'
mean? Why is an empty string falsy?''

\textbf{Ternary operator --- limit its use:}

The ternary \texttt{condition\ ?\ valueIfTrue\ :\ valueIfFalse} is
concise but can be hard to follow when read aloud:

\begin{Shaded}
\begin{Highlighting}[]
\KeywordTok{const} status \OperatorTok{=} score \OperatorTok{\textgreater{}=} \CodeNumTok{60} \OperatorTok{?} \StringTok{"pass"} \OperatorTok{:} \StringTok{"fail"}\OperatorTok{;}
\end{Highlighting}
\end{Shaded}

Read: \emph{``const status gets: if score greater-than-or-equal 60,
quote pass, else quote fail''}

Teach it but establish this rule: ``If the ternary is easy to read aloud
in one breath, use it. If it requires re-reading, use if-else.''

\paragraph{Thought Experiments}\label{thought-experiments-21}

\begin{enumerate}
\def\labelenumi{\arabic{enumi}.}
\item
  \emph{``JavaScript's \texttt{==} performs type coercion before
  comparing. TypeScript (a typed superset of JavaScript) was partly
  created to address this. What does TypeScript's existence tell you
  about the real-world cost of JavaScript's permissiveness? What did
  developers gain from TypeScript that JavaScript didn't provide?''}
\item
  \emph{``Consider: \texttt{if\ (userName)}
  vs.~\texttt{if\ (userName\ !==\ \textquotesingle{}\textquotesingle{}\ \&\&\ userName\ !==\ null\ \&\&\ userName\ !==\ undefined)}.
  The first uses JavaScript's truthiness; the second is explicit. In a
  code review, which would you prefer seeing? What are the readability
  and correctness tradeoffs?''}
\item
  \emph{``JavaScript's optional chaining (\texttt{?.}) and nullish
  coalescing (\texttt{??}) were added to the language to handle
  null/undefined more gracefully. \texttt{user?.address?.city} returns
  \texttt{undefined} instead of throwing if \texttt{user} or
  \texttt{address} is null. Why did this feature take decades to add to
  a language created in 1995?''}
\item
  \emph{``The \texttt{switch} statement falls through by default in
  JavaScript. Modern JavaScript has been moving away from
  \texttt{switch} toward mapping objects for dispatch tables.
  \texttt{const\ actions\ =\ \{\ \textquotesingle{}add\textquotesingle{}:\ addFunc,\ \textquotesingle{}delete\textquotesingle{}:\ deleteFunc\ \};\ actions{[}command{]}?.();}.
  What are the accessibility implications of this pattern vs.~a
  \texttt{switch} statement when read with a screen reader?''}
\end{enumerate}

\begin{center}\rule{0.5\linewidth}{0.5pt}\end{center}

\subsection{8. Unit 4: Arrays and
Objects}\label{unit-4-arrays-and-objects}

\subsection{Lesson 4.1: Arrays and Functional
Methods}\label{lesson-4.1-arrays-and-functional-methods}

\paragraph{Instructional Cues}\label{instructional-cues-15}

\textbf{Introduce \texttt{map}, \texttt{filter}, \texttt{reduce} as a
trio:}


\marginnote{\texttt{map}, \texttt{filter}, and \texttt{reduce} originate in functional programming (Lisp, Haskell). JavaScript borrowed them from these traditions. Teaching the trio together as a conceptual unit---transform, select, aggregate---is more durable than teaching individual method signatures.}

\emph{``These three methods each take a function and apply it to every
element in an array. \texttt{map} transforms elements. \texttt{filter}
keeps or removes elements. \texttt{reduce} combines elements into one
value. Together they cover most array processing.''}

Then give one shared example that uses all three:

\begin{Shaded}
\begin{Highlighting}[]
\KeywordTok{const} students \OperatorTok{=} [
  \{ \}\CodeTypeTok{name}\OperatorTok{:} \StringTok{"Alice"}\OperatorTok{,} \CodeTypeTok{score}\OperatorTok{:} \CodeNumTok{88} \}\OperatorTok{,}
  \{ \CodeTypeTok{name}\OperatorTok{:} \StringTok{"Bob"}\OperatorTok{,} \CodeTypeTok{score}\OperatorTok{:} \CodeNumTok{62} \}\OperatorTok{,}
  \{ \CodeTypeTok{name}\OperatorTok{:} \StringTok{"Charlie"}\OperatorTok{,} \CodeTypeTok{score}\OperatorTok{:} \CodeNumTok{95} \}\OperatorTok{,}
  \{ \CodeTypeTok{name}\OperatorTok{:} \StringTok{"Diana"}\OperatorTok{,} \CodeTypeTok{score}\OperatorTok{:} \CodeNumTok{71} \}\OperatorTok{,}
]\OperatorTok{;}
 
\CommentTok{// 1. Filter passing students (score \textgreater{}= 70)}
\KeywordTok{const} passing \OperatorTok{=} students\OperatorTok{.}\CodeFuncTok{filter}((s) \KeywordTok{=\textgreater{}} s\OperatorTok{.}\CodeTypeTok{score} \OperatorTok{\textgreater{}=} \CodeNumTok{70})\OperatorTok{;}
 
\CommentTok{// 2. Map to get just names}
\KeywordTok{const} passingNames \OperatorTok{=} passing\OperatorTok{.}\CodeFuncTok{map}((s) \KeywordTok{=\textgreater{}} s\OperatorTok{.}\CodeTypeTok{name})\OperatorTok{;}
 
\CommentTok{// 3. Reduce to build a report string}
\KeywordTok{const} report \OperatorTok{=} passingNames\OperatorTok{.}\CodeFuncTok{reduce}(
  (acc\OperatorTok{,} name) \KeywordTok{=\textgreater{}} acc \OperatorTok{+} name \OperatorTok{+} \StringTok{", "}\OperatorTok{,}
  \StringTok{"Passing: "}\OperatorTok{,}
)\OperatorTok{;}
 
console\OperatorTok{.}\CodeFuncTok{log}(report)\OperatorTok{;} \CommentTok{// "Passing: Alice, Charlie, Diana, "}
console\OperatorTok{.}\CodeFuncTok{log}(\StringTok{"Count:"}\OperatorTok{,} passing\OperatorTok{.}\CodeTypeTok{length})\OperatorTok{;}
\end{Highlighting}
\end{Shaded}

\textbf{Accessible printing for arrays of objects:}

\begin{Shaded}
\begin{Highlighting}[]
\CommentTok{// Accessible: numbered, labeled fields}
students\OperatorTok{.}\CodeFuncTok{forEach}((student\OperatorTok{,} index) \KeywordTok{=\textgreater{}} \{
  console\OperatorTok{.}\CodeFuncTok{log}(\StringTok{\textasciigrave{}}\StringTok{\$\{}index \OperatorTok{+} \CodeNumTok{1}\StringTok{\}}\StringTok{. }\StringTok{\$\{}student\OperatorTok{.}\CodeTypeTok{name}\StringTok{\}}\StringTok{: }\StringTok{\$\{}student\OperatorTok{.}\CodeTypeTok{score}\StringTok{\}}\StringTok{\textasciigrave{}})\OperatorTok{;}
)\OperatorTok{;}
 
\CommentTok{// Less accessible: raw object output}
console\OperatorTok{.}\CodeFuncTok{log}(students)\OperatorTok{;} \CommentTok{// Produces cluttered bracket notation when read aloud}
\end{Highlighting}
\end{Shaded}

\textbf{Chaining --- build up incrementally:}

Show chaining as a composition of steps, each of which you can trace
separately:

\begin{Shaded}
\begin{Highlighting}[]
\KeywordTok{const} result \OperatorTok{=} numbers
  \OperatorTok{.}\CodeFuncTok{filter}((n) \KeywordTok{=\textgreater{}} n \OperatorTok{\%} \CodeNumTok{2} \OperatorTok{===} \CodeNumTok{0}) \CommentTok{// Step 1: keep evens}
  \OperatorTok{.}\CodeFuncTok{map}((n) \KeywordTok{=\textgreater{}} n \OperatorTok{**} \CodeNumTok{2}) \CommentTok{// Step 2: square each}
  \OperatorTok{.}\CodeFuncTok{reduce}((sum\OperatorTok{,} n) \KeywordTok{=\textgreater{}} sum \OperatorTok{+} n\OperatorTok{,} \CodeNumTok{0})\OperatorTok{;} \CommentTok{// Step 3: sum}
 
\CommentTok{// DEBUG: break the chain to trace each step}
\KeywordTok{const} step1 \OperatorTok{=} numbers\OperatorTok{.}\CodeFuncTok{filter}((n) \KeywordTok{=\textgreater{}} n \OperatorTok{\%} \CodeNumTok{2} \OperatorTok{===} \CodeNumTok{0})\OperatorTok{;}
console\OperatorTok{.}\CodeFuncTok{log}(\StringTok{"After filter:"}\OperatorTok{,} step1)\OperatorTok{;}
\KeywordTok{const} step2 \OperatorTok{=} step1\OperatorTok{.}\CodeFuncTok{map}((n) \KeywordTok{=\textgreater{}} n \OperatorTok{**} \CodeNumTok{2})\OperatorTok{;}
console\OperatorTok{.}\CodeFuncTok{log}(\StringTok{"After map:"}\OperatorTok{,} step2)\OperatorTok{;}
\KeywordTok{const} step3 \OperatorTok{=} step2\OperatorTok{.}\CodeFuncTok{reduce}((sum\OperatorTok{,} n) \KeywordTok{=\textgreater{}} sum \OperatorTok{+} n\OperatorTok{,} \CodeNumTok{0})\OperatorTok{;}
console\OperatorTok{.}\CodeFuncTok{log}(\StringTok{"Final sum:"}\OperatorTok{,} step3)\OperatorTok{;}
\end{Highlighting}
\end{Shaded}

The debug version (step-by-step) is more accessible during development.
Consolidate to the chain after confirming correctness.

\paragraph{Thought Experiments}\label{thought-experiments-22}

\begin{enumerate}
\def\labelenumi{\arabic{enumi}.}
\item
  \emph{``\texttt{map}, \texttt{filter}, and \texttt{reduce} come from
  functional programming --- a paradigm based on mathematical functions
  that have no side effects. JavaScript arrays also have
  \texttt{forEach} which is for side effects only. What is the
  philosophical distinction between \texttt{map} and \texttt{forEach}?
  Why does this distinction matter for reasoning about code
  correctness?''}
\item
  \emph{``The \texttt{sort} method in JavaScript sorts in-place and
  mutates the original array. \texttt{filter}, \texttt{map}, and
  \texttt{reduce} return new arrays, leaving the original unchanged. Why
  the asymmetry? What bugs could arise from \texttt{sort}'s mutation?
  How would you sort without mutating:
  \texttt{{[}...arr{]}.sort(...)}?''}
\item
  \_``Functional composition:
  \texttt{const\ double\ =\ n\ =\textgreater{}\ n\ \_\ 2;\ const\ addOne\ =\ n\ =\textgreater{}\ n\ +\ 1;\ const\ doubleAndAdd\ =\ n\ =\textgreater{}\ addOne(double(n))}.
  How does function composition scale? What is the advantage of
  composing small, focused functions over writing one large function
  that does everything?''*
\item
  \emph{``The spread operator \texttt{{[}...arr{]}} creates a shallow
  copy of an array. For an array of primitives \texttt{{[}1,\ 2,\ 3{]}},
  the spread copy is fully independent. For an array of objects
  \texttt{{[}\{name:\ \textquotesingle{}A\textquotesingle{}\},\ \{name:\ \textquotesingle{}B\textquotesingle{}\}{]}},
  the objects inside are still shared. What does `shallow copy' mean?
  When does this distinction cause bugs?''}
\end{enumerate}

\begin{center}\rule{0.5\linewidth}{0.5pt}\end{center}

\subsection{Lesson 4.2: Objects}\label{lesson-4.2-objects}

\paragraph{Instructional Cues}\label{instructional-cues-16}

\textbf{Destructuring --- introduce before objects are complex:}

\begin{Shaded}
\begin{Highlighting}[]
\KeywordTok{const} person \OperatorTok{=} \{ \}\CodeTypeTok{name}\OperatorTok{:} \StringTok{"Alice"}\OperatorTok{,} \CodeTypeTok{age}\OperatorTok{:} \CodeNumTok{15}\OperatorTok{,} \CodeTypeTok{city}\OperatorTok{:} \StringTok{"Portland"} \}\OperatorTok{;}
 
\CommentTok{// Without destructuring:}
\KeywordTok{const} name \OperatorTok{=} person\OperatorTok{.}\CodeTypeTok{name}\OperatorTok{;}
\KeywordTok{const} age \OperatorTok{=} person\OperatorTok{.}\CodeTypeTok{age}\OperatorTok{;}
 
\CommentTok{// With destructuring:}
\KeywordTok{const} \{ name\OperatorTok{,} age \} \OperatorTok{=} person\OperatorTok{;}
 
\CommentTok{// In function parameter:}
\KeywordTok{function} \CodeFuncTok{introduce}(\{ name\OperatorTok{,} age ) \{
  console\OperatorTok{.}\CodeFuncTok{log}(\StringTok{\textasciigrave{}Hi, I\textquotesingle{}m }\StringTok{\$\{}name\StringTok{\}}\StringTok{, }\StringTok{\$\{}age\StringTok{\}}\StringTok{ years old.\textasciigrave{}})\OperatorTok{;}
\}
\CodeFuncTok{introduce}(person)\OperatorTok{;}
\end{Highlighting}
\end{Shaded}

 object access code
, which in turn makes
\sidenote{Destructuring also appears in function parameters (\texttt{function f(\{ name, age \})}), in \texttt{for...of} loops, and in import statements (\texttt{import \{ useState \} from 'react'}). Introducing the concept once and pointing out these variations as they arise is more efficient than re-teaching each form separately.}
programs cleaner to read and hear.

\textbf{\texttt{Object.keys()}, \texttt{Object.values()},
\texttt{Object.entries()} as navigation tools:}


\marginnote{\texttt{Object.keys()}, \texttt{.values()}, and \texttt{.entries()} are the JS equivalent}\amn{\textbf{Accessibility Tip:} When you receive an unfamiliar object, run \texttt{console.log(Object.keys(obj).join(", "))} to hear all its property names. This replaces the visual object inspection that sighted developers do by expanding an object in the browser console.}
\marginnote{\texttt{Object.keys()}, \texttt{.values()}, and \texttt{.entries()} are the JS equivalent of Python's \texttt{dict.keys()}, \texttt{.values()}, and \texttt{.items()}. Students moving between the two languages should be shown this correspondence explicitly.}

\begin{Shaded}
\begin{Highlighting}[]
\KeywordTok{const} config \OperatorTok{=} \{ \}\CodeTypeTok{theme}\OperatorTok{:} \StringTok{"dark"}\OperatorTok{,} \CodeTypeTok{fontSize}\OperatorTok{:} \CodeNumTok{14}\OperatorTok{,} \CodeTypeTok{notifications}\OperatorTok{:} \KeywordTok{true} \}\OperatorTok{;}
 
\CommentTok{// What properties does this object have?}
console\OperatorTok{.}\CodeFuncTok{log}(Object\OperatorTok{.}\CodeFuncTok{keys}(config)\OperatorTok{.}\CodeFuncTok{join}(\StringTok{", "}))\OperatorTok{;}
 
\CommentTok{// What are all the values?}
Object\OperatorTok{.}\CodeFuncTok{values}(config)\OperatorTok{.}\CodeFuncTok{forEach}((v) \KeywordTok{=\textgreater{}} console\OperatorTok{.}\CodeFuncTok{log}(v))\OperatorTok{;}
 
\CommentTok{// Both together:}
Object\OperatorTok{.}\CodeFuncTok{entries}(config)\OperatorTok{.}\CodeFuncTok{forEach}(([key\OperatorTok{,} value]) \KeywordTok{=\textgreater{}} \{
  console\OperatorTok{.}\CodeFuncTok{log}(\StringTok{\textasciigrave{}}\StringTok{\$\{}key\StringTok{\}}\StringTok{: }\StringTok{\$\{}value\StringTok{\}}\StringTok{\textasciigrave{}})\OperatorTok{;}
)\OperatorTok{;}
\end{Highlighting}
\end{Shaded}

\textbf{Optional chaining for safe nested access:}

\begin{Shaded}
\begin{Highlighting}[]
\KeywordTok{const} user \OperatorTok{=} \{ \CodeTypeTok{profile}\OperatorTok{:} \{ \}\CodeTypeTok{address}\OperatorTok{:} \KeywordTok{null} \} \}\OperatorTok{;}
 
\CommentTok{// Without optional chaining {-}{-}{-} crashes if address is null}
\KeywordTok{const} city \OperatorTok{=} user\OperatorTok{.}\CodeTypeTok{profile}\OperatorTok{.}\CodeTypeTok{address}\OperatorTok{.}\CodeTypeTok{city}\OperatorTok{;} \CommentTok{// TypeError!}
 
\CommentTok{// With optional chaining {-}{-}{-} returns undefined instead}
\KeywordTok{const} city \OperatorTok{=} user\OperatorTok{?.}\CodeTypeTok{profile}\OperatorTok{?.}\CodeTypeTok{address}\OperatorTok{?.}\CodeTypeTok{city}\OperatorTok{;} \CommentTok{// undefined, no crash}
 
\CommentTok{// With nullish coalescing for a default}
\KeywordTok{const} city \OperatorTok{=} user\OperatorTok{?.}\CodeTypeTok{profile}\OperatorTok{?.}\CodeTypeTok{address}\OperatorTok{?.}\CodeTypeTok{city} \OperatorTok{??} \StringTok{"Unknown city"}\OperatorTok{;}
\end{Highlighting}
\end{Shaded}

This pattern is essential for working with real-world data where fields
may be missing.

\paragraph{Thought Experiments}\label{thought-experiments-23}

\begin{enumerate}
\def\labelenumi{\arabic{enumi}.}
\item
  \emph{``JavaScript objects can have properties added or removed at any
  time.
  \texttt{user.newField\ =\ \textquotesingle{}value\textquotesingle{}}
  just works. Java objects have a fixed set of fields defined at compile
  time. Which approach is better for representing data from an API that
  might add new fields in future versions? Which is better for ensuring
  all required fields are present?''}
\item
  \emph{``\texttt{Object.freeze(obj)} makes an object immutable --- you
  can't add, remove, or modify properties. This mimics the behavior of a
  Java \texttt{final} field or Python's \texttt{frozenset}. When would
  immutable objects be useful? What patterns become possible when you
  can guarantee an object won't change?''}
\item
  \emph{``Prototype-based inheritance means JavaScript objects can
  inherit from other objects directly, without classes. Classes
  (introduced in ES6) are `syntactic sugar' over this prototype system.
  What does this mean? What can you do with prototypes that you cannot
  easily express with class syntax?''}
\item
  \emph{``JSON (JavaScript Object Notation) is the universal data
  exchange format. It looks almost like JavaScript object literals but
  has stricter rules: all keys must be double-quoted strings, no
  functions, no comments. Why must JSON be stricter than JavaScript
  object literals? What security problems would arise if JSON allowed
  functions?''}
\end{enumerate}

\begin{center}\rule{0.5\linewidth}{0.5pt}\end{center}

\subsection{9. Unit 5: Strings}\label{unit-5-strings}

\paragraph{Instructional Cues}\label{instructional-cues-17}

\textbf{Template literals for multiline strings:}

\begin{Shaded}
\begin{Highlighting}[]
\KeywordTok{const} menu \OperatorTok{=} \StringTok{\textasciigrave{}}
\StringTok{{-}{-}{-} Main Menu {-}{-}{-}}
\StringTok{1. Add item}
\StringTok{2. View all}
\StringTok{3. Search}
\StringTok{4. Delete}
\StringTok{5. Quit}
\StringTok{{-}{-}{-}{-}{-}{-}{-}{-}{-}{-}{-}{-}{-}{-}{-}{-}{-}\textasciigrave{}}\OperatorTok{;}
 
console\OperatorTok{.}\CodeFuncTok{log}(menu)\OperatorTok{;}
\end{Highlighting}
\end{Shaded}

Template literals with newlines are much cleaner for building multiline
output than repeated \texttt{console.log()} calls. Teach this pattern
for menus and reports.

\textbf{Regular expression basics --- connect to real need:}

Before introducing regex syntax, establish the problem it solves:

\emph{``You have user input that might have extra spaces, uppercase, or
punctuation. You want just the words. Without regex, you'd call
\texttt{.trim().toLowerCase().replace(\textquotesingle{},\textquotesingle{},\ \textquotesingle{}\textquotesingle{}).replace(\textquotesingle{}.\textquotesingle{},\ \textquotesingle{}\textquotesingle{})...}
for every punctuation mark. Regex lets you express `any non-letter
character' once.''}

\begin{Shaded}
\begin{Highlighting}[]
\KeywordTok{const} messy \OperatorTok{=} \StringTok{"  Hello, World!   How are you?  "}\OperatorTok{;}
\KeywordTok{const} words \OperatorTok{=} messy
  \OperatorTok{.}\CodeFuncTok{trim}()
  \OperatorTok{.}\CodeFuncTok{toLowerCase}()
  \OperatorTok{.}\CodeFuncTok{split}(\StringTok{/}\StringTok{\textbackslash{}W+}\StringTok{/})
  \OperatorTok{.}\CodeFuncTok{filter}((w) \KeywordTok{=\textgreater{}} w\OperatorTok{.}\CodeTypeTok{length} \OperatorTok{\textgreater{}} \CodeNumTok{0})\OperatorTok{;}
console\OperatorTok{.}\CodeFuncTok{log}(words)\OperatorTok{;} \CommentTok{// [\textquotesingle{}hello\textquotesingle{}, \textquotesingle{}world\textquotesingle{}, \textquotesingle{}how\textquotesingle{}, \textquotesingle{}are\textquotesingle{}, \textquotesingle{}you\textquotesingle{}]}
\end{Highlighting}
\end{Shaded}

Teach \texttt{\textbackslash{}w+} (word characters),
\texttt{\textbackslash{}d+} (digits), \texttt{\textbackslash{}s+}
(whitespace), and the \texttt{g} and \texttt{i} flags. Don't attempt to
teach full regex --- a basic vocabulary is enough for data cleaning
tasks.

\paragraph{Thought Experiments}\label{thought-experiments-24}

\begin{enumerate}
\def\labelenumi{\arabic{enumi}.}
\item
  \emph{``Template literals use backticks \texttt{\textasciigrave{}}
  while regular strings use \texttt{"} or \texttt{\textquotesingle{}}.
  Backticks on a keyboard are typically in the upper-left corner. For
  touch typists, backticks are less natural than quotes. What keyboard
  shortcuts or text expansions might a blind programmer use to reduce
  the friction of frequently-needed but awkward characters?''}
\item
  \emph{``String immutability in JavaScript (same as Python) means every
  \texttt{.toUpperCase()}, \texttt{.replace()}, and \texttt{.slice()}
  creates a new string. For a loop that builds a long string
  incrementally, what is the efficient alternative? (Hint: consider
  \texttt{Array.push()} + \texttt{.join()} as a form of
  StringBuilder.)''}
\end{enumerate}

\begin{center}\rule{0.5\linewidth}{0.5pt}\end{center}

\subsection{10. Unit 6: Functions}\label{unit-6-functions-2}

\subsection{Extended Instructional
Cues}\label{extended-instructional-cues-4}

\textbf{The four function forms --- teach one style, acknowledge the
rest:}

\begin{Shaded}
\begin{Highlighting}[]
\CommentTok{// 1. Function declaration (hoisted {-}{-}{-} can be called before declaration)}
\KeywordTok{function} \CodeFuncTok{add}(a\OperatorTok{,} b) \{
  \KeywordTok{return} a \OperatorTok{+} b\OperatorTok{;}
\}
 
\CommentTok{// 2. Function expression (not hoisted)}
\KeywordTok{const} multiply \OperatorTok{=} \KeywordTok{function} (a\OperatorTok{,} b) \{
  \KeywordTok{return} a \OperatorTok{*} b\OperatorTok{;}
\}\OperatorTok{;}
 
\CommentTok{// 3. Arrow function {-}{-}{-} SHORT BODY (implicit return)}
\KeywordTok{const} square \OperatorTok{=} (n) \KeywordTok{=\textgreater{}} n \OperatorTok{*} n\OperatorTok{;}
 
\CommentTok{// 4. Arrow function {-}{-}{-} LONG BODY (explicit return required)}
\KeywordTok{const} absolute \OperatorTok{=} (n) \KeywordTok{=\textgreater{}} \{
  \KeywordTok{if} (n \OperatorTok{\textless{}} \CodeNumTok{0}) \KeywordTok{return} \OperatorTok{{-}}n\OperatorTok{;}
  \KeywordTok{return} n\OperatorTok{;}
\}\OperatorTok{;}
\end{Highlighting}
\end{Shaded}

\textbf{Teach style 3 (short arrow) and style 4 (long arrow) as the
primary forms.} Acknowledge 1 and 2 exist and how to read them. Do not
attempt to teach all four simultaneously.

\textbf{The implicit return trap --- a critical accessibility issue:}

\begin{Shaded}
\begin{Highlighting}[]
\CommentTok{// These DO return a value:}
\KeywordTok{const} double \OperatorTok{=} (n) \KeywordTok{=\textgreater{}} n \OperatorTok{*} \CodeNumTok{2}\OperatorTok{;} \CommentTok{// [x] implicit return}
\KeywordTok{const} greet \OperatorTok{=} (name) \KeywordTok{=\textgreater{}} \StringTok{\textasciigrave{}Hi }\StringTok{\$\{}name\StringTok{\}}\StringTok{\textasciigrave{}}\OperatorTok{;} \CommentTok{// [x] implicit return}
 
\CommentTok{// This does NOT return a value (the curly brace requires explicit return):}
\KeywordTok{const} broken \OperatorTok{=} (n) \KeywordTok{=\textgreater{}} \{
  n \OperatorTok{*} \CodeNumTok{2}\OperatorTok{;} \CommentTok{// No return statement! Returns undefined.}
\}\OperatorTok{;}
 
\CommentTok{// Fixed:}
\KeywordTok{const} fixed \OperatorTok{=} (n) \KeywordTok{=\textgreater{}} \{
  \KeywordTok{return} n \OperatorTok{*} \CodeNumTok{2}\OperatorTok{;}
\}\OperatorTok{;}
\end{Highlighting}
\end{Shaded}

When heard aloud, \texttt{n\ =\textgreater{}\ \{\ n\ *\ 2;\ \}} and
\texttt{n\ =\textgreater{}\ \{\ return\ n\ *\ 2;\ \}} sound different
--- the \texttt{return} keyword is present in one and absent in the
other. Teach students to listen for \texttt{return} inside block-bodied
arrow functions.

\textbf{Higher-order functions --- connect to previous array methods:}

\begin{Shaded}
\begin{Highlighting}[]
\CommentTok{// map() takes a function as argument}
\KeywordTok{const} doubled \OperatorTok{=} [\CodeNumTok{1}\OperatorTok{,} \CodeNumTok{2}\OperatorTok{,} \CodeNumTok{3}]\OperatorTok{.}\CodeFuncTok{map}((n) \KeywordTok{=\textgreater{}} n \OperatorTok{*} \CodeNumTok{2})\OperatorTok{;}
 
\CommentTok{// We can also pass a named function}
\KeywordTok{const} double \OperatorTok{=} (n) \KeywordTok{=\textgreater{}} n \OperatorTok{*} \CodeNumTok{2}\OperatorTok{;}
\KeywordTok{const} doubled2 \OperatorTok{=} [\CodeNumTok{1}\OperatorTok{,} \CodeNumTok{2}\OperatorTok{,} \CodeNumTok{3}]\OperatorTok{.}\CodeFuncTok{map}(double)\OperatorTok{;}
 
\CommentTok{// This works because map expects a function {-}{-}{-} it doesn\textquotesingle{}t care if it\textquotesingle{}s}
\CommentTok{// named or anonymous, arrow or traditional}
\end{Highlighting}
\end{Shaded}

\emph{``Every time you used \texttt{array.map(...)}, you were passing a
function as an argument. Higher-order functions --- functions that take
other functions --- are normal JavaScript.''}

\textbf{Closures --- establish the mental model before the code:}

\emph{``A closure is a function that remembers the environment where it
was created, even after that environment is gone. Imagine a letter in an
envelope. The letter is the function. The envelope is the environment
--- it contains variables the function needs. Even if the person who
wrote the letter is gone, the letter still has access to everything in
its envelope.''}

\begin{Shaded}
\begin{Highlighting}[]
\KeywordTok{function} \CodeFuncTok{makeCounter}(startValue \OperatorTok{=} \CodeNumTok{0}) \{
  \KeywordTok{let} count \OperatorTok{=} startValue\OperatorTok{;} \CommentTok{// This is the "envelope"}
 
  \KeywordTok{return} \{
    \CodeTypeTok{increment}\OperatorTok{:} () \KeywordTok{=\textgreater{}} \OperatorTok{++}count\OperatorTok{,}
    \CodeTypeTok{decrement}\OperatorTok{:} () \KeywordTok{=\textgreater{}} \OperatorTok{{-}{-}}count\OperatorTok{,}
    \CodeTypeTok{reset}\OperatorTok{:} () \KeywordTok{=\textgreater{}} \{
      count \OperatorTok{=} startValue\OperatorTok{;}
    \}\OperatorTok{,}
    \CodeTypeTok{getCount}\OperatorTok{:} () \KeywordTok{=\textgreater{}} count\OperatorTok{,}
  \}\OperatorTok{;}
\}
 
\KeywordTok{const} counter \OperatorTok{=} \CodeFuncTok{makeCounter}(\CodeNumTok{10})\OperatorTok{;}
counter\OperatorTok{.}\CodeFuncTok{increment}()\OperatorTok{;}
counter\OperatorTok{.}\CodeFuncTok{increment}()\OperatorTok{;}
counter\OperatorTok{.}\CodeFuncTok{increment}()\OperatorTok{;}
console\OperatorTok{.}\CodeFuncTok{log}(counter\OperatorTok{.}\CodeFuncTok{getCount}())\OperatorTok{;} \CommentTok{// 13 {-}{-}{-} count persists!}
\end{Highlighting}
\end{Shaded}

\paragraph{Thought Experiments}\label{thought-experiments-25}

\begin{enumerate}
\def\labelenumi{\arabic{enumi}.}
\item
  \emph{``Arrow functions don't have their own \texttt{this} --- they
  inherit \texttt{this} from the surrounding scope. Regular functions
  create their own \texttt{this}. This was a major source of bugs in
  pre-ES6 JavaScript. How many of the following do you think involved
  confusing \texttt{this}? (a) Typos, (b) Variable scope errors, (c)
  Async timing bugs, (d) \texttt{this} binding bugs. In JavaScript
  specifically, (d) is disproportionately common --- why?''}
\item
  \emph{``Pure functions always produce the same output for the same
  input and have no side effects. \texttt{Math.sin(x)} is pure.
  \texttt{console.log(x)} is not (it has a side effect). Why are pure
  functions easier to test? Why are pure functions easier to reason
  about in a screen-reader debugging session?''}
\item
  \emph{``Memoization caches the results of function calls so repeated
  calls with the same arguments return the cached result. Write a
  description of a memoized \texttt{fibonacci(n)} function. What is
  saved? What is the tradeoff? When would memoization be harmful rather
  than helpful?''}
\item
  \emph{``JavaScript is a multi-paradigm language: you can write
  imperative code (explicit step-by-step), object-oriented code (classes
  and methods), or functional code (pure functions and immutable data).
  For the accessible data manager project, which paradigm produces the
  most maintainable code? The most testable? The most readable when
  heard aloud?''}
\end{enumerate}

\begin{center}\rule{0.5\linewidth}{0.5pt}\end{center}

\subsection{11. Unit 7: Classes and OOP}\label{unit-7-classes-and-oop}

\paragraph{Instructional Cues}\label{instructional-cues-18}

\textbf{\texttt{toString()} is mandatory:}

\begin{Shaded}
\begin{Highlighting}[]
\KeywordTok{class} BankAccount \{
  \#balance\OperatorTok{;}
 
  \CodeFuncTok{constructor}(owner\OperatorTok{,} initialBalance \OperatorTok{=} \CodeNumTok{0}) \{
    \KeywordTok{this}\OperatorTok{.}\CodeTypeTok{owner} \OperatorTok{=} owner\OperatorTok{;}
    \KeywordTok{this}\OperatorTok{.}\#\CodeTypeTok{balance} \OperatorTok{=} initialBalance\OperatorTok{;}
  \}
 
  \CodeFuncTok{toString}() \{
    \KeywordTok{return} \StringTok{\textasciigrave{}BankAccount[owner=}\StringTok{\$\{}\KeywordTok{this}\OperatorTok{.}\CodeTypeTok{owner}\StringTok{\}}\StringTok{, balance=$}\StringTok{\$\{}\KeywordTok{this}\OperatorTok{.}\#\CodeTypeTok{balance}\OperatorTok{.}\CodeFuncTok{toFixed}(\CodeNumTok{2})\StringTok{\}}\StringTok{]\textasciigrave{}}\OperatorTok{;}
  \}
\}
 
\KeywordTok{const} account \OperatorTok{=} \KeywordTok{new} \CodeFuncTok{BankAccount}(\StringTok{"Alice"}\OperatorTok{,} \CodeNumTok{1000})\OperatorTok{;}
console\OperatorTok{.}\CodeFuncTok{log}(account\OperatorTok{.}\CodeFuncTok{toString}())\OperatorTok{;} \CommentTok{// BankAccount[owner=Alice, balance=$1000.00]}
console\OperatorTok{.}\CodeFuncTok{log}(\StringTok{\textasciigrave{}}\StringTok{\$\{}account\StringTok{\}}\StringTok{\textasciigrave{}})\OperatorTok{;} \CommentTok{// Also calls toString() implicitly}
\end{Highlighting}
\end{Shaded}

Require \texttt{toString()} in every class. Explain: ``When your screen
reader reads \texttt{console.log(myAccount)}, it reads
\texttt{toString()}. Without it, you hear \texttt{{[}object\ Object{]}}
--- completely useless. Define \texttt{toString()} to hear meaningful
information.''

\textbf{Private fields \texttt{\#} --- establish as the default:}

\begin{Shaded}
\begin{Highlighting}[]
\KeywordTok{class} Dog \{
  \#name\OperatorTok{;} \CommentTok{// Private {-}{-}{-} cannot be accessed outside the class}
  \#age\OperatorTok{;}
 
  \CodeFuncTok{constructor}(name\OperatorTok{,} age) \{
    \KeywordTok{this}\OperatorTok{.}\#\CodeTypeTok{name} \OperatorTok{=} name\OperatorTok{;}
    \KeywordTok{this}\OperatorTok{.}\#\CodeTypeTok{age} \OperatorTok{=} age\OperatorTok{;}
  \}
 
  \KeywordTok{get} \CodeFuncTok{name}() \{
    \KeywordTok{return} \KeywordTok{this}\OperatorTok{.}\#\CodeTypeTok{name}\OperatorTok{;}
  \} \CommentTok{// Getter {-}{-}{-} controlled access}
  \KeywordTok{get} \CodeFuncTok{age}() \{
    \KeywordTok{return} \KeywordTok{this}\OperatorTok{.}\#\CodeTypeTok{age}\OperatorTok{;}
  \}
 
  \KeywordTok{set} \CodeFuncTok{age}(value) \{
    \KeywordTok{if} (value \OperatorTok{\textless{}} \CodeNumTok{0}) \KeywordTok{throw} \KeywordTok{new} Error(\StringTok{"Age cannot be negative"})\OperatorTok{;}
    \KeywordTok{this}\OperatorTok{.}\#\CodeTypeTok{age} \OperatorTok{=} value\OperatorTok{;}
  \}
 
  \CodeFuncTok{bark}() \{
    console\OperatorTok{.}\CodeFuncTok{log}(\StringTok{\textasciigrave{}}\StringTok{\$\{}\KeywordTok{this}\OperatorTok{.}\#\CodeTypeTok{name}\StringTok{\}}\StringTok{ says: Woof!\textasciigrave{}})\OperatorTok{;}
  \}
 
  \CodeFuncTok{toString}() \{
    \KeywordTok{return} \StringTok{\textasciigrave{}Dog[name=}\StringTok{\$\{}\KeywordTok{this}\OperatorTok{.}\#\CodeTypeTok{name}\StringTok{\}}\StringTok{, age=}\StringTok{\$\{}\KeywordTok{this}\OperatorTok{.}\#\CodeTypeTok{age}\StringTok{\}}\StringTok{]\textasciigrave{}}\OperatorTok{;}
  \}
\}
\end{Highlighting}
\end{Shaded}

\emph{``Private fields starting with \texttt{\#} are truly private ---
not just conventional privacy. \texttt{myDog.\#name} outside the class
throws a SyntaxError. This is real encapsulation.''}

\paragraph{Thought Experiments}\label{thought-experiments-26}

\begin{enumerate}
\def\labelenumi{\arabic{enumi}.}
\item
  \emph{``JavaScript classes are `syntactic sugar' over prototype-based
  inheritance. TypeScript adds static typing on top of JavaScript
  classes. With three layers (prototypes -\textgreater{} ES6 classes
  -\textgreater{} TypeScript classes), is JavaScript's OOP model more or
  less trustworthy than Java's, which has only one layer? What does
  `syntactic sugar' imply about the behavior of the underlying
  mechanism?''}
\item
  \emph{``The \texttt{class} keyword was added to JavaScript in 2015, 20
  years after the language was created. In those 20 years, JavaScript
  developers built massive applications without classes using only
  prototypes and closures. What does this tell you about the role of
  syntax vs.~capability in programming languages?''}
\item
  \emph{``Mixins are a way to compose behavior from multiple sources
  without inheritance. Instead of \texttt{class\ A\ extends\ B}, you
  copy methods from B onto A's prototype. JavaScript supports this; Java
  does not (without interfaces). When would mixins be more flexible than
  single inheritance? What problems could they cause?''}
\end{enumerate}

\begin{center}\rule{0.5\linewidth}{0.5pt}\end{center}

\subsection{12. Unit 8: Files, Modules, and Error
Handling}\label{unit-8-files-modules-and-error-handling}

\subsection{Lesson 8.1: Files}\label{lesson-8.1-files}

\paragraph{Instructional Cues}\label{instructional-cues-19}

\textbf{\texttt{\_\_dirname} and path construction:}

\begin{Shaded}
\begin{Highlighting}[]
\KeywordTok{const} path \OperatorTok{=} \CodeAccentTok{require}(\StringTok{"path"})\OperatorTok{;}
\KeywordTok{const} fs \OperatorTok{=} \CodeAccentTok{require}(\StringTok{"fs"})\OperatorTok{;}
 
\CommentTok{// \_\_dirname is the directory of the current file}
\KeywordTok{const} dataFile \OperatorTok{=} path\OperatorTok{.}\CodeFuncTok{join}(\_\_dirname\OperatorTok{,} \StringTok{"data"}\OperatorTok{,} \StringTok{"contacts.json"})\OperatorTok{;}
 
\CommentTok{// Path joining prevents slash vs. backslash issues on Windows/macOS/Linux}
console\OperatorTok{.}\CodeFuncTok{log}(\StringTok{"Data file location:"}\OperatorTok{,} dataFile)\OperatorTok{;}
\end{Highlighting}
\end{Shaded}

For a student who cannot visually navigate a file tree,
\texttt{console.log(\_\_dirname)} and \texttt{console.log(dataFile)} are
essential diagnostic tools. Teach them to print file paths as the first
debugging step when a file isn't found.

\textbf{\texttt{fs.existsSync()} before file operations:}

\begin{Shaded}
\begin{Highlighting}[]
\KeywordTok{if} (fs\OperatorTok{.}\CodeFuncTok{existsSync}(dataFile)) \{
  \KeywordTok{const} data \OperatorTok{=} fs\OperatorTok{.}\CodeFuncTok{readFileSync}(dataFile\OperatorTok{,} \StringTok{"utf8"})\OperatorTok{;}
  \KeywordTok{return} JSON\OperatorTok{.}\CodeFuncTok{parse}(data)\OperatorTok{;}
\} \KeywordTok{else} \{
  console\OperatorTok{.}\CodeFuncTok{log}(\StringTok{"No data file found. Starting with empty data."})\OperatorTok{;}
  \KeywordTok{return} []\OperatorTok{;}
\}
\end{Highlighting}
\end{Shaded}

Always check existence before reading. The alternative --- catching
\texttt{ENOENT} in a try-catch --- works but puts normal ``first run''
behavior in the error path.

\subsection{Lesson 8.3: CommonJS
Modules}\label{lesson-8.3-commonjs-modules}

\textbf{Module design for accessibility:}

A well-designed module has a clear exports interface. Teach students to
think of their module's exports as its ``public API'' --- what other
files can use:

\begin{Shaded}
\begin{Highlighting}[]
\CommentTok{// contacts.js {-}{-}{-} exports only what other files need}
\KeywordTok{const} path \OperatorTok{=} \CodeAccentTok{require}(\StringTok{"path"})\OperatorTok{;}
\KeywordTok{const} fs \OperatorTok{=} \CodeAccentTok{require}(\StringTok{"fs"})\OperatorTok{;}
 
\KeywordTok{const} DATA\_FILE \OperatorTok{=} path\OperatorTok{.}\CodeFuncTok{join}(\_\_dirname\OperatorTok{,} \StringTok{".."}\OperatorTok{,} \StringTok{"data"}\OperatorTok{,} \StringTok{"contacts.json"})\OperatorTok{;}
 
\KeywordTok{function} \CodeFuncTok{loadContacts}() \{
  \CommentTok{/* ... */}
\}
\KeywordTok{function} \CodeFuncTok{saveContacts}(contacts) \{
  \CommentTok{/* ... */}
\}
\KeywordTok{function} \CodeFuncTok{findContact}(contacts\OperatorTok{,} name) \{
  \CommentTok{/* ... */}
\}
 
\CommentTok{// Export only the public interface}
module\OperatorTok{.}\CodeTypeTok{exports} \OperatorTok{=} \{ loadContacts\OperatorTok{,} saveContacts\OperatorTok{,} findContact \}\OperatorTok{;}
\end{Highlighting}
\end{Shaded}

\emph{``Everything not in \texttt{module.exports} is private to this
file. Other files can use \texttt{loadContacts} and
\texttt{saveContacts}, but they cannot directly access
\texttt{DATA\_FILE}. This is encapsulation at the module level.''}

\paragraph{Thought Experiments}\label{thought-experiments-27}

\begin{enumerate}
\def\labelenumi{\arabic{enumi}.}
\item
  \emph{``Node.js's CommonJS (\texttt{require}) and the ES Module
  standard (\texttt{import}) are two different module systems in
  JavaScript. This creates confusion --- some packages use one, some use
  the other. TypeScript can use both. Why did two module systems emerge?
  What does this history say about how standards develop in open-source
  ecosystems?''}
\item
  \emph{``Synchronous file I/O (\texttt{readFileSync}) blocks execution
  until the file is read. Asynchronous file I/O (\texttt{readFile} or
  \texttt{fs.promises.readFile}) allows other code to run while waiting.
  For a startup script that reads configuration before doing anything
  else, which makes more sense? For a web server that handles thousands
  of requests, which matters more?''}
\item
  \emph{``Custom error classes allow code to throw and catch specific,
  named errors:
  \texttt{throw\ new\ ValidationError(\textquotesingle{}email\textquotesingle{},\ \textquotesingle{}Invalid\ format\textquotesingle{})}.
  What is the advantage of this over throwing
  \texttt{new\ Error(\textquotesingle{}Validation\ failed:\ email\ field\ invalid\ format\textquotesingle{})}?
  Think about code that catches and handles the error.''}
\end{enumerate}

\begin{center}\rule{0.5\linewidth}{0.5pt}\end{center}

\subsection{13. Unit 9: Collections, Analysis, and
Testing}\label{unit-9-collections-analysis-and-testing}

\subsection{Jest Testing}\label{jest-testing}

\paragraph{Instructional Cues}\label{instructional-cues-20}

\textbf{Jest output is accessible:}

Run \texttt{npx\ jest\ -\/-verbose} for maximum accessibility. The
\texttt{-\/-verbose} flag produces:

\begin{verbatim}
PASS src/math.test.js
  add function
    [x] adds two positive numbers (3ms)
    [x] adds negative numbers (1ms)
    [x] handles zero (0ms)
  getGrade function
    [x] returns A for score >= 90 (0ms)
    [x] returns F for score < 60 (0ms)

Test Suites: 1 passed, 1 total
Tests:       5 passed, 5 total
\end{verbatim}

Each test name is announced. Failures include file path, test name,
expected vs.~received values, and a diff. All of this is plain text ---
fully screen-reader-accessible.

\textbf{Add to \texttt{package.json} scripts:}

\begin{Shaded}
\begin{Highlighting}[]
\StringTok{"scripts"}\ErrorTok{:} \CodeFuncTok{\{}
    \CodeTypeTok{"test"}\CodeFuncTok{:} \StringTok{"jest {-}{-}verbose"}\CodeFuncTok{,}
    \CodeTypeTok{"test:watch"}\CodeFuncTok{:} \StringTok{"jest {-}{-}verbose {-}{-}watch"}
\CodeFuncTok{\}}
\end{Highlighting}
\end{Shaded}

\texttt{npm\ test} runs all tests. \texttt{npm\ run\ test:watch} reruns
tests whenever a file changes --- useful during active development.

\textbf{Test naming convention:}

\begin{Shaded}
\begin{Highlighting}[]
\CodeFuncTok{describe}(\StringTok{"BankAccount"}\OperatorTok{,} () \KeywordTok{=\textgreater{}} \{
  \CodeFuncTok{describe}(\StringTok{"deposit"}\OperatorTok{,} () \KeywordTok{=\textgreater{}} \{
    \CodeFuncTok{test}(\StringTok{"increases balance by deposited amount"}\OperatorTok{,} () \KeywordTok{=\textgreater{}} \{
      \CommentTok{/* ... */}
    \})\OperatorTok{;}
    \CodeFuncTok{test}(\StringTok{"rejects negative amounts"}\OperatorTok{,} () \KeywordTok{=\textgreater{}} \{
      \CommentTok{/* ... */}
    \})\OperatorTok{;}
    \CodeFuncTok{test}(\StringTok{"rejects zero amount"}\OperatorTok{,} () \KeywordTok{=\textgreater{}} \{
      \CommentTok{/* ... */}
    \})\OperatorTok{;}
  )\OperatorTok{;}
  \CodeFuncTok{describe}(\StringTok{"withdraw"}\OperatorTok{,} () \KeywordTok{=\textgreater{}} \{
    \CodeFuncTok{test}(\StringTok{"decreases balance by withdrawn amount"}\OperatorTok{,} () \KeywordTok{=\textgreater{}} \{\}
      \CommentTok{/* ... */}
    \})\OperatorTok{;}
    \CodeFuncTok{test}(\StringTok{"throws InsufficientFundsError when balance is too low"}\OperatorTok{,} () \KeywordTok{=\textgreater{}} \{
      \CommentTok{/* ... */}
    \})\OperatorTok{;}
  )\OperatorTok{;}
)\}\OperatorTok{;}
\end{Highlighting}
\end{Shaded}

The nested \texttt{describe} / \texttt{test} structure, when announced
by Jest with \texttt{-\/-verbose}, reads as complete sentences:
``BankAccount \textgreater{} deposit \textgreater{} increases balance by
deposited amount.''

\paragraph{Thought Experiments}\label{thought-experiments-28}

\begin{enumerate}
\def\labelenumi{\arabic{enumi}.}
\item
  \emph{``Jest's \texttt{expect(value).toEqual(expected)} performs deep
  comparison --- it checks that two objects have the same structure and
  values, not that they are the same object in memory. Why is deep
  comparison usually what you want when testing? When would you
  specifically want reference equality (\texttt{toBe})?''}
\item
  \emph{``Snapshot testing in Jest captures the output of a component or
  function the first time, then checks future runs against that
  snapshot. What are the advantages? What are the risks? When might a
  snapshot become `stale' and mislead you into thinking code is correct
  when it isn't?''}
\item
  \emph{``Code coverage measures what percentage of your code is
  executed by tests. 100\% coverage is a goal but not a guarantee of
  correctness. Design a function with 100\% test coverage that still has
  a bug. (Hint: think about boundary conditions --- what if the test
  covers the case where x=5 but not x=0?)''}
\end{enumerate}

\begin{center}\rule{0.5\linewidth}{0.5pt}\end{center}

\subsection{14. Units 10--11: Projects and
Capstone}\label{units-1011-projects-and-capstone}

\subsection{JavaScript Project
Structure}\label{javascript-project-structure}

\begin{verbatim}
my-data-manager/
|--- src/
|   |--- models/
|   |   +--- contact.js
|   |--- storage/
|   |   +--- dataManager.js
|   |--- ui/
|   |   +--- consoleUI.js
|   +--- utils/
|       |--- input.js         (the readline helper)
|       +--- validation.js    (reusable validation functions)
|--- data/
|   +--- contacts.json
|--- tests/
|   |--- models/
|   |   +--- contact.test.js
|   +--- storage/
|       +--- dataManager.test.js
|--- .eslintrc.json
|--- package.json
+--- index.js
\end{verbatim}

\textbf{\texttt{package.json} scripts --- teach as accessibility tools:}

\begin{Shaded}
\begin{Highlighting}[]
\CodeFuncTok{\{}
  \CodeTypeTok{"scripts"}\CodeFuncTok{:} \CodeFuncTok{\{}
    \CodeTypeTok{"start"}\CodeFuncTok{:} \StringTok{"node index.js"}\CodeFuncTok{,}
    \CodeTypeTok{"test"}\CodeFuncTok{:} \StringTok{"jest {-}{-}verbose"}\CodeFuncTok{,}
    \CodeTypeTok{"test:watch"}\CodeFuncTok{:} \StringTok{"jest {-}{-}verbose {-}{-}watch"}\CodeFuncTok{,}
    \CodeTypeTok{"lint"}\CodeFuncTok{:} \StringTok{"eslint src/"}\CodeFuncTok{,}
    \CodeTypeTok{"lint:fix"}\CodeFuncTok{:} \StringTok{"eslint src/ {-}{-}fix"}
  \CodeFuncTok{\}}
\CodeFuncTok{\}}
\end{Highlighting}
\end{Shaded}

\texttt{npm\ start} --- consistent command regardless of entry point
filename. \texttt{npm\ test} --- runs all tests. \texttt{npm\ run\ lint}
--- checks for errors without running the program.

\textbf{\texttt{npm\ run\ lint} as a pre-run check:} Teach students to
run \texttt{npm\ run\ lint} before running the program. ESLint catches
undefined variables, unused variables, and \texttt{==} before the
program runs --- much faster than finding runtime errors.

\subsection{Accessible Output Design for
Capstone}\label{accessible-output-design-for-capstone}

Require students to design their program's output to be navigable and
self-describing:

\begin{Shaded}
\begin{Highlighting}[]
\CommentTok{// Navigation landmark prints}
\KeywordTok{function} \CodeFuncTok{printSeparator}(title) \{
  \KeywordTok{const} line \OperatorTok{=} \StringTok{"{-}"}\OperatorTok{.}\CodeFuncTok{repeat}(\CodeNumTok{40})\OperatorTok{;}
  console\OperatorTok{.}\CodeFuncTok{log}(line)\OperatorTok{;}
  \KeywordTok{if} (title) console\OperatorTok{.}\CodeFuncTok{log}(\StringTok{\textasciigrave{}  }\StringTok{\$\{}title\OperatorTok{.}\CodeFuncTok{toUpperCase}()\StringTok{\}}\StringTok{\textasciigrave{}})\OperatorTok{;}
  console\OperatorTok{.}\CodeFuncTok{log}(line)\OperatorTok{;}
\}
 
\CommentTok{// Numbered list with clear field labels}
\KeywordTok{function} \CodeFuncTok{printContactList}(contacts) \{
  \KeywordTok{if} (contacts\OperatorTok{.}\CodeTypeTok{length} \OperatorTok{===} \CodeNumTok{0}) \{\}
    console\OperatorTok{.}\CodeFuncTok{log}(\StringTok{"No contacts found."})\OperatorTok{;}
    \KeywordTok{return}\OperatorTok{;}
  \}
  \CodeFuncTok{printSeparator}(\StringTok{"contacts"})\OperatorTok{;}
  contacts\OperatorTok{.}\CodeFuncTok{forEach}((c\OperatorTok{,} i) \KeywordTok{=\textgreater{}} \{
    console\OperatorTok{.}\CodeFuncTok{log}(\StringTok{\textasciigrave{}}\StringTok{\$\{}i \OperatorTok{+} \CodeNumTok{1}\StringTok{\}}\StringTok{. }\StringTok{\$\{}c\OperatorTok{.}\CodeTypeTok{name}\StringTok{\}}\StringTok{\textasciigrave{}})\OperatorTok{;}
    console\OperatorTok{.}\CodeFuncTok{log}(\StringTok{\textasciigrave{}   Phone: }\StringTok{\$\{}c\OperatorTok{.}\CodeTypeTok{phone} \OperatorTok{||} \StringTok{"not provided"}\StringTok{\}}\StringTok{\textasciigrave{}})\OperatorTok{;}
    console\OperatorTok{.}\CodeFuncTok{log}(\StringTok{\textasciigrave{}   Email: }\StringTok{\$\{}c\OperatorTok{.}\CodeTypeTok{email} \OperatorTok{||} \StringTok{"not provided"}\StringTok{\}}\StringTok{\textasciigrave{}})\OperatorTok{;}
  )\OperatorTok{;}
  console\OperatorTok{.}\CodeFuncTok{log}(\StringTok{\textasciigrave{}Total: }\StringTok{\$\{}contacts\OperatorTok{.}\CodeTypeTok{length}\StringTok{\}}\StringTok{ contact(s)\textasciigrave{}})\OperatorTok{;}
\}
 
\CommentTok{// Confirmation messages {-}{-}{-} every action must produce one}
\KeywordTok{function} \CodeFuncTok{confirmAction}(action\OperatorTok{,} details) \{
  console\OperatorTok{.}\CodeFuncTok{log}(\StringTok{\textasciigrave{}[x] }\StringTok{\$\{}action\StringTok{\}}\StringTok{: }\StringTok{\$\{}details\StringTok{\}}\StringTok{\textasciigrave{}})\OperatorTok{;}
\}
\end{Highlighting}
\end{Shaded}

\paragraph{Thought Experiments}\label{thought-experiments-29}

\begin{enumerate}
\def\labelenumi{\arabic{enumi}.}
\item
  \emph{``Your capstone is a command-line program. Eventually, it could
  become a web application. What would need to change? What would
  transfer directly? Write a two-column list: `Can reuse as-is' and
  `Must rewrite.'\,''}
\item
  \emph{``The accessible design principle `no visual-only cues for
  critical information' means every status, error, or confirmation must
  be communicated through text, not color or icon alone. Where in your
  program have you relied on visual formatting? How would you replace
  it?''}
\item
  \emph{``Open-source your capstone project. Write a README that
  explains: what the program does, how to install it, how to run it, and
  specifically how to use it with a screen reader. What information
  would a sighted user find obvious but a blind user needs explicitly
  stated?''}
\end{enumerate}

\begin{center}\rule{0.5\linewidth}{0.5pt}\end{center}

\subsection{15. Bridging from Python/Java to
JavaScript}\label{bridging-from-pythonjava-to-javascript}

\subsection{From Python to
JavaScript}\label{from-python-to-javascript}

\begin{fullwidth}{\def\LTcaptype{none} % do not increment counter
\begin{longtable}[]{@{}
  >{\raggedright\arraybackslash}p{(\linewidth - 4\tabcolsep) * \real{0.2745}}
  >{\raggedright\arraybackslash}p{(\linewidth - 4\tabcolsep) * \real{0.4020}}
  >{\raggedright\arraybackslash}p{(\linewidth - 4\tabcolsep) * \real{0.3235}}@{}}
\toprule\noalign{}
\begin{minipage}[b]{\linewidth}\raggedright
Python
\end{minipage} & \begin{minipage}[b]{\linewidth}\raggedright
JavaScript
\end{minipage} & \begin{minipage}[b]{\linewidth}\raggedright
Key Difference
\end{minipage} \\
\midrule\noalign{}
\endhead
\bottomrule\noalign{}
\endlastfoot
\texttt{print("hello")} & \texttt{console.log("hello")} & Method
vs.~function \\
\texttt{input("prompt")} & \texttt{await\ askQuestion("prompt")} & Async
in JavaScript \\
\texttt{f"Hello\ \{name\}"} &
\texttt{\textasciigrave{}Hello\ \$\{name\}\textasciigrave{}} & f-string
vs.~template literal \\
\texttt{def\ func(x):} &
\texttt{const\ func\ =\ x\ =\textgreater{}\ ...} & Arrow function
style \\
\texttt{True}/\texttt{False} & \texttt{true}/\texttt{false} & Lowercase
in JavaScript \\
\texttt{None} & \texttt{null} and \texttt{undefined} & JS has two
``nothing'' values \\
\texttt{my\_list} & \texttt{myList} & JavaScript uses camelCase \\
\texttt{len(x)} & \texttt{x.length} & Property, not function \\
\texttt{x\ is\ None} & \texttt{x\ ===\ null} or
\texttt{x\ ===\ undefined} & Different null checks \\
\texttt{import\ json} &
\texttt{const\ fs\ =\ require(\textquotesingle{}fs\textquotesingle{})} &
Built-in \texttt{JSON} object, no import \\
\texttt{if\ \_\_name\_\_\ ==\ \textquotesingle{}\_\_main\_\_\textquotesingle{}:}
& No direct equivalent (just call \texttt{main()}) & \\
\end{longtable}
}
\end{fullwidth}

\textbf{Thought experiments for Python-\textgreater JS students:}

\begin{enumerate}
\def\labelenumi{\arabic{enumi}.}
\item
  \emph{``In Python, you use spaces to define blocks (indentation). In
  JavaScript, you use \texttt{\{\}}. Python's blocks are invisible ---
  they're just the alignment of characters. JavaScript's blocks are
  explicitly marked. Which feels more natural? Which provides clearer
  information when read aloud?''}
\item
  \emph{``Python's \texttt{None} and JavaScript's
  \texttt{null}/\texttt{undefined} both represent `nothing,' but they
  behave differently. Python has one `nothing'; JavaScript has two. When
  is \texttt{null} appropriate? When is \texttt{undefined} appropriate?
  Under what circumstances would you encounter one vs.~the other?''}
\item
  \emph{``Python's type system and JavaScript's both use dynamic typing,
  but JavaScript's is more permissive --- it coerces types rather than
  raising errors. After using both, which approach do you find more
  helpful for learning? Which is more appropriate for production
  software?''}
\end{enumerate}

\begin{center}\rule{0.5\linewidth}{0.5pt}\end{center}

\subsection{16. Assessment Strategies}\label{assessment-strategies-1}

\subsection{Alternative Assessments for
JavaScript}\label{alternative-assessments-for-javascript}

\textbf{Console output analysis:} Provide a JavaScript program and 3--4
sample inputs. Ask the student to predict what the program prints for
each input before running it. This tests reading comprehension without
requiring writing.

\textbf{Linting-first review:} Submit code that is free of all ESLint
warnings. A student who achieves this has internalized JavaScript's best
practices at the syntax level.

\textbf{Refactoring exercises:} Give a working but poorly structured
program. Ask the student to refactor it to use functions, modules, or
classes appropriately. Assess the quality of the architecture, not just
whether it runs.

\textbf{Test-first implementation:} Give the student a set of failing
tests. Ask them to write the implementation that makes the tests pass.
This tests code comprehension and problem-solving without a blank-page
start.

\begin{center}\rule{0.5\linewidth}{0.5pt}\end{center}

\subsection{17. Supporting Diverse
Needs}\label{supporting-diverse-needs}

\subsection{For Students Who Are New to Async
Programming}\label{for-students-who-are-new-to-async-programming}

Async/await is conceptually hard. For students who struggle:

\begin{enumerate}
\def\labelenumi{\arabic{enumi}.}
\tightlist
\item
  Provide the complete \texttt{input.js} helper and instruct them to use
  it without modifying it.
\item
  For the first several lessons, treat
  \texttt{async\ function\ main()\ \{\ ...\ \}\ main();} as boilerplate,
  not content.
\item
  Introduce the conceptual explanation of async/await only after
  students are comfortable with the rest of the language.
\item
  Consider deferring async until after basic functions are mastered.
\end{enumerate}

\subsection{For Students with Additional Cognitive
Considerations}\label{for-students-with-additional-cognitive-considerations}

\begin{itemize}
\tightlist
\item
  ESLint provides immediate text feedback for common errors --- use it
  as a learning tool. Hearing ``no-unused-vars'' and then explaining
  what that means is a teaching moment triggered by the student's own
  code.
\item
  JavaScript's many syntactic variations can be overwhelming. Enforce a
  single consistent style (arrow functions, \texttt{const}/\texttt{let},
  \texttt{===}) and defer alternatives.
\item
  Use the REPL liberally for short experiments so the student gets
  immediate feedback without file management overhead.
\end{itemize}

\begin{center}\rule{0.5\linewidth}{0.5pt}\end{center}

\subsection{18. Alternative and Supplementary
Resources}\label{alternative-and-supplementary-resources-1}

\subsection{Textbooks and Reading}\label{textbooks-and-reading-1}

\begin{fullwidth}{\def\LTcaptype{none} % do not increment counter
\begin{longtable}[]{@{}
  >{\raggedright\arraybackslash}p{(\linewidth - 6\tabcolsep) * \real{0.3529}}
  >{\raggedright\arraybackslash}p{(\linewidth - 6\tabcolsep) * \real{0.1618}}
  >{\raggedright\arraybackslash}p{(\linewidth - 6\tabcolsep) * \real{0.0956}}
  >{\raggedright\arraybackslash}p{(\linewidth - 6\tabcolsep) * \real{0.3897}}@{}}
\toprule\noalign{}
\begin{minipage}[b]{\linewidth}\raggedright
Resource
\end{minipage} & \begin{minipage}[b]{\linewidth}\raggedright
Level
\end{minipage} & \begin{minipage}[b]{\linewidth}\raggedright
Format
\end{minipage} & \begin{minipage}[b]{\linewidth}\raggedright
Accessibility
\end{minipage} \\
\midrule\noalign{}
\endhead
\bottomrule\noalign{}
\endlastfoot
\textbf{Eloquent JavaScript} (Haverbeke, free online) &
Beginner--Intermediate & Free HTML & Excellent screen-reader-friendly
HTML \\
\textbf{You Don't Know JS} (Simpson, free on GitHub) & Intermediate &
Free HTML/PDF & In-depth on closures, scope, types \\
\textbf{JavaScript: The Good Parts} (Crockford) & Intermediate & Print &
Short; focuses on the reliable subset of JS \\
\textbf{JavaScript: The Definitive Guide} (Flanagan) & Reference &
Print/ebook & Comprehensive; good for looking up specifics \\
\textbf{MDN Web Docs} (developer.mozilla.org) & Reference & Free HTML &
Excellent HTML accessibility; definitive JS reference \\
\textbf{Node.js Documentation} (nodejs.org/docs) & Reference & Free HTML
& Official Node.js API reference \\
\textbf{JavaScript.info} & Beginner--Advanced & Free HTML & Very
well-organized; plain HTML; comprehensive \\
\end{longtable}
}
\end{fullwidth}

\textbf{javascript.info is especially recommended} --- it's a
comprehensive free tutorial organized from basics to advanced topics,
all in clean, accessible HTML. It's often more readable than MDN for
learning (vs.~reference), and screen readers navigate it very well.

\subsection{Online Courses and
Platforms}\label{online-courses-and-platforms-1}

\begin{fullwidth}{\def\LTcaptype{none} % do not increment counter
\begin{longtable}[]{@{}
  >{\raggedright\arraybackslash}p{(\linewidth - 6\tabcolsep) * \real{0.1984}}
  >{\raggedright\arraybackslash}p{(\linewidth - 6\tabcolsep) * \real{0.3095}}
  >{\raggedright\arraybackslash}p{(\linewidth - 6\tabcolsep) * \real{0.1032}}
  >{\raggedright\arraybackslash}p{(\linewidth - 6\tabcolsep) * \real{0.3889}}@{}}
\toprule\noalign{}
\begin{minipage}[b]{\linewidth}\raggedright
Platform
\end{minipage} & \begin{minipage}[b]{\linewidth}\raggedright
Course
\end{minipage} & \begin{minipage}[b]{\linewidth}\raggedright
Accessibility
\end{minipage} & \begin{minipage}[b]{\linewidth}\raggedright
Notes
\end{minipage} \\
\midrule\noalign{}
\endhead
\bottomrule\noalign{}
\endlastfoot
\textbf{freeCodeCamp} & JavaScript Algorithms \& Data Structures &
Medium & Certification path; some UI inaccessible \\
\textbf{Codecademy} & Learn JavaScript & Partial & Editor accessible;
some UI elements not \\
\textbf{Exercism.io} & JavaScript Track & High & Terminal-based; mentor
feedback \\
\textbf{JavaScript30} (Wesbos) & 30 vanilla JS projects & Medium &
Video-based; transcripts available \\
\textbf{edX} & Various JavaScript courses & High & Video transcripts;
accessible exercises \\
\textbf{Scrimba} & Learn JavaScript & Low & Interactive video; poor
screen reader support \\
\textbf{Khan Academy} & JavaScript course & Medium & Browser-based
editor; variable accessibility \\
\textbf{NodeSchool.io} & Node.js workshops & High & Command-line
workshopper format; fully accessible \\
\end{longtable}
}
\end{fullwidth}

\textbf{NodeSchool.io is especially recommended} --- workshops are run
entirely in the terminal via a command-line tool. Students work through
exercises, get feedback from the tool, and progress at their own pace.
The terminal-first design makes it highly accessible.

\subsection{Video and Audio
Resources}\label{video-and-audio-resources}

\begin{fullwidth}{\def\LTcaptype{none} % do not increment counter
\begin{longtable}[]{@{}
  >{\raggedright\arraybackslash}p{(\linewidth - 4\tabcolsep) * \real{0.4240}}
  >{\raggedright\arraybackslash}p{(\linewidth - 4\tabcolsep) * \real{0.3360}}
  >{\raggedright\arraybackslash}p{(\linewidth - 4\tabcolsep) * \real{0.2400}}@{}}
\toprule\noalign{}
\begin{minipage}[b]{\linewidth}\raggedright
Resource
\end{minipage} & \begin{minipage}[b]{\linewidth}\raggedright
Description
\end{minipage} & \begin{minipage}[b]{\linewidth}\raggedright
Accessibility
\end{minipage} \\
\midrule\noalign{}
\endhead
\bottomrule\noalign{}
\endlastfoot
\textbf{Traversy Media} (YouTube) & JavaScript crash courses and
projects & Captions available \\
\textbf{The Net Ninja} (YouTube) & In-depth JavaScript series & Captions
available \\
\textbf{Academind} (YouTube) & Deep dives into JS/Node & Captions
available \\
\textbf{Kevin Powell} (YouTube) & JavaScript and web fundamentals &
Captions; clear explanations \\
\textbf{CS50W} (Harvard) & Web programming with Python and JavaScript &
Full transcripts; high quality \\
\textbf{JavaScript: Understanding the Weird Parts} (Udemy) & Deep JS
internals & Captions \\
\end{longtable}
}
\end{fullwidth}

\subsection{JavaScript-Specific Accessibility
Tools}\label{javascript-specific-accessibility-tools}

\begin{fullwidth}{\def\LTcaptype{none} % do not increment counter
\begin{longtable}[]{@{}
  >{\raggedright\arraybackslash}p{(\linewidth - 4\tabcolsep) * \real{0.1778}}
  >{\raggedright\arraybackslash}p{(\linewidth - 4\tabcolsep) * \real{0.3667}}
  >{\raggedright\arraybackslash}p{(\linewidth - 4\tabcolsep) * \real{0.4556}}@{}}
\toprule\noalign{}
\begin{minipage}[b]{\linewidth}\raggedright
Tool
\end{minipage} & \begin{minipage}[b]{\linewidth}\raggedright
Purpose
\end{minipage} & \begin{minipage}[b]{\linewidth}\raggedright
Notes
\end{minipage} \\
\midrule\noalign{}
\endhead
\bottomrule\noalign{}
\endlastfoot
\textbf{ESLint} & Code quality and error prevention & Command-line;
fully accessible \\
\textbf{Prettier} & Consistent code formatting & Run on save; no screen
interaction needed \\
\textbf{Jest} & Testing framework & Accessible output;
\texttt{-\/-verbose} flag \\
\textbf{jshint.com} & Online JS linter & Plain HTML; basic linting \\
\textbf{node-inspect} & Node.js debugger & Terminal-based debugging \\
\textbf{0x} & Node.js profiler & Terminal output \\
\textbf{nvm} & Node version manager & Command-line; accessible \\
\end{longtable}
}
\end{fullwidth}

\subsection{Communities}\label{communities}

\begin{fullwidth}{\def\LTcaptype{none} % do not increment counter
\begin{longtable}[]{@{}
  >{\raggedright\arraybackslash}p{(\linewidth - 2\tabcolsep) * \real{0.4545}}
  >{\raggedright\arraybackslash}p{(\linewidth - 2\tabcolsep) * \real{0.5455}}@{}}
\toprule\noalign{}
\begin{minipage}[b]{\linewidth}\raggedright
Community
\end{minipage} & \begin{minipage}[b]{\linewidth}\raggedright
Relevance
\end{minipage} \\
\midrule\noalign{}
\endhead
\bottomrule\noalign{}
\endlastfoot
\textbf{Node.js Community} (nodejs.org/community) & Official Node.js
community resources \\
\textbf{JavaScript Weekly} (newsletter) & Weekly digest of JS news;
plain HTML email \\
\textbf{Stack Overflow JavaScript} & Q\&A community; plain HTML;
excellent accessibility \\
\textbf{Dev.to} & Developer blog platform; reasonable accessibility \\
\textbf{GitHub Discussions} on popular JS projects & Community support;
text-based \\
\textbf{NFB Jernigan Institute} & Resources for blind STEM students \\
\textbf{SIGACCESS} & Research on accessible computing including
programming \\
\end{longtable}
}
\end{fullwidth}

\begin{center}\rule{0.5\linewidth}{0.5pt}\end{center}

\subsection{19. Quick Reference Sheets}\label{quick-reference-sheets-1}

\subsection{JavaScript vs Python vs Java at a
Glance}\label{javascript-vs-python-vs-java-at-a-glance}

\begin{fullwidth}{\def\LTcaptype{none} % do not increment counter
\begin{longtable}[]{@{}
  >{\raggedright\arraybackslash}p{(\linewidth - 6\tabcolsep) * \real{0.1628}}
  >{\raggedright\arraybackslash}p{(\linewidth - 6\tabcolsep) * \real{0.1860}}
  >{\raggedright\arraybackslash}p{(\linewidth - 6\tabcolsep) * \real{0.3488}}
  >{\raggedright\arraybackslash}p{(\linewidth - 6\tabcolsep) * \real{0.3023}}@{}}
\toprule\noalign{}
\begin{minipage}[b]{\linewidth}\raggedright
Concept
\end{minipage} & \begin{minipage}[b]{\linewidth}\raggedright
Python
\end{minipage} & \begin{minipage}[b]{\linewidth}\raggedright
Java
\end{minipage} & \begin{minipage}[b]{\linewidth}\raggedright
JavaScript
\end{minipage} \\
\midrule\noalign{}
\endhead
\bottomrule\noalign{}
\endlastfoot
Print & \texttt{print("x")} & \texttt{System.out.println("x")} &
\texttt{console.log("x")} \\
Variable & \texttt{x\ =\ 5} & \texttt{int\ x\ =\ 5;} &
\texttt{const\ x\ =\ 5;} \\
String format & \texttt{f"Hi\ \{name\}"} &
\texttt{String.format("Hi\ \%s",\ name)} &
\texttt{\textasciigrave{}Hi\ \$\{name\}\textasciigrave{}} \\
User input & \texttt{input("?\ ")} & \texttt{scanner.nextLine()} &
\texttt{await\ askQuestion("?\ ")} \\
Array/List & \texttt{{[}1,\ 2,\ 3{]}} &
\texttt{int{[}{]}\ arr\ =\ \{1,2,3\}} & \texttt{{[}1,\ 2,\ 3{]}} \\
Dictionary/Map & \texttt{\{"a":\ 1\}} &
\texttt{new\ HashMap\textless{}\textgreater{}()} &
\texttt{\{\ a:\ 1\ \}} \\
Function & \texttt{def\ f(x):} & \texttt{public\ int\ f(int\ x)\ \{} &
\texttt{const\ f\ =\ x\ =\textgreater{}\ ...} \\
Class & \texttt{class\ Dog:} & \texttt{public\ class\ Dog\ \{} &
\texttt{class\ Dog\ \{} \\
Exception & \texttt{try/except} & \texttt{try/catch} &
\texttt{try/catch} \\
Null & \texttt{None} & \texttt{null} &
\texttt{null}/\texttt{undefined} \\
Boolean & \texttt{True}/\texttt{False} & \texttt{true}/\texttt{false} &
\texttt{true}/\texttt{false} \\
Length & \texttt{len(x)} & \texttt{x.length} & \texttt{x.length} \\
For each & \texttt{for\ x\ in\ list:} &
\texttt{for\ (T\ x\ :\ list)\ \{} &
\texttt{for\ (const\ x\ of\ list)\ \{} \\
Import & \texttt{import\ math} & \texttt{import\ java.util.*;} &
\texttt{const\ x\ =\ require(\textquotesingle{}./x\textquotesingle{})} \\
\end{longtable}
}
\end{fullwidth}

\subsection{Accessible Output
Patterns}\label{accessible-output-patterns-1}

\begin{Shaded}
\begin{Highlighting}[]
\CommentTok{// Standard separator}
console\OperatorTok{.}\CodeFuncTok{log}(\StringTok{"{-}"}\OperatorTok{.}\CodeFuncTok{repeat}(\CodeNumTok{40}))\OperatorTok{;}
 
\CommentTok{// Section header}
console\OperatorTok{.}\CodeFuncTok{log}(\StringTok{\textasciigrave{}}\StringTok{\textbackslash{}n}\StringTok{=== }\StringTok{\$\{}title\OperatorTok{.}\CodeFuncTok{toUpperCase}()\StringTok{\}}\StringTok{ ===\textasciigrave{}})\OperatorTok{;}
 
\CommentTok{// Numbered list with details}
items\OperatorTok{.}\CodeFuncTok{forEach}((item\OperatorTok{,} i) \KeywordTok{=\textgreater{}} \{
  console\OperatorTok{.}\CodeFuncTok{log}(\StringTok{\textasciigrave{}}\StringTok{\$\{}i \OperatorTok{+} \CodeNumTok{1}\StringTok{\}}\StringTok{. }\StringTok{\$\{}item\OperatorTok{.}\CodeTypeTok{name}\StringTok{\}}\StringTok{\textasciigrave{}})\OperatorTok{;}
  console\OperatorTok{.}\CodeFuncTok{log}(\StringTok{\textasciigrave{}   }\StringTok{\$\{}item\OperatorTok{.}\CodeTypeTok{detail}\StringTok{\}}\StringTok{\textasciigrave{}})\OperatorTok{;}
)\}\OperatorTok{;}
console\OperatorTok{.}\CodeFuncTok{log}(\StringTok{\textasciigrave{}Total: }\StringTok{\$\{}items\OperatorTok{.}\CodeTypeTok{length}\StringTok{\}}\StringTok{\textasciigrave{}})\OperatorTok{;}
 
\CommentTok{// Menu}
\KeywordTok{const} MENU \OperatorTok{=} \StringTok{\textasciigrave{}}
\StringTok{{-}{-}{-} Main Menu {-}{-}{-}}
\StringTok{1. Add item}
\StringTok{2. View all items}
\StringTok{3. Search}
\StringTok{4. Edit item}
\StringTok{5. Delete item}
\StringTok{6. Save}
\StringTok{7. Quit}
\StringTok{{-}{-}{-}{-}{-}{-}{-}{-}{-}{-}{-}{-}{-}{-}{-}{-}{-}\textasciigrave{}}\OperatorTok{;}
 
\CommentTok{// Confirmation}
console\OperatorTok{.}\CodeFuncTok{log}(\StringTok{\textasciigrave{}[x] Added: "}\StringTok{\$\{}name\StringTok{\}}\StringTok{"\textasciigrave{}})\OperatorTok{;}
console\OperatorTok{.}\CodeFuncTok{log}(\StringTok{\textasciigrave{}[x] Deleted: "}\StringTok{\$\{}name\StringTok{\}}\StringTok{"\textasciigrave{}})\OperatorTok{;}
console\OperatorTok{.}\CodeFuncTok{log}(\StringTok{\textasciigrave{}[x] Saved to }\StringTok{\$\{}filename\StringTok{\}}\StringTok{\textasciigrave{}})\OperatorTok{;}
 
\CommentTok{// Error messages}
console\OperatorTok{.}\CodeFuncTok{log}(\StringTok{\textasciigrave{}Error: Name cannot be empty. Please try again.\textasciigrave{}})\OperatorTok{;}
console\OperatorTok{.}\CodeFuncTok{log}(\StringTok{\textasciigrave{}Error: "}\StringTok{\$\{}filename\StringTok{\}}\StringTok{" not found. Starting with empty data.\textasciigrave{}})\OperatorTok{;}
console\OperatorTok{.}\CodeFuncTok{log}(\StringTok{\textasciigrave{}Error: Task number must be between 1 and }\StringTok{\$\{}count\StringTok{\}}\StringTok{.\textasciigrave{}})\OperatorTok{;}
\end{Highlighting}
\end{Shaded}

\subsection{Common JavaScript Error Quick
Reference}\label{common-javascript-error-quick-reference}

\begin{fullwidth}{\def\LTcaptype{none} % do not increment counter
\begin{longtable}[]{@{}
  >{\raggedright\arraybackslash}p{(\linewidth - 4\tabcolsep) * \real{0.3680}}
  >{\raggedright\arraybackslash}p{(\linewidth - 4\tabcolsep) * \real{0.2640}}
  >{\raggedright\arraybackslash}p{(\linewidth - 4\tabcolsep) * \real{0.3680}}@{}}
\toprule\noalign{}
\begin{minipage}[b]{\linewidth}\raggedright
Error
\end{minipage} & \begin{minipage}[b]{\linewidth}\raggedright
Common Cause
\end{minipage} & \begin{minipage}[b]{\linewidth}\raggedright
First Check
\end{minipage} \\
\midrule\noalign{}
\endhead
\bottomrule\noalign{}
\endlastfoot
\texttt{ReferenceError:\ x\ is\ not\ defined} & Typo or not declared &
Spelling; scope \\
\texttt{TypeError:\ Cannot\ read\ property\ of\ undefined} & Accessing
property on \texttt{undefined} & Which variable is undefined; optional
chaining \\
\texttt{TypeError:\ x\ is\ not\ a\ function} & Calling non-function &
\texttt{typeof\ x} before calling \\
\texttt{SyntaxError:\ Unexpected\ token} & Syntax mistake & Line number;
check brackets, commas \\
\texttt{SyntaxError:\ Cannot\ use\ import\ statement} & ES module syntax
in CommonJS & Use \texttt{require()} or add
\texttt{"type":\ "module"} \\
\texttt{Error:\ ENOENT:\ no\ such\ file\ or\ directory} & File not found
& \texttt{console.log(\_\_dirname)} and check path \\
\texttt{Error:\ EACCES:\ permission\ denied} & No write permission &
Check file permissions \\
\texttt{RangeError:\ Maximum\ call\ stack\ size\ exceeded} & Infinite
recursion & Check recursive function's base case \\
\texttt{UnhandledPromiseRejectionWarning} & Missing \texttt{await} or
\texttt{.catch()} & Add \texttt{try/catch} around \texttt{await}
calls \\
\texttt{NaN} (not an error, but wrong value) & Arithmetic on non-number
& \texttt{console.log(typeof\ variable)} before math \\
\end{longtable}
}
\end{fullwidth}

\subsection{Input Helper Module (Copy into every
project)}\label{input-helper-module-copy-into-every-project}

\begin{Shaded}
\begin{Highlighting}[]
\CommentTok{// utils/input.js}
\KeywordTok{const} readline \OperatorTok{=} \CodeAccentTok{require}(\StringTok{"readline"})\OperatorTok{;}
 
\KeywordTok{const} rl \OperatorTok{=} readline\OperatorTok{.}\CodeFuncTok{createInterface}(\{\}
  \CodeTypeTok{input}\OperatorTok{:} process\OperatorTok{.}\CodeTypeTok{stdin}\OperatorTok{,}
  \CodeTypeTok{output}\OperatorTok{:} process\OperatorTok{.}\CodeTypeTok{stdout}\OperatorTok{,}
\})\OperatorTok{;}
 
\CommentTok{/**}
\CommentTok{ * Prompts the user and returns their response as a string.}
\CommentTok{ * }\CommentTok{@param}\CommentTok{ }\CommentTok{\{string\}}\CommentTok{ prompt {-} The question to display}
\CommentTok{ * }\CommentTok{@returns}\CommentTok{ \{Promise}\KeywordTok{\textless{}string\textgreater{}}\CommentTok{\} The user\textquotesingle{}s trimmed response}
\CommentTok{ */}
\KeywordTok{function} \CodeFuncTok{askQuestion}(prompt) \{
  \KeywordTok{return} \KeywordTok{new} Promise((resolve) \KeywordTok{=\textgreater{}} \{
    rl\OperatorTok{.}\CodeFuncTok{question}(prompt\OperatorTok{,} (answer) \KeywordTok{=\textgreater{}} \CodeFuncTok{resolve}(answer))\OperatorTok{;}
  )\OperatorTok{;}
\}
 
\CommentTok{/**}
\CommentTok{ * Prompts until a valid number is entered.}
\CommentTok{ * }\CommentTok{@param}\CommentTok{ }\CommentTok{\{string\}}\CommentTok{ prompt {-} The question to display}
\CommentTok{ * }\CommentTok{@returns}\CommentTok{ \{Promise}\KeywordTok{\textless{}number\textgreater{}}\CommentTok{\} A valid number}
\CommentTok{ */}
\KeywordTok{async} \KeywordTok{function} \CodeFuncTok{askNumber}(prompt) \{
  \KeywordTok{while} (\KeywordTok{true}) \{\}
    \KeywordTok{const} input \OperatorTok{=} \KeywordTok{await} \CodeFuncTok{askQuestion}(prompt)\OperatorTok{;}
    \KeywordTok{const} num \OperatorTok{=} Number(input\OperatorTok{.}\CodeFuncTok{trim}())\OperatorTok{;}
    \KeywordTok{if} (\OperatorTok{!}\CodeAccentTok{isNaN}(num)) \KeywordTok{return} num\OperatorTok{;}
    console\OperatorTok{.}\CodeFuncTok{log}(\StringTok{"Invalid input: please enter a number."})\OperatorTok{;}
  \}
\}
 
\CommentTok{/**}
\CommentTok{ * Closes the readline interface. Call at the end of every program.}
\CommentTok{ */}
\KeywordTok{function} \CodeFuncTok{close}() \{
  rl\OperatorTok{.}\CodeFuncTok{close}()\OperatorTok{;}
\}
 
module\OperatorTok{.}\CodeTypeTok{exports} \OperatorTok{=} \{ askQuestion\OperatorTok{,} askNumber\OperatorTok{,} close \}\OperatorTok{;}
\end{Highlighting}
\end{Shaded}

\newpage

\section{Unit 0: Setting Up Your JavaScript Development
Environment}\label{unit-0-setting-up-your-javascript-development-environment}

\begin{center}\rule{0.5\linewidth}{0.5pt}\end{center}

\subsection{Lesson 0.1: Installing
Node.js}\label{lesson-0.1-installing-node.js-1}

Node.js lets you run JavaScript on your computer (outside a browser)
using the terminal.

\subsection{Installing Node.js}\label{installing-node.js}

Download \textbf{Node.js 18 LTS} from \url{https://nodejs.org/}.

\paragraph{Windows}\label{windows-2}

\begin{enumerate}
\def\labelenumi{\arabic{enumi}.}
\tightlist
\item
  Download the \texttt{.msi} installer.
\item
  \textbf{CRUCIAL STEP:} Ensure ``Add to PATH'' is checked during
  installation.
\item
  Activate ``Install'' and ``Finish''.
\end{enumerate}

\paragraph{macOS}\label{macos-2}

\begin{enumerate}
\def\labelenumi{\arabic{enumi}.}
\tightlist
\item
  Download the \texttt{.pkg} installer and follow the prompts.
\item
  Or use Homebrew: \texttt{brew\ install\ node}
\end{enumerate}

\paragraph{Linux (Debian/Ubuntu)}\label{linux-debianubuntu-1}

\begin{Shaded}
\begin{Highlighting}[]
curl \CodeTypeTok{{-}fsSL} https://deb.nodesource.com/setup\_18.x \KeywordTok{|} \CodeFuncTok{sudo} \CodeTypeTok{{-}E} bash \CodeTypeTok{{-}}
\CodeFuncTok{sudo} apt install nodejs
\end{Highlighting}
\end{Shaded}

\subsection{Verify Installation}\label{verify-installation}

\begin{Shaded}
\begin{Highlighting}[]
node \CodeTypeTok{{-}{-}version}     \CommentTok{\# Should show v18.x.x or later}
npm \CodeTypeTok{{-}{-}version}      \CommentTok{\# Should show npm version}
\end{Highlighting}
\end{Shaded}

\subsection{Skills Checklist}\label{skills-checklist-52}

\begin{itemize}
\tightlist
\item[$\square$]
  Running \texttt{node\ -\/-version} shows the version number.
\item[$\square$]
  Running \texttt{npm\ -\/-version} shows the npm version.
\end{itemize}

\begin{center}\rule{0.5\linewidth}{0.5pt}\end{center}

\subsection{Lesson 0.2: Setting Up VS
Code}\label{lesson-0.2-setting-up-vs-code}

\begin{enumerate}
\def\labelenumi{\arabic{enumi}.}
\tightlist
\item
  Install VS Code from \url{https://code.visualstudio.com/}.
\item
  Press \texttt{Ctrl\ +\ Shift\ +\ X} to open Extensions. Install:

  \begin{itemize}
  \tightlist
  \item
    \textbf{JavaScript (ES6) code snippets} --- Handy code shortcuts.
  \item
    \textbf{ESLint} --- Catches common errors.
  \end{itemize}
\item
  Enable Screen Reader Mode: \texttt{Ctrl\ +\ Shift\ +\ P}
  -\textgreater{} ``Toggle Screen Reader Optimized''.
\end{enumerate}

\subsection{VS Code Shortcuts for
JavaScript}\label{vs-code-shortcuts-for-javascript}

\begin{fullwidth}{\def\LTcaptype{none} % do not increment counter
\begin{longtable}[]{@{}
  >{\raggedright\arraybackslash}p{(\linewidth - 2\tabcolsep) * \real{0.3966}}
  >{\raggedright\arraybackslash}p{(\linewidth - 2\tabcolsep) * \real{0.6034}}@{}}
\toprule\noalign{}
\begin{minipage}[b]{\linewidth}\raggedright
Action
\end{minipage} & \begin{minipage}[b]{\linewidth}\raggedright
Shortcut
\end{minipage} \\
\midrule\noalign{}
\endhead
\bottomrule\noalign{}
\endlastfoot
Run file (via terminal) &
\texttt{Ctrl\ +\ \textbackslash{}\textasciigrave{}\textasciigrave{}\ then}node
filename.js\texttt{\textbar{}\ \textbar{}\ Command\ Palette\ \ \ \ \ \ \ \ \ \textbar{}}Ctrl
+ Shift +
P\texttt{\textbar{}\ \textbar{}\ Problems\ View\ \ \ \ \ \ \ \ \ \ \ \textbar{}}Ctrl
+ Shift +
M\texttt{\textbar{}\ \textbar{}\ Next\ Error\ \ \ \ \ \ \ \ \ \ \ \ \ \ \textbar{}}F8\texttt{\textbar{}\ \textbar{}\ Format\ Code\ \ \ \ \ \ \ \ \ \ \ \ \ \textbar{}}Shift
+ Alt +
F\texttt{\textbar{}\ \textbar{}\ Go\ to\ Definition\ \ \ \ \ \ \ \ \textbar{}}F12` \\
\end{longtable}
}
\end{fullwidth}

\begin{center}\rule{0.5\linewidth}{0.5pt}\end{center}

\subsection{Lesson 0.3: Your First JavaScript
File}\label{lesson-0.3-your-first-javascript-file}

\subsection{Hello World}\label{hello-world}

Create \texttt{hello.js}:

\begin{Shaded}
\begin{Highlighting}[]
console\OperatorTok{.}\CodeFuncTok{log}(\StringTok{"Hello, World!"})\OperatorTok{;}
console\OperatorTok{.}\CodeFuncTok{log}(\StringTok{"Welcome to JavaScript!"})\OperatorTok{;}
\end{Highlighting}
\end{Shaded}

Run it:

\begin{Shaded}
\begin{Highlighting}[]
node hello.js
\end{Highlighting}
\end{Shaded}

\subsection{The Node.js REPL}\label{the-node.js-repl}

For quick experiments, open an interactive session:

\begin{Shaded}
\begin{Highlighting}[]
node
\end{Highlighting}
\end{Shaded}

Type JavaScript directly and press Enter to execute. Type \texttt{.exit}
to quit.

\subsection{Key Differences from
Python}\label{key-differences-from-python}

\begin{fullwidth}{\def\LTcaptype{none} % do not increment counter
\begin{longtable}[]{@{}
  >{\raggedright\arraybackslash}p{(\linewidth - 2\tabcolsep) * \real{0.3714}}
  >{\raggedright\arraybackslash}p{(\linewidth - 2\tabcolsep) * \real{0.6286}}@{}}
\toprule\noalign{}
\begin{minipage}[b]{\linewidth}\raggedright
Python
\end{minipage} & \begin{minipage}[b]{\linewidth}\raggedright
JavaScript
\end{minipage} \\
\midrule\noalign{}
\endhead
\bottomrule\noalign{}
\endlastfoot
\texttt{print()} & \texttt{console.log()} \\
\texttt{\#} for comments & \texttt{//} for single-line, \texttt{/*\ */}
for multi-line \\
No semicolons required & Semicolons optional but recommended \\
Indentation defines blocks & Curly braces \texttt{\{\}} define blocks \\
\texttt{def} for functions & \texttt{function} or
\texttt{=\textgreater{}} for functions \\
\texttt{True}/\texttt{False} & \texttt{true}/\texttt{false}
(lowercase) \\
\texttt{None} & \texttt{null} and \texttt{undefined} \\
\texttt{len(x)} & \texttt{x.length} \\
\end{longtable}
}
\end{fullwidth}

\subsection{Skills Checklist}\label{skills-checklist-53}

\begin{itemize}
\tightlist
\item[$\square$]
  I can create a \texttt{.js} file and run it with
  \texttt{node\ filename.js}.
\item[$\square$]
  I can use the Node.js REPL for quick experiments.
\end{itemize}

\begin{center}\rule{0.5\linewidth}{0.5pt}\end{center}

\section{Unit 1: Foundations --- Getting Started with
JavaScript}\label{unit-1-foundations-getting-started-with-javascript}

\subsection{Overview}\label{overview-12}

This unit covers JavaScript fundamentals: variables, data types, and
basic output.

\begin{center}\rule{0.5\linewidth}{0.5pt}\end{center}

\subsection{Lesson 1.1: Welcome to Computer
Science}\label{lesson-1.1-welcome-to-computer-science-2}

\emph{(See Python Unit 1.1 for conceptual content --- hardware,
software, algorithms, and problem-solving are language-agnostic.)}

\subsection{JavaScript's Role}\label{javascripts-role}


\marginnote{JavaScript was designed in ten days. Many of its quirks---\texttt{typeof null}, loose equality, hoisting---are direct consequences of that original sprint. Node.js, created by Ryan Dahl, moved JavaScript to the server in 2009 and changed its role from scripting language to general-purpose runtime.}

\begin{center}\rule{0.5\linewidth}{0.5pt}\end{center}

\subsection{Lesson 1.2: Variables and
Output}\label{lesson-1.2-variables-and-output}

\subsection{\texorpdfstring{\texttt{console.log()}}{console.log()}}\label{console.log}\amn{\textbf{Accessibility Tip:} Always include a label in \texttt{console.log()} calls: \texttt{console.log("Result:", value)} rather than \texttt{console.log(value)}. The label is spoken first, giving context before the value is announced. This prevents the common confusion of hearing a number with no idea what it represents.}
\subsection{\texorpdfstring{\texttt{console.log()}}{console.log()}}\label{console.log}

\begin{Shaded}
\begin{Highlighting}[]
console\OperatorTok{.}\CodeFuncTok{log}(\StringTok{"Hello, World!"})\OperatorTok{;}
console\OperatorTok{.}\CodeFuncTok{log}(\CodeNumTok{42})\OperatorTok{;}
console\OperatorTok{.}\CodeFuncTok{log}(\KeywordTok{true})\OperatorTok{;}
console\OperatorTok{.}\CodeFuncTok{log}(\StringTok{"Value:"}\OperatorTok{,} \CodeNumTok{42}\OperatorTok{,} \StringTok{"and"}\OperatorTok{,} \StringTok{"text"})\OperatorTok{;}
\end{Highlighting}
\end{Shaded}

\subsection{Declaring Variables}\label{declaring-variables}\amn{\textbf{Accessibility Tip:} Use \texttt{const} for every variable by default. If you later try to reassign a \texttt{const}, JavaScript throws a \texttt{TypeError} immediately, giving you an audible signal that you have made an error. Change to \texttt{let} only when you genuinely need to reassign.}
\subsection{Declaring Variables}\label{declaring-variables}

JavaScript has three ways to declare variables:

\begin{fullwidth}{\def\LTcaptype{none} % do not increment counter
\begin{longtable}[]{@{}
  >{\raggedright\arraybackslash}p{(\linewidth - 6\tabcolsep) * \real{0.1129}}
  >{\raggedright\arraybackslash}p{(\linewidth - 6\tabcolsep) * \real{0.1290}}
  >{\raggedright\arraybackslash}p{(\linewidth - 6\tabcolsep) * \real{0.1935}}
  >{\raggedright\arraybackslash}p{(\linewidth - 6\tabcolsep) * \real{0.5645}}@{}}
\toprule\noalign{}
\begin{minipage}[b]{\linewidth}\raggedright
Keyword
\end{minipage} & \begin{minipage}[b]{\linewidth}\raggedright
Scope
\end{minipage} & \begin{minipage}[b]{\linewidth}\raggedright
Reassignable
\end{minipage} & \begin{minipage}[b]{\linewidth}\raggedright
Use When
\end{minipage} \\
\midrule\noalign{}
\endhead
\bottomrule\noalign{}
\endlastfoot
\texttt{const} & Block & No & Value won't change (default choice) \\
\texttt{let} & Block & Yes & Value will change \\
\texttt{var} & Function & Yes & Legacy --- avoid in modern code \\
\end{longtable}
}
\end{fullwidth}

\begin{Shaded}
\begin{Highlighting}[]
\KeywordTok{const} \CodeNumTok{PI} \OperatorTok{=} \CodeNumTok{3.14159}\OperatorTok{;} \CommentTok{// Cannot be reassigned}
\KeywordTok{let} score \OperatorTok{=} \CodeNumTok{100}\OperatorTok{;} \CommentTok{// Can change}
score \OperatorTok{=} \CodeNumTok{150}\OperatorTok{;} \CommentTok{// OK}
 
\CommentTok{// var is function{-}scoped (legacy)}
\KeywordTok{var} oldWay \OperatorTok{=} \StringTok{"avoid this"}\OperatorTok{;}
\end{Highlighting}
\end{Shaded}


\marginnote{\texttt{var} is still valid JavaScript and appears in most legacy code.}\amn{\textbf{Accessibility Tip:} Use \texttt{const} for every variable declaration unless you know the value will change. If you try to reassign a \texttt{const}, JavaScript throws a \texttt{TypeError} immediately---an audible error that helps you catch accidental reassignments that would otherwise be silent bugs with \texttt{let} or \texttt{var}.}
\marginnote{\texttt{var} is still valid JavaScript and appears in most legacy code. Students will encounter it in tutorials and Stack Overflow answers. Teach \texttt{const}/\texttt{let} as the current standard but make sure students can \emph{read} \texttt{var} without confusion.}

\subsection{Template Literals}\label{template-literals}

\begin{Shaded}
\begin{Highlighting}[]
\KeywordTok{const} name \OperatorTok{=} \StringTok{"Alice"}\OperatorTok{;}
\KeywordTok{const} age \OperatorTok{=} \CodeNumTok{15}\OperatorTok{;}
 
\CommentTok{// Template literal (backticks {-}{-}{-} preferred)}
console\OperatorTok{.}\CodeFuncTok{log}(\StringTok{\textasciigrave{}Hello, }\StringTok{\$\{}name\StringTok{\}}\StringTok{! You are }\StringTok{\$\{}age\StringTok{\}}\StringTok{ years old.\textasciigrave{}})\OperatorTok{;}
 
\CommentTok{// Concatenation (older style)}
console\OperatorTok{.}\CodeFuncTok{log}(\StringTok{"Hello, "} \OperatorTok{+} name \OperatorTok{+} \StringTok{"! You are "} \OperatorTok{+} age \OperatorTok{+} \StringTok{" years old."})\OperatorTok{;}
\end{Highlighting}
\end{Shaded}

\subsection{Skills Checklist}\label{skills-checklist-54}

\begin{itemize}
\tightlist
\item[$\square$]
  I can use \texttt{console.log()} to display output.
\item[$\square$]
  I can declare variables using \texttt{const} and \texttt{let}.
\item[$\square$]
  I can use template literals with \texttt{\$\{\}}.
\end{itemize}

\begin{center}\rule{0.5\linewidth}{0.5pt}\end{center}

\subsection{Lesson 1.3: Data Types}\label{lesson-1.3-data-types}

\subsection{Primitive Types}\label{primitive-types}

\begin{Shaded}
\begin{Highlighting}[]
\CommentTok{// Number (integers and floats use the same type)}
\KeywordTok{let} integer \OperatorTok{=} \CodeNumTok{42}\OperatorTok{;}
\KeywordTok{let} decimal \OperatorTok{=} \CodeNumTok{3.14}\OperatorTok{;}
\KeywordTok{let} negative \OperatorTok{=} \OperatorTok{{-}}\CodeNumTok{10}\OperatorTok{;}
 
\CommentTok{// String}
\KeywordTok{let} name \OperatorTok{=} \StringTok{"Alice"}\OperatorTok{;}
\KeywordTok{let} greeting \OperatorTok{=} \StringTok{"Hello"}\OperatorTok{;}
\KeywordTok{let} multiLine \OperatorTok{=} \StringTok{\textasciigrave{}Line 1}
\StringTok{Line 2\textasciigrave{}}\OperatorTok{;}
 
\CommentTok{// Boolean}
\KeywordTok{let} isRaining \OperatorTok{=} \KeywordTok{false}\OperatorTok{;}
\KeywordTok{let} isLoggedIn \OperatorTok{=} \KeywordTok{true}\OperatorTok{;}
 
\CommentTok{// null (intentional absence of value)}
\KeywordTok{let} emptyValue \OperatorTok{=} \KeywordTok{null}\OperatorTok{;}
 
\CommentTok{// undefined (value not assigned)}
\KeywordTok{let} notAssigned\OperatorTok{;}
console\OperatorTok{.}\CodeFuncTok{log}(notAssigned)\OperatorTok{;} \CommentTok{// undefined}
 
\CommentTok{// Check type}
console\OperatorTok{.}\CodeFuncTok{log}(\KeywordTok{typeof} \CodeNumTok{42})\OperatorTok{;} \CommentTok{// "number"}
console\OperatorTok{.}\CodeFuncTok{log}(\KeywordTok{typeof} \StringTok{"hello"})\OperatorTok{;} \CommentTok{// "string"}
console\OperatorTok{.}\CodeFuncTok{log}(\KeywordTok{typeof} \KeywordTok{true})\OperatorTok{;} \CommentTok{// "boolean"}
console\OperatorTok{.}\CodeFuncTok{log}(\KeywordTok{typeof} \KeywordTok{null})\OperatorTok{;} \CommentTok{// "object" (known quirk!)}
console\OperatorTok{.}\CodeFuncTok{log}(\KeywordTok{typeof} \KeywordTok{undefined})\OperatorTok{;} \CommentTok{// "undefined"}
\end{Highlighting}
\end{Shaded}
\sidenote{\texttt{typeof~null~===~"object"} is a decades-old bug in JavaScript that was never fixed for backward compatibility. The correct test for null is \texttt{value~===~null} (strict equality), not \texttt{typeof}. This is one of several JavaScript ``gotchas'' documented in depth by Douglas Crockford's \emph{JavaScript: The Good Parts}.}

\subsection{Arithmetic Operators}\label{arithmetic-operators}

\begin{Shaded}
\begin{Highlighting}[]
console\OperatorTok{.}\CodeFuncTok{log}(\CodeNumTok{5} \OperatorTok{+} \CodeNumTok{3})\OperatorTok{;} \CommentTok{// 8}
console\OperatorTok{.}\CodeFuncTok{log}(\CodeNumTok{10} \OperatorTok{{-}} \CodeNumTok{4})\OperatorTok{;} \CommentTok{// 6}
console\OperatorTok{.}\CodeFuncTok{log}(\CodeNumTok{7} \OperatorTok{*} \CodeNumTok{6})\OperatorTok{;} \CommentTok{// 42}
console\OperatorTok{.}\CodeFuncTok{log}(\CodeNumTok{15} \OperatorTok{/} \CodeNumTok{4})\OperatorTok{;} \CommentTok{// 3.75  (always decimal, unlike Python 2 or Java)}
console\OperatorTok{.}\CodeFuncTok{log}(\CodeNumTok{10} \OperatorTok{\%} \CodeNumTok{3})\OperatorTok{;} \CommentTok{// 1  (remainder)}
console\OperatorTok{.}\CodeFuncTok{log}(\CodeNumTok{2} \OperatorTok{**} \CodeNumTok{10})\OperatorTok{;} \CommentTok{// 1024  (exponentiation)}
 
\CommentTok{// Math.floor for integer division equivalent}
console\OperatorTok{.}\CodeFuncTok{log}(Math\OperatorTok{.}\CodeFuncTok{floor}(\CodeNumTok{15} \OperatorTok{/} \CodeNumTok{4}))\OperatorTok{;} \CommentTok{// 3}
\end{Highlighting}
\end{Shaded}

\subsection{Type Conversion}\label{type-conversion}

\begin{Shaded}
\begin{Highlighting}[]
\CommentTok{// String to Number}
\KeywordTok{let} str \OperatorTok{=} \StringTok{"123"}\OperatorTok{;}
\KeywordTok{let} num \OperatorTok{=} Number(str)\OperatorTok{;} \CommentTok{// 123}
\KeywordTok{let} numParsed \OperatorTok{=} \CodeAccentTok{parseInt}(str)\OperatorTok{;} \CommentTok{// 123}
\KeywordTok{let} floatParsed \OperatorTok{=} \CodeAccentTok{parseFloat}(\StringTok{"3.14"})\OperatorTok{;} \CommentTok{// 3.14}
 
\CommentTok{// Number to String}
\KeywordTok{let} n \OperatorTok{=} \CodeNumTok{42}\OperatorTok{;}
\KeywordTok{let} s \OperatorTok{=} String(n)\OperatorTok{;} \CommentTok{// "42"}
\KeywordTok{let} s2 \OperatorTok{=} n\OperatorTok{.}\CodeFuncTok{toString}()\OperatorTok{;} \CommentTok{// "42"}
\KeywordTok{let} s3 \OperatorTok{=} \StringTok{\textasciigrave{}}\StringTok{\$\{}n\StringTok{\}}\StringTok{\textasciigrave{}}\OperatorTok{;} \CommentTok{// "42" (template literal)}
 
\CommentTok{// Checking}
console\OperatorTok{.}\CodeFuncTok{log}(\CodeAccentTok{isNaN}(\StringTok{"abc"}))\OperatorTok{;} \CommentTok{// true {-}{-}{-} "not a number"}
console\OperatorTok{.}\CodeFuncTok{log}(\CodeAccentTok{isNaN}(\CodeNumTok{42}))\OperatorTok{;} \CommentTok{// false}
console\OperatorTok{.}\CodeFuncTok{log}(Number(\StringTok{"abc"}))\OperatorTok{;} \CommentTok{// NaN}
\end{Highlighting}
\end{Shaded}

\subsection{Constants}\label{constants}

\begin{Shaded}
\begin{Highlighting}[]
\KeywordTok{const} MAX\_SIZE \OperatorTok{=} \CodeNumTok{100}\OperatorTok{;}
\KeywordTok{const} APP\_NAME \OperatorTok{=} \StringTok{"My App"}\OperatorTok{;}
\KeywordTok{const} \CodeNumTok{PI} \OperatorTok{=} Math\OperatorTok{.}\CodeNumTok{PI}\OperatorTok{;} \CommentTok{// Use Math.PI {-}{-}{-} already defined}
\end{Highlighting}
\end{Shaded}

\subsection{Examples}\label{examples-14}

\textbf{Example 1.3.1: Personal Information}

\begin{Shaded}
\begin{Highlighting}[]
\KeywordTok{const} myName \OperatorTok{=} \StringTok{"Jamie"}\OperatorTok{;}
\KeywordTok{const} favoriteNumber \OperatorTok{=} \CodeNumTok{7}\OperatorTok{;}
\KeywordTok{const} heightInMeters \OperatorTok{=} \CodeNumTok{1.75}\OperatorTok{;}
\KeywordTok{const} likesPizza \OperatorTok{=} \KeywordTok{true}\OperatorTok{;}
 
console\OperatorTok{.}\CodeFuncTok{log}(\StringTok{\textasciigrave{}Name: }\StringTok{\$\{}myName\StringTok{\}}\StringTok{\textasciigrave{}})\OperatorTok{;}
console\OperatorTok{.}\CodeFuncTok{log}(\StringTok{\textasciigrave{}Favorite number: }\StringTok{\$\{}favoriteNumber\StringTok{\}}\StringTok{\textasciigrave{}})\OperatorTok{;}
console\OperatorTok{.}\CodeFuncTok{log}(\StringTok{\textasciigrave{}Height: }\StringTok{\$\{}heightInMeters\StringTok{\}}\StringTok{ meters\textasciigrave{}})\OperatorTok{;}
console\OperatorTok{.}\CodeFuncTok{log}(\StringTok{\textasciigrave{}Likes pizza: }\StringTok{\$\{}likesPizza\StringTok{\}}\StringTok{\textasciigrave{}})\OperatorTok{;}
\end{Highlighting}
\end{Shaded}

\textbf{Example 1.3.2: Simple Math}

\begin{Shaded}
\begin{Highlighting}[]
\KeywordTok{const} num1 \OperatorTok{=} \CodeNumTok{15}\OperatorTok{;}
\KeywordTok{const} num2 \OperatorTok{=} \CodeNumTok{7}\OperatorTok{;}
 
console\OperatorTok{.}\CodeFuncTok{log}(\StringTok{\textasciigrave{}Sum: }\StringTok{\$\{}num1 \OperatorTok{+} num2\StringTok{\}}\StringTok{\textasciigrave{}})\OperatorTok{;}
console\OperatorTok{.}\CodeFuncTok{log}(\StringTok{\textasciigrave{}Difference: }\StringTok{\$\{}num1 \OperatorTok{{-}} num2\StringTok{\}}\StringTok{\textasciigrave{}})\OperatorTok{;}
console\OperatorTok{.}\CodeFuncTok{log}(\StringTok{\textasciigrave{}Product: }\StringTok{\$\{}num1 \OperatorTok{*} num2\StringTok{\}}\StringTok{\textasciigrave{}})\OperatorTok{;}
console\OperatorTok{.}\CodeFuncTok{log}(\StringTok{\textasciigrave{}Quotient: }\StringTok{\$\{}(num1 \OperatorTok{/} num2)\OperatorTok{.}\CodeFuncTok{toFixed}(\CodeNumTok{2})\StringTok{\}}\StringTok{\textasciigrave{}})\OperatorTok{;}
console\OperatorTok{.}\CodeFuncTok{log}(\StringTok{\textasciigrave{}Remainder: }\StringTok{\$\{}num1 \OperatorTok{\%} num2\StringTok{\}}\StringTok{\textasciigrave{}})\OperatorTok{;}
\end{Highlighting}
\end{Shaded}

\subsection{Activities}\label{activities-12}

\textbf{Activity 1.3.1: Area Calculator}

\begin{Shaded}
\begin{Highlighting}[]
\KeywordTok{const} length \OperatorTok{=} \CodeNumTok{12}\OperatorTok{;}
\KeywordTok{const} width \OperatorTok{=} \CodeNumTok{5}\OperatorTok{;}
\KeywordTok{const} height \OperatorTok{=} \CodeNumTok{8}\OperatorTok{;}
 
\KeywordTok{const} rectangleArea \OperatorTok{=} length \OperatorTok{*} width\OperatorTok{;}
\KeywordTok{const} boxVolume \OperatorTok{=} length \OperatorTok{*} width \OperatorTok{*} height\OperatorTok{;}
\KeywordTok{const} circleArea \OperatorTok{=} Math\OperatorTok{.}\CodeNumTok{PI} \OperatorTok{*} (length \OperatorTok{/} \CodeNumTok{2}) \OperatorTok{**} \CodeNumTok{2}\OperatorTok{;}
 
console\OperatorTok{.}\CodeFuncTok{log}(\StringTok{\textasciigrave{}Rectangle area: }\StringTok{\$\{}rectangleArea\StringTok{\}}\StringTok{\textasciigrave{}})\OperatorTok{;}
console\OperatorTok{.}\CodeFuncTok{log}(\StringTok{\textasciigrave{}Box volume: }\StringTok{\$\{}boxVolume\StringTok{\}}\StringTok{\textasciigrave{}})\OperatorTok{;}
console\OperatorTok{.}\CodeFuncTok{log}(\StringTok{\textasciigrave{}Circle area: }\StringTok{\$\{}circleArea\OperatorTok{.}\CodeFuncTok{toFixed}(\CodeNumTok{2})\StringTok{\}}\StringTok{\textasciigrave{}})\OperatorTok{;}
\end{Highlighting}
\end{Shaded}

\subsection{Skills Checklist}\label{skills-checklist-55}

\begin{itemize}
\tightlist
\item[$\square$]
  I can identify JavaScript's primitive data types.
\item[$\square$]
  I can perform arithmetic operations in JavaScript.
\item[$\square$]
  I understand that JavaScript's \texttt{/} always gives a decimal
  result.
\item[$\square$]
  I can convert between types using \texttt{Number()},
  \texttt{String()}, \texttt{parseInt()}.
\item[$\square$]
  I understand \texttt{const} vs.~\texttt{let}.
\end{itemize}

\newpage

\section{Unit 2: Numbers, Input, and the Math
Object}\label{unit-2-numbers-input-and-the-math-object}

\begin{center}\rule{0.5\linewidth}{0.5pt}\end{center}

\subsection{Lesson 2.1: The Math
Object}\label{lesson-2.1-the-math-object}

JavaScript's built-in \texttt{Math} object provides mathematical
functions and constants:

\begin{Shaded}
\begin{Highlighting}[]
\CommentTok{// Constants}
console\OperatorTok{.}\CodeFuncTok{log}(Math\OperatorTok{.}\CodeNumTok{PI})\OperatorTok{;} \CommentTok{// 3.14159265358979...}
console\OperatorTok{.}\CodeFuncTok{log}(Math\OperatorTok{.}\CodeNumTok{E})\OperatorTok{;} \CommentTok{// 2.71828...}
 
\CommentTok{// Functions}
console\OperatorTok{.}\CodeFuncTok{log}(Math\OperatorTok{.}\CodeFuncTok{sqrt}(\CodeNumTok{25}))\OperatorTok{;} \CommentTok{// 5}
console\OperatorTok{.}\CodeFuncTok{log}(Math\OperatorTok{.}\CodeFuncTok{pow}(\CodeNumTok{2}\OperatorTok{,} \CodeNumTok{10}))\OperatorTok{;} \CommentTok{// 1024}
console\OperatorTok{.}\CodeFuncTok{log}(Math\OperatorTok{.}\CodeFuncTok{abs}(\OperatorTok{{-}}\CodeNumTok{7.5}))\OperatorTok{;} \CommentTok{// 7.5}
console\OperatorTok{.}\CodeFuncTok{log}(Math\OperatorTok{.}\CodeFuncTok{ceil}(\CodeNumTok{4.2}))\OperatorTok{;} \CommentTok{// 5}
console\OperatorTok{.}\CodeFuncTok{log}(Math\OperatorTok{.}\CodeFuncTok{floor}(\CodeNumTok{4.8}))\OperatorTok{;} \CommentTok{// 4}
console\OperatorTok{.}\CodeFuncTok{log}(Math\OperatorTok{.}\CodeFuncTok{round}(\CodeNumTok{4.5}))\OperatorTok{;} \CommentTok{// 5}
console\OperatorTok{.}\CodeFuncTok{log}(Math\OperatorTok{.}\CodeFuncTok{round}(\CodeNumTok{4.4}))\OperatorTok{;} \CommentTok{// 4}
console\OperatorTok{.}\CodeFuncTok{log}(Math\OperatorTok{.}\CodeFuncTok{max}(\CodeNumTok{3}\OperatorTok{,} \CodeNumTok{7}\OperatorTok{,} \CodeNumTok{1}\OperatorTok{,} \CodeNumTok{9}\OperatorTok{,} \CodeNumTok{5}))\OperatorTok{;} \CommentTok{// 9}
console\OperatorTok{.}\CodeFuncTok{log}(Math\OperatorTok{.}\CodeFuncTok{min}(\CodeNumTok{3}\OperatorTok{,} \CodeNumTok{7}\OperatorTok{,} \CodeNumTok{1}\OperatorTok{,} \CodeNumTok{9}\OperatorTok{,} \CodeNumTok{5}))\OperatorTok{;} \CommentTok{// 1}
 
\CommentTok{// Random number between 0 (inclusive) and 1 (exclusive)}
console\OperatorTok{.}\CodeFuncTok{log}(Math\OperatorTok{.}\CodeFuncTok{random}())\OperatorTok{;}
 
\CommentTok{// Random integer between min and max (inclusive)}
\KeywordTok{function} \CodeFuncTok{randomInt}(min\OperatorTok{,} max) \{
  \KeywordTok{return} Math\OperatorTok{.}\CodeFuncTok{floor}(Math\OperatorTok{.}\CodeFuncTok{random}() \OperatorTok{*} (max \OperatorTok{{-}} min \OperatorTok{+} \CodeNumTok{1})) \OperatorTok{+} min\OperatorTok{;}
\}
console\OperatorTok{.}\CodeFuncTok{log}(\CodeFuncTok{randomInt}(\CodeNumTok{1}\OperatorTok{,} \CodeNumTok{6}))\OperatorTok{;} \CommentTok{// Simulates a dice roll}
\end{Highlighting}
\end{Shaded}

\begin{center}\rule{0.5\linewidth}{0.5pt}\end{center}

\subsection{Lesson 2.2: Getting User Input in
Node.js}\label{lesson-2.2-getting-user-input-in-node.js}

Getting user input in Node.js requires the built-in \texttt{readline}
module.

\subsection{\texorpdfstring{Using \texttt{readline} (Standard
Library)}{Using readline (Standard Library)}}\label{using-readline-standard-library}

\begin{Shaded}
\begin{Highlighting}[]
\KeywordTok{const} readline \OperatorTok{=} \CodeAccentTok{require}(\StringTok{"readline"})\OperatorTok{;}
 
\KeywordTok{const} rl \OperatorTok{=} readline\OperatorTok{.}\CodeFuncTok{createInterface}(\{
  \CodeTypeTok{input}\OperatorTok{:} process\OperatorTok{.}\CodeTypeTok{stdin}\OperatorTok{,}
  \CodeTypeTok{output}\OperatorTok{:} process\OperatorTok{.}\CodeTypeTok{stdout}\OperatorTok{,}
\})\OperatorTok{;}
 
rl\OperatorTok{.}\CodeFuncTok{question}(\StringTok{"What is your name? "}\OperatorTok{,} (name) \KeywordTok{=\textgreater{}} \{
  console\OperatorTok{.}\CodeFuncTok{log}(\StringTok{\textasciigrave{}Hello, }\StringTok{\$\{}name\StringTok{\}}\StringTok{!\textasciigrave{}})\OperatorTok{;}
  rl\OperatorTok{.}\CodeFuncTok{close}()\OperatorTok{;}
)\OperatorTok{;}
\end{Highlighting}
\end{Shaded}

\subsection{Promisified Input (Modern
Approach)}\label{promisified-input-modern-approach}

\begin{Shaded}
\begin{Highlighting}[]
\KeywordTok{const} readline \OperatorTok{=} \CodeAccentTok{require}(\StringTok{"readline"})\OperatorTok{;}
 
\KeywordTok{const} rl \OperatorTok{=} readline\OperatorTok{.}\CodeFuncTok{createInterface}(\{\}
  \CodeTypeTok{input}\OperatorTok{:} process\OperatorTok{.}\CodeTypeTok{stdin}\OperatorTok{,}
  \CodeTypeTok{output}\OperatorTok{:} process\OperatorTok{.}\CodeTypeTok{stdout}\OperatorTok{,}
\})\OperatorTok{;}
 
\KeywordTok{function} \CodeFuncTok{askQuestion}(prompt) \{
  \KeywordTok{return} \KeywordTok{new} Promise((resolve) \KeywordTok{=\textgreater{}} \{
    rl\OperatorTok{.}\CodeFuncTok{question}(prompt\OperatorTok{,} (answer) \KeywordTok{=\textgreater{}} \{
      \CodeFuncTok{resolve}(answer)\OperatorTok{;}
    \})\OperatorTok{;}
  )\OperatorTok{;}
\}
 
\KeywordTok{async} \KeywordTok{function} \CodeFuncTok{main}() \{
  \KeywordTok{const} name \OperatorTok{=} \KeywordTok{await} \CodeFuncTok{askQuestion}(\StringTok{"What is your name? "})\OperatorTok{;}
  \KeywordTok{const} ageStr \OperatorTok{=} \KeywordTok{await} \CodeFuncTok{askQuestion}(\StringTok{"How old are you? "})\OperatorTok{;}
  \KeywordTok{const} age \OperatorTok{=} \CodeAccentTok{parseInt}(ageStr)\OperatorTok{;}
 
  console\OperatorTok{.}\CodeFuncTok{log}(\StringTok{\textasciigrave{}Hello, }\StringTok{\$\{}name\StringTok{\}}\StringTok{!\textasciigrave{}})\OperatorTok{;}
  console\OperatorTok{.}\CodeFuncTok{log}(\StringTok{\textasciigrave{}Next year you will be }\StringTok{\$\{}age \OperatorTok{+} \CodeNumTok{1}\StringTok{\}}\StringTok{ years old.\textasciigrave{}})\OperatorTok{;}
 
  rl\OperatorTok{.}\CodeFuncTok{close}()\OperatorTok{;}
\}
 
\CodeFuncTok{main}()\OperatorTok{;}
\end{Highlighting}
\end{Shaded}

\subsection{Helper Module for Input
(Reusable)}\label{helper-module-for-input-reusable}

Create \texttt{input.js} once and reuse across projects:

\begin{Shaded}
\begin{Highlighting}[]
\CommentTok{// input.js {-}{-}{-} save this file in your project}
\KeywordTok{const} readline \OperatorTok{=} \CodeAccentTok{require}(\StringTok{"readline"})\OperatorTok{;}
 
\KeywordTok{const} rl \OperatorTok{=} readline\OperatorTok{.}\CodeFuncTok{createInterface}(\{\}
  \CodeTypeTok{input}\OperatorTok{:} process\OperatorTok{.}\CodeTypeTok{stdin}\OperatorTok{,}
  \CodeTypeTok{output}\OperatorTok{:} process\OperatorTok{.}\CodeTypeTok{stdout}\OperatorTok{,}
\})\OperatorTok{;}
 
\KeywordTok{function} \CodeFuncTok{askQuestion}(prompt) \{
  \KeywordTok{return} \KeywordTok{new} Promise((resolve) \KeywordTok{=\textgreater{}} \{
    rl\OperatorTok{.}\CodeFuncTok{question}(prompt\OperatorTok{,} resolve)\OperatorTok{;}
  \})\OperatorTok{;}
\}
 
\KeywordTok{function} \CodeFuncTok{close}() \{
  rl\OperatorTok{.}\CodeFuncTok{close}()\OperatorTok{;}
\}
 
module\OperatorTok{.}\CodeTypeTok{exports} \OperatorTok{=} \{ askQuestion\OperatorTok{,} close \}\OperatorTok{;}
\end{Highlighting}
\end{Shaded}

Usage in your main file:

\begin{Shaded}
\begin{Highlighting}[]
\CommentTok{// main.js}
\KeywordTok{const} \{ askQuestion\OperatorTok{,} close \} \OperatorTok{=} \CodeAccentTok{require}(\StringTok{"./input"})\OperatorTok{;}
 
\KeywordTok{async} \KeywordTok{function} \CodeFuncTok{main}() \{
  \KeywordTok{const} name \OperatorTok{=} \KeywordTok{await} \CodeFuncTok{askQuestion}(\StringTok{"Enter your name: "})\OperatorTok{;}
  \KeywordTok{const} numStr \OperatorTok{=} \KeywordTok{await} \CodeFuncTok{askQuestion}(\StringTok{"Enter a number: "})\OperatorTok{;}
  \KeywordTok{const} num \OperatorTok{=} \CodeAccentTok{parseFloat}(numStr)\OperatorTok{;}
 
  console\OperatorTok{.}\CodeFuncTok{log}(\StringTok{\textasciigrave{}Hello, }\StringTok{\$\{}name\StringTok{\}}\StringTok{!\textasciigrave{}})\OperatorTok{;}
  console\OperatorTok{.}\CodeFuncTok{log}(\StringTok{\textasciigrave{}Your number squared is: }\StringTok{\$\{}num \OperatorTok{**} \CodeNumTok{2}\StringTok{\}}\StringTok{\textasciigrave{}})\OperatorTok{;}
 
  \CodeFuncTok{close}()\OperatorTok{;}
\}
 
\CodeFuncTok{main}()\OperatorTok{;}
\end{Highlighting}
\end{Shaded}

\subsection{Examples}\label{examples-15}

\textbf{Example 2.2.1: Temperature Converter}

\begin{Shaded}
\begin{Highlighting}[]
\KeywordTok{const} \{ askQuestion\OperatorTok{,} close \} \OperatorTok{=} \CodeAccentTok{require}(\StringTok{"./input"})\OperatorTok{;}
 
\KeywordTok{async} \KeywordTok{function} \CodeFuncTok{main}() \{
  \KeywordTok{const} fahrenheitStr \OperatorTok{=} \KeywordTok{await} \CodeFuncTok{askQuestion}(\StringTok{"Enter temperature in Fahrenheit: "})\OperatorTok{;}
  \KeywordTok{const} fahrenheit \OperatorTok{=} \CodeAccentTok{parseFloat}(fahrenheitStr)\OperatorTok{;}
 
  \KeywordTok{const} celsius \OperatorTok{=} ((fahrenheit \OperatorTok{{-}} \CodeNumTok{32}) \OperatorTok{*} \CodeNumTok{5}) \OperatorTok{/} \CodeNumTok{9}\OperatorTok{;}
 
  console\OperatorTok{.}\CodeFuncTok{log}(\StringTok{\textasciigrave{}}\StringTok{\$\{}fahrenheit\StringTok{\}}\StringTok{°F is }\StringTok{\$\{}celsius\OperatorTok{.}\CodeFuncTok{toFixed}(\CodeNumTok{1})\StringTok{\}}\StringTok{°C\textasciigrave{}})\OperatorTok{;}
  \CodeFuncTok{close}()\OperatorTok{;}
\}
 
\CodeFuncTok{main}()\OperatorTok{;}
\end{Highlighting}
\end{Shaded}

\textbf{Example 2.2.2: Circle Calculator}

\begin{Shaded}
\begin{Highlighting}[]
\KeywordTok{const} \{ askQuestion\OperatorTok{,} close \} \OperatorTok{=} \CodeAccentTok{require}(\StringTok{"./input"})\OperatorTok{;}
 
\KeywordTok{async} \KeywordTok{function} \CodeFuncTok{main}() \{
  \KeywordTok{const} radiusStr \OperatorTok{=} \KeywordTok{await} \CodeFuncTok{askQuestion}(\StringTok{"Enter the radius: "})\OperatorTok{;}
  \KeywordTok{const} radius \OperatorTok{=} \CodeAccentTok{parseFloat}(radiusStr)\OperatorTok{;}
 
  \KeywordTok{const} area \OperatorTok{=} Math\OperatorTok{.}\CodeNumTok{PI} \OperatorTok{*} radius \OperatorTok{**} \CodeNumTok{2}\OperatorTok{;}
  \KeywordTok{const} circumference \OperatorTok{=} \CodeNumTok{2} \OperatorTok{*} Math\OperatorTok{.}\CodeNumTok{PI} \OperatorTok{*} radius\OperatorTok{;}
 
  console\OperatorTok{.}\CodeFuncTok{log}(\StringTok{\textasciigrave{}Area: }\StringTok{\$\{}area\OperatorTok{.}\CodeFuncTok{toFixed}(\CodeNumTok{2})\StringTok{\}}\StringTok{\textasciigrave{}})\OperatorTok{;}
  console\OperatorTok{.}\CodeFuncTok{log}(\StringTok{\textasciigrave{}Circumference: }\StringTok{\$\{}circumference\OperatorTok{.}\CodeFuncTok{toFixed}(\CodeNumTok{2})\StringTok{\}}\StringTok{\textasciigrave{}})\OperatorTok{;}
  \CodeFuncTok{close}()\OperatorTok{;}
\}
 
\CodeFuncTok{main}()\OperatorTok{;}
\end{Highlighting}
\end{Shaded}

\subsection{Skills Checklist}\label{skills-checklist-56}

\begin{itemize}
\tightlist
\item[$\square$]
  I can use \texttt{Math} object methods and constants.
\item[$\square$]
  I can get user input using the \texttt{readline} module.
\item[$\square$]
  I understand \texttt{async}/\texttt{await} for asynchronous input.
\item[$\square$]
  I can convert string input to numbers using \texttt{parseInt()} and
  \texttt{parseFloat()}.
\end{itemize}

\begin{center}\rule{0.5\linewidth}{0.5pt}\end{center}

\section{Unit 3: Control Flow --- Decisions and
Loops}\label{unit-3-control-flow-decisions-and-loops-2}

\begin{center}\rule{0.5\linewidth}{0.5pt}\end{center}

\subsection{\texorpdfstring{Lesson 3.1: Boolean Expressions and
\texttt{if}
Statements}{Lesson 3.1: Boolean Expressions and if Statements}}\label{lesson-3.1-boolean-expressions-and-if-statements-1}

\subsection{Comparison Operators}\label{comparison-operators}

\begin{fullwidth}{\def\LTcaptype{none} % do not increment counter
\begin{longtable}[]{@{}lll@{}}
\toprule\noalign{}
Operator & Meaning & Note \\
\midrule\noalign{}
\endhead
\bottomrule\noalign{}
\endlastfoot
\texttt{===} & Strict equal & Checks value AND type (preferred) \\
\texttt{!==} & Strict not equal & Checks value AND type (preferred) \\
\texttt{==} & Loose equal & Type coercion (avoid) \\
\texttt{!=} & Loose not equal & Type coercion (avoid) \\
\texttt{\textless{}} & Less than & \\
\texttt{\textgreater{}} & Greater than & \\
\texttt{\textless{}=} & Less than or equal & \\
\texttt{\textgreater{}=} & Greater than or equal & \\
\end{longtable}
}
\end{fullwidth}

\textbf{Always use \texttt{===} and \texttt{!==}} --- \texttt{==} does
unexpected type conversions.

\subsection{Logical Operators}\label{logical-operators}

\begin{fullwidth}{\def\LTcaptype{none} % do not increment counter
\begin{longtable}[]{@{}lll@{}}
\toprule\noalign{}
Operator & Meaning & Example \\
\midrule\noalign{}
\endhead
\bottomrule\noalign{}
\endlastfoot
\texttt{\&\&} & AND & \texttt{age\ \textgreater{}=\ 18\ \&\&\ hasId} \\
\texttt{\textbackslash{}\textbar{}\textbackslash{}\textbar{}} & OR &
\texttt{day\ ===\ \textquotesingle{}Sat\textquotesingle{}\ \textbackslash{}\textbar{}\textbackslash{}\textbar{}\ day\ ===\ \textquotesingle{}Sun\textquotesingle{}} \\
\texttt{!} & NOT & \texttt{!isRaining} \\
\end{longtable}
}
\end{fullwidth}

\subsection{\texorpdfstring{\texttt{if-else\ if-else}}{if-else if-else}}\label{if-else-if-else}

\begin{Shaded}
\begin{Highlighting}[]
\KeywordTok{const} score \OperatorTok{=} \CodeNumTok{85}\OperatorTok{;}
 
\KeywordTok{if} (score \OperatorTok{\textgreater{}=} \CodeNumTok{90}) \{
  console\OperatorTok{.}\CodeFuncTok{log}(\StringTok{"Grade: A"})\OperatorTok{;}
\} \KeywordTok{else} \KeywordTok{if} (score \OperatorTok{\textgreater{}=} \CodeNumTok{80}) \{
  console\OperatorTok{.}\CodeFuncTok{log}(\StringTok{"Grade: B"})\OperatorTok{;}
\} \KeywordTok{else} \KeywordTok{if} (score \OperatorTok{\textgreater{}=} \CodeNumTok{70}) \{
  console\OperatorTok{.}\CodeFuncTok{log}(\StringTok{"Grade: C"})\OperatorTok{;}
\} \KeywordTok{else} \{
  console\OperatorTok{.}\CodeFuncTok{log}(\StringTok{"Grade: F"})\OperatorTok{;}
\}
\end{Highlighting}
\end{Shaded}

\subsection{Ternary Operator}\label{ternary-operator}

\begin{Shaded}
\begin{Highlighting}[]
\KeywordTok{const} age \OperatorTok{=} \CodeNumTok{20}\OperatorTok{;}
\KeywordTok{const} status \OperatorTok{=} age \OperatorTok{\textgreater{}=} \CodeNumTok{18} \OperatorTok{?} \StringTok{"adult"} \OperatorTok{:} \StringTok{"minor"}\OperatorTok{;}
console\OperatorTok{.}\CodeFuncTok{log}(status)\OperatorTok{;} \CommentTok{// "adult"}
\end{Highlighting}
\end{Shaded}

\subsection{\texorpdfstring{\texttt{switch}
Statement}{switch Statement}}\label{switch-statement}

\begin{Shaded}
\begin{Highlighting}[]
\KeywordTok{const} day \OperatorTok{=} \StringTok{"Wednesday"}\OperatorTok{;}
\KeywordTok{switch} (day) \{
  \KeywordTok{case} \StringTok{"Monday"}\OperatorTok{:}
  \KeywordTok{case} \StringTok{"Tuesday"}\OperatorTok{:}
  \KeywordTok{case} \StringTok{"Wednesday"}\OperatorTok{:}
  \KeywordTok{case} \StringTok{"Thursday"}\OperatorTok{:}
  \KeywordTok{case} \StringTok{"Friday"}\OperatorTok{:}
    console\OperatorTok{.}\CodeFuncTok{log}(\StringTok{"Weekday"})\OperatorTok{;}
    \KeywordTok{break}\OperatorTok{;}
  \KeywordTok{case} \StringTok{"Saturday"}\OperatorTok{:}
  \KeywordTok{case} \StringTok{"Sunday"}\OperatorTok{:}
    console\OperatorTok{.}\CodeFuncTok{log}(\StringTok{"Weekend"})\OperatorTok{;}
    \KeywordTok{break}\OperatorTok{;}
  \KeywordTok{default}\OperatorTok{:}
    console\OperatorTok{.}\CodeFuncTok{log}(\StringTok{"Unknown day"})\OperatorTok{;}
\}
\end{Highlighting}
\end{Shaded}

\subsection{Truthy and Falsy Values}\label{truthy-and-falsy-values}

In JavaScript, the following values are \textbf{falsy} (evaluate to
\texttt{false} in conditions):

\begin{itemize}
\tightlist
\item
  \texttt{false}, \texttt{0}, \texttt{""} (empty string), \texttt{null},
  \texttt{undefined}, \texttt{NaN}
\end{itemize}

Everything else is \textbf{truthy}.

\begin{Shaded}
\begin{Highlighting}[]
\KeywordTok{const} name \OperatorTok{=} \StringTok{""}\OperatorTok{;}
\KeywordTok{if} (name) \{
  console\OperatorTok{.}\CodeFuncTok{log}(\StringTok{"Name is set"})\OperatorTok{;}
\} \KeywordTok{else} \{
  console\OperatorTok{.}\CodeFuncTok{log}(\StringTok{"Name is empty"})\OperatorTok{;} \CommentTok{// This runs}
\}
\end{Highlighting}
\end{Shaded}

\subsection{Examples}\label{examples-16}

\textbf{Example 3.1.1: Grade Calculator}

\begin{Shaded}
\begin{Highlighting}[]
\KeywordTok{const} \{ askQuestion\OperatorTok{,} close \} \OperatorTok{=} \CodeAccentTok{require}(\StringTok{"./input"})\OperatorTok{;}
 
\KeywordTok{async} \KeywordTok{function} \CodeFuncTok{main}() \{
  \KeywordTok{const} scoreStr \OperatorTok{=} \KeywordTok{await} \CodeFuncTok{askQuestion}(\StringTok{"Enter your score (0{-}100): "})\OperatorTok{;}
  \KeywordTok{const} score \OperatorTok{=} \CodeAccentTok{parseInt}(scoreStr)\OperatorTok{;}
 
  \KeywordTok{let} grade\OperatorTok{;}
  \KeywordTok{if} (score \OperatorTok{\textgreater{}=} \CodeNumTok{90}) grade \OperatorTok{=} \StringTok{"A"}\OperatorTok{;}
  \KeywordTok{else} \KeywordTok{if} (score \OperatorTok{\textgreater{}=} \CodeNumTok{80}) grade \OperatorTok{=} \StringTok{"B"}\OperatorTok{;}
  \KeywordTok{else} \KeywordTok{if} (score \OperatorTok{\textgreater{}=} \CodeNumTok{70}) grade \OperatorTok{=} \StringTok{"C"}\OperatorTok{;}
  \KeywordTok{else} \KeywordTok{if} (score \OperatorTok{\textgreater{}=} \CodeNumTok{60}) grade \OperatorTok{=} \StringTok{"D"}\OperatorTok{;}
  \KeywordTok{else} grade \OperatorTok{=} \StringTok{"F"}\OperatorTok{;}
 
  console\OperatorTok{.}\CodeFuncTok{log}(\StringTok{\textasciigrave{}Your grade: }\StringTok{\$\{}grade\StringTok{\}}\StringTok{\textasciigrave{}})\OperatorTok{;}
  \CodeFuncTok{close}()\OperatorTok{;}
\}
 
\CodeFuncTok{main}()\OperatorTok{;}
\end{Highlighting}
\end{Shaded}

\subsection{Skills Checklist}\label{skills-checklist-57}

\begin{itemize}
\tightlist
\item[$\square$]
  I use \texttt{===} and \texttt{!==} for comparisons in JavaScript.
\item[$\square$]
  I can write \texttt{if}, \texttt{else\ if}, and \texttt{else}
  statements.
\item[$\square$]
  I understand truthy and falsy values.
\item[$\square$]
  I can write a \texttt{switch} statement.
\end{itemize}

\begin{center}\rule{0.5\linewidth}{0.5pt}\end{center}

\subsection{\texorpdfstring{Lesson 3.2: \texttt{for}
Loops}{Lesson 3.2: for Loops}}\label{lesson-3.2-for-loops}

\begin{Shaded}
\begin{Highlighting}[]
\CommentTok{// Basic for loop}
\KeywordTok{for} (\KeywordTok{let} i \OperatorTok{=} \CodeNumTok{0}\OperatorTok{;} i \OperatorTok{\textless{}} \CodeNumTok{5}\OperatorTok{;} i\OperatorTok{++}) \{
  console\OperatorTok{.}\CodeFuncTok{log}(\StringTok{\textasciigrave{}i = }\StringTok{\$\{}i\StringTok{\}}\StringTok{\textasciigrave{}})\OperatorTok{;}
\}
 
\CommentTok{// Counting down}
\KeywordTok{for} (\KeywordTok{let} i \OperatorTok{=} \CodeNumTok{10}\OperatorTok{;} i \OperatorTok{\textgreater{}} \CodeNumTok{0}\OperatorTok{;} i\OperatorTok{{-}{-}}) \{
  console\OperatorTok{.}\CodeFuncTok{log}(i)\OperatorTok{;}
\}
console\OperatorTok{.}\CodeFuncTok{log}(\StringTok{"Blast off!"})\OperatorTok{;}
 
\CommentTok{// for...of {-}{-}{-} iterate over array elements (like Python\textquotesingle{}s for{-}in)}
\KeywordTok{const} fruits \OperatorTok{=} [\StringTok{"apple"}\OperatorTok{,} \StringTok{"banana"}\OperatorTok{,} \StringTok{"cherry"}]\OperatorTok{;}
\KeywordTok{for} (\KeywordTok{const} fruit \KeywordTok{of} fruits) \{
  console\OperatorTok{.}\CodeFuncTok{log}(\StringTok{\textasciigrave{}I like }\StringTok{\$\{}fruit\StringTok{\}}\StringTok{.\textasciigrave{}})\OperatorTok{;}
\}
 
\CommentTok{// for...in {-}{-}{-} iterate over object keys}
\KeywordTok{const} person \OperatorTok{=} \{ \CodeTypeTok{name}\OperatorTok{:} \StringTok{"Alice"}\OperatorTok{,} \CodeTypeTok{age}\OperatorTok{:} \CodeNumTok{15}\OperatorTok{,} \CodeTypeTok{city}\OperatorTok{:} \StringTok{"Portland"} \}\OperatorTok{;}
\KeywordTok{for} (\KeywordTok{const} key \KeywordTok{in} person) \{
  console\OperatorTok{.}\CodeFuncTok{log}(\StringTok{\textasciigrave{}}\StringTok{\$\{}key\StringTok{\}}\StringTok{: }\StringTok{\$\{}person[key]\StringTok{\}}\StringTok{\textasciigrave{}})\OperatorTok{;}
\}
\end{Highlighting}
\end{Shaded}

\subsection{\texorpdfstring{\texttt{Array.forEach()}}{Array.forEach()}}\label{array.foreach}

\begin{Shaded}
\begin{Highlighting}[]
\KeywordTok{const} numbers \OperatorTok{=} [\CodeNumTok{1}\OperatorTok{,} \CodeNumTok{2}\OperatorTok{,} \CodeNumTok{3}\OperatorTok{,} \CodeNumTok{4}\OperatorTok{,} \CodeNumTok{5}]\OperatorTok{;}
 
numbers\OperatorTok{.}\CodeFuncTok{forEach}((num) \KeywordTok{=\textgreater{}} \{
  console\OperatorTok{.}\CodeFuncTok{log}(num \OperatorTok{*} \CodeNumTok{2})\OperatorTok{;}
\})\OperatorTok{;}
 
\CommentTok{// With index}
numbers\OperatorTok{.}\CodeFuncTok{forEach}((num\OperatorTok{,} index) \KeywordTok{=\textgreater{}} \{
  console\OperatorTok{.}\CodeFuncTok{log}(\StringTok{\textasciigrave{}Index }\StringTok{\$\{}index\StringTok{\}}\StringTok{: }\StringTok{\$\{}num\StringTok{\}}\StringTok{\textasciigrave{}})\OperatorTok{;}
)\OperatorTok{;}
\end{Highlighting}
\end{Shaded}

\subsection{Examples}\label{examples-17}

\textbf{Example 3.2.1: Multiplication Table}

\begin{Shaded}
\begin{Highlighting}[]
\KeywordTok{const} n \OperatorTok{=} \CodeNumTok{5}\OperatorTok{;}
console\OperatorTok{.}\CodeFuncTok{log}(\StringTok{\textasciigrave{}Multiplication table for }\StringTok{\$\{}n\StringTok{\}}\StringTok{:\textasciigrave{}})\OperatorTok{;}
\KeywordTok{for} (\KeywordTok{let} i \OperatorTok{=} \CodeNumTok{1}\OperatorTok{;} i \OperatorTok{\textless{}=} \CodeNumTok{10}\OperatorTok{;} i\OperatorTok{++}) \{
  console\OperatorTok{.}\CodeFuncTok{log}(\StringTok{\textasciigrave{}}\StringTok{\$\{}n\StringTok{\}}\StringTok{ x }\StringTok{\$\{}i\StringTok{\}}\StringTok{ = }\StringTok{\$\{}n \OperatorTok{*} i\StringTok{\}}\StringTok{\textasciigrave{}})\OperatorTok{;}
\}
\end{Highlighting}
\end{Shaded}

\subsection{Skills Checklist}\label{skills-checklist-58}

\begin{itemize}
\tightlist
\item[$\square$]
  I can write a \texttt{for} loop with initialization, condition, and
  update.
\item[$\square$]
  I can use \texttt{for...of} to iterate over arrays.
\item[$\square$]
  I can use \texttt{for...in} to iterate over object keys.
\item[$\square$]
  I can use \texttt{Array.forEach()}.
\end{itemize}

\begin{center}\rule{0.5\linewidth}{0.5pt}\end{center}

\subsection{\texorpdfstring{Lesson 3.3: \texttt{while}
Loops}{Lesson 3.3: while Loops}}\label{lesson-3.3-while-loops}

\begin{Shaded}
\begin{Highlighting}[]
\KeywordTok{let} count \OperatorTok{=} \CodeNumTok{1}\OperatorTok{;}
\KeywordTok{while} (count \OperatorTok{\textless{}=} \CodeNumTok{5}) \{
  console\OperatorTok{.}\CodeFuncTok{log}(\StringTok{\textasciigrave{}Count: }\StringTok{\$\{}count\StringTok{\}}\StringTok{\textasciigrave{}})\OperatorTok{;}
  count\OperatorTok{++;}
\}
 
\CommentTok{// do{-}while {-}{-}{-} runs at least once}
\KeywordTok{let} num\OperatorTok{;}
\KeywordTok{do} \{
  num \OperatorTok{=} Math\OperatorTok{.}\CodeFuncTok{floor}(Math\OperatorTok{.}\CodeFuncTok{random}() \OperatorTok{*} \CodeNumTok{10})\OperatorTok{;}
  console\OperatorTok{.}\CodeFuncTok{log}(\StringTok{\textasciigrave{}Rolled: }\StringTok{\$\{}num\StringTok{\}}\StringTok{\textasciigrave{}})\OperatorTok{;}
\ \}\KeywordTok{while} (num \OperatorTok{\textless{}} \CodeNumTok{5})\OperatorTok{;}
\end{Highlighting}
\end{Shaded}

\subsection{\texorpdfstring{\texttt{break} and
\texttt{continue}}{break and continue}}\label{break-and-continue}

\begin{Shaded}
\begin{Highlighting}[]
\CommentTok{// break {-}{-}{-} exit entirely}
\KeywordTok{for} (\KeywordTok{let} i \OperatorTok{=} \CodeNumTok{0}\OperatorTok{;} i \OperatorTok{\textless{}} \CodeNumTok{10}\OperatorTok{;} i\OperatorTok{++}) \{
  \KeywordTok{if} (i \OperatorTok{===} \CodeNumTok{5}) \KeywordTok{break}\OperatorTok{;}
  console\OperatorTok{.}\CodeFuncTok{log}(i)\OperatorTok{;} \CommentTok{// Prints 0{-}4}
\}
 
\CommentTok{// continue {-}{-}{-} skip iteration}
\KeywordTok{for} (\KeywordTok{let} i \OperatorTok{=} \CodeNumTok{0}\OperatorTok{;} i \OperatorTok{\textless{}} \CodeNumTok{10}\OperatorTok{;} i\OperatorTok{++}) \{\}
  \KeywordTok{if} (i \OperatorTok{\%} \CodeNumTok{2} \OperatorTok{===} \CodeNumTok{0}) \KeywordTok{continue}\OperatorTok{;}
  console\OperatorTok{.}\CodeFuncTok{log}(i)\OperatorTok{;} \CommentTok{// Prints 1, 3, 5, 7, 9}
\}
\end{Highlighting}
\end{Shaded}

\subsection{Skills Checklist}\label{skills-checklist-59}

\begin{itemize}
\tightlist
\item[$\square$]
  I can write \texttt{while} and \texttt{do-while} loops.
\item[$\square$]
  I can use \texttt{break} and \texttt{continue} appropriately.
\end{itemize}

\begin{center}\rule{0.5\linewidth}{0.5pt}\end{center}

\section{Unit 4: Data Structures --- Arrays and
Objects}\label{unit-4-data-structures-arrays-and-objects}

\begin{center}\rule{0.5\linewidth}{0.5pt}\end{center}

\subsection{Lesson 4.1: Arrays}\label{lesson-4.1-arrays}\amn{\textbf{Accessibility Tip:} To inspect a JavaScript array accessibly, use: \texttt{arr.forEach((item, i) => console.log(`\$\{i+1\}. \$\{item\}`))}. This prints each item numbered on its own line. Never use \texttt{console.log(arr)} alone---the bracket notation it produces is hard to navigate when spoken by a screen reader.}
\subsection{Lesson 4.1: Arrays}\label{lesson-4.1-arrays}

\subsection{Creating and Accessing
Arrays}\label{creating-and-accessing-arrays}

\begin{Shaded}
\begin{Highlighting}[]
\KeywordTok{const} colors \OperatorTok{=} [\StringTok{"red"}\OperatorTok{,} \StringTok{"blue"}\OperatorTok{,} \StringTok{"green"}\OperatorTok{,} \StringTok{"yellow"}]\OperatorTok{;}
 
\CommentTok{// Access}
console\OperatorTok{.}\CodeFuncTok{log}(colors[\CodeNumTok{0}])\OperatorTok{;} \CommentTok{// "red"}
console\OperatorTok{.}\CodeFuncTok{log}(colors[colors\OperatorTok{.}\CodeTypeTok{length} \OperatorTok{{-}} \CodeNumTok{1}])\OperatorTok{;} \CommentTok{// "yellow"}
 
\CommentTok{// Modify}
colors[\CodeNumTok{1}] \OperatorTok{=} \StringTok{"purple"}\OperatorTok{;}
 
\CommentTok{// Check}
console\OperatorTok{.}\CodeFuncTok{log}(colors\OperatorTok{.}\CodeTypeTok{length})\OperatorTok{;} \CommentTok{// 4}
console\OperatorTok{.}\CodeFuncTok{log}(colors\OperatorTok{.}\CodeFuncTok{includes}(\StringTok{"red"}))\OperatorTok{;} \CommentTok{// true}
console\OperatorTok{.}\CodeFuncTok{log}(colors\OperatorTok{.}\CodeFuncTok{indexOf}(\StringTok{"green"}))\OperatorTok{;} \CommentTok{// 2}
\end{Highlighting}
\end{Shaded}

\subsection{Array Methods}\label{array-methods}

\begin{Shaded}
\begin{Highlighting}[]
\KeywordTok{const} fruits \OperatorTok{=} [\StringTok{"apple"}\OperatorTok{,} \StringTok{"banana"}]\OperatorTok{;}
 
\CommentTok{// Adding}
fruits\OperatorTok{.}\CodeFuncTok{push}(\StringTok{"cherry"})\OperatorTok{;} \CommentTok{// Add to end {-}\textgreater{} ["apple", "banana", "cherry"]}
fruits\OperatorTok{.}\CodeFuncTok{unshift}(\StringTok{"avocado"})\OperatorTok{;} \CommentTok{// Add to start {-}\textgreater{} ["avocado", "apple", "banana", "cherry"]}
 
\CommentTok{// Removing}
\KeywordTok{const} last \OperatorTok{=} fruits\OperatorTok{.}\CodeFuncTok{pop}()\OperatorTok{;} \CommentTok{// Remove and return last}
\KeywordTok{const} first \OperatorTok{=} fruits\OperatorTok{.}\CodeFuncTok{shift}()\OperatorTok{;} \CommentTok{// Remove and return first}
 
\CommentTok{// Slicing (non{-}destructive)}
\KeywordTok{const} sub \OperatorTok{=} fruits\OperatorTok{.}\CodeFuncTok{slice}(\CodeNumTok{1}\OperatorTok{,} \CodeNumTok{3})\OperatorTok{;} \CommentTok{// Returns new array, original unchanged}
 
\CommentTok{// Splicing (destructive {-}{-}{-} modifies original)}
fruits\OperatorTok{.}\CodeFuncTok{splice}(\CodeNumTok{1}\OperatorTok{,} \CodeNumTok{1})\OperatorTok{;} \CommentTok{// Remove 1 element at index 1}
fruits\OperatorTok{.}\CodeFuncTok{splice}(\CodeNumTok{1}\OperatorTok{,} \CodeNumTok{0}\OperatorTok{,} \StringTok{"mango"})\OperatorTok{;} \CommentTok{// Insert "mango" at index 1}
 
\CommentTok{// Joining}
\KeywordTok{const} str \OperatorTok{=} fruits\OperatorTok{.}\CodeFuncTok{join}(\StringTok{", "})\OperatorTok{;} \CommentTok{// "apple, mango, cherry"}
 
\CommentTok{// Sorting}
\KeywordTok{const} nums \OperatorTok{=} [\CodeNumTok{3}\OperatorTok{,} \CodeNumTok{1}\OperatorTok{,} \CodeNumTok{4}\OperatorTok{,} \CodeNumTok{1}\OperatorTok{,} \CodeNumTok{5}\OperatorTok{,} \CodeNumTok{9}\OperatorTok{,} \CodeNumTok{2}\OperatorTok{,} \CodeNumTok{6}]\OperatorTok{;}
nums\OperatorTok{.}\CodeFuncTok{sort}((a\OperatorTok{,} b) \KeywordTok{=\textgreater{}} a \OperatorTok{{-}} b)\OperatorTok{;} \CommentTok{// Ascending: [1, 1, 2, 3, 4, 5, 6, 9]}
nums\OperatorTok{.}\CodeFuncTok{sort}((a\OperatorTok{,} b) \KeywordTok{=\textgreater{}} b \OperatorTok{{-}} a)\OperatorTok{;} \CommentTok{// Descending}
 
\CommentTok{// Reversing}
\KeywordTok{const} reversed \OperatorTok{=} [\OperatorTok{...}fruits]\OperatorTok{.}\CodeFuncTok{reverse}()\OperatorTok{;} \CommentTok{// Use spread to avoid mutating original}
 
\CommentTok{// Concatenating}
\KeywordTok{const} combined \OperatorTok{=} fruits\OperatorTok{.}\CodeFuncTok{concat}(nums)\OperatorTok{;}
\end{Highlighting}
\end{Shaded}

\subsection{Functional Array Methods}\label{functional-array-methods}

These are powerful for data transformation:

\begin{Shaded}
\begin{Highlighting}[]
\KeywordTok{const} numbers \OperatorTok{=} [\CodeNumTok{1}\OperatorTok{,} \CodeNumTok{2}\OperatorTok{,} \CodeNumTok{3}\OperatorTok{,} \CodeNumTok{4}\OperatorTok{,} \CodeNumTok{5}\OperatorTok{,} \CodeNumTok{6}\OperatorTok{,} \CodeNumTok{7}\OperatorTok{,} \CodeNumTok{8}\OperatorTok{,} \CodeNumTok{9}\OperatorTok{,} \CodeNumTok{10}]\OperatorTok{;}
 
\CommentTok{// map {-}{-}{-} transform each element, returns new array}
\KeywordTok{const} doubled \OperatorTok{=} numbers\OperatorTok{.}\CodeFuncTok{map}((n) \KeywordTok{=\textgreater{}} n \OperatorTok{*} \CodeNumTok{2})\OperatorTok{;}
console\OperatorTok{.}\CodeFuncTok{log}(doubled)\OperatorTok{;} \CommentTok{// [2, 4, 6, 8, 10, 12, 14, 16, 18, 20]}
 
\CommentTok{// filter {-}{-}{-} keep elements matching condition}
\KeywordTok{const} evens \OperatorTok{=} numbers\OperatorTok{.}\CodeFuncTok{filter}((n) \KeywordTok{=\textgreater{}} n \OperatorTok{\%} \CodeNumTok{2} \OperatorTok{===} \CodeNumTok{0})\OperatorTok{;}
console\OperatorTok{.}\CodeFuncTok{log}(evens)\OperatorTok{;} \CommentTok{// [2, 4, 6, 8, 10]}
 
\CommentTok{// reduce {-}{-}{-} combine elements into single value}
\KeywordTok{const} sum \OperatorTok{=} numbers\OperatorTok{.}\CodeFuncTok{reduce}((accumulator\OperatorTok{,} current) \KeywordTok{=\textgreater{}} accumulator \OperatorTok{+} current\OperatorTok{,} \CodeNumTok{0})\OperatorTok{;}
console\OperatorTok{.}\CodeFuncTok{log}(sum)\OperatorTok{;} \CommentTok{// 55}
 
\CommentTok{// find {-}{-}{-} return first matching element}
\KeywordTok{const} firstOver5 \OperatorTok{=} numbers\OperatorTok{.}\CodeFuncTok{find}((n) \KeywordTok{=\textgreater{}} n \OperatorTok{\textgreater{}} \CodeNumTok{5})\OperatorTok{;}
console\OperatorTok{.}\CodeFuncTok{log}(firstOver5)\OperatorTok{;} \CommentTok{// 6}
 
\CommentTok{// some {-}{-}{-} true if any element matches}
\KeywordTok{const} hasEven \OperatorTok{=} numbers\OperatorTok{.}\CodeFuncTok{some}((n) \KeywordTok{=\textgreater{}} n \OperatorTok{\%} \CodeNumTok{2} \OperatorTok{===} \CodeNumTok{0})\OperatorTok{;}
console\OperatorTok{.}\CodeFuncTok{log}(hasEven)\OperatorTok{;} \CommentTok{// true}
 
\CommentTok{// every {-}{-}{-} true if all elements match}
\KeywordTok{const} allPositive \OperatorTok{=} numbers\OperatorTok{.}\CodeFuncTok{every}((n) \KeywordTok{=\textgreater{}} n \OperatorTok{\textgreater{}} \CodeNumTok{0})\OperatorTok{;}
console\OperatorTok{.}\CodeFuncTok{log}(allPositive)\OperatorTok{;} \CommentTok{// true}
 
\CommentTok{// Chaining}
\KeywordTok{const} result \OperatorTok{=} numbers
  \OperatorTok{.}\CodeFuncTok{filter}((n) \KeywordTok{=\textgreater{}} n \OperatorTok{\%} \CodeNumTok{2} \OperatorTok{===} \CodeNumTok{0})
  \OperatorTok{.}\CodeFuncTok{map}((n) \KeywordTok{=\textgreater{}} n \OperatorTok{**} \CodeNumTok{2})
  \OperatorTok{.}\CodeFuncTok{reduce}((sum\OperatorTok{,} n) \KeywordTok{=\textgreater{}} sum \OperatorTok{+} n\OperatorTok{,} \CodeNumTok{0})\OperatorTok{;}
console\OperatorTok{.}\CodeFuncTok{log}(result)\OperatorTok{;} \CommentTok{// 220  (4+16+36+64+100)}
\end{Highlighting}
\end{Shaded}

\subsection{Examples}\label{examples-18}

\textbf{Example 4.1.1: Array Analysis}

\begin{Shaded}
\begin{Highlighting}[]
\KeywordTok{const} scores \OperatorTok{=} [\CodeNumTok{85}\OperatorTok{,} \CodeNumTok{92}\OperatorTok{,} \CodeNumTok{78}\OperatorTok{,} \CodeNumTok{95}\OperatorTok{,} \CodeNumTok{88}\OperatorTok{,} \CodeNumTok{72}\OperatorTok{,} \CodeNumTok{90}\OperatorTok{,} \CodeNumTok{83}]\OperatorTok{;}
 
\KeywordTok{const} average \OperatorTok{=} scores\OperatorTok{.}\CodeFuncTok{reduce}((sum\OperatorTok{,} s) \KeywordTok{=\textgreater{}} sum \OperatorTok{+} s\OperatorTok{,} \CodeNumTok{0}) \OperatorTok{/} scores\OperatorTok{.}\CodeTypeTok{length}\OperatorTok{;}
\KeywordTok{const} max \OperatorTok{=} Math\OperatorTok{.}\CodeFuncTok{max}(\OperatorTok{...}scores)\OperatorTok{;}
\KeywordTok{const} min \OperatorTok{=} Math\OperatorTok{.}\CodeFuncTok{min}(\OperatorTok{...}scores)\OperatorTok{;}
\KeywordTok{const} passing \OperatorTok{=} scores\OperatorTok{.}\CodeFuncTok{filter}((s) \KeywordTok{=\textgreater{}} s \OperatorTok{\textgreater{}=} \CodeNumTok{80})\OperatorTok{;}
 
console\OperatorTok{.}\CodeFuncTok{log}(\StringTok{\textasciigrave{}Scores: }\StringTok{\$\{}scores\OperatorTok{.}\CodeFuncTok{join}(\StringTok{", "})\StringTok{\}}\StringTok{\textasciigrave{}})\OperatorTok{;}
console\OperatorTok{.}\CodeFuncTok{log}(\StringTok{\textasciigrave{}Average: }\StringTok{\$\{}average\OperatorTok{.}\CodeFuncTok{toFixed}(\CodeNumTok{1})\StringTok{\}}\StringTok{\textasciigrave{}})\OperatorTok{;}
console\OperatorTok{.}\CodeFuncTok{log}(\StringTok{\textasciigrave{}Highest: }\StringTok{\$\{}max\StringTok{\}}\StringTok{\textasciigrave{}})\OperatorTok{;}
console\OperatorTok{.}\CodeFuncTok{log}(\StringTok{\textasciigrave{}Lowest: }\StringTok{\$\{}min\StringTok{\}}\StringTok{\textasciigrave{}})\OperatorTok{;}
console\OperatorTok{.}\CodeFuncTok{log}(\StringTok{\textasciigrave{}Passing (\textgreater{}=80): }\StringTok{\$\{}passing\OperatorTok{.}\CodeTypeTok{length}\StringTok{\}}\StringTok{ students\textasciigrave{}})\OperatorTok{;}
\end{Highlighting}
\end{Shaded}

\subsection{Skills Checklist}\label{skills-checklist-60}

\begin{itemize}
\tightlist
\item[$\square$]
  I can create, access, and modify arrays.
\item[$\square$]
  I can use \texttt{push}, \texttt{pop}, \texttt{shift},
  \texttt{unshift}, \texttt{splice}, and \texttt{slice}.
\item[$\square$]
  I can use \texttt{map}, \texttt{filter}, and \texttt{reduce}.
\item[$\square$]
  I can sort arrays using a comparator function.
\end{itemize}

\begin{center}\rule{0.5\linewidth}{0.5pt}\end{center}

\subsection{Lesson 4.2: Objects}\label{lesson-4.2-objects-1}

\subsection{Creating and Using
Objects}\label{creating-and-using-objects}

Objects in JavaScript are like Python dictionaries and Java Maps
combined with class instances.

\begin{Shaded}
\begin{Highlighting}[]
\CommentTok{// Object literal}
\KeywordTok{const} studentProfile \OperatorTok{=} \{
  \CodeTypeTok{name}\OperatorTok{:} \StringTok{"Emily"}\OperatorTok{,}
  \CodeTypeTok{age}\OperatorTok{:} \CodeNumTok{13}\OperatorTok{,}
  \CodeTypeTok{grade}\OperatorTok{:} \StringTok{"7th"}\OperatorTok{,}
  \CodeTypeTok{favoriteSubject}\OperatorTok{:} \StringTok{"Science"}\OperatorTok{,}
\}\OperatorTok{;}
 
\CommentTok{// Access properties}
console\OperatorTok{.}\CodeFuncTok{log}(studentProfile\OperatorTok{.}\CodeTypeTok{name})\OperatorTok{;} \CommentTok{// "Emily" (dot notation)}
console\OperatorTok{.}\CodeFuncTok{log}(studentProfile[\StringTok{"age"}])\OperatorTok{;} \CommentTok{// 13 (bracket notation)}
 
\CommentTok{// Add/modify}
studentProfile\OperatorTok{.}\CodeTypeTok{city} \OperatorTok{=} \StringTok{"Springfield"}\OperatorTok{;}
studentProfile\OperatorTok{.}\CodeTypeTok{age} \OperatorTok{=} \CodeNumTok{14}\OperatorTok{;}
 
\CommentTok{// Delete}
\KeywordTok{delete} studentProfile\OperatorTok{.}\CodeTypeTok{grade}\OperatorTok{;}
 
\CommentTok{// Check if property exists}
console\OperatorTok{.}\CodeFuncTok{log}(\StringTok{"name"} \KeywordTok{in} studentProfile)\OperatorTok{;} \CommentTok{// true}
console\OperatorTok{.}\CodeFuncTok{log}(studentProfile\OperatorTok{.}\CodeFuncTok{hasOwnProperty}(\StringTok{"city"}))\OperatorTok{;} \CommentTok{// true}
 
\CommentTok{// Get keys, values, entries}
console\OperatorTok{.}\CodeFuncTok{log}(Object\OperatorTok{.}\CodeFuncTok{keys}(studentProfile))\OperatorTok{;}
console\OperatorTok{.}\CodeFuncTok{log}(Object\OperatorTok{.}\CodeFuncTok{values}(studentProfile))\OperatorTok{;}
console\OperatorTok{.}\CodeFuncTok{log}(Object\OperatorTok{.}\CodeFuncTok{entries}(studentProfile))\OperatorTok{;}
\end{Highlighting}
\end{Shaded}

\subsection{Iterating Over Objects}\label{iterating-over-objects}

\begin{Shaded}
\begin{Highlighting}[]
\KeywordTok{const} inventory \OperatorTok{=} \{ \CodeTypeTok{apples}\OperatorTok{:} \CodeNumTok{50}\OperatorTok{,} \CodeTypeTok{bananas}\OperatorTok{:} \CodeNumTok{30}\OperatorTok{,} \CodeTypeTok{oranges}\OperatorTok{:} \CodeNumTok{40} \}\OperatorTok{;}
 
\CommentTok{// for...in}
\KeywordTok{for} (\KeywordTok{const} key \KeywordTok{in} inventory) \{
  console\OperatorTok{.}\CodeFuncTok{log}(\StringTok{\textasciigrave{}}\StringTok{\$\{}key\StringTok{\}}\StringTok{: }\StringTok{\$\{}inventory[key]\StringTok{\}}\StringTok{\textasciigrave{}})\OperatorTok{;}
\}
 
\CommentTok{// Object.entries() with destructuring}
\KeywordTok{for} (\KeywordTok{const} [key\OperatorTok{,} value] \KeywordTok{of} Object\OperatorTok{.}\CodeFuncTok{entries}(inventory)) \{
  console\OperatorTok{.}\CodeFuncTok{log}(\StringTok{\textasciigrave{}}\StringTok{\$\{}key\StringTok{\}}\StringTok{: }\StringTok{\$\{}value\StringTok{\}}\StringTok{ in stock\textasciigrave{}})\OperatorTok{;}
\}
\end{Highlighting}
\end{Shaded}

\subsection{Nested Objects and
Arrays}\label{nested-objects-and-arrays}

\begin{Shaded}
\begin{Highlighting}[]
\KeywordTok{const} library \OperatorTok{=} \{
  \CodeTypeTok{name}\OperatorTok{:} \StringTok{"City Library"}\OperatorTok{,}
  \CodeTypeTok{location}\OperatorTok{:} \StringTok{"Downtown"}\OperatorTok{,}
  \CodeTypeTok{books}\OperatorTok{:} [
    \{ \CodeTypeTok{title}\OperatorTok{:} \StringTok{"Dune"}\OperatorTok{,} \CodeTypeTok{author}\OperatorTok{:} \StringTok{"Frank Herbert"}\OperatorTok{,} \CodeTypeTok{available}\OperatorTok{:} \KeywordTok{true} \}\OperatorTok{,}
    \{ \CodeTypeTok{title}\OperatorTok{:} \StringTok{"1984"}\OperatorTok{,} \CodeTypeTok{author}\OperatorTok{:} \StringTok{"George Orwell"}\OperatorTok{,} \CodeTypeTok{available}\OperatorTok{:} \KeywordTok{false} \}\OperatorTok{,}
  ]\OperatorTok{,}
\}\OperatorTok{;}
 
\CommentTok{// Access nested data}
console\OperatorTok{.}\CodeFuncTok{log}(library\OperatorTok{.}\CodeTypeTok{books}[\CodeNumTok{0}]\OperatorTok{.}\CodeTypeTok{title})\OperatorTok{;} \CommentTok{// "Dune"}
console\OperatorTok{.}\CodeFuncTok{log}(library\OperatorTok{.}\CodeTypeTok{books}[\CodeNumTok{1}]\OperatorTok{.}\CodeTypeTok{author})\OperatorTok{;} \CommentTok{// "George Orwell"}
 
\CommentTok{// Filter available books}
\KeywordTok{const} available \OperatorTok{=} library\OperatorTok{.}\CodeTypeTok{books}\OperatorTok{.}\CodeFuncTok{filter}((book) \KeywordTok{=\textgreater{}} book\OperatorTok{.}\CodeTypeTok{available})\OperatorTok{;}
console\OperatorTok{.}\CodeFuncTok{log}(\StringTok{\textasciigrave{}Available: }\StringTok{\$\{}available\OperatorTok{.}\CodeFuncTok{map}((b) \KeywordTok{=\textgreater{}} b\OperatorTok{.}\CodeTypeTok{title})\OperatorTok{.}\CodeFuncTok{join}(\StringTok{", "})\StringTok{\}}\StringTok{\textasciigrave{}})\OperatorTok{;}
\end{Highlighting}
\end{Shaded}

\subsection{Destructuring}\label{destructuring}

\begin{Shaded}
\begin{Highlighting}[]
\CommentTok{// Object destructuring}
\KeywordTok{const} \{ name\OperatorTok{,} age \} \OperatorTok{=} studentProfile\OperatorTok{;}
console\OperatorTok{.}\CodeFuncTok{log}(name\OperatorTok{,} age)\OperatorTok{;} \CommentTok{// "Emily" 13}
 
\CommentTok{// With renaming}
\KeywordTok{const} \{ \CodeTypeTok{name}\OperatorTok{:} studentName\OperatorTok{,} \CodeTypeTok{favoriteSubject}\OperatorTok{:} subject \} \OperatorTok{=} studentProfile\OperatorTok{;}
 
\CommentTok{// Array destructuring}
\KeywordTok{const} [first\OperatorTok{,} second\OperatorTok{,} \OperatorTok{...}rest] \OperatorTok{=} [\CodeNumTok{1}\OperatorTok{,} \CodeNumTok{2}\OperatorTok{,} \CodeNumTok{3}\OperatorTok{,} \CodeNumTok{4}\OperatorTok{,} \CodeNumTok{5}]\OperatorTok{;}
console\OperatorTok{.}\CodeFuncTok{log}(first)\OperatorTok{;} \CommentTok{// 1}
console\OperatorTok{.}\CodeFuncTok{log}(rest)\OperatorTok{;} \CommentTok{// [3, 4, 5]}
 
\CommentTok{// In function parameters}
\KeywordTok{function} \CodeFuncTok{printStudent}(\{ name\OperatorTok{,} age\OperatorTok{,} grade ) \{
  console\OperatorTok{.}\CodeFuncTok{log}(\StringTok{\textasciigrave{}}\StringTok{\$\{}name\StringTok{\}}\StringTok{ | Age: }\StringTok{\$\{}age\StringTok{\}}\StringTok{ | Grade: }\StringTok{\$\{}grade\StringTok{\}}\StringTok{\textasciigrave{}})\OperatorTok{;}
\}
\CodeFuncTok{printStudent}(studentProfile)\OperatorTok{;}
\end{Highlighting}
\end{Shaded}

\subsection{Spread and Rest
Operators}\label{spread-and-rest-operators}

\begin{Shaded}
\begin{Highlighting}[]
\CommentTok{// Spread {-}{-}{-} expand array or object}
\KeywordTok{const} arr1 \OperatorTok{=} [\CodeNumTok{1}\OperatorTok{,} \CodeNumTok{2}\OperatorTok{,} \CodeNumTok{3}]\OperatorTok{;}
\KeywordTok{const} arr2 \OperatorTok{=} [\CodeNumTok{4}\OperatorTok{,} \CodeNumTok{5}\OperatorTok{,} \CodeNumTok{6}]\OperatorTok{;}
\KeywordTok{const} combined \OperatorTok{=} [\OperatorTok{...}arr1\OperatorTok{,} \OperatorTok{...}arr2]\OperatorTok{;} \CommentTok{// [1, 2, 3, 4, 5, 6]}
 
\KeywordTok{const} obj1 \OperatorTok{=} \{ \}\CodeTypeTok{a}\OperatorTok{:} \CodeNumTok{1}\OperatorTok{,} \CodeTypeTok{b}\OperatorTok{:} \CodeNumTok{2} \}\OperatorTok{;}
\KeywordTok{const} obj2 \OperatorTok{=} \{ \CodeTypeTok{c}\OperatorTok{:} \CodeNumTok{3}\OperatorTok{,} \CodeTypeTok{d}\OperatorTok{:} \CodeNumTok{4} \}\OperatorTok{;}
\KeywordTok{const} merged \OperatorTok{=} \{ \OperatorTok{...}obj1\OperatorTok{,} \OperatorTok{...}obj2 \}\OperatorTok{;} \CommentTok{// \{ a: 1, b: 2, c: 3, d: 4 \}}
 
\CommentTok{// Copy without mutating}
\KeywordTok{const} originalArr \OperatorTok{=} [\CodeNumTok{1}\OperatorTok{,} \CodeNumTok{2}\OperatorTok{,} \CodeNumTok{3}]\OperatorTok{;}
\KeywordTok{const} copyArr \OperatorTok{=} [\OperatorTok{...}originalArr]\OperatorTok{;}
 
\KeywordTok{const} originalObj \OperatorTok{=} \{ \CodeTypeTok{name}\OperatorTok{:} \StringTok{"Alice"} \}\OperatorTok{;}
\KeywordTok{const} copyObj \OperatorTok{=} \{ \OperatorTok{...}originalObj \}\OperatorTok{;}
\end{Highlighting}
\end{Shaded}

\subsection{Skills Checklist}\label{skills-checklist-61}

\begin{itemize}
\tightlist
\item[$\square$]
  I can create, access, modify, and delete object properties.
\item[$\square$]
  I can iterate over objects using \texttt{for...in} and
  \texttt{Object.entries()}.
\item[$\square$]
  I can work with nested objects and arrays.
\item[$\square$]
  I can use destructuring assignment.
\item[$\square$]
  I can use spread and rest operators.
\end{itemize}

\newpage

\section{Unit 5: Strings and String
Methods}\label{unit-5-strings-and-string-methods}

\begin{center}\rule{0.5\linewidth}{0.5pt}\end{center}

\subsection{Lesson 5.1: Working with
Strings}\label{lesson-5.1-working-with-strings}

\begin{Shaded}
\begin{Highlighting}[]
\KeywordTok{const} name \OperatorTok{=} \StringTok{"Alice Smith"}\OperatorTok{;}
 
\CommentTok{// Properties}
console\OperatorTok{.}\CodeFuncTok{log}(name\OperatorTok{.}\CodeTypeTok{length})\OperatorTok{;} \CommentTok{// 11}
 
\CommentTok{// Accessing characters}
console\OperatorTok{.}\CodeFuncTok{log}(name[\CodeNumTok{0}])\OperatorTok{;} \CommentTok{// "A"}
console\OperatorTok{.}\CodeFuncTok{log}(name\OperatorTok{.}\CodeFuncTok{charAt}(\CodeNumTok{0}))\OperatorTok{;} \CommentTok{// "A"}
console\OperatorTok{.}\CodeFuncTok{log}(name[name\OperatorTok{.}\CodeTypeTok{length} \OperatorTok{{-}} \CodeNumTok{1}])\OperatorTok{;} \CommentTok{// "h"}
 
\CommentTok{// Substrings}
console\OperatorTok{.}\CodeFuncTok{log}(name\OperatorTok{.}\CodeFuncTok{slice}(\CodeNumTok{6}))\OperatorTok{;} \CommentTok{// "Smith"}
console\OperatorTok{.}\CodeFuncTok{log}(name\OperatorTok{.}\CodeFuncTok{slice}(\CodeNumTok{0}\OperatorTok{,} \CodeNumTok{5}))\OperatorTok{;} \CommentTok{// "Alice"}
console\OperatorTok{.}\CodeFuncTok{log}(name\OperatorTok{.}\CodeFuncTok{substring}(\CodeNumTok{6}\OperatorTok{,} \CodeNumTok{11}))\OperatorTok{;} \CommentTok{// "Smith"}
 
\CommentTok{// Case}
console\OperatorTok{.}\CodeFuncTok{log}(name\OperatorTok{.}\CodeFuncTok{toUpperCase}())\OperatorTok{;} \CommentTok{// "ALICE SMITH"}
console\OperatorTok{.}\CodeFuncTok{log}(name\OperatorTok{.}\CodeFuncTok{toLowerCase}())\OperatorTok{;} \CommentTok{// "alice smith"}
 
\CommentTok{// Trimming}
\KeywordTok{const} padded \OperatorTok{=} \StringTok{"  hello  "}\OperatorTok{;}
console\OperatorTok{.}\CodeFuncTok{log}(padded\OperatorTok{.}\CodeFuncTok{trim}())\OperatorTok{;} \CommentTok{// "hello"}
console\OperatorTok{.}\CodeFuncTok{log}(padded\OperatorTok{.}\CodeFuncTok{trimStart}())\OperatorTok{;} \CommentTok{// "hello  "}
console\OperatorTok{.}\CodeFuncTok{log}(padded\OperatorTok{.}\CodeFuncTok{trimEnd}())\OperatorTok{;} \CommentTok{// "  hello"}
 
\CommentTok{// Searching}
console\OperatorTok{.}\CodeFuncTok{log}(name\OperatorTok{.}\CodeFuncTok{indexOf}(\StringTok{"Smith"}))\OperatorTok{;} \CommentTok{// 6}
console\OperatorTok{.}\CodeFuncTok{log}(name\OperatorTok{.}\CodeFuncTok{lastIndexOf}(\StringTok{"l"}))\OperatorTok{;} \CommentTok{// 3}
console\OperatorTok{.}\CodeFuncTok{log}(name\OperatorTok{.}\CodeFuncTok{includes}(\StringTok{"Alice"}))\OperatorTok{;} \CommentTok{// true}
console\OperatorTok{.}\CodeFuncTok{log}(name\OperatorTok{.}\CodeFuncTok{startsWith}(\StringTok{"Ali"}))\OperatorTok{;} \CommentTok{// true}
console\OperatorTok{.}\CodeFuncTok{log}(name\OperatorTok{.}\CodeFuncTok{endsWith}(\StringTok{"ith"}))\OperatorTok{;} \CommentTok{// true}
 
\CommentTok{// Replacing}
\KeywordTok{const} result \OperatorTok{=} name\OperatorTok{.}\CodeFuncTok{replace}(\StringTok{"Alice"}\OperatorTok{,} \StringTok{"Bob"})\OperatorTok{;} \CommentTok{// "Bob Smith"}
\KeywordTok{const} allReplaced \OperatorTok{=} \StringTok{"a,b,a,c,a"}\OperatorTok{.}\CodeFuncTok{replaceAll}(\StringTok{"a"}\OperatorTok{,} \StringTok{"x"})\OperatorTok{;} \CommentTok{// "x,b,x,c,x"}
 
\CommentTok{// Splitting}
\KeywordTok{const} sentence \OperatorTok{=} \StringTok{"the quick brown fox"}\OperatorTok{;}
\KeywordTok{const} words \OperatorTok{=} sentence\OperatorTok{.}\CodeFuncTok{split}(\StringTok{" "})\OperatorTok{;}
console\OperatorTok{.}\CodeFuncTok{log}(words)\OperatorTok{;} \CommentTok{// ["the", "quick", "brown", "fox"]}
 
\KeywordTok{const} csv \OperatorTok{=} \StringTok{"Alice,15,Portland"}\OperatorTok{;}
\KeywordTok{const} parts \OperatorTok{=} csv\OperatorTok{.}\CodeFuncTok{split}(\StringTok{","})\OperatorTok{;}
console\OperatorTok{.}\CodeFuncTok{log}(parts)\OperatorTok{;} \CommentTok{// ["Alice", "15", "Portland"]}
 
\CommentTok{// Repeating and Padding}
console\OperatorTok{.}\CodeFuncTok{log}(\StringTok{"ha"}\OperatorTok{.}\CodeFuncTok{repeat}(\CodeNumTok{3}))\OperatorTok{;} \CommentTok{// "hahaha"}
console\OperatorTok{.}\CodeFuncTok{log}(\StringTok{"5"}\OperatorTok{.}\CodeFuncTok{padStart}(\CodeNumTok{3}\OperatorTok{,} \StringTok{"0"}))\OperatorTok{;} \CommentTok{// "005"}
console\OperatorTok{.}\CodeFuncTok{log}(\StringTok{"5"}\OperatorTok{.}\CodeFuncTok{padEnd}(\CodeNumTok{3}\OperatorTok{,} \StringTok{"0"}))\OperatorTok{;} \CommentTok{// "500"}
 
\CommentTok{// Template literals (multi{-}line)}
\KeywordTok{const} multiLine \OperatorTok{=} \StringTok{\textasciigrave{}}
\StringTok{Name: }\StringTok{\$\{}name\StringTok{\}}
\StringTok{Length: }\StringTok{\$\{}name\OperatorTok{.}\CodeTypeTok{length}\StringTok{\}}
\StringTok{Uppercase: }\StringTok{\$\{}name\OperatorTok{.}\CodeFuncTok{toUpperCase}()\StringTok{\}}
\StringTok{\textasciigrave{}}\OperatorTok{;}
console\OperatorTok{.}\CodeFuncTok{log}(multiLine)\OperatorTok{;}
\end{Highlighting}
\end{Shaded}

\subsection{Regular Expressions
(Basics)}\label{regular-expressions-basics}

\begin{Shaded}
\begin{Highlighting}[]
\KeywordTok{const} text \OperatorTok{=} \StringTok{"Hello World! Hello JavaScript!"}\OperatorTok{;}
 
\CommentTok{// Test if pattern exists}
console\OperatorTok{.}\CodeFuncTok{log}(\StringTok{/hello/i}\OperatorTok{.}\CodeFuncTok{test}(text))\OperatorTok{;} \CommentTok{// true (i = case{-}insensitive)}
 
\CommentTok{// Find matches}
\KeywordTok{const} matches \OperatorTok{=} text\OperatorTok{.}\CodeFuncTok{match}(\StringTok{/Hello/g})\OperatorTok{;} \CommentTok{// g = global (find all)}
console\OperatorTok{.}\CodeFuncTok{log}(matches)\OperatorTok{;} \CommentTok{// ["Hello", "Hello"]}
 
\CommentTok{// Replace with regex}
\KeywordTok{const} cleaned \OperatorTok{=} text\OperatorTok{.}\CodeFuncTok{replace}(\StringTok{/}\StringTok{[!]}\StringTok{/g}\OperatorTok{,} \StringTok{"."})\OperatorTok{;} \CommentTok{// "Hello World. Hello JavaScript."}
 
\CommentTok{// Split with regex}
\KeywordTok{const} words2 \OperatorTok{=} \StringTok{"one  two   three"}\OperatorTok{.}\CodeFuncTok{split}(\StringTok{/}\StringTok{\textbackslash{}s+}\StringTok{/})\OperatorTok{;}
console\OperatorTok{.}\CodeFuncTok{log}(words2)\OperatorTok{;} \CommentTok{// ["one", "two", "three"]}
\end{Highlighting}
\end{Shaded}

\subsection{Skills Checklist}\label{skills-checklist-62}

\begin{itemize}
\tightlist
\item[$\square$]
  I can use common \texttt{String} methods in JavaScript.
\item[$\square$]
  I can use template literals for multi-line strings and interpolation.
\item[$\square$]
  I can split strings into arrays and join arrays into strings.
\item[$\square$]
  I understand basic regular expressions.
\end{itemize}

\begin{center}\rule{0.5\linewidth}{0.5pt}\end{center}

\section{Unit 6: Functions}\label{unit-6-functions-3}

\begin{center}\rule{0.5\linewidth}{0.5pt}\end{center}

\subsection{Lesson 6.1: Defining
Functions}\label{lesson-6.1-defining-functions-1}

JavaScript has multiple ways to define functions:

\subsection{Function Declaration}\label{function-declaration}\amn{\textbf{Accessibility Tip:} Write a comment above every function describing what it does, its parameters, and its return value. When IntelliSense shows the function in a completion list, it displays this comment as documentation that is read aloud by the screen reader.}
\subsection{Function Declaration}\label{function-declaration}\amn{\textbf{Accessibility Tip:} Write a JSDoc comment above every function: \texttt{/** @param \{number\} x - The input. @returns \{number\} The result. */}. IntelliSense reads this aloud when you type the function name, giving you parameter and return type information without having to navigate to the definition.}
\subsection{Function Declaration}\label{function-declaration}

\begin{Shaded}
\begin{Highlighting}[]
\KeywordTok{function} \CodeFuncTok{greetUser}() \{
  console\OperatorTok{.}\CodeFuncTok{log}(\StringTok{"Hello, JavaScript programmer!"})\OperatorTok{;}
\}
 
\CodeFuncTok{greetUser}()\OperatorTok{;}
\CodeFuncTok{greetUser}()\OperatorTok{;}
\end{Highlighting}
\end{Shaded}

\subsection{Function Expression}\label{function-expression}

\begin{Shaded}
\begin{Highlighting}[]
\KeywordTok{const} greetByName \OperatorTok{=} \KeywordTok{function} (name) \{
  console\OperatorTok{.}\CodeFuncTok{log}(\StringTok{\textasciigrave{}Hello, }\StringTok{\$\{}name\StringTok{\}}\StringTok{!\textasciigrave{}})\OperatorTok{;}
\}\OperatorTok{;}
 
\CodeFuncTok{greetByName}(\StringTok{"Alice"})\OperatorTok{;}
\end{Highlighting}
\end{Shaded}

\subsection{Arrow Functions (ES6) --- Modern Preferred
Style}\label{arrow-functions-es6-modern-preferred-style}

\begin{Shaded}
\begin{Highlighting}[]
\CommentTok{// Full arrow function}
\KeywordTok{const} add \OperatorTok{=} (a\OperatorTok{,} b) \KeywordTok{=\textgreater{}} \{
  \KeywordTok{return} a \OperatorTok{+} b\OperatorTok{;}
\}\OperatorTok{;}
 
\CommentTok{// Concise body (implicit return for single expressions)}
\KeywordTok{const} addShort \OperatorTok{=} (a\OperatorTok{,} b) \KeywordTok{=\textgreater{}} a \OperatorTok{+} b\OperatorTok{;}
 
\CommentTok{// Single parameter (no parentheses needed)}
\KeywordTok{const} double \OperatorTok{=} (n) \KeywordTok{=\textgreater{}} n \OperatorTok{*} \CodeNumTok{2}\OperatorTok{;}
 
\CommentTok{// No parameters}
\KeywordTok{const} sayHello \OperatorTok{=} () \KeywordTok{=\textgreater{}} console\OperatorTok{.}\CodeFuncTok{log}(\StringTok{"Hello!"})\OperatorTok{;}
 
console\OperatorTok{.}\CodeFuncTok{log}(\CodeFuncTok{add}(\CodeNumTok{5}\OperatorTok{,} \CodeNumTok{3}))\OperatorTok{;} \CommentTok{// 8}
console\OperatorTok{.}\CodeFuncTok{log}(\CodeFuncTok{addShort}(\CodeNumTok{5}\OperatorTok{,} \CodeNumTok{3}))\OperatorTok{;} \CommentTok{// 8}
console\OperatorTok{.}\CodeFuncTok{log}(\CodeFuncTok{double}(\CodeNumTok{7}))\OperatorTok{;} \CommentTok{// 14}
\CodeFuncTok{sayHello}()\OperatorTok{;}
\end{Highlighting}
\end{Shaded}

\subsection{Default Parameters}\label{default-parameters}

\begin{Shaded}
\begin{Highlighting}[]
\KeywordTok{function} \CodeFuncTok{makeShirt}(size\OperatorTok{,} message \OperatorTok{=} \StringTok{"I love JavaScript"}) \{
  console\OperatorTok{.}\CodeFuncTok{log}(\StringTok{\textasciigrave{}Making a }\StringTok{\$\{}size\StringTok{\}}\StringTok{{-}sized shirt: \textquotesingle{}}\StringTok{\$\{}message\StringTok{\}}\StringTok{\textquotesingle{}\textasciigrave{}})\OperatorTok{;}
\}
 
\CodeFuncTok{makeShirt}(\StringTok{"large"})\OperatorTok{;} \CommentTok{// Uses default}
\CodeFuncTok{makeShirt}(\StringTok{"medium"}\OperatorTok{,} \StringTok{"Hello World"})\OperatorTok{;} \CommentTok{// Custom message}
\end{Highlighting}
\end{Shaded}

\subsection{Rest Parameters}\label{rest-parameters}

\begin{Shaded}
\begin{Highlighting}[]
\KeywordTok{function} \CodeFuncTok{sum}(\OperatorTok{...}numbers) \{
  \KeywordTok{return} numbers\OperatorTok{.}\CodeFuncTok{reduce}((total\OperatorTok{,} n) \KeywordTok{=\textgreater{}} total \OperatorTok{+} n\OperatorTok{,} \CodeNumTok{0})\OperatorTok{;}
\}
 
console\OperatorTok{.}\CodeFuncTok{log}(\CodeFuncTok{sum}(\CodeNumTok{1}\OperatorTok{,} \CodeNumTok{2}\OperatorTok{,} \CodeNumTok{3}))\OperatorTok{;} \CommentTok{// 6}
console\OperatorTok{.}\CodeFuncTok{log}(\CodeFuncTok{sum}(\CodeNumTok{10}\OperatorTok{,} \CodeNumTok{20}\OperatorTok{,} \CodeNumTok{30}\OperatorTok{,} \CodeNumTok{40}))\OperatorTok{;} \CommentTok{// 100}
\end{Highlighting}
\end{Shaded}

\subsection{Examples}\label{examples-19}

\textbf{Example 6.1.1: Calculator Functions}

\begin{Shaded}
\begin{Highlighting}[]
\KeywordTok{const} add \OperatorTok{=} (a\OperatorTok{,} b) \KeywordTok{=\textgreater{}} a \OperatorTok{+} b\OperatorTok{;}
\KeywordTok{const} subtract \OperatorTok{=} (a\OperatorTok{,} b) \KeywordTok{=\textgreater{}} a \OperatorTok{{-}} b\OperatorTok{;}
\KeywordTok{const} multiply \OperatorTok{=} (a\OperatorTok{,} b) \KeywordTok{=\textgreater{}} a \OperatorTok{*} b\OperatorTok{;}
\KeywordTok{const} divide \OperatorTok{=} (a\OperatorTok{,} b) \KeywordTok{=\textgreater{}} \{
  \KeywordTok{if} (b \OperatorTok{===} \CodeNumTok{0}) \KeywordTok{throw} \KeywordTok{new} Error(\StringTok{"Cannot divide by zero"})\OperatorTok{;}
  \KeywordTok{return} a \OperatorTok{/} b\OperatorTok{;}
\}\OperatorTok{;}
 
\KeywordTok{const} circleArea \OperatorTok{=} (radius) \KeywordTok{=\textgreater{}} Math\OperatorTok{.}\CodeNumTok{PI} \OperatorTok{*} radius \OperatorTok{**} \CodeNumTok{2}\OperatorTok{;}
\KeywordTok{const} getGrade \OperatorTok{=} (score) \KeywordTok{=\textgreater{}} \{
  \KeywordTok{if} (score \OperatorTok{\textgreater{}=} \CodeNumTok{90}) \KeywordTok{return} \StringTok{"A"}\OperatorTok{;}
  \KeywordTok{if} (score \OperatorTok{\textgreater{}=} \CodeNumTok{80}) \KeywordTok{return} \StringTok{"B"}\OperatorTok{;}
  \KeywordTok{if} (score \OperatorTok{\textgreater{}=} \CodeNumTok{70}) \KeywordTok{return} \StringTok{"C"}\OperatorTok{;}
  \KeywordTok{if} (score \OperatorTok{\textgreater{}=} \CodeNumTok{60}) \KeywordTok{return} \StringTok{"D"}\OperatorTok{;}
  \KeywordTok{return} \StringTok{"F"}\OperatorTok{;}
\}\OperatorTok{;}
 
console\OperatorTok{.}\CodeFuncTok{log}(\CodeFuncTok{add}(\CodeNumTok{5}\OperatorTok{,} \CodeNumTok{3}))\OperatorTok{;}
console\OperatorTok{.}\CodeFuncTok{log}(\CodeFuncTok{circleArea}(\CodeNumTok{5})\OperatorTok{.}\CodeFuncTok{toFixed}(\CodeNumTok{2}))\OperatorTok{;}
console\OperatorTok{.}\CodeFuncTok{log}(\CodeFuncTok{getGrade}(\CodeNumTok{85}))\OperatorTok{;}
\end{Highlighting}
\end{Shaded}

\textbf{Example 6.1.2: Higher-Order Functions}

\begin{Shaded}
\begin{Highlighting}[]
\CommentTok{// Functions that take or return other functions}
 
\CommentTok{// Map using our own function}
\KeywordTok{function} \CodeFuncTok{doubleAll}(numbers) \{
  \KeywordTok{return} numbers\OperatorTok{.}\CodeFuncTok{map}((n) \KeywordTok{=\textgreater{}} n \OperatorTok{*} \CodeNumTok{2})\OperatorTok{;}
\}
 
\CommentTok{// Function that returns a function (closure)}
\KeywordTok{function} \CodeFuncTok{makeAdder}(x) \{
  \KeywordTok{return} (y) \KeywordTok{=\textgreater{}} x \OperatorTok{+} y\OperatorTok{;}
\}
 
\KeywordTok{const} addFive \OperatorTok{=} \CodeFuncTok{makeAdder}(\CodeNumTok{5})\OperatorTok{;}
\KeywordTok{const} addTen \OperatorTok{=} \CodeFuncTok{makeAdder}(\CodeNumTok{10})\OperatorTok{;}
 
console\OperatorTok{.}\CodeFuncTok{log}(\CodeFuncTok{addFive}(\CodeNumTok{3}))\OperatorTok{;} \CommentTok{// 8}
console\OperatorTok{.}\CodeFuncTok{log}(\CodeFuncTok{addTen}(\CodeNumTok{3}))\OperatorTok{;} \CommentTok{// 13}
console\OperatorTok{.}\CodeFuncTok{log}(\CodeFuncTok{doubleAll}([\CodeNumTok{1}\OperatorTok{,} \CodeNumTok{2}\OperatorTok{,} \CodeNumTok{3}\OperatorTok{,} \CodeNumTok{4}\OperatorTok{,} \CodeNumTok{5}]))\OperatorTok{;}
\end{Highlighting}
\end{Shaded}

\subsection{Closures}\label{closures}

\begin{Shaded}
\begin{Highlighting}[]
\KeywordTok{function} \CodeFuncTok{makeCounter}() \{
  \KeywordTok{let} count \OperatorTok{=} \CodeNumTok{0}\OperatorTok{;} \CommentTok{// Private variable}
 
  \KeywordTok{return} \{
    \CodeTypeTok{increment}\OperatorTok{:} () \KeywordTok{=\textgreater{}} \OperatorTok{++}count\OperatorTok{,}
    \CodeTypeTok{decrement}\OperatorTok{:} () \KeywordTok{=\textgreater{}} \OperatorTok{{-}{-}}count\OperatorTok{,}
    \CodeTypeTok{getCount}\OperatorTok{:} () \KeywordTok{=\textgreater{}} count\OperatorTok{,}
  \}\OperatorTok{;}
\}
 
\KeywordTok{const} counter \OperatorTok{=} \CodeFuncTok{makeCounter}()\OperatorTok{;}
counter\OperatorTok{.}\CodeFuncTok{increment}()\OperatorTok{;}
counter\OperatorTok{.}\CodeFuncTok{increment}()\OperatorTok{;}
counter\OperatorTok{.}\CodeFuncTok{increment}()\OperatorTok{;}
counter\OperatorTok{.}\CodeFuncTok{decrement}()\OperatorTok{;}
console\OperatorTok{.}\CodeFuncTok{log}(counter\OperatorTok{.}\CodeFuncTok{getCount}())\OperatorTok{;} \CommentTok{// 2}
\end{Highlighting}
\end{Shaded}

\subsection{Skills Checklist}\label{skills-checklist-63}

\begin{itemize}
\tightlist
\item[$\square$]
  I can define functions using declarations, expressions, and arrow
  syntax.
\item[$\square$]
  I can use default and rest parameters.
\item[$\square$]
  I understand closures and how they capture variables.
\item[$\square$]
  I can pass functions as arguments and return them.
\end{itemize}

\begin{center}\rule{0.5\linewidth}{0.5pt}\end{center}

\section{Unit 7: Classes and Object-Oriented
Programming}\label{unit-7-classes-and-object-oriented-programming-2}

\begin{center}\rule{0.5\linewidth}{0.5pt}\end{center}

\amn{\textbf{Accessibility Tip:} Override \texttt{toString()} on every class you create. When \texttt{console.log(myObject)} is called, JavaScript uses \texttt{toString()} for the output. Without it you hear \texttt{[object Object]}. Return a descriptive string: \texttt{return `Dog[name=\$\{this.name\}, age=\$\{this.age\}]`;}}
\subsection{Lesson 7.1: Classes in
JavaScript}\label{lesson-7.1-classes-in-javascript}

JavaScript uses prototype-based inheritance, but ES6 \texttt{class}
syntax provides a familiar OOP style.

\subsection{Class Syntax}\label{class-syntax-1}

\begin{Shaded}
\begin{Highlighting}[]
\KeywordTok{class} Dog \{
  \CommentTok{// Instance variables declared in constructor}
  \CodeFuncTok{constructor}(name\OperatorTok{,} age) \{
    \KeywordTok{this}\OperatorTok{.}\CodeTypeTok{name} \OperatorTok{=} name\OperatorTok{;}
    \KeywordTok{this}\OperatorTok{.}\CodeTypeTok{age} \OperatorTok{=} age\OperatorTok{;}
  \}
 
  \CommentTok{// Methods}
  \CodeFuncTok{sit}() \{
    console\OperatorTok{.}\CodeFuncTok{log}(\StringTok{\textasciigrave{}}\StringTok{\$\{}\KeywordTok{this}\OperatorTok{.}\CodeTypeTok{name}\StringTok{\}}\StringTok{ is now sitting.\textasciigrave{}})\OperatorTok{;}
  \}
 
  \CodeFuncTok{rollOver}() \{
    console\OperatorTok{.}\CodeFuncTok{log}(\StringTok{\textasciigrave{}}\StringTok{\$\{}\KeywordTok{this}\OperatorTok{.}\CodeTypeTok{name}\StringTok{\}}\StringTok{ rolled over!\textasciigrave{}})\OperatorTok{;}
  \}
 
  \CodeFuncTok{getOlder}(years) \{
    \KeywordTok{this}\OperatorTok{.}\CodeTypeTok{age} \OperatorTok{+=} years\OperatorTok{;}
    console\OperatorTok{.}\CodeFuncTok{log}(\StringTok{\textasciigrave{}}\StringTok{\$\{}\KeywordTok{this}\OperatorTok{.}\CodeTypeTok{name}\StringTok{\}}\StringTok{ is now }\StringTok{\$\{}\KeywordTok{this}\OperatorTok{.}\CodeTypeTok{age}\StringTok{\}}\StringTok{ years old.\textasciigrave{}})\OperatorTok{;}
  \}
 
  \CommentTok{// Getter}
  \KeywordTok{get} \CodeFuncTok{info}() \{
    \KeywordTok{return} \StringTok{\textasciigrave{}}\StringTok{\$\{}\KeywordTok{this}\OperatorTok{.}\CodeTypeTok{name}\StringTok{\}}\StringTok{ (}\StringTok{\$\{}\KeywordTok{this}\OperatorTok{.}\CodeTypeTok{age}\StringTok{\}}\StringTok{ years old)\textasciigrave{}}\OperatorTok{;}
  \}
 
  \CommentTok{// toString equivalent}
  \CodeFuncTok{toString}() \{
    \KeywordTok{return} \StringTok{\textasciigrave{}Dog \{ name: \textquotesingle{}}\StringTok{\$\{}\KeywordTok{this}\OperatorTok{.}\CodeTypeTok{name}\StringTok{\}}\StringTok{\textquotesingle{}, age: }\StringTok{\$\{}\KeywordTok{this}\OperatorTok{.}\CodeTypeTok{age}\StringTok{\}}\StringTok{ \}\textasciigrave{}}\OperatorTok{;}
  \}
\}
 
\CommentTok{// Create instances}
\KeywordTok{const} myDog \OperatorTok{=} \KeywordTok{new} \CodeFuncTok{Dog}(\StringTok{"Buddy"}\OperatorTok{,} \CodeNumTok{3})\OperatorTok{;}
\KeywordTok{const} yourDog \OperatorTok{=} \KeywordTok{new} \CodeFuncTok{Dog}(\StringTok{"Lucy"}\OperatorTok{,} \CodeNumTok{1})\OperatorTok{;}
 
myDog\OperatorTok{.}\CodeFuncTok{sit}()\OperatorTok{;}
myDog\OperatorTok{.}\CodeFuncTok{getOlder}(\CodeNumTok{2})\OperatorTok{;}
console\OperatorTok{.}\CodeFuncTok{log}(myDog\OperatorTok{.}\CodeTypeTok{info})\OperatorTok{;} \CommentTok{// Uses getter}
console\OperatorTok{.}\CodeFuncTok{log}(myDog\OperatorTok{.}\CodeFuncTok{toString}())\OperatorTok{;}
\end{Highlighting}
\end{Shaded}

\subsection{Private Fields (ES2022)}\label{private-fields-es2022}

\begin{Shaded}
\begin{Highlighting}[]
\KeywordTok{class} BankAccount \{
  \#balance\OperatorTok{;} \CommentTok{// Private field {-}{-}{-} only accessible within the class}
 
  \CodeFuncTok{constructor}(initialBalance) \{
    \KeywordTok{this}\OperatorTok{.}\#\CodeTypeTok{balance} \OperatorTok{=} initialBalance\OperatorTok{;}
  \}
 
  \CodeFuncTok{deposit}(amount) \{
    \KeywordTok{if} (amount \OperatorTok{\textgreater{}} \CodeNumTok{0}) \{
      \KeywordTok{this}\OperatorTok{.}\#\CodeTypeTok{balance} \OperatorTok{+=} amount\OperatorTok{;}
      console\OperatorTok{.}\CodeFuncTok{log}(\StringTok{\textasciigrave{}Deposited $}\StringTok{\$\{}amount\StringTok{\}}\StringTok{. New balance: $}\StringTok{\$\{}\KeywordTok{this}\OperatorTok{.}\#\CodeTypeTok{balance}\StringTok{\}}\StringTok{\textasciigrave{}})\OperatorTok{;}
    \}
  \}
 
  \CodeFuncTok{withdraw}(amount) \{
    \KeywordTok{if} (amount \OperatorTok{\textgreater{}} \KeywordTok{this}\OperatorTok{.}\#\CodeTypeTok{balance}) \{
      console\OperatorTok{.}\CodeFuncTok{log}(\StringTok{"Insufficient funds!"})\OperatorTok{;}
    \} \KeywordTok{else} \{
      \KeywordTok{this}\OperatorTok{.}\#\CodeTypeTok{balance} \OperatorTok{{-}=} amount\OperatorTok{;}
      console\OperatorTok{.}\CodeFuncTok{log}(\StringTok{\textasciigrave{}Withdrew $}\StringTok{\$\{}amount\StringTok{\}}\StringTok{. New balance: $}\StringTok{\$\{}\KeywordTok{this}\OperatorTok{.}\#\CodeTypeTok{balance}\StringTok{\}}\StringTok{\textasciigrave{}})\OperatorTok{;}
    \}
  \}
 
  \KeywordTok{get} \CodeFuncTok{balance}() \{
    \KeywordTok{return} \KeywordTok{this}\OperatorTok{.}\#\CodeTypeTok{balance}\OperatorTok{;}
  \}
\}
 
\KeywordTok{const} account \OperatorTok{=} \KeywordTok{new} \CodeFuncTok{BankAccount}(\CodeNumTok{1000})\OperatorTok{;}
account\OperatorTok{.}\CodeFuncTok{deposit}(\CodeNumTok{500})\OperatorTok{;}
account\OperatorTok{.}\CodeFuncTok{withdraw}(\CodeNumTok{200})\OperatorTok{;}
console\OperatorTok{.}\CodeFuncTok{log}(\StringTok{\textasciigrave{}Balance: $}\StringTok{\$\{}account\OperatorTok{.}\CodeTypeTok{balance}\StringTok{\}}\StringTok{\textasciigrave{}})\OperatorTok{;}
\end{Highlighting}
\end{Shaded}

\subsection{Static Methods}\label{static-methods}

\begin{Shaded}
\begin{Highlighting}[]
\KeywordTok{class} MathHelper \{
  \KeywordTok{static} \CodeFuncTok{add}(a\OperatorTok{,} b) \{
    \KeywordTok{return} a \OperatorTok{+} b\OperatorTok{;}
  \}
  \KeywordTok{static} \CodeFuncTok{multiply}(a\OperatorTok{,} b) \{
    \KeywordTok{return} a \OperatorTok{*} b\OperatorTok{;}
  \}
  \KeywordTok{static} \CodeNumTok{PI} \OperatorTok{=} \CodeNumTok{3.14159265358979}\OperatorTok{;}
\}
 
\CommentTok{// Call without creating an instance}
console\OperatorTok{.}\CodeFuncTok{log}(MathHelper\OperatorTok{.}\CodeFuncTok{add}(\CodeNumTok{5}\OperatorTok{,} \CodeNumTok{3}))\OperatorTok{;}
console\OperatorTok{.}\CodeFuncTok{log}(MathHelper\OperatorTok{.}\CodeTypeTok{PI})\OperatorTok{;}
\end{Highlighting}
\end{Shaded}

\begin{center}\rule{0.5\linewidth}{0.5pt}\end{center}

\subsection{Lesson 7.2: Inheritance}\label{lesson-7.2-inheritance}

\begin{Shaded}
\begin{Highlighting}[]
\KeywordTok{class} Vehicle \{
  \CodeFuncTok{constructor}(make\OperatorTok{,} model\OperatorTok{,} year) \{
    \KeywordTok{this}\OperatorTok{.}\CodeTypeTok{make} \OperatorTok{=} make\OperatorTok{;}
    \KeywordTok{this}\OperatorTok{.}\CodeTypeTok{model} \OperatorTok{=} model\OperatorTok{;}
    \KeywordTok{this}\OperatorTok{.}\CodeTypeTok{year} \OperatorTok{=} year\OperatorTok{;}
    \KeywordTok{this}\OperatorTok{.}\CodeTypeTok{odometer} \OperatorTok{=} \CodeNumTok{0}\OperatorTok{;}
  \}
 
  \CodeFuncTok{getDescriptiveName}() \{
    \KeywordTok{return} \StringTok{\textasciigrave{}}\StringTok{\$\{}\KeywordTok{this}\OperatorTok{.}\CodeTypeTok{year}\StringTok{\}}\StringTok{ }\StringTok{\$\{}\KeywordTok{this}\OperatorTok{.}\CodeTypeTok{make}\StringTok{\}}\StringTok{ }\StringTok{\$\{}\KeywordTok{this}\OperatorTok{.}\CodeTypeTok{model}\StringTok{\}}\StringTok{\textasciigrave{}}\OperatorTok{;}
  \}
 
  \CodeFuncTok{readOdometer}() \{
    console\OperatorTok{.}\CodeFuncTok{log}(\StringTok{\textasciigrave{}This vehicle has }\StringTok{\$\{}\KeywordTok{this}\OperatorTok{.}\CodeTypeTok{odometer}\StringTok{\}}\StringTok{ miles on it.\textasciigrave{}})\OperatorTok{;}
  \}
\}
 
\KeywordTok{class} Car \KeywordTok{extends} Vehicle \{
  \CodeFuncTok{constructor}(make\OperatorTok{,} model\OperatorTok{,} year) \{
    \KeywordTok{super}(make\OperatorTok{,} model\OperatorTok{,} year)\OperatorTok{;} \CommentTok{// Call parent constructor}
    \KeywordTok{this}\OperatorTok{.}\CodeTypeTok{numDoors} \OperatorTok{=} \CodeNumTok{4}\OperatorTok{;}
  \}
 
  \CodeFuncTok{readOdometer}() \{
    \CommentTok{// Override parent method}
    console\OperatorTok{.}\CodeFuncTok{log}(\StringTok{\textasciigrave{}This car has }\StringTok{\$\{}\KeywordTok{this}\OperatorTok{.}\CodeTypeTok{odometer}\StringTok{\}}\StringTok{ miles on it.\textasciigrave{}})\OperatorTok{;}
  \}
\}
 
\KeywordTok{class} ElectricCar \KeywordTok{extends} Car \{
  \#batterySize\OperatorTok{;}
 
  \CodeFuncTok{constructor}(make\OperatorTok{,} model\OperatorTok{,} year) \{
    \KeywordTok{super}(make\OperatorTok{,} model\OperatorTok{,} year)\OperatorTok{;}
    \KeywordTok{this}\OperatorTok{.}\#\CodeTypeTok{batterySize} \OperatorTok{=} \CodeNumTok{75}\OperatorTok{;}
  \}
 
  \CodeFuncTok{describeBattery}() \{
    console\OperatorTok{.}\CodeFuncTok{log}(\StringTok{\textasciigrave{}This car has a }\StringTok{\$\{}\KeywordTok{this}\OperatorTok{.}\#\CodeTypeTok{batterySize}\StringTok{\}}\StringTok{{-}kWh battery.\textasciigrave{}})\OperatorTok{;}
  \}
\}
 
\KeywordTok{const} myEV \OperatorTok{=} \KeywordTok{new} \CodeFuncTok{ElectricCar}(\StringTok{"Ford"}\OperatorTok{,} \StringTok{"Mustang Mach{-}E"}\OperatorTok{,} \CodeNumTok{2023})\OperatorTok{;}
console\OperatorTok{.}\CodeFuncTok{log}(myEV\OperatorTok{.}\CodeFuncTok{getDescriptiveName}())\OperatorTok{;} \CommentTok{// Inherited from Vehicle}
myEV\OperatorTok{.}\CodeFuncTok{readOdometer}()\OperatorTok{;} \CommentTok{// Overridden in Car}
myEV\OperatorTok{.}\CodeFuncTok{describeBattery}()\OperatorTok{;} \CommentTok{// Defined in ElectricCar}
\end{Highlighting}
\end{Shaded}

\subsection{Skills Checklist}\label{skills-checklist-64}

\begin{itemize}
\tightlist
\item[$\square$]
  I can define a JavaScript class with a constructor and methods.
\item[$\square$]
  I can use \texttt{extends} to create a subclass.
\item[$\square$]
  I can use \texttt{super()} to call the parent constructor.
\item[$\square$]
  I can use \texttt{\#} for private fields.
\item[$\square$]
  I can define static methods and properties.
\end{itemize}

\begin{center}\rule{0.5\linewidth}{0.5pt}\end{center}

\section{Unit 8: Files, Modules, and Error
Handling}\label{unit-8-files-modules-and-error-handling-1}

\begin{center}\rule{0.5\linewidth}{0.5pt}\end{center}

\subsection{Lesson 8.1: Reading and Writing
Files}\label{lesson-8.1-reading-and-writing-files}
\amn{\textbf{Accessibility Tip:} Wrap all file and async operations in \texttt{try}/\texttt{catch} and always log the error: \texttt{catch (e) \{ console.log("Error:", e.message); \}}. Without this, failed operations are silently swallowed---the program continues with missing data and no audio signal that anything went wrong.}

\begin{Shaded}
\begin{Highlighting}[]
\KeywordTok{const} fs \OperatorTok{=} \CodeAccentTok{require}(\StringTok{"fs"})\OperatorTok{;}
\KeywordTok{const} path \OperatorTok{=} \CodeAccentTok{require}(\StringTok{"path"})\OperatorTok{;}
 
\CommentTok{// Synchronous reading (simple but blocks execution)}
\KeywordTok{try} \{
  \KeywordTok{const} content \OperatorTok{=} fs\OperatorTok{.}\CodeFuncTok{readFileSync}(\StringTok{"my\_notes.txt"}\OperatorTok{,} \StringTok{"utf8"})\OperatorTok{;}
  console\OperatorTok{.}\CodeFuncTok{log}(content)\OperatorTok{;}
\} \KeywordTok{catch} (err) \{
  console\OperatorTok{.}\CodeFuncTok{log}(\StringTok{\textasciigrave{}Error: }\StringTok{\$\{}err\OperatorTok{.}\CodeTypeTok{message}\StringTok{\}}\StringTok{\textasciigrave{}})\OperatorTok{;}
\}
 
\CommentTok{// Reading line by line}
\KeywordTok{const} lines \OperatorTok{=} content\OperatorTok{.}\CodeFuncTok{split}(\StringTok{"}\StringTok{\textbackslash{}n}\StringTok{"})\OperatorTok{;}
lines\OperatorTok{.}\CodeFuncTok{forEach}((line\OperatorTok{,} index) \KeywordTok{=\textgreater{}} \{
  console\OperatorTok{.}\CodeFuncTok{log}(\StringTok{\textasciigrave{}Line }\StringTok{\$\{}index \OperatorTok{+} \CodeNumTok{1}\StringTok{\}}\StringTok{: }\StringTok{\$\{}line\StringTok{\}}\StringTok{\textasciigrave{}})\OperatorTok{;}
)\OperatorTok{;}
 
\CommentTok{// Writing (overwrites existing content)}
\KeywordTok{try} \{
  fs\OperatorTok{.}\CodeFuncTok{writeFileSync}(\StringTok{"output.txt"}\OperatorTok{,} \StringTok{"Hello, Node.js file I/O!}\StringTok{\textbackslash{}n}\StringTok{"})\OperatorTok{;}
  console\OperatorTok{.}\CodeFuncTok{log}(\StringTok{"File written successfully."})\OperatorTok{;}
\ \}\KeywordTok{catch} (err) \{
  console\OperatorTok{.}\CodeFuncTok{log}(\StringTok{\textasciigrave{}Error writing file: }\StringTok{\$\{}err\OperatorTok{.}\CodeTypeTok{message}\StringTok{\}}\StringTok{\textasciigrave{}})\OperatorTok{;}
\}
 
\CommentTok{// Appending}
\KeywordTok{try} \{
  fs\OperatorTok{.}\CodeFuncTok{appendFileSync}(\StringTok{"output.txt"}\OperatorTok{,} \StringTok{"This line is appended.}\StringTok{\textbackslash{}n}\StringTok{"})\OperatorTok{;}
  console\OperatorTok{.}\CodeFuncTok{log}(\StringTok{"Line appended."})\OperatorTok{;}
\ \}\KeywordTok{catch} (err) \{
  console\OperatorTok{.}\CodeFuncTok{log}(\StringTok{\textasciigrave{}Error appending: }\StringTok{\$\{}err\OperatorTok{.}\CodeTypeTok{message}\StringTok{\}}\StringTok{\textasciigrave{}})\OperatorTok{;}
\}
 
\CommentTok{// Check if file exists}
\KeywordTok{if} (fs\OperatorTok{.}\CodeFuncTok{existsSync}(\StringTok{"my\_notes.txt"})) \{\}
  console\OperatorTok{.}\CodeFuncTok{log}(\StringTok{"File exists."})\OperatorTok{;}
\}
\end{Highlighting}
\end{Shaded}

\subsection{Asynchronous File I/O (Modern
Approach)}\label{asynchronous-file-io-modern-approach}

\begin{Shaded}
\begin{Highlighting}[]
\KeywordTok{const} fs \OperatorTok{=} \CodeAccentTok{require}(\StringTok{"fs"})\OperatorTok{.}\CodeTypeTok{promises}\OperatorTok{;} \CommentTok{// Promise{-}based fs}
 
\KeywordTok{async} \KeywordTok{function} \CodeFuncTok{readFile}() \{
  \KeywordTok{try} \{
    \KeywordTok{const} content \OperatorTok{=} \KeywordTok{await} fs\OperatorTok{.}\CodeFuncTok{readFile}(\StringTok{"my\_notes.txt"}\OperatorTok{,} \StringTok{"utf8"})\OperatorTok{;}
    console\OperatorTok{.}\CodeFuncTok{log}(content)\OperatorTok{;}
  \} \KeywordTok{catch} (err) \{
    console\OperatorTok{.}\CodeFuncTok{log}(\StringTok{\textasciigrave{}Error reading file: }\StringTok{\$\{}err\OperatorTok{.}\CodeTypeTok{message}\StringTok{\}}\StringTok{\textasciigrave{}})\OperatorTok{;}
  \}
\}
 
\KeywordTok{async} \KeywordTok{function} \CodeFuncTok{writeFile}() \{
  \KeywordTok{try} \{
    \KeywordTok{await} fs\OperatorTok{.}\CodeFuncTok{writeFile}(\StringTok{"output.txt"}\OperatorTok{,} \StringTok{"Async write!}\StringTok{\textbackslash{}n}\StringTok{"})\OperatorTok{;}
    \KeywordTok{await} fs\OperatorTok{.}\CodeFuncTok{appendFile}(\StringTok{"output.txt"}\OperatorTok{,} \StringTok{"Appended!}\StringTok{\textbackslash{}n}\StringTok{"})\OperatorTok{;}
    console\OperatorTok{.}\CodeFuncTok{log}(\StringTok{"Files written."})\OperatorTok{;}
  \ \KeywordTok{catch} (err) \{
    console\OperatorTok{.}\CodeFuncTok{log}(\StringTok{\textasciigrave{}Error: }\StringTok{\$\{}err\OperatorTok{.}\CodeTypeTok{message}\StringTok{\}}\StringTok{\textasciigrave{}})\OperatorTok{;}
  \}
\}
 
\CodeFuncTok{readFile}()\OperatorTok{;}
\CodeFuncTok{writeFile}()\OperatorTok{;}
\end{Highlighting}
\end{Shaded}

\begin{center}\rule{0.5\linewidth}{0.5pt}\end{center}

\subsection{Lesson 8.2: Error Handling}\label{lesson-8.2-error-handling}

\begin{Shaded}
\begin{Highlighting}[]
\CommentTok{// try{-}catch{-}finally}
\KeywordTok{try} \{
  \KeywordTok{const} result \OperatorTok{=} \CodeNumTok{10} \OperatorTok{/} \CodeNumTok{0}\OperatorTok{;} \CommentTok{// Doesn\textquotesingle{}t throw in JS (gives Infinity)}
  \KeywordTok{const} parsed \OperatorTok{=} JSON\OperatorTok{.}\CodeFuncTok{parse}(\StringTok{"invalid json"})\OperatorTok{;} \CommentTok{// This throws}
\ \}\KeywordTok{catch} (err) \{
  console\OperatorTok{.}\CodeFuncTok{log}(\StringTok{\textasciigrave{}Error type: }\StringTok{\$\{}err\OperatorTok{.}\CodeTypeTok{constructor}\OperatorTok{.}\CodeTypeTok{name}\StringTok{\}}\StringTok{\textasciigrave{}})\OperatorTok{;}
  console\OperatorTok{.}\CodeFuncTok{log}(\StringTok{\textasciigrave{}Error message: }\StringTok{\$\{}err\OperatorTok{.}\CodeTypeTok{message}\StringTok{\}}\StringTok{\textasciigrave{}})\OperatorTok{;}
\ \}\KeywordTok{finally} \{
  console\OperatorTok{.}\CodeFuncTok{log}(\StringTok{"This always runs."})\OperatorTok{;}
\}
 
\CommentTok{// Throwing custom errors}
\KeywordTok{function} \CodeFuncTok{divide}(a\OperatorTok{,} b) \{
  \KeywordTok{if} (b \OperatorTok{===} \CodeNumTok{0}) \{
    \KeywordTok{throw} \KeywordTok{new} Error(\StringTok{"Division by zero is not allowed."})\OperatorTok{;}
  \}
  \KeywordTok{return} a \OperatorTok{/} b\OperatorTok{;}
\}
 
\KeywordTok{try} \{
  console\OperatorTok{.}\CodeFuncTok{log}(\CodeFuncTok{divide}(\CodeNumTok{10}\OperatorTok{,} \CodeNumTok{2}))\OperatorTok{;} \CommentTok{// 5}
  console\OperatorTok{.}\CodeFuncTok{log}(\CodeFuncTok{divide}(\CodeNumTok{10}\OperatorTok{,} \CodeNumTok{0}))\OperatorTok{;} \CommentTok{// Throws}
\} \KeywordTok{catch} (err) \{
  console\OperatorTok{.}\CodeFuncTok{log}(\StringTok{\textasciigrave{}Caught: }\StringTok{\$\{}err\OperatorTok{.}\CodeTypeTok{message}\StringTok{\}}\StringTok{\textasciigrave{}})\OperatorTok{;}
\}
\end{Highlighting}
\end{Shaded}

\subsection{Custom Error Classes}\label{custom-error-classes}

\begin{Shaded}
\begin{Highlighting}[]
\KeywordTok{class} ValidationError \KeywordTok{extends} Error \{
  \CodeFuncTok{constructor}(field\OperatorTok{,} message) \{
    \KeywordTok{super}(message)\OperatorTok{;}
    \KeywordTok{this}\OperatorTok{.}\CodeTypeTok{name} \OperatorTok{=} \StringTok{"ValidationError"}\OperatorTok{;}
    \KeywordTok{this}\OperatorTok{.}\CodeTypeTok{field} \OperatorTok{=} field\OperatorTok{;}
  \}
\}
 
\KeywordTok{class} InsufficientFundsError \KeywordTok{extends} Error \{
  \CodeFuncTok{constructor}(needed) \{
    \KeywordTok{super}(\StringTok{\textasciigrave{}Insufficient funds. Need $}\StringTok{\$\{}needed\StringTok{\}}\StringTok{ more.\textasciigrave{}})\OperatorTok{;}
    \KeywordTok{this}\OperatorTok{.}\CodeTypeTok{name} \OperatorTok{=} \StringTok{"InsufficientFundsError"}\OperatorTok{;}
    \KeywordTok{this}\OperatorTok{.}\CodeTypeTok{needed} \OperatorTok{=} needed\OperatorTok{;}
  \}
\}
 
\KeywordTok{try} \{
  \KeywordTok{throw} \KeywordTok{new} \CodeFuncTok{ValidationError}(\StringTok{"email"}\OperatorTok{,} \StringTok{"Invalid email format."})\OperatorTok{;}
\} \KeywordTok{catch} (err) \{
  \KeywordTok{if} (err \KeywordTok{instanceof} ValidationError) \{
    console\OperatorTok{.}\CodeFuncTok{log}(\StringTok{\textasciigrave{}Validation failed on \textquotesingle{}}\StringTok{\$\{}err\OperatorTok{.}\CodeTypeTok{field}\StringTok{\}}\StringTok{\textquotesingle{}: }\StringTok{\$\{}err\OperatorTok{.}\CodeTypeTok{message}\StringTok{\}}\StringTok{\textasciigrave{}})\OperatorTok{;}
  \ \KeywordTok{else} \{
    \KeywordTok{throw} err\OperatorTok{;} \CommentTok{// Re{-}throw if not our custom error}
  \}
\}
\end{Highlighting}
\end{Shaded}

\begin{center}\rule{0.5\linewidth}{0.5pt}\end{center}

\subsection{Lesson 8.3: JSON in
Node.js}\label{lesson-8.3-json-in-node.js}

\begin{Shaded}
\begin{Highlighting}[]
\KeywordTok{const} fs \OperatorTok{=} \CodeAccentTok{require}(\StringTok{"fs"})\OperatorTok{;}
 
\CommentTok{// JavaScript object to JSON string}
\KeywordTok{const} settings \OperatorTok{=} \{\}
  \CodeTypeTok{username}\OperatorTok{:} \StringTok{"coder\_kid"}\OperatorTok{,}
  \CodeTypeTok{theme}\OperatorTok{:} \StringTok{"dark"}\OperatorTok{,}
  \CodeTypeTok{notificationsOn}\OperatorTok{:} \KeywordTok{true}\OperatorTok{,}
  \CodeTypeTok{favoriteNumbers}\OperatorTok{:} [\CodeNumTok{7}\OperatorTok{,} \CodeNumTok{13}\OperatorTok{,} \CodeNumTok{42}]\OperatorTok{,}
\}\OperatorTok{;}
 
\KeywordTok{const} jsonString \OperatorTok{=} JSON\OperatorTok{.}\CodeFuncTok{stringify}(settings\OperatorTok{,} \KeywordTok{null}\OperatorTok{,} \CodeNumTok{2})\OperatorTok{;} \CommentTok{// Pretty{-}printed}
console\OperatorTok{.}\CodeFuncTok{log}(jsonString)\OperatorTok{;}
 
\CommentTok{// Save to file}
fs\OperatorTok{.}\CodeFuncTok{writeFileSync}(\StringTok{"settings.json"}\OperatorTok{,} jsonString)\OperatorTok{;}
 
\CommentTok{// Load from file}
\KeywordTok{try} \{
  \KeywordTok{const} loaded \OperatorTok{=} JSON\OperatorTok{.}\CodeFuncTok{parse}(fs\OperatorTok{.}\CodeFuncTok{readFileSync}(\StringTok{"settings.json"}\OperatorTok{,} \StringTok{"utf8"}))\OperatorTok{;}
  console\OperatorTok{.}\CodeFuncTok{log}(\StringTok{\textasciigrave{}Username: }\StringTok{\$\{}loaded\OperatorTok{.}\CodeTypeTok{username}\StringTok{\}}\StringTok{\textasciigrave{}})\OperatorTok{;}
  console\OperatorTok{.}\CodeFuncTok{log}(\StringTok{\textasciigrave{}Theme: }\StringTok{\$\{}loaded\OperatorTok{.}\CodeTypeTok{theme}\StringTok{\}}\StringTok{\textasciigrave{}})\OperatorTok{;}
  console\OperatorTok{.}\CodeFuncTok{log}(\StringTok{\textasciigrave{}Numbers: }\StringTok{\$\{}loaded\OperatorTok{.}\CodeTypeTok{favoriteNumbers}\StringTok{\}}\StringTok{\textasciigrave{}})\OperatorTok{;}
\ \KeywordTok{catch} (err) \{
  \KeywordTok{if} (err \KeywordTok{instanceof} SyntaxError) \{\}
    console\OperatorTok{.}\CodeFuncTok{log}(\StringTok{"Error: JSON file is corrupted."})\OperatorTok{;}
  \} \KeywordTok{else} \KeywordTok{if} (err\OperatorTok{.}\CodeTypeTok{code} \OperatorTok{===} \StringTok{"ENOENT"}) \{
    console\OperatorTok{.}\CodeFuncTok{log}(\StringTok{"Error: File not found."})\OperatorTok{;}
  \} \KeywordTok{else} \{
    console\OperatorTok{.}\CodeFuncTok{log}(\StringTok{\textasciigrave{}Error: }\StringTok{\$\{}err\OperatorTok{.}\CodeTypeTok{message}\StringTok{\}}\StringTok{\textasciigrave{}})\OperatorTok{;}
  \}
\}
\end{Highlighting}
\end{Shaded}

\subsection{Skills Checklist}\label{skills-checklist-65}

\begin{itemize}
\tightlist
\item[$\square$]
  I can read and write files using the \texttt{fs} module.
\item[$\square$]
  I can use \texttt{try-catch-finally} to handle errors.
\item[$\square$]
  I can throw custom errors using \texttt{Error} or custom classes.
\item[$\square$]
  I can serialize/deserialize data using \texttt{JSON.stringify} and
  \texttt{JSON.parse}.
\end{itemize}

\newpage

\section{Unit 9: Modules, Collections, and
Testing}\label{unit-9-modules-collections-and-testing}

\begin{center}\rule{0.5\linewidth}{0.5pt}\end{center}

\subsection{Lesson 9.1: JavaScript
Modules}\amn{\textbf{Accessibility Tip:} Run Jest with \texttt{npx jest -{}-verbose} to get one result line per test. Each line reads ``PASS'' or ``FAIL'' followed by the test name. Navigate the output with arrow keys after running. Verbose mode is far more screen-reader-friendly than the default dot notation.}
\subsection{Lesson 9.1: JavaScript
Modules}\label{lesson-9.1-javascript-modules}

\subsection{CommonJS (Node.js
default)}\label{commonjs-node.js-default}

\begin{Shaded}
\begin{Highlighting}[]
\CommentTok{// math.js {-}{-}{-} exporting}
\KeywordTok{const} add \OperatorTok{=} (a\OperatorTok{,} b) \KeywordTok{=\textgreater{}} a \OperatorTok{+} b\OperatorTok{;}
\KeywordTok{const} multiply \OperatorTok{=} (a\OperatorTok{,} b) \KeywordTok{=\textgreater{}} a \OperatorTok{*} b\OperatorTok{;}
\KeywordTok{const} \CodeNumTok{PI} \OperatorTok{=} Math\OperatorTok{.}\CodeNumTok{PI}\OperatorTok{;}
 
module\OperatorTok{.}\CodeTypeTok{exports} \OperatorTok{=} \{ add\OperatorTok{,} multiply\OperatorTok{,} \CodeNumTok{PI} \}\OperatorTok{;}
 
\CommentTok{// Or export individually:}
\CommentTok{// module.exports.add = add;}
\end{Highlighting}
\end{Shaded}

\begin{Shaded}
\begin{Highlighting}[]
\CommentTok{// main.js {-}{-}{-} importing}
\KeywordTok{const} \{ add\OperatorTok{,} multiply\OperatorTok{,} \CodeNumTok{PI} \} \OperatorTok{=} \CodeAccentTok{require}(\StringTok{"./math"})\OperatorTok{;}
 
console\OperatorTok{.}\CodeFuncTok{log}(\CodeFuncTok{add}(\CodeNumTok{5}\OperatorTok{,} \CodeNumTok{3}))\OperatorTok{;} \CommentTok{// 8}
console\OperatorTok{.}\CodeFuncTok{log}(\CodeFuncTok{multiply}(\CodeNumTok{4}\OperatorTok{,} \CodeNumTok{6}))\OperatorTok{;} \CommentTok{// 24}
console\OperatorTok{.}\CodeFuncTok{log}(\CodeNumTok{PI}\OperatorTok{.}\CodeFuncTok{toFixed}(\CodeNumTok{5}))\OperatorTok{;} \CommentTok{// 3.14159}
\end{Highlighting}
\end{Shaded}

\subsection{\texorpdfstring{ES Modules (Modern --- requires
\texttt{"type":\ "module"} in
package.json)}{ES Modules (Modern --- requires "type": "module" in package.json)}}\label{es-modules-modern-requires-type-module-in-package.json}

\begin{Shaded}
\begin{Highlighting}[]
\CommentTok{// math.mjs}
\CodeAccentTok{export} \KeywordTok{const} add \OperatorTok{=} (a\OperatorTok{,} b) \KeywordTok{=\textgreater{}} a \OperatorTok{+} b\OperatorTok{;}
\CodeAccentTok{export} \KeywordTok{const} multiply \OperatorTok{=} (a\OperatorTok{,} b) \KeywordTok{=\textgreater{}} a \OperatorTok{*} b\OperatorTok{;}
\CodeAccentTok{export} \KeywordTok{const} \CodeNumTok{PI} \OperatorTok{=} Math\OperatorTok{.}\CodeNumTok{PI}\OperatorTok{;}
\CodeAccentTok{export} \CodeAccentTok{default} \KeywordTok{function} \CodeFuncTok{greet}(name) \{
  \KeywordTok{return} \StringTok{\textasciigrave{}Hello, }\StringTok{\$\{}name\StringTok{\}}\StringTok{!\textasciigrave{}}\OperatorTok{;}
\}
\end{Highlighting}
\end{Shaded}

\begin{Shaded}
\begin{Highlighting}[]
\CommentTok{// main.mjs}
\CodeAccentTok{import} greet\OperatorTok{,} \{ add\OperatorTok{,} multiply\OperatorTok{,} \CodeNumTok{PI} \} \CodeAccentTok{from} \StringTok{"./math.mjs"}\OperatorTok{;}
 
console\OperatorTok{.}\CodeFuncTok{log}(\CodeFuncTok{greet}(\StringTok{"Alice"}))\OperatorTok{;}
console\OperatorTok{.}\CodeFuncTok{log}(\CodeFuncTok{add}(\CodeNumTok{5}\OperatorTok{,} \CodeNumTok{3}))\OperatorTok{;}
\end{Highlighting}
\end{Shaded}

\begin{center}\rule{0.5\linewidth}{0.5pt}\end{center}

\subsection{Lesson 9.2: Maps and Sets}\label{lesson-9.2-maps-and-sets}

\subsection{\texorpdfstring{\texttt{Map} --- Key-Value (Better than
plain objects for
maps)}{Map --- Key-Value (Better than plain objects for maps)}}\label{map-key-value-better-than-plain-objects-for-maps}

\begin{Shaded}
\begin{Highlighting}[]
\KeywordTok{const} inventory \OperatorTok{=} \KeywordTok{new} Map()\OperatorTok{;}
 
\CommentTok{// Adding}
inventory\OperatorTok{.}\CodeFuncTok{set}(\StringTok{"apples"}\OperatorTok{,} \CodeNumTok{50})\OperatorTok{;}
inventory\OperatorTok{.}\CodeFuncTok{set}(\StringTok{"bananas"}\OperatorTok{,} \CodeNumTok{30})\OperatorTok{;}
inventory\OperatorTok{.}\CodeFuncTok{set}(\StringTok{"oranges"}\OperatorTok{,} \CodeNumTok{40})\OperatorTok{;}
 
\CommentTok{// Access}
console\OperatorTok{.}\CodeFuncTok{log}(inventory\OperatorTok{.}\CodeFuncTok{get}(\StringTok{"apples"}))\OperatorTok{;} \CommentTok{// 50}
console\OperatorTok{.}\CodeFuncTok{log}(inventory\OperatorTok{.}\CodeFuncTok{has}(\StringTok{"bananas"}))\OperatorTok{;} \CommentTok{// true}
console\OperatorTok{.}\CodeFuncTok{log}(inventory\OperatorTok{.}\CodeTypeTok{size})\OperatorTok{;} \CommentTok{// 3}
 
\CommentTok{// Delete}
inventory\OperatorTok{.}\CodeFuncTok{delete}(\StringTok{"bananas"})\OperatorTok{;}
 
\CommentTok{// Iterate}
\KeywordTok{for} (\KeywordTok{const} [key\OperatorTok{,} value] \KeywordTok{of} inventory) \{
  console\OperatorTok{.}\CodeFuncTok{log}(\StringTok{\textasciigrave{}}\StringTok{\$\{}key\StringTok{\}}\StringTok{: }\StringTok{\$\{}value\StringTok{\}}\StringTok{\textasciigrave{}})\OperatorTok{;}
\}
 
inventory\OperatorTok{.}\CodeFuncTok{forEach}((value\OperatorTok{,} key) \KeywordTok{=\textgreater{}} \{
  console\OperatorTok{.}\CodeFuncTok{log}(\StringTok{\textasciigrave{}}\StringTok{\$\{}key\StringTok{\}}\StringTok{ = }\StringTok{\$\{}value\StringTok{\}}\StringTok{\textasciigrave{}})\OperatorTok{;}
)\OperatorTok{;}
 
console\OperatorTok{.}\CodeFuncTok{log}([\OperatorTok{...}inventory\OperatorTok{.}\CodeFuncTok{keys}()])\OperatorTok{;}
console\OperatorTok{.}\CodeFuncTok{log}([\OperatorTok{...}inventory\OperatorTok{.}\CodeFuncTok{values}()])\OperatorTok{;}
console\OperatorTok{.}\CodeFuncTok{log}([\OperatorTok{...}inventory\OperatorTok{.}\CodeFuncTok{entries}()])\OperatorTok{;}
\end{Highlighting}
\end{Shaded}

\subsection{\texorpdfstring{\texttt{Set} --- Unique
Values}{Set --- Unique Values}}\label{set-unique-values}

\begin{Shaded}
\begin{Highlighting}[]
\KeywordTok{const} uniqueWords \OperatorTok{=} \KeywordTok{new} Set()\OperatorTok{;}
uniqueWords\OperatorTok{.}\CodeFuncTok{add}(\StringTok{"apple"})\OperatorTok{;}
uniqueWords\OperatorTok{.}\CodeFuncTok{add}(\StringTok{"banana"})\OperatorTok{;}
uniqueWords\OperatorTok{.}\CodeFuncTok{add}(\StringTok{"apple"})\OperatorTok{;} \CommentTok{// Duplicate {-}{-}{-} ignored}
 
console\OperatorTok{.}\CodeFuncTok{log}(uniqueWords\OperatorTok{.}\CodeTypeTok{size})\OperatorTok{;} \CommentTok{// 2}
console\OperatorTok{.}\CodeFuncTok{log}(uniqueWords\OperatorTok{.}\CodeFuncTok{has}(\StringTok{"banana"}))\OperatorTok{;} \CommentTok{// true}
uniqueWords\OperatorTok{.}\CodeFuncTok{delete}(\StringTok{"banana"})\OperatorTok{;}
 
\CommentTok{// Convert array to set (remove duplicates)}
\KeywordTok{const} nums \OperatorTok{=} [\CodeNumTok{1}\OperatorTok{,} \CodeNumTok{2}\OperatorTok{,} \CodeNumTok{2}\OperatorTok{,} \CodeNumTok{3}\OperatorTok{,} \CodeNumTok{3}\OperatorTok{,} \CodeNumTok{3}\OperatorTok{,} \CodeNumTok{4}]\OperatorTok{;}
\KeywordTok{const} unique \OperatorTok{=} [\OperatorTok{...}\KeywordTok{new} Set(nums)]\OperatorTok{;}
console\OperatorTok{.}\CodeFuncTok{log}(unique)\OperatorTok{;} \CommentTok{// [1, 2, 3, 4]}
 
\CommentTok{// Iterate}
\KeywordTok{for} (\KeywordTok{const} item \KeywordTok{of} uniqueWords) \{\}
  console\OperatorTok{.}\CodeFuncTok{log}(item)\OperatorTok{;}
\}
\end{Highlighting}
\end{Shaded}

\begin{center}\rule{0.5\linewidth}{0.5pt}\end{center}

\subsection{Lesson 9.3: Word Frequency
Analysis}\label{lesson-9.3-word-frequency-analysis}

\begin{Shaded}
\begin{Highlighting}[]
\KeywordTok{const} fs \OperatorTok{=} \CodeAccentTok{require}(\StringTok{"fs"})\OperatorTok{;}
 
\KeywordTok{function} \CodeFuncTok{countWords}(text) \{
  \KeywordTok{const} words \OperatorTok{=} text
    \OperatorTok{.}\CodeFuncTok{toLowerCase}()
    \OperatorTok{.}\CodeFuncTok{replace}(\StringTok{/}\StringTok{[\^{}a{-}z\textbackslash{}s]}\StringTok{/g}\OperatorTok{,} \StringTok{""})
    \OperatorTok{.}\CodeFuncTok{trim}()
    \OperatorTok{.}\CodeFuncTok{split}(\StringTok{/}\StringTok{\textbackslash{}s+}\StringTok{/})
    \OperatorTok{.}\CodeFuncTok{filter}((word) \KeywordTok{=\textgreater{}} word\OperatorTok{.}\CodeTypeTok{length} \OperatorTok{\textgreater{}} \CodeNumTok{0})\OperatorTok{;}
 
  \KeywordTok{const} frequency \OperatorTok{=} \KeywordTok{new} Map()\OperatorTok{;}
  \KeywordTok{for} (\KeywordTok{const} word \KeywordTok{of} words) \{
    frequency\OperatorTok{.}\CodeFuncTok{set}(word\OperatorTok{,} (frequency\OperatorTok{.}\CodeFuncTok{get}(word) \OperatorTok{||} \CodeNumTok{0}) \OperatorTok{+} \CodeNumTok{1})\OperatorTok{;}
  \}
 
  \KeywordTok{return} frequency\OperatorTok{;}
\}
 
\KeywordTok{const} text \OperatorTok{=} \StringTok{"the quick brown fox jumps over the lazy dog the quick brown fox"}\OperatorTok{;}
\KeywordTok{const} freq \OperatorTok{=} \CodeFuncTok{countWords}(text)\OperatorTok{;}
 
\CommentTok{// Sort by frequency (most common first)}
\KeywordTok{const} sorted \OperatorTok{=} [\OperatorTok{...}freq\OperatorTok{.}\CodeFuncTok{entries}()]\OperatorTok{.}\CodeFuncTok{sort}((a\OperatorTok{,} b) \KeywordTok{=\textgreater{}} b[\CodeNumTok{1}] \OperatorTok{{-}} a[\CodeNumTok{1}])\OperatorTok{;}
 
console\OperatorTok{.}\CodeFuncTok{log}(\StringTok{"Word Frequencies:"})\OperatorTok{;}
\KeywordTok{for} (\KeywordTok{const} [word\OperatorTok{,} count] \KeywordTok{of} sorted) \{
  console\OperatorTok{.}\CodeFuncTok{log}(\StringTok{\textasciigrave{}\textquotesingle{}}\StringTok{\$\{}word\StringTok{\}}\StringTok{\textquotesingle{}: }\StringTok{\$\{}count\StringTok{\}}\StringTok{\textasciigrave{}})\OperatorTok{;}
\}
 
\CommentTok{// Top 5 words}
console\OperatorTok{.}\CodeFuncTok{log}(\StringTok{"}\StringTok{\textbackslash{}n}\StringTok{Top 5 words:"})\OperatorTok{;}
sorted\OperatorTok{.}\CodeFuncTok{slice}(\CodeNumTok{0}\OperatorTok{,} \CodeNumTok{5})\OperatorTok{.}\CodeFuncTok{forEach}(([word\OperatorTok{,} count]) \KeywordTok{=\textgreater{}} \{
  console\OperatorTok{.}\CodeFuncTok{log}(\StringTok{\textasciigrave{}  }\StringTok{\$\{}word\StringTok{\}}\StringTok{: }\StringTok{\$\{}count\StringTok{\}}\StringTok{\textasciigrave{}})\OperatorTok{;}
)\OperatorTok{;}
\end{Highlighting}
\end{Shaded}

\begin{center}\rule{0.5\linewidth}{0.5pt}\end{center}

\subsection{Lesson 9.4: Working with CSV
Files}\label{lesson-9.4-working-with-csv-files}

\begin{Shaded}
\begin{Highlighting}[]
\KeywordTok{const} fs \OperatorTok{=} \CodeAccentTok{require}(\StringTok{"fs"})\OperatorTok{;}
 
\KeywordTok{function} \CodeFuncTok{parseCSV}(content) \{
  \KeywordTok{const} lines \OperatorTok{=} content\OperatorTok{.}\CodeFuncTok{trim}()\OperatorTok{.}\CodeFuncTok{split}(\StringTok{"}\StringTok{\textbackslash{}n}\StringTok{"})\OperatorTok{;}
  \KeywordTok{const} headers \OperatorTok{=} lines[\CodeNumTok{0}]\OperatorTok{.}\CodeFuncTok{split}(\StringTok{","})\OperatorTok{.}\CodeFuncTok{map}((h) \KeywordTok{=\textgreater{}} h\OperatorTok{.}\CodeFuncTok{trim}())\OperatorTok{;}
 
  \KeywordTok{return} lines\OperatorTok{.}\CodeFuncTok{slice}(\CodeNumTok{1})\OperatorTok{.}\CodeFuncTok{map}((line) \KeywordTok{=\textgreater{}} \{
    \KeywordTok{const} values \OperatorTok{=} line\OperatorTok{.}\CodeFuncTok{split}(\StringTok{","})\OperatorTok{.}\CodeFuncTok{map}((v) \KeywordTok{=\textgreater{}} v\OperatorTok{.}\CodeFuncTok{trim}())\OperatorTok{;}
    \KeywordTok{const} record \OperatorTok{=} \{\}\OperatorTok{;}
    headers\OperatorTok{.}\CodeFuncTok{forEach}((header\OperatorTok{,} index) \KeywordTok{=\textgreater{}} \{\}
      record[header] \OperatorTok{=} values[index]\OperatorTok{;}
    \})\OperatorTok{;}
    \KeywordTok{return} record\OperatorTok{;}
  )\OperatorTok{;}
\}
 
\KeywordTok{function} \CodeFuncTok{generateCSV}(records) \{
  \KeywordTok{if} (records\OperatorTok{.}\CodeTypeTok{length} \OperatorTok{===} \CodeNumTok{0}) \KeywordTok{return} \StringTok{""}\OperatorTok{;}
  \KeywordTok{const} headers \OperatorTok{=} Object\OperatorTok{.}\CodeFuncTok{keys}(records[\CodeNumTok{0}])\OperatorTok{;}
  \KeywordTok{const} headerLine \OperatorTok{=} headers\OperatorTok{.}\CodeFuncTok{join}(\StringTok{","})\OperatorTok{;}
  \KeywordTok{const} dataLines \OperatorTok{=} records\OperatorTok{.}\CodeFuncTok{map}((record) \KeywordTok{=\textgreater{}}
    headers\OperatorTok{.}\CodeFuncTok{map}((h) \KeywordTok{=\textgreater{}} record[h])\OperatorTok{.}\CodeFuncTok{join}(\StringTok{","})\OperatorTok{,}
  )\OperatorTok{;}
  \KeywordTok{return} [headerLine\OperatorTok{,} \OperatorTok{...}dataLines]\OperatorTok{.}\CodeFuncTok{join}(\StringTok{"}\StringTok{\textbackslash{}n}\StringTok{"})\OperatorTok{;}
\}
 
\CommentTok{// Read and parse CSV}
\KeywordTok{try} \{
  \KeywordTok{const} content \OperatorTok{=} fs\OperatorTok{.}\CodeFuncTok{readFileSync}(\StringTok{"students.csv"}\OperatorTok{,} \StringTok{"utf8"})\OperatorTok{;}
  \KeywordTok{const} students \OperatorTok{=} \CodeFuncTok{parseCSV}(content)\OperatorTok{;}
 
  console\OperatorTok{.}\CodeFuncTok{log}(\StringTok{\textasciigrave{}Loaded }\StringTok{\$\{}students\OperatorTok{.}\CodeTypeTok{length}\StringTok{\}}\StringTok{ students:\textasciigrave{}})\OperatorTok{;}
  students\OperatorTok{.}\CodeFuncTok{forEach}((student) \KeywordTok{=\textgreater{}} \{
    console\OperatorTok{.}\CodeFuncTok{log}(
      \StringTok{\textasciigrave{}  Name: }\StringTok{\$\{}student\OperatorTok{.}\CodeTypeTok{name}\StringTok{\}}\StringTok{, Age: }\StringTok{\$\{}student\OperatorTok{.}\CodeTypeTok{age}\StringTok{\}}\StringTok{, Grade: }\StringTok{\$\{}student\OperatorTok{.}\CodeTypeTok{grade}\StringTok{\}}\StringTok{\textasciigrave{}}\OperatorTok{,}
    )\OperatorTok{;}
  )\}\OperatorTok{;}
\ \}\KeywordTok{catch} (err) \{
  \KeywordTok{if} (err\OperatorTok{.}\CodeTypeTok{code} \OperatorTok{===} \StringTok{"ENOENT"}) \{
    \CommentTok{// Create a sample CSV for demonstration}
    \KeywordTok{const} sampleData \OperatorTok{=} \StringTok{\textasciigrave{}name,age,grade}
\StringTok{Alice,15,A}
\StringTok{Bob,14,B}
\StringTok{Charlie,16,A\textasciigrave{}}\OperatorTok{;}
    fs\OperatorTok{.}\CodeFuncTok{writeFileSync}(\StringTok{"students.csv"}\OperatorTok{,} sampleData)\OperatorTok{;}
    console\OperatorTok{.}\CodeFuncTok{log}(\StringTok{"Created sample students.csv"})\OperatorTok{;}
  \}
\}
\end{Highlighting}
\end{Shaded}

\begin{center}\rule{0.5\linewidth}{0.5pt}\end{center}

\subsection{Lesson 9.5: Testing with
Jest}\label{lesson-9.5-testing-with-jest}

\subsection{Setting Up Jest}\label{setting-up-jest}

\begin{Shaded}
\begin{Highlighting}[]
npm init \CodeTypeTok{{-}y}
npm install \CodeTypeTok{{-}{-}save{-}dev} jest
\end{Highlighting}
\end{Shaded}

Add to \texttt{package.json}:

\begin{Shaded}
\begin{Highlighting}[]
\CodeFuncTok{\{}
  \CodeTypeTok{"scripts"}\CodeFuncTok{:} \CodeFuncTok{\{}
    \CodeTypeTok{"test"}\CodeFuncTok{:} \StringTok{"jest"}
  \CodeFuncTok{\}}
\CodeFuncTok{\}}
\end{Highlighting}
\end{Shaded}

\subsection{Writing Tests}\label{writing-tests}

\begin{Shaded}
\begin{Highlighting}[]
\CommentTok{// math.js}
\KeywordTok{const} add \OperatorTok{=} (a\OperatorTok{,} b) \KeywordTok{=\textgreater{}} a \OperatorTok{+} b\OperatorTok{;}
\KeywordTok{const} getGrade \OperatorTok{=} (score) \KeywordTok{=\textgreater{}} \{
  \KeywordTok{if} (score \OperatorTok{\textgreater{}=} \CodeNumTok{90}) \KeywordTok{return} \StringTok{"A"}\OperatorTok{;}
  \KeywordTok{if} (score \OperatorTok{\textgreater{}=} \CodeNumTok{80}) \KeywordTok{return} \StringTok{"B"}\OperatorTok{;}
  \KeywordTok{if} (score \OperatorTok{\textgreater{}=} \CodeNumTok{70}) \KeywordTok{return} \StringTok{"C"}\OperatorTok{;}
  \KeywordTok{return} \StringTok{"F"}\OperatorTok{;}
\}\OperatorTok{;}
\KeywordTok{const} circleArea \OperatorTok{=} (radius) \KeywordTok{=\textgreater{}} Math\OperatorTok{.}\CodeNumTok{PI} \OperatorTok{*} radius \OperatorTok{**} \CodeNumTok{2}\OperatorTok{;}
 
module\OperatorTok{.}\CodeTypeTok{exports} \OperatorTok{=} \{ add\}\OperatorTok{,} getGrade\OperatorTok{,} circleArea \}\OperatorTok{;}
\end{Highlighting}
\end{Shaded}

\begin{Shaded}
\begin{Highlighting}[]
\CommentTok{// math.test.js}
\KeywordTok{const} \{ add\OperatorTok{,} getGrade\OperatorTok{,} circleArea \} \OperatorTok{=} \CodeAccentTok{require}(\StringTok{"./math"})\OperatorTok{;}
 
\CodeFuncTok{describe}(\StringTok{"add function"}\OperatorTok{,} () \KeywordTok{=\textgreater{}} \{
  \CodeFuncTok{test}(\StringTok{"adds two positive numbers"}\OperatorTok{,} () \KeywordTok{=\textgreater{}} \{
    \CodeFuncTok{expect}(\CodeFuncTok{add}(\CodeNumTok{5}\OperatorTok{,} \CodeNumTok{3}))\OperatorTok{.}\CodeFuncTok{toBe}(\CodeNumTok{8})\OperatorTok{;}
  \})\OperatorTok{;}
 
  \CodeFuncTok{test}(\StringTok{"adds negative numbers"}\OperatorTok{,} () \KeywordTok{=\textgreater{}} \{
    \CodeFuncTok{expect}(\CodeFuncTok{add}(\OperatorTok{{-}}\CodeNumTok{5}\OperatorTok{,} \OperatorTok{{-}}\CodeNumTok{3}))\OperatorTok{.}\CodeFuncTok{toBe}(\OperatorTok{{-}}\CodeNumTok{8})\OperatorTok{;}
  )\OperatorTok{;}
 
  \CodeFuncTok{test}(\StringTok{"adds zero"}\OperatorTok{,} () \KeywordTok{=\textgreater{}} \{\}
    \CodeFuncTok{expect}(\CodeFuncTok{add}(\CodeNumTok{0}\OperatorTok{,} \CodeNumTok{7}))\OperatorTok{.}\CodeFuncTok{toBe}(\CodeNumTok{7})\OperatorTok{;}
  \})\OperatorTok{;}
)\OperatorTok{;}
 
\CodeFuncTok{describe}(\StringTok{"getGrade function"}\OperatorTok{,} () \KeywordTok{=\textgreater{}} \{
  \CodeFuncTok{test}(\StringTok{"returns A for score \textgreater{}= 90"}\OperatorTok{,} () \KeywordTok{=\textgreater{}} \{\}
    \CodeFuncTok{expect}(\CodeFuncTok{getGrade}(\CodeNumTok{95}))\OperatorTok{.}\CodeFuncTok{toBe}(\StringTok{"A"})\OperatorTok{;}
    \CodeFuncTok{expect}(\CodeFuncTok{getGrade}(\CodeNumTok{90}))\OperatorTok{.}\CodeFuncTok{toBe}(\StringTok{"A"})\OperatorTok{;}
  \})\OperatorTok{;}
 
  \CodeFuncTok{test}(\StringTok{"returns B for score 80{-}89"}\OperatorTok{,} () \KeywordTok{=\textgreater{}} \{
    \CodeFuncTok{expect}(\CodeFuncTok{getGrade}(\CodeNumTok{85}))\OperatorTok{.}\CodeFuncTok{toBe}(\StringTok{"B"})\OperatorTok{;}
  \})\OperatorTok{;}
 
  \CodeFuncTok{test}(\StringTok{"returns F for score \textless{} 60"}\OperatorTok{,} () \KeywordTok{=\textgreater{}} \{
    \CodeFuncTok{expect}(\CodeFuncTok{getGrade}(\CodeNumTok{55}))\OperatorTok{.}\CodeFuncTok{toBe}(\StringTok{"F"})\OperatorTok{;}
  \})\OperatorTok{;}
)\OperatorTok{;}
 
\CodeFuncTok{describe}(\StringTok{"circleArea function"}\OperatorTok{,} () \KeywordTok{=\textgreater{}} \{
  \CodeFuncTok{test}(\StringTok{"calculates area of circle"}\OperatorTok{,} () \KeywordTok{=\textgreater{}} \{\}
    \CodeFuncTok{expect}(\CodeFuncTok{circleArea}(\CodeNumTok{5}))\OperatorTok{.}\CodeFuncTok{toBeCloseTo}(\CodeNumTok{78.54}\OperatorTok{,} \CodeNumTok{1})\OperatorTok{;}
  \})\OperatorTok{;}
)\OperatorTok{;}
\end{Highlighting}
\end{Shaded}

Run tests:

\begin{Shaded}
\begin{Highlighting}[]
npm test
\end{Highlighting}
\end{Shaded}

\subsection{Common Jest Matchers}\label{common-jest-matchers}

\begin{fullwidth}{\def\LTcaptype{none} % do not increment counter
\begin{longtable}[]{@{}ll@{}}
\toprule\noalign{}
Matcher & Description \\
\midrule\noalign{}
\endhead
\bottomrule\noalign{}
\endlastfoot
\texttt{toBe(value)} & Strict equality (\texttt{===}) \\
\texttt{toEqual(value)} & Deep equality (for objects/arrays) \\
\texttt{toBeCloseTo(num,\ decimals)} & Floating-point comparison \\
\texttt{toBeTruthy()} & Check if truthy \\
\texttt{toBeFalsy()} & Check if falsy \\
\texttt{toBeNull()} & Check if null \\
\texttt{toBeUndefined()} & Check if undefined \\
\texttt{toContain(item)} & Array/string contains item \\
\texttt{toHaveLength(n)} & Check length \\
\texttt{toThrow(error)} & Check if function throws \\
\end{longtable}
}
\end{fullwidth}

\subsection{Skills Checklist}\label{skills-checklist-66}

\begin{itemize}
\tightlist
\item[$\square$]
  I can create and import modules using CommonJS.
\item[$\square$]
  I can use \texttt{Map} and \texttt{Set} for collections.
\item[$\square$]
  I can write and run Jest tests.
\item[$\square$]
  I can use common Jest matchers.
\end{itemize}

\begin{center}\rule{0.5\linewidth}{0.5pt}\end{center}

\section{Unit 10: Project-Based
Learning}\label{unit-10-project-based-learning-1}

\begin{center}\rule{0.5\linewidth}{0.5pt}\end{center}

\subsection{Lesson 10.1: Text-Based Adventure
Game}\label{lesson-10.1-text-based-adventure-game}

\begin{Shaded}
\begin{Highlighting}[]
\KeywordTok{class} Room \{
  \CodeFuncTok{constructor}(name\OperatorTok{,} description) \{
    \KeywordTok{this}\OperatorTok{.}\CodeTypeTok{name} \OperatorTok{=} name\OperatorTok{;}
    \KeywordTok{this}\OperatorTok{.}\CodeTypeTok{description} \OperatorTok{=} description\OperatorTok{;}
    \KeywordTok{this}\OperatorTok{.}\CodeTypeTok{exits} \OperatorTok{=} \{\}\OperatorTok{;}
    \KeywordTok{this}\OperatorTok{.}\CodeTypeTok{items} \OperatorTok{=} []\OperatorTok{;}
  \}
 
  \CodeFuncTok{addExit}(direction\OperatorTok{,} room) \{\}
    \KeywordTok{this}\OperatorTok{.}\CodeTypeTok{exits}[direction\OperatorTok{.}\CodeFuncTok{toLowerCase}()] \OperatorTok{=} room\OperatorTok{;}
  \}
 
  \CodeFuncTok{addItem}(item) \{
    \KeywordTok{this}\OperatorTok{.}\CodeTypeTok{items}\OperatorTok{.}\CodeFuncTok{push}(item)\OperatorTok{;}
  \}
 
  \CodeFuncTok{describe}() \{
    \KeywordTok{let} desc \OperatorTok{=} \StringTok{\textasciigrave{}{-}{-}{-} }\StringTok{\$\{}\KeywordTok{this}\OperatorTok{.}\CodeTypeTok{name}\StringTok{\}}\StringTok{ {-}{-}{-}}\StringTok{\textbackslash{}n$\{}\KeywordTok{this}\OperatorTok{.}\CodeTypeTok{description}\StringTok{\}\textbackslash{}n}\StringTok{\textasciigrave{}}\OperatorTok{;}
    \KeywordTok{if} (\KeywordTok{this}\OperatorTok{.}\CodeTypeTok{items}\OperatorTok{.}\CodeTypeTok{length} \OperatorTok{\textgreater{}} \CodeNumTok{0}) \{
      desc \OperatorTok{+=} \StringTok{\textasciigrave{}You see: }\StringTok{\$\{}\KeywordTok{this}\OperatorTok{.}\CodeTypeTok{items}\OperatorTok{.}\CodeFuncTok{join}(\StringTok{", "})\StringTok{\}\textbackslash{}n}\StringTok{\textasciigrave{}}\OperatorTok{;}
    \}
    \KeywordTok{const} exitList \OperatorTok{=} Object\OperatorTok{.}\CodeFuncTok{keys}(\KeywordTok{this}\OperatorTok{.}\CodeTypeTok{exits})\OperatorTok{.}\CodeFuncTok{join}(\StringTok{", "})\OperatorTok{;}
    desc \OperatorTok{+=} \StringTok{\textasciigrave{}Exits: }\StringTok{\$\{}exitList\StringTok{\}}\StringTok{\textasciigrave{}}\OperatorTok{;}
    \KeywordTok{return} desc\OperatorTok{;}
  \}
\}
 
\KeywordTok{class} Player \{
  \CodeFuncTok{constructor}(name\OperatorTok{,} startRoom) \{
    \KeywordTok{this}\OperatorTok{.}\CodeTypeTok{name} \OperatorTok{=} name\OperatorTok{;}
    \KeywordTok{this}\OperatorTok{.}\CodeTypeTok{currentRoom} \OperatorTok{=} startRoom\OperatorTok{;}
    \KeywordTok{this}\OperatorTok{.}\CodeTypeTok{inventory} \OperatorTok{=} []\OperatorTok{;}
    \KeywordTok{this}\OperatorTok{.}\CodeTypeTok{health} \OperatorTok{=} \CodeNumTok{100}\OperatorTok{;}
  \}
 
  \CodeFuncTok{move}(direction) \{
    \KeywordTok{const} nextRoom \OperatorTok{=} \KeywordTok{this}\OperatorTok{.}\CodeTypeTok{currentRoom}\OperatorTok{.}\CodeTypeTok{exits}[direction\OperatorTok{.}\CodeFuncTok{toLowerCase}()]\OperatorTok{;}
    \KeywordTok{if} (nextRoom) \{
      \KeywordTok{this}\OperatorTok{.}\CodeTypeTok{currentRoom} \OperatorTok{=} nextRoom\OperatorTok{;}
      \KeywordTok{return} \KeywordTok{true}\OperatorTok{;}
    \}
    \KeywordTok{return} \KeywordTok{false}\OperatorTok{;}
  \}
 
  \CodeFuncTok{takeItem}(itemName) \{
    \KeywordTok{const} index \OperatorTok{=} \KeywordTok{this}\OperatorTok{.}\CodeTypeTok{currentRoom}\OperatorTok{.}\CodeTypeTok{items}\OperatorTok{.}\CodeFuncTok{findIndex}(
      (item) \KeywordTok{=\textgreater{}} item\OperatorTok{.}\CodeFuncTok{toLowerCase}() \OperatorTok{===} itemName\OperatorTok{.}\CodeFuncTok{toLowerCase}()\OperatorTok{,}
    )\OperatorTok{;}
    \KeywordTok{if} (index \OperatorTok{!==} \OperatorTok{{-}}\CodeNumTok{1}) \{
      \KeywordTok{const} item \OperatorTok{=} \KeywordTok{this}\OperatorTok{.}\CodeTypeTok{currentRoom}\OperatorTok{.}\CodeTypeTok{items}\OperatorTok{.}\CodeFuncTok{splice}(index\OperatorTok{,} \CodeNumTok{1})[\CodeNumTok{0}]\OperatorTok{;}
      \KeywordTok{this}\OperatorTok{.}\CodeTypeTok{inventory}\OperatorTok{.}\CodeFuncTok{push}(item)\OperatorTok{;}
      \KeywordTok{return} \KeywordTok{true}\OperatorTok{;}
    \}
    \KeywordTok{return} \KeywordTok{false}\OperatorTok{;}
  \}
 
  \CodeFuncTok{getStatus}() \{
    \KeywordTok{const} inv \OperatorTok{=} \KeywordTok{this}\OperatorTok{.}\CodeTypeTok{inventory}\OperatorTok{.}\CodeTypeTok{length} \OperatorTok{\textgreater{}} \CodeNumTok{0} \OperatorTok{?} \KeywordTok{this}\OperatorTok{.}\CodeTypeTok{inventory}\OperatorTok{.}\CodeFuncTok{join}(\StringTok{", "}) \OperatorTok{:} \StringTok{"empty"}\OperatorTok{;}
    \KeywordTok{return} \StringTok{\textasciigrave{}}\StringTok{\$\{}\KeywordTok{this}\OperatorTok{.}\CodeTypeTok{name}\StringTok{\}}\StringTok{ | Health: }\StringTok{\$\{}\KeywordTok{this}\OperatorTok{.}\CodeTypeTok{health}\StringTok{\}}\StringTok{ | Inventory: }\StringTok{\$\{}inv\StringTok{\}}\StringTok{\textasciigrave{}}\OperatorTok{;}
  \}
\}
 
\KeywordTok{async} \KeywordTok{function} \CodeFuncTok{playGame}() \{
  \KeywordTok{const} \{ askQuestion\OperatorTok{,} close \} \OperatorTok{=} \CodeAccentTok{require}(\StringTok{"./input"})\OperatorTok{;}
 
  \CommentTok{// Build world}
  \KeywordTok{const} cave \OperatorTok{=} \KeywordTok{new} \CodeFuncTok{Room}(\StringTok{"Dark Cave"}\OperatorTok{,} \StringTok{"You are in a dimly lit cave."})\OperatorTok{;}
  \KeywordTok{const} tunnel \OperatorTok{=} \KeywordTok{new} \CodeFuncTok{Room}(
    \StringTok{"Stone Tunnel"}\OperatorTok{,}
    \StringTok{"A narrow tunnel runs north and south."}\OperatorTok{,}
  )\OperatorTok{;}
  \KeywordTok{const} clearing \OperatorTok{=} \KeywordTok{new} \CodeFuncTok{Room}(
    \StringTok{"Sunny Clearing"}\OperatorTok{,}
    \StringTok{"A bright clearing with tall trees."}\OperatorTok{,}
  )\OperatorTok{;}
 
  cave\OperatorTok{.}\CodeFuncTok{addExit}(\StringTok{"north"}\OperatorTok{,} tunnel)\OperatorTok{;}
  cave\OperatorTok{.}\CodeFuncTok{addItem}(\StringTok{"torch"})\OperatorTok{;}
  tunnel\OperatorTok{.}\CodeFuncTok{addExit}(\StringTok{"south"}\OperatorTok{,} cave)\OperatorTok{;}
  tunnel\OperatorTok{.}\CodeFuncTok{addExit}(\StringTok{"north"}\OperatorTok{,} clearing)\OperatorTok{;}
  clearing\OperatorTok{.}\CodeFuncTok{addExit}(\StringTok{"south"}\OperatorTok{,} tunnel)\OperatorTok{;}
  clearing\OperatorTok{.}\CodeFuncTok{addItem}(\StringTok{"magic key"})\OperatorTok{;}
 
  \KeywordTok{const} player \OperatorTok{=} \KeywordTok{new} \CodeFuncTok{Player}(\StringTok{"Hero"}\OperatorTok{,} cave)\OperatorTok{;}
 
  console\OperatorTok{.}\CodeFuncTok{log}(\StringTok{"=== Text Adventure ==="})\OperatorTok{;}
  console\OperatorTok{.}\CodeFuncTok{log}(\StringTok{"Commands: go [direction], take [item], inventory, look, quit}\StringTok{\textbackslash{}n}\StringTok{"})\OperatorTok{;}
 
  \KeywordTok{while} (\KeywordTok{true}) \{
    console\OperatorTok{.}\CodeFuncTok{log}(player\OperatorTok{.}\CodeTypeTok{currentRoom}\OperatorTok{.}\CodeFuncTok{describe}())\OperatorTok{;}
    console\OperatorTok{.}\CodeFuncTok{log}(player\OperatorTok{.}\CodeFuncTok{getStatus}())\OperatorTok{;}
 
    \KeywordTok{const} input \OperatorTok{=} (\KeywordTok{await} \CodeFuncTok{askQuestion}(\StringTok{"}\StringTok{\textbackslash{}n}\StringTok{\textgreater{} "}))\OperatorTok{.}\CodeFuncTok{trim}()\OperatorTok{.}\CodeFuncTok{toLowerCase}()\OperatorTok{;}
    \KeywordTok{const} [command\OperatorTok{,} \OperatorTok{...}argParts] \OperatorTok{=} input\OperatorTok{.}\CodeFuncTok{split}(\StringTok{/}\StringTok{\textbackslash{}s+}\StringTok{/})\OperatorTok{;}
    \KeywordTok{const} argument \OperatorTok{=} argParts\OperatorTok{.}\CodeFuncTok{join}(\StringTok{" "})\OperatorTok{;}
 
    console\OperatorTok{.}\CodeFuncTok{log}(\StringTok{""})\OperatorTok{;}
 
    \KeywordTok{switch} (command) \{
      \KeywordTok{case} \StringTok{"go"}\OperatorTok{:}
        \KeywordTok{if} (\OperatorTok{!}argument) \{
          console\OperatorTok{.}\CodeFuncTok{log}(\StringTok{"Go where?"})\OperatorTok{;}
        \} \KeywordTok{else} \KeywordTok{if} (\OperatorTok{!}player\OperatorTok{.}\CodeFuncTok{move}(argument)) \{
          console\OperatorTok{.}\CodeFuncTok{log}(\StringTok{"You can\textquotesingle{}t go that way."})\OperatorTok{;}
        \}
        \KeywordTok{break}\OperatorTok{;}
      \KeywordTok{case} \StringTok{"take"}\OperatorTok{:}
        \KeywordTok{if} (player\OperatorTok{.}\CodeFuncTok{takeItem}(argument)) \{
          console\OperatorTok{.}\CodeFuncTok{log}(\StringTok{\textasciigrave{}You picked up the }\StringTok{\$\{}argument\StringTok{\}}\StringTok{.\textasciigrave{}})\OperatorTok{;}
        \ \KeywordTok{else} \{
          console\OperatorTok{.}\CodeFuncTok{log}(\StringTok{\textasciigrave{}There\textquotesingle{}s no }\StringTok{\$\{}argument\StringTok{\}}\StringTok{ here.\textasciigrave{}})\OperatorTok{;}
        \}
        \KeywordTok{break}\OperatorTok{;}
      \KeywordTok{case} \StringTok{"inventory"}\OperatorTok{:}
        console\OperatorTok{.}\CodeFuncTok{log}(\StringTok{\textasciigrave{}Inventory: }\StringTok{\$\{}player\OperatorTok{.}\CodeTypeTok{inventory}\OperatorTok{.}\CodeFuncTok{join}(\StringTok{", "}) \OperatorTok{||} \StringTok{"empty"}\StringTok{\}}\StringTok{\textasciigrave{}})\OperatorTok{;}
        \KeywordTok{break}\OperatorTok{;}
      \KeywordTok{case} \StringTok{"look"}\OperatorTok{:}
        \CommentTok{// Description already shown at top}
        \KeywordTok{break}\OperatorTok{;}
      \KeywordTok{case} \StringTok{"quit"}\OperatorTok{:}
        console\OperatorTok{.}\CodeFuncTok{log}(\StringTok{"Thanks for playing!"})\OperatorTok{;}
        \CodeFuncTok{close}()\OperatorTok{;}
        \KeywordTok{return}\OperatorTok{;}
      \KeywordTok{default}\OperatorTok{:}
        console\OperatorTok{.}\CodeFuncTok{log}(\StringTok{\textasciigrave{}Unknown command: \textquotesingle{}}\StringTok{\$\{}command\StringTok{\}}\StringTok{\textquotesingle{}\textasciigrave{}})\OperatorTok{;}
    \}
  \}
\}
 
\CodeFuncTok{playGame}()\OperatorTok{;}
\end{Highlighting}
\end{Shaded}

\begin{center}\rule{0.5\linewidth}{0.5pt}\end{center}

\subsection{Lesson 10.2: Command-Line To-Do
Manager}\label{lesson-10.2-command-line-to-do-manager}

\begin{Shaded}
\begin{Highlighting}[]
\KeywordTok{const} fs \OperatorTok{=} \CodeAccentTok{require}(\StringTok{"fs"})\OperatorTok{;}
\KeywordTok{const} \{ askQuestion\}\OperatorTok{,} close \} \OperatorTok{=} \CodeAccentTok{require}(\StringTok{"./input"})\OperatorTok{;}
 
\KeywordTok{const} FILE\_PATH \OperatorTok{=} \StringTok{"todos.json"}\OperatorTok{;}
 
\KeywordTok{function} \CodeFuncTok{loadTodos}() \{
  \KeywordTok{try} \{
    \KeywordTok{const} data \OperatorTok{=} fs\OperatorTok{.}\CodeFuncTok{readFileSync}(FILE\_PATH\OperatorTok{,} \StringTok{"utf8"})\OperatorTok{;}
    \KeywordTok{return} JSON\OperatorTok{.}\CodeFuncTok{parse}(data)\OperatorTok{;}
  \} \KeywordTok{catch} (err) \{
    \KeywordTok{if} (err\OperatorTok{.}\CodeTypeTok{code} \OperatorTok{===} \StringTok{"ENOENT"}) \{
      \KeywordTok{return} []\OperatorTok{;} \CommentTok{// File doesn\textquotesingle{}t exist yet}
    \}
    \KeywordTok{if} (err \KeywordTok{instanceof} SyntaxError) \{
      console\OperatorTok{.}\CodeFuncTok{log}(\StringTok{"Warning: Corrupted todos file. Starting fresh."})\OperatorTok{;}
      \KeywordTok{return} []\OperatorTok{;}
    \}
    \KeywordTok{throw} err\OperatorTok{;}
  \}
\}
 
\KeywordTok{function} \CodeFuncTok{saveTodos}(todos) \{
  fs\OperatorTok{.}\CodeFuncTok{writeFileSync}(FILE\_PATH\OperatorTok{,} JSON\OperatorTok{.}\CodeFuncTok{stringify}(todos\OperatorTok{,} \KeywordTok{null}\OperatorTok{,} \CodeNumTok{2}))\OperatorTok{;}
\}
 
\KeywordTok{function} \CodeFuncTok{addTask}(todos\OperatorTok{,} description) \{
  todos\OperatorTok{.}\CodeFuncTok{push}(\{
    description\OperatorTok{,}
    \CodeTypeTok{completed}\OperatorTok{:} \KeywordTok{false}\OperatorTok{,}
    \CodeTypeTok{createdAt}\OperatorTok{:} \KeywordTok{new} Date()\OperatorTok{.}\CodeFuncTok{toISOString}()\OperatorTok{,}
  \})\OperatorTok{;}
  console\OperatorTok{.}\CodeFuncTok{log}(\StringTok{\textasciigrave{}Task added: "}\StringTok{\$\{}description\StringTok{\}}\StringTok{"\textasciigrave{}})\OperatorTok{;}
\}
 
\KeywordTok{function} \CodeFuncTok{viewTasks}(todos) \{
  \KeywordTok{if} (todos\OperatorTok{.}\CodeTypeTok{length} \OperatorTok{===} \CodeNumTok{0}) \{
    console\OperatorTok{.}\CodeFuncTok{log}(\StringTok{"No tasks in your list."})\OperatorTok{;}
    \KeywordTok{return}\OperatorTok{;}
  \}
  console\OperatorTok{.}\CodeFuncTok{log}(\StringTok{"{-}{-}{-} Your To{-}Do List {-}{-}{-}"})\OperatorTok{;}
  todos\OperatorTok{.}\CodeFuncTok{forEach}((task\OperatorTok{,} index) \KeywordTok{=\textgreater{}} \{
    \KeywordTok{const} status \OperatorTok{=} task\OperatorTok{.}\CodeTypeTok{completed} \OperatorTok{?} \StringTok{"[X]"} \OperatorTok{:} \StringTok{"[ ]"}\OperatorTok{;}
    console\OperatorTok{.}\CodeFuncTok{log}(\StringTok{\textasciigrave{}}\StringTok{\$\{}index \OperatorTok{+} \CodeNumTok{1}\StringTok{\}}\StringTok{. }\StringTok{\$\{}status\StringTok{\}}\StringTok{ }\StringTok{\$\{}task\OperatorTok{.}\CodeTypeTok{description}\StringTok{\}}\StringTok{\textasciigrave{}})\OperatorTok{;}
  )\OperatorTok{;}
  console\OperatorTok{.}\CodeFuncTok{log}(\StringTok{"{-}{-}{-}{-}{-}{-}{-}{-}{-}{-}{-}{-}{-}{-}{-}{-}{-}{-}{-}{-}{-}{-}{-}"})\OperatorTok{;}
\}
 
\KeywordTok{function} \CodeFuncTok{markComplete}(todos\OperatorTok{,} taskNumber) \{
  \KeywordTok{const} index \OperatorTok{=} \CodeAccentTok{parseInt}(taskNumber) \OperatorTok{{-}} \CodeNumTok{1}\OperatorTok{;}
  \KeywordTok{if} (index \OperatorTok{\textgreater{}=} \CodeNumTok{0} \OperatorTok{\&\&} index \OperatorTok{\textless{}} todos\OperatorTok{.}\CodeTypeTok{length}) \{
    todos[index]\OperatorTok{.}\CodeTypeTok{completed} \OperatorTok{=} \KeywordTok{true}\OperatorTok{;}
    console\OperatorTok{.}\CodeFuncTok{log}(\StringTok{\textasciigrave{}Completed: "}\StringTok{\$\{}todos[index]\OperatorTok{.}\CodeTypeTok{description}\StringTok{\}}\StringTok{"\textasciigrave{}})\OperatorTok{;}
  \ \}\KeywordTok{else} \{
    console\OperatorTok{.}\CodeFuncTok{log}(\StringTok{"Invalid task number."})\OperatorTok{;}
  \}
\}
 
\KeywordTok{function} \CodeFuncTok{deleteTask}(todos\OperatorTok{,} taskNumber) \{
  \KeywordTok{const} index \OperatorTok{=} \CodeAccentTok{parseInt}(taskNumber) \OperatorTok{{-}} \CodeNumTok{1}\OperatorTok{;}
  \KeywordTok{if} (index \OperatorTok{\textgreater{}=} \CodeNumTok{0} \OperatorTok{\&\&} index \OperatorTok{\textless{}} todos\OperatorTok{.}\CodeTypeTok{length}) \{
    \KeywordTok{const} removed \OperatorTok{=} todos\OperatorTok{.}\CodeFuncTok{splice}(index\OperatorTok{,} \CodeNumTok{1})[\CodeNumTok{0}]\OperatorTok{;}
    console\OperatorTok{.}\CodeFuncTok{log}(\StringTok{\textasciigrave{}Deleted: "}\StringTok{\$\{}removed\OperatorTok{.}\CodeTypeTok{description}\StringTok{\}}\StringTok{"\textasciigrave{}})\OperatorTok{;}
  \ \KeywordTok{else} \{\}
    console\OperatorTok{.}\CodeFuncTok{log}(\StringTok{"Invalid task number."})\OperatorTok{;}
  \}
\}
 
\KeywordTok{async} \KeywordTok{function} \CodeFuncTok{main}() \{
  \KeywordTok{let} todos \OperatorTok{=} \CodeFuncTok{loadTodos}()\OperatorTok{;}
  console\OperatorTok{.}\CodeFuncTok{log}(\StringTok{"=== To{-}Do Manager ==="})\OperatorTok{;}
 
  \KeywordTok{while} (\KeywordTok{true}) \{
    console\OperatorTok{.}\CodeFuncTok{log}(
      \StringTok{"}\StringTok{\textbackslash{}n}\StringTok{1. Add Task  2. View Tasks  3. Mark Complete  4. Delete  5. Quit"}\OperatorTok{,}
    )\OperatorTok{;}
    \KeywordTok{const} choice \OperatorTok{=} (\KeywordTok{await} \CodeFuncTok{askQuestion}(\StringTok{"Choice: "}))\OperatorTok{.}\CodeFuncTok{trim}()\OperatorTok{;}
 
    \KeywordTok{switch} (choice) \{
      \KeywordTok{case} \StringTok{"1"}\OperatorTok{:}
        \KeywordTok{const} desc \OperatorTok{=} \KeywordTok{await} \CodeFuncTok{askQuestion}(\StringTok{"Task description: "})\OperatorTok{;}
        \CodeFuncTok{addTask}(todos\OperatorTok{,} desc)\OperatorTok{;}
        \CodeFuncTok{saveTodos}(todos)\OperatorTok{;}
        \KeywordTok{break}\OperatorTok{;}
      \KeywordTok{case} \StringTok{"2"}\OperatorTok{:}
        \CodeFuncTok{viewTasks}(todos)\OperatorTok{;}
        \KeywordTok{break}\OperatorTok{;}
      \KeywordTok{case} \StringTok{"3"}\OperatorTok{:}
        \CodeFuncTok{viewTasks}(todos)\OperatorTok{;}
        \KeywordTok{const} completeNum \OperatorTok{=} \KeywordTok{await} \CodeFuncTok{askQuestion}(\StringTok{"Task number: "})\OperatorTok{;}
        \CodeFuncTok{markComplete}(todos\OperatorTok{,} completeNum)\OperatorTok{;}
        \CodeFuncTok{saveTodos}(todos)\OperatorTok{;}
        \KeywordTok{break}\OperatorTok{;}
      \KeywordTok{case} \StringTok{"4"}\OperatorTok{:}
        \CodeFuncTok{viewTasks}(todos)\OperatorTok{;}
        \KeywordTok{const} deleteNum \OperatorTok{=} \KeywordTok{await} \CodeFuncTok{askQuestion}(\StringTok{"Task number to delete: "})\OperatorTok{;}
        \CodeFuncTok{deleteTask}(todos\OperatorTok{,} deleteNum)\OperatorTok{;}
        \CodeFuncTok{saveTodos}(todos)\OperatorTok{;}
        \KeywordTok{break}\OperatorTok{;}
      \KeywordTok{case} \StringTok{"5"}\OperatorTok{:}
        console\OperatorTok{.}\CodeFuncTok{log}(\StringTok{"Goodbye!"})\OperatorTok{;}
        \CodeFuncTok{close}()\OperatorTok{;}
        \KeywordTok{return}\OperatorTok{;}
      \KeywordTok{default}\OperatorTok{:}
        console\OperatorTok{.}\CodeFuncTok{log}(\StringTok{"Invalid choice."})\OperatorTok{;}
    \}
  \}
\}
 
\CodeFuncTok{main}()\OperatorTok{;}
\end{Highlighting}
\end{Shaded}

\begin{center}\rule{0.5\linewidth}{0.5pt}\end{center}

\section{Unit 11: Capstone Project --- Personal Data Manager in
JavaScript}\label{unit-11-capstone-project-personal-data-manager-in-javascript}

\subsection{Overview}\label{overview-13}

Build a fully functional, accessible command-line personal data manager
using Node.js.

\subsection{Project Options}\label{project-options}

\begin{itemize}
\tightlist
\item
  \textbf{Contact Manager:} name, phone, email, notes
\item
  \textbf{Book Tracker:} title, author, genre, read status
\item
  \textbf{Recipe Collection:} name, ingredients, instructions, cuisine
  type
\item
  \textbf{Budget Tracker:} category, amount, date, description
\end{itemize}

\subsection{Required Features}\label{required-features}

\begin{enumerate}
\def\labelenumi{\arabic{enumi}.}
\tightlist
\item
  \textbf{Add item} --- prompt for all attributes and save.
\item
  \textbf{View all items} --- display clearly with numbers.
\item
  \textbf{Search/filter} --- find items by keyword.
\item
  \textbf{Edit item} --- select and modify specific fields.
\item
  \textbf{Delete item} --- remove by number.
\item
  \textbf{Persistence} --- save/load to JSON file.
\end{enumerate}

\subsection{Project Structure}\label{project-structure}

\begin{verbatim}
my-data-manager/
|--- src/
|   |--- models/
|   |   +--- contact.js      (data model/class)
|   |--- storage/
|   |   +--- dataManager.js  (JSON file I/O)
|   |--- ui/
|   |   +--- consoleUI.js    (menu and display)
|   +--- utils/
|       +--- input.js        (reusable input helper)
|--- data/
|   +--- contacts.json       (data storage)
|--- tests/
|   +--- contact.test.js     (Jest tests)
|--- package.json
|--- README.md
+--- index.js                (entry point)
\end{verbatim}

\subsection{Accessibility Guidelines}\label{accessibility-guidelines}

\begin{enumerate}
\def\labelenumi{\arabic{enumi}.}
\tightlist
\item
  Every prompt must clearly describe what input is expected.
\item
  Print a confirmation message after every action.
\item
  Number all listed items for easy selection.
\item
  Provide clear, specific error messages.
\item
  Include a \texttt{help} command listing all available options.
\item
  Use consistent separators (\texttt{-\/-\/-}) to group output.
\end{enumerate}

\subsection{Assessment Criteria}\label{assessment-criteria-1}

\subsection{Functionality (60\%)}\label{functionality-60-1}

\begin{itemize}
\tightlist
\item[$\square$]
  All CRUD operations (Create, Read, Update, Delete) work correctly.
\item[$\square$]
  Data persists to JSON between program runs.
\item[$\square$]
  Search works by at least one field.
\item[$\square$]
  Program handles invalid input without crashing.
\end{itemize}

\subsection{Code Quality (20\%)}\label{code-quality-20}

\begin{itemize}
\tightlist
\item[$\square$]
  Functions are small with single responsibilities.
\item[$\square$]
  Consistent use of \texttt{const}/\texttt{let}, arrow functions,
  template literals.
\item[$\square$]
  Meaningful names for all variables and functions.
\item[$\square$]
  Code is organized into modules.
\item[$\square$]
  Includes JSDoc comments for key functions.
\end{itemize}

\subsection{Accessibility \& UX (15\%)}\label{accessibility-ux-15}

\begin{itemize}
\tightlist
\item[$\square$]
  All prompts are descriptive and clear.
\item[$\square$]
  Confirmation messages after every action.
\item[$\square$]
  Consistent, predictable output format.
\item[$\square$]
  Help system available.
\end{itemize}

\subsection{Documentation \& Tests
(5\%)}\label{documentation-tests-5}

\begin{itemize}
\tightlist
\item[$\square$]
  \texttt{README.md} with setup and usage instructions.
\item[$\square$]
  At least 5 Jest unit tests for key functions.
\end{itemize}

\subsection{Deliverables}\label{deliverables-1}

\begin{enumerate}
\def\labelenumi{\arabic{enumi}.}
\tightlist
\item
  All source code files.
\item
  \texttt{data/} folder with example JSON data.
\item
  \texttt{README.md} including:

  \begin{itemize}
  \tightlist
  \item
    Installation steps: \texttt{npm\ install}
  \item
    Run instructions: \texttt{node\ index.js}
  \item
    Available commands list.
  \item
    Screen reader usage tips.
  \item
    Reflection on challenges faced.
  \end{itemize}
\end{enumerate}

\cleardoublepage

\part{JAVA Curriculum}

\FloatBarrier\chapter*{Introduction to Computer Science with Java}

\subsection{For Visually Impaired
Students}\label{for-visually-impaired-students-1}

\subsection{Course Overview}\label{course-overview-1}

This curriculum introduces computer science concepts using Java, adapted
for students who rely on screen readers. All exercises are text-based
and terminal-focused. No graphical output is required.

\textbf{Target Audience:} Students with no prior programming experience
using screen readers (NVDA, JAWS, VoiceOver, or Windows Narrator).

\begin{center}\rule{0.5\linewidth}{0.5pt}\end{center}

\subsection{Tools Required}\label{tools-required-1}

\begin{itemize}
\tightlist
\item
  \textbf{Java Development Kit (JDK) 17+} ---
  \url{https://adoptium.net/}
\item
  \textbf{Visual Studio Code} with the \textbf{Extension Pack for Java}
  by Microsoft
\item
  A screen reader (NVDA, JAWS, or VoiceOver)
\end{itemize}

\subsection{Accessibility Notes}\label{accessibility-notes-1}

\begin{itemize}
\tightlist
\item
  All program output is text-based (console/terminal).
\item
  VS Code's Screen Reader Optimized Mode works well with Java.
\item
  Configure your screen reader to announce indentation changes --- Java
  uses braces \texttt{\{\}} to define blocks, but consistent indentation
  is still essential for readability.
\item
  Use \texttt{Ctrl\ +\ Shift\ +\ P} -\textgreater{} ``Toggle Screen
  Reader Optimized'' in VS Code.
\item
  Use \texttt{F8} / \texttt{Shift+F8} to navigate between errors.
\end{itemize}

\begin{center}\rule{0.5\linewidth}{0.5pt}\end{center}

\subsection{Units}\label{units-1}

\begin{fullwidth}{\def\LTcaptype{none} % do not increment counter
\begin{longtable}[]{@{}ll@{}}
\toprule\noalign{}
Unit & Title \\
\midrule\noalign{}
\endhead
\bottomrule\noalign{}
\endlastfoot
Unit 0 & Setting Up Your Java Development Environment \\
Unit 1 & Foundations of Computing and Getting Started with Java \\
Unit 2 & Writing Simple Programs and Numeric Computations \\
Unit 3 & Control Flow --- Decisions and Loops \\
Unit 4 & Data Structures --- Arrays and ArrayLists \\
Unit 5 & User Input and String Processing \\
Unit 6 & Methods (Functions) \\
Unit 7 & Classes and Object-Oriented Programming \\
Unit 8 & Files and Exception Handling \\
Unit 9 & Collections and Basic Data Analysis \\
Unit 10 & Project-Based Learning \\
Unit 11 & Capstone Project \\
\end{longtable}
}
\end{fullwidth}

\begin{center}\rule{0.5\linewidth}{0.5pt}\end{center}

\subsection{References}\label{references-1}

\begin{itemize}
\tightlist
\item
  Eck, D.J. (2022). \emph{Introduction to Programming Using Java} (9th
  ed.). Free PDF available at \url{https://math.hws.edu/javanotes/}
\item
  Horstmann, C.S. (2018). \emph{Big Java: Early Objects} (6th ed.).
  Wiley.
\end{itemize}

\newpage

\section{Teacher's Guide: Introduction to Computer Science with
Java}\label{teachers-guide-introduction-to-computer-science-with-java}

\subsection{For Students with Visual Impairments --- Complete Instructor
Edition}\label{for-students-with-visual-impairments-complete-instructor-edition-2}

\begin{center}\rule{0.5\linewidth}{0.5pt}\end{center}

\subsection{Table of Contents}\label{table-of-contents-2}

\begin{enumerate}
\def\labelenumi{\arabic{enumi}.}
\tightlist
\item
  \hyperref[philosophy-and-why-java]{Philosophy and Why Java}
\item
  \hyperref[java-specific-accessibility-challenges]{Java-Specific
  Accessibility Challenges}
\item
  \hyperref[classroom-setup-and-configuration]{Classroom Setup and
  Configuration}
\item
  \hyperref[spoken-java-conventions]{Spoken Java Conventions}
\item
  \hyperref[unit-01-setup-and-foundations]{Unit 0--1: Setup and
  Foundations}
\item
  \hyperref[unit-2-io-and-the-math-class]{Unit 2: I/O and the Math
  Class}
\item
  \hyperref[unit-3-control-flow]{Unit 3: Control Flow}
\item
  \hyperref[unit-4-arrays-and-arraylists]{Unit 4: Arrays and ArrayLists}
\item
  \hyperref[unit-5-strings-and-string-processing]{Unit 5: Strings and
  String Processing}
\item
  \hyperref[unit-6-methods]{Unit 6: Methods}
\item
  \hyperref[unit-7-classes-and-oop]{Unit 7: Classes and OOP}
\item
  \hyperref[unit-8-files-and-exception-handling]{Unit 8: Files and
  Exception Handling}
\item
  \hyperref[unit-9-collections-analysis-and-testing]{Unit 9:
  Collections, Analysis, and Testing}
\item
  \hyperref[units-1011-projects-and-capstone]{Units 10--11: Projects and
  Capstone}
\item
  \hyperref[bridging-from-python-to-java]{Bridging from Python to Java}
\item
  \hyperref[assessment-strategies]{Assessment Strategies}
\item
  \hyperref[supporting-diverse-needs]{Supporting Diverse Needs}
\item
  \hyperref[alternative-and-supplementary-resources]{Alternative and
  Supplementary Resources}
\item
  \hyperref[quick-reference-sheets]{Quick Reference Sheets}
\end{enumerate}

\begin{center}\rule{0.5\linewidth}{0.5pt}\end{center}

\subsection{1. Philosophy and Why Java}\label{philosophy-and-why-java}

\subsection{The Case for Teaching Java to Blind
Students}\label{the-case-for-teaching-java-to-blind-students}

Java is not the easiest language for beginners, and its verbosity
creates overhead that requires specific pedagogical attention. Despite
this, Java has genuine advantages for visually impaired learners that
are often overlooked:

\textbf{Compile-time error detection.} Java's compiler catches type
errors, name errors, and structural mistakes before any code runs. A
blind student does not need to search through runtime output for subtle
type mismatches --- the compiler announces them precisely and
immediately with file name, line number, and error description.
\marginnote{Java's compiler output follows a consistent, parseable format:}\amn{\textbf{Accessibility Tip:} Java compiler errors follow the format \texttt{Filename.java:line: error: message}. Navigate to the terminal after compilation, read the first error (line number and message), use \texttt{Ctrl+G} in VS Code to jump to that line, fix it, then recompile. Address one error at a time---fixing the first often resolves several others.}
\marginnote{Java's compiler output follows a consistent, parseable format: \texttt{filename.java:linenum: error: message}. This predictability makes it far easier to navigate with a screen reader than Python's multi-line tracebacks or JavaScript's runtime surprises. Teach students the compiler output format explicitly on day one.}

\textbf{Strong, static typing is auditorily informative.} When
IntelliSense in VS Code suggests methods, it includes the return type:
``\texttt{double\ getBalance()} --- returns a double.''
\sidenote{For blind students, typed IntelliSense completions are a form of accessible API documentation. When typing \texttt{account.}, the student hears method names \emph{and} their signatures: ``getBalance, returns double.'' This is information a sighted programmer gets from hovering---converted to audio automatically.} This type
information helps a screen reader user understand what a method does
without having to navigate to its definition.

\textbf{Enforced code structure.} Every Java program requires a class;
every executable program requires a \texttt{main} method with a specific
signature. This structure is more overhead to type but less ambiguity
about how programs are organized. For a student who cannot visually scan
a file for structure, predictable structural conventions reduce
cognitive load.

\textbf{Exception specification.} Java's checked exceptions force
programmers to declare which exceptions a method can throw. This means a
screen reader user reading a method signature gets information about
failure modes that Python provides only through documentation.

\textbf{Enterprise relevance.} Java remains a dominant language in
enterprise software, Android development, and many backend systems. A
student who learns Java learns patterns that transfer directly to
professional employment.

\begin{center}\rule{0.5\linewidth}{0.5pt}\end{center}

\subsection{2. Java-Specific Accessibility
Challenges}\label{java-specific-accessibility-challenges}

\subsection{The Verbosity Problem}\label{the-verbosity-problem}

Java requires significantly more text than Python to accomplish the same
task:

\begin{Shaded}
\begin{Highlighting}[]
\CommentTok{// Java: Hello, World}
\KeywordTok{public} \KeywordTok{class} HelloWorld \OperatorTok{\{}
    \KeywordTok{public} \CodeTypeTok{static} \CodeTypeTok{void} \CodeFuncTok{main}\OperatorTok{(}String\OperatorTok{[]} args\OperatorTok{)} \OperatorTok{\{}
        System\OperatorTok{.}\CodeFuncTok{out}\OperatorTok{.}\CodeFuncTok{println}\OperatorTok{(}\StringTok{"Hello, World!"}\OperatorTok{);}
    \OperatorTok{\}}
\OperatorTok{\}}
\end{Highlighting}
\end{Shaded}

\begin{Shaded}
\begin{Highlighting}[]
\CommentTok{\# Python: Hello, World}
print(\StringTok{"Hello, World!"})
\end{Highlighting}
\end{Shaded}

For a sighted student, extra text is a visual annoyance. For a blind
student, it is additional audio before they reach the meaningful code.
\textbf{Mitigation strategies:}

\begin{enumerate}
\def\labelenumi{\arabic{enumi}.}
\item
  \textbf{Code templates.} Provide a starter file the student copies for
  every new program. They modify only the method body for early lessons.
\sidenote{A single \texttt{Starter.java} template that the student renames and modifies is far more accessible than typing the class/main boilerplate from scratch. Create the template on day one, commit it to git, and have students \texttt{cp Starter.java NewProgram.java} for every exercise.}
  This reduces the typing burden and prevents structural errors while
  the student is learning concepts.
\item
  \textbf{JShell (Java REPL).} Java includes an interactive shell since
  Java 9: run \texttt{jshell} in the terminal. JShell strips away class
  and method boilerplate. For exploration and experimentation, JShell is
  far more accessible than creating files:

\begin{verbatim}
jshell> int x = 5;
x ==> 5
jshell> System.out.println(x * 2);
10
\end{verbatim}
\item
  \textbf{Progressive disclosure.} In early lessons, treat
  \texttt{public\ static\ void\ main(String{[}{]}\ args)\ \{} as a
  formula, not something to understand. something to understand.
Students learn to recognize it
\marginnote{``Formula before explanation'' is a proven strategy in VI programming education. The student hears the boilerplate, learns to type it, builds muscle memory, and only later learns what each word means. Attempting to explain \texttt{public static void main} before the student has written a working program wastes instructional time that should go to concepts.}
  by ear. Explain each word only when the concept it relates to (access
  modifiers, static methods, arrays) is being taught.
\end{enumerate}

\subsection{Integer Division
Surprises}\label{integer-division-surprises}

Java's \texttt{/} operator performs integer division when both operands
are integers:

\begin{itemize}
\tightlist
\item
  \texttt{15\ /\ 4} = \texttt{3} (not \texttt{3.75})
\item
  \texttt{15.0\ /\ 4} = \texttt{3.75}
\item
  \texttt{(double)\ 15\ /\ 4} = \texttt{3.75}
\end{itemize}

This is silent --- there is no error or warning. A blind student may not
catch a calculation that silently discards its decimal portion. catch a calculation that silently discards its decimal portion.
Teach
the rule explicitly:
\marginnote{Java's integer division is in the same category as JavaScript's type coercion:}\amn{\textbf{Accessibility Tip:} When dividing in Java, if you want a decimal result, cast at least one operand: \texttt{(double) numerator / denominator}. Integer division silently discards the decimal portion with no warning. Test every division expression in JShell first to verify it produces the expected type of result.}
\marginnote{Java's integer division is in the same category as JavaScript's type coercion: a silent wrong-answer failure with no error. Establish the rule as a mantra: ``If you want a decimal, cast one operand to \texttt{double}.'' Require students to say it aloud when writing any division expression in early units.} \textbf{if you want a decimal result, at least one
operand must be a \texttt{double}.}

\subsection{The String Equality Trap}\label{the-string-equality-trap}

\texttt{==} on Strings compares memory addresses, not content. Two
strings with identical characters in different memory locations are
\texttt{==} false but \texttt{.equals()} true. This never appears in
Python; blind students coming from Python are not conditioned to suspect
\texttt{==} as a source of String bugs.

\textbf{Teaching approach:} Show the broken behavior explicitly. Use
\texttt{new\ String("hello")\ ==\ new\ String("hello")} which is
guaranteed to be \texttt{false} even though the strings are identical.
Then show \texttt{.equals()} which is \texttt{true}. Make
\texttt{.equals()} for Strings a non-negotiable rule from day one.
\marginnote{The \texttt{==} vs \texttt{.equals()} distinction is Java's most common beginner trap.}\amn{\textbf{Accessibility Tip:} In Java, always use \texttt{str1.equals(str2)} to compare strings, never \texttt{str1 == str2}. The \texttt{==} operator compares memory addresses, not content---two strings with identical characters may be \texttt{==} false but \texttt{.equals()} true. This bug is silent and produces wrong output with no error message.}
\marginnote{The \texttt{==} vs \texttt{.equals()} distinction is Java's most common beginner trap. Python students are especially vulnerable because Python's \texttt{==} works correctly on strings. Frame it as: ``In Python, \texttt{==} compares values. In Java, \texttt{==} compares identity---memory address. For Strings, always use \texttt{.equals()}.'' Then never allow \texttt{==} on a String in student code.}

\subsection{Null Pointer Exceptions}\label{null-pointer-exceptions}

\texttt{NullPointerException} is the most common Java runtime error. It
occurs when you call a method or access a field on a variable that holds
\texttt{null} (no object reference). Unlike Python's \texttt{AttributeError}
which gives useful context, Java's NPE message is often terse. Teaching
students to recognize and debug NPEs is essential.
\sidenote{A \texttt{NullPointerException} with no message is among the least accessible errors in Java.}\amn{\textbf{Accessibility Tip:} Before calling any method on an object reference, check that it is not null: \texttt{if (obj != null) \{ obj.method(); \}}. A \texttt{NullPointerException} crashes the program immediately and is easy to miss in terminal output if you have navigated away. Null checks on any variable received from user input, file reads, or method returns prevent this entirely.}
\sidenote{A \texttt{NullPointerException} with no message is among the least accessible errors in Java. In Java~14+, NullPointerExceptions include a helpful message identifying which variable was null. Ensure students are using JDK~14 or later, and teach them to read the full stack trace rather than stopping at the exception name.}

\begin{center}\rule{0.5\linewidth}{0.5pt}\end{center}

\subsection{3. Classroom Setup and
Configuration}\label{classroom-setup-and-configuration-1}

\subsection{IDE Configuration}\label{ide-configuration}

\textbf{VS Code with Extension Pack for Java:} This is the recommended
setup. The Extension Pack includes Language Support for Java (Red Hat),
Debugger for Java, Java Test Runner, Maven for Java, and Project Manager
for Java.

Alternative IDEs to consider:

\begin{fullwidth}{\def\LTcaptype{none} % do not increment counter
\begin{longtable}[]{@{}
  >{\raggedright\arraybackslash}p{(\linewidth - 4\tabcolsep) * \real{0.2090}}
  >{\raggedright\arraybackslash}p{(\linewidth - 4\tabcolsep) * \real{0.0970}}
  >{\raggedright\arraybackslash}p{(\linewidth - 4\tabcolsep) * \real{0.6940}}@{}}
\toprule\noalign{}
\begin{minipage}[b]{\linewidth}\raggedright
IDE
\end{minipage} & \begin{minipage}[b]{\linewidth}\raggedright
Accessibility
\end{minipage} & \begin{minipage}[b]{\linewidth}\raggedright
Notes
\end{minipage} \\
\midrule\noalign{}
\endhead
\bottomrule\noalign{}
\endlastfoot
\textbf{VS Code + Extension Pack} & High & Recommended. Full
accessibility via keyboard. \\
\textbf{IntelliJ IDEA Community} & Medium--High & Powerful but complex.
More keyboard shortcuts to learn. JetBrains has an accessibility
guide. \\
\textbf{Eclipse} & Medium & Older UI; can be navigated but more complex
accessibility model. \\
\textbf{BlueJ} & Low--Medium & Designed for teaching; visual design
features may impede accessibility. \\
\textbf{DrJava} & Low--Medium & Simple; not actively maintained. \\
\textbf{jGRASP} & Low & Visualization-focused; primarily visual. \\
\textbf{Greenfoot} & Low & Visual/game-focused; not accessible for
screen reader use. \\
\end{longtable}
}
\end{fullwidth}

\textbf{VS Code Java Specific Settings:}

\begin{Shaded}
\begin{Highlighting}[]
\CodeFuncTok{\{}
  \CodeTypeTok{"java.completion.enabled"}\CodeFuncTok{:} \KeywordTok{true}\CodeFuncTok{,}
  \CodeTypeTok{"java.format.enabled"}\CodeFuncTok{:} \KeywordTok{true}\CodeFuncTok{,}
  \CodeTypeTok{"editor.formatOnSave"}\CodeFuncTok{:} \KeywordTok{true}\CodeFuncTok{,}
  \CodeTypeTok{"java.saveActions.organizeImports"}\CodeFuncTok{:} \KeywordTok{true}\CodeFuncTok{,}
  \CodeTypeTok{"java.completion.importOrder"}\CodeFuncTok{:} [\StringTok{"java"}, \StringTok{"javax"}, \StringTok{"com"}, \StringTok{""}]\CodeFuncTok{,}
  \CodeTypeTok{"editor.suggestOnTriggerCharacters"}\CodeFuncTok{:} \KeywordTok{true}
\CodeFuncTok{\}}
\end{Highlighting}
\end{Shaded}

 announce available methods.
This allows audio exploration of an object's
API.
\marginnote{For blind students, IntelliSense triggered by typing \texttt{.} is a form of audio API documentation.}\amn{\textbf{Accessibility Tip:} Type a variable name followed by a dot (\texttt{myList.}) and wait for IntelliSense to open. Use the arrow keys to navigate the completion list and hear each method name and its return type. This is how you explore what a Java object can do without memorizing or looking up documentation.}
\marginnote{For blind students, IntelliSense triggered by typing \texttt{.} is a form of audio API documentation. The student does not need to memorize method names---they can ``browse'' an object's interface by listening to the completion list. This is a first-class accessibility feature, not a shortcut.}

\subsection{JDK Installation
Verification}\label{jdk-installation-verification}

\begin{Shaded}
\begin{Highlighting}[]
java \CodeTypeTok{{-}{-}version}      \CommentTok{\# Runtime version}
javac \CodeTypeTok{{-}{-}version}     \CommentTok{\# Compiler version}
jshell              \CommentTok{\# Start REPL (Ctrl+D or /exit to quit)}
\end{Highlighting}
\end{Shaded}

All three must work before proceeding. If \texttt{jshell} is not found,
the student may have installed only the JRE (runtime), not the full JDK.
The JDK must be installed.

\begin{center}\rule{0.5\linewidth}{0.5pt}\end{center}

\subsection{4. Spoken Java Conventions}\label{spoken-java-conventions}

Extend the Python symbol conventions with Java-specific additions:

\begin{fullwidth}{\def\LTcaptype{none} % do not increment counter
\begin{longtable}[]{@{}ll@{}}
\toprule\noalign{}
Symbol/Pattern & Spoken Form \\
\midrule\noalign{}
\endhead
\bottomrule\noalign{}
\endlastfoot
\texttt{public} & ``public'' \\
\texttt{private} & ``private'' \\
\texttt{static} & ``static'' \\
\texttt{void} & ``void'' (returns nothing) \\
\texttt{new} & ``new'' \\
\texttt{this.} & ``this dot'' \\
\texttt{super.} & ``super dot'' \\
\texttt{@Override} & ``at-override annotation'' \\
\texttt{//\ comment} & ``double-slash comment'' \\
\texttt{/*\ ...\ */} & ``block comment'' \\
\texttt{/**\ ...\ */} & ``JavaDoc comment'' \\
\texttt{String{[}{]}} & ``String array'' \\
\texttt{ArrayList\textless{}String\textgreater{}} & ``ArrayList of
String'' \\
\texttt{HashMap\textless{}String,\ Integer\textgreater{}} & ``HashMap
from String to Integer'' \\
\texttt{throws\ IOException} & ``throws IOException'' \\
\texttt{System.out.println()} & ``sysout'' or ``system-out-println'' \\
\texttt{System.out.printf()} & ``system-out-printf'' \\
\texttt{Integer.parseInt()} & ``Integer-parse-int'' \\
\texttt{(int)\ x} & ``cast to int'' \\
\texttt{null} & ``null'' \\
\texttt{NullPointerException} & ``null pointer exception'' \\
\texttt{ArrayIndexOutOfBoundsException} & ``array index out of
bounds'' \\
\end{longtable}
}
\end{fullwidth}
\marginnote{``Sysout'' is the universal Java shorthand.}\amn{\textbf{Accessibility Tip:} In VS Code with the Java extension, type \texttt{sysout} and press Tab to expand it to \texttt{System.out.println()}. This live template saves keystrokes and reduces the chance of typing the verbose output statement incorrectly. Set up similar snippets for \texttt{Scanner} and \texttt{ArrayList} boilerplate.}
\marginnote{``Sysout'' is the universal Java shorthand. Most IDEs include a \texttt{sysout} live template that expands to \texttt{System.out.println()} on Tab. Teaching this shorthand reduces the cognitive overhead of Java's verbose output syntax.}

\textbf{Reading a method signature:}

\begin{Shaded}
\begin{Highlighting}[]
\KeywordTok{public} \CodeTypeTok{double} \CodeFuncTok{calculateArea}\OperatorTok{(}\CodeTypeTok{double} radius\OperatorTok{)} \KeywordTok{throws} IllegalArgumentException
\end{Highlighting}
\end{Shaded}

\emph{``public double method calculateArea, takes a double called
radius, may throw IllegalArgumentException''}

\begin{center}\rule{0.5\linewidth}{0.5pt}\end{center}

\subsection{5. Unit 0--1: Setup and
Foundations}\label{unit-01-setup-and-foundations-1}

\subsection{Lesson 0.1: Installing the
JDK}\label{lesson-0.1-installing-the-jdk}

\paragraph{Instructional Cues}\label{instructional-cues-21}

\begin{itemize}
\tightlist
\item
  \textbf{Use Adoptium (Eclipse Temurin) rather than Oracle JDK.} The
  Oracle JDK requires account creation and license agreement navigation. license agreement navigation.
\sidenote{Oracle's JDK download page requires navigating a login dialog}\amn{\textbf{Accessibility Tip:} Download the JDK from \texttt{adoptium.net} rather than Oracle's website. Adoptium provides direct download links with no account creation or licence acceptance dialog. Choose the LTS version (currently 21) for the most stable and long-supported option.}
\sidenote{Oracle's JDK download page requires navigating a login dialog and license acceptance UI that has historically been difficult with screen readers. Adoptium (\texttt{adoptium.net}) provides a straightforward download page with direct links to installers, no account required.}
  Adoptium is fully open-source and simpler to download.
\item
  \textbf{Verify the download before running.} Adoptium provides SHA
  checksums. Teaching students to verify downloads is a security habit
  worth establishing.
\item
   Ensure this is
  checked.
Many Java tools (Maven, Gradle, some IDEs) require
  \texttt{JAVA\_HOME} to be set.
\sidenote{\texttt{JAVA\_HOME} is an environment variable that tells build tools (Maven, Gradle, ant) where the JDK is installed. If it is not set, \texttt{mvn} and \texttt{gradle} commands will fail with cryptic errors. Verify with \texttt{echo \%JAVA\_HOME\%} (Windows) or \texttt{echo \$JAVA\_HOME} (macOS/Linux) immediately after installation.}
\end{itemize}

\paragraph{JShell as the First
Environment}\label{jshell-as-the-first-environment}

Before any file-based Java, spend time in JShell:

\begin{verbatim}
$ jshell
|  Welcome to JShell -- Version 17.0.x
|  For an introduction type: /help intro

jshell> 5 + 3
$1 ==> 8

jshell> "Hello " + "World"
$2 ==> "Hello World"

jshell> int x = 42;
x ==> 42

jshell> System.out.println("Value: " + x);
Value: 42
\end{verbatim}


\marginnote{JShell is Java's answer to the Python REPL.}\amn{\textbf{Accessibility Tip:} Start every new Java concept in JShell before writing a full class. Type an expression or statement and hear the result immediately. JShell eliminates the compile-run-navigate cycle, giving immediate audio feedback equivalent to the Python REPL.}
\marginnote{JShell is Java's answer to the Python REPL. For VI students, this is the recommended first Java environment for the same reason the Python REPL is: immediate audio feedback with no file navigation, compilation step, or run command.}

JShell commands to teach:

\begin{itemize}
\tightlist
\item
  \texttt{/list} --- show all code entered so far
\item
  \texttt{/edit} --- open in editor
\item
  \texttt{/vars} --- show all variables
\item
  \texttt{/methods} --- show all methods defined
\item
  \texttt{/exit} --- quit JShell
\end{itemize}

\paragraph{Thought Experiments}\label{thought-experiments-30}

\begin{enumerate}
\def\labelenumi{\arabic{enumi}.}
\item
  \emph{``Python is interpreted --- each line runs immediately. Java
  compiles to bytecode first, then runs. What does this mean for the
  kind of errors you encounter, and when? If you have a SyntaxError on
  line 50, in Python you might not discover it until the program reaches
  line 50. In Java, when do you discover it?''}
\item
  \emph{``The Java Virtual Machine (JVM) allows the same Java bytecode
  to run on Windows, macOS, and Linux. The JVM is different on each
  platform, but the bytecode is the same. What does this tell you about
  how Java achieves portability? What alternative approaches could a
  language designer take?''}
\item
  \emph{``JShell is a Read-Eval-Print Loop (REPL). Python has one;
  JavaScript has one; many languages have one. But Java didn't have one
  until 2016. Why might a language's designers add a REPL 20 years after
  the language was created? What did programmers do without it?''}
\item
  \emph{``Java was originally designed for embedded systems (like
  televisions), then became dominant for web servers, then for Android
  apps. Each time it succeeded in a new domain, it required extending
  the language. What does this history suggest about the relationship
  between a language's design and the problems it eventually solves?''}
\end{enumerate}

\begin{center}\rule{0.5\linewidth}{0.5pt}\end{center}

\subsection{Lesson 1.2: Your First Java
Program}\label{lesson-1.2-your-first-java-program}

\paragraph{Instructional Cues}\label{instructional-cues-22}

\textbf{Anatomy reading exercise --- before typing any code:} Read the
Hello World program aloud, symbol by symbol, using the pronunciation
table:

\begin{Shaded}
\begin{Highlighting}[]
\KeywordTok{public} \KeywordTok{class} HelloWorld \OperatorTok{\{}
    \KeywordTok{public} \CodeTypeTok{static} \CodeTypeTok{void} \CodeFuncTok{main}\OperatorTok{(}String\OperatorTok{[]} args\OperatorTok{)} \OperatorTok{\{}
        System\OperatorTok{.}\CodeFuncTok{out}\OperatorTok{.}\CodeFuncTok{println}\OperatorTok{(}\StringTok{"Hello, World!"}\OperatorTok{);}
    \OperatorTok{\}}
\OperatorTok{\}}
\end{Highlighting}
\end{Shaded}

\emph{``public class HelloWorld open-brace newline four-spaces public
static void main open-paren String array args close-paren open-brace
newline eight-spaces System dot out dot println open-paren quote Hello
comma World! close-quote close-paren semicolon newline four-spaces
close-brace newline close-brace''}

This exercise establishes pattern recognition for the boilerplate that
the student will type hundreds of times. Making it auditorily familiar
early reduces cognitive load later.

\textbf{Three \texttt{System.out} methods:}

\begin{Shaded}
\begin{Highlighting}[]
System\OperatorTok{.}\CodeFuncTok{out}\OperatorTok{.}\CodeFuncTok{println}\OperatorTok{(}\StringTok{"With newline at end"}\OperatorTok{);}     \CommentTok{// Most common}
System\OperatorTok{.}\CodeFuncTok{out}\OperatorTok{.}\CodeFuncTok{print}\OperatorTok{(}\StringTok{"No newline at end"}\OperatorTok{);}          \CommentTok{// For prompts}
System\OperatorTok{.}\CodeFuncTok{out}\OperatorTok{.}\CodeFuncTok{printf}\OperatorTok{(}\StringTok{"Pi is }\StringTok{\%.2f\%n}\StringTok{"}\OperatorTok{,} Math\OperatorTok{.}\CodeFuncTok{PI}\OperatorTok{);}    \CommentTok{// Formatted}
\end{Highlighting}
\end{Shaded}

For accessible programs, \texttt{printf} is extremely valuable --- it
produces consistent, predictable output formatting that reads cleanly
with a screen reader.

\paragraph{Thought Experiments}\label{thought-experiments-31}

\begin{enumerate}
\def\labelenumi{\arabic{enumi}.}
\item
  \emph{``Every Java program requires a \texttt{public\ class} whose
  name matches the file name. Why does Java enforce this one-to-one
  correspondence? What would break if you could put multiple public
  classes in one file?''}
\item
  \emph{``The \texttt{main} method signature is exactly
  \texttt{public\ static\ void\ main(String{[}{]}\ args)}. It cannot be
  \texttt{public\ void\ main()} or
  \texttt{static\ void\ main(String{[}{]}\ args)}. Why is this signature
  fixed? What does each word contribute?''}
\item
  \emph{``Java uses curly braces \texttt{\{\}} to delimit code blocks;
  Python uses indentation. From an accessibility perspective --- for a
  blind programmer --- which approach provides better information about
  structure when read aloud? Consider that braces can be read as `open
  block' and `close block,' while indentation changes require screen
  reader configuration.''}
\item
  \emph{``\texttt{System.out.println()} is long to type but
  descriptively named. Python's \texttt{print()} is short. What naming
  philosophy does each represent? When does conciseness become
  ambiguity?''}
\end{enumerate}

\begin{center}\rule{0.5\linewidth}{0.5pt}\end{center}

\subsection{Lesson 1.3: Variables and Primitive Data
Types}\label{lesson-1.3-variables-and-primitive-data-types}

\paragraph{Instructional Cues}\label{instructional-cues-23}

\textbf{Type declaration as specification:} Frame Java's type
declarations positively: \emph{``When you write
\texttt{int\ score\ =\ 85}, you're telling Java: this variable will
always hold a whole number. Java then protects you --- if you
accidentally try to put a name or a decimal into it, the compiler tells
you immediately.''}

\textbf{The eight primitive types --- prioritize four:} In this
curriculum, students need: \texttt{int}, \texttt{double},
\texttt{boolean}, \texttt{char}. The others (\texttt{byte},
\texttt{short}, \texttt{long}, \texttt{float}) exist but are rarely
needed in beginner programs. Introduce the four primary types fully
before acknowledging the others exist.

\textbf{Type casting demonstration:}

\begin{Shaded}
\begin{Highlighting}[]
\CommentTok{// Widening (automatic)}
\CodeTypeTok{int} i \OperatorTok{=} \CodeNumTok{100}\OperatorTok{;}
\CodeTypeTok{double} d \OperatorTok{=} i\OperatorTok{;}   \CommentTok{// 100 {-}\textgreater{} 100.0, automatic}
 
\CommentTok{// Narrowing (explicit, may lose data)}
\CodeTypeTok{double} pi \OperatorTok{=} \CodeNumTok{3.14159}\OperatorTok{;}
\CodeTypeTok{int} truncated \OperatorTok{=} \OperatorTok{(}\CodeTypeTok{int}\OperatorTok{)} pi\OperatorTok{;}  \CommentTok{// 3.14159 {-}\textgreater{} 3, data LOST}
 
System\OperatorTok{.}\CodeFuncTok{out}\OperatorTok{.}\CodeFuncTok{println}\OperatorTok{(}\StringTok{"Truncated: "} \OperatorTok{+} truncated\OperatorTok{);}  \CommentTok{// 3, not 4!}
\end{Highlighting}
\end{Shaded}

The truncation behavior (not rounding) is consistently surprising.
Demonstrate it with \texttt{(int)\ 3.9} producing \texttt{3}, and ask:
``Is this rounding or truncating? What is the difference? When would
truncation cause a bug?''

\paragraph{Non-Obvious Accessibility
Strategies}\label{non-obvious-accessibility-strategies-4}

\begin{itemize}
\item
  \textbf{\texttt{System.out.printf()} for consistent numeric output:}
  \texttt{printf("Balance:\ \$\%.2f\%n",\ balance)} always produces
  exactly 2 decimal places. Without formatting, \texttt{double} values
  may produce varying numbers of decimal places (\texttt{3.0},
  \texttt{3.14}, \texttt{3.1415926535...}) that read differently each
  time. Consistent formatting makes output more predictable.
\item
  \textbf{\texttt{Scanner.hasNext\_\_\_()} methods:} Before calling
  \texttt{scanner.nextInt()}, check \texttt{scanner.hasNextInt()}. If
  the user typed a non-integer, this returns \texttt{false} and prevents
  the exception, allowing graceful re-prompting. For a blind student who
  may not immediately catch an exception in terminal output, this is
  defensive programming with real accessibility value.
\end{itemize}

\paragraph{Thought Experiments}\label{thought-experiments-32}

\begin{enumerate}
\def\labelenumi{\arabic{enumi}.}
\item
  \emph{``Java's \texttt{int} holds values from about -2.1 billion to
  +2.1 billion. What happens if you exceed this limit (integer
  overflow)? The value silently wraps around to the opposite extreme.
  Name a real-world system where integer overflow has caused a
  documented failure.''}
\item
  \emph{``Java has \texttt{int} (primitive) and \texttt{Integer}
  (object). They hold the same values but behave differently.
  \texttt{int} cannot be null; \texttt{Integer} can. \texttt{int} cannot
  be used in a \texttt{List\textless{}int\textgreater{}} --- you need
  \texttt{List\textless{}Integer\textgreater{}}. Why would a language
  designer make this distinction? What does it cost? What does it
  protect against?''}
\item
  \emph{``Java's \texttt{char} holds a single Unicode character.
  \texttt{char\ c\ =\ \textquotesingle{}A\textquotesingle{}} stores the
  letter A. But \texttt{char} is also a numeric type ---
  \texttt{(int)\ \textquotesingle{}A\textquotesingle{}} is \texttt{65}.
  What does this dual nature (character and number) tell you about how
  computers store text internally?''}
\end{enumerate}

\begin{center}\rule{0.5\linewidth}{0.5pt}\end{center}

\subsection{6. Unit 2: I/O and the Math
Class}\label{unit-2-io-and-the-math-class}

\subsection{Lesson 2.2: Scanner for User
Input}\label{lesson-2.2-scanner-for-user-input}

\paragraph{Instructional Cues}\label{instructional-cues-24}

\textbf{The Scanner as a wrapper --- explain the layers:}

\begin{verbatim}
Keyboard
  V produces bytes
System.in (InputStream)
  V raw byte stream
new Scanner(System.in)
  V wraps with text-reading methods
scanner.nextLine(), scanner.nextInt(), etc.
\end{verbatim}

This layered model explains why \texttt{Scanner} is needed rather than
reading from \texttt{System.in} directly. It also foreshadows file
reading where \texttt{new\ Scanner(new\ File(...))} replaces
\texttt{new\ Scanner(System.in)}.

\textbf{\texttt{nextInt()} / \texttt{nextLine()} pitfall --- mandatory
demonstration:}

\begin{Shaded}
\begin{Highlighting}[]
Scanner scanner \OperatorTok{=} \KeywordTok{new} Scanner\OperatorTok{(}System\OperatorTok{.}\CodeFuncTok{in}\OperatorTok{);}
System\OperatorTok{.}\CodeFuncTok{out}\OperatorTok{.}\CodeFuncTok{print}\OperatorTok{(}\StringTok{"Enter your age: "}\OperatorTok{);}
\CodeTypeTok{int} age \OperatorTok{=} scanner\OperatorTok{.}\CodeFuncTok{nextInt}\OperatorTok{();}          \CommentTok{// Reads "25", leaves "\textbackslash{}n" in buffer}
System\OperatorTok{.}\CodeFuncTok{out}\OperatorTok{.}\CodeFuncTok{print}\OperatorTok{(}\StringTok{"Enter your name: "}\OperatorTok{);}
String name \OperatorTok{=} scanner\OperatorTok{.}\CodeFuncTok{nextLine}\OperatorTok{();}     \CommentTok{// Reads the leftover "\textbackslash{}n" {-}{-}{-} gets EMPTY string!}
System\OperatorTok{.}\CodeFuncTok{out}\OperatorTok{.}\CodeFuncTok{println}\OperatorTok{(}\StringTok{"Hello "} \OperatorTok{+} name \OperatorTok{+} \StringTok{", age "} \OperatorTok{+} age\OperatorTok{);}
\CommentTok{// Output: "Hello , age 25" {-}{-}{-} name is empty!}
\end{Highlighting}
\end{Shaded}

The fix:

\begin{Shaded}
\begin{Highlighting}[]
\CodeTypeTok{int} age \OperatorTok{=} scanner\OperatorTok{.}\CodeFuncTok{nextInt}\OperatorTok{();}
scanner\OperatorTok{.}\CodeFuncTok{nextLine}\OperatorTok{();}          \CommentTok{// Consume the leftover newline}
String name \OperatorTok{=} scanner\OperatorTok{.}\CodeFuncTok{nextLine}\OperatorTok{();}   \CommentTok{// Now reads properly}
\end{Highlighting}
\end{Shaded}

Run the broken version first. The student should hear the incorrect
output and understand why before seeing the fix.

\paragraph{Thought Experiments}\label{thought-experiments-33}

\begin{enumerate}
\def\labelenumi{\arabic{enumi}.}
\item
  \emph{``The \texttt{Scanner} reads from a `stream' of input. A stream
  is a sequence of data that flows from a source to a destination.
  \texttt{System.in} is the keyboard stream. What other streams exist in
  a typical program? (Think about file streams, network streams, the
  output stream.)''}
\item
  \emph{``After \texttt{scanner.nextInt()}, there's a leftover newline
  in the buffer. What IS a buffer? Why does the operating system buffer
  keyboard input rather than delivering each character the moment a key
  is pressed?''}
\item
  \emph{``In GUI (graphical) programs, users fill out forms with text
  fields and click buttons. There's no \texttt{scanner.nextLine()}. What
  mechanism do GUI programs use to collect user input? How does
  event-driven GUI input differ from sequential console input?''}
\end{enumerate}

\begin{center}\rule{0.5\linewidth}{0.5pt}\end{center}

\subsection{7. Unit 3: Control Flow}\label{unit-3-control-flow-1}

\subsection{Lessons 3.1--3.3}\label{lessons-3.13.3-1}

\paragraph{Instructional Cues}\label{instructional-cues-25}

\textbf{String comparison --- immediate practice:}

\begin{Shaded}
\begin{Highlighting}[]
String day \OperatorTok{=} \StringTok{"Wednesday"}\OperatorTok{;}
 
\CommentTok{// Wrong {-}{-}{-} compares memory addresses}
\KeywordTok{if} \OperatorTok{(}day \OperatorTok{==} \StringTok{"Wednesday"}\OperatorTok{)} \OperatorTok{\{}
    System\OperatorTok{.}\CodeFuncTok{out}\OperatorTok{.}\CodeFuncTok{println}\OperatorTok{(}\StringTok{"Middle of the week"}\OperatorTok{);}  \CommentTok{// May or may not print}
\OperatorTok{\}}
 
\CommentTok{// Correct {-}{-}{-} compares content}
\KeywordTok{if} \OperatorTok{(}day\OperatorTok{.}\CodeFuncTok{equals}\OperatorTok{(}\StringTok{"Wednesday"}\OperatorTok{))} \OperatorTok{\{}
    System\OperatorTok{.}\CodeFuncTok{out}\OperatorTok{.}\CodeFuncTok{println}\OperatorTok{(}\StringTok{"Middle of the week"}\OperatorTok{);}  \CommentTok{// Always correct}
\OperatorTok{\}}
 
\CommentTok{// Case{-}insensitive}
\KeywordTok{if} \OperatorTok{(}day\OperatorTok{.}\CodeFuncTok{equalsIgnoreCase}\OperatorTok{(}\StringTok{"wednesday"}\OperatorTok{))} \OperatorTok{\{}
    System\OperatorTok{.}\CodeFuncTok{out}\OperatorTok{.}\CodeFuncTok{println}\OperatorTok{(}\StringTok{"This works too"}\OperatorTok{);}
\OperatorTok{\}}
\end{Highlighting}
\end{Shaded}

Create a practice file with 10 comparisons, half using \texttt{==} and
half using \texttt{.equals()}. Ask the student to predict which ones
will work correctly before running.

\textbf{Switch expressions (Java 14+):}

Modern Java supports switch expressions that return a value:

\begin{Shaded}
\begin{Highlighting}[]
String dayType \OperatorTok{=} \KeywordTok{switch} \OperatorTok{(}day\OperatorTok{)} \OperatorTok{\{}
    \KeywordTok{case} \StringTok{"Monday"}\OperatorTok{,} \StringTok{"Tuesday"}\OperatorTok{,} \StringTok{"Wednesday"}\OperatorTok{,} \StringTok{"Thursday"}\OperatorTok{,} \StringTok{"Friday"} \OperatorTok{{-}\textgreater{}} \StringTok{"Weekday"}\OperatorTok{;}
    \KeywordTok{case} \StringTok{"Saturday"}\OperatorTok{,} \StringTok{"Sunday"} \OperatorTok{{-}\textgreater{}} \StringTok{"Weekend"}\OperatorTok{;}
    \KeywordTok{default} \OperatorTok{{-}\textgreater{}} \StringTok{"Unknown"}\OperatorTok{;}
\OperatorTok{\};}
System\OperatorTok{.}\CodeFuncTok{out}\OperatorTok{.}\CodeFuncTok{println}\OperatorTok{(}day \OperatorTok{+} \StringTok{" is a "} \OperatorTok{+} dayType\OperatorTok{);}
\end{Highlighting}
\end{Shaded}

This eliminates the fallthrough problem and is more concise. Teach both
traditional \texttt{switch} (for reading legacy code) and switch
expressions (for writing new code).

\paragraph{Thought Experiments}\label{thought-experiments-34}

\begin{enumerate}
\def\labelenumi{\arabic{enumi}.}
\item
  \emph{``Java uses \texttt{\&\&} for AND, \texttt{\textbar{}\textbar{}}
  for OR, \texttt{!} for NOT. Python uses the words \texttt{and},
  \texttt{or}, \texttt{not}. Which approach is more readable when heard
  aloud by a screen reader? Is there a meaningful accessibility argument
  for one over the other?''}
\item
  \emph{``The enhanced for loop \texttt{for\ (String\ item\ :\ list)} is
  simpler to write than the traditional
  \texttt{for\ (int\ i\ =\ 0;\ i\ \textless{}\ list.size();\ i++)}. But
  sometimes you need the index. Name three operations you can do with
  the traditional loop that you cannot do (directly) with the enhanced
  loop.''}
\item
  \emph{``Java's \texttt{switch} statement traditionally falls through
  from one case to the next unless you write \texttt{break}. This
  `feature' has caused countless bugs. Modern Java (14+) has switch
  expressions that don't fall through. What does this change tell you
  about how programming language design evolves in response to
  real-world usage patterns?''}
\item
  \emph{``Short-circuit evaluation:
  \texttt{if\ (obj\ !=\ null\ \&\&\ obj.getValue()\ \textgreater{}\ 0)}
  --- if \texttt{obj} is null, the second condition is never evaluated.
  This prevents a NullPointerException. How does this behavior make
  certain coding patterns possible that would otherwise crash?''}
\end{enumerate}

\begin{center}\rule{0.5\linewidth}{0.5pt}\end{center}

\subsection{8. Unit 4: Arrays and
ArrayLists}\label{unit-4-arrays-and-arraylists}

\subsection{Lesson 4.1: Arrays}\label{lesson-4.1-arrays-1}

\paragraph{Instructional Cues}\label{instructional-cues-26}

\textbf{The fixed-size limitation --- demonstrate, don't just describe:}

\begin{Shaded}
\begin{Highlighting}[]
String\OperatorTok{[]} fruits \OperatorTok{=} \KeywordTok{new} String\OperatorTok{[}\CodeNumTok{3}\OperatorTok{];}
fruits\OperatorTok{[}\CodeNumTok{0}\OperatorTok{]} \OperatorTok{=} \StringTok{"apple"}\OperatorTok{;}
fruits\OperatorTok{[}\CodeNumTok{1}\OperatorTok{]} \OperatorTok{=} \StringTok{"banana"}\OperatorTok{;}
fruits\OperatorTok{[}\CodeNumTok{2}\OperatorTok{]} \OperatorTok{=} \StringTok{"cherry"}\OperatorTok{;}
fruits\OperatorTok{[}\CodeNumTok{3}\OperatorTok{]} \OperatorTok{=} \StringTok{"date"}\OperatorTok{;}   \CommentTok{// ArrayIndexOutOfBoundsException!}
\end{Highlighting}
\end{Shaded}

Let the student hear the exception and analyze the message:
\texttt{Index\ 3\ out\ of\ bounds\ for\ length\ 3}. Ask: ``The message
gives you two numbers. What are they? What does each tell you?''

\textbf{\texttt{Arrays.toString()} as mandatory print method:}

The first time \texttt{print(fruits)} produces
\texttt{{[}Ljava.lang.String;@7a5d012b}, stop and explain: ``This is the
memory address of the array, not its contents. Java arrays don't have a
built-in `print nicely' behavior. Always use
\texttt{Arrays.toString(fruits)} or loop through manually.''

\textbf{The \texttt{for} vs \texttt{for-each} tradeoff --- make it
concrete:}

\begin{Shaded}
\begin{Highlighting}[]
\CommentTok{// for{-}each: simple, clear, but no index}
\KeywordTok{for} \OperatorTok{(}String fruit \OperatorTok{:} fruits\OperatorTok{)} \OperatorTok{\{}
    System\OperatorTok{.}\CodeFuncTok{out}\OperatorTok{.}\CodeFuncTok{println}\OperatorTok{(}fruit\OperatorTok{);}
\OperatorTok{\}}
 
\CommentTok{// traditional: more complex, but gives index}
\KeywordTok{for} \OperatorTok{(}\CodeTypeTok{int} i \OperatorTok{=} \CodeNumTok{0}\OperatorTok{;} i \OperatorTok{\textless{}} fruits\OperatorTok{.}\CodeFuncTok{length}\OperatorTok{;} i\OperatorTok{++)} \OperatorTok{\{}
    System\OperatorTok{.}\CodeFuncTok{out}\OperatorTok{.}\CodeFuncTok{println}\OperatorTok{(}i \OperatorTok{+} \StringTok{": "} \OperatorTok{+} fruits\OperatorTok{[}i\OperatorTok{]);}
\OperatorTok{\}}
 
\CommentTok{// The numbered version is more accessible!}
\end{Highlighting}
\end{Shaded}

For blind students, numbered output (showing both index and value) is
generally more useful than unnumbered output. Teach the traditional
\texttt{for} loop for array printing specifically because it enables
indexed output.

\paragraph{Thought Experiments}\label{thought-experiments-35}

\begin{enumerate}
\def\labelenumi{\arabic{enumi}.}
\item
  \emph{``Arrays are stored in contiguous memory --- all elements in a
  row, one after another. This is why accessing by index is fast: the
  address of element \texttt{i} is just the start address plus
  \texttt{i\ ×\ element\ size}. Why does this contiguous storage make
  inserting into the middle of an array slow? What must happen to all
  the elements after the insertion point?''}
\item
  \emph{``If you create
  \texttt{String{[}{]}\ arr\ =\ new\ String{[}5{]}}, what are the
  initial values of \texttt{arr{[}0{]}} through \texttt{arr{[}4{]}}? Now
  create \texttt{int{[}{]}\ nums\ =\ new\ int{[}5{]}} --- what are the
  initial values? Why are the defaults different between String arrays
  and int arrays?''}
\item
  \emph{``Two-dimensional arrays in Java are `arrays of arrays.' This
  means each row can have a different length --- a `jagged array.' When
  would a jagged array be more appropriate than a rectangular 2D array?
  Give a real-world data example.''}
\end{enumerate}

\begin{center}\rule{0.5\linewidth}{0.5pt}\end{center}

\subsection{Lesson 4.2: ArrayLists}\label{lesson-4.2-arraylists}

\paragraph{Instructional Cues}\label{instructional-cues-27}

\textbf{Autoboxing --- explain the invisibility:}

\begin{Shaded}
\begin{Highlighting}[]
ArrayList\OperatorTok{\textless{}}Integer\OperatorTok{\textgreater{}} numbers \OperatorTok{=} \KeywordTok{new} ArrayList\OperatorTok{\textless{}\textgreater{}();}
numbers\OperatorTok{.}\CodeFuncTok{add}\OperatorTok{(}\CodeNumTok{42}\OperatorTok{);}        \CommentTok{// 42 (int) is auto{-}boxed to Integer(42)}
\CodeTypeTok{int} first \OperatorTok{=} numbers\OperatorTok{.}\CodeFuncTok{get}\OperatorTok{(}\CodeNumTok{0}\OperatorTok{);}  \CommentTok{// Integer(42) is auto{-}unboxed to 42}
\end{Highlighting}
\end{Shaded}

This happens invisibly but has a performance cost. For the course, tell
students: ``Use \texttt{ArrayList\textless{}Integer\textgreater{}} when
you need a flexible list of integers. The boxing/unboxing is automatic
but slightly less efficient than an array of \texttt{int}.''

\textbf{Removal by index vs.~by value:}

\begin{Shaded}
\begin{Highlighting}[]
ArrayList\OperatorTok{\textless{}}String\OperatorTok{\textgreater{}} list \OperatorTok{=} \KeywordTok{new} ArrayList\OperatorTok{\textless{}\textgreater{}(}Arrays\OperatorTok{.}\CodeFuncTok{asList}\OperatorTok{(}\StringTok{"a"}\OperatorTok{,} \StringTok{"b"}\OperatorTok{,} \StringTok{"c"}\OperatorTok{,} \StringTok{"b"}\OperatorTok{,} \StringTok{"d"}\OperatorTok{));}
list\OperatorTok{.}\CodeFuncTok{remove}\OperatorTok{(}\CodeNumTok{1}\OperatorTok{);}            \CommentTok{// Removes by INDEX {-}{-}{-} removes "b" at position 1}
list\OperatorTok{.}\CodeFuncTok{remove}\OperatorTok{(}\StringTok{"b"}\OperatorTok{);}          \CommentTok{// Removes by VALUE {-}{-}{-} removes first "b" found}
list\OperatorTok{.}\CodeFuncTok{removeAll}\OperatorTok{(}Arrays\OperatorTok{.}\CodeFuncTok{asList}\OperatorTok{(}\StringTok{"b"}\OperatorTok{,} \StringTok{"c"}\OperatorTok{));}  \CommentTok{// Removes all matching values}
\end{Highlighting}
\end{Shaded}

Ask the student: ``If I call \texttt{list.remove(0)} on a list of
integers \texttt{{[}10,\ 20,\ 30{]}}, does this remove the element at
index 0 (which is \texttt{10}) or remove the value \texttt{0} (which
doesn't exist)? How does Java decide?'' (Answer: removes index 0,
because \texttt{0} is an \texttt{int}, not an \texttt{Integer}.)

\paragraph{Thought Experiments}\label{thought-experiments-36}

\begin{enumerate}
\def\labelenumi{\arabic{enumi}.}
\item
  \emph{``An \texttt{ArrayList} doubles its internal capacity when it
  runs out of space. If you start with capacity 16 and keep adding
  elements, when does reallocation occur? Trace the capacity: 16
  -\textgreater{} 32 -\textgreater{} 64 -\textgreater{} 128\ldots{} What
  is the total cost (in element moves) of adding 1000 elements to an
  initially empty ArrayList this way? Is this efficient?''}
\item
  \emph{``Java generics
  \texttt{ArrayList\textless{}String\textgreater{}} enforce type safety
  at compile time. Python lists can hold any types mixed together. When
  is a typed list (Java style) better? When is a mixed list (Python
  style) more flexible? Which approach produces more reliable
  programs?''}
\end{enumerate}

\begin{center}\rule{0.5\linewidth}{0.5pt}\end{center}

\subsection{9. Unit 5: Strings and String
Processing}\label{unit-5-strings-and-string-processing}

\paragraph{Instructional Cues}\label{instructional-cues-28}

\textbf{String immutability --- the Java version:}

In Java, Strings are immutable for the same reason as Python. But Java's
\texttt{String} object is also \textbf{interned} --- the JVM may reuse
the same \texttt{String} object for two variables holding the same
literal value. This is why \texttt{==} sometimes works for String
comparison (when both refer to the same interned object) but is not
reliable.

\textbf{\texttt{StringBuilder} motivation:}

\begin{Shaded}
\begin{Highlighting}[]
\CommentTok{// Inefficient {-}{-}{-} creates 100,000 String objects}
String result \OperatorTok{=} \StringTok{""}\OperatorTok{;}
\KeywordTok{for} \OperatorTok{(}\CodeTypeTok{int} i \OperatorTok{=} \CodeNumTok{0}\OperatorTok{;} i \OperatorTok{\textless{}} \CodeNumTok{100000}\OperatorTok{;} i\OperatorTok{++)} \OperatorTok{\{}
    result \OperatorTok{+=} \StringTok{"x"}\OperatorTok{;}
\OperatorTok{\}}
 
\CommentTok{// Efficient {-}{-}{-} modifies a single buffer}
StringBuilder sb \OperatorTok{=} \KeywordTok{new} StringBuilder\OperatorTok{();}
\KeywordTok{for} \OperatorTok{(}\CodeTypeTok{int} i \OperatorTok{=} \CodeNumTok{0}\OperatorTok{;} i \OperatorTok{\textless{}} \CodeNumTok{100000}\OperatorTok{;} i\OperatorTok{++)} \OperatorTok{\{}
    sb\OperatorTok{.}\CodeFuncTok{append}\OperatorTok{(}\StringTok{"x"}\OperatorTok{);}
\OperatorTok{\}}
String result \OperatorTok{=} sb\OperatorTok{.}\CodeFuncTok{toString}\OperatorTok{();}
\end{Highlighting}
\end{Shaded}

The \texttt{+} in a loop creates a new String object on every iteration,
then discards the old one. For small loops this is fine; for large loops
it's a significant performance problem. \texttt{StringBuilder} maintains
a mutable buffer.

\paragraph{Thought Experiments}\label{thought-experiments-37}

\begin{enumerate}
\def\labelenumi{\arabic{enumi}.}
\item
  \emph{``String concatenation with \texttt{+} is convenient but
  expensive in loops. In Python, the equivalent is also discouraged ---
  use \texttt{\textquotesingle{}\textquotesingle{}.join(list)} instead.
  What does this common advice across two languages tell you about how
  immutable strings interact with loops?''}
\item
  \emph{``Java's \texttt{String.format()} and \texttt{printf()} use
  format specifiers like \texttt{\%s}, \texttt{\%d}, \texttt{\%f}.
  Python uses f-strings with \texttt{\{\}}. What are the accessibility
  tradeoffs of each approach? Which reads more naturally when heard
  aloud? Which is easier to type?''}
\item
  \emph{``Regular expressions are a tiny language inside a programming
  language. \texttt{\textbackslash{}\textbackslash{}d+} matches one or
  more digits. \texttt{{[}a-zA-Z{]}+} matches one or more letters. Why
  do programmers use regular expressions despite their difficult-to-read
  syntax? What would the alternative be for complex pattern matching?''}
\end{enumerate}

\begin{center}\rule{0.5\linewidth}{0.5pt}\end{center}

\subsection{10. Unit 6: Methods}\label{unit-6-methods}

\paragraph{Instructional Cues}\label{instructional-cues-29}

\textbf{Method overloading --- auditory navigation:}

\begin{Shaded}
\begin{Highlighting}[]
\KeywordTok{public} \CodeTypeTok{static} \CodeTypeTok{int} \CodeFuncTok{add}\OperatorTok{(}\CodeTypeTok{int} a\OperatorTok{,} \CodeTypeTok{int} b\OperatorTok{)} \OperatorTok{\{} \KeywordTok{return} a \OperatorTok{+} b\OperatorTok{;} \OperatorTok{\}}
\KeywordTok{public} \CodeTypeTok{static} \CodeTypeTok{double} \CodeFuncTok{add}\OperatorTok{(}\CodeTypeTok{double} a\OperatorTok{,} \CodeTypeTok{double} b\OperatorTok{)} \OperatorTok{\{} \KeywordTok{return} a \OperatorTok{+} b\OperatorTok{;} \OperatorTok{\}}
\KeywordTok{public} \CodeTypeTok{static} \CodeTypeTok{int} \CodeFuncTok{add}\OperatorTok{(}\CodeTypeTok{int} a\OperatorTok{,} \CodeTypeTok{int} b\OperatorTok{,} \CodeTypeTok{int} c\OperatorTok{)} \OperatorTok{\{} \KeywordTok{return} a \OperatorTok{+} b \OperatorTok{+} c\OperatorTok{;} \OperatorTok{\}}
\end{Highlighting}
\end{Shaded}

When IntelliSense shows overloads in VS Code, pressing
\texttt{Shift+Alt+F2} (Accessible View) lists them all. Teach students
to use this to discover available overloads.

\textbf{\texttt{void} vs.~return type --- make the distinction
tangible:}

\begin{Shaded}
\begin{Highlighting}[]
\CommentTok{// void: for side effects only}
\KeywordTok{public} \CodeTypeTok{static} \CodeTypeTok{void} \CodeFuncTok{printGreeting}\OperatorTok{(}String name\OperatorTok{)} \OperatorTok{\{}
    System\OperatorTok{.}\CodeFuncTok{out}\OperatorTok{.}\CodeFuncTok{println}\OperatorTok{(}\StringTok{"Hello, "} \OperatorTok{+} name \OperatorTok{+} \StringTok{"!"}\OperatorTok{);}
\OperatorTok{\}}
 
\CommentTok{// Return type: for computation}
\KeywordTok{public} \CodeTypeTok{static} \CodeTypeTok{double} \CodeFuncTok{calculateTip}\OperatorTok{(}\CodeTypeTok{double} billAmount\OperatorTok{,} \CodeTypeTok{double} rate\OperatorTok{)} \OperatorTok{\{}
    \KeywordTok{return} billAmount \OperatorTok{*} rate\OperatorTok{;}
\OperatorTok{\}}
 
\CommentTok{// Using them:}
\CodeFuncTok{printGreeting}\OperatorTok{(}\StringTok{"Alice"}\OperatorTok{);}                      \CommentTok{// Standalone call}
\CodeTypeTok{double} tip \OperatorTok{=} \CodeFuncTok{calculateTip}\OperatorTok{(}\CodeNumTok{50.00}\OperatorTok{,} \CodeNumTok{0.18}\OperatorTok{);}      \CommentTok{// Must capture return value}
System\OperatorTok{.}\CodeFuncTok{out}\OperatorTok{.}\CodeFuncTok{printf}\OperatorTok{(}\StringTok{"Tip: $}\StringTok{\%.2f\%n}\StringTok{"}\OperatorTok{,} tip\OperatorTok{);}
\end{Highlighting}
\end{Shaded}

\emph{``\texttt{void} methods exist to DO something. Return-type methods
exist to CALCULATE something and give you the result. Ask: `Do I need to
use the output later in my code?' If yes, it needs a return type.''}

\paragraph{Thought Experiments}\label{thought-experiments-38}

\begin{enumerate}
\def\labelenumi{\arabic{enumi}.}
\item
  \emph{``Java methods can throw exceptions:
  \texttt{public\ double\ sqrt(double\ n)\ throws\ IllegalArgumentException}.
  This signature tells callers that the method might fail in a specific
  way. Python has no equivalent syntax. What is the programmer's
  obligation when they see \texttt{throws} in a method signature?''}
\item
  \emph{``Static methods belong to the class, not to any instance:
  \texttt{Math.sqrt(25)} --- you don't need a \texttt{Math} object.
  Instance methods belong to an object: \texttt{myString.length()}. What
  is the conceptual difference? When should a method be static
  vs.~instance?''}
\item
  \emph{``Method parameters in Java are `pass by value' for primitives
  and `pass by reference' for objects. For primitives: the method gets a
  copy, and modifying the parameter doesn't affect the original. For
  objects: the method gets the reference, and modifications to the
  object's fields DO affect the original. Design a test to confirm this
  behavior for both cases.''}
\item
  \_``Recursion: a method that calls itself.
  \texttt{int\ factorial(int\ n)\ \{\ return\ n\ ==\ 0\ ?\ 1\ :\ n\ \_\ factorial(n-1);\ \}}.
  Every recursive solution has an iterative equivalent. When is the
  recursive version clearer? What is the risk of deep recursion (many
  nested calls)?''*
\end{enumerate}

\begin{center}\rule{0.5\linewidth}{0.5pt}\end{center}

\subsection{11. Unit 7: Classes and OOP}\label{unit-7-classes-and-oop-1}

\subsection{Extended Instructional
Cues}\label{extended-instructional-cues-5}

\textbf{The mandatory \texttt{\_\_str\_\_} equivalent in Java ---
\texttt{toString()}:}

Every student-defined class must override \texttt{toString()}. Without
it, \texttt{System.out.println(myObject)} produces an unreadable string.
The \texttt{@Override} annotation makes this explicit:

\begin{Shaded}
\begin{Highlighting}[]
\CodeTypeTok{@Override}
\KeywordTok{public} String \CodeFuncTok{toString}\OperatorTok{()} \OperatorTok{\{}
    \KeywordTok{return} String\OperatorTok{.}\CodeFuncTok{format}\OperatorTok{(}\StringTok{"Dog[name=}\StringTok{\%s}\StringTok{, age=}\StringTok{\%d}\StringTok{]"}\OperatorTok{,} name\OperatorTok{,} age\OperatorTok{);}
\OperatorTok{\}}
\end{Highlighting}
\end{Shaded}

Require this in every class from Unit 7 onward. Explain: ``When your
screen reader reads \texttt{System.out.println(myDog)}, it reads
whatever \texttt{toString()} returns. If you don't define
\texttt{toString()}, you'll hear a memory address. That's useless.''

\textbf{JavaDoc comments as accessible documentation:}

Teach this format for every method:

\begin{Shaded}
\begin{Highlighting}[]
\CommentTok{/**}
 \CommentTok{*} Calculates the area of this circle\CommentTok{.}
 \CommentTok{*}
 \CommentTok{*} \CommentTok{@return} the area as a double\CommentTok{,} in square units
 \CommentTok{*} \CommentTok{@throws}\CommentTok{ }\CommentTok{IllegalArgumentException}\CommentTok{ }if radius is negative
 \CommentTok{*/}
\KeywordTok{public} \CodeTypeTok{double} \CodeFuncTok{getArea}\OperatorTok{()} \OperatorTok{\{}
    \KeywordTok{if} \OperatorTok{(}radius \OperatorTok{\textless{}} \CodeNumTok{0}\OperatorTok{)} \KeywordTok{throw} \KeywordTok{new} IllegalArgumentException\OperatorTok{(}\StringTok{"Radius cannot be negative"}\OperatorTok{);}
    \KeywordTok{return} Math\OperatorTok{.}\CodeFuncTok{PI} \OperatorTok{*} radius \OperatorTok{*} radius\OperatorTok{;}
\OperatorTok{\}}
\end{Highlighting}
\end{Shaded}

When IntelliSense shows this method in VS Code, pressing
\texttt{Shift+Alt+F2} reads the JavaDoc aloud. This makes the student's
own code navigable via audio --- a profound accessibility benefit.

\textbf{The VS Code symbol navigator:} \texttt{Ctrl+Shift+O} opens a
dropdown listing all methods in the current file. Type a few letters to
filter. This allows jumping to any method by name --- much more
efficient than arrow-key navigation through a long file.

\textbf{Access modifiers --- enforce from day one:}

\begin{fullwidth}{\def\LTcaptype{none} % do not increment counter
\begin{longtable}[]{@{}
  >{\raggedright\arraybackslash}p{(\linewidth - 4\tabcolsep) * \real{0.3864}}
  >{\raggedright\arraybackslash}p{(\linewidth - 4\tabcolsep) * \real{0.1250}}
  >{\raggedright\arraybackslash}p{(\linewidth - 4\tabcolsep) * \real{0.4886}}@{}}
\toprule\noalign{}
\begin{minipage}[b]{\linewidth}\raggedright
Modifier
\end{minipage} & \begin{minipage}[b]{\linewidth}\raggedright
Rule
\end{minipage} & \begin{minipage}[b]{\linewidth}\raggedright
Reason
\end{minipage} \\
\midrule\noalign{}
\endhead
\bottomrule\noalign{}
\endlastfoot
\texttt{private} on fields & Always & Encapsulation; prevents direct
manipulation \\
\texttt{public} on getters/setters & When needed & Controlled access \\
\texttt{private} on helper methods & Yes & Internal implementation
detail \\
\texttt{public} on methods the user calls & Yes & Interface to the
class \\
\end{longtable}
}
\end{fullwidth}

Establish this as a non-negotiable habit, not a preference.
\emph{``Private fields, public interface. Every time.''}

\paragraph{Thought Experiments}\label{thought-experiments-39}

\begin{enumerate}
\def\labelenumi{\arabic{enumi}.}
\item
  \emph{``Encapsulation means the internal state of an object is
  private. Only its public methods can access or modify it. Why would
  you want to prevent direct access to \texttt{account.balance}? What
  bugs could occur if balance were public?''}
\item
  \emph{``Inheritance creates a hierarchy. A \texttt{SavingsAccount}
  extends \texttt{BankAccount}. Every \texttt{SavingsAccount} is-a
  \texttt{BankAccount}, but not every \texttt{BankAccount} is-a
  \texttt{SavingsAccount}. What are the type-safety implications of
  this? If you have a \texttt{BankAccount} variable, can it hold a
  \texttt{SavingsAccount} object?''}
\item
  \emph{``Java interfaces define contracts --- a list of methods that
  any implementing class must provide. A class can implement multiple
  interfaces but can only extend one class. Why this asymmetry? What
  problem does single inheritance solve? What does it prevent?''}
\item
  \emph{``Design Pattern: the Factory Method creates objects without
  specifying the exact class. Instead of \texttt{new\ Dog()}, you'd call
  \texttt{AnimalFactory.create("dog")}. What is the advantage of this
  indirection? When would it be useful?''}
\item
  \emph{``Object-oriented programming was supposed to model the real
  world. A \texttt{Dog} is-a \texttt{Animal}. But real-world
  relationships are complex --- a \texttt{FlyingFish} is both a
  \texttt{Fish} and can fly like a \texttt{Bird}. How does Java's
  single-inheritance model handle this situation? What alternatives
  exist?''}
\item
  \emph{``Polymorphism means one interface, many implementations.
  \texttt{Animal{[}{]}\ animals\ =\ \{new\ Dog(),\ new\ Cat(),\ new\ Bird()\};\ for\ (Animal\ a\ :\ animals)\ \{\ a.speak();\ \}}
  --- each animal speaks differently. Why is this more flexible than
  checking \texttt{instanceof} and calling different methods
  explicitly?''}
\end{enumerate}

\paragraph{Alternative Resources}\label{alternative-resources-6}

\begin{itemize}
\tightlist
\item
  \textbf{``Head First Design Patterns'' by Freeman, Robson et al.:}
  Visual book, but audio-describable concepts. Covers fundamental OOP
  design patterns. For advanced students.
\item
  \textbf{``Effective Java'' by Joshua Bloch:} The definitive book on
  Java best practices. Dense but invaluable for understanding WHY Java
  is designed the way it is.
\item
  \textbf{Oracle Java Tutorials (docs.oracle.com/javase/tutorial/):}
  Comprehensive, plain HTML, fully accessible. Covers OOP concepts with
  extensive examples.
\end{itemize}

\begin{center}\rule{0.5\linewidth}{0.5pt}\end{center}

\subsection{12. Unit 8: Files and Exception
Handling}\label{unit-8-files-and-exception-handling}

\paragraph{Instructional Cues}\label{instructional-cues-30}

\textbf{Checked vs.~unchecked exceptions --- the core distinction:}

\begin{fullwidth}{\def\LTcaptype{none} % do not increment counter
\begin{longtable}[]{@{}
  >{\raggedright\arraybackslash}p{(\linewidth - 6\tabcolsep) * \real{0.1600}}
  >{\raggedright\arraybackslash}p{(\linewidth - 6\tabcolsep) * \real{0.4200}}
  >{\raggedright\arraybackslash}p{(\linewidth - 6\tabcolsep) * \real{0.1250}}
  >{\raggedright\arraybackslash}p{(\linewidth - 6\tabcolsep) * \real{0.2950}}@{}}
\toprule\noalign{}
\begin{minipage}[b]{\linewidth}\raggedright
Type
\end{minipage} & \begin{minipage}[b]{\linewidth}\raggedright
Examples
\end{minipage} & \begin{minipage}[b]{\linewidth}\raggedright
Must Handle?
\end{minipage} & \begin{minipage}[b]{\linewidth}\raggedright
Why?
\end{minipage} \\
\midrule\noalign{}
\endhead
\bottomrule\noalign{}
\endlastfoot
\textbf{Checked} & \texttt{IOException}, \texttt{FileNotFoundException},
\texttt{SQLException} & Yes --- compiler enforces & These are
recoverable; programmer should handle them \\
\textbf{Unchecked (RuntimeException)} & \texttt{NullPointerException},
\texttt{ArrayIndexOutOfBoundsException},
\texttt{IllegalArgumentException} & No --- optional & Usually indicate
programming errors; may not be recoverable \\
\end{longtable}
}
\end{fullwidth}

\emph{``Checked exceptions are like warning labels: `This operation can
fail. You must decide what to do if it does.' Unchecked exceptions are
usually bugs --- the program did something wrong.''}

\textbf{Try-with-resources --- teach first, manual close second:}

\begin{Shaded}
\begin{Highlighting}[]
\CommentTok{// Modern: try{-}with{-}resources (always teach this first)}
\KeywordTok{try} \OperatorTok{(}BufferedReader reader \OperatorTok{=} \KeywordTok{new} BufferedReader\OperatorTok{(}\KeywordTok{new} FileReader\OperatorTok{(}\StringTok{"data.txt"}\OperatorTok{)))} \OperatorTok{\{}
    String line\OperatorTok{;}
    \KeywordTok{while} \OperatorTok{((}line \OperatorTok{=} reader\OperatorTok{.}\CodeFuncTok{readLine}\OperatorTok{())} \OperatorTok{!=} \KeywordTok{null}\OperatorTok{)} \OperatorTok{\{}
        System\OperatorTok{.}\CodeFuncTok{out}\OperatorTok{.}\CodeFuncTok{println}\OperatorTok{(}line\OperatorTok{);}
    \OperatorTok{\}}
\OperatorTok{\}} \KeywordTok{catch} \OperatorTok{(}IOException e\OperatorTok{)} \OperatorTok{\{}
    System\OperatorTok{.}\CodeFuncTok{out}\OperatorTok{.}\CodeFuncTok{println}\OperatorTok{(}\StringTok{"Error reading file: "} \OperatorTok{+} e\OperatorTok{.}\CodeFuncTok{getMessage}\OperatorTok{());}
\OperatorTok{\}}
\CommentTok{// reader is automatically closed here, even if an exception occurred}
 
\CommentTok{// Legacy: manual close (show this so students understand what try{-}with{-}resources does)}
BufferedReader reader \OperatorTok{=} \KeywordTok{null}\OperatorTok{;}
\KeywordTok{try} \OperatorTok{\{}
    reader \OperatorTok{=} \KeywordTok{new} BufferedReader\OperatorTok{(}\KeywordTok{new} FileReader\OperatorTok{(}\StringTok{"data.txt"}\OperatorTok{));}
    \CommentTok{// ... use reader}
\OperatorTok{\}} \KeywordTok{catch} \OperatorTok{(}IOException e\OperatorTok{)} \OperatorTok{\{}
    System\OperatorTok{.}\CodeFuncTok{out}\OperatorTok{.}\CodeFuncTok{println}\OperatorTok{(}\StringTok{"Error: "} \OperatorTok{+} e\OperatorTok{.}\CodeFuncTok{getMessage}\OperatorTok{());}
\OperatorTok{\}} \KeywordTok{finally} \OperatorTok{\{}
    \KeywordTok{if} \OperatorTok{(}reader \OperatorTok{!=} \KeywordTok{null}\OperatorTok{)} \OperatorTok{\{}
        \KeywordTok{try} \OperatorTok{\{} reader\OperatorTok{.}\CodeFuncTok{close}\OperatorTok{();} \OperatorTok{\}} \KeywordTok{catch} \OperatorTok{(}IOException e\OperatorTok{)} \OperatorTok{\{} \CommentTok{/* ignore */} \OperatorTok{\}}
    \OperatorTok{\}}
\OperatorTok{\}}
\end{Highlighting}
\end{Shaded}

Show the legacy version to explain why try-with-resources exists and
what it simplifies.

\paragraph{Thought Experiments}\label{thought-experiments-40}

\begin{enumerate}
\def\labelenumi{\arabic{enumi}.}
\item
  \emph{``Checked exceptions force you to handle failures at compile
  time. Critics argue this creates `exception chaining hell' --- every
  method signature becomes littered with \texttt{throws} clauses.
  Defenders argue it makes failure modes explicit and programs more
  reliable. Which side do you find more persuasive? Why?''}
\item
  \emph{``Java's \texttt{Path} and \texttt{Files} classes
  (java.nio.file) are a modern replacement for the older \texttt{File}
  class. The older \texttt{File} class's \texttt{delete()} method
  returns a \texttt{boolean} (true/false) instead of throwing an
  exception. What is the danger of reporting failure via a return value?
  What might a programmer accidentally do wrong?''}
\item
  \emph{``A \texttt{NullPointerException} is called `The Billion Dollar
  Mistake' by its inventor Tony Hoare. Explain why null references cause
  so many bugs. Newer languages (Kotlin, Swift) have built-in null
  safety. What would Java code look like if null didn't exist?''}
\end{enumerate}

\begin{center}\rule{0.5\linewidth}{0.5pt}\end{center}

\subsection{13. Unit 9: Collections, Analysis, and
Testing}\label{unit-9-collections-analysis-and-testing-1}

\subsection{JUnit Testing Deep Dive}\label{junit-testing-deep-dive}

\paragraph{Instructional Cues}\label{instructional-cues-31}

\textbf{Test output is inherently accessible:} JUnit produces plain text
output: \texttt{.} for each passing test, \texttt{F} for failures, with
descriptive failure messages. This is one of the most
screen-reader-friendly outputs in the entire course. Let students
experience testing by running a test suite and listening to the results
before writing their own tests.

\textbf{Test naming convention --- sentence style:}

\begin{Shaded}
\begin{Highlighting}[]
\CodeTypeTok{@Test}
\KeywordTok{public} \CodeTypeTok{void} \CodeFuncTok{calculateArea\_withPositiveRadius\_returnsCorrectValue}\OperatorTok{()} \OperatorTok{\{} \OperatorTok{\}}
 
\CodeTypeTok{@Test}
\KeywordTok{public} \CodeTypeTok{void} \CodeFuncTok{calculateArea\_withZeroRadius\_returnsZero}\OperatorTok{()} \OperatorTok{\{} \OperatorTok{\}}
 
\CodeTypeTok{@Test}
\KeywordTok{public} \CodeTypeTok{void} \CodeFuncTok{calculateArea\_withNegativeRadius\_throwsException}\OperatorTok{()} \OperatorTok{\{} \OperatorTok{\}}
\end{Highlighting}
\end{Shaded}

The naming convention \texttt{methodName\_condition\_expectedBehavior}
makes test names read as complete sentences when the test runner
announces them. This is both good practice and an accessibility benefit
--- each test announcement tells you exactly what was tested.

\textbf{\texttt{assertThrows} --- test that exceptions work correctly:}

\begin{Shaded}
\begin{Highlighting}[]
\CodeTypeTok{@Test}
\KeywordTok{public} \CodeTypeTok{void} \CodeFuncTok{withdraw\_withInsufficientFunds\_throwsException}\OperatorTok{()} \OperatorTok{\{}
    BankAccount account \OperatorTok{=} \KeywordTok{new} \CodeFuncTok{BankAccount}\OperatorTok{(}\CodeNumTok{100.0}\OperatorTok{);}
    \CodeFuncTok{assertThrows}\OperatorTok{(}InsufficientFundsException\OperatorTok{.}\CodeFuncTok{class}\OperatorTok{,} \OperatorTok{()} \OperatorTok{{-}\textgreater{}} account\OperatorTok{.}\CodeFuncTok{withdraw}\OperatorTok{(}\CodeNumTok{200.0}\OperatorTok{));}
\OperatorTok{\}}
\end{Highlighting}
\end{Shaded}

Testing for exceptions is as important as testing for correct outputs.
Many beginners only test the ``happy path.''

\paragraph{Thought Experiments}\label{thought-experiments-41}

\begin{enumerate}
\def\labelenumi{\arabic{enumi}.}
\item
  \emph{``Test-driven development (TDD) says: write the test first, then
  write the implementation. How does writing a test first change the way
  you think about a method? What is forced to be clear before you write
  a line of implementation?''}
\item
  \emph{``Mocking means replacing real dependencies (like a database or
  file system) with fake versions for testing. Why would you want to
  test a method that reads from a database without using a real
  database? What complexity does mocking introduce?''}
\item
  \emph{``A project has 100\% test coverage --- tests execute every
  line. Is the software correct? Name a bug that 100\% coverage would
  not catch. (Hint: think about a bug that only appears with certain
  combinations of inputs.)''}
\end{enumerate}

\begin{center}\rule{0.5\linewidth}{0.5pt}\end{center}

\subsection{14. Units 10--11: Projects and
Capstone}\label{units-1011-projects-and-capstone-1}

\subsection{Java Project
Architecture}\label{java-project-architecture}

For the capstone, encourage a clean package structure:

\begin{verbatim}
src/
|--- com/
|   +--- student/
|       +--- datamanager/
|           |--- model/
|           |   +--- Contact.java
|           |--- storage/
|           |   |--- DataManager.java
|           |   +--- JsonStorage.java
|           |--- ui/
|           |   +--- ConsoleUI.java
|           +--- Main.java
test/
+--- com/
    +--- student/
        +--- datamanager/
            |--- model/
            |   +--- ContactTest.java
            +--- storage/
                +--- DataManagerTest.java
\end{verbatim}

Navigating packages in VS Code: \texttt{Ctrl+Shift+E} for Explorer
(shows package tree), \texttt{Ctrl+P} to open any file by name,
\texttt{Ctrl+Shift+O} for symbols within the current file.

\subsection{IDE-Generated Code
Accessibility}\label{ide-generated-code-accessibility}

VS Code and IntelliJ can generate getters, setters, constructors, and
toString. Teach these as time-savers, but require students to read and
understand all generated code:

\begin{itemize}
\tightlist
\item
  \texttt{Ctrl+Shift+P} -\textgreater{} ``Java: Generate Getters and
  Setters''
\item
  \texttt{Ctrl+Shift+P} -\textgreater{} ``Java: Override/Implement
  Methods''
\item
  \texttt{Ctrl+Shift+P} -\textgreater{} ``Java: Generate toString''
\end{itemize}

Generated code is correct but often not styled for accessibility
(missing JavaDoc, generic names). Require students to add JavaDoc and
verify \texttt{toString()} output is meaningful when heard.

\paragraph{Thought Experiments}\label{thought-experiments-42}

\begin{enumerate}
\def\labelenumi{\arabic{enumi}.}
\item
  \emph{``Java in enterprise software is often criticized for having too
  much `boilerplate' --- code that exists for structural reasons rather
  than to express logic. How does this boilerplate affect the experience
  of maintaining a large Java codebase? What tools and patterns does the
  Java ecosystem provide to reduce it?''}
\item
  \emph{``Your capstone application stores data in a file. If your team
  adds a new field to the data class, existing saved files no longer
  match the new format. This is a `schema migration' problem. How would
  you design your storage layer to handle this gracefully?''}
\end{enumerate}

\begin{center}\rule{0.5\linewidth}{0.5pt}\end{center}

\subsection{15. Bridging from Python to
Java}\label{bridging-from-python-to-java}

For courses that teach Python before Java, use these comparison thought
experiments to bridge the languages:

\textbf{Syntax bridge table:}

\begin{fullwidth}{\def\LTcaptype{none} % do not increment counter
\begin{longtable}[]{@{}
  >{\raggedright\arraybackslash}p{(\linewidth - 4\tabcolsep) * \real{0.1638}}
  >{\raggedright\arraybackslash}p{(\linewidth - 4\tabcolsep) * \real{0.4310}}
  >{\raggedright\arraybackslash}p{(\linewidth - 4\tabcolsep) * \real{0.4052}}@{}}
\toprule\noalign{}
\begin{minipage}[b]{\linewidth}\raggedright
Python
\end{minipage} & \begin{minipage}[b]{\linewidth}\raggedright
Java
\end{minipage} & \begin{minipage}[b]{\linewidth}\raggedright
Notes
\end{minipage} \\
\midrule\noalign{}
\endhead
\bottomrule\noalign{}
\endlastfoot
\texttt{print("Hello")} & \texttt{System.out.println("Hello")} & \\
\texttt{x\ =\ 5} & \texttt{int\ x\ =\ 5;} & Java requires type
declaration and semicolon \\
\texttt{if\ condition:} & \texttt{if\ (condition)\ \{} & Parens around
condition; brace instead of colon \\
\texttt{elif} & \texttt{else\ if} & Two words in Java \\
\texttt{for\ item\ in\ list:} & \texttt{for\ (Type\ item\ :\ list)\ \{}
& Java requires type \\
\texttt{def\ func(x):} & \texttt{public\ static\ void\ func(int\ x)\ \{}
& Return type and access modifier required \\
\texttt{return\ x} & \texttt{return\ x;} & Semicolon \\
\texttt{\#\ comment} & \texttt{//\ comment} & \\
\texttt{None} & \texttt{null} & \\
\texttt{True}/\texttt{False} & \texttt{true}/\texttt{false} & Lowercase
in Java \\
\texttt{len(list)} & \texttt{list.size()} & Method call, not function \\
\texttt{list.append(x)} & \texttt{list.add(x)} & \\
\texttt{dict{[}"key"{]}} & \texttt{map.get("key")} & \\
\texttt{try/except} & \texttt{try/catch} & \\
\texttt{import\ math} & \texttt{import\ java.lang.Math;} (or use
\texttt{Math.} directly) & \\
\end{longtable}
}
\end{fullwidth}

\textbf{Thought experiments for bridging:}

\begin{enumerate}
\def\labelenumi{\arabic{enumi}.}
\item
  \emph{``In Python you write \texttt{x\ =\ 5} and Python figures out
  the type. In Java you write \texttt{int\ x\ =\ 5} and declare the type
  explicitly. Which approach do you prefer after using both? Which one
  would you choose for a large team project? For a quick experiment?''}
\item
  \emph{``Python lists can hold any mixture of types:
  \texttt{{[}1,\ \textquotesingle{}hello\textquotesingle{},\ True,\ 3.14{]}}.
  Java's \texttt{ArrayList\textless{}String\textgreater{}} can only hold
  Strings. What bugs does Python's flexibility enable? What bugs does it
  prevent compared to Java?''}
\item
  \emph{``Python functions are first-class objects --- you can store
  them in variables, pass them as arguments. In Java (before lambdas),
  this was much more complex. Java added lambdas in 2014. What does it
  mean for a language to add a major feature 19 years after it was
  created?''}
\end{enumerate}

\begin{center}\rule{0.5\linewidth}{0.5pt}\end{center}

\subsection{16. Assessment Strategies}\label{assessment-strategies-2}

\subsection{Alternative Assessment
Forms}\label{alternative-assessment-forms}

\textbf{Code review presentations:} The student explains their code out
loud, as if presenting to a team. The teacher asks questions about
design decisions. This tests understanding more authentically than
multiple-choice questions.

\textbf{Bug-hunting assessments:} Provide a working program with 5--8
intentional bugs of different types. The student must find and fix them
using only their screen reader and development tools. This directly
tests real-world debugging skill.

\textbf{Design walkthroughs:} Give the student a problem and ask them to
design the solution --- class structure, methods, data structures ---
without writing any code. Assess the quality of the design, not the
typing speed.

\textbf{Pair programming sessions:} Two students work on one computer.
One dictates code; the other types. Roles swap every 10--15 minutes.
This surfaces communication and reasoning skills that solo assessments
miss.

\begin{center}\rule{0.5\linewidth}{0.5pt}\end{center}

\subsection{17. Supporting Diverse
Needs}\label{supporting-diverse-needs-1}

\subsection{For Braille Display
Users}\label{for-braille-display-users}

Java's verbose syntax can be particularly burdensome on a braille
display with limited width. Strategies:

\begin{itemize}
\tightlist
\item
  Keep lines under 79 characters (PEP 8 applies to style; the Java
  equivalent is generally under 100).
\item
  Use shorter variable names for very long class hierarchies where the
  full name is unwieldy.
\item
  Encourage use of IntelliJ IDEA, which has better braille display
  integration than VS Code in some configurations.
\end{itemize}

\subsection{For Students New to Screen
Readers}\label{for-students-new-to-screen-readers}

If a student is still developing screen reader proficiency:

\begin{itemize}
\tightlist
\item
  Begin with JShell instead of file-based programming. Single-line
  interactions provide immediate feedback.
\item
  Use pair work more extensively until independent navigation is fluent.
\item
  Reduce the number of new concepts per session until screen reader
  fluency stabilizes.
\end{itemize}

\begin{center}\rule{0.5\linewidth}{0.5pt}\end{center}

\subsection{18. Alternative and Supplementary
Resources}\label{alternative-and-supplementary-resources-2}

\subsection{Textbooks and Reading}\label{textbooks-and-reading-2}

\begin{fullwidth}{\def\LTcaptype{none} % do not increment counter
\begin{longtable}[]{@{}
  >{\raggedright\arraybackslash}p{(\linewidth - 6\tabcolsep) * \real{0.3789}}
  >{\raggedright\arraybackslash}p{(\linewidth - 6\tabcolsep) * \real{0.1366}}
  >{\raggedright\arraybackslash}p{(\linewidth - 6\tabcolsep) * \real{0.0807}}
  >{\raggedright\arraybackslash}p{(\linewidth - 6\tabcolsep) * \real{0.4037}}@{}}
\toprule\noalign{}
\begin{minipage}[b]{\linewidth}\raggedright
Resource
\end{minipage} & \begin{minipage}[b]{\linewidth}\raggedright
Level
\end{minipage} & \begin{minipage}[b]{\linewidth}\raggedright
Format
\end{minipage} & \begin{minipage}[b]{\linewidth}\raggedright
Accessibility
\end{minipage} \\
\midrule\noalign{}
\endhead
\bottomrule\noalign{}
\endlastfoot
\textbf{Introduction to Programming Using Java} (Eck, free online) &
Beginner--Intermediate & Free HTML/PDF & HTML version excellent;
comprehensive \\
\textbf{Big Java: Early Objects} (Horstmann) & Beginner & Print/ebook &
Standard university text; clear explanations \\
\textbf{Effective Java, 3rd Ed.} (Bloch) & Intermediate--Advanced &
Print/ebook & Essential for understanding Java's design philosophy \\
\textbf{Head First Java} (Sierra, Bates) & Beginner & Print & Very
visual; but text portions accessible; engaging writing style \\
\textbf{Java: The Complete Reference} (Schildt) & Reference &
Print/ebook & Encyclopedic; good for looking up specific features \\
\textbf{Oracle Java Tutorials} (docs.oracle.com) & Beginner--Advanced &
Free HTML & Excellent HTML accessibility; definitive \\
\textbf{Java Documentation} (javadoc) & Reference & Free HTML & Standard
library docs; fully screen-reader-accessible \\
\end{longtable}
}
\end{fullwidth}

\subsection{Online Courses and
Platforms}\label{online-courses-and-platforms-2}

\begin{fullwidth}{\def\LTcaptype{none} % do not increment counter
\begin{longtable}[]{@{}
  >{\raggedright\arraybackslash}p{(\linewidth - 6\tabcolsep) * \real{0.2791}}
  >{\raggedright\arraybackslash}p{(\linewidth - 6\tabcolsep) * \real{0.2713}}
  >{\raggedright\arraybackslash}p{(\linewidth - 6\tabcolsep) * \real{0.1008}}
  >{\raggedright\arraybackslash}p{(\linewidth - 6\tabcolsep) * \real{0.3488}}@{}}
\toprule\noalign{}
\begin{minipage}[b]{\linewidth}\raggedright
Platform
\end{minipage} & \begin{minipage}[b]{\linewidth}\raggedright
Course
\end{minipage} & \begin{minipage}[b]{\linewidth}\raggedright
Accessibility
\end{minipage} & \begin{minipage}[b]{\linewidth}\raggedright
Notes
\end{minipage} \\
\midrule\noalign{}
\endhead
\bottomrule\noalign{}
\endlastfoot
\textbf{edX / MITx 6.092} & Introduction to Programming in Java & High &
MIT course; video transcripts available \\
\textbf{Coursera} & Java Programming (Duke) & Medium & Free audit
available \\
\textbf{Codecademy} & Learn Java & Partial & Some UI inaccessible;
editor readable \\
\textbf{MOOC.fi} (University of Helsinki) & Java Programming MOOC & High
& Text-based exercises; excellent accessibility \\
\textbf{CodingBat Java} & Java practice problems & High & Plain HTML; no
complex UI \\
\textbf{Exercism Java Track} & Structured exercises with mentoring &
High & Terminal-based; excellent \\
\textbf{LeetCode} & Algorithm problems & Medium & Code editor
accessible; UI moderate \\
\end{longtable}
}
\end{fullwidth}

\textbf{MOOC.fi is especially recommended} for blind students --- the
University of Helsinki's Java course has no-frills HTML exercises and
very clear progression. It is one of the most accessible Java learning
platforms available.

\subsection{Video and Audio
Resources}\label{video-and-audio-resources-1}

\begin{fullwidth}{\def\LTcaptype{none} % do not increment counter
\begin{longtable}[]{@{}
  >{\raggedright\arraybackslash}p{(\linewidth - 4\tabcolsep) * \real{0.4182}}
  >{\raggedright\arraybackslash}p{(\linewidth - 4\tabcolsep) * \real{0.3545}}
  >{\raggedright\arraybackslash}p{(\linewidth - 4\tabcolsep) * \real{0.2273}}@{}}
\toprule\noalign{}
\begin{minipage}[b]{\linewidth}\raggedright
Resource
\end{minipage} & \begin{minipage}[b]{\linewidth}\raggedright
Description
\end{minipage} & \begin{minipage}[b]{\linewidth}\raggedright
Accessibility
\end{minipage} \\
\midrule\noalign{}
\endhead
\bottomrule\noalign{}
\endlastfoot
\textbf{CS50 AP} (Harvard) & AP Computer Science A course using Java &
Captions and transcripts \\
\textbf{Derek Banas Java Tutorials} (YouTube) & Concise
concept-by-concept tutorials & Captions available \\
\textbf{Cave of Programming Java} (YouTube/website) & Beginner Java
series & Captions; text summaries \\
\textbf{thenewboston Java} (YouTube) & Large beginner series & Captions
available \\
\textbf{Coding with John} (YouTube) & Clear explanations of Java
concepts & Good pacing for beginners \\
\end{longtable}
}
\end{fullwidth}

\begin{center}\rule{0.5\linewidth}{0.5pt}\end{center}

\subsection{19. Quick Reference Sheets}\label{quick-reference-sheets-2}

\subsection{Java Compilation and
Running}\label{java-compilation-and-running}

\begin{Shaded}
\begin{Highlighting}[]
\CommentTok{\# Compile a single file}
javac HelloWorld.java    \CommentTok{\# Creates HelloWorld.class}
 
\CommentTok{\# Run compiled program}
java HelloWorld          \CommentTok{\# No .class extension!}
 
\CommentTok{\# Compile with output directory}
javac \CodeTypeTok{{-}d} bin src/HelloWorld.java
 
\CommentTok{\# Run JShell}
jshell                   \CommentTok{\# Interactive REPL}
/exit                    \CommentTok{\# Exit JShell}
 
\CommentTok{\# Run tests (if using Maven)}
mvn test
 
\CommentTok{\# Run tests (if using Gradle)}
./gradlew test
\end{Highlighting}
\end{Shaded}

\subsection{Error Type Quick
Reference}\label{error-type-quick-reference-1}

\begin{fullwidth}{\def\LTcaptype{none} % do not increment counter
\begin{longtable}[]{@{}
  >{\raggedright\arraybackslash}p{(\linewidth - 4\tabcolsep) * \real{0.4276}}
  >{\raggedright\arraybackslash}p{(\linewidth - 4\tabcolsep) * \real{0.3448}}
  >{\raggedright\arraybackslash}p{(\linewidth - 4\tabcolsep) * \real{0.2276}}@{}}
\toprule\noalign{}
\begin{minipage}[b]{\linewidth}\raggedright
Error Type
\end{minipage} & \begin{minipage}[b]{\linewidth}\raggedright
Common Cause
\end{minipage} & \begin{minipage}[b]{\linewidth}\raggedright
First Check
\end{minipage} \\
\midrule\noalign{}
\endhead
\bottomrule\noalign{}
\endlastfoot
\texttt{cannot\ find\ symbol} & Misspelled name or missing import &
Spelling; check imports \\
\texttt{incompatible\ types} & Assigning wrong type or missing cast &
Type of both sides of assignment \\
\texttt{\textquotesingle{};\textquotesingle{}\ expected} & Missing
semicolon & End of the line indicated \\
\texttt{reached\ end\ of\ file\ while\ parsing} & Missing \texttt{\}} &
Count open braces vs close braces \\
\texttt{non-static\ method\ cannot\ be\ referenced\ from\ a\ static\ context}
& Calling instance method from \texttt{main} without object & Need to
create an object first \\
\texttt{NullPointerException} & Calling method on null reference & Which
variable is null \\
\texttt{ArrayIndexOutOfBoundsException} & Index beyond array length &
Print length and the index used \\
\texttt{ClassCastException} & Invalid cast & Check \texttt{instanceof}
before casting \\
\texttt{StackOverflowError} & Infinite recursion & Verify base case \\
\texttt{NumberFormatException} & \texttt{parseInt()} on non-numeric
string & Check input string before parsing \\
\end{longtable}
}
\end{fullwidth}

\subsection{Scanner Method Reference}\label{scanner-method-reference}

\begin{Shaded}
\begin{Highlighting}[]
Scanner scanner \OperatorTok{=} \KeywordTok{new} Scanner\OperatorTok{(}System\OperatorTok{.}\CodeFuncTok{in}\OperatorTok{);}
 
String line   \OperatorTok{=} scanner\OperatorTok{.}\CodeFuncTok{nextLine}\OperatorTok{();}     \CommentTok{// Entire line as String}
String word   \OperatorTok{=} scanner\OperatorTok{.}\CodeFuncTok{next}\OperatorTok{();}         \CommentTok{// Single token (whitespace{-}delimited)}
\CodeTypeTok{int} n         \OperatorTok{=} scanner\OperatorTok{.}\CodeFuncTok{nextInt}\OperatorTok{();}      \CommentTok{// Integer}
\CodeTypeTok{double} d      \OperatorTok{=} scanner\OperatorTok{.}\CodeFuncTok{nextDouble}\OperatorTok{();}   \CommentTok{// Double}
\CodeTypeTok{boolean} b     \OperatorTok{=} scanner\OperatorTok{.}\CodeFuncTok{nextBoolean}\OperatorTok{();}  \CommentTok{// Boolean (true/false)}
 
\CommentTok{// Safe versions with type checking}
\KeywordTok{if} \OperatorTok{(}scanner\OperatorTok{.}\CodeFuncTok{hasNextInt}\OperatorTok{())} \OperatorTok{\{}
    \CodeTypeTok{int} n \OperatorTok{=} scanner\OperatorTok{.}\CodeFuncTok{nextInt}\OperatorTok{();}
\OperatorTok{\}}
 
\CommentTok{// After nextInt(), consume leftover newline:}
scanner\OperatorTok{.}\CodeFuncTok{nextLine}\OperatorTok{();}  \CommentTok{// ALWAYS do this before reading a String}
 
\CommentTok{// From file}
Scanner fileScanner \OperatorTok{=} \KeywordTok{new} Scanner\OperatorTok{(}\KeywordTok{new} File\OperatorTok{(}\StringTok{"data.txt"}\OperatorTok{));}
\CommentTok{// Same methods work; wrap in try{-}catch for FileNotFoundException}
 
scanner\OperatorTok{.}\CodeFuncTok{close}\OperatorTok{();}  \CommentTok{// Always close when done}
\end{Highlighting}
\end{Shaded}

\subsection{Accessible Output
Template}\label{accessible-output-template}

\begin{Shaded}
\begin{Highlighting}[]
\CommentTok{// Numbered list output}
System\OperatorTok{.}\CodeFuncTok{out}\OperatorTok{.}\CodeFuncTok{println}\OperatorTok{(}\StringTok{"{-}{-}{-} Contacts {-}{-}{-}"}\OperatorTok{);}
\KeywordTok{for} \OperatorTok{(}\CodeTypeTok{int} i \OperatorTok{=} \CodeNumTok{0}\OperatorTok{;} i \OperatorTok{\textless{}} contacts\OperatorTok{.}\CodeFuncTok{size}\OperatorTok{();} i\OperatorTok{++)} \OperatorTok{\{}
    System\OperatorTok{.}\CodeFuncTok{out}\OperatorTok{.}\CodeFuncTok{printf}\OperatorTok{(}\StringTok{"}\StringTok{\%d}\StringTok{. }\StringTok{\%s\%n}\StringTok{"}\OperatorTok{,} i \OperatorTok{+} \CodeNumTok{1}\OperatorTok{,} contacts\OperatorTok{.}\CodeFuncTok{get}\OperatorTok{(}i\OperatorTok{));}
\OperatorTok{\}}
System\OperatorTok{.}\CodeFuncTok{out}\OperatorTok{.}\CodeFuncTok{printf}\OperatorTok{(}\StringTok{"Total: }\StringTok{\%d}\StringTok{ contacts}\StringTok{\%n}\StringTok{"}\OperatorTok{,} contacts\OperatorTok{.}\CodeFuncTok{size}\OperatorTok{());}
 
\CommentTok{// Menu}
System\OperatorTok{.}\CodeFuncTok{out}\OperatorTok{.}\CodeFuncTok{println}\OperatorTok{(}\StringTok{"}\StringTok{\textbackslash{}n}\StringTok{{-}{-}{-} Main Menu {-}{-}{-}"}\OperatorTok{);}
System\OperatorTok{.}\CodeFuncTok{out}\OperatorTok{.}\CodeFuncTok{println}\OperatorTok{(}\StringTok{"1. Add contact"}\OperatorTok{);}
System\OperatorTok{.}\CodeFuncTok{out}\OperatorTok{.}\CodeFuncTok{println}\OperatorTok{(}\StringTok{"2. View all contacts"}\OperatorTok{);}
System\OperatorTok{.}\CodeFuncTok{out}\OperatorTok{.}\CodeFuncTok{println}\OperatorTok{(}\StringTok{"3. Search contacts"}\OperatorTok{);}
System\OperatorTok{.}\CodeFuncTok{out}\OperatorTok{.}\CodeFuncTok{println}\OperatorTok{(}\StringTok{"4. Delete contact"}\OperatorTok{);}
System\OperatorTok{.}\CodeFuncTok{out}\OperatorTok{.}\CodeFuncTok{println}\OperatorTok{(}\StringTok{"5. Save and quit"}\OperatorTok{);}
System\OperatorTok{.}\CodeFuncTok{out}\OperatorTok{.}\CodeFuncTok{println}\OperatorTok{(}\StringTok{"{-}{-}{-}{-}{-}{-}{-}{-}{-}{-}{-}{-}{-}{-}{-}{-}{-}"}\OperatorTok{);}
System\OperatorTok{.}\CodeFuncTok{out}\OperatorTok{.}\CodeFuncTok{print}\OperatorTok{(}\StringTok{"Enter option (1{-}5): "}\OperatorTok{);}
 
\CommentTok{// Confirmation messages}
System\OperatorTok{.}\CodeFuncTok{out}\OperatorTok{.}\CodeFuncTok{printf}\OperatorTok{(}\StringTok{"Contact \textquotesingle{}}\StringTok{\%s}\StringTok{\textquotesingle{} added successfully.}\StringTok{\%n}\StringTok{"}\OperatorTok{,} name\OperatorTok{);}
System\OperatorTok{.}\CodeFuncTok{out}\OperatorTok{.}\CodeFuncTok{printf}\OperatorTok{(}\StringTok{"Found }\StringTok{\%d}\StringTok{ matching records.}\StringTok{\%n}\StringTok{"}\OperatorTok{,} count\OperatorTok{);}
System\OperatorTok{.}\CodeFuncTok{out}\OperatorTok{.}\CodeFuncTok{println}\OperatorTok{(}\StringTok{"Data saved to contacts.json."}\OperatorTok{);}
System\OperatorTok{.}\CodeFuncTok{out}\OperatorTok{.}\CodeFuncTok{println}\OperatorTok{(}\StringTok{"Error: Name field cannot be empty. Please try again."}\OperatorTok{);}
\end{Highlighting}
\end{Shaded}

\newpage

\section{Unit 0: Setting Up Your Java Development
Environment}\label{unit-0-setting-up-your-java-development-environment}

\subsection{Overview}\label{overview-14}

This unit walks through installing Java and configuring VS Code for
accessible Java development.

\begin{center}\rule{0.5\linewidth}{0.5pt}\end{center}

\subsection{Lesson 0.1: Installing the Java Development Kit
(JDK)}\label{lesson-0.1-installing-the-java-development-kit-jdk}

\subsection{What Is the JDK?}\label{what-is-the-jdk}

The \textbf{JDK (Java Development Kit)} includes everything you need to
write, compile, and run Java programs:

\begin{itemize}
\tightlist
\item
  \textbf{JDK} --- Tools for developers (compiler, debugger).
\item
  \textbf{JVM (Java Virtual Machine)} --- Runs compiled Java programs.
\item
  \textbf{JRE (Java Runtime Environment)} --- Libraries and runtime
  support.
\end{itemize}

\subsection{Installing the JDK}\label{installing-the-jdk}

\textbf{Recommended:} Download Eclipse Temurin (OpenJDK) from
\url{https://adoptium.net/}.

Choose \textbf{JDK 17 LTS} (Long Term Support) for stability.

\paragraph{Windows}\label{windows-3}

\begin{enumerate}
\def\labelenumi{\arabic{enumi}.}
\tightlist
\item
  Download the \texttt{.msi} installer.
\item
  \textbf{CRUCIAL STEP:} In the installer, ensure ``Set JAVA\_HOME
  variable'' and ``Add to PATH'' are checked.
\item
  Activate ``Install'' and then ``Finish''.
\end{enumerate}

\paragraph{macOS}\label{macos-3}

\begin{enumerate}
\def\labelenumi{\arabic{enumi}.}
\tightlist
\item
  Download the \texttt{.pkg} installer.
\item
  Follow the standard macOS installation steps.
\item
  The JDK is automatically added to your PATH.
\end{enumerate}

\paragraph{Linux (Debian/Ubuntu)}\label{linux-debianubuntu-2}

\begin{Shaded}
\begin{Highlighting}[]
\CodeFuncTok{sudo} apt update
\CodeFuncTok{sudo} apt install openjdk{-}17{-}jdk
\end{Highlighting}
\end{Shaded}

\subsection{Verify Installation}\label{verify-installation-1}

Open a terminal and run:

\begin{Shaded}
\begin{Highlighting}[]
java \CodeTypeTok{{-}{-}version}
javac \CodeTypeTok{{-}{-}version}
\end{Highlighting}
\end{Shaded}

You should hear the version number (e.g., ``openjdk 17.0.x'').

\subsection{Skills Checklist}\label{skills-checklist-67}

\begin{itemize}
\tightlist
\item[$\square$]
  I have installed JDK 17 or later.
\item[$\square$]
  Running \texttt{java\ -\/-version} in the terminal shows the version
  number.
\item[$\square$]
  Running \texttt{javac\ -\/-version} in the terminal shows the compiler
  version.
\end{itemize}

\begin{center}\rule{0.5\linewidth}{0.5pt}\end{center}

\subsection{Lesson 0.2: Installing VS Code and the Java Extension
Pack}\label{lesson-0.2-installing-vs-code-and-the-java-extension-pack}

\begin{enumerate}
\def\labelenumi{\arabic{enumi}.}
\tightlist
\item
  Install VS Code from \url{https://code.visualstudio.com/}.
\item
  Open VS Code and press \texttt{Ctrl\ +\ Shift\ +\ X} to open
  Extensions.
\item
  Search for \textbf{``Extension Pack for Java''} by Microsoft. Install
  it. This installs: Language Support for Java, Debugger for Java, Java
  Test Runner, and Maven support.
\end{enumerate}

\subsection{Configure Screen Reader
Mode}\label{configure-screen-reader-mode}

\begin{itemize}
\tightlist
\item
  Press \texttt{Ctrl\ +\ Shift\ +\ P} -\textgreater{} type ``Toggle
  Screen Reader Optimized'' -\textgreater{} press Enter.
\end{itemize}

\subsection{VS Code Keyboard Shortcuts for
Java}\label{vs-code-keyboard-shortcuts-for-java}

\begin{fullwidth}{\def\LTcaptype{none} % do not increment counter
\begin{longtable}[]{@{}ll@{}}
\toprule\noalign{}
Action & Shortcut \\
\midrule\noalign{}
\endhead
\bottomrule\noalign{}
\endlastfoot
Run Java file & \texttt{Ctrl\ +\ F5} \\
Debug Java file & \texttt{F5} \\
Command Palette & \texttt{Ctrl\ +\ Shift\ +\ P} \\
Problems View & \texttt{Ctrl\ +\ Shift\ +\ M} \\
Next Error & \texttt{F8} \\
Go to Definition & \texttt{F12} \\
Format Code & \texttt{Shift\ +\ Alt\ +\ F} \\
\end{longtable}
}
\end{fullwidth}

\subsection{Skills Checklist}\label{skills-checklist-68}

\begin{itemize}
\tightlist
\item[$\square$]
  VS Code is installed with the Extension Pack for Java.
\item[$\square$]
  I can toggle Screen Reader Optimized Mode.
\item[$\square$]
  I can navigate errors using \texttt{F8}.
\end{itemize}

\begin{center}\rule{0.5\linewidth}{0.5pt}\end{center}

\subsection{Lesson 0.3: Your First Java
Project}\label{lesson-0.3-your-first-java-project}

\subsection{Creating a Project}\label{creating-a-project}

\begin{enumerate}
\def\labelenumi{\arabic{enumi}.}
\tightlist
\item
  Press \texttt{Ctrl\ +\ Shift\ +\ P} -\textgreater{} ``Java: Create
  Java Project''.
\item
  Select ``No build tools''.
\item
  Choose a folder and name the project \texttt{HelloWorld}.
\item
  VS Code will create a \texttt{src} folder with a sample file.
\end{enumerate}

\subsection{Writing Hello World}\label{writing-hello-world}\amn{\textbf{Accessibility Tip:} In VS Code, press \texttt{Ctrl+F5} to compile and run a Java file. Output appears in the Terminal. Use \texttt{Ctrl+`} to focus the terminal, then arrow keys to read output. If you see a compiler error, note the line number in the message and press \texttt{Ctrl+G} in the editor to jump there.}
\subsection{Writing Hello World}\label{writing-hello-world}

Create a file \texttt{HelloWorld.java} in the \texttt{src} folder:

\begin{Shaded}
\begin{Highlighting}[]
\KeywordTok{public} \KeywordTok{class} HelloWorld \OperatorTok{\{}
    \KeywordTok{public} \CodeTypeTok{static} \CodeTypeTok{void} \CodeFuncTok{main}\OperatorTok{(}String\OperatorTok{[]} args\OperatorTok{)} \OperatorTok{\{}
        System\OperatorTok{.}\CodeFuncTok{out}\OperatorTok{.}\CodeFuncTok{println}\OperatorTok{(}\StringTok{"Hello, World!"}\OperatorTok{);}
    \OperatorTok{\}}
\OperatorTok{\}}
\end{Highlighting}
\end{Shaded}

\subsection{Running the Program}\label{running-the-program}

\begin{itemize}
\tightlist
\item
  Press \texttt{Ctrl\ +\ F5} to run without debugging.
\item
  Or open the terminal (`Ctrl + ``) and type:
\end{itemize}

\begin{Shaded}
\begin{Highlighting}[]
javac HelloWorld.java
java HelloWorld
\end{Highlighting}
\end{Shaded}

\subsection{Key Differences from
Python}\label{key-differences-from-python-1}

\begin{fullwidth}{\def\LTcaptype{none} % do not increment counter
\begin{longtable}[]{@{}ll@{}}
\toprule\noalign{}
Python & Java \\
\midrule\noalign{}
\endhead
\bottomrule\noalign{}
\endlastfoot
No required structure & Every program needs a \texttt{class} \\
No semicolons & Statements end with \texttt{;} \\
Indentation defines blocks & Curly braces \texttt{\{\}} define blocks \\
\texttt{print()} & \texttt{System.out.println()} \\
Interpreted & Compiled (\texttt{.java} -\textgreater{} \texttt{.class}
-\textgreater{} run) \\
Dynamic typing & Static typing (declare variable types) \\
\end{longtable}
}
\end{fullwidth}

\subsection{Skills Checklist}\label{skills-checklist-69}

\begin{itemize}
\tightlist
\item[$\square$]
  I can create a new Java project in VS Code.
\item[$\square$]
  I can write and run a Hello World program.
\item[$\square$]
  I understand the basic structure of a Java program.
\end{itemize}

\begin{center}\rule{0.5\linewidth}{0.5pt}\end{center}

\section{Unit 1: Foundations of Computing and Getting Started with
Java}\label{unit-1-foundations-of-computing-and-getting-started-with-java}

\subsection{Overview}\label{overview-15}

This unit introduces computer science concepts and the fundamentals of
Java programming.

\textbf{Objectives:}

\begin{itemize}
\tightlist
\item
  Understand hardware vs.~software.
\item
  Write and run basic Java programs.
\item
  Use variables and data types.
\end{itemize}

\begin{center}\rule{0.5\linewidth}{0.5pt}\end{center}

\subsection{Lesson 1.1: Welcome to Computer
Science}\label{lesson-1.1-welcome-to-computer-science-3}

\emph{(See Python Unit 1.1 for conceptual content --- hardware,
software, algorithms, and problem-solving are language-agnostic.)}

\subsection{Java-Specific Notes}\label{java-specific-notes}

Java was created by James Gosling at Sun Microsystems in 1995. Key
principles:

\begin{itemize}
\tightlist
\item
  \textbf{``Write Once, Run Anywhere''} --- Java programs compile to
  bytecode that runs on any JVM.
\item
  \textbf{Strongly typed} --- Every variable must have a declared type.
\item
  \textbf{Object-Oriented} --- Everything lives inside classes.
\end{itemize}

\begin{center}\rule{0.5\linewidth}{0.5pt}\end{center}

\subsection{Lesson 1.2: Your First Java
Program}\label{lesson-1.2-your-first-java-program-1}

\subsection{Anatomy of a Java
Program}\label{anatomy-of-a-java-program}

\begin{Shaded}
\begin{Highlighting}[]
\CommentTok{// This is a comment}
\KeywordTok{public} \KeywordTok{class} HelloWorld \OperatorTok{\{}                    \CommentTok{// Class declaration}
    \KeywordTok{public} \CodeTypeTok{static} \CodeTypeTok{void} \CodeFuncTok{main}\OperatorTok{(}String\OperatorTok{[]} args\OperatorTok{)} \OperatorTok{\{} \CommentTok{// Main method {-}{-}{-} program entry point}
        System\OperatorTok{.}\CodeFuncTok{out}\OperatorTok{.}\CodeFuncTok{println}\OperatorTok{(}\StringTok{"Hello, World!"}\OperatorTok{);} \CommentTok{// Print a line to the console}
    \OperatorTok{\}}
\OperatorTok{\}}
\end{Highlighting}
\end{Shaded}

\textbf{Key rules:}

\begin{itemize}
\tightlist
\item
  The file name \textbf{must match} the class name:
  \texttt{HelloWorld.java}.
\item
  Every Java program starts execution from
  \texttt{public\ static\ void\ main(String{[}{]}\ args)}.
\item
  Statements end with a semicolon \texttt{;}.
\item
  Blocks of code are surrounded by curly braces \texttt{\{\}}.
\end{itemize}

\subsection{Output Methods}\label{output-methods}

\begin{Shaded}
\begin{Highlighting}[]
System\OperatorTok{.}\CodeFuncTok{out}\OperatorTok{.}\CodeFuncTok{println}\OperatorTok{(}\StringTok{"This prints with a newline at the end."}\OperatorTok{);}
System\OperatorTok{.}\CodeFuncTok{out}\OperatorTok{.}\CodeFuncTok{print}\OperatorTok{(}\StringTok{"This prints without a newline."}\OperatorTok{);}
System\OperatorTok{.}\CodeFuncTok{out}\OperatorTok{.}\CodeFuncTok{printf}\OperatorTok{(}\StringTok{"My name is }\StringTok{\%s}\StringTok{ and I am }\StringTok{\%d}\StringTok{ years old.}\StringTok{\%n}\StringTok{"}\OperatorTok{,} \StringTok{"Alice"}\OperatorTok{,} \CodeNumTok{15}\OperatorTok{);}
\end{Highlighting}
\end{Shaded}

\subsection{Examples}\label{examples-20}

\textbf{Example 1.2.1: Simple Greeting}

\begin{Shaded}
\begin{Highlighting}[]
\KeywordTok{public} \KeywordTok{class} SimpleGreeting \OperatorTok{\{}
    \KeywordTok{public} \CodeTypeTok{static} \CodeTypeTok{void} \CodeFuncTok{main}\OperatorTok{(}String\OperatorTok{[]} args\OperatorTok{)} \OperatorTok{\{}
        System\OperatorTok{.}\CodeFuncTok{out}\OperatorTok{.}\CodeFuncTok{println}\OperatorTok{(}\StringTok{"Hello, World!"}\OperatorTok{);}
        System\OperatorTok{.}\CodeFuncTok{out}\OperatorTok{.}\CodeFuncTok{println}\OperatorTok{(}\StringTok{"Welcome to Java programming!"}\OperatorTok{);}
    \OperatorTok{\}}
\OperatorTok{\}}
\end{Highlighting}
\end{Shaded}

\textbf{Example 1.2.2: Using printf for Formatted Output}

\begin{Shaded}
\begin{Highlighting}[]
\KeywordTok{public} \KeywordTok{class} FormattedOutput \OperatorTok{\{}
    \KeywordTok{public} \CodeTypeTok{static} \CodeTypeTok{void} \CodeFuncTok{main}\OperatorTok{(}String\OperatorTok{[]} args\OperatorTok{)} \OperatorTok{\{}
        System\OperatorTok{.}\CodeFuncTok{out}\OperatorTok{.}\CodeFuncTok{printf}\OperatorTok{(}\StringTok{"Pi is approximately }\StringTok{\%.2f\%n}\StringTok{"}\OperatorTok{,} \CodeNumTok{3.14159}\OperatorTok{);}
        System\OperatorTok{.}\CodeFuncTok{out}\OperatorTok{.}\CodeFuncTok{printf}\OperatorTok{(}\StringTok{"}\StringTok{\%{-}15s}\StringTok{ }\StringTok{\%5d\%n}\StringTok{"}\OperatorTok{,} \StringTok{"Apples:"}\OperatorTok{,} \CodeNumTok{42}\OperatorTok{);}
    \OperatorTok{\}}
\OperatorTok{\}}
\end{Highlighting}
\end{Shaded}

\subsection{Skills Checklist}\label{skills-checklist-70}

\begin{itemize}
\tightlist
\item[$\square$]
  I can write a basic Java program with a \texttt{main} method.
\item[$\square$]
  I can use \texttt{System.out.println()} to display output.
\item[$\square$]
  I understand that Java files must match their class name.
\end{itemize}

\begin{center}\rule{0.5\linewidth}{0.5pt}\end{center}

\subsection{Lesson 1.3: Variables and Primitive Data
Types}\label{lesson-1.3-variables-and-primitive-data-types-1}

\subsection{Primitive Data Types}\label{primitive-data-types}

\begin{fullwidth}{\def\LTcaptype{none} % do not increment counter
\begin{longtable}[]{@{}
  >{\raggedright\arraybackslash}p{(\linewidth - 6\tabcolsep) * \real{0.1304}}
  >{\raggedright\arraybackslash}p{(\linewidth - 6\tabcolsep) * \real{0.3188}}
  >{\raggedright\arraybackslash}p{(\linewidth - 6\tabcolsep) * \real{0.0870}}
  >{\raggedright\arraybackslash}p{(\linewidth - 6\tabcolsep) * \real{0.4638}}@{}}
\toprule\noalign{}
\begin{minipage}[b]{\linewidth}\raggedright
Type
\end{minipage} & \begin{minipage}[b]{\linewidth}\raggedright
Description
\end{minipage} & \begin{minipage}[b]{\linewidth}\raggedright
Size
\end{minipage} & \begin{minipage}[b]{\linewidth}\raggedright
Example
\end{minipage} \\
\midrule\noalign{}
\endhead
\bottomrule\noalign{}
\endlastfoot
\texttt{int} & Whole numbers & 32-bit & \texttt{int\ age\ =\ 25;} \\
\texttt{long} & Large whole numbers & 64-bit &
\texttt{long\ population\ =\ 8000000000L;} \\
\texttt{double} & Decimal numbers & 64-bit &
\texttt{double\ pi\ =\ 3.14159;} \\
\texttt{float} & Decimal (less precise) & 32-bit &
\texttt{float\ price\ =\ 9.99f;} \\
\texttt{boolean} & True or False & 1-bit &
\texttt{boolean\ isRaining\ =\ false;} \\
\texttt{char} & Single character & 16-bit &
\texttt{char\ grade\ =\ \textquotesingle{}A\textquotesingle{};} \\
\end{longtable}
}
\end{fullwidth}

\subsection{String (Reference Type)}\label{string-reference-type}

\begin{Shaded}
\begin{Highlighting}[]
String name \OperatorTok{=} \StringTok{"Alice"}\OperatorTok{;}
String message \OperatorTok{=} \StringTok{"Hello, "} \OperatorTok{+} name \OperatorTok{+} \StringTok{"!"}\OperatorTok{;}  \CommentTok{// Concatenation}
\end{Highlighting}
\end{Shaded}

\subsection{Variable Declaration and
Assignment}\label{variable-declaration-and-assignment}

\begin{Shaded}
\begin{Highlighting}[]
\CommentTok{// Declare and assign in one line}
\CodeTypeTok{int} score \OperatorTok{=} \CodeNumTok{100}\OperatorTok{;}
\CodeTypeTok{double} price \OperatorTok{=} \CodeNumTok{4.99}\OperatorTok{;}
\CodeTypeTok{boolean} isActive \OperatorTok{=} \KeywordTok{true}\OperatorTok{;}
String studentName \OperatorTok{=} \StringTok{"Charlie"}\OperatorTok{;}
 
\CommentTok{// Declare first, assign later}
\CodeTypeTok{int} total\OperatorTok{;}
total \OperatorTok{=} \CodeNumTok{50} \OperatorTok{+} \CodeNumTok{25}\OperatorTok{;}
\end{Highlighting}
\end{Shaded}

\subsection{Constants}\label{constants-1}

\begin{Shaded}
\begin{Highlighting}[]
\CodeTypeTok{final} \CodeTypeTok{double} PI \OperatorTok{=} \CodeNumTok{3.14159265358979}\OperatorTok{;}
\CodeTypeTok{final} \CodeTypeTok{int} MAX\_SIZE \OperatorTok{=} \CodeNumTok{100}\OperatorTok{;}
\end{Highlighting}
\end{Shaded}

\subsection{Type Conversion}\label{type-conversion-1}

\begin{Shaded}
\begin{Highlighting}[]
\CommentTok{// Widening (automatic {-}{-}{-} safe)}
\CodeTypeTok{int} intValue \OperatorTok{=} \CodeNumTok{42}\OperatorTok{;}
\CodeTypeTok{double} doubleValue \OperatorTok{=} intValue\OperatorTok{;}   \CommentTok{// 42 {-}\textgreater{} 42.0}
 
\CommentTok{// Narrowing (explicit cast {-}{-}{-} may lose data)}
\CodeTypeTok{double} d \OperatorTok{=} \CodeNumTok{9.99}\OperatorTok{;}
\CodeTypeTok{int} i \OperatorTok{=} \OperatorTok{(}\CodeTypeTok{int}\OperatorTok{)} d\OperatorTok{;}                 \CommentTok{// 9.99 {-}\textgreater{} 9}
 
\CommentTok{// String to number}
\CodeTypeTok{int} parsed \OperatorTok{=} Integer\OperatorTok{.}\CodeFuncTok{parseInt}\OperatorTok{(}\StringTok{"123"}\OperatorTok{);}
\CodeTypeTok{double} parsedD \OperatorTok{=} Double\OperatorTok{.}\CodeFuncTok{parseDouble}\OperatorTok{(}\StringTok{"3.14"}\OperatorTok{);}
 
\CommentTok{// Number to String}
String s \OperatorTok{=} String\OperatorTok{.}\CodeFuncTok{valueOf}\OperatorTok{(}\CodeNumTok{42}\OperatorTok{);}
String s2 \OperatorTok{=} Integer\OperatorTok{.}\CodeFuncTok{toString}\OperatorTok{(}\CodeNumTok{42}\OperatorTok{);}
\end{Highlighting}
\end{Shaded}

\subsection{Arithmetic Operators}\label{arithmetic-operators-1}

\begin{fullwidth}{\def\LTcaptype{none} % do not increment counter
\begin{longtable}[]{@{}lll@{}}
\toprule\noalign{}
Operator & Description & Example \\
\midrule\noalign{}
\endhead
\bottomrule\noalign{}
\endlastfoot
\texttt{+} & Addition & \texttt{5\ +\ 3} = \texttt{8} \\
\texttt{-} & Subtraction & \texttt{10\ -\ 4} = \texttt{6} \\
\texttt{*} & Multiplication & \texttt{7\ *\ 6} = \texttt{42} \\
\texttt{/} & Division & \texttt{15\ /\ 4} = \texttt{3} (integer
division!) \\
\texttt{\%} & Modulo & \texttt{10\ \%\ 3} = \texttt{1} \\
\texttt{++} & Increment & \texttt{count++} \\
\texttt{-\/-} & Decrement & \texttt{count-\/-} \\
\end{longtable}
}
\end{fullwidth}

\textbf{Important:} Integer division in Java truncates:
\texttt{15\ /\ 4} = \texttt{3}, not \texttt{3.75}. To get a decimal
result, use \texttt{double}: \texttt{15.0\ /\ 4} = \texttt{3.75}.

\subsection{Examples}\label{examples-21}

\textbf{Example 1.3.1: Variables and Arithmetic}

\begin{Shaded}
\begin{Highlighting}[]
\KeywordTok{public} \KeywordTok{class} VariablesDemo \OperatorTok{\{}
    \KeywordTok{public} \CodeTypeTok{static} \CodeTypeTok{void} \CodeFuncTok{main}\OperatorTok{(}String\OperatorTok{[]} args\OperatorTok{)} \OperatorTok{\{}
        \CodeTypeTok{int} appleCount \OperatorTok{=} \CodeNumTok{5}\OperatorTok{;}
        \CodeTypeTok{int} orangeCount \OperatorTok{=} \CodeNumTok{3}\OperatorTok{;}
        \CodeTypeTok{int} totalFruit \OperatorTok{=} appleCount \OperatorTok{+} orangeCount\OperatorTok{;}
        System\OperatorTok{.}\CodeFuncTok{out}\OperatorTok{.}\CodeFuncTok{println}\OperatorTok{(}\StringTok{"You have "} \OperatorTok{+} totalFruit \OperatorTok{+} \StringTok{" fruits in total."}\OperatorTok{);}
 
        \CodeTypeTok{double} pricePerItem \OperatorTok{=} \CodeNumTok{2.50}\OperatorTok{;}
        \CodeTypeTok{int} quantity \OperatorTok{=} \CodeNumTok{4}\OperatorTok{;}
        \CodeTypeTok{double} totalCost \OperatorTok{=} pricePerItem \OperatorTok{*} quantity\OperatorTok{;}
        System\OperatorTok{.}\CodeFuncTok{out}\OperatorTok{.}\CodeFuncTok{println}\OperatorTok{(}\StringTok{"The total cost is: "} \OperatorTok{+} totalCost\OperatorTok{);}
 
        \CommentTok{// Integer division}
        \CodeTypeTok{int} a \OperatorTok{=} \CodeNumTok{10}\OperatorTok{,} b \OperatorTok{=} \CodeNumTok{3}\OperatorTok{;}
        System\OperatorTok{.}\CodeFuncTok{out}\OperatorTok{.}\CodeFuncTok{println}\OperatorTok{(}\StringTok{"10 / 3 = "} \OperatorTok{+} \OperatorTok{(}a \OperatorTok{/} b\OperatorTok{));}       \CommentTok{// 3}
        System\OperatorTok{.}\CodeFuncTok{out}\OperatorTok{.}\CodeFuncTok{println}\OperatorTok{(}\StringTok{"10.0 / 3 = "} \OperatorTok{+} \OperatorTok{(}\CodeNumTok{10.0} \OperatorTok{/} b\OperatorTok{));}  \CommentTok{// 3.3333...}
        System\OperatorTok{.}\CodeFuncTok{out}\OperatorTok{.}\CodeFuncTok{println}\OperatorTok{(}\StringTok{"10 \% 3 = "} \OperatorTok{+} \OperatorTok{(}a \OperatorTok{\%} b\OperatorTok{));}        \CommentTok{// 1}
    \OperatorTok{\}}
\OperatorTok{\}}
\end{Highlighting}
\end{Shaded}

\subsection{Activities}\label{activities-13}

\textbf{Activity 1.3.1: Personal Information}

\begin{Shaded}
\begin{Highlighting}[]
\KeywordTok{public} \KeywordTok{class} PersonalInfo \OperatorTok{\{}
    \KeywordTok{public} \CodeTypeTok{static} \CodeTypeTok{void} \CodeFuncTok{main}\OperatorTok{(}String\OperatorTok{[]} args\OperatorTok{)} \OperatorTok{\{}
        String myName \OperatorTok{=} \StringTok{"Jamie"}\OperatorTok{;}
        \CodeTypeTok{int} favoriteNumber \OperatorTok{=} \CodeNumTok{7}\OperatorTok{;}
        \CodeTypeTok{double} heightInMeters \OperatorTok{=} \CodeNumTok{1.75}\OperatorTok{;}
        \CodeTypeTok{boolean} likesPizza \OperatorTok{=} \KeywordTok{true}\OperatorTok{;}
 
        System\OperatorTok{.}\CodeFuncTok{out}\OperatorTok{.}\CodeFuncTok{println}\OperatorTok{(}\StringTok{"My name is: "} \OperatorTok{+} myName\OperatorTok{);}
        System\OperatorTok{.}\CodeFuncTok{out}\OperatorTok{.}\CodeFuncTok{println}\OperatorTok{(}\StringTok{"Favorite number: "} \OperatorTok{+} favoriteNumber\OperatorTok{);}
        System\OperatorTok{.}\CodeFuncTok{out}\OperatorTok{.}\CodeFuncTok{printf}\OperatorTok{(}\StringTok{"Height: }\StringTok{\%.2f}\StringTok{ meters}\StringTok{\%n}\StringTok{"}\OperatorTok{,} heightInMeters\OperatorTok{);}
        System\OperatorTok{.}\CodeFuncTok{out}\OperatorTok{.}\CodeFuncTok{println}\OperatorTok{(}\StringTok{"Likes pizza: "} \OperatorTok{+} likesPizza\OperatorTok{);}
    \OperatorTok{\}}
\OperatorTok{\}}
\end{Highlighting}
\end{Shaded}

\textbf{Activity 1.3.2: Simple Math}

\begin{Shaded}
\begin{Highlighting}[]
\KeywordTok{public} \KeywordTok{class} SimpleMath \OperatorTok{\{}
    \KeywordTok{public} \CodeTypeTok{static} \CodeTypeTok{void} \CodeFuncTok{main}\OperatorTok{(}String\OperatorTok{[]} args\OperatorTok{)} \OperatorTok{\{}
        \CodeTypeTok{int} num1 \OperatorTok{=} \CodeNumTok{15}\OperatorTok{;}
        \CodeTypeTok{int} num2 \OperatorTok{=} \CodeNumTok{7}\OperatorTok{;}
 
        System\OperatorTok{.}\CodeFuncTok{out}\OperatorTok{.}\CodeFuncTok{println}\OperatorTok{(}\StringTok{"Sum: "} \OperatorTok{+} \OperatorTok{(}num1 \OperatorTok{+} num2\OperatorTok{));}
        System\OperatorTok{.}\CodeFuncTok{out}\OperatorTok{.}\CodeFuncTok{println}\OperatorTok{(}\StringTok{"Difference: "} \OperatorTok{+} \OperatorTok{(}num1 \OperatorTok{{-}} num2\OperatorTok{));}
        System\OperatorTok{.}\CodeFuncTok{out}\OperatorTok{.}\CodeFuncTok{println}\OperatorTok{(}\StringTok{"Product: "} \OperatorTok{+} \OperatorTok{(}num1 \OperatorTok{*} num2\OperatorTok{));}
        System\OperatorTok{.}\CodeFuncTok{out}\OperatorTok{.}\CodeFuncTok{printf}\OperatorTok{(}\StringTok{"Quotient: }\StringTok{\%.2f\%n}\StringTok{"}\OperatorTok{,} \OperatorTok{(}\CodeTypeTok{double}\OperatorTok{)} num1 \OperatorTok{/} num2\OperatorTok{);}
        System\OperatorTok{.}\CodeFuncTok{out}\OperatorTok{.}\CodeFuncTok{println}\OperatorTok{(}\StringTok{"Remainder: "} \OperatorTok{+} \OperatorTok{(}num1 \OperatorTok{\%} num2\OperatorTok{));}
    \OperatorTok{\}}
\OperatorTok{\}}
\end{Highlighting}
\end{Shaded}

\subsection{Skills Checklist}\label{skills-checklist-71}

\begin{itemize}
\tightlist
\item[$\square$]
  I can declare variables with the correct Java types.
\item[$\square$]
  I can perform basic arithmetic operations in Java.
\item[$\square$]
  I understand the difference between integer and floating-point
  division.
\item[$\square$]
  I can convert between types using casting and parse methods.
\item[$\square$]
  I can use \texttt{final} to declare constants.
\end{itemize}

\newpage

\section{Unit 2: Writing Simple Programs and
Input/Output}\label{unit-2-writing-simple-programs-and-inputoutput}

\subsection{Overview}\label{overview-16}

This unit covers getting user input, type conversions, and working with
the Math library in Java.

\begin{center}\rule{0.5\linewidth}{0.5pt}\end{center}

\subsection{Lesson 2.1: The Software Development
Process}\label{lesson-2.1-the-software-development-process-2}\amn{\textbf{Accessibility Tip:} Create your \texttt{Scanner} once at the top of \texttt{main}: \texttt{Scanner scanner = new Scanner(System.in);}. Always print a prompt with \texttt{System.out.print()} (not \texttt{println}) immediately before each \texttt{scanner.nextLine()} call so the screen reader announces what you should type.}
\subsection{Lesson 2.1: The Software Development
Process}\label{lesson-2.1-the-software-development-process-2}

\emph{(See Python Unit 2.1 --- the same design principles apply:
analyze, design, implement, test.)}

\begin{center}\rule{0.5\linewidth}{0.5pt}\end{center}

\subsection{Lesson 2.2: Getting User Input with
Scanner}\label{lesson-2.2-getting-user-input-with-scanner}

\subsection{\texorpdfstring{The \texttt{Scanner}
Class}{The Scanner Class}}\label{the-scanner-class}

Java uses the \texttt{Scanner} class to read input from the user:

\begin{Shaded}
\begin{Highlighting}[]
\KeywordTok{import} \CodeAccentTok{java}\OperatorTok{.}\CodeAccentTok{util}\OperatorTok{.}\CodeAccentTok{Scanner}\OperatorTok{;}
 
\KeywordTok{public} \KeywordTok{class} UserInput \OperatorTok{\{}
    \KeywordTok{public} \CodeTypeTok{static} \CodeTypeTok{void} \CodeFuncTok{main}\OperatorTok{(}String\OperatorTok{[]} args\OperatorTok{)} \OperatorTok{\{}
        Scanner scanner \OperatorTok{=} \KeywordTok{new} Scanner\OperatorTok{(}System\OperatorTok{.}\CodeFuncTok{in}\OperatorTok{);}
 
        System\OperatorTok{.}\CodeFuncTok{out}\OperatorTok{.}\CodeFuncTok{print}\OperatorTok{(}\StringTok{"What is your name? "}\OperatorTok{);}
        String name \OperatorTok{=} scanner\OperatorTok{.}\CodeFuncTok{nextLine}\OperatorTok{();}
 
        System\OperatorTok{.}\CodeFuncTok{out}\OperatorTok{.}\CodeFuncTok{print}\OperatorTok{(}\StringTok{"How old are you? "}\OperatorTok{);}
        \CodeTypeTok{int} age \OperatorTok{=} scanner\OperatorTok{.}\CodeFuncTok{nextInt}\OperatorTok{();}
 
        System\OperatorTok{.}\CodeFuncTok{out}\OperatorTok{.}\CodeFuncTok{println}\OperatorTok{(}\StringTok{"Hello, "} \OperatorTok{+} name \OperatorTok{+} \StringTok{"!"}\OperatorTok{);}
        System\OperatorTok{.}\CodeFuncTok{out}\OperatorTok{.}\CodeFuncTok{println}\OperatorTok{(}\StringTok{"Next year you will be "} \OperatorTok{+} \OperatorTok{(}age \OperatorTok{+} \CodeNumTok{1}\OperatorTok{)} \OperatorTok{+} \StringTok{" years old."}\OperatorTok{);}
 
        scanner\OperatorTok{.}\CodeFuncTok{close}\OperatorTok{();}
    \OperatorTok{\}}
\OperatorTok{\}}
\end{Highlighting}
\end{Shaded}

\subsection{Scanner Methods}\label{scanner-methods}

\begin{fullwidth}{\def\LTcaptype{none} % do not increment counter
\begin{longtable}[]{@{}
  >{\raggedright\arraybackslash}p{(\linewidth - 4\tabcolsep) * \real{0.1852}}
  >{\raggedright\arraybackslash}p{(\linewidth - 4\tabcolsep) * \real{0.3704}}
  >{\raggedright\arraybackslash}p{(\linewidth - 4\tabcolsep) * \real{0.4444}}@{}}
\toprule\noalign{}
\begin{minipage}[b]{\linewidth}\raggedright
Method
\end{minipage} & \begin{minipage}[b]{\linewidth}\raggedright
Reads
\end{minipage} & \begin{minipage}[b]{\linewidth}\raggedright
Example
\end{minipage} \\
\midrule\noalign{}
\endhead
\bottomrule\noalign{}
\endlastfoot
\texttt{nextLine()} & Entire line as String &
\texttt{String\ name\ =\ scanner.nextLine();} \\
\texttt{next()} & Single word (up to whitespace) &
\texttt{String\ word\ =\ scanner.next();} \\
\texttt{nextInt()} & Integer & \texttt{int\ n\ =\ scanner.nextInt();} \\
\texttt{nextDouble()} & Double &
\texttt{double\ d\ =\ scanner.nextDouble();} \\
\texttt{nextBoolean()} & Boolean &
\texttt{boolean\ b\ =\ scanner.nextBoolean();} \\
\end{longtable}
}
\end{fullwidth}

\subsection{\texorpdfstring{Common Pitfall: \texttt{nextInt()} +
\texttt{nextLine()}}{Common Pitfall: nextInt() + nextLine()}}\label{common-pitfall-nextint-nextline}

After calling \texttt{nextInt()}, there is a leftover newline character.
Call \texttt{scanner.nextLine()} to consume it before reading another
string:

\begin{Shaded}
\begin{Highlighting}[]
\CodeTypeTok{int} age \OperatorTok{=} scanner\OperatorTok{.}\CodeFuncTok{nextInt}\OperatorTok{();}
scanner\OperatorTok{.}\CodeFuncTok{nextLine}\OperatorTok{();}            \CommentTok{// Consume leftover newline}
String name \OperatorTok{=} scanner\OperatorTok{.}\CodeFuncTok{nextLine}\OperatorTok{();}  \CommentTok{// Now reads correctly}
\end{Highlighting}
\end{Shaded}

\subsection{Examples}\label{examples-22}

\textbf{Example 2.2.1: Temperature Converter}

\begin{Shaded}
\begin{Highlighting}[]
\KeywordTok{import} \CodeAccentTok{java}\OperatorTok{.}\CodeAccentTok{util}\OperatorTok{.}\CodeAccentTok{Scanner}\OperatorTok{;}
 
\KeywordTok{public} \KeywordTok{class} TempConverter \OperatorTok{\{}
    \KeywordTok{public} \CodeTypeTok{static} \CodeTypeTok{void} \CodeFuncTok{main}\OperatorTok{(}String\OperatorTok{[]} args\OperatorTok{)} \OperatorTok{\{}
        Scanner scanner \OperatorTok{=} \KeywordTok{new} Scanner\OperatorTok{(}System\OperatorTok{.}\CodeFuncTok{in}\OperatorTok{);}
 
        System\OperatorTok{.}\CodeFuncTok{out}\OperatorTok{.}\CodeFuncTok{print}\OperatorTok{(}\StringTok{"Enter temperature in Fahrenheit: "}\OperatorTok{);}
        \CodeTypeTok{double} fahrenheit \OperatorTok{=} scanner\OperatorTok{.}\CodeFuncTok{nextDouble}\OperatorTok{();}
 
        \CodeTypeTok{double} celsius \OperatorTok{=} \OperatorTok{(}fahrenheit \OperatorTok{{-}} \CodeNumTok{32}\OperatorTok{)} \OperatorTok{*} \CodeNumTok{5.0} \OperatorTok{/} \CodeNumTok{9.0}\OperatorTok{;}
 
        System\OperatorTok{.}\CodeFuncTok{out}\OperatorTok{.}\CodeFuncTok{printf}\OperatorTok{(}\StringTok{"}\StringTok{\%.1f}\StringTok{°F is }\StringTok{\%.1f}\StringTok{°C}\StringTok{\%n}\StringTok{"}\OperatorTok{,} fahrenheit\OperatorTok{,} celsius\OperatorTok{);}
        scanner\OperatorTok{.}\CodeFuncTok{close}\OperatorTok{();}
    \OperatorTok{\}}
\OperatorTok{\}}
\end{Highlighting}
\end{Shaded}

\textbf{Example 2.2.2: Circle Calculator}

\begin{Shaded}
\begin{Highlighting}[]
\KeywordTok{import} \CodeAccentTok{java}\OperatorTok{.}\CodeAccentTok{util}\OperatorTok{.}\CodeAccentTok{Scanner}\OperatorTok{;}
 
\KeywordTok{public} \KeywordTok{class} CircleCalculator \OperatorTok{\{}
    \KeywordTok{public} \CodeTypeTok{static} \CodeTypeTok{void} \CodeFuncTok{main}\OperatorTok{(}String\OperatorTok{[]} args\OperatorTok{)} \OperatorTok{\{}
        Scanner scanner \OperatorTok{=} \KeywordTok{new} Scanner\OperatorTok{(}System\OperatorTok{.}\CodeFuncTok{in}\OperatorTok{);}
 
        System\OperatorTok{.}\CodeFuncTok{out}\OperatorTok{.}\CodeFuncTok{print}\OperatorTok{(}\StringTok{"Enter the radius: "}\OperatorTok{);}
        \CodeTypeTok{double} radius \OperatorTok{=} scanner\OperatorTok{.}\CodeFuncTok{nextDouble}\OperatorTok{();}
 
        \CodeTypeTok{double} area \OperatorTok{=} Math\OperatorTok{.}\CodeFuncTok{PI} \OperatorTok{*} radius \OperatorTok{*} radius\OperatorTok{;}
        \CodeTypeTok{double} circumference \OperatorTok{=} \CodeNumTok{2} \OperatorTok{*} Math\OperatorTok{.}\CodeFuncTok{PI} \OperatorTok{*} radius\OperatorTok{;}
 
        System\OperatorTok{.}\CodeFuncTok{out}\OperatorTok{.}\CodeFuncTok{printf}\OperatorTok{(}\StringTok{"Area: }\StringTok{\%.2f\%n}\StringTok{"}\OperatorTok{,} area\OperatorTok{);}
        System\OperatorTok{.}\CodeFuncTok{out}\OperatorTok{.}\CodeFuncTok{printf}\OperatorTok{(}\StringTok{"Circumference: }\StringTok{\%.2f\%n}\StringTok{"}\OperatorTok{,} circumference\OperatorTok{);}
        scanner\OperatorTok{.}\CodeFuncTok{close}\OperatorTok{();}
    \OperatorTok{\}}
\OperatorTok{\}}
\end{Highlighting}
\end{Shaded}

\subsection{Skills Checklist}\label{skills-checklist-72}

\begin{itemize}
\tightlist
\item[$\square$]
  I can import and create a \texttt{Scanner} object.
\item[$\square$]
  I can read different types of input from the user.
\item[$\square$]
  I can handle the \texttt{nextInt()}/\texttt{nextLine()} newline issue.
\end{itemize}

\begin{center}\rule{0.5\linewidth}{0.5pt}\end{center}

\subsection{Lesson 2.3: The Math Class}\label{lesson-2.3-the-math-class}

\begin{Shaded}
\begin{Highlighting}[]
\CommentTok{// Math class provides static methods {-}{-}{-} no import needed}
\CodeTypeTok{double} squareRoot \OperatorTok{=} Math\OperatorTok{.}\CodeFuncTok{sqrt}\OperatorTok{(}\CodeNumTok{25.0}\OperatorTok{);}     \CommentTok{// 5.0}
\CodeTypeTok{double} power \OperatorTok{=} Math\OperatorTok{.}\CodeFuncTok{pow}\OperatorTok{(}\CodeNumTok{2}\OperatorTok{,} \CodeNumTok{10}\OperatorTok{);}           \CommentTok{// 1024.0}
\CodeTypeTok{double} absolute \OperatorTok{=} Math\OperatorTok{.}\CodeFuncTok{abs}\OperatorTok{({-}}\CodeNumTok{7.5}\OperatorTok{);}         \CommentTok{// 7.5}
\CodeTypeTok{int} ceiling \OperatorTok{=} \OperatorTok{(}\CodeTypeTok{int}\OperatorTok{)} Math\OperatorTok{.}\CodeFuncTok{ceil}\OperatorTok{(}\CodeNumTok{4.2}\OperatorTok{);}       \CommentTok{// 5}
\CodeTypeTok{int} floor \OperatorTok{=} \OperatorTok{(}\CodeTypeTok{int}\OperatorTok{)} Math\OperatorTok{.}\CodeFuncTok{floor}\OperatorTok{(}\CodeNumTok{4.8}\OperatorTok{);}        \CommentTok{// 4}
\CodeTypeTok{double} rounded \OperatorTok{=} Math\OperatorTok{.}\CodeFuncTok{round}\OperatorTok{(}\CodeNumTok{4.6}\OperatorTok{);}         \CommentTok{// 5}
\CodeTypeTok{double} maxVal \OperatorTok{=} Math\OperatorTok{.}\CodeFuncTok{max}\OperatorTok{(}\CodeNumTok{10}\OperatorTok{,} \CodeNumTok{20}\OperatorTok{);}         \CommentTok{// 20}
\CodeTypeTok{double} minVal \OperatorTok{=} Math\OperatorTok{.}\CodeFuncTok{min}\OperatorTok{(}\CodeNumTok{10}\OperatorTok{,} \CodeNumTok{20}\OperatorTok{);}         \CommentTok{// 10}
\CodeTypeTok{double} piValue \OperatorTok{=} Math\OperatorTok{.}\CodeFuncTok{PI}\OperatorTok{;}                 \CommentTok{// 3.14159...}
\CodeTypeTok{double} eValue \OperatorTok{=} Math\OperatorTok{.}\CodeFuncTok{E}\OperatorTok{;}                   \CommentTok{// 2.71828...}
\CodeTypeTok{double} sinVal \OperatorTok{=} Math\OperatorTok{.}\CodeFuncTok{sin}\OperatorTok{(}Math\OperatorTok{.}\CodeFuncTok{PI} \OperatorTok{/} \CodeNumTok{2}\OperatorTok{);}   \CommentTok{// 1.0}
\CodeTypeTok{double} logVal \OperatorTok{=} Math\OperatorTok{.}\CodeFuncTok{log}\OperatorTok{(}Math\OperatorTok{.}\CodeFuncTok{E}\OperatorTok{);}        \CommentTok{// 1.0  (natural log)}
\CodeTypeTok{double} log10 \OperatorTok{=} Math\OperatorTok{.}\CodeFuncTok{log10}\OperatorTok{(}\CodeNumTok{100}\OperatorTok{);}          \CommentTok{// 2.0}
\end{Highlighting}
\end{Shaded}

\subsection{Skills Checklist}\label{skills-checklist-73}

\begin{itemize}
\tightlist
\item[$\square$]
  I can use common \texttt{Math} class methods.
\item[$\square$]
  I can use \texttt{Math.PI} for circle calculations.
\end{itemize}

\begin{center}\rule{0.5\linewidth}{0.5pt}\end{center}

\section{Unit 3: Control Flow --- Decisions and
Loops}\label{unit-3-control-flow-decisions-and-loops-3}

\subsection{Overview}\label{overview-17}

This unit introduces conditional statements and loops in Java.

\textbf{Objectives:}

\begin{itemize}
\tightlist
\item
  Use \texttt{if}, \texttt{else\ if}, and \texttt{else} for
  decision-making.
\item
  Use \texttt{for}, \texttt{while}, and \texttt{do-while} loops.
\item
  Use \texttt{break} and \texttt{continue}.
\end{itemize}

\begin{center}\rule{0.5\linewidth}{0.5pt}\end{center}

\subsection{\texorpdfstring{Lesson 3.1: Boolean Expressions and
\texttt{if}
Statements}{Lesson 3.1: Boolean Expressions and if Statements}}\label{lesson-3.1-boolean-expressions-and-if-statements-2}

\subsection{Comparison Operators}\label{comparison-operators-1}

\begin{fullwidth}{\def\LTcaptype{none} % do not increment counter
\begin{longtable}[]{@{}lll@{}}
\toprule\noalign{}
Operator & Meaning & Example \\
\midrule\noalign{}
\endhead
\bottomrule\noalign{}
\endlastfoot
\texttt{==} & Equal to & \texttt{age\ ==\ 18} \\
\texttt{!=} & Not equal to & \texttt{score\ !=\ 0} \\
\texttt{\textless{}} & Less than & \texttt{temp\ \textless{}\ 0} \\
\texttt{\textgreater{}} & Greater than &
\texttt{score\ \textgreater{}\ 90} \\
\texttt{\textless{}=} & Less than or equal &
\texttt{age\ \textless{}=\ 17} \\
\texttt{\textgreater{}=} & Greater than or equal &
\texttt{score\ \textgreater{}=\ 60} \\
\end{longtable}
}
\end{fullwidth}

\subsection{Logical Operators}\label{logical-operators-1}

\begin{fullwidth}{\def\LTcaptype{none} % do not increment counter
\begin{longtable}[]{@{}
  >{\raggedright\arraybackslash}p{(\linewidth - 4\tabcolsep) * \real{0.1111}}
  >{\raggedright\arraybackslash}p{(\linewidth - 4\tabcolsep) * \real{0.3056}}
  >{\raggedright\arraybackslash}p{(\linewidth - 4\tabcolsep) * \real{0.5833}}@{}}
\toprule\noalign{}
\begin{minipage}[b]{\linewidth}\raggedright
Operator
\end{minipage} & \begin{minipage}[b]{\linewidth}\raggedright
Meaning
\end{minipage} & \begin{minipage}[b]{\linewidth}\raggedright
Example
\end{minipage} \\
\midrule\noalign{}
\endhead
\bottomrule\noalign{}
\endlastfoot
\texttt{\&\&} & AND (both true) &
\texttt{age\ \textgreater{}=\ 18\ \&\&\ hasId\ ==\ true} \\
\texttt{\textbackslash{}\textbar{}\textbackslash{}\textbar{}} & OR (at
least one true) &
\texttt{day.equals("Sat")\ \textbackslash{}\textbar{}\textbackslash{}\textbar{}\ day.equals("Sun")} \\
\texttt{!} & NOT (negates) & \texttt{!isRaining} \\
\end{longtable}
}
\end{fullwidth}

\subsection{\texorpdfstring{\texttt{if-else\ if-else}}{if-else if-else}}\label{if-else-if-else-1}

\begin{Shaded}
\begin{Highlighting}[]
\CodeTypeTok{int} score \OperatorTok{=} \CodeNumTok{85}\OperatorTok{;}
 
\KeywordTok{if} \OperatorTok{(}score \OperatorTok{\textgreater{}=} \CodeNumTok{90}\OperatorTok{)} \OperatorTok{\{}
    System\OperatorTok{.}\CodeFuncTok{out}\OperatorTok{.}\CodeFuncTok{println}\OperatorTok{(}\StringTok{"Grade: A"}\OperatorTok{);}
\OperatorTok{\}} \KeywordTok{else} \KeywordTok{if} \OperatorTok{(}score \OperatorTok{\textgreater{}=} \CodeNumTok{80}\OperatorTok{)} \OperatorTok{\{}
    System\OperatorTok{.}\CodeFuncTok{out}\OperatorTok{.}\CodeFuncTok{println}\OperatorTok{(}\StringTok{"Grade: B"}\OperatorTok{);}
\OperatorTok{\}} \KeywordTok{else} \KeywordTok{if} \OperatorTok{(}score \OperatorTok{\textgreater{}=} \CodeNumTok{70}\OperatorTok{)} \OperatorTok{\{}
    System\OperatorTok{.}\CodeFuncTok{out}\OperatorTok{.}\CodeFuncTok{println}\OperatorTok{(}\StringTok{"Grade: C"}\OperatorTok{);}
\OperatorTok{\}} \KeywordTok{else} \OperatorTok{\{}
    System\OperatorTok{.}\CodeFuncTok{out}\OperatorTok{.}\CodeFuncTok{println}\OperatorTok{(}\StringTok{"Grade: F"}\OperatorTok{);}
\OperatorTok{\}}
\end{Highlighting}
\end{Shaded}

\subsection{String Comparison}\label{string-comparison}

\textbf{Important:} Never use \texttt{==} to compare Strings. Use
\texttt{.equals()}:

\begin{Shaded}
\begin{Highlighting}[]
String color \OperatorTok{=} \StringTok{"blue"}\OperatorTok{;}
\KeywordTok{if} \OperatorTok{(}color\OperatorTok{.}\CodeFuncTok{equals}\OperatorTok{(}\StringTok{"blue"}\OperatorTok{))} \OperatorTok{\{}
    System\OperatorTok{.}\CodeFuncTok{out}\OperatorTok{.}\CodeFuncTok{println}\OperatorTok{(}\StringTok{"It\textquotesingle{}s blue!"}\OperatorTok{);}
\OperatorTok{\}}
\KeywordTok{if} \OperatorTok{(}color\OperatorTok{.}\CodeFuncTok{equalsIgnoreCase}\OperatorTok{(}\StringTok{"BLUE"}\OperatorTok{))} \OperatorTok{\{}
    System\OperatorTok{.}\CodeFuncTok{out}\OperatorTok{.}\CodeFuncTok{println}\OperatorTok{(}\StringTok{"Case{-}insensitive match!"}\OperatorTok{);}
\OperatorTok{\}}
\end{Highlighting}
\end{Shaded}

\subsection{\texorpdfstring{\texttt{switch}
Statement}{switch Statement}}\label{switch-statement-1}

\begin{Shaded}
\begin{Highlighting}[]
\CodeTypeTok{int} day \OperatorTok{=} \CodeNumTok{3}\OperatorTok{;}
\KeywordTok{switch} \OperatorTok{(}day\OperatorTok{)} \OperatorTok{\{}
    \KeywordTok{case} \CodeNumTok{1}\OperatorTok{:}
        System\OperatorTok{.}\CodeFuncTok{out}\OperatorTok{.}\CodeFuncTok{println}\OperatorTok{(}\StringTok{"Monday"}\OperatorTok{);}
        \KeywordTok{break}\OperatorTok{;}
    \KeywordTok{case} \CodeNumTok{2}\OperatorTok{:}
        System\OperatorTok{.}\CodeFuncTok{out}\OperatorTok{.}\CodeFuncTok{println}\OperatorTok{(}\StringTok{"Tuesday"}\OperatorTok{);}
        \KeywordTok{break}\OperatorTok{;}
    \KeywordTok{case} \CodeNumTok{3}\OperatorTok{:}
        System\OperatorTok{.}\CodeFuncTok{out}\OperatorTok{.}\CodeFuncTok{println}\OperatorTok{(}\StringTok{"Wednesday"}\OperatorTok{);}
        \KeywordTok{break}\OperatorTok{;}
    \KeywordTok{default}\OperatorTok{:}
        System\OperatorTok{.}\CodeFuncTok{out}\OperatorTok{.}\CodeFuncTok{println}\OperatorTok{(}\StringTok{"Other day"}\OperatorTok{);}
\OperatorTok{\}}
\end{Highlighting}
\end{Shaded}

\subsection{Examples}\label{examples-23}

\textbf{Example 3.1.1: Age Check}

\begin{Shaded}
\begin{Highlighting}[]
\KeywordTok{import} \CodeAccentTok{java}\OperatorTok{.}\CodeAccentTok{util}\OperatorTok{.}\CodeAccentTok{Scanner}\OperatorTok{;}
 
\KeywordTok{public} \KeywordTok{class} AgeCheck \OperatorTok{\{}
    \KeywordTok{public} \CodeTypeTok{static} \CodeTypeTok{void} \CodeFuncTok{main}\OperatorTok{(}String\OperatorTok{[]} args\OperatorTok{)} \OperatorTok{\{}
        Scanner scanner \OperatorTok{=} \KeywordTok{new} Scanner\OperatorTok{(}System\OperatorTok{.}\CodeFuncTok{in}\OperatorTok{);}
        System\OperatorTok{.}\CodeFuncTok{out}\OperatorTok{.}\CodeFuncTok{print}\OperatorTok{(}\StringTok{"Enter your age: "}\OperatorTok{);}
        \CodeTypeTok{int} age \OperatorTok{=} scanner\OperatorTok{.}\CodeFuncTok{nextInt}\OperatorTok{();}
 
        \KeywordTok{if} \OperatorTok{(}age \OperatorTok{\textgreater{}=} \CodeNumTok{18}\OperatorTok{)} \OperatorTok{\{}
            System\OperatorTok{.}\CodeFuncTok{out}\OperatorTok{.}\CodeFuncTok{println}\OperatorTok{(}\StringTok{"You are an adult."}\OperatorTok{);}
        \OperatorTok{\}} \KeywordTok{else} \OperatorTok{\{}
            System\OperatorTok{.}\CodeFuncTok{out}\OperatorTok{.}\CodeFuncTok{println}\OperatorTok{(}\StringTok{"You are a minor."}\OperatorTok{);}
        \OperatorTok{\}}
        scanner\OperatorTok{.}\CodeFuncTok{close}\OperatorTok{();}
    \OperatorTok{\}}
\OperatorTok{\}}
\end{Highlighting}
\end{Shaded}

\textbf{Example 3.1.2: Grade Calculator}

\begin{Shaded}
\begin{Highlighting}[]
\KeywordTok{import} \CodeAccentTok{java}\OperatorTok{.}\CodeAccentTok{util}\OperatorTok{.}\CodeAccentTok{Scanner}\OperatorTok{;}
 
\KeywordTok{public} \KeywordTok{class} GradeCalculator \OperatorTok{\{}
    \KeywordTok{public} \CodeTypeTok{static} \CodeTypeTok{void} \CodeFuncTok{main}\OperatorTok{(}String\OperatorTok{[]} args\OperatorTok{)} \OperatorTok{\{}
        Scanner scanner \OperatorTok{=} \KeywordTok{new} Scanner\OperatorTok{(}System\OperatorTok{.}\CodeFuncTok{in}\OperatorTok{);}
        System\OperatorTok{.}\CodeFuncTok{out}\OperatorTok{.}\CodeFuncTok{print}\OperatorTok{(}\StringTok{"Enter your score (0{-}100): "}\OperatorTok{);}
        \CodeTypeTok{int} score \OperatorTok{=} scanner\OperatorTok{.}\CodeFuncTok{nextInt}\OperatorTok{();}
 
        String grade\OperatorTok{;}
        \KeywordTok{if} \OperatorTok{(}score \OperatorTok{\textgreater{}=} \CodeNumTok{90}\OperatorTok{)} \OperatorTok{\{}
            grade \OperatorTok{=} \StringTok{"A"}\OperatorTok{;}
        \OperatorTok{\}} \KeywordTok{else} \KeywordTok{if} \OperatorTok{(}score \OperatorTok{\textgreater{}=} \CodeNumTok{80}\OperatorTok{)} \OperatorTok{\{}
            grade \OperatorTok{=} \StringTok{"B"}\OperatorTok{;}
        \OperatorTok{\}} \KeywordTok{else} \KeywordTok{if} \OperatorTok{(}score \OperatorTok{\textgreater{}=} \CodeNumTok{70}\OperatorTok{)} \OperatorTok{\{}
            grade \OperatorTok{=} \StringTok{"C"}\OperatorTok{;}
        \OperatorTok{\}} \KeywordTok{else} \KeywordTok{if} \OperatorTok{(}score \OperatorTok{\textgreater{}=} \CodeNumTok{60}\OperatorTok{)} \OperatorTok{\{}
            grade \OperatorTok{=} \StringTok{"D"}\OperatorTok{;}
        \OperatorTok{\}} \KeywordTok{else} \OperatorTok{\{}
            grade \OperatorTok{=} \StringTok{"F"}\OperatorTok{;}
        \OperatorTok{\}}
 
        System\OperatorTok{.}\CodeFuncTok{out}\OperatorTok{.}\CodeFuncTok{println}\OperatorTok{(}\StringTok{"Your grade: "} \OperatorTok{+} grade\OperatorTok{);}
        scanner\OperatorTok{.}\CodeFuncTok{close}\OperatorTok{();}
    \OperatorTok{\}}
\OperatorTok{\}}
\end{Highlighting}
\end{Shaded}

\subsection{Skills Checklist}\label{skills-checklist-74}

\begin{itemize}
\tightlist
\item[$\square$]
  I can write \texttt{if}, \texttt{else\ if}, and \texttt{else}
  statements in Java.
\item[$\square$]
  I can use \texttt{\&\&}, \texttt{\textbar{}\textbar{}}, and \texttt{!}
  to combine conditions.
\item[$\square$]
  I use \texttt{.equals()} (not \texttt{==}) to compare Strings.
\item[$\square$]
  I can write a \texttt{switch} statement.
\end{itemize}

\begin{center}\rule{0.5\linewidth}{0.5pt}\end{center}

\subsection{\texorpdfstring{Lesson 3.2: \texttt{for}
Loops}{Lesson 3.2: for Loops}}\label{lesson-3.2-for-loops-1}

\begin{Shaded}
\begin{Highlighting}[]
\CommentTok{// Basic for loop}
\KeywordTok{for} \OperatorTok{(}\CodeTypeTok{int} i \OperatorTok{=} \CodeNumTok{0}\OperatorTok{;} i \OperatorTok{\textless{}} \CodeNumTok{5}\OperatorTok{;} i\OperatorTok{++)} \OperatorTok{\{}
    System\OperatorTok{.}\CodeFuncTok{out}\OperatorTok{.}\CodeFuncTok{println}\OperatorTok{(}\StringTok{"i = "} \OperatorTok{+} i\OperatorTok{);}
\OperatorTok{\}}
 
\CommentTok{// Counting down}
\KeywordTok{for} \OperatorTok{(}\CodeTypeTok{int} i \OperatorTok{=} \CodeNumTok{10}\OperatorTok{;} i \OperatorTok{\textgreater{}} \CodeNumTok{0}\OperatorTok{;} i\OperatorTok{{-}{-})} \OperatorTok{\{}
    System\OperatorTok{.}\CodeFuncTok{out}\OperatorTok{.}\CodeFuncTok{println}\OperatorTok{(}i\OperatorTok{);}
\OperatorTok{\}}
System\OperatorTok{.}\CodeFuncTok{out}\OperatorTok{.}\CodeFuncTok{println}\OperatorTok{(}\StringTok{"Blast off!"}\OperatorTok{);}
 
\CommentTok{// Stepping by 2}
\KeywordTok{for} \OperatorTok{(}\CodeTypeTok{int} i \OperatorTok{=} \CodeNumTok{0}\OperatorTok{;} i \OperatorTok{\textless{}=} \CodeNumTok{10}\OperatorTok{;} i \OperatorTok{+=} \CodeNumTok{2}\OperatorTok{)} \OperatorTok{\{}
    System\OperatorTok{.}\CodeFuncTok{out}\OperatorTok{.}\CodeFuncTok{println}\OperatorTok{(}i\OperatorTok{);}
\OperatorTok{\}}
\end{Highlighting}
\end{Shaded}

\subsection{\texorpdfstring{Enhanced \texttt{for} Loop
(for-each)}{Enhanced for Loop (for-each)}}\label{enhanced-for-loop-for-each}

\begin{Shaded}
\begin{Highlighting}[]
String\OperatorTok{[]} fruits \OperatorTok{=} \OperatorTok{\{}\StringTok{"apple"}\OperatorTok{,} \StringTok{"banana"}\OperatorTok{,} \StringTok{"cherry"}\OperatorTok{\};}
\KeywordTok{for} \OperatorTok{(}String fruit \OperatorTok{:} fruits\OperatorTok{)} \OperatorTok{\{}
    System\OperatorTok{.}\CodeFuncTok{out}\OperatorTok{.}\CodeFuncTok{println}\OperatorTok{(}\StringTok{"I like "} \OperatorTok{+} fruit \OperatorTok{+} \StringTok{"."}\OperatorTok{);}
\OperatorTok{\}}
 
\CodeTypeTok{int}\OperatorTok{[]} numbers \OperatorTok{=} \OperatorTok{\{}\CodeNumTok{1}\OperatorTok{,} \CodeNumTok{2}\OperatorTok{,} \CodeNumTok{3}\OperatorTok{,} \CodeNumTok{4}\OperatorTok{,} \CodeNumTok{5}\OperatorTok{\};}
\CodeTypeTok{int} sum \OperatorTok{=} \CodeNumTok{0}\OperatorTok{;}
\KeywordTok{for} \OperatorTok{(}\CodeTypeTok{int} num \OperatorTok{:} numbers\OperatorTok{)} \OperatorTok{\{}
    sum \OperatorTok{+=} num\OperatorTok{;}
\OperatorTok{\}}
System\OperatorTok{.}\CodeFuncTok{out}\OperatorTok{.}\CodeFuncTok{println}\OperatorTok{(}\StringTok{"Sum: "} \OperatorTok{+} sum\OperatorTok{);}
\end{Highlighting}
\end{Shaded}

\subsection{Examples}\label{examples-24}

\textbf{Example 3.2.1: Multiplication Table}

\begin{Shaded}
\begin{Highlighting}[]
\KeywordTok{public} \KeywordTok{class} MultiplicationTable \OperatorTok{\{}
    \KeywordTok{public} \CodeTypeTok{static} \CodeTypeTok{void} \CodeFuncTok{main}\OperatorTok{(}String\OperatorTok{[]} args\OperatorTok{)} \OperatorTok{\{}
        \CodeTypeTok{int} n \OperatorTok{=} \CodeNumTok{5}\OperatorTok{;}
        System\OperatorTok{.}\CodeFuncTok{out}\OperatorTok{.}\CodeFuncTok{println}\OperatorTok{(}\StringTok{"Multiplication table for "} \OperatorTok{+} n \OperatorTok{+} \StringTok{":"}\OperatorTok{);}
        \KeywordTok{for} \OperatorTok{(}\CodeTypeTok{int} i \OperatorTok{=} \CodeNumTok{1}\OperatorTok{;} i \OperatorTok{\textless{}=} \CodeNumTok{10}\OperatorTok{;} i\OperatorTok{++)} \OperatorTok{\{}
            System\OperatorTok{.}\CodeFuncTok{out}\OperatorTok{.}\CodeFuncTok{println}\OperatorTok{(}n \OperatorTok{+} \StringTok{" x "} \OperatorTok{+} i \OperatorTok{+} \StringTok{" = "} \OperatorTok{+} \OperatorTok{(}n \OperatorTok{*} i\OperatorTok{));}
        \OperatorTok{\}}
    \OperatorTok{\}}
\OperatorTok{\}}
\end{Highlighting}
\end{Shaded}

\subsection{Skills Checklist}\label{skills-checklist-75}

\begin{itemize}
\tightlist
\item[$\square$]
  I can write a basic \texttt{for} loop with initialization, condition,
  and update.
\item[$\square$]
  I can use an enhanced \texttt{for} loop to iterate over an array.
\end{itemize}

\begin{center}\rule{0.5\linewidth}{0.5pt}\end{center}

\subsection{\texorpdfstring{Lesson 3.3: \texttt{while} and
\texttt{do-while}
Loops}{Lesson 3.3: while and do-while Loops}}\label{lesson-3.3-while-and-do-while-loops}

\subsection{\texorpdfstring{\texttt{while}
Loop}{while Loop}}\label{while-loop}\amn{\textbf{Accessibility Tip:} In Java, every block is defined by braces \texttt{\{\}}---not indentation. Indent consistently (4 spaces) so that when your screen reader announces indent level changes, they correspond to block entry and exit. Add a comment at closing braces of long blocks: \texttt{\} // end while} to help you navigate.}
\subsection{\texorpdfstring{\texttt{while}
Loop}{while Loop}}\label{while-loop}

\begin{Shaded}
\begin{Highlighting}[]
\CodeTypeTok{int} count \OperatorTok{=} \CodeNumTok{1}\OperatorTok{;}
\KeywordTok{while} \OperatorTok{(}count \OperatorTok{\textless{}=} \CodeNumTok{5}\OperatorTok{)} \OperatorTok{\{}
    System\OperatorTok{.}\CodeFuncTok{out}\OperatorTok{.}\CodeFuncTok{println}\OperatorTok{(}\StringTok{"Count: "} \OperatorTok{+} count\OperatorTok{);}
    count\OperatorTok{++;}
\OperatorTok{\}}
\end{Highlighting}
\end{Shaded}

\subsection{\texorpdfstring{\texttt{do-while}
Loop}{do-while Loop}}\label{do-while-loop}

The \texttt{do-while} loop runs at least once before checking the
condition:

\begin{Shaded}
\begin{Highlighting}[]
Scanner scanner \OperatorTok{=} \KeywordTok{new} Scanner\OperatorTok{(}System\OperatorTok{.}\CodeFuncTok{in}\OperatorTok{);}
String input\OperatorTok{;}
\KeywordTok{do} \OperatorTok{\{}
    System\OperatorTok{.}\CodeFuncTok{out}\OperatorTok{.}\CodeFuncTok{print}\OperatorTok{(}\StringTok{"Enter a positive number: "}\OperatorTok{);}
    \CodeTypeTok{int} number \OperatorTok{=} scanner\OperatorTok{.}\CodeFuncTok{nextInt}\OperatorTok{();}
    System\OperatorTok{.}\CodeFuncTok{out}\OperatorTok{.}\CodeFuncTok{println}\OperatorTok{(}\StringTok{"You entered: "} \OperatorTok{+} number\OperatorTok{);}
    System\OperatorTok{.}\CodeFuncTok{out}\OperatorTok{.}\CodeFuncTok{print}\OperatorTok{(}\StringTok{"Continue? (yes/no): "}\OperatorTok{);}
    input \OperatorTok{=} scanner\OperatorTok{.}\CodeFuncTok{next}\OperatorTok{();}
\OperatorTok{\}} \KeywordTok{while} \OperatorTok{(}input\OperatorTok{.}\CodeFuncTok{equals}\OperatorTok{(}\StringTok{"yes"}\OperatorTok{));}
\end{Highlighting}
\end{Shaded}

\subsection{\texorpdfstring{\texttt{break} and
\texttt{continue}}{break and continue}}\label{break-and-continue-1}

\begin{Shaded}
\begin{Highlighting}[]
\CommentTok{// break {-}{-}{-} exit the loop entirely}
\KeywordTok{for} \OperatorTok{(}\CodeTypeTok{int} i \OperatorTok{=} \CodeNumTok{0}\OperatorTok{;} i \OperatorTok{\textless{}} \CodeNumTok{10}\OperatorTok{;} i\OperatorTok{++)} \OperatorTok{\{}
    \KeywordTok{if} \OperatorTok{(}i \OperatorTok{==} \CodeNumTok{5}\OperatorTok{)} \KeywordTok{break}\OperatorTok{;}
    System\OperatorTok{.}\CodeFuncTok{out}\OperatorTok{.}\CodeFuncTok{println}\OperatorTok{(}i\OperatorTok{);}   \CommentTok{// Prints 0, 1, 2, 3, 4}
\OperatorTok{\}}
 
\CommentTok{// continue {-}{-}{-} skip the current iteration}
\KeywordTok{for} \OperatorTok{(}\CodeTypeTok{int} i \OperatorTok{=} \CodeNumTok{0}\OperatorTok{;} i \OperatorTok{\textless{}} \CodeNumTok{10}\OperatorTok{;} i\OperatorTok{++)} \OperatorTok{\{}
    \KeywordTok{if} \OperatorTok{(}i \OperatorTok{\%} \CodeNumTok{2} \OperatorTok{==} \CodeNumTok{0}\OperatorTok{)} \KeywordTok{continue}\OperatorTok{;}
    System\OperatorTok{.}\CodeFuncTok{out}\OperatorTok{.}\CodeFuncTok{println}\OperatorTok{(}i\OperatorTok{);}   \CommentTok{// Prints 1, 3, 5, 7, 9}
\OperatorTok{\}}
\end{Highlighting}
\end{Shaded}

\subsection{Skills Checklist}\label{skills-checklist-76}

\begin{itemize}
\tightlist
\item[$\square$]
  I can write a \texttt{while} loop.
\item[$\square$]
  I can write a \texttt{do-while} loop.
\item[$\square$]
  I can use \texttt{break} to exit a loop early.
\item[$\square$]
  I can use \texttt{continue} to skip an iteration.
\end{itemize}

\begin{center}\rule{0.5\linewidth}{0.5pt}\end{center}

\section{Unit 4: Data Structures --- Arrays and
ArrayLists}\label{unit-4-data-structures-arrays-and-arraylists}

\subsection{Overview}\label{overview-18}

This unit introduces Java's built-in array type and the flexible
\texttt{ArrayList} class.

\begin{center}\rule{0.5\linewidth}{0.5pt}\end{center}

\subsection{Lesson 4.1: Arrays}\label{lesson-4.1-arrays-2}

\subsection{Declaring and Initializing
Arrays}\label{declaring-and-initializing-arrays}

\begin{Shaded}
\begin{Highlighting}[]
\CommentTok{// Declare and allocate}
\CodeTypeTok{int}\OperatorTok{[]} numbers \OperatorTok{=} \KeywordTok{new} \CodeTypeTok{int}\OperatorTok{[}\CodeNumTok{5}\OperatorTok{];}            \CommentTok{// [0, 0, 0, 0, 0]}
 
\CommentTok{// Declare and initialize with values}
String\OperatorTok{[]} colors \OperatorTok{=} \OperatorTok{\{}\StringTok{"red"}\OperatorTok{,} \StringTok{"blue"}\OperatorTok{,} \StringTok{"green"}\OperatorTok{,} \StringTok{"yellow"}\OperatorTok{\};}
 
\CommentTok{// Access elements (zero{-}based index)}
System\OperatorTok{.}\CodeFuncTok{out}\OperatorTok{.}\CodeFuncTok{println}\OperatorTok{(}colors\OperatorTok{[}\CodeNumTok{0}\OperatorTok{]);}         \CommentTok{// "red"}
System\OperatorTok{.}\CodeFuncTok{out}\OperatorTok{.}\CodeFuncTok{println}\OperatorTok{(}colors\OperatorTok{[}colors\OperatorTok{.}\CodeFuncTok{length} \OperatorTok{{-}} \CodeNumTok{1}\OperatorTok{]);}  \CommentTok{// "yellow"}
 
\CommentTok{// Modify an element}
colors\OperatorTok{[}\CodeNumTok{1}\OperatorTok{]} \OperatorTok{=} \StringTok{"purple"}\OperatorTok{;}
 
\CommentTok{// Loop through array}
\KeywordTok{for} \OperatorTok{(}\CodeTypeTok{int} i \OperatorTok{=} \CodeNumTok{0}\OperatorTok{;} i \OperatorTok{\textless{}} numbers\OperatorTok{.}\CodeFuncTok{length}\OperatorTok{;} i\OperatorTok{++)} \OperatorTok{\{}
    numbers\OperatorTok{[}i\OperatorTok{]} \OperatorTok{=} i \OperatorTok{*} \CodeNumTok{2}\OperatorTok{;}
    System\OperatorTok{.}\CodeFuncTok{out}\OperatorTok{.}\CodeFuncTok{println}\OperatorTok{(}\StringTok{"numbers["} \OperatorTok{+} i \OperatorTok{+} \StringTok{"] = "} \OperatorTok{+} numbers\OperatorTok{[}i\OperatorTok{]);}
\OperatorTok{\}}
\end{Highlighting}
\end{Shaded}

\subsection{Common Array Operations}\label{common-array-operations}

\begin{Shaded}
\begin{Highlighting}[]
\KeywordTok{import} \CodeAccentTok{java}\OperatorTok{.}\CodeAccentTok{util}\OperatorTok{.}\CodeAccentTok{Arrays}\OperatorTok{;}
 
\CodeTypeTok{int}\OperatorTok{[]} data \OperatorTok{=} \OperatorTok{\{}\CodeNumTok{5}\OperatorTok{,} \CodeNumTok{2}\OperatorTok{,} \CodeNumTok{9}\OperatorTok{,} \CodeNumTok{1}\OperatorTok{,} \CodeNumTok{7}\OperatorTok{,} \CodeNumTok{3}\OperatorTok{\};}
 
\CommentTok{// Sort}
Arrays\OperatorTok{.}\CodeFuncTok{sort}\OperatorTok{(}data\OperatorTok{);}
System\OperatorTok{.}\CodeFuncTok{out}\OperatorTok{.}\CodeFuncTok{println}\OperatorTok{(}Arrays\OperatorTok{.}\CodeFuncTok{toString}\OperatorTok{(}data\OperatorTok{));}   \CommentTok{// [1, 2, 3, 5, 7, 9]}
 
\CommentTok{// Search (array must be sorted first)}
\CodeTypeTok{int} index \OperatorTok{=} Arrays\OperatorTok{.}\CodeFuncTok{binarySearch}\OperatorTok{(}data\OperatorTok{,} \CodeNumTok{7}\OperatorTok{);}
System\OperatorTok{.}\CodeFuncTok{out}\OperatorTok{.}\CodeFuncTok{println}\OperatorTok{(}\StringTok{"Found 7 at index: "} \OperatorTok{+} index\OperatorTok{);}
 
\CommentTok{// Copy}
\CodeTypeTok{int}\OperatorTok{[]} copy \OperatorTok{=} Arrays\OperatorTok{.}\CodeFuncTok{copyOf}\OperatorTok{(}data\OperatorTok{,} data\OperatorTok{.}\CodeFuncTok{length}\OperatorTok{);}
\CodeTypeTok{int}\OperatorTok{[]} partial \OperatorTok{=} Arrays\OperatorTok{.}\CodeFuncTok{copyOfRange}\OperatorTok{(}data\OperatorTok{,} \CodeNumTok{1}\OperatorTok{,} \CodeNumTok{4}\OperatorTok{);}  \CommentTok{// [2, 3, 5]}
\end{Highlighting}
\end{Shaded}

\subsection{Skills Checklist}\label{skills-checklist-77}

\begin{itemize}
\tightlist
\item[$\square$]
  I can declare and initialize an array in Java.
\item[$\square$]
  I can access and modify array elements by index.
\item[$\square$]
  I can iterate over an array using a \texttt{for} loop.
\item[$\square$]
  I can use \texttt{Arrays.sort()} and \texttt{Arrays.toString()}.
\end{itemize}

\begin{center}\rule{0.5\linewidth}{0.5pt}\end{center}

\amn{\textbf{Accessibility Tip:} Prefer \texttt{ArrayList} over arrays. To inspect contents accessibly, iterate with index and print one item per line: \texttt{for (int i = 0; i < list.size(); i++) System.out.println((i+1) + ". " + list.get(i));}. Numbered output is easy to navigate line-by-line with a screen reader.}
\subsection{Lesson 4.2: ArrayLists}\label{lesson-4.2-arraylists-1}

\texttt{ArrayList} is a resizable array that can grow or shrink
dynamically.

\begin{Shaded}
\begin{Highlighting}[]
\KeywordTok{import} \CodeAccentTok{java}\OperatorTok{.}\CodeAccentTok{util}\OperatorTok{.}\CodeAccentTok{ArrayList}\OperatorTok{;}
 
ArrayList\OperatorTok{\textless{}}String\OperatorTok{\textgreater{}} shoppingList \OperatorTok{=} \KeywordTok{new} ArrayList\OperatorTok{\textless{}\textgreater{}();}
 
\CommentTok{// Add items}
shoppingList\OperatorTok{.}\CodeFuncTok{add}\OperatorTok{(}\StringTok{"milk"}\OperatorTok{);}
shoppingList\OperatorTok{.}\CodeFuncTok{add}\OperatorTok{(}\StringTok{"bread"}\OperatorTok{);}
shoppingList\OperatorTok{.}\CodeFuncTok{add}\OperatorTok{(}\StringTok{"eggs"}\OperatorTok{);}
shoppingList\OperatorTok{.}\CodeFuncTok{add}\OperatorTok{(}\CodeNumTok{1}\OperatorTok{,} \StringTok{"butter"}\OperatorTok{);}   \CommentTok{// Insert at index 1}
 
\CommentTok{// Access items}
System\OperatorTok{.}\CodeFuncTok{out}\OperatorTok{.}\CodeFuncTok{println}\OperatorTok{(}shoppingList\OperatorTok{.}\CodeFuncTok{get}\OperatorTok{(}\CodeNumTok{0}\OperatorTok{));}          \CommentTok{// "milk"}
System\OperatorTok{.}\CodeFuncTok{out}\OperatorTok{.}\CodeFuncTok{println}\OperatorTok{(}shoppingList\OperatorTok{.}\CodeFuncTok{size}\OperatorTok{());}           \CommentTok{// 4}
 
\CommentTok{// Remove items}
shoppingList\OperatorTok{.}\CodeFuncTok{remove}\OperatorTok{(}\StringTok{"bread"}\OperatorTok{);}                     \CommentTok{// Remove by value}
shoppingList\OperatorTok{.}\CodeFuncTok{remove}\OperatorTok{(}\CodeNumTok{0}\OperatorTok{);}                           \CommentTok{// Remove by index}
 
\CommentTok{// Check membership}
\CodeTypeTok{boolean} hasMilk \OperatorTok{=} shoppingList\OperatorTok{.}\CodeFuncTok{contains}\OperatorTok{(}\StringTok{"milk"}\OperatorTok{);}
 
\CommentTok{// Loop through}
\KeywordTok{for} \OperatorTok{(}String item \OperatorTok{:} shoppingList\OperatorTok{)} \OperatorTok{\{}
    System\OperatorTok{.}\CodeFuncTok{out}\OperatorTok{.}\CodeFuncTok{println}\OperatorTok{(}\StringTok{"{-} "} \OperatorTok{+} item\OperatorTok{);}
\OperatorTok{\}}
\end{Highlighting}
\end{Shaded}

\subsection{ArrayList vs.~Array}\label{arraylist-vs.-array}

\begin{fullwidth}{\def\LTcaptype{none} % do not increment counter
\begin{longtable}[]{@{}lll@{}}
\toprule\noalign{}
Feature & \texttt{Array} & \texttt{ArrayList} \\
\midrule\noalign{}
\endhead
\bottomrule\noalign{}
\endlastfoot
Size & Fixed & Dynamic \\
Type & Primitive or Object & Objects only \\
Syntax & \texttt{int{[}{]}} &
\texttt{ArrayList\textless{}Integer\textgreater{}} \\
\texttt{.length} / \texttt{.size()} & \texttt{.length} &
\texttt{.size()} \\
Performance & Slightly faster & More flexible \\
\end{longtable}
}
\end{fullwidth}

\subsection{Skills Checklist}\label{skills-checklist-78}

\begin{itemize}
\tightlist
\item[$\square$]
  I can create an \texttt{ArrayList} with a generic type.
\item[$\square$]
  I can add, access, remove, and search items in an \texttt{ArrayList}.
\item[$\square$]
  I can explain when to use an array vs.~an \texttt{ArrayList}.
\end{itemize}

\begin{center}\rule{0.5\linewidth}{0.5pt}\end{center}

\subsection{Lesson 4.3: Multidimensional
Arrays}\label{lesson-4.3-multidimensional-arrays}

\begin{Shaded}
\begin{Highlighting}[]
\CommentTok{// 2D array (matrix)}
\CodeTypeTok{int}\OperatorTok{[][]} matrix \OperatorTok{=} \OperatorTok{\{}
    \OperatorTok{\{}\CodeNumTok{1}\OperatorTok{,} \CodeNumTok{2}\OperatorTok{,} \CodeNumTok{3}\OperatorTok{\},}
    \OperatorTok{\{}\CodeNumTok{4}\OperatorTok{,} \CodeNumTok{5}\OperatorTok{,} \CodeNumTok{6}\OperatorTok{\},}
    \OperatorTok{\{}\CodeNumTok{7}\OperatorTok{,} \CodeNumTok{8}\OperatorTok{,} \CodeNumTok{9}\OperatorTok{\}}
\OperatorTok{\};}
 
\CommentTok{// Access element at row 1, column 2}
System\OperatorTok{.}\CodeFuncTok{out}\OperatorTok{.}\CodeFuncTok{println}\OperatorTok{(}matrix\OperatorTok{[}\CodeNumTok{1}\OperatorTok{][}\CodeNumTok{2}\OperatorTok{]);}   \CommentTok{// 6}
 
\CommentTok{// Loop through 2D array}
\KeywordTok{for} \OperatorTok{(}\CodeTypeTok{int} row \OperatorTok{=} \CodeNumTok{0}\OperatorTok{;} row \OperatorTok{\textless{}} matrix\OperatorTok{.}\CodeFuncTok{length}\OperatorTok{;} row\OperatorTok{++)} \OperatorTok{\{}
    \KeywordTok{for} \OperatorTok{(}\CodeTypeTok{int} col \OperatorTok{=} \CodeNumTok{0}\OperatorTok{;} col \OperatorTok{\textless{}} matrix\OperatorTok{[}row\OperatorTok{].}\CodeFuncTok{length}\OperatorTok{;} col\OperatorTok{++)} \OperatorTok{\{}
        System\OperatorTok{.}\CodeFuncTok{out}\OperatorTok{.}\CodeFuncTok{print}\OperatorTok{(}matrix\OperatorTok{[}row\OperatorTok{][}col\OperatorTok{]} \OperatorTok{+} \StringTok{" "}\OperatorTok{);}
    \OperatorTok{\}}
    System\OperatorTok{.}\CodeFuncTok{out}\OperatorTok{.}\CodeFuncTok{println}\OperatorTok{();}
\OperatorTok{\}}
\end{Highlighting}
\end{Shaded}

\subsection{Skills Checklist}\label{skills-checklist-79}

\begin{itemize}
\tightlist
\item[$\square$]
  I can declare and initialize a 2D array.
\item[$\square$]
  I can access elements using row and column indices.
\item[$\square$]
  I can iterate through a 2D array using nested loops.
\end{itemize}

\newpage

\section{Unit 5: Strings and String
Processing}\label{unit-5-strings-and-string-processing-1}

\subsection{Overview}\label{overview-19}

Java's \texttt{String} class is rich with methods for text manipulation.

\begin{center}\rule{0.5\linewidth}{0.5pt}\end{center}

\subsection{Lesson 5.1: Working with
Strings}\label{lesson-5.1-working-with-strings-1}

\subsection{String Basics}\label{string-basics}

\begin{Shaded}
\begin{Highlighting}[]
String name \OperatorTok{=} \StringTok{"Alice Smith"}\OperatorTok{;}
 
\CommentTok{// Length}
System\OperatorTok{.}\CodeFuncTok{out}\OperatorTok{.}\CodeFuncTok{println}\OperatorTok{(}name\OperatorTok{.}\CodeFuncTok{length}\OperatorTok{());}           \CommentTok{// 11}
 
\CommentTok{// Accessing characters}
System\OperatorTok{.}\CodeFuncTok{out}\OperatorTok{.}\CodeFuncTok{println}\OperatorTok{(}name\OperatorTok{.}\CodeFuncTok{charAt}\OperatorTok{(}\CodeNumTok{0}\OperatorTok{));}          \CommentTok{// \textquotesingle{}A\textquotesingle{}}
System\OperatorTok{.}\CodeFuncTok{out}\OperatorTok{.}\CodeFuncTok{println}\OperatorTok{(}name\OperatorTok{.}\CodeFuncTok{charAt}\OperatorTok{(}name\OperatorTok{.}\CodeFuncTok{length}\OperatorTok{()} \OperatorTok{{-}} \CodeNumTok{1}\OperatorTok{));}  \CommentTok{// \textquotesingle{}h\textquotesingle{}}
 
\CommentTok{// Substrings}
System\OperatorTok{.}\CodeFuncTok{out}\OperatorTok{.}\CodeFuncTok{println}\OperatorTok{(}name\OperatorTok{.}\CodeFuncTok{substring}\OperatorTok{(}\CodeNumTok{6}\OperatorTok{));}       \CommentTok{// "Smith"}
System\OperatorTok{.}\CodeFuncTok{out}\OperatorTok{.}\CodeFuncTok{println}\OperatorTok{(}name\OperatorTok{.}\CodeFuncTok{substring}\OperatorTok{(}\CodeNumTok{0}\OperatorTok{,} \CodeNumTok{5}\OperatorTok{));}    \CommentTok{// "Alice"}
 
\CommentTok{// Case conversion}
System\OperatorTok{.}\CodeFuncTok{out}\OperatorTok{.}\CodeFuncTok{println}\OperatorTok{(}name\OperatorTok{.}\CodeFuncTok{toUpperCase}\OperatorTok{());}      \CommentTok{// "ALICE SMITH"}
System\OperatorTok{.}\CodeFuncTok{out}\OperatorTok{.}\CodeFuncTok{println}\OperatorTok{(}name\OperatorTok{.}\CodeFuncTok{toLowerCase}\OperatorTok{());}      \CommentTok{// "alice smith"}
 
\CommentTok{// Trimming whitespace}
String padded \OperatorTok{=} \StringTok{"  hello  "}\OperatorTok{;}
System\OperatorTok{.}\CodeFuncTok{out}\OperatorTok{.}\CodeFuncTok{println}\OperatorTok{(}padded\OperatorTok{.}\CodeFuncTok{trim}\OperatorTok{());}           \CommentTok{// "hello"}
System\OperatorTok{.}\CodeFuncTok{out}\OperatorTok{.}\CodeFuncTok{println}\OperatorTok{(}padded\OperatorTok{.}\CodeFuncTok{strip}\OperatorTok{());}          \CommentTok{// "hello" (better Unicode support)}
 
\CommentTok{// Searching}
System\OperatorTok{.}\CodeFuncTok{out}\OperatorTok{.}\CodeFuncTok{println}\OperatorTok{(}name\OperatorTok{.}\CodeFuncTok{indexOf}\OperatorTok{(}\StringTok{"Smith"}\OperatorTok{));}   \CommentTok{// 6}
System\OperatorTok{.}\CodeFuncTok{out}\OperatorTok{.}\CodeFuncTok{println}\OperatorTok{(}name\OperatorTok{.}\CodeFuncTok{contains}\OperatorTok{(}\StringTok{"Alice"}\OperatorTok{));}  \CommentTok{// true}
System\OperatorTok{.}\CodeFuncTok{out}\OperatorTok{.}\CodeFuncTok{println}\OperatorTok{(}name\OperatorTok{.}\CodeFuncTok{startsWith}\OperatorTok{(}\StringTok{"Ali"}\OperatorTok{));}  \CommentTok{// true}
System\OperatorTok{.}\CodeFuncTok{out}\OperatorTok{.}\CodeFuncTok{println}\OperatorTok{(}name\OperatorTok{.}\CodeFuncTok{endsWith}\OperatorTok{(}\StringTok{"ith"}\OperatorTok{));}    \CommentTok{// true}
 
\CommentTok{// Replacing}
String result \OperatorTok{=} name\OperatorTok{.}\CodeFuncTok{replace}\OperatorTok{(}\StringTok{"Alice"}\OperatorTok{,} \StringTok{"Bob"}\OperatorTok{);}   \CommentTok{// "Bob Smith"}
String clean \OperatorTok{=} \StringTok{"a,b,c"}\OperatorTok{.}\CodeFuncTok{replace}\OperatorTok{(}\StringTok{","}\OperatorTok{,} \StringTok{" "}\OperatorTok{);}       \CommentTok{// "a b c"}
 
\CommentTok{// Splitting}
String sentence \OperatorTok{=} \StringTok{"the quick brown fox"}\OperatorTok{;}
String\OperatorTok{[]} words \OperatorTok{=} sentence\OperatorTok{.}\CodeFuncTok{split}\OperatorTok{(}\StringTok{" "}\OperatorTok{);}
\KeywordTok{for} \OperatorTok{(}String word \OperatorTok{:} words\OperatorTok{)} \OperatorTok{\{}
    System\OperatorTok{.}\CodeFuncTok{out}\OperatorTok{.}\CodeFuncTok{println}\OperatorTok{(}word\OperatorTok{);}
\OperatorTok{\}}
 
\CommentTok{// Checking equality}
String s1 \OperatorTok{=} \StringTok{"hello"}\OperatorTok{;}
String s2 \OperatorTok{=} \StringTok{"hello"}\OperatorTok{;}
System\OperatorTok{.}\CodeFuncTok{out}\OperatorTok{.}\CodeFuncTok{println}\OperatorTok{(}s1\OperatorTok{.}\CodeFuncTok{equals}\OperatorTok{(}s2\OperatorTok{));}           \CommentTok{// true (correct)}
System\OperatorTok{.}\CodeFuncTok{out}\OperatorTok{.}\CodeFuncTok{println}\OperatorTok{(}s1\OperatorTok{.}\CodeFuncTok{equalsIgnoreCase}\OperatorTok{(}\StringTok{"HELLO"}\OperatorTok{));} \CommentTok{// true}
 
\CommentTok{// Converting other types to String}
String fromInt \OperatorTok{=} String\OperatorTok{.}\CodeFuncTok{valueOf}\OperatorTok{(}\CodeNumTok{42}\OperatorTok{);}
String fromDouble \OperatorTok{=} String\OperatorTok{.}\CodeFuncTok{valueOf}\OperatorTok{(}\CodeNumTok{3.14}\OperatorTok{);}
String formatted \OperatorTok{=} String\OperatorTok{.}\CodeFuncTok{format}\OperatorTok{(}\StringTok{"Name: }\StringTok{\%{-}10s}\StringTok{ Age: }\StringTok{\%3d}\StringTok{"}\OperatorTok{,} \StringTok{"Alice"}\OperatorTok{,} \CodeNumTok{15}\OperatorTok{);}
\end{Highlighting}
\end{Shaded}

\subsection{\texorpdfstring{\texttt{StringBuilder} for Efficient
Concatenation}{StringBuilder for Efficient Concatenation}}\label{stringbuilder-for-efficient-concatenation}

Use \texttt{StringBuilder} when building strings in a loop (more
efficient than \texttt{+}):

\begin{Shaded}
\begin{Highlighting}[]
StringBuilder sb \OperatorTok{=} \KeywordTok{new} StringBuilder\OperatorTok{();}
\KeywordTok{for} \OperatorTok{(}\CodeTypeTok{int} i \OperatorTok{=} \CodeNumTok{1}\OperatorTok{;} i \OperatorTok{\textless{}=} \CodeNumTok{5}\OperatorTok{;} i\OperatorTok{++)} \OperatorTok{\{}
    sb\OperatorTok{.}\CodeFuncTok{append}\OperatorTok{(}\StringTok{"Item "}\OperatorTok{).}\CodeFuncTok{append}\OperatorTok{(}i\OperatorTok{).}\CodeFuncTok{append}\OperatorTok{(}\StringTok{"}\StringTok{\textbackslash{}n}\StringTok{"}\OperatorTok{);}
\OperatorTok{\}}
System\OperatorTok{.}\CodeFuncTok{out}\OperatorTok{.}\CodeFuncTok{println}\OperatorTok{(}sb\OperatorTok{.}\CodeFuncTok{toString}\OperatorTok{());}
\end{Highlighting}
\end{Shaded}

\subsection{Examples}\label{examples-25}

\textbf{Example 5.1.1: String Analysis}

\begin{Shaded}
\begin{Highlighting}[]
\KeywordTok{import} \CodeAccentTok{java}\OperatorTok{.}\CodeAccentTok{util}\OperatorTok{.}\CodeAccentTok{Scanner}\OperatorTok{;}
 
\KeywordTok{public} \KeywordTok{class} StringAnalysis \OperatorTok{\{}
    \KeywordTok{public} \CodeTypeTok{static} \CodeTypeTok{void} \CodeFuncTok{main}\OperatorTok{(}String\OperatorTok{[]} args\OperatorTok{)} \OperatorTok{\{}
        Scanner scanner \OperatorTok{=} \KeywordTok{new} Scanner\OperatorTok{(}System\OperatorTok{.}\CodeFuncTok{in}\OperatorTok{);}
        System\OperatorTok{.}\CodeFuncTok{out}\OperatorTok{.}\CodeFuncTok{print}\OperatorTok{(}\StringTok{"Enter a sentence: "}\OperatorTok{);}
        String sentence \OperatorTok{=} scanner\OperatorTok{.}\CodeFuncTok{nextLine}\OperatorTok{();}
 
        String\OperatorTok{[]} words \OperatorTok{=} sentence\OperatorTok{.}\CodeFuncTok{trim}\OperatorTok{().}\CodeFuncTok{split}\OperatorTok{(}\StringTok{"}\StringTok{\textbackslash{}\textbackslash{}}\StringTok{s+"}\OperatorTok{);}
 
        System\OperatorTok{.}\CodeFuncTok{out}\OperatorTok{.}\CodeFuncTok{println}\OperatorTok{(}\StringTok{"Original: "} \OperatorTok{+} sentence\OperatorTok{);}
        System\OperatorTok{.}\CodeFuncTok{out}\OperatorTok{.}\CodeFuncTok{println}\OperatorTok{(}\StringTok{"Character count: "} \OperatorTok{+} sentence\OperatorTok{.}\CodeFuncTok{length}\OperatorTok{());}
        System\OperatorTok{.}\CodeFuncTok{out}\OperatorTok{.}\CodeFuncTok{println}\OperatorTok{(}\StringTok{"Word count: "} \OperatorTok{+} words\OperatorTok{.}\CodeFuncTok{length}\OperatorTok{);}
        System\OperatorTok{.}\CodeFuncTok{out}\OperatorTok{.}\CodeFuncTok{println}\OperatorTok{(}\StringTok{"Uppercase: "} \OperatorTok{+} sentence\OperatorTok{.}\CodeFuncTok{toUpperCase}\OperatorTok{());}
        System\OperatorTok{.}\CodeFuncTok{out}\OperatorTok{.}\CodeFuncTok{println}\OperatorTok{(}\StringTok{"Lowercase: "} \OperatorTok{+} sentence\OperatorTok{.}\CodeFuncTok{toLowerCase}\OperatorTok{());}
 
        scanner\OperatorTok{.}\CodeFuncTok{close}\OperatorTok{();}
    \OperatorTok{\}}
\OperatorTok{\}}
\end{Highlighting}
\end{Shaded}

\textbf{Example 5.1.2: Word Frequency Counter}

\begin{Shaded}
\begin{Highlighting}[]
\KeywordTok{import} \CodeAccentTok{java}\OperatorTok{.}\CodeAccentTok{util}\OperatorTok{.}\CodeAccentTok{HashMap}\OperatorTok{;}
\KeywordTok{import} \CodeAccentTok{java}\OperatorTok{.}\CodeAccentTok{util}\OperatorTok{.}\CodeAccentTok{Map}\OperatorTok{;}
 
\KeywordTok{public} \KeywordTok{class} WordFrequency \OperatorTok{\{}
    \KeywordTok{public} \CodeTypeTok{static} \CodeTypeTok{void} \CodeFuncTok{main}\OperatorTok{(}String\OperatorTok{[]} args\OperatorTok{)} \OperatorTok{\{}
        String text \OperatorTok{=} \StringTok{"the quick brown fox jumps over the lazy dog the quick brown fox"}\OperatorTok{;}
        String\OperatorTok{[]} words \OperatorTok{=} text\OperatorTok{.}\CodeFuncTok{toLowerCase}\OperatorTok{().}\CodeFuncTok{split}\OperatorTok{(}\StringTok{"}\StringTok{\textbackslash{}\textbackslash{}}\StringTok{s+"}\OperatorTok{);}
 
        Map\OperatorTok{\textless{}}String\OperatorTok{,} Integer\OperatorTok{\textgreater{}} frequency \OperatorTok{=} \KeywordTok{new} HashMap\OperatorTok{\textless{}\textgreater{}();}
        \KeywordTok{for} \OperatorTok{(}String word \OperatorTok{:} words\OperatorTok{)} \OperatorTok{\{}
            frequency\OperatorTok{.}\CodeFuncTok{put}\OperatorTok{(}word\OperatorTok{,} frequency\OperatorTok{.}\CodeFuncTok{getOrDefault}\OperatorTok{(}word\OperatorTok{,} \CodeNumTok{0}\OperatorTok{)} \OperatorTok{+} \CodeNumTok{1}\OperatorTok{);}
        \OperatorTok{\}}
 
        System\OperatorTok{.}\CodeFuncTok{out}\OperatorTok{.}\CodeFuncTok{println}\OperatorTok{(}\StringTok{"Word Frequencies:"}\OperatorTok{);}
        \KeywordTok{for} \OperatorTok{(}Map\OperatorTok{.}\CodeFuncTok{Entry}\OperatorTok{\textless{}}String\OperatorTok{,} Integer\OperatorTok{\textgreater{}} entry \OperatorTok{:} frequency\OperatorTok{.}\CodeFuncTok{entrySet}\OperatorTok{())} \OperatorTok{\{}
            System\OperatorTok{.}\CodeFuncTok{out}\OperatorTok{.}\CodeFuncTok{println}\OperatorTok{(}\StringTok{"\textquotesingle{}"} \OperatorTok{+} entry\OperatorTok{.}\CodeFuncTok{getKey}\OperatorTok{()} \OperatorTok{+} \StringTok{"\textquotesingle{}: "} \OperatorTok{+} entry\OperatorTok{.}\CodeFuncTok{getValue}\OperatorTok{());}
        \OperatorTok{\}}
    \OperatorTok{\}}
\OperatorTok{\}}
\end{Highlighting}
\end{Shaded}

\subsection{Skills Checklist}\label{skills-checklist-80}

\begin{itemize}
\tightlist
\item[$\square$]
  I can use common \texttt{String} methods in Java.
\item[$\square$]
  I can compare strings using \texttt{.equals()}.
\item[$\square$]
  I can split a string into an array of words.
\item[$\square$]
  I can use \texttt{StringBuilder} for efficient string building.
\end{itemize}

\begin{center}\rule{0.5\linewidth}{0.5pt}\end{center}

\section{Unit 6: Methods (Functions)}\label{unit-6-methods-functions}

\subsection{Overview}\label{overview-20}

Java organizes reusable code into \textbf{methods} (the Java equivalent
of Python functions).

\begin{center}\rule{0.5\linewidth}{0.5pt}\end{center}

\subsection{Lesson 6.1: Defining and Calling
Methods}\amn{\textbf{Accessibility Tip:} Write a Javadoc comment above every method: \texttt{/** Returns the area of a circle. @param r radius @return area */}. VS Code's IntelliSense reads this comment aloud when you type the method name. This is how you document your methods for audio navigation.}
\subsection{Lesson 6.1: Defining and Calling
Methods}\label{lesson-6.1-defining-and-calling-methods}

\subsection{Method Syntax}\label{method-syntax}

\begin{Shaded}
\begin{Highlighting}[]
accessModifier returnType \CodeFuncTok{methodName}\OperatorTok{(}paramType paramName\OperatorTok{,} \KeywordTok{...}\OperatorTok{)} \OperatorTok{\{}
    \CommentTok{// method body}
    \KeywordTok{return} value\OperatorTok{;}  \CommentTok{// if returnType is not void}
\OperatorTok{\}}
\end{Highlighting}
\end{Shaded}

\subsection{Examples}\label{examples-26}

\textbf{Example 6.1.1: Simple Greeting Method}

\begin{Shaded}
\begin{Highlighting}[]
\KeywordTok{public} \KeywordTok{class} GreetingDemo \OperatorTok{\{}
 
    \KeywordTok{public} \CodeTypeTok{static} \CodeTypeTok{void} \CodeFuncTok{greetUser}\OperatorTok{()} \OperatorTok{\{}
        System\OperatorTok{.}\CodeFuncTok{out}\OperatorTok{.}\CodeFuncTok{println}\OperatorTok{(}\StringTok{"Hello, Java programmer!"}\OperatorTok{);}
    \OperatorTok{\}}
 
    \KeywordTok{public} \CodeTypeTok{static} \CodeTypeTok{void} \CodeFuncTok{greetByName}\OperatorTok{(}String name\OperatorTok{)} \OperatorTok{\{}
        System\OperatorTok{.}\CodeFuncTok{out}\OperatorTok{.}\CodeFuncTok{println}\OperatorTok{(}\StringTok{"Hello, "} \OperatorTok{+} name \OperatorTok{+} \StringTok{"!"}\OperatorTok{);}
    \OperatorTok{\}}
 
    \KeywordTok{public} \CodeTypeTok{static} \CodeTypeTok{void} \CodeFuncTok{main}\OperatorTok{(}String\OperatorTok{[]} args\OperatorTok{)} \OperatorTok{\{}
        \CodeFuncTok{greetUser}\OperatorTok{();}
        \CodeFuncTok{greetByName}\OperatorTok{(}\StringTok{"Alice"}\OperatorTok{);}
        \CodeFuncTok{greetByName}\OperatorTok{(}\StringTok{"Bob"}\OperatorTok{);}
    \OperatorTok{\}}
\OperatorTok{\}}
\end{Highlighting}
\end{Shaded}

\textbf{Example 6.1.2: Return Values}

\begin{Shaded}
\begin{Highlighting}[]
\KeywordTok{public} \KeywordTok{class} Calculator \OperatorTok{\{}
 
    \KeywordTok{public} \CodeTypeTok{static} \CodeTypeTok{int} \CodeFuncTok{addNumbers}\OperatorTok{(}\CodeTypeTok{int} a\OperatorTok{,} \CodeTypeTok{int} b\OperatorTok{)} \OperatorTok{\{}
        \KeywordTok{return} a \OperatorTok{+} b\OperatorTok{;}
    \OperatorTok{\}}
 
    \KeywordTok{public} \CodeTypeTok{static} \CodeTypeTok{double} \CodeFuncTok{calculateCircleArea}\OperatorTok{(}\CodeTypeTok{double} radius\OperatorTok{)} \OperatorTok{\{}
        \KeywordTok{return} Math\OperatorTok{.}\CodeFuncTok{PI} \OperatorTok{*} radius \OperatorTok{*} radius\OperatorTok{;}
    \OperatorTok{\}}
 
    \KeywordTok{public} \CodeTypeTok{static} String \CodeFuncTok{getGrade}\OperatorTok{(}\CodeTypeTok{int} score\OperatorTok{)} \OperatorTok{\{}
        \KeywordTok{if} \OperatorTok{(}score \OperatorTok{\textgreater{}=} \CodeNumTok{90}\OperatorTok{)} \KeywordTok{return} \StringTok{"A"}\OperatorTok{;}
        \KeywordTok{else} \KeywordTok{if} \OperatorTok{(}score \OperatorTok{\textgreater{}=} \CodeNumTok{80}\OperatorTok{)} \KeywordTok{return} \StringTok{"B"}\OperatorTok{;}
        \KeywordTok{else} \KeywordTok{if} \OperatorTok{(}score \OperatorTok{\textgreater{}=} \CodeNumTok{70}\OperatorTok{)} \KeywordTok{return} \StringTok{"C"}\OperatorTok{;}
        \KeywordTok{else} \KeywordTok{if} \OperatorTok{(}score \OperatorTok{\textgreater{}=} \CodeNumTok{60}\OperatorTok{)} \KeywordTok{return} \StringTok{"D"}\OperatorTok{;}
        \KeywordTok{else} \KeywordTok{return} \StringTok{"F"}\OperatorTok{;}
    \OperatorTok{\}}
 
    \KeywordTok{public} \CodeTypeTok{static} \CodeTypeTok{void} \CodeFuncTok{main}\OperatorTok{(}String\OperatorTok{[]} args\OperatorTok{)} \OperatorTok{\{}
        \CodeTypeTok{int} sum \OperatorTok{=} \CodeFuncTok{addNumbers}\OperatorTok{(}\CodeNumTok{5}\OperatorTok{,} \CodeNumTok{3}\OperatorTok{);}
        System\OperatorTok{.}\CodeFuncTok{out}\OperatorTok{.}\CodeFuncTok{println}\OperatorTok{(}\StringTok{"Sum: "} \OperatorTok{+} sum\OperatorTok{);}
 
        \CodeTypeTok{double} area \OperatorTok{=} \CodeFuncTok{calculateCircleArea}\OperatorTok{(}\CodeNumTok{5.0}\OperatorTok{);}
        System\OperatorTok{.}\CodeFuncTok{out}\OperatorTok{.}\CodeFuncTok{printf}\OperatorTok{(}\StringTok{"Circle area: }\StringTok{\%.2f\%n}\StringTok{"}\OperatorTok{,} area\OperatorTok{);}
 
        System\OperatorTok{.}\CodeFuncTok{out}\OperatorTok{.}\CodeFuncTok{println}\OperatorTok{(}\StringTok{"Grade for 85: "} \OperatorTok{+} \CodeFuncTok{getGrade}\OperatorTok{(}\CodeNumTok{85}\OperatorTok{));}
    \OperatorTok{\}}
\OperatorTok{\}}
\end{Highlighting}
\end{Shaded}

\subsection{Method Overloading}\label{method-overloading}

Java allows multiple methods with the same name but different
parameters:

\begin{Shaded}
\begin{Highlighting}[]
\KeywordTok{public} \CodeTypeTok{static} \CodeTypeTok{int} \CodeFuncTok{add}\OperatorTok{(}\CodeTypeTok{int} a\OperatorTok{,} \CodeTypeTok{int} b\OperatorTok{)} \OperatorTok{\{} \KeywordTok{return} a \OperatorTok{+} b\OperatorTok{;} \OperatorTok{\}}
\KeywordTok{public} \CodeTypeTok{static} \CodeTypeTok{double} \CodeFuncTok{add}\OperatorTok{(}\CodeTypeTok{double} a\OperatorTok{,} \CodeTypeTok{double} b\OperatorTok{)} \OperatorTok{\{} \KeywordTok{return} a \OperatorTok{+} b\OperatorTok{;} \OperatorTok{\}}
\KeywordTok{public} \CodeTypeTok{static} \CodeTypeTok{int} \CodeFuncTok{add}\OperatorTok{(}\CodeTypeTok{int} a\OperatorTok{,} \CodeTypeTok{int} b\OperatorTok{,} \CodeTypeTok{int} c\OperatorTok{)} \OperatorTok{\{} \KeywordTok{return} a \OperatorTok{+} b \OperatorTok{+} c\OperatorTok{;} \OperatorTok{\}}
\end{Highlighting}
\end{Shaded}

\subsection{Passing Arrays to
Methods}\label{passing-arrays-to-methods}

\begin{Shaded}
\begin{Highlighting}[]
\KeywordTok{public} \CodeTypeTok{static} \CodeTypeTok{int} \CodeFuncTok{sum}\OperatorTok{(}\CodeTypeTok{int}\OperatorTok{[]} numbers\OperatorTok{)} \OperatorTok{\{}
    \CodeTypeTok{int} total \OperatorTok{=} \CodeNumTok{0}\OperatorTok{;}
    \KeywordTok{for} \OperatorTok{(}\CodeTypeTok{int} num \OperatorTok{:} numbers\OperatorTok{)} \OperatorTok{\{}
        total \OperatorTok{+=} num\OperatorTok{;}
    \OperatorTok{\}}
    \KeywordTok{return} total\OperatorTok{;}
\OperatorTok{\}}
 
\CommentTok{// Calling:}
\CodeTypeTok{int}\OperatorTok{[]} data \OperatorTok{=} \OperatorTok{\{}\CodeNumTok{1}\OperatorTok{,} \CodeNumTok{2}\OperatorTok{,} \CodeNumTok{3}\OperatorTok{,} \CodeNumTok{4}\OperatorTok{,} \CodeNumTok{5}\OperatorTok{\};}
System\OperatorTok{.}\CodeFuncTok{out}\OperatorTok{.}\CodeFuncTok{println}\OperatorTok{(}\StringTok{"Sum: "} \OperatorTok{+} \CodeFuncTok{sum}\OperatorTok{(}data\OperatorTok{));}
\end{Highlighting}
\end{Shaded}

\subsection{Varargs (Variable
Arguments)}\label{varargs-variable-arguments}

\begin{Shaded}
\begin{Highlighting}[]
\KeywordTok{public} \CodeTypeTok{static} \CodeTypeTok{int} \CodeFuncTok{sumAll}\OperatorTok{(}\CodeTypeTok{int}\KeywordTok{...} numbers\OperatorTok{)} \OperatorTok{\{}
    \CodeTypeTok{int} total \OperatorTok{=} \CodeNumTok{0}\OperatorTok{;}
    \KeywordTok{for} \OperatorTok{(}\CodeTypeTok{int} num \OperatorTok{:} numbers\OperatorTok{)} \OperatorTok{\{}
        total \OperatorTok{+=} num\OperatorTok{;}
    \OperatorTok{\}}
    \KeywordTok{return} total\OperatorTok{;}
\OperatorTok{\}}
 
\CommentTok{// Calling:}
System\OperatorTok{.}\CodeFuncTok{out}\OperatorTok{.}\CodeFuncTok{println}\OperatorTok{(}\CodeFuncTok{sumAll}\OperatorTok{(}\CodeNumTok{1}\OperatorTok{,} \CodeNumTok{2}\OperatorTok{,} \CodeNumTok{3}\OperatorTok{));}          \CommentTok{// 6}
System\OperatorTok{.}\CodeFuncTok{out}\OperatorTok{.}\CodeFuncTok{println}\OperatorTok{(}\CodeFuncTok{sumAll}\OperatorTok{(}\CodeNumTok{10}\OperatorTok{,} \CodeNumTok{20}\OperatorTok{,} \CodeNumTok{30}\OperatorTok{,} \CodeNumTok{40}\OperatorTok{));}    \CommentTok{// 100}
\end{Highlighting}
\end{Shaded}

\subsection{Skills Checklist}\label{skills-checklist-81}

\begin{itemize}
\tightlist
\item[$\square$]
  I can define a method with a return type and parameters.
\item[$\square$]
  I can call a method and use its return value.
\item[$\square$]
  I understand the difference between \texttt{void} and
  non-\texttt{void} methods.
\item[$\square$]
  I can overload methods with different parameter lists.
\item[$\square$]
  I can pass arrays to methods.
\end{itemize}

\begin{center}\rule{0.5\linewidth}{0.5pt}\end{center}

\section{Unit 7: Classes and Object-Oriented
Programming}\label{unit-7-classes-and-object-oriented-programming-3}

\subsection{Overview}\label{overview-21}

Java is fundamentally object-oriented. This unit covers defining and
using classes.

\begin{center}\rule{0.5\linewidth}{0.5pt}\end{center}

\subsection{Lesson 7.1: Why Learn About
Classes?}\label{lesson-7.1-why-learn-about-classes-1}
\amn{\textbf{Accessibility Tip:} Override \texttt{toString()} on every class you create so that \texttt{System.out.println(myObject)} produces a readable description rather than a memory address like \texttt{Dog@7852e922}. Return a clear string: \texttt{\detokenize{return "Dog[name=" + name + ", age=" + age + "]";}}}
\emph{(See Python Unit 7.1 for conceptual content --- OOP concepts are
language-agnostic.)}

Java is a \textbf{class-based} language --- every Java program is a
class. Understanding OOP is essential for Java.

\begin{center}\rule{0.5\linewidth}{0.5pt}\end{center}

\subsection{Lesson 7.2: Creating and Using a
Class}\label{lesson-7.2-creating-and-using-a-class-1}

\subsection{Class Syntax}\label{class-syntax-2}

\begin{Shaded}
\begin{Highlighting}[]
\KeywordTok{public} \KeywordTok{class} Dog \OperatorTok{\{}
    \CommentTok{// Instance variables (fields)}
    \KeywordTok{private} String name\OperatorTok{;}
    \KeywordTok{private} \CodeTypeTok{int} age\OperatorTok{;}
 
    \CommentTok{// Constructor}
    \KeywordTok{public} \CodeFuncTok{Dog}\OperatorTok{(}String name\OperatorTok{,} \CodeTypeTok{int} age\OperatorTok{)} \OperatorTok{\{}
        \KeywordTok{this}\OperatorTok{.}\CodeFuncTok{name} \OperatorTok{=} name\OperatorTok{;}
        \KeywordTok{this}\OperatorTok{.}\CodeFuncTok{age} \OperatorTok{=} age\OperatorTok{;}
    \OperatorTok{\}}
 
    \CommentTok{// Methods}
    \KeywordTok{public} \CodeTypeTok{void} \CodeFuncTok{sit}\OperatorTok{()} \OperatorTok{\{}
        System\OperatorTok{.}\CodeFuncTok{out}\OperatorTok{.}\CodeFuncTok{println}\OperatorTok{(}name \OperatorTok{+} \StringTok{" is now sitting."}\OperatorTok{);}
    \OperatorTok{\}}
 
    \KeywordTok{public} \CodeTypeTok{void} \CodeFuncTok{rollOver}\OperatorTok{()} \OperatorTok{\{}
        System\OperatorTok{.}\CodeFuncTok{out}\OperatorTok{.}\CodeFuncTok{println}\OperatorTok{(}name \OperatorTok{+} \StringTok{" rolled over!"}\OperatorTok{);}
    \OperatorTok{\}}
 
    \KeywordTok{public} \CodeTypeTok{void} \CodeFuncTok{getOlder}\OperatorTok{(}\CodeTypeTok{int} years\OperatorTok{)} \OperatorTok{\{}
        \KeywordTok{this}\OperatorTok{.}\CodeFuncTok{age} \OperatorTok{+=} years\OperatorTok{;}
        System\OperatorTok{.}\CodeFuncTok{out}\OperatorTok{.}\CodeFuncTok{println}\OperatorTok{(}name \OperatorTok{+} \StringTok{" is now "} \OperatorTok{+} age \OperatorTok{+} \StringTok{" years old."}\OperatorTok{);}
    \OperatorTok{\}}
 
    \CommentTok{// Getters and Setters}
    \KeywordTok{public} String \CodeFuncTok{getName}\OperatorTok{()} \OperatorTok{\{} \KeywordTok{return} name\OperatorTok{;} \OperatorTok{\}}
    \KeywordTok{public} \CodeTypeTok{int} \CodeFuncTok{getAge}\OperatorTok{()} \OperatorTok{\{} \KeywordTok{return} age\OperatorTok{;} \OperatorTok{\}}
    \KeywordTok{public} \CodeTypeTok{void} \CodeFuncTok{setAge}\OperatorTok{(}\CodeTypeTok{int} age\OperatorTok{)} \OperatorTok{\{} \KeywordTok{this}\OperatorTok{.}\CodeFuncTok{age} \OperatorTok{=} age\OperatorTok{;} \OperatorTok{\}}
 
    \CodeTypeTok{@Override}
    \KeywordTok{public} String \CodeFuncTok{toString}\OperatorTok{()} \OperatorTok{\{}
        \KeywordTok{return} \StringTok{"Dog\{name=\textquotesingle{}"} \OperatorTok{+} name \OperatorTok{+} \StringTok{"\textquotesingle{}, age="} \OperatorTok{+} age \OperatorTok{+} \StringTok{"\}"}\OperatorTok{;}
    \OperatorTok{\}}
\OperatorTok{\}}
\end{Highlighting}
\end{Shaded}

\subsection{Using the Class}\label{using-the-class}

\begin{Shaded}
\begin{Highlighting}[]
\KeywordTok{public} \KeywordTok{class} DogDemo \OperatorTok{\{}
    \KeywordTok{public} \CodeTypeTok{static} \CodeTypeTok{void} \CodeFuncTok{main}\OperatorTok{(}String\OperatorTok{[]} args\OperatorTok{)} \OperatorTok{\{}
        Dog myDog \OperatorTok{=} \KeywordTok{new} \CodeFuncTok{Dog}\OperatorTok{(}\StringTok{"Buddy"}\OperatorTok{,} \CodeNumTok{3}\OperatorTok{);}
        Dog yourDog \OperatorTok{=} \KeywordTok{new} \CodeFuncTok{Dog}\OperatorTok{(}\StringTok{"Lucy"}\OperatorTok{,} \CodeNumTok{1}\OperatorTok{);}
 
        System\OperatorTok{.}\CodeFuncTok{out}\OperatorTok{.}\CodeFuncTok{println}\OperatorTok{(}\StringTok{"My dog\textquotesingle{}s name is "} \OperatorTok{+} myDog\OperatorTok{.}\CodeFuncTok{getName}\OperatorTok{()} \OperatorTok{+} \StringTok{"."}\OperatorTok{);}
        myDog\OperatorTok{.}\CodeFuncTok{sit}\OperatorTok{();}
        myDog\OperatorTok{.}\CodeFuncTok{rollOver}\OperatorTok{();}
        myDog\OperatorTok{.}\CodeFuncTok{getOlder}\OperatorTok{(}\CodeNumTok{2}\OperatorTok{);}
 
        System\OperatorTok{.}\CodeFuncTok{out}\OperatorTok{.}\CodeFuncTok{println}\OperatorTok{(}myDog\OperatorTok{);}
    \OperatorTok{\}}
\OperatorTok{\}}
\end{Highlighting}
\end{Shaded}

\subsection{Access Modifiers}\label{access-modifiers}

\begin{fullwidth}{\def\LTcaptype{none} % do not increment counter
\begin{longtable}[]{@{}ll@{}}
\toprule\noalign{}
Modifier & Access Level \\
\midrule\noalign{}
\endhead
\bottomrule\noalign{}
\endlastfoot
\texttt{public} & Accessible from anywhere \\
\texttt{private} & Only within the same class \\
\texttt{protected} & Within the class and subclasses \\
(default) & Within the same package \\
\end{longtable}
}
\end{fullwidth}

\textbf{Best practice:} Make instance variables \texttt{private} and
provide \texttt{public} getter/setter methods --- this is
\textbf{encapsulation}.

\subsection{Skills Checklist}\label{skills-checklist-82}

\begin{itemize}
\tightlist
\item[$\square$]
  I can define a Java class with fields, a constructor, and methods.
\item[$\square$]
  I can create instances of a class using \texttt{new}.
\item[$\square$]
  I understand the purpose of \texttt{this} in Java.
\item[$\square$]
  I can write getter and setter methods.
\item[$\square$]
  I understand access modifiers (\texttt{public}, \texttt{private}).
\end{itemize}

\begin{center}\rule{0.5\linewidth}{0.5pt}\end{center}

\subsection{Lesson 7.3: Inheritance}\label{lesson-7.3-inheritance}

\begin{Shaded}
\begin{Highlighting}[]
\KeywordTok{public} \KeywordTok{class} Vehicle \OperatorTok{\{}
    \KeywordTok{protected} String make\OperatorTok{;}
    \KeywordTok{protected} String model\OperatorTok{;}
    \KeywordTok{protected} \CodeTypeTok{int} year\OperatorTok{;}
 
    \KeywordTok{public} \CodeFuncTok{Vehicle}\OperatorTok{(}String make\OperatorTok{,} String model\OperatorTok{,} \CodeTypeTok{int} year\OperatorTok{)} \OperatorTok{\{}
        \KeywordTok{this}\OperatorTok{.}\CodeFuncTok{make} \OperatorTok{=} make\OperatorTok{;}
        \KeywordTok{this}\OperatorTok{.}\CodeFuncTok{model} \OperatorTok{=} model\OperatorTok{;}
        \KeywordTok{this}\OperatorTok{.}\CodeFuncTok{year} \OperatorTok{=} year\OperatorTok{;}
    \OperatorTok{\}}
 
    \KeywordTok{public} String \CodeFuncTok{getDescriptiveName}\OperatorTok{()} \OperatorTok{\{}
        \KeywordTok{return} year \OperatorTok{+} \StringTok{" "} \OperatorTok{+} make \OperatorTok{+} \StringTok{" "} \OperatorTok{+} model\OperatorTok{;}
    \OperatorTok{\}}
 
    \KeywordTok{public} \CodeTypeTok{void} \CodeFuncTok{readOdometer}\OperatorTok{()} \OperatorTok{\{}
        System\OperatorTok{.}\CodeFuncTok{out}\OperatorTok{.}\CodeFuncTok{println}\OperatorTok{(}\StringTok{"No odometer reading available."}\OperatorTok{);}
    \OperatorTok{\}}
\OperatorTok{\}}
 
\KeywordTok{public} \KeywordTok{class} Car \KeywordTok{extends} Vehicle \OperatorTok{\{}
    \KeywordTok{private} \CodeTypeTok{int} odometer\OperatorTok{;}
 
    \KeywordTok{public} \CodeFuncTok{Car}\OperatorTok{(}String make\OperatorTok{,} String model\OperatorTok{,} \CodeTypeTok{int} year\OperatorTok{)} \OperatorTok{\{}
        \KeywordTok{super}\OperatorTok{(}make\OperatorTok{,} model\OperatorTok{,} year\OperatorTok{);}   \CommentTok{// Call parent constructor}
        \KeywordTok{this}\OperatorTok{.}\CodeFuncTok{odometer} \OperatorTok{=} \CodeNumTok{0}\OperatorTok{;}
    \OperatorTok{\}}
 
    \CodeTypeTok{@Override}
    \KeywordTok{public} \CodeTypeTok{void} \CodeFuncTok{readOdometer}\OperatorTok{()} \OperatorTok{\{}
        System\OperatorTok{.}\CodeFuncTok{out}\OperatorTok{.}\CodeFuncTok{println}\OperatorTok{(}\StringTok{"This car has "} \OperatorTok{+} odometer \OperatorTok{+} \StringTok{" miles on it."}\OperatorTok{);}
    \OperatorTok{\}}
 
    \KeywordTok{public} \CodeTypeTok{void} \CodeFuncTok{updateOdometer}\OperatorTok{(}\CodeTypeTok{int} miles\OperatorTok{)} \OperatorTok{\{}
        \KeywordTok{if} \OperatorTok{(}miles \OperatorTok{\textgreater{}=} odometer\OperatorTok{)} \OperatorTok{\{}
            odometer \OperatorTok{=} miles\OperatorTok{;}
        \OperatorTok{\}} \KeywordTok{else} \OperatorTok{\{}
            System\OperatorTok{.}\CodeFuncTok{out}\OperatorTok{.}\CodeFuncTok{println}\OperatorTok{(}\StringTok{"Error: Cannot roll back odometer."}\OperatorTok{);}
        \OperatorTok{\}}
    \OperatorTok{\}}
\OperatorTok{\}}
 
\KeywordTok{public} \KeywordTok{class} ElectricCar \KeywordTok{extends} Car \OperatorTok{\{}
    \KeywordTok{private} \CodeTypeTok{int} batterySize\OperatorTok{;}
 
    \KeywordTok{public} \CodeFuncTok{ElectricCar}\OperatorTok{(}String make\OperatorTok{,} String model\OperatorTok{,} \CodeTypeTok{int} year\OperatorTok{)} \OperatorTok{\{}
        \KeywordTok{super}\OperatorTok{(}make\OperatorTok{,} model\OperatorTok{,} year\OperatorTok{);}
        \KeywordTok{this}\OperatorTok{.}\CodeFuncTok{batterySize} \OperatorTok{=} \CodeNumTok{75}\OperatorTok{;}
    \OperatorTok{\}}
 
    \KeywordTok{public} \CodeTypeTok{void} \CodeFuncTok{describeBattery}\OperatorTok{()} \OperatorTok{\{}
        System\OperatorTok{.}\CodeFuncTok{out}\OperatorTok{.}\CodeFuncTok{println}\OperatorTok{(}\StringTok{"This car has a "} \OperatorTok{+} batterySize \OperatorTok{+} \StringTok{"{-}kWh battery."}\OperatorTok{);}
    \OperatorTok{\}}
\OperatorTok{\}}
\end{Highlighting}
\end{Shaded}

\subsection{Abstract Classes and
Interfaces}\label{abstract-classes-and-interfaces}

\begin{Shaded}
\begin{Highlighting}[]
\CommentTok{// Abstract class {-}{-}{-} cannot be instantiated}
\KeywordTok{public} \KeywordTok{abstract} \KeywordTok{class} Shape \OperatorTok{\{}
    \KeywordTok{public} \KeywordTok{abstract} \CodeTypeTok{double} \CodeFuncTok{area}\OperatorTok{();}      \CommentTok{// Abstract method {-}{-}{-} no body}
    \KeywordTok{public} \KeywordTok{abstract} \CodeTypeTok{double} \CodeFuncTok{perimeter}\OperatorTok{();} \CommentTok{// Must be implemented by subclasses}
 
    \KeywordTok{public} \CodeTypeTok{void} \CodeFuncTok{printInfo}\OperatorTok{()} \OperatorTok{\{}           \CommentTok{// Concrete method {-}{-}{-} has body}
        System\OperatorTok{.}\CodeFuncTok{out}\OperatorTok{.}\CodeFuncTok{printf}\OperatorTok{(}\StringTok{"Area: }\StringTok{\%.2f}\StringTok{, Perimeter: }\StringTok{\%.2f\%n}\StringTok{"}\OperatorTok{,} \CodeFuncTok{area}\OperatorTok{(),} \CodeFuncTok{perimeter}\OperatorTok{());}
    \OperatorTok{\}}
\OperatorTok{\}}
 
\KeywordTok{public} \KeywordTok{class} Circle \KeywordTok{extends} Shape \OperatorTok{\{}
    \KeywordTok{private} \CodeTypeTok{double} radius\OperatorTok{;}
 
    \KeywordTok{public} \CodeFuncTok{Circle}\OperatorTok{(}\CodeTypeTok{double} radius\OperatorTok{)} \OperatorTok{\{} \KeywordTok{this}\OperatorTok{.}\CodeFuncTok{radius} \OperatorTok{=} radius\OperatorTok{;} \OperatorTok{\}}
 
    \CodeTypeTok{@Override}
    \KeywordTok{public} \CodeTypeTok{double} \CodeFuncTok{area}\OperatorTok{()} \OperatorTok{\{} \KeywordTok{return} Math\OperatorTok{.}\CodeFuncTok{PI} \OperatorTok{*} radius \OperatorTok{*} radius\OperatorTok{;} \OperatorTok{\}}
 
    \CodeTypeTok{@Override}
    \KeywordTok{public} \CodeTypeTok{double} \CodeFuncTok{perimeter}\OperatorTok{()} \OperatorTok{\{} \KeywordTok{return} \CodeNumTok{2} \OperatorTok{*} Math\OperatorTok{.}\CodeFuncTok{PI} \OperatorTok{*} radius\OperatorTok{;} \OperatorTok{\}}
\OperatorTok{\}}
 
\CommentTok{// Interface {-}{-}{-} defines a contract}
\KeywordTok{public} \KeywordTok{interface} Drawable \OperatorTok{\{}
    \CodeTypeTok{void} \CodeFuncTok{draw}\OperatorTok{();}
    \KeywordTok{default} \CodeTypeTok{void} \CodeFuncTok{resize}\OperatorTok{(}\CodeTypeTok{double} factor\OperatorTok{)} \OperatorTok{\{}
        System\OperatorTok{.}\CodeFuncTok{out}\OperatorTok{.}\CodeFuncTok{println}\OperatorTok{(}\StringTok{"Resizing by factor "} \OperatorTok{+} factor\OperatorTok{);}
    \OperatorTok{\}}
\OperatorTok{\}}
\end{Highlighting}
\end{Shaded}

\subsection{Skills Checklist}\label{skills-checklist-83}

\begin{itemize}
\tightlist
\item[$\square$]
  I can create a subclass that extends a parent class.
\item[$\square$]
  I can use \texttt{super()} to call the parent constructor.
\item[$\square$]
  I can override methods using \texttt{@Override}.
\item[$\square$]
  I understand abstract classes and interfaces.
\end{itemize}

\begin{center}\rule{0.5\linewidth}{0.5pt}\end{center}

\section{Unit 8: Files and Exception
Handling}\label{unit-8-files-and-exception-handling-1}

\subsection{Overview}\label{overview-22}

This unit covers reading and writing files in Java and handling
exceptions.

\begin{center}\rule{0.5\linewidth}{0.5pt}\end{center}

\subsection{Lesson 8.1: Reading from
Files}\amn{\textbf{Accessibility Tip:} Always catch specific exception types: \texttt{catch (IOException e)} not bare \texttt{catch (Exception e)}. Print the message in every catch block: \texttt{System.out.println("Error: " + e.getMessage());}. A silent catch block hides bugs that are impossible to debug without visual scanning.}
\subsection{Lesson 8.1: Reading from
Files}\label{lesson-8.1-reading-from-files}

\begin{Shaded}
\begin{Highlighting}[]
\KeywordTok{import} \CodeAccentTok{java}\OperatorTok{.}\CodeAccentTok{io}\OperatorTok{.*;}
\KeywordTok{import} \CodeAccentTok{java}\OperatorTok{.}\CodeAccentTok{util}\OperatorTok{.*;}
 
\KeywordTok{public} \KeywordTok{class} FileReaderDemo \OperatorTok{\{}
    \KeywordTok{public} \CodeTypeTok{static} \CodeTypeTok{void} \CodeFuncTok{main}\OperatorTok{(}String\OperatorTok{[]} args\OperatorTok{)} \OperatorTok{\{}
        \CommentTok{// Reading entire file using Scanner}
        \KeywordTok{try} \OperatorTok{(}Scanner fileScanner \OperatorTok{=} \KeywordTok{new} Scanner\OperatorTok{(}\KeywordTok{new} File\OperatorTok{(}\StringTok{"my\_notes.txt"}\OperatorTok{)))} \OperatorTok{\{}
            \KeywordTok{while} \OperatorTok{(}fileScanner\OperatorTok{.}\CodeFuncTok{hasNextLine}\OperatorTok{())} \OperatorTok{\{}
                String line \OperatorTok{=} fileScanner\OperatorTok{.}\CodeFuncTok{nextLine}\OperatorTok{();}
                System\OperatorTok{.}\CodeFuncTok{out}\OperatorTok{.}\CodeFuncTok{println}\OperatorTok{(}line\OperatorTok{);}
            \OperatorTok{\}}
        \OperatorTok{\}} \KeywordTok{catch} \OperatorTok{(}FileNotFoundException e\OperatorTok{)} \OperatorTok{\{}
            System\OperatorTok{.}\CodeFuncTok{out}\OperatorTok{.}\CodeFuncTok{println}\OperatorTok{(}\StringTok{"Error: File not found {-} "} \OperatorTok{+} e\OperatorTok{.}\CodeFuncTok{getMessage}\OperatorTok{());}
        \OperatorTok{\}}
 
        \CommentTok{// Reading with BufferedReader (more efficient for large files)}
        \KeywordTok{try} \OperatorTok{(}BufferedReader reader \OperatorTok{=} \KeywordTok{new} BufferedReader\OperatorTok{(}\KeywordTok{new} FileReader\OperatorTok{(}\StringTok{"my\_notes.txt"}\OperatorTok{)))} \OperatorTok{\{}
            String line\OperatorTok{;}
            \KeywordTok{while} \OperatorTok{((}line \OperatorTok{=} reader\OperatorTok{.}\CodeFuncTok{readLine}\OperatorTok{())} \OperatorTok{!=} \KeywordTok{null}\OperatorTok{)} \OperatorTok{\{}
                System\OperatorTok{.}\CodeFuncTok{out}\OperatorTok{.}\CodeFuncTok{println}\OperatorTok{(}line\OperatorTok{);}
            \OperatorTok{\}}
        \OperatorTok{\}} \KeywordTok{catch} \OperatorTok{(}IOException e\OperatorTok{)} \OperatorTok{\{}
            System\OperatorTok{.}\CodeFuncTok{out}\OperatorTok{.}\CodeFuncTok{println}\OperatorTok{(}\StringTok{"Error reading file: "} \OperatorTok{+} e\OperatorTok{.}\CodeFuncTok{getMessage}\OperatorTok{());}
        \OperatorTok{\}}
 
        \CommentTok{// Reading all lines into a List}
        \KeywordTok{try} \OperatorTok{\{}
            List\OperatorTok{\textless{}}String\OperatorTok{\textgreater{}} lines \OperatorTok{=} \KeywordTok{new} java\OperatorTok{.}\CodeFuncTok{nio}\OperatorTok{.}\CodeFuncTok{file}\OperatorTok{.}\CodeFuncTok{Files}\OperatorTok{.}\CodeFuncTok{readAllLines}\OperatorTok{(}
                java\OperatorTok{.}\CodeFuncTok{nio}\OperatorTok{.}\CodeFuncTok{file}\OperatorTok{.}\CodeFuncTok{Paths}\OperatorTok{.}\CodeFuncTok{get}\OperatorTok{(}\StringTok{"my\_notes.txt"}\OperatorTok{)}
            \OperatorTok{);}
            \KeywordTok{for} \OperatorTok{(}String line \OperatorTok{:} lines\OperatorTok{)} \OperatorTok{\{}
                System\OperatorTok{.}\CodeFuncTok{out}\OperatorTok{.}\CodeFuncTok{println}\OperatorTok{(}line\OperatorTok{);}
            \OperatorTok{\}}
        \OperatorTok{\}} \KeywordTok{catch} \OperatorTok{(}IOException e\OperatorTok{)} \OperatorTok{\{}
            System\OperatorTok{.}\CodeFuncTok{out}\OperatorTok{.}\CodeFuncTok{println}\OperatorTok{(}\StringTok{"Error: "} \OperatorTok{+} e\OperatorTok{.}\CodeFuncTok{getMessage}\OperatorTok{());}
        \OperatorTok{\}}
    \OperatorTok{\}}
\OperatorTok{\}}
\end{Highlighting}
\end{Shaded}

\subsection{\texorpdfstring{Simpler approach using
\texttt{java.nio.file}}{Simpler approach using java.nio.file}}\label{simpler-approach-using-java.nio.file}

\begin{Shaded}
\begin{Highlighting}[]
\KeywordTok{import} \CodeAccentTok{java}\OperatorTok{.}\CodeAccentTok{nio}\OperatorTok{.}\CodeAccentTok{file}\OperatorTok{.*;}
\KeywordTok{import} \CodeAccentTok{java}\OperatorTok{.}\CodeAccentTok{io}\OperatorTok{.}\CodeAccentTok{IOException}\OperatorTok{;}
\KeywordTok{import} \CodeAccentTok{java}\OperatorTok{.}\CodeAccentTok{util}\OperatorTok{.}\CodeAccentTok{List}\OperatorTok{;}
 
\KeywordTok{public} \KeywordTok{class} ModernFileReader \OperatorTok{\{}
    \KeywordTok{public} \CodeTypeTok{static} \CodeTypeTok{void} \CodeFuncTok{main}\OperatorTok{(}String\OperatorTok{[]} args\OperatorTok{)} \OperatorTok{\{}
        Path filePath \OperatorTok{=} Paths\OperatorTok{.}\CodeFuncTok{get}\OperatorTok{(}\StringTok{"my\_notes.txt"}\OperatorTok{);}
 
        \KeywordTok{try} \OperatorTok{\{}
            \CommentTok{// Read all lines}
            List\OperatorTok{\textless{}}String\OperatorTok{\textgreater{}} lines \OperatorTok{=} Files\OperatorTok{.}\CodeFuncTok{readAllLines}\OperatorTok{(}filePath\OperatorTok{);}
            \KeywordTok{for} \OperatorTok{(}String line \OperatorTok{:} lines\OperatorTok{)} \OperatorTok{\{}
                System\OperatorTok{.}\CodeFuncTok{out}\OperatorTok{.}\CodeFuncTok{println}\OperatorTok{(}line\OperatorTok{);}
            \OperatorTok{\}}
 
            \CommentTok{// Read entire file as one string}
            String content \OperatorTok{=} Files\OperatorTok{.}\CodeFuncTok{readString}\OperatorTok{(}filePath\OperatorTok{);}
            System\OperatorTok{.}\CodeFuncTok{out}\OperatorTok{.}\CodeFuncTok{println}\OperatorTok{(}content\OperatorTok{);}
 
        \OperatorTok{\}} \KeywordTok{catch} \OperatorTok{(}IOException e\OperatorTok{)} \OperatorTok{\{}
            System\OperatorTok{.}\CodeFuncTok{out}\OperatorTok{.}\CodeFuncTok{println}\OperatorTok{(}\StringTok{"Error reading file: "} \OperatorTok{+} e\OperatorTok{.}\CodeFuncTok{getMessage}\OperatorTok{());}
        \OperatorTok{\}}
    \OperatorTok{\}}
\OperatorTok{\}}
\end{Highlighting}
\end{Shaded}

\begin{center}\rule{0.5\linewidth}{0.5pt}\end{center}

\subsection{Lesson 8.2: Writing to
Files}\label{lesson-8.2-writing-to-files}

\begin{Shaded}
\begin{Highlighting}[]
\KeywordTok{import} \CodeAccentTok{java}\OperatorTok{.}\CodeAccentTok{io}\OperatorTok{.*;}
\KeywordTok{import} \CodeAccentTok{java}\OperatorTok{.}\CodeAccentTok{nio}\OperatorTok{.}\CodeAccentTok{file}\OperatorTok{.*;}
 
\KeywordTok{public} \KeywordTok{class} FileWriterDemo \OperatorTok{\{}
    \KeywordTok{public} \CodeTypeTok{static} \CodeTypeTok{void} \CodeFuncTok{main}\OperatorTok{(}String\OperatorTok{[]} args\OperatorTok{)} \OperatorTok{\{}
        \CommentTok{// Writing with PrintWriter (creates new or overwrites)}
        \KeywordTok{try} \OperatorTok{(}PrintWriter writer \OperatorTok{=} \KeywordTok{new} PrintWriter\OperatorTok{(}\StringTok{"output.txt"}\OperatorTok{))} \OperatorTok{\{}
            writer\OperatorTok{.}\CodeFuncTok{println}\OperatorTok{(}\StringTok{"Hello, Java file I/O!"}\OperatorTok{);}
            writer\OperatorTok{.}\CodeFuncTok{println}\OperatorTok{(}\StringTok{"Writing is easy."}\OperatorTok{);}
            writer\OperatorTok{.}\CodeFuncTok{printf}\OperatorTok{(}\StringTok{"Value: }\StringTok{\%d\%n}\StringTok{"}\OperatorTok{,} \CodeNumTok{42}\OperatorTok{);}
        \OperatorTok{\}} \KeywordTok{catch} \OperatorTok{(}FileNotFoundException e\OperatorTok{)} \OperatorTok{\{}
            System\OperatorTok{.}\CodeFuncTok{out}\OperatorTok{.}\CodeFuncTok{println}\OperatorTok{(}\StringTok{"Error: Could not create file."}\OperatorTok{);}
        \OperatorTok{\}}
 
        \CommentTok{// Appending with FileWriter(true = append mode)}
        \KeywordTok{try} \OperatorTok{(}PrintWriter appender \OperatorTok{=} \KeywordTok{new} PrintWriter\OperatorTok{(}\KeywordTok{new} FileWriter\OperatorTok{(}\StringTok{"output.txt"}\OperatorTok{,} \KeywordTok{true}\OperatorTok{)))} \OperatorTok{\{}
            appender\OperatorTok{.}\CodeFuncTok{println}\OperatorTok{(}\StringTok{"This line is appended."}\OperatorTok{);}
        \OperatorTok{\}} \KeywordTok{catch} \OperatorTok{(}IOException e\OperatorTok{)} \OperatorTok{\{}
            System\OperatorTok{.}\CodeFuncTok{out}\OperatorTok{.}\CodeFuncTok{println}\OperatorTok{(}\StringTok{"Error appending to file."}\OperatorTok{);}
        \OperatorTok{\}}
 
        \CommentTok{// Modern NIO approach}
        \KeywordTok{try} \OperatorTok{\{}
            String content \OperatorTok{=} \StringTok{"Hello from NIO!}\StringTok{\textbackslash{}n}\StringTok{Second line."}\OperatorTok{;}
            Files\OperatorTok{.}\CodeFuncTok{writeString}\OperatorTok{(}Paths\OperatorTok{.}\CodeFuncTok{get}\OperatorTok{(}\StringTok{"modern\_output.txt"}\OperatorTok{),} content\OperatorTok{);}
 
            \CommentTok{// Append}
            Files\OperatorTok{.}\CodeFuncTok{writeString}\OperatorTok{(}Paths\OperatorTok{.}\CodeFuncTok{get}\OperatorTok{(}\StringTok{"modern\_output.txt"}\OperatorTok{),}
                \StringTok{"}\StringTok{\textbackslash{}n}\StringTok{Appended line."}\OperatorTok{,}
                StandardOpenOption\OperatorTok{.}\CodeFuncTok{APPEND}\OperatorTok{);}
        \OperatorTok{\}} \KeywordTok{catch} \OperatorTok{(}IOException e\OperatorTok{)} \OperatorTok{\{}
            System\OperatorTok{.}\CodeFuncTok{out}\OperatorTok{.}\CodeFuncTok{println}\OperatorTok{(}\StringTok{"Error: "} \OperatorTok{+} e\OperatorTok{.}\CodeFuncTok{getMessage}\OperatorTok{());}
        \OperatorTok{\}}
    \OperatorTok{\}}
\OperatorTok{\}}
\end{Highlighting}
\end{Shaded}

\begin{center}\rule{0.5\linewidth}{0.5pt}\end{center}

\subsection{Lesson 8.3: Exception
Handling}\label{lesson-8.3-exception-handling}

\subsection{\texorpdfstring{\texttt{try-catch-finally}}{try-catch-finally}}\label{try-catch-finally}

\begin{Shaded}
\begin{Highlighting}[]
\KeywordTok{public} \KeywordTok{class} ExceptionHandling \OperatorTok{\{}
    \KeywordTok{public} \CodeTypeTok{static} \CodeTypeTok{void} \CodeFuncTok{main}\OperatorTok{(}String\OperatorTok{[]} args\OperatorTok{)} \OperatorTok{\{}
        \CommentTok{// Basic try{-}catch}
        \KeywordTok{try} \OperatorTok{\{}
            \CodeTypeTok{int} result \OperatorTok{=} \CodeNumTok{10} \OperatorTok{/} \CodeNumTok{0}\OperatorTok{;}
        \OperatorTok{\}} \KeywordTok{catch} \OperatorTok{(}ArithmeticException e\OperatorTok{)} \OperatorTok{\{}
            System\OperatorTok{.}\CodeFuncTok{out}\OperatorTok{.}\CodeFuncTok{println}\OperatorTok{(}\StringTok{"Error: "} \OperatorTok{+} e\OperatorTok{.}\CodeFuncTok{getMessage}\OperatorTok{());}
        \OperatorTok{\}}
 
        \CommentTok{// Multiple catch blocks}
        \KeywordTok{try} \OperatorTok{\{}
            String input \OperatorTok{=} \StringTok{"abc"}\OperatorTok{;}
            \CodeTypeTok{int} number \OperatorTok{=} Integer\OperatorTok{.}\CodeFuncTok{parseInt}\OperatorTok{(}input\OperatorTok{);}
        \OperatorTok{\}} \KeywordTok{catch} \OperatorTok{(}NumberFormatException e\OperatorTok{)} \OperatorTok{\{}
            System\OperatorTok{.}\CodeFuncTok{out}\OperatorTok{.}\CodeFuncTok{println}\OperatorTok{(}\StringTok{"Invalid number format: "} \OperatorTok{+} e\OperatorTok{.}\CodeFuncTok{getMessage}\OperatorTok{());}
        \OperatorTok{\}} \KeywordTok{catch} \OperatorTok{(}Exception e\OperatorTok{)} \OperatorTok{\{}
            System\OperatorTok{.}\CodeFuncTok{out}\OperatorTok{.}\CodeFuncTok{println}\OperatorTok{(}\StringTok{"Unexpected error: "} \OperatorTok{+} e\OperatorTok{.}\CodeFuncTok{getMessage}\OperatorTok{());}
        \OperatorTok{\}}
 
        \CommentTok{// try{-}catch{-}finally}
        PrintWriter writer \OperatorTok{=} \KeywordTok{null}\OperatorTok{;}
        \KeywordTok{try} \OperatorTok{\{}
            writer \OperatorTok{=} \KeywordTok{new} PrintWriter\OperatorTok{(}\StringTok{"data.txt"}\OperatorTok{);}
            writer\OperatorTok{.}\CodeFuncTok{println}\OperatorTok{(}\StringTok{"Some data"}\OperatorTok{);}
        \OperatorTok{\}} \KeywordTok{catch} \OperatorTok{(}FileNotFoundException e\OperatorTok{)} \OperatorTok{\{}
            System\OperatorTok{.}\CodeFuncTok{out}\OperatorTok{.}\CodeFuncTok{println}\OperatorTok{(}\StringTok{"Could not create file."}\OperatorTok{);}
        \OperatorTok{\}} \KeywordTok{finally} \OperatorTok{\{}
            \KeywordTok{if} \OperatorTok{(}writer \OperatorTok{!=} \KeywordTok{null}\OperatorTok{)} \OperatorTok{\{}
                writer\OperatorTok{.}\CodeFuncTok{close}\OperatorTok{();}
                System\OperatorTok{.}\CodeFuncTok{out}\OperatorTok{.}\CodeFuncTok{println}\OperatorTok{(}\StringTok{"File closed."}\OperatorTok{);}
            \OperatorTok{\}}
        \OperatorTok{\}}
 
        \CommentTok{// try{-}with{-}resources (automatically closes resources)}
        \KeywordTok{try} \OperatorTok{(}PrintWriter pw \OperatorTok{=} \KeywordTok{new} PrintWriter\OperatorTok{(}\StringTok{"data2.txt"}\OperatorTok{))} \OperatorTok{\{}
            pw\OperatorTok{.}\CodeFuncTok{println}\OperatorTok{(}\StringTok{"Better approach!"}\OperatorTok{);}
        \OperatorTok{\}} \KeywordTok{catch} \OperatorTok{(}FileNotFoundException e\OperatorTok{)} \OperatorTok{\{}
            System\OperatorTok{.}\CodeFuncTok{out}\OperatorTok{.}\CodeFuncTok{println}\OperatorTok{(}\StringTok{"Could not create file."}\OperatorTok{);}
        \OperatorTok{\}}
    \OperatorTok{\}}
\OperatorTok{\}}
\end{Highlighting}
\end{Shaded}

\subsection{Common Java Exceptions}\label{common-java-exceptions}

\begin{fullwidth}{\def\LTcaptype{none} % do not increment counter
\begin{longtable}[]{@{}
  >{\raggedright\arraybackslash}p{(\linewidth - 2\tabcolsep) * \real{0.3516}}
  >{\raggedright\arraybackslash}p{(\linewidth - 2\tabcolsep) * \real{0.6484}}@{}}
\toprule\noalign{}
\begin{minipage}[b]{\linewidth}\raggedright
Exception
\end{minipage} & \begin{minipage}[b]{\linewidth}\raggedright
When it occurs
\end{minipage} \\
\midrule\noalign{}
\endhead
\bottomrule\noalign{}
\endlastfoot
\texttt{ArithmeticException} & Division by zero \\
\texttt{NumberFormatException} & Invalid number conversion (e.g.,
\texttt{Integer.parseInt("abc")}) \\
\texttt{ArrayIndexOutOfBoundsException} & Array index out of range \\
\texttt{NullPointerException} & Calling method on null reference \\
\texttt{FileNotFoundException} & File doesn't exist \\
\texttt{IOException} & General I/O error \\
\texttt{ClassCastException} & Invalid type cast \\
\end{longtable}
}
\end{fullwidth}

\subsection{Custom Exceptions}\label{custom-exceptions}

\begin{Shaded}
\begin{Highlighting}[]
\KeywordTok{public} \KeywordTok{class} InsufficientFundsException \KeywordTok{extends} Exception \OperatorTok{\{}
    \KeywordTok{private} \CodeTypeTok{double} amount\OperatorTok{;}
 
    \KeywordTok{public} \CodeFuncTok{InsufficientFundsException}\OperatorTok{(}\CodeTypeTok{double} amount\OperatorTok{)} \OperatorTok{\{}
        \KeywordTok{super}\OperatorTok{(}\StringTok{"Insufficient funds: need "} \OperatorTok{+} amount \OperatorTok{+} \StringTok{" more."}\OperatorTok{);}
        \KeywordTok{this}\OperatorTok{.}\CodeFuncTok{amount} \OperatorTok{=} amount\OperatorTok{;}
    \OperatorTok{\}}
 
    \KeywordTok{public} \CodeTypeTok{double} \CodeFuncTok{getAmount}\OperatorTok{()} \OperatorTok{\{} \KeywordTok{return} amount\OperatorTok{;} \OperatorTok{\}}
\OperatorTok{\}}
 
\CommentTok{// Usage:}
\KeywordTok{public} \CodeTypeTok{void} \CodeFuncTok{withdraw}\OperatorTok{(}\CodeTypeTok{double} amount\OperatorTok{)} \KeywordTok{throws} InsufficientFundsException \OperatorTok{\{}
    \KeywordTok{if} \OperatorTok{(}amount \OperatorTok{\textgreater{}} balance\OperatorTok{)} \OperatorTok{\{}
        \KeywordTok{throw} \KeywordTok{new} \CodeFuncTok{InsufficientFundsException}\OperatorTok{(}amount \OperatorTok{{-}} balance\OperatorTok{);}
    \OperatorTok{\}}
    balance \OperatorTok{{-}=} amount\OperatorTok{;}
\OperatorTok{\}}
\end{Highlighting}
\end{Shaded}

\subsection{Skills Checklist}\label{skills-checklist-84}

\begin{itemize}
\tightlist
\item[$\square$]
  I can read from a file using \texttt{Scanner} or
  \texttt{BufferedReader}.
\item[$\square$]
  I can write to a file using \texttt{PrintWriter}.
\item[$\square$]
  I can use \texttt{try-catch-finally} to handle exceptions.
\item[$\square$]
  I can use try-with-resources for automatic resource cleanup.
\item[$\square$]
  I know common Java exception types.
\end{itemize}

\newpage

\section{Unit 9: Collections, Data Analysis, and
Testing}\label{unit-9-collections-data-analysis-and-testing}

\subsection{Overview}\label{overview-23}

This unit introduces Java's Collections Framework, basic data analysis,
and unit testing with JUnit.

\begin{center}\rule{0.5\linewidth}{0.5pt}\end{center}

\subsection{Lesson 9.1: Java Collections
Framework}\amn{\textbf{Accessibility Tip:} Run JUnit tests from the terminal with \texttt{mvn test} or \texttt{./gradlew test}. Use the \texttt{-{}-verbose} flag where available to get one result line per test. Each line names the test and its result---far more navigable than summary dots when using a screen reader.}
\subsection{Lesson 9.1: Java Collections
Framework}\label{lesson-9.1-java-collections-framework}

\subsection{Key Interfaces and
Classes}\label{key-interfaces-and-classes}

\begin{verbatim}
Collection
|--- List
|   |--- ArrayList
|   +--- LinkedList
+--- Set
    |--- HashSet
    +--- TreeSet

Map
|--- HashMap
+--- TreeMap
\end{verbatim}

\subsection{\texorpdfstring{\texttt{HashMap} --- Key-Value
Storage}{HashMap --- Key-Value Storage}}\label{hashmap-key-value-storage}

\begin{Shaded}
\begin{Highlighting}[]
\KeywordTok{import} \CodeAccentTok{java}\OperatorTok{.}\CodeAccentTok{util}\OperatorTok{.}\CodeAccentTok{HashMap}\OperatorTok{;}
\KeywordTok{import} \CodeAccentTok{java}\OperatorTok{.}\CodeAccentTok{util}\OperatorTok{.}\CodeAccentTok{Map}\OperatorTok{;}
 
HashMap\OperatorTok{\textless{}}String\OperatorTok{,} Integer\OperatorTok{\textgreater{}} inventory \OperatorTok{=} \KeywordTok{new} HashMap\OperatorTok{\textless{}\textgreater{}();}
 
\CommentTok{// Add entries}
inventory\OperatorTok{.}\CodeFuncTok{put}\OperatorTok{(}\StringTok{"apples"}\OperatorTok{,} \CodeNumTok{50}\OperatorTok{);}
inventory\OperatorTok{.}\CodeFuncTok{put}\OperatorTok{(}\StringTok{"bananas"}\OperatorTok{,} \CodeNumTok{30}\OperatorTok{);}
inventory\OperatorTok{.}\CodeFuncTok{put}\OperatorTok{(}\StringTok{"oranges"}\OperatorTok{,} \CodeNumTok{40}\OperatorTok{);}
 
\CommentTok{// Access}
System\OperatorTok{.}\CodeFuncTok{out}\OperatorTok{.}\CodeFuncTok{println}\OperatorTok{(}inventory\OperatorTok{.}\CodeFuncTok{get}\OperatorTok{(}\StringTok{"apples"}\OperatorTok{));}           \CommentTok{// 50}
System\OperatorTok{.}\CodeFuncTok{out}\OperatorTok{.}\CodeFuncTok{println}\OperatorTok{(}inventory\OperatorTok{.}\CodeFuncTok{getOrDefault}\OperatorTok{(}\StringTok{"grapes"}\OperatorTok{,} \CodeNumTok{0}\OperatorTok{));} \CommentTok{// 0}
 
\CommentTok{// Check}
System\OperatorTok{.}\CodeFuncTok{out}\OperatorTok{.}\CodeFuncTok{println}\OperatorTok{(}inventory\OperatorTok{.}\CodeFuncTok{containsKey}\OperatorTok{(}\StringTok{"bananas"}\OperatorTok{));}   \CommentTok{// true}
System\OperatorTok{.}\CodeFuncTok{out}\OperatorTok{.}\CodeFuncTok{println}\OperatorTok{(}inventory\OperatorTok{.}\CodeFuncTok{containsValue}\OperatorTok{(}\CodeNumTok{30}\OperatorTok{));}         \CommentTok{// true}
 
\CommentTok{// Update}
inventory\OperatorTok{.}\CodeFuncTok{put}\OperatorTok{(}\StringTok{"apples"}\OperatorTok{,} \CodeNumTok{55}\OperatorTok{);}
 
\CommentTok{// Remove}
inventory\OperatorTok{.}\CodeFuncTok{remove}\OperatorTok{(}\StringTok{"bananas"}\OperatorTok{);}
 
\CommentTok{// Iterate}
\KeywordTok{for} \OperatorTok{(}Map\OperatorTok{.}\CodeFuncTok{Entry}\OperatorTok{\textless{}}String\OperatorTok{,} Integer\OperatorTok{\textgreater{}} entry \OperatorTok{:} inventory\OperatorTok{.}\CodeFuncTok{entrySet}\OperatorTok{())} \OperatorTok{\{}
    System\OperatorTok{.}\CodeFuncTok{out}\OperatorTok{.}\CodeFuncTok{println}\OperatorTok{(}entry\OperatorTok{.}\CodeFuncTok{getKey}\OperatorTok{()} \OperatorTok{+} \StringTok{": "} \OperatorTok{+} entry\OperatorTok{.}\CodeFuncTok{getValue}\OperatorTok{());}
\OperatorTok{\}}
 
\CommentTok{// Keys only}
\KeywordTok{for} \OperatorTok{(}String key \OperatorTok{:} inventory\OperatorTok{.}\CodeFuncTok{keySet}\OperatorTok{())} \OperatorTok{\{}
    System\OperatorTok{.}\CodeFuncTok{out}\OperatorTok{.}\CodeFuncTok{println}\OperatorTok{(}key\OperatorTok{);}
\OperatorTok{\}}
 
\CommentTok{// Values only}
\KeywordTok{for} \OperatorTok{(}\CodeTypeTok{int} value \OperatorTok{:} inventory\OperatorTok{.}\CodeFuncTok{values}\OperatorTok{())} \OperatorTok{\{}
    System\OperatorTok{.}\CodeFuncTok{out}\OperatorTok{.}\CodeFuncTok{println}\OperatorTok{(}value\OperatorTok{);}
\OperatorTok{\}}
 
System\OperatorTok{.}\CodeFuncTok{out}\OperatorTok{.}\CodeFuncTok{println}\OperatorTok{(}\StringTok{"Size: "} \OperatorTok{+} inventory\OperatorTok{.}\CodeFuncTok{size}\OperatorTok{());}
\end{Highlighting}
\end{Shaded}

\subsection{\texorpdfstring{\texttt{HashSet} --- Unique
Elements}{HashSet --- Unique Elements}}\label{hashset-unique-elements}

\begin{Shaded}
\begin{Highlighting}[]
\KeywordTok{import} \CodeAccentTok{java}\OperatorTok{.}\CodeAccentTok{util}\OperatorTok{.}\CodeAccentTok{HashSet}\OperatorTok{;}
 
HashSet\OperatorTok{\textless{}}String\OperatorTok{\textgreater{}} uniqueWords \OperatorTok{=} \KeywordTok{new} HashSet\OperatorTok{\textless{}\textgreater{}();}
uniqueWords\OperatorTok{.}\CodeFuncTok{add}\OperatorTok{(}\StringTok{"apple"}\OperatorTok{);}
uniqueWords\OperatorTok{.}\CodeFuncTok{add}\OperatorTok{(}\StringTok{"banana"}\OperatorTok{);}
uniqueWords\OperatorTok{.}\CodeFuncTok{add}\OperatorTok{(}\StringTok{"apple"}\OperatorTok{);}   \CommentTok{// Duplicate {-}{-}{-} ignored}
 
System\OperatorTok{.}\CodeFuncTok{out}\OperatorTok{.}\CodeFuncTok{println}\OperatorTok{(}uniqueWords\OperatorTok{.}\CodeFuncTok{size}\OperatorTok{());}          \CommentTok{// 2}
System\OperatorTok{.}\CodeFuncTok{out}\OperatorTok{.}\CodeFuncTok{println}\OperatorTok{(}uniqueWords\OperatorTok{.}\CodeFuncTok{contains}\OperatorTok{(}\StringTok{"banana"}\OperatorTok{));} \CommentTok{// true}
\end{Highlighting}
\end{Shaded}

\subsection{\texorpdfstring{\texttt{TreeMap} --- Sorted
Keys}{TreeMap --- Sorted Keys}}\label{treemap-sorted-keys}

\begin{Shaded}
\begin{Highlighting}[]
\KeywordTok{import} \CodeAccentTok{java}\OperatorTok{.}\CodeAccentTok{util}\OperatorTok{.}\CodeAccentTok{TreeMap}\OperatorTok{;}
 
TreeMap\OperatorTok{\textless{}}String\OperatorTok{,} Integer\OperatorTok{\textgreater{}} sorted \OperatorTok{=} \KeywordTok{new} TreeMap\OperatorTok{\textless{}\textgreater{}();}
sorted\OperatorTok{.}\CodeFuncTok{put}\OperatorTok{(}\StringTok{"banana"}\OperatorTok{,} \CodeNumTok{2}\OperatorTok{);}
sorted\OperatorTok{.}\CodeFuncTok{put}\OperatorTok{(}\StringTok{"apple"}\OperatorTok{,} \CodeNumTok{1}\OperatorTok{);}
sorted\OperatorTok{.}\CodeFuncTok{put}\OperatorTok{(}\StringTok{"cherry"}\OperatorTok{,} \CodeNumTok{3}\OperatorTok{);}
 
\CommentTok{// Iterates in alphabetical key order}
\KeywordTok{for} \OperatorTok{(}Map\OperatorTok{.}\CodeFuncTok{Entry}\OperatorTok{\textless{}}String\OperatorTok{,} Integer\OperatorTok{\textgreater{}} e \OperatorTok{:} sorted\OperatorTok{.}\CodeFuncTok{entrySet}\OperatorTok{())} \OperatorTok{\{}
    System\OperatorTok{.}\CodeFuncTok{out}\OperatorTok{.}\CodeFuncTok{println}\OperatorTok{(}e\OperatorTok{.}\CodeFuncTok{getKey}\OperatorTok{()} \OperatorTok{+} \StringTok{": "} \OperatorTok{+} e\OperatorTok{.}\CodeFuncTok{getValue}\OperatorTok{());}
\OperatorTok{\}}
\end{Highlighting}
\end{Shaded}

\subsection{Sorting with Collections}\label{sorting-with-collections}

\begin{Shaded}
\begin{Highlighting}[]
\KeywordTok{import} \CodeAccentTok{java}\OperatorTok{.}\CodeAccentTok{util}\OperatorTok{.*;}
 
List\OperatorTok{\textless{}}Integer\OperatorTok{\textgreater{}} numbers \OperatorTok{=} \KeywordTok{new} ArrayList\OperatorTok{\textless{}\textgreater{}(}Arrays\OperatorTok{.}\CodeFuncTok{asList}\OperatorTok{(}\CodeNumTok{5}\OperatorTok{,} \CodeNumTok{2}\OperatorTok{,} \CodeNumTok{9}\OperatorTok{,} \CodeNumTok{1}\OperatorTok{,} \CodeNumTok{7}\OperatorTok{));}
Collections\OperatorTok{.}\CodeFuncTok{sort}\OperatorTok{(}numbers\OperatorTok{);}
System\OperatorTok{.}\CodeFuncTok{out}\OperatorTok{.}\CodeFuncTok{println}\OperatorTok{(}numbers\OperatorTok{);}   \CommentTok{// [1, 2, 5, 7, 9]}
 
Collections\OperatorTok{.}\CodeFuncTok{sort}\OperatorTok{(}numbers\OperatorTok{,} Collections\OperatorTok{.}\CodeFuncTok{reverseOrder}\OperatorTok{());}
System\OperatorTok{.}\CodeFuncTok{out}\OperatorTok{.}\CodeFuncTok{println}\OperatorTok{(}numbers\OperatorTok{);}   \CommentTok{// [9, 7, 5, 2, 1]}
 
\CodeTypeTok{int} max \OperatorTok{=} Collections\OperatorTok{.}\CodeFuncTok{max}\OperatorTok{(}numbers\OperatorTok{);}
\CodeTypeTok{int} min \OperatorTok{=} Collections\OperatorTok{.}\CodeFuncTok{min}\OperatorTok{(}numbers\OperatorTok{);}
\end{Highlighting}
\end{Shaded}

\begin{center}\rule{0.5\linewidth}{0.5pt}\end{center}

\subsection{Lesson 9.2: Word Frequency
Analysis}\label{lesson-9.2-word-frequency-analysis}

\begin{Shaded}
\begin{Highlighting}[]
\KeywordTok{import} \CodeAccentTok{java}\OperatorTok{.}\CodeAccentTok{util}\OperatorTok{.*;}
 
\KeywordTok{public} \KeywordTok{class} WordFrequency \OperatorTok{\{}
    \KeywordTok{public} \CodeTypeTok{static} Map\OperatorTok{\textless{}}String\OperatorTok{,} Integer\OperatorTok{\textgreater{}} \CodeFuncTok{countWords}\OperatorTok{(}String text\OperatorTok{)} \OperatorTok{\{}
        Map\OperatorTok{\textless{}}String\OperatorTok{,} Integer\OperatorTok{\textgreater{}} frequency \OperatorTok{=} \KeywordTok{new} HashMap\OperatorTok{\textless{}\textgreater{}();}
        String\OperatorTok{[]} words \OperatorTok{=} text\OperatorTok{.}\CodeFuncTok{toLowerCase}\OperatorTok{()}
                             \OperatorTok{.}\CodeFuncTok{replaceAll}\OperatorTok{(}\StringTok{"[\^{}a{-}z}\StringTok{\textbackslash{}\textbackslash{}}\StringTok{s]"}\OperatorTok{,} \StringTok{""}\OperatorTok{)}
                             \OperatorTok{.}\CodeFuncTok{trim}\OperatorTok{()}
                             \OperatorTok{.}\CodeFuncTok{split}\OperatorTok{(}\StringTok{"}\StringTok{\textbackslash{}\textbackslash{}}\StringTok{s+"}\OperatorTok{);}
 
        \KeywordTok{for} \OperatorTok{(}String word \OperatorTok{:} words\OperatorTok{)} \OperatorTok{\{}
            frequency\OperatorTok{.}\CodeFuncTok{put}\OperatorTok{(}word\OperatorTok{,} frequency\OperatorTok{.}\CodeFuncTok{getOrDefault}\OperatorTok{(}word\OperatorTok{,} \CodeNumTok{0}\OperatorTok{)} \OperatorTok{+} \CodeNumTok{1}\OperatorTok{);}
        \OperatorTok{\}}
        \KeywordTok{return} frequency\OperatorTok{;}
    \OperatorTok{\}}
 
    \KeywordTok{public} \CodeTypeTok{static} \CodeTypeTok{void} \CodeFuncTok{main}\OperatorTok{(}String\OperatorTok{[]} args\OperatorTok{)} \OperatorTok{\{}
        String text \OperatorTok{=} \StringTok{"the quick brown fox jumps over the lazy dog the quick brown fox"}\OperatorTok{;}
        Map\OperatorTok{\textless{}}String\OperatorTok{,} Integer\OperatorTok{\textgreater{}} freq \OperatorTok{=} \CodeFuncTok{countWords}\OperatorTok{(}text\OperatorTok{);}
 
        \CommentTok{// Sort by frequency (most common first)}
        List\OperatorTok{\textless{}}Map\OperatorTok{.}\CodeFuncTok{Entry}\OperatorTok{\textless{}}String\OperatorTok{,} Integer\OperatorTok{\textgreater{}\textgreater{}} entries \OperatorTok{=} \KeywordTok{new} ArrayList\OperatorTok{\textless{}\textgreater{}(}freq\OperatorTok{.}\CodeFuncTok{entrySet}\OperatorTok{());}
        entries\OperatorTok{.}\CodeFuncTok{sort}\OperatorTok{((}a\OperatorTok{,} b\OperatorTok{)} \OperatorTok{{-}\textgreater{}} b\OperatorTok{.}\CodeFuncTok{getValue}\OperatorTok{()} \OperatorTok{{-}} a\OperatorTok{.}\CodeFuncTok{getValue}\OperatorTok{());}
 
        System\OperatorTok{.}\CodeFuncTok{out}\OperatorTok{.}\CodeFuncTok{println}\OperatorTok{(}\StringTok{"Word Frequencies:"}\OperatorTok{);}
        \KeywordTok{for} \OperatorTok{(}Map\OperatorTok{.}\CodeFuncTok{Entry}\OperatorTok{\textless{}}String\OperatorTok{,} Integer\OperatorTok{\textgreater{}} entry \OperatorTok{:} entries\OperatorTok{)} \OperatorTok{\{}
            System\OperatorTok{.}\CodeFuncTok{out}\OperatorTok{.}\CodeFuncTok{printf}\OperatorTok{(}\StringTok{"\textquotesingle{}}\StringTok{\%s}\StringTok{\textquotesingle{}: }\StringTok{\%d\%n}\StringTok{"}\OperatorTok{,} entry\OperatorTok{.}\CodeFuncTok{getKey}\OperatorTok{(),} entry\OperatorTok{.}\CodeFuncTok{getValue}\OperatorTok{());}
        \OperatorTok{\}}
    \OperatorTok{\}}
\OperatorTok{\}}
\end{Highlighting}
\end{Shaded}

\begin{center}\rule{0.5\linewidth}{0.5pt}\end{center}

\subsection{Lesson 9.3: Working with CSV
Files}\label{lesson-9.3-working-with-csv-files}

\begin{Shaded}
\begin{Highlighting}[]
\KeywordTok{import} \CodeAccentTok{java}\OperatorTok{.}\CodeAccentTok{io}\OperatorTok{.*;}
\KeywordTok{import} \CodeAccentTok{java}\OperatorTok{.}\CodeAccentTok{util}\OperatorTok{.*;}
 
\KeywordTok{public} \KeywordTok{class} CSVReader \OperatorTok{\{}
    \KeywordTok{public} \CodeTypeTok{static} \CodeTypeTok{void} \CodeFuncTok{main}\OperatorTok{(}String\OperatorTok{[]} args\OperatorTok{)} \OperatorTok{\{}
        String filename \OperatorTok{=} \StringTok{"students.csv"}\OperatorTok{;}
 
        \KeywordTok{try} \OperatorTok{(}BufferedReader reader \OperatorTok{=} \KeywordTok{new} BufferedReader\OperatorTok{(}\KeywordTok{new} FileReader\OperatorTok{(}filename\OperatorTok{)))} \OperatorTok{\{}
            String headerLine \OperatorTok{=} reader\OperatorTok{.}\CodeFuncTok{readLine}\OperatorTok{();}
            String\OperatorTok{[]} headers \OperatorTok{=} headerLine\OperatorTok{.}\CodeFuncTok{split}\OperatorTok{(}\StringTok{","}\OperatorTok{);}
            System\OperatorTok{.}\CodeFuncTok{out}\OperatorTok{.}\CodeFuncTok{println}\OperatorTok{(}\StringTok{"Columns: "} \OperatorTok{+} Arrays\OperatorTok{.}\CodeFuncTok{toString}\OperatorTok{(}headers\OperatorTok{));}
            System\OperatorTok{.}\CodeFuncTok{out}\OperatorTok{.}\CodeFuncTok{println}\OperatorTok{(}\StringTok{"{-}{-}{-}"}\OperatorTok{);}
 
            String line\OperatorTok{;}
            \KeywordTok{while} \OperatorTok{((}line \OperatorTok{=} reader\OperatorTok{.}\CodeFuncTok{readLine}\OperatorTok{())} \OperatorTok{!=} \KeywordTok{null}\OperatorTok{)} \OperatorTok{\{}
                String\OperatorTok{[]} data \OperatorTok{=} line\OperatorTok{.}\CodeFuncTok{split}\OperatorTok{(}\StringTok{","}\OperatorTok{);}
                \KeywordTok{if} \OperatorTok{(}data\OperatorTok{.}\CodeFuncTok{length} \OperatorTok{==} headers\OperatorTok{.}\CodeFuncTok{length}\OperatorTok{)} \OperatorTok{\{}
                    String name \OperatorTok{=} data\OperatorTok{[}\CodeNumTok{0}\OperatorTok{].}\CodeFuncTok{trim}\OperatorTok{();}
                    \CodeTypeTok{int} age \OperatorTok{=} Integer\OperatorTok{.}\CodeFuncTok{parseInt}\OperatorTok{(}data\OperatorTok{[}\CodeNumTok{1}\OperatorTok{].}\CodeFuncTok{trim}\OperatorTok{());}
                    String grade \OperatorTok{=} data\OperatorTok{[}\CodeNumTok{2}\OperatorTok{].}\CodeFuncTok{trim}\OperatorTok{();}
                    System\OperatorTok{.}\CodeFuncTok{out}\OperatorTok{.}\CodeFuncTok{printf}\OperatorTok{(}\StringTok{"Name: }\StringTok{\%{-}15s}\StringTok{ Age: }\StringTok{\%3d}\StringTok{  Grade: }\StringTok{\%s\%n}\StringTok{"}\OperatorTok{,}
                                      name\OperatorTok{,} age\OperatorTok{,} grade\OperatorTok{);}
                \OperatorTok{\}}
            \OperatorTok{\}}
        \OperatorTok{\}} \KeywordTok{catch} \OperatorTok{(}FileNotFoundException e\OperatorTok{)} \OperatorTok{\{}
            System\OperatorTok{.}\CodeFuncTok{out}\OperatorTok{.}\CodeFuncTok{println}\OperatorTok{(}\StringTok{"Error: File not found {-} "} \OperatorTok{+} filename\OperatorTok{);}
        \OperatorTok{\}} \KeywordTok{catch} \OperatorTok{(}IOException e\OperatorTok{)} \OperatorTok{\{}
            System\OperatorTok{.}\CodeFuncTok{out}\OperatorTok{.}\CodeFuncTok{println}\OperatorTok{(}\StringTok{"Error reading file: "} \OperatorTok{+} e\OperatorTok{.}\CodeFuncTok{getMessage}\OperatorTok{());}
        \OperatorTok{\}} \KeywordTok{catch} \OperatorTok{(}NumberFormatException e\OperatorTok{)} \OperatorTok{\{}
            System\OperatorTok{.}\CodeFuncTok{out}\OperatorTok{.}\CodeFuncTok{println}\OperatorTok{(}\StringTok{"Error: Invalid number format in CSV."}\OperatorTok{);}
        \OperatorTok{\}}
    \OperatorTok{\}}
\OperatorTok{\}}
\end{Highlighting}
\end{Shaded}

\begin{center}\rule{0.5\linewidth}{0.5pt}\end{center}

\subsection{Lesson 9.4: Unit Testing with
JUnit}\label{lesson-9.4-unit-testing-with-junit}

\subsection{Adding JUnit to Your
Project}\label{adding-junit-to-your-project}

If using VS Code with Maven, add to \texttt{pom.xml}:

\begin{Shaded}
\begin{Highlighting}[]
\textless{}\KeywordTok{dependency}\textgreater{}
    \textless{}\KeywordTok{groupId}\textgreater{}org.junit.jupiter\textless{}/\KeywordTok{groupId}\textgreater{}
    \textless{}\KeywordTok{artifactId}\textgreater{}junit{-}jupiter\textless{}/\KeywordTok{artifactId}\textgreater{}
    \textless{}\KeywordTok{version}\textgreater{}5.10.0\textless{}/\KeywordTok{version}\textgreater{}
    \textless{}\KeywordTok{scope}\textgreater{}test\textless{}/\KeywordTok{scope}\textgreater{}
\textless{}/\KeywordTok{dependency}\textgreater{}
\end{Highlighting}
\end{Shaded}

\subsection{Writing Tests}\label{writing-tests-1}

\begin{Shaded}
\begin{Highlighting}[]
\CommentTok{// MyMath.java}
\KeywordTok{public} \KeywordTok{class} MyMath \OperatorTok{\{}
    \KeywordTok{public} \CodeTypeTok{static} \CodeTypeTok{int} \CodeFuncTok{add}\OperatorTok{(}\CodeTypeTok{int} a\OperatorTok{,} \CodeTypeTok{int} b\OperatorTok{)} \OperatorTok{\{} \KeywordTok{return} a \OperatorTok{+} b\OperatorTok{;} \OperatorTok{\}}
    \KeywordTok{public} \CodeTypeTok{static} \CodeTypeTok{double} \CodeFuncTok{circleArea}\OperatorTok{(}\CodeTypeTok{double} radius\OperatorTok{)} \OperatorTok{\{} \KeywordTok{return} Math\OperatorTok{.}\CodeFuncTok{PI} \OperatorTok{*} radius \OperatorTok{*} radius\OperatorTok{;} \OperatorTok{\}}
    \KeywordTok{public} \CodeTypeTok{static} String \CodeFuncTok{getGrade}\OperatorTok{(}\CodeTypeTok{int} score\OperatorTok{)} \OperatorTok{\{}
        \KeywordTok{if} \OperatorTok{(}score \OperatorTok{\textgreater{}=} \CodeNumTok{90}\OperatorTok{)} \KeywordTok{return} \StringTok{"A"}\OperatorTok{;}
        \KeywordTok{if} \OperatorTok{(}score \OperatorTok{\textgreater{}=} \CodeNumTok{80}\OperatorTok{)} \KeywordTok{return} \StringTok{"B"}\OperatorTok{;}
        \KeywordTok{if} \OperatorTok{(}score \OperatorTok{\textgreater{}=} \CodeNumTok{70}\OperatorTok{)} \KeywordTok{return} \StringTok{"C"}\OperatorTok{;}
        \KeywordTok{return} \StringTok{"F"}\OperatorTok{;}
    \OperatorTok{\}}
\OperatorTok{\}}
\end{Highlighting}
\end{Shaded}

\begin{Shaded}
\begin{Highlighting}[]
\CommentTok{// TestMyMath.java}
\KeywordTok{import} \CodeAccentTok{org}\OperatorTok{.}\CodeAccentTok{junit}\OperatorTok{.}\CodeAccentTok{jupiter}\OperatorTok{.}\CodeAccentTok{api}\OperatorTok{.*;}
\KeywordTok{import} \KeywordTok{static}\CodeAccentTok{ org}\OperatorTok{.}\CodeAccentTok{junit}\OperatorTok{.}\CodeAccentTok{jupiter}\OperatorTok{.}\CodeAccentTok{api}\OperatorTok{.}\CodeAccentTok{Assertions}\OperatorTok{.*;}
 
\KeywordTok{public} \KeywordTok{class} TestMyMath \OperatorTok{\{}
 
    \CodeTypeTok{@Test}
    \KeywordTok{public} \CodeTypeTok{void} \CodeFuncTok{testAddPositiveNumbers}\OperatorTok{()} \OperatorTok{\{}
        \CodeFuncTok{assertEquals}\OperatorTok{(}\CodeNumTok{8}\OperatorTok{,} MyMath\OperatorTok{.}\CodeFuncTok{add}\OperatorTok{(}\CodeNumTok{5}\OperatorTok{,} \CodeNumTok{3}\OperatorTok{));}
    \OperatorTok{\}}
 
    \CodeTypeTok{@Test}
    \KeywordTok{public} \CodeTypeTok{void} \CodeFuncTok{testAddNegativeNumbers}\OperatorTok{()} \OperatorTok{\{}
        \CodeFuncTok{assertEquals}\OperatorTok{({-}}\CodeNumTok{8}\OperatorTok{,} MyMath\OperatorTok{.}\CodeFuncTok{add}\OperatorTok{({-}}\CodeNumTok{5}\OperatorTok{,} \OperatorTok{{-}}\CodeNumTok{3}\OperatorTok{));}
    \OperatorTok{\}}
 
    \CodeTypeTok{@Test}
    \KeywordTok{public} \CodeTypeTok{void} \CodeFuncTok{testCircleArea}\OperatorTok{()} \OperatorTok{\{}
        \CodeTypeTok{double} expected \OperatorTok{=} Math\OperatorTok{.}\CodeFuncTok{PI} \OperatorTok{*} \CodeNumTok{25}\OperatorTok{;}
        \CodeFuncTok{assertEquals}\OperatorTok{(}expected\OperatorTok{,} MyMath\OperatorTok{.}\CodeFuncTok{circleArea}\OperatorTok{(}\CodeNumTok{5.0}\OperatorTok{),} \CodeNumTok{0.0001}\OperatorTok{);}
    \OperatorTok{\}}
 
    \CodeTypeTok{@Test}
    \KeywordTok{public} \CodeTypeTok{void} \CodeFuncTok{testGradeA}\OperatorTok{()} \OperatorTok{\{}
        \CodeFuncTok{assertEquals}\OperatorTok{(}\StringTok{"A"}\OperatorTok{,} MyMath\OperatorTok{.}\CodeFuncTok{getGrade}\OperatorTok{(}\CodeNumTok{95}\OperatorTok{));}
    \OperatorTok{\}}
 
    \CodeTypeTok{@Test}
    \KeywordTok{public} \CodeTypeTok{void} \CodeFuncTok{testGradeF}\OperatorTok{()} \OperatorTok{\{}
        \CodeFuncTok{assertEquals}\OperatorTok{(}\StringTok{"F"}\OperatorTok{,} MyMath\OperatorTok{.}\CodeFuncTok{getGrade}\OperatorTok{(}\CodeNumTok{55}\OperatorTok{));}
    \OperatorTok{\}}
\OperatorTok{\}}
\end{Highlighting}
\end{Shaded}

\subsection{Common JUnit Assertions}\label{common-junit-assertions}

\begin{fullwidth}{\def\LTcaptype{none} % do not increment counter
\begin{longtable}[]{@{}
  >{\raggedright\arraybackslash}p{(\linewidth - 2\tabcolsep) * \real{0.6269}}
  >{\raggedright\arraybackslash}p{(\linewidth - 2\tabcolsep) * \real{0.3731}}@{}}
\toprule\noalign{}
\begin{minipage}[b]{\linewidth}\raggedright
Method
\end{minipage} & \begin{minipage}[b]{\linewidth}\raggedright
Description
\end{minipage} \\
\midrule\noalign{}
\endhead
\bottomrule\noalign{}
\endlastfoot
\texttt{assertEquals(expected,\ actual)} & Check equality \\
\texttt{assertTrue(condition)} & Check condition is true \\
\texttt{assertFalse(condition)} & Check condition is false \\
\texttt{assertNull(object)} & Check object is null \\
\texttt{assertNotNull(object)} & Check object is not null \\
\texttt{assertThrows(ExceptionClass,\ executable)} & Check exception is
thrown \\
\end{longtable}
}
\end{fullwidth}

\subsection{Skills Checklist}\label{skills-checklist-85}

\begin{itemize}
\tightlist
\item[$\square$]
  I can use \texttt{HashMap} for key-value storage in Java.
\item[$\square$]
  I can iterate through a map's entries, keys, and values.
\item[$\square$]
  I can write unit tests using JUnit 5.
\item[$\square$]
  I can use common assertion methods.
\end{itemize}

\begin{center}\rule{0.5\linewidth}{0.5pt}\end{center}

\section{Unit 10: Project-Based
Learning}\label{unit-10-project-based-learning-2}

\subsection{Overview}\label{overview-24}

This unit guides students through building text-based Java applications.

\begin{center}\rule{0.5\linewidth}{0.5pt}\end{center}

\subsection{Lesson 10.1: Designing a Text-Based Adventure
Game}\label{lesson-10.1-designing-a-text-based-adventure-game}

\subsection{Room and Player
Structure}\label{room-and-player-structure}

\begin{Shaded}
\begin{Highlighting}[]
\KeywordTok{import} \CodeAccentTok{java}\OperatorTok{.}\CodeAccentTok{util}\OperatorTok{.*;}
 
\KeywordTok{public} \KeywordTok{class} Room \OperatorTok{\{}
    \KeywordTok{private} String name\OperatorTok{;}
    \KeywordTok{private} String description\OperatorTok{;}
    \KeywordTok{private} Map\OperatorTok{\textless{}}String\OperatorTok{,} Room\OperatorTok{\textgreater{}} exits\OperatorTok{;}
    \KeywordTok{private} List\OperatorTok{\textless{}}String\OperatorTok{\textgreater{}} items\OperatorTok{;}
 
    \KeywordTok{public} \CodeFuncTok{Room}\OperatorTok{(}String name\OperatorTok{,} String description\OperatorTok{)} \OperatorTok{\{}
        \KeywordTok{this}\OperatorTok{.}\CodeFuncTok{name} \OperatorTok{=} name\OperatorTok{;}
        \KeywordTok{this}\OperatorTok{.}\CodeFuncTok{description} \OperatorTok{=} description\OperatorTok{;}
        \KeywordTok{this}\OperatorTok{.}\CodeFuncTok{exits} \OperatorTok{=} \KeywordTok{new} HashMap\OperatorTok{\textless{}\textgreater{}();}
        \KeywordTok{this}\OperatorTok{.}\CodeFuncTok{items} \OperatorTok{=} \KeywordTok{new} ArrayList\OperatorTok{\textless{}\textgreater{}();}
    \OperatorTok{\}}
 
    \KeywordTok{public} \CodeTypeTok{void} \CodeFuncTok{addExit}\OperatorTok{(}String direction\OperatorTok{,} Room room\OperatorTok{)} \OperatorTok{\{}
        exits\OperatorTok{.}\CodeFuncTok{put}\OperatorTok{(}direction\OperatorTok{.}\CodeFuncTok{toLowerCase}\OperatorTok{(),} room\OperatorTok{);}
    \OperatorTok{\}}
 
    \KeywordTok{public} \CodeTypeTok{void} \CodeFuncTok{addItem}\OperatorTok{(}String item\OperatorTok{)} \OperatorTok{\{}
        items\OperatorTok{.}\CodeFuncTok{add}\OperatorTok{(}item\OperatorTok{);}
    \OperatorTok{\}}
 
    \KeywordTok{public} Room \CodeFuncTok{getExit}\OperatorTok{(}String direction\OperatorTok{)} \OperatorTok{\{}
        \KeywordTok{return} exits\OperatorTok{.}\CodeFuncTok{get}\OperatorTok{(}direction\OperatorTok{.}\CodeFuncTok{toLowerCase}\OperatorTok{());}
    \OperatorTok{\}}
 
    \KeywordTok{public} String \CodeFuncTok{getDescription}\OperatorTok{()} \OperatorTok{\{}
        StringBuilder sb \OperatorTok{=} \KeywordTok{new} StringBuilder\OperatorTok{();}
        sb\OperatorTok{.}\CodeFuncTok{append}\OperatorTok{(}\StringTok{"{-}{-}{-} "}\OperatorTok{).}\CodeFuncTok{append}\OperatorTok{(}name\OperatorTok{).}\CodeFuncTok{append}\OperatorTok{(}\StringTok{" {-}{-}{-}}\StringTok{\textbackslash{}n}\StringTok{"}\OperatorTok{);}
        sb\OperatorTok{.}\CodeFuncTok{append}\OperatorTok{(}description\OperatorTok{).}\CodeFuncTok{append}\OperatorTok{(}\StringTok{"}\StringTok{\textbackslash{}n}\StringTok{"}\OperatorTok{);}
        \KeywordTok{if} \OperatorTok{(!}items\OperatorTok{.}\CodeFuncTok{isEmpty}\OperatorTok{())} \OperatorTok{\{}
            sb\OperatorTok{.}\CodeFuncTok{append}\OperatorTok{(}\StringTok{"You see: "}\OperatorTok{).}\CodeFuncTok{append}\OperatorTok{(}String\OperatorTok{.}\CodeFuncTok{join}\OperatorTok{(}\StringTok{", "}\OperatorTok{,} items\OperatorTok{)).}\CodeFuncTok{append}\OperatorTok{(}\StringTok{"}\StringTok{\textbackslash{}n}\StringTok{"}\OperatorTok{);}
        \OperatorTok{\}}
        sb\OperatorTok{.}\CodeFuncTok{append}\OperatorTok{(}\StringTok{"Exits: "}\OperatorTok{).}\CodeFuncTok{append}\OperatorTok{(}String\OperatorTok{.}\CodeFuncTok{join}\OperatorTok{(}\StringTok{", "}\OperatorTok{,} exits\OperatorTok{.}\CodeFuncTok{keySet}\OperatorTok{()));}
        \KeywordTok{return} sb\OperatorTok{.}\CodeFuncTok{toString}\OperatorTok{();}
    \OperatorTok{\}}
 
    \KeywordTok{public} String \CodeFuncTok{getName}\OperatorTok{()} \OperatorTok{\{} \KeywordTok{return} name\OperatorTok{;} \OperatorTok{\}}
    \KeywordTok{public} List\OperatorTok{\textless{}}String\OperatorTok{\textgreater{}} \CodeFuncTok{getItems}\OperatorTok{()} \OperatorTok{\{} \KeywordTok{return} items\OperatorTok{;} \OperatorTok{\}}
\OperatorTok{\}}
\end{Highlighting}
\end{Shaded}

\begin{Shaded}
\begin{Highlighting}[]
\KeywordTok{import} \CodeAccentTok{java}\OperatorTok{.}\CodeAccentTok{util}\OperatorTok{.*;}
 
\KeywordTok{public} \KeywordTok{class} Player \OperatorTok{\{}
    \KeywordTok{private} String name\OperatorTok{;}
    \KeywordTok{private} Room currentRoom\OperatorTok{;}
    \KeywordTok{private} List\OperatorTok{\textless{}}String\OperatorTok{\textgreater{}} inventory\OperatorTok{;}
    \KeywordTok{private} \CodeTypeTok{int} health\OperatorTok{;}
 
    \KeywordTok{public} \CodeFuncTok{Player}\OperatorTok{(}String name\OperatorTok{,} Room startRoom\OperatorTok{)} \OperatorTok{\{}
        \KeywordTok{this}\OperatorTok{.}\CodeFuncTok{name} \OperatorTok{=} name\OperatorTok{;}
        \KeywordTok{this}\OperatorTok{.}\CodeFuncTok{currentRoom} \OperatorTok{=} startRoom\OperatorTok{;}
        \KeywordTok{this}\OperatorTok{.}\CodeFuncTok{inventory} \OperatorTok{=} \KeywordTok{new} ArrayList\OperatorTok{\textless{}\textgreater{}();}
        \KeywordTok{this}\OperatorTok{.}\CodeFuncTok{health} \OperatorTok{=} \CodeNumTok{100}\OperatorTok{;}
    \OperatorTok{\}}
 
    \KeywordTok{public} \CodeTypeTok{boolean} \CodeFuncTok{move}\OperatorTok{(}String direction\OperatorTok{)} \OperatorTok{\{}
        Room nextRoom \OperatorTok{=} currentRoom\OperatorTok{.}\CodeFuncTok{getExit}\OperatorTok{(}direction\OperatorTok{);}
        \KeywordTok{if} \OperatorTok{(}nextRoom \OperatorTok{!=} \KeywordTok{null}\OperatorTok{)} \OperatorTok{\{}
            currentRoom \OperatorTok{=} nextRoom\OperatorTok{;}
            \KeywordTok{return} \KeywordTok{true}\OperatorTok{;}
        \OperatorTok{\}}
        \KeywordTok{return} \KeywordTok{false}\OperatorTok{;}
    \OperatorTok{\}}
 
    \KeywordTok{public} \CodeTypeTok{boolean} \CodeFuncTok{takeItem}\OperatorTok{(}String itemName\OperatorTok{)} \OperatorTok{\{}
        List\OperatorTok{\textless{}}String\OperatorTok{\textgreater{}} roomItems \OperatorTok{=} currentRoom\OperatorTok{.}\CodeFuncTok{getItems}\OperatorTok{();}
        \KeywordTok{for} \OperatorTok{(}String item \OperatorTok{:} roomItems\OperatorTok{)} \OperatorTok{\{}
            \KeywordTok{if} \OperatorTok{(}item\OperatorTok{.}\CodeFuncTok{equalsIgnoreCase}\OperatorTok{(}itemName\OperatorTok{))} \OperatorTok{\{}
                inventory\OperatorTok{.}\CodeFuncTok{add}\OperatorTok{(}item\OperatorTok{);}
                roomItems\OperatorTok{.}\CodeFuncTok{remove}\OperatorTok{(}item\OperatorTok{);}
                \KeywordTok{return} \KeywordTok{true}\OperatorTok{;}
            \OperatorTok{\}}
        \OperatorTok{\}}
        \KeywordTok{return} \KeywordTok{false}\OperatorTok{;}
    \OperatorTok{\}}
 
    \KeywordTok{public} String \CodeFuncTok{getStatus}\OperatorTok{()} \OperatorTok{\{}
        \KeywordTok{return} \StringTok{"Player: "} \OperatorTok{+} name \OperatorTok{+} \StringTok{" | Health: "} \OperatorTok{+} health \OperatorTok{+}
               \StringTok{" | Inventory: "} \OperatorTok{+} \OperatorTok{(}inventory\OperatorTok{.}\CodeFuncTok{isEmpty}\OperatorTok{()} \OperatorTok{?} \StringTok{"empty"} \OperatorTok{:} String\OperatorTok{.}\CodeFuncTok{join}\OperatorTok{(}\StringTok{", "}\OperatorTok{,} inventory\OperatorTok{));}
    \OperatorTok{\}}
 
    \KeywordTok{public} Room \CodeFuncTok{getCurrentRoom}\OperatorTok{()} \OperatorTok{\{} \KeywordTok{return} currentRoom\OperatorTok{;} \OperatorTok{\}}
    \KeywordTok{public} List\OperatorTok{\textless{}}String\OperatorTok{\textgreater{}} \CodeFuncTok{getInventory}\OperatorTok{()} \OperatorTok{\{} \KeywordTok{return} inventory\OperatorTok{;} \OperatorTok{\}}
\OperatorTok{\}}
\end{Highlighting}
\end{Shaded}

\begin{Shaded}
\begin{Highlighting}[]
\KeywordTok{import} \CodeAccentTok{java}\OperatorTok{.}\CodeAccentTok{util}\OperatorTok{.}\CodeAccentTok{Scanner}\OperatorTok{;}
 
\KeywordTok{public} \KeywordTok{class} AdventureGame \OperatorTok{\{}
    \KeywordTok{public} \CodeTypeTok{static} \CodeTypeTok{void} \CodeFuncTok{main}\OperatorTok{(}String\OperatorTok{[]} args\OperatorTok{)} \OperatorTok{\{}
        \CommentTok{// Build the world}
        Room cave \OperatorTok{=} \KeywordTok{new} \CodeFuncTok{Room}\OperatorTok{(}\StringTok{"Dark Cave"}\OperatorTok{,} \StringTok{"You are in a dimly lit cave."}\OperatorTok{);}
        Room tunnel \OperatorTok{=} \KeywordTok{new} \CodeFuncTok{Room}\OperatorTok{(}\StringTok{"Stone Tunnel"}\OperatorTok{,} \StringTok{"A narrow tunnel leads north and south."}\OperatorTok{);}
        Room clearing \OperatorTok{=} \KeywordTok{new} \CodeFuncTok{Room}\OperatorTok{(}\StringTok{"Sunny Clearing"}\OperatorTok{,} \StringTok{"A bright clearing with tall trees."}\OperatorTok{);}
 
        cave\OperatorTok{.}\CodeFuncTok{addExit}\OperatorTok{(}\StringTok{"north"}\OperatorTok{,} tunnel\OperatorTok{);}
        cave\OperatorTok{.}\CodeFuncTok{addItem}\OperatorTok{(}\StringTok{"torch"}\OperatorTok{);}
        tunnel\OperatorTok{.}\CodeFuncTok{addExit}\OperatorTok{(}\StringTok{"south"}\OperatorTok{,} cave\OperatorTok{);}
        tunnel\OperatorTok{.}\CodeFuncTok{addExit}\OperatorTok{(}\StringTok{"north"}\OperatorTok{,} clearing\OperatorTok{);}
        clearing\OperatorTok{.}\CodeFuncTok{addExit}\OperatorTok{(}\StringTok{"south"}\OperatorTok{,} tunnel\OperatorTok{);}
 
        Player player \OperatorTok{=} \KeywordTok{new} \CodeFuncTok{Player}\OperatorTok{(}\StringTok{"Hero"}\OperatorTok{,} cave\OperatorTok{);}
        Scanner scanner \OperatorTok{=} \KeywordTok{new} Scanner\OperatorTok{(}System\OperatorTok{.}\CodeFuncTok{in}\OperatorTok{);}
 
        System\OperatorTok{.}\CodeFuncTok{out}\OperatorTok{.}\CodeFuncTok{println}\OperatorTok{(}\StringTok{"=== Text Adventure ==="}\OperatorTok{);}
        System\OperatorTok{.}\CodeFuncTok{out}\OperatorTok{.}\CodeFuncTok{println}\OperatorTok{(}\StringTok{"Commands: go [direction], take [item], inventory, look, quit"}\OperatorTok{);}
        System\OperatorTok{.}\CodeFuncTok{out}\OperatorTok{.}\CodeFuncTok{println}\OperatorTok{();}
 
        \KeywordTok{while} \OperatorTok{(}\KeywordTok{true}\OperatorTok{)} \OperatorTok{\{}
            System\OperatorTok{.}\CodeFuncTok{out}\OperatorTok{.}\CodeFuncTok{println}\OperatorTok{(}player\OperatorTok{.}\CodeFuncTok{getCurrentRoom}\OperatorTok{().}\CodeFuncTok{getDescription}\OperatorTok{());}
            System\OperatorTok{.}\CodeFuncTok{out}\OperatorTok{.}\CodeFuncTok{println}\OperatorTok{(}player\OperatorTok{.}\CodeFuncTok{getStatus}\OperatorTok{());}
            System\OperatorTok{.}\CodeFuncTok{out}\OperatorTok{.}\CodeFuncTok{print}\OperatorTok{(}\StringTok{"}\StringTok{\textbackslash{}n}\StringTok{\textgreater{} "}\OperatorTok{);}
 
            String input \OperatorTok{=} scanner\OperatorTok{.}\CodeFuncTok{nextLine}\OperatorTok{().}\CodeFuncTok{trim}\OperatorTok{().}\CodeFuncTok{toLowerCase}\OperatorTok{();}
            String\OperatorTok{[]} parts \OperatorTok{=} input\OperatorTok{.}\CodeFuncTok{split}\OperatorTok{(}\StringTok{"}\StringTok{\textbackslash{}\textbackslash{}}\StringTok{s+"}\OperatorTok{,} \CodeNumTok{2}\OperatorTok{);}
            String command \OperatorTok{=} parts\OperatorTok{[}\CodeNumTok{0}\OperatorTok{];}
            String argument \OperatorTok{=} parts\OperatorTok{.}\CodeFuncTok{length} \OperatorTok{\textgreater{}} \CodeNumTok{1} \OperatorTok{?} parts\OperatorTok{[}\CodeNumTok{1}\OperatorTok{]} \OperatorTok{:} \StringTok{""}\OperatorTok{;}
 
            \KeywordTok{switch} \OperatorTok{(}command\OperatorTok{)} \OperatorTok{\{}
                \KeywordTok{case} \StringTok{"go"}\OperatorTok{:}
                    \KeywordTok{if} \OperatorTok{(}argument\OperatorTok{.}\CodeFuncTok{isEmpty}\OperatorTok{())} \OperatorTok{\{}
                        System\OperatorTok{.}\CodeFuncTok{out}\OperatorTok{.}\CodeFuncTok{println}\OperatorTok{(}\StringTok{"Go where?"}\OperatorTok{);}
                    \OperatorTok{\}} \KeywordTok{else} \KeywordTok{if} \OperatorTok{(!}player\OperatorTok{.}\CodeFuncTok{move}\OperatorTok{(}argument\OperatorTok{))} \OperatorTok{\{}
                        System\OperatorTok{.}\CodeFuncTok{out}\OperatorTok{.}\CodeFuncTok{println}\OperatorTok{(}\StringTok{"You can\textquotesingle{}t go that way."}\OperatorTok{);}
                    \OperatorTok{\}}
                    \KeywordTok{break}\OperatorTok{;}
                \KeywordTok{case} \StringTok{"take"}\OperatorTok{:}
                    \KeywordTok{if} \OperatorTok{(}player\OperatorTok{.}\CodeFuncTok{takeItem}\OperatorTok{(}argument\OperatorTok{))} \OperatorTok{\{}
                        System\OperatorTok{.}\CodeFuncTok{out}\OperatorTok{.}\CodeFuncTok{println}\OperatorTok{(}\StringTok{"You picked up the "} \OperatorTok{+} argument \OperatorTok{+} \StringTok{"."}\OperatorTok{);}
                    \OperatorTok{\}} \KeywordTok{else} \OperatorTok{\{}
                        System\OperatorTok{.}\CodeFuncTok{out}\OperatorTok{.}\CodeFuncTok{println}\OperatorTok{(}\StringTok{"There\textquotesingle{}s no "} \OperatorTok{+} argument \OperatorTok{+} \StringTok{" here."}\OperatorTok{);}
                    \OperatorTok{\}}
                    \KeywordTok{break}\OperatorTok{;}
                \KeywordTok{case} \StringTok{"inventory"}\OperatorTok{:}
                    System\OperatorTok{.}\CodeFuncTok{out}\OperatorTok{.}\CodeFuncTok{println}\OperatorTok{(}\StringTok{"Inventory: "} \OperatorTok{+} player\OperatorTok{.}\CodeFuncTok{getInventory}\OperatorTok{());}
                    \KeywordTok{break}\OperatorTok{;}
                \KeywordTok{case} \StringTok{"look"}\OperatorTok{:}
                    \CommentTok{// Room description already shown at top of loop}
                    \KeywordTok{break}\OperatorTok{;}
                \KeywordTok{case} \StringTok{"quit"}\OperatorTok{:}
                    System\OperatorTok{.}\CodeFuncTok{out}\OperatorTok{.}\CodeFuncTok{println}\OperatorTok{(}\StringTok{"Thanks for playing!"}\OperatorTok{);}
                    scanner\OperatorTok{.}\CodeFuncTok{close}\OperatorTok{();}
                    \KeywordTok{return}\OperatorTok{;}
                \KeywordTok{default}\OperatorTok{:}
                    System\OperatorTok{.}\CodeFuncTok{out}\OperatorTok{.}\CodeFuncTok{println}\OperatorTok{(}\StringTok{"Unknown command: "} \OperatorTok{+} command\OperatorTok{);}
            \OperatorTok{\}}
            System\OperatorTok{.}\CodeFuncTok{out}\OperatorTok{.}\CodeFuncTok{println}\OperatorTok{();}
        \OperatorTok{\}}
    \OperatorTok{\}}
\OperatorTok{\}}
\end{Highlighting}
\end{Shaded}

\begin{center}\rule{0.5\linewidth}{0.5pt}\end{center}

\subsection{Lesson 10.2: Command-Line To-Do
Manager}\label{lesson-10.2-command-line-to-do-manager-1}

\begin{Shaded}
\begin{Highlighting}[]
\KeywordTok{import} \CodeAccentTok{java}\OperatorTok{.}\CodeAccentTok{util}\OperatorTok{.*;}
\KeywordTok{import} \CodeAccentTok{java}\OperatorTok{.}\CodeAccentTok{io}\OperatorTok{.*;}
\KeywordTok{import} \CodeAccentTok{com}\OperatorTok{.}\CodeAccentTok{google}\OperatorTok{.}\CodeAccentTok{gson}\OperatorTok{.*;}   \CommentTok{// Add Gson library for JSON support}
 
\KeywordTok{public} \KeywordTok{class} TodoManager \OperatorTok{\{}
    \KeywordTok{private} List\OperatorTok{\textless{}}Map\OperatorTok{\textless{}}String\OperatorTok{,} Object\OperatorTok{\textgreater{}\textgreater{}} todos \OperatorTok{=} \KeywordTok{new} ArrayList\OperatorTok{\textless{}\textgreater{}();}
    \KeywordTok{private} \CodeTypeTok{final} String FILE\_PATH \OperatorTok{=} \StringTok{"todos.json"}\OperatorTok{;}
 
    \KeywordTok{public} \CodeTypeTok{void} \CodeFuncTok{addTask}\OperatorTok{(}String description\OperatorTok{)} \OperatorTok{\{}
        Map\OperatorTok{\textless{}}String\OperatorTok{,} Object\OperatorTok{\textgreater{}} task \OperatorTok{=} \KeywordTok{new} HashMap\OperatorTok{\textless{}\textgreater{}();}
        task\OperatorTok{.}\CodeFuncTok{put}\OperatorTok{(}\StringTok{"description"}\OperatorTok{,} description\OperatorTok{);}
        task\OperatorTok{.}\CodeFuncTok{put}\OperatorTok{(}\StringTok{"completed"}\OperatorTok{,} \KeywordTok{false}\OperatorTok{);}
        todos\OperatorTok{.}\CodeFuncTok{add}\OperatorTok{(}task\OperatorTok{);}
        System\OperatorTok{.}\CodeFuncTok{out}\OperatorTok{.}\CodeFuncTok{println}\OperatorTok{(}\StringTok{"Task added: "} \OperatorTok{+} description\OperatorTok{);}
    \OperatorTok{\}}
 
    \KeywordTok{public} \CodeTypeTok{void} \CodeFuncTok{viewTasks}\OperatorTok{()} \OperatorTok{\{}
        \KeywordTok{if} \OperatorTok{(}todos\OperatorTok{.}\CodeFuncTok{isEmpty}\OperatorTok{())} \OperatorTok{\{}
            System\OperatorTok{.}\CodeFuncTok{out}\OperatorTok{.}\CodeFuncTok{println}\OperatorTok{(}\StringTok{"No tasks in your list."}\OperatorTok{);}
            \KeywordTok{return}\OperatorTok{;}
        \OperatorTok{\}}
        System\OperatorTok{.}\CodeFuncTok{out}\OperatorTok{.}\CodeFuncTok{println}\OperatorTok{(}\StringTok{"{-}{-}{-} Your To{-}Do List {-}{-}{-}"}\OperatorTok{);}
        \KeywordTok{for} \OperatorTok{(}\CodeTypeTok{int} i \OperatorTok{=} \CodeNumTok{0}\OperatorTok{;} i \OperatorTok{\textless{}} todos\OperatorTok{.}\CodeFuncTok{size}\OperatorTok{();} i\OperatorTok{++)} \OperatorTok{\{}
            Map\OperatorTok{\textless{}}String\OperatorTok{,} Object\OperatorTok{\textgreater{}} task \OperatorTok{=} todos\OperatorTok{.}\CodeFuncTok{get}\OperatorTok{(}i\OperatorTok{);}
            String status \OperatorTok{=} \OperatorTok{(}Boolean\OperatorTok{)} task\OperatorTok{.}\CodeFuncTok{get}\OperatorTok{(}\StringTok{"completed"}\OperatorTok{)} \OperatorTok{?} \StringTok{"[X]"} \OperatorTok{:} \StringTok{"[ ]"}\OperatorTok{;}
            System\OperatorTok{.}\CodeFuncTok{out}\OperatorTok{.}\CodeFuncTok{printf}\OperatorTok{(}\StringTok{"}\StringTok{\%d}\StringTok{. }\StringTok{\%s}\StringTok{ }\StringTok{\%s\%n}\StringTok{"}\OperatorTok{,} i \OperatorTok{+} \CodeNumTok{1}\OperatorTok{,} status\OperatorTok{,} task\OperatorTok{.}\CodeFuncTok{get}\OperatorTok{(}\StringTok{"description"}\OperatorTok{));}
        \OperatorTok{\}}
        System\OperatorTok{.}\CodeFuncTok{out}\OperatorTok{.}\CodeFuncTok{println}\OperatorTok{(}\StringTok{"{-}{-}{-}{-}{-}{-}{-}{-}{-}{-}{-}{-}{-}{-}{-}{-}{-}{-}{-}{-}{-}{-}{-}"}\OperatorTok{);}
    \OperatorTok{\}}
 
    \KeywordTok{public} \CodeTypeTok{void} \CodeFuncTok{markComplete}\OperatorTok{(}\CodeTypeTok{int} taskNumber\OperatorTok{)} \OperatorTok{\{}
        \CodeTypeTok{int} index \OperatorTok{=} taskNumber \OperatorTok{{-}} \CodeNumTok{1}\OperatorTok{;}
        \KeywordTok{if} \OperatorTok{(}index \OperatorTok{\textgreater{}=} \CodeNumTok{0} \OperatorTok{\&\&} index \OperatorTok{\textless{}} todos\OperatorTok{.}\CodeFuncTok{size}\OperatorTok{())} \OperatorTok{\{}
            todos\OperatorTok{.}\CodeFuncTok{get}\OperatorTok{(}index\OperatorTok{).}\CodeFuncTok{put}\OperatorTok{(}\StringTok{"completed"}\OperatorTok{,} \KeywordTok{true}\OperatorTok{);}
            System\OperatorTok{.}\CodeFuncTok{out}\OperatorTok{.}\CodeFuncTok{println}\OperatorTok{(}\StringTok{"Task marked complete: "} \OperatorTok{+} todos\OperatorTok{.}\CodeFuncTok{get}\OperatorTok{(}index\OperatorTok{).}\CodeFuncTok{get}\OperatorTok{(}\StringTok{"description"}\OperatorTok{));}
        \OperatorTok{\}} \KeywordTok{else} \OperatorTok{\{}
            System\OperatorTok{.}\CodeFuncTok{out}\OperatorTok{.}\CodeFuncTok{println}\OperatorTok{(}\StringTok{"Invalid task number."}\OperatorTok{);}
        \OperatorTok{\}}
    \OperatorTok{\}}
 
    \KeywordTok{public} \CodeTypeTok{static} \CodeTypeTok{void} \CodeFuncTok{main}\OperatorTok{(}String\OperatorTok{[]} args\OperatorTok{)} \OperatorTok{\{}
        TodoManager manager \OperatorTok{=} \KeywordTok{new} \CodeFuncTok{TodoManager}\OperatorTok{();}
        Scanner scanner \OperatorTok{=} \KeywordTok{new} Scanner\OperatorTok{(}System\OperatorTok{.}\CodeFuncTok{in}\OperatorTok{);}
 
        System\OperatorTok{.}\CodeFuncTok{out}\OperatorTok{.}\CodeFuncTok{println}\OperatorTok{(}\StringTok{"=== To{-}Do Manager ==="}\OperatorTok{);}
 
        \KeywordTok{while} \OperatorTok{(}\KeywordTok{true}\OperatorTok{)} \OperatorTok{\{}
            System\OperatorTok{.}\CodeFuncTok{out}\OperatorTok{.}\CodeFuncTok{println}\OperatorTok{(}\StringTok{"}\StringTok{\textbackslash{}n}\StringTok{1. Add Task  2. View Tasks  3. Mark Complete  4. Quit"}\OperatorTok{);}
            System\OperatorTok{.}\CodeFuncTok{out}\OperatorTok{.}\CodeFuncTok{print}\OperatorTok{(}\StringTok{"Choice: "}\OperatorTok{);}
            String choice \OperatorTok{=} scanner\OperatorTok{.}\CodeFuncTok{nextLine}\OperatorTok{().}\CodeFuncTok{trim}\OperatorTok{();}
 
            \KeywordTok{switch} \OperatorTok{(}choice\OperatorTok{)} \OperatorTok{\{}
                \KeywordTok{case} \StringTok{"1"}\OperatorTok{:}
                    System\OperatorTok{.}\CodeFuncTok{out}\OperatorTok{.}\CodeFuncTok{print}\OperatorTok{(}\StringTok{"Task description: "}\OperatorTok{);}
                    manager\OperatorTok{.}\CodeFuncTok{addTask}\OperatorTok{(}scanner\OperatorTok{.}\CodeFuncTok{nextLine}\OperatorTok{());}
                    \KeywordTok{break}\OperatorTok{;}
                \KeywordTok{case} \StringTok{"2"}\OperatorTok{:}
                    manager\OperatorTok{.}\CodeFuncTok{viewTasks}\OperatorTok{();}
                    \KeywordTok{break}\OperatorTok{;}
                \KeywordTok{case} \StringTok{"3"}\OperatorTok{:}
                    manager\OperatorTok{.}\CodeFuncTok{viewTasks}\OperatorTok{();}
                    System\OperatorTok{.}\CodeFuncTok{out}\OperatorTok{.}\CodeFuncTok{print}\OperatorTok{(}\StringTok{"Task number: "}\OperatorTok{);}
                    \KeywordTok{try} \OperatorTok{\{}
                        manager\OperatorTok{.}\CodeFuncTok{markComplete}\OperatorTok{(}Integer\OperatorTok{.}\CodeFuncTok{parseInt}\OperatorTok{(}scanner\OperatorTok{.}\CodeFuncTok{nextLine}\OperatorTok{()));}
                    \OperatorTok{\}} \KeywordTok{catch} \OperatorTok{(}NumberFormatException e\OperatorTok{)} \OperatorTok{\{}
                        System\OperatorTok{.}\CodeFuncTok{out}\OperatorTok{.}\CodeFuncTok{println}\OperatorTok{(}\StringTok{"Invalid input."}\OperatorTok{);}
                    \OperatorTok{\}}
                    \KeywordTok{break}\OperatorTok{;}
                \KeywordTok{case} \StringTok{"4"}\OperatorTok{:}
                    System\OperatorTok{.}\CodeFuncTok{out}\OperatorTok{.}\CodeFuncTok{println}\OperatorTok{(}\StringTok{"Goodbye!"}\OperatorTok{);}
                    scanner\OperatorTok{.}\CodeFuncTok{close}\OperatorTok{();}
                    \KeywordTok{return}\OperatorTok{;}
                \KeywordTok{default}\OperatorTok{:}
                    System\OperatorTok{.}\CodeFuncTok{out}\OperatorTok{.}\CodeFuncTok{println}\OperatorTok{(}\StringTok{"Invalid choice."}\OperatorTok{);}
            \OperatorTok{\}}
        \OperatorTok{\}}
    \OperatorTok{\}}
\OperatorTok{\}}
\end{Highlighting}
\end{Shaded}

\begin{center}\rule{0.5\linewidth}{0.5pt}\end{center}

\section{Unit 11: Capstone Project --- Personal Data Manager in
Java}\label{unit-11-capstone-project-personal-data-manager-in-java}

\subsection{Overview}\label{overview-25}

Students build a fully functional command-line personal data manager in
Java.

\subsection{Project Options}\label{project-options-1}

\begin{itemize}
\tightlist
\item
  \textbf{Contact Manager:} name, phone, email, notes
\item
  \textbf{Book Tracker:} title, author, genre, read status
\item
  \textbf{Inventory System:} item name, quantity, location, notes
\item
  \textbf{To-Do App:} description, due date, priority, status
\end{itemize}

\subsection{Required Features}\label{required-features-1}

\begin{enumerate}
\def\labelenumi{\arabic{enumi}.}
\tightlist
\item
  \textbf{Add item} --- prompt for all attributes and save to list.
\item
  \textbf{View all items} --- display all stored items clearly.
\item
  \textbf{Edit item} --- select by number/name and modify.
\item
  \textbf{Delete item} --- remove selected item.
\item
  \textbf{Search/filter} --- find items by keyword.
\item
  \textbf{Persistence} --- save/load to a JSON or CSV file.
\end{enumerate}

\subsection{Architecture
Recommendations}\label{architecture-recommendations}

\begin{verbatim}
src/
|--- model/
|   +--- Contact.java        (data model class)
|--- storage/
|   +--- DataManager.java    (file I/O)
|--- ui/
|   +--- ConsoleUI.java      (user interface)
+--- Main.java               (entry point)
\end{verbatim}

\subsection{Assessment Criteria}\label{assessment-criteria-2}

\subsection{Functionality (60\%)}\label{functionality-60-2}

\begin{itemize}
\tightlist
\item[$\square$]
  Add, view, edit, and delete items work correctly.
\item[$\square$]
  Data persists between program runs.
\item[$\square$]
  Program handles invalid input without crashing.
\end{itemize}

\subsection{Code Quality (20\%)}\label{code-quality-20-1}

\begin{itemize}
\tightlist
\item[$\square$]
  Methods are small and have a single responsibility.
\item[$\square$]
  Classes follow OOP principles (encapsulation).
\item[$\square$]
  Meaningful names for all variables, methods, and classes.
\item[$\square$]
  Appropriate comments and Javadoc.
\end{itemize}

\subsection{Accessibility \& UX (15\%)}\label{accessibility-ux-15-1}

\begin{itemize}
\tightlist
\item[$\square$]
  Clear, descriptive prompts for all input.
\item[$\square$]
  Explicit confirmation messages after every action.
\item[$\square$]
  Consistent output format.
\item[$\square$]
  Informative error messages.
\end{itemize}

\subsection{Documentation (5\%)}\label{documentation-5}

\begin{itemize}
\tightlist
\item[$\square$]
  \texttt{README.md} with setup and usage instructions.
\item[$\square$]
  Reflection on challenges and solutions.
\end{itemize}

\end{document}
